% !TEX root = cubic-moduli-descent.tex
\documentclass[11pt,a4paper]{ltjsarticle}
\usepackage{amsmath,amssymb,amsthm,mathtools,bm,mathrsfs}
\usepackage{luatexja-fontspec}
\usepackage{enumitem,booktabs,longtable,array,needspace}
\usepackage{xcolor}
\usepackage{etoolbox,xr-hyper}
\usepackage[colorlinks=true,linkcolor=blue!45!black,citecolor=blue!45!black,urlcolor=blue!45!black]{hyperref}
\DeclareMathSizes{9}{9}{6}{5}
\DeclareMathSizes{9.85501}{9}{6}{5}
\IfFileExists{/usr/share/fonts/opentype/noto/NotoSerifCJK-Regular.ttc}{%
  \setmainjfont[Path=/usr/share/fonts/opentype/noto/]{NotoSerifCJK-Regular.ttc}%
}{\setmainjfont{HaranoAjiMincho-Regular}}
\setlength{\textwidth}{15.6cm}
\setlength{\oddsidemargin}{0.25cm}
\setlength{\evensidemargin}{0.25cm}
\setlength{\textheight}{22.6cm}
\setlength{\topmargin}{-0.2cm}
\setlength{\parindent}{1em}
\setlength{\parskip}{1pt}
\numberwithin{equation}{section}
\newtheorem{theorem}{定理}[section]
\newtheorem{maintheorem}{主定理}
\renewcommand{\themaintheorem}{\Alph{maintheorem}}
\newtheorem{lemma}[theorem]{補題}
\newtheorem{proposition}[theorem]{命題}
\newtheorem{corollary}[theorem]{系}
\newtheorem{remark}[theorem]{注意}
\newtheorem{definition}[theorem]{定義}
\newtheorem{example}[theorem]{例}
\newtheorem{conjecture}[theorem]{予想}
\renewcommand{\proofname}{証明}
% 三稿で共通に用いる記号。本文と証明はそれぞれ一つの稿に置く。
\providecommand{\ZZ}{\mathbb Z}
\providecommand{\cK}{\mathcal K}
\providecommand{\cP}{\mathcal P}
\providecommand{\cD}{\mathcal D}
\providecommand{\cC}{\mathcal C}
\providecommand{\cE}{\mathcal E}
\providecommand{\cF}{\mathcal F}
\providecommand{\cM}{\mathcal M}
\providecommand{\fm}{\mathfrak m}
\providecommand{\Perf}{\operatorname{Perf}}
\providecommand{\Ku}{\operatorname{Ku}}
\providecommand{\MF}{\operatorname{MF}}
\providecommand{\HH}{\operatorname{HH}}
\providecommand{\Dsg}{D_{\mathrm{sg}}}
\providecommand{\RHom}{\operatorname{RHom}}
\providecommand{\REnd}{\operatorname{REnd}}
\providecommand{\Ind}{\operatorname{Ind}}
\providecommand{\Def}{\operatorname{Def}}
\providecommand{\thick}{\operatorname{thick}}
\providecommand{\Spf}{\operatorname{Spf}}
\providecommand{\Id}{\operatorname{Id}}
\providecommand{\sym}{\operatorname{sym}}
\providecommand{\Cone}{\operatorname{Cone}}
\providecommand{\Pic}{\operatorname{Pic}}
\providecommand{\otL}{\mathbin{\otimes^{\mathbf L}}}
\newtheorem{problem}[theorem]{問題}

\newcommand{\C}{\mathbb C}
\newcommand{\CC}{\mathbb C}
\newcommand{\PP}{\mathbb P}
\newcommand{\OO}{\mathcal O}
\newcommand{\Sym}{\operatorname{Sym}}
\newcommand{\Spec}{\operatorname{Spec}}
\newcommand{\Proj}{\operatorname{Proj}}
\newcommand{\Hom}{\operatorname{Hom}}
\newcommand{\End}{\operatorname{End}}
\newcommand{\Ext}{\operatorname{Ext}}
\newcommand{\Aut}{\operatorname{Aut}}
\newcommand{\GL}{\operatorname{GL}}
\newcommand{\PGL}{\operatorname{PGL}}
\newcommand{\Gr}{\operatorname{Gr}}
\newcommand{\rank}{\operatorname{rank}}
\newcommand{\qgr}{\operatorname{qgr}}
\newcommand{\QGr}{\operatorname{QGr}}

% 他稿の番号には R (可換復元), M (モジュライ・降下), S (混合族) を添える。
% clean build では scripts/build_papers.py が三稿の aux を順に生成・収束させる。
\newcommand{\SetPaper}[1]{%
  \def\CurrentPaper{#1}%
  \ifstrequal{#1}{main}{}{\externaldocument[main-][nocite]{build/nc-cubic}[nc-cubic.pdf]}%
  \ifstrequal{#1}{moduli}{}{\externaldocument[moduli-][nocite]{build/cubic-moduli-descent}[cubic-moduli-descent.pdf]}%
  \ifstrequal{#1}{mixed}{}{\externaldocument[mixed-][nocite]{build/mixed-sklyanin-cubics}[mixed-sklyanin-cubics.pdf]}%
}
\makeatletter
\def\PaperTag@main{R}
\def\PaperTag@moduli{M}
\def\PaperTag@mixed{S}
\newcommand{\PaperRefText}[2]{\textnormal{[\csname PaperTag@#1\endcsname,~\ref{#1-#2}]}}
\newcommand{\PaperEqRefText}[2]{\textnormal{[\csname PaperTag@#1\endcsname,~\eqref{#1-#2}]}}
\DeclareRobustCommand{\paperref}[2]{\ifmmode\text{\PaperRefText{#1}{#2}}\else\PaperRefText{#1}{#2}\fi}
\DeclareRobustCommand{\papereqref}[2]{\ifmmode\text{\PaperEqRefText{#1}{#2}}\else\PaperEqRefText{#1}{#2}\fi}
\DeclareRobustCommand{\paperpageref}[2]{\textnormal{[\csname PaperTag@#1\endcsname,~p.~\pageref{#1-#2}]}}
\newcommand{\PaperReferenceGuide}{%
  本稿と関連する二稿の参照では、R を可換復元稿\cite{NCCubicRealization}、
  M をモジュライ・降下稿\cite{NCCubicModuliDescent}、
  S を混合族稿\cite{NCCubicMixedSklyanin}の略号とする。
  他稿の結果は、例えば定理 [R, 2.1] のように稿の略号とその稿内の番号を併記する。
  略号のない番号は本稿内の参照である。
}
\makeatother

\SetPaper{moduli}
\title{非可換三次の構造対象付きモジュライと比較核の降下}
\author{}
\date{2026年11月22日・研究稿 v50}
\hypersetup{pdftitle={非可換三次の構造対象付きモジュライと比較核の降下},
pdfsubject={表示認識・Poissonモジュライ・Picard障害・周期位相}}
\begin{document}
\maketitle
\begin{abstract}
非可換三次の構造対象付き圏から、次数付き表示、中心作用、点対応を
認識する条件と、比較核を族および係数体上へ降下する障害を調べる。
良い局所一ブロック型と満支持楕円型について、Veronese 切断環からの復元、
表示 stack と全構造対象保存同値の比較、相対点関手と Poisson 対への射を構成する。
さらに、可換復元の明記した入力の下で、比較核の Picard 障害と周期位相を求め、
点型、中心作用付き、通常の圏同値で最小比較次数が異なることを示す。
全内在的対象の分類や、任意の非可換三次の可換 Kuznetsov 実現は主張しない。
\end{abstract}
\setcounter{tocdepth}{1}
\tableofcontents

% !TEX root = ../cubic-moduli-descent.tex
\section{序論}
\label{sec:moduli-introduction}

\subsection{構造対象から何を認識し、どの体へ降ろせるか}

非可換三次を次数付き代数の中心三次商 $A=B/(w)$ で表示すると、
点曲線、中心、次数一 shift、残余圏という複数の幾何的データが現れる。
本稿の問題は、そのうち何が構造対象を備えた圏
$(\QGr A,\mathcal O_A)$ から内在的に決まり、
表示の選択を変えても族や比較核として降下するか、ということである。
次数一 shift や中心作用を最初から構造対象付き圏の標識に含めず、
回復できる条件と回復を妨げる例を分けて扱う。

対象は、良い局所一ブロック型の三次 $n$-fold と、
$k\ge2$ 個の楕円 Sklyanin 因子を持つ満支持型の三次 $(3k-2)$-fold である。
各定理に明記する PBW 性、楕円因子、普遍中心、相対双対、有限捩れ除外などの
仮定の下で、表示 stack と全 $\mathcal O$-同値の比較を行う。
全ての内在的な非可換三次がこれらの表示で覆われるとは主張しない。

比較核については、幾何学的な圏同値の存在、核の同型類の定義体、
その類を実際の核として実現する scalar 障害、
rotation と指定周期を保存する位相障害を区別する。
中心作用を保つ比較や点型の比較は、通常の Kuznetsov 圏の同値に追加した条件である。
これらの条件を外すと最小比較次数が変わる例も扱う。

\subsection{主要結果と既存理論との関係}

第一に、Serre 関手から得られる Veronese 切断環とその中心を用い、
次数一の表示と全 $\mathcal O$-同値を認識する。
一ブロック型では良い局所モジュライと中心三次の torsor を構成し、
満支持型では第三 Veronese の自己同型持上げと全慣性を求める
（定理\ref{thm:nfm-local-moduli}、\ref{mbca:thm:full-inertia}、
系\ref{mbrr:cor:stack-isomorphism}）。
Artin--Zhang の ample pair と非可換 Proj の理論\cite{LP-AZ}、
Sklyanin 代数と中心拡大の理論\cite{GRSklyaninATV,LBSV}を基礎に、
本稿では普遍中心、余法加群、飽和冪、相対双対を用いて、
指定した中心商における全同値と降下を具体化する。

第二に、相対点関手と中心三次直線を復元し、
良い一ブロック型から Poisson 対のモジュライへの射を構成する。
trace--Hessian 処方が選ぶ三次と残余圏から回復する三次の差を明示し、
代数的補正、点曲線による尺度の選択、有限被覆上の比較を調べる
（定理\ref{thm:point-omarked-descent}、\ref{thm:point-trace-hessian}、
第\ref{sec:poisson-inverse}節、第\ref{sec:nfm-residual-th-obstruction}節）。
Poisson 構造の古典的分類は未公表研究稿\cite{NCCubicA}、
代数的量子化の存在は別稿\cite{CubicQuantization}を入力とする。
本稿が扱うのは、その入力の下でのモジュライ射、補正、相対性と降下である。
点スキームと Poisson 零スキームを同一視する素朴な期待の限界は、
可換復元稿の注意\paperref{main}{ppz:rem:naive-expectation}と具体例を参照されたい。

第三に、比較核の Picard 障害と周期位相を求める。
満支持型の Picard 補正核、全中心同変核の判定、
一ブロック型の点型核とその三分点被覆について、
核の存在と一般係数体上の最小比較次数を区別して記述する
（第\ref{mbpd:sec:picard-descent}節、定理\ref{mbfc:thm:full-criterion}、
定理\ref{main:nfm-third-division-monodromy}）。
Picard gerbe、Kummer 理論、有限群の二重同変化\cite{MEDElagin}や
可換多様体の核の理論\cite{MBFCOrlovAbelian}そのものは既知の入力である。
本稿では指定した中心作用・次数・周期の下で障害を明示し、
その適用範囲と、標識を忘れたときに残る問題を示す。
これらの適用について、文献上の新規性の網羅的確認は未完了である。

\subsection{他稿との分担と主張の範囲}

完備行列分解の移送、一般点の可換復元、および分離型の実際の
Kuznetsov 比較は可換復元稿の結果を用いる
（定理\paperref{main}{cft:thm:complete}、\paperref{main}{gdr:thm:generic}、
\paperref{main}{csc:thm:comparison}）。
本稿はそれらの比較を入力として、表示からの認識、モジュライ上の自然性、
比較核の型と定義体を精密化する。
指定した混合族の中心・特異点・滑らかさの具体的計算は混合族の別稿で扱う。

第\ref{sec:moduli-overview}節に、対象範囲と仮定を保った主要結果の全体像を置き、
以後の各節で証明する。有限計算と証明書を使う箇所では、
その計算が確認する主張と、文献・紙上の証明に依存する主張を分ける。
全 $\mathcal O$-marked モジュライの構成、全 intrinsic 成分の表示被覆、
任意の滑らかな非可換三次の可換 Kuznetsov 実現は未解決である。

\PaperReferenceGuide

% !TEX root = ../cubic-moduli-descent.tex
\section{構造対象付きモジュライと降下の主要結果}
\label{sec:moduli-overview}

本節では、後続の証明に用いる仮定と主要な結論を、
満支持複数ブロック型、局所一ブロック型、比較核の降下の順にまとめる。
各主張の直後に詳しい定理と証明の所在を示す。
全 intrinsic 対象の表示分類や、任意の滑らかな非可換三次の
可換 Kuznetsov 実現を結論するものではない。

\subsection{満支持複数ブロック型で成立すること}
$n=3k-2$, $k\geq2$ とし、$S_i$ は滑らかな点曲線 $E_i$ と
平行移動 $\sigma_i=t_{p_i}$ を持つ三次元楕円 Sklyanin 代数とする。
$p_i,2p_i,3p_i\ne0$、$j(E_i)\ne0,1728$、$j(E_i)$ は相異なると仮定する。
各正則中心三次を $\Omega_i$ とし、
\[
 B=\bigotimes_{i=1}^k S_i,\qquad
 A=B/\left(\sum_i a_i\Omega_i\right),\qquad a_i\ne0
\]
と置く。定理\ref{mbs:thm:presentation-stack}は、共通 grading scalar を除いた
全次数付き表示の滑らかな既約 DM stack $\mathfrak P_k$ を構成する。
その次元は $2k$、幾何点の慣性は
$(\prod_i\mathcal H_{27,i})/\Delta\mu_3$、位数は $3^{3k-1}$ である。
二ブロック三次四様体では、次元四と位数 $243$ を得る。
この結果は相対 Jacobian による各因子の表示と、
次数付き代数からのブロック・中心三次直線の認識を用いる。

一方、点スキームは $\coprod_iE_i$ であり、その双線型イデアルは
実際の二次関係より $9\binom{k}{2}$ 次元大きい
（定理\paperref{main}{tbp:thm:point-graph}、系\paperref{main}{tbp:cor:all-blocks}）。
従って点曲線を知るだけではブロック間の交換関係を選べない。
ブロック別 $\mu_3$ による Zhang twist は、
同じ $(3k-3)$-Veronese と同じ $\OO$-marked 圏を持ちながら、
元の可換 tensor 積表示と次数付き同型でない例を与える。
Zhang の圏同値は既知の入力\cite{TBPZhangTwist}であり、
ここでは満支持型の中心商に適用して次数一復元の具体的な限界を示す。

相対双対では基底直線束
$\mathscr D=\det(\bigoplus_iW_i)\otimes\ell^\vee$ を保持する。
定理\paperref{main}{mbsd:thm:serre}と系\paperref{main}{mbsd:cor:saturation}により、
Serre 関手は $(-3k+3)[3k-2]\otimes\mathscr D$ であり、全整数の切断は飽和する。
全構造対象保存同値が認識する反標準切断環は
$\bigoplus_{q\geq0}A_{(3k-3)q}\otimes\mathscr D^{-q}$ である。
さらに $jp_i\ne0$ $(1\leq j\leq3k-3)$ の有限個の条件の下で、
第三 Veronese の中心を直接計算し、中心平坦性を証明する。
定理\ref{mbv:thm:all-o-third}は fpqc 層として
\[
 \operatorname{Eq}_{\OO}(\operatorname{QGr}A,\operatorname{QGr}A')
 \simeq
 \underline{\operatorname{Isom}}_{\rm gr}(A^{(3)},(A')^{(3)})/\mathbf G_m
\]
を与える。二ブロックでは第三 Veronese は $55$ 生成元と $2619$ 個の独立な
二次関係を持ち、固定模型間の全同値層は有限表示のスキームで表される。
中心一次部分を殺した商 $D=C/(Z(C)_1)$ の点加群は、各因子の第三 Veronese に由来する。
その純成分のイデアル $I_i$ と余法加群 $I_i/I_i^2$ は、元の shift の逆を復元する。
従って、可換 tensor 表示の幾何点では全 $\OO$-同値から
$(E_i,L_i,\sigma_i)$ と次数付き同型類が決まる
（定理\ref{mbiq:thm:triple-recovery}、系\ref{mbiq:cor:commuting-o-rigidity}）。
本稿ではさらに、各 $I_i$ の冪の右飽和から $D$ の全多重次数を回復する。
各剰余次数加群の次数付き自己準同型環が $K$ であることを使うと、
$D=(\bigotimes_iT_i)^{(3)}$ の全自己同型は $\bigotimes_iT_i$ へ持ち上がる
（定理\ref{mbca:thm:center-automorphisms}）。
一方、中心完全交叉の AS--Cohen--Macaulay 性により
$\overline M_j=(E(j))^{(3)}$ の負次数 $\Hom,\Ext^1$ が消える。
これを全中心変数による thickening に適用し、正次数の次数別安定性で
形式的同型を実際の次数付き同型へ戻す
（補題\ref{mbcr:lem:graded-lifting}）。
この二つの議論を接続した定理\ref{mbca:thm:full-inertia}は
\[
 \Aut_{\rm gr}(A^{(3)})/K^\times=\mathcal K_k,
 \qquad
 \Aut_{\OO}(\QGr A)=\mathcal K_k
\]
を与える。後者には上記の有限個の捩れ条件を課し、前者には冒頭の基本条件だけを使う。
中心還元の射影的な核も自明である（定理\ref{mbcr:thm:reduction-rigidity}）。
したがって二ブロックの全慣性の位数は正確に $243$ となる。
さらに、一次中心変形が自明なら中心変数の符号反転が還元核を与えることから、
二ブロックでは $\kappa_C\ne0$ を得る（系\ref{mbcr:cor:nonzero-class}）。

可換三次への支配的特殊化と関係空間の generic flatness による独立の構成から、
非空開上の全 $\OO$-局所型は滑らかな既約 DM stack $\mathfrak M_k^\circ$ となり、
次元は $2k$ である（定理\ref{mbg:thm:generic-stack}）。
さらに、射影枠を使って $\mathfrak P_k^\circ\to\mathfrak M_k^\circ$ の表現可能性を示し、
標数零の一般滑らかさと有限表示降下から、非空開上で有限 \'etale とする。
全慣性の一致によって階数は一となるため、この射は同型である
（系\ref{mbrr:cor:stack-isomorphism}）。
従ってこの開では、非被約底を含む全ての試験基底上で
全 $\OO$-同値の層と次数付き同型のスカラー商が一致する。
全慣性の中心は、順序付きブロック被覆上で
$(\mu_3^k)/\mu_3^{\rm diag}$ と内在的に回復される。
全境界や任意の抽象 Abel 圏変形を覆うことは含めない。

共通スカラーを戻した表示 gerbe の類 $\beta$ は、各非空開で正確に位数三である
（定理\ref{mbsd2:thm:gerbe-period}）。
上界は関係直線 $\ell$ の重み三、下界は実表示慣性の commutator に含まれる
共通 $\mu_3$ による。順序付き被覆では重み一の最小ベクトル束階数も三である。
$\OO(j)$ の重みは $j$ であり、例外的成分は
$\operatorname{Perf}(\mathfrak P_k,-j\beta)$ として降下する。
RHom と対象の重みの相殺により compact 残余 projector 自体は降下する
（定理\ref{mbsd2:thm:projector-descent}）。
標準 rotation は重み一であり、内部 shift を $\rho$ へ移す MF/window 比較がある範囲では、関係直線を含む周期
$\rho^3\otimes\ell^\vee\simeq[2]$ と、方程式を標識した後の周期を区別する
（命題\ref{mbsd2:prop:line-period}）。
これは line-valued MF の既知の翻訳規約
\cite[Definitions~1.1, 1.3]{MBS2PolishchukVaintrob}を本表示へ適用したものである。

中心群 $Z_0$ の作用は、共通 scalar の自然同型を用いて合成則ごと降下する
（定理\paperref{main}{mbza:thm:central-action}）。
一方、その線型的持上げの選択は
$\{(a_i)\in\mathbf F_3^k:\sum_i a_i=1\}$ という
$\widehat Z_0$-torsor であり、基準変更は rotation の指標 twist となる。
ブロック置換に不変な選択は $3\nmid k$ のときに限り存在する
（系\paperref{main}{mbza:cor:permutation-obstruction}）。
この有限群の障害は中心作用自体を妨げず、表示 gerbe の $\beta$ とも別である。
群作用・同変化・双対群作用の一般論は Elagin\cite{MEDElagin}、
随伴と半直交分解の同変性は Shinder\cite{MBZAShinder}を入力とし、
ここでは scalar 自明化と全ての指標変換を明示する。

因子の中心係数を忘れる射は $Z_0$-gerbe であり、
rigidification $\mathfrak P_k\mathbin{/\mkern-6mu/}Z_0$ は
楕円 Sklyanin 因子の組の stack $\mathfrak B_k$ と同型になる
（命題\ref{mbfm:prop:rigidification}）。
これは Romagny\cite[Theorem~5.1]{MBFMRomagny}の一般的な操作を、
係数 torus 上の三乗写像によって具体化したものである。
平面三次の組の stack $\mathfrak E_k$ への射
$\mathfrak B_k\to\mathfrak E_k$ は表現可能・滑らか・全射で相対次元 $k$ となり、
曲線積の普遍族を引き戻せる（命題\ref{mbfm:prop:point-map}、
命題\ref{mbfm:prop:product-family}）。
中心作用付き圏の同変化は、この rigidification 上の圏としても記述できるが、
それだけで曲線積との同値や表示 gerbe の中立化を結論しない。

根については、$\mathsf s=(1)$ としたとき
$\mathsf r\OO\simeq\OO(1)$ を満たす自己同値は一意に
$\mathsf r\simeq\mathsf s\Phi_h$、$h\in\mathcal K_k$ と書け、
$\mathsf r^3\simeq\mathsf s^3$ も自動的に成立する
（系\ref{mbca:cor:marked-roots}）。
その自然同型類は $3^{3k-1}$ 個、全 $\OO$-群による共役類は
$(11^k+2\cdot8^k)/3$ 個であり、二・三ブロックでは $83,785$ となる。
この関手の分類には $6p_i\ne0$ を追加しない。
一つの次数付き自己同型の冪で定まる全代数的 Zhang twist の表示を
直接分類する定理\ref{mbr:thm:algebraic-twist-groupoid}では、この追加条件を使う。
両端が可換 tensor 積表示である一般 twisting system の剛性も保持する
（定理\ref{mbr:thm:commuting-twist-rigidity}）。
根の三乗同型そのものや coherence は標識しておらず、
$\mathsf r^3\simeq\mathsf s^3$ だけから $\mathsf r\OO$ を決定する問題は残る。

Veronese の自己同型持上げについて、Aitzhanova--Makar-Limanov--Umirbaev
\cite[Theorem~1, Corollaries~4, 5]{MBCAAMU}は自由 Poisson 代数と多項式環を扱う。
同論文の $d$-graded は次数を $d$ で剰余化した意味であり、通常の次数保存とは区別する。
また Chan--Young--Zhang\cite[Theorems~0.1, 0.2]{MBCACYZ}は、
パラメータが非自明な1の冪根である skew polynomial 環の Veronese を判別式で扱う。
これらは本稿の楕円 tensor 中心商への直接の持上げ定理ではない。
本稿の認識は、余法加群・飽和冪・負次数 Ext と中心平坦性を組み合わせる。
THCR の理論、Heisenberg 群、次数付き持上げの一般原理自体を新規成果とはせず、
本件への組合せと適用については初出の網羅的確認を継続する。

多ブロック一変数族では、既存の定理\paperref{main}{cft:thm:complete}と
定理\paperref{main}{gdr:thm:generic}を入力として、
$G=(\mu_3)^{k-1}$ 不変な低 rotation 積の積分空間から
分裂した復元方程式 $F_1+\cdots+F_k$ を元の関数体上に得る。
同変核は crossed product 上の perfect 加群として coherence ごと有限型底へ降ろし、
その閉点において有限拡大上の $G$ 同変・rotation・指定周期保存比較を得る
（定理\paperref{main}{med:thm:multiblock-generic}）。
これは指定した形式核の全 jet の代数化ではない。
可換側では Ballard--Favero--Katzarkov の Morita 積公式
\cite[Definition~3.50, Corollary~5.18]{NFBFKKernels}を
$\Gamma=\{(\lambda_i):\lambda_i^3\text{ が共通}\}$ に適用する。
その Orlov 比較\cite[Theorem~6.13]{NFBFKKernels}と合わせると、
全 $k$ で $\operatorname{Ku}(V(\sum_iF_i))^{Z_0}\simeq D^b(\prod_iC_i)$ となる
（定理\paperref{main}{mbrd:thm:split-product}）。
二曲線の場合の categorical cover は既に
\cite[Example~4.7]{MBRDBFKFurther}にあり、Hu\cite{TBHu}は幾何的比較を与える。
本稿では全 $k$ への既知定理の反復適用を明記し、
中心作用の coherent な同定と非可換一般点の有限核降下へ接続する。
可換側の積公式そのものを新しい一般定理とは数えない。
明示的な Jacobian tensor 族ではさらに、
$J_i=g_i+\epsilon_i^2\det\operatorname{Hess}(g_i)/24$ を入力の点三次とすると、
復元模型を $\sum_i c_i\lambda_iJ_i$ と取れる
（定理\paperref{main}{mpc:thm:point-model}）。
Kontsevich のグラフ構成\cite[\S\S5.1,6.1--6.3]{NFMKontsevich}の直和性と、
既存の一ブロック全次数恒等式を用いる。
点三次式の既知の計算\cite[Theorem~7.2]{MoriSmith}と、
本稿の同変残余比較への接続を区別する。

この点三次への対応には、中心線の尺度を含む自然な形がある。
$S_i$ の最大可換商の三次部分 $Q_{i,3}$ が直線束で、
中心三次線 $C_i$ と点三次線 $J_i$ からの射がともに可逆である非空開では、
その商を介した同型 $\theta_i:C_i\simeq J_i$ が全次数付き同型と底変更に自然である
（命題\ref{mbpl:prop:line-comparison}、命題\ref{mbpl:prop:nonempty-open}）。
相対 theta 群の枠を選ぶ滑らかな被覆上で
$\theta_i(\Omega_{H,i})=3(a_i-b_i)F_{H,i}$ と書ける。
既知の Hesse 中心式\cite[\S2, p.~3]{GRSklyaninDLL}の非零性と、
分母のない最大可換商の恒等式を用いて、この比較を基本楕円開全体へ一意に延長する
（定理\ref{mbpl:thm:global-line-comparison}）。
相対被覆には Mumford--Shin の theta 群の理論
\cite{MBPLMumfordTheta,MBPLShinWeil}を用い、
普遍被約重なりで降下した後に任意の非被約底へ引き戻す。
従って中心関係直線 $\ell$ を点三次の分離和へ写すことにより、
滑らかな射影三次の stack への射 $\Phi_{\rm pt}$ を得る。
標的の Brauer--Severi 表示を許せば、
$\beta=\Phi_{\rm pt}^*\beta_{\rm cub}$ と $\ell\simeq\ell_{\rm pt}$ を
同時に、しかも自然に実現する（定理\ref{mbpl:thm:commutative-family}）。
この延長は、全 $\OO$-認識自体の開を広げるものではない。
また、二捩れ境界では比較係数が零となるため、同じ線同型を延長できない。

可換分離和そのものについては、滑らかな準射影基底 $S$ 上で
$f_i\in\Sym^3W_i\otimes L$ を共通の係数直線 $L$ に値を取る滑らかな三次とすると、
相対 Orlov 比較\cite[Corollary~3.4]{MBRPBDFIK}と実 dg Koszul 模型から
\[
 \Perf(E_1\times_S\cdots\times_SE_k)
 \simeq\Ku_{\Perf}(V(\textstyle\sum_i f_i)/S)^{Z_0}
\]
を実際の perfect 核で構成する（定理\paperref{main}{mbrp:thm:relative-product}）。
共通の周期直線を含む Hom 複体の外積を直接照合し、
最小 $A_\infty$ 模型に関する同文献の予想は用いない。
核の表現は To\"en\cite[Theorem~8.15]{MBRPToenMorita}を
滑らか proper な相対 scheme に適用して得る。
構成済みの核・逆核・projector を引き戻すことで、
準コンパクト・分離的 noetherian 底への perfect 圏の比較も従う。
底変更の平坦性や底の被約性は要求しない
（系\paperref{main}{mbrp:cor:base-change}）。
この相対可換比較は既知の Orlov--Thom--Sebastiani 法の本設定への相対的な実施であり、
非可換側からの相対核を新たに存在させるものではない。
体上の非可換実表示には、さらに直接の外積法を適用できる。
THCR の多重斉次イデアルを使って全 $k$ の tensor 積の noetherian 性を証明する。
Keeler\cite[Theorem~2.7, Corollary~3.5]{MBNPKeelerMultiple}の二因子定理を
形式的に反復せず、NC-ampleness と多重次数の端に関する帰納を用いる。
各因子の有限 Koszul 解消の曲率付き stabilization を tensor し、
全体の augmentation stabilization と同定する。
Burke--Stevenson\cite[Theorem~6.4]{LP-BS}の window と合わせて
$(\mathcal K_A^{\mathrm{line}})^{Z_0}\simeq\Perf(\prod_iE_i)$ を得る
（定理\paperref{main}{mbnp:thm:product}）。
有限群の双対同変化と既存の coherent 飽和判定により、
$\mathcal K_A^{\mathrm{line}}=\mathcal K_A^{\mathrm{coh}}$ も従う。
ここでは一変数変形や形式比較の存在を仮定しない。
さらに Keeler の相対 Fujita 消滅
\cite[Theorem~1.5, Definition~2.9, erratum Theorem~E2.5]{NFMKeelerAmple}を使うと、
相対次数が全方向で十分大きい線束の消滅を一様に得られる。
この消滅と同じイデアル層の安定化により、全因子数の noetherian 性を
任意の noetherian 基底上で示す。
滑らかな基底では Burke--Stevenson の次数零環に関する仮定も満たされ、
有限 Koszul--Shamash 外積と window を相対化できる。
その結果、順序付き実表示上で
\[
 \mathcal K_{A/S}^{Z_0}\simeq\Perf(E_1\times_S\cdots\times_SE_k/S)
\]
を smooth proper な相対 compact 圏の同値として構成する
（定理\paperref{main}{mbrn:thm:relative-product}）。
非正則な底への引戻しでは、この構成済み perfect 模型を用い、
絶対的な全 coherent 特異圏との一致を要求しない
（系\paperref{main}{mbrn:cor:base-change}）。
可換積側も、相対核の底変更により元の体上で比較できる
（系\paperref{main}{mbrp:cor:comparison-field}）。
通常の Ku 自体は、積圏の双対指標作用による同変化として記述できるが、
積圏そのものとは同一視しない（系\paperref{main}{mbrd:cor:dual-action}）。
非可換側の因子 rotation は
$R_{i,p_i}=T_{\OO_{E_i}}\circ(t_{p_i})_*\circ(-\otimes L_i)$ である。
$p_i$ を表す実線束の tensor により可換側の rotation と共役し、
指定周期の差を三乗根で補正する。
各曲線の次数三の閉点とこの補正から、
順序付き実表示では次数高々 $3^{2k}$ の有限拡大後に
$\Ku_{\mathrm{coh}}(A)\simeq\Ku(X_{\mathrm{pt}})$ を、
中心作用・rotation・指定周期を保って構成できる
（定理\paperref{main}{mbrc:thm:finite-point-comparison}）。
特に代数閉体上の全ての基本領域の点に適用できる。
従って、同じ偏極点曲線の組を持つ幾何点では、
標識を忘れた残余圏の同値類は一定である
（系\paperref{main}{mbrc:cor:point-fiber-constancy}）。
これは比較関手の一意性や相対的な自明化を意味しない。
この体上の比較を、全 $\OO$-基底上の核の大域的 section、
あるいは指定形式核の全 jet の代数化と同一視しない。
この大域問題について、両側の smooth proper な族と指定周期を固定する。
残余圏の周期には引き続き MF/window 比較を要求する。
局所的に rotation を intertwine する
scalar-simple な核の局所型を固定し、負の自己 Ext 消滅と底変更を仮定する。
核の実現 gerbe の類を $\alpha_P$、関係直線の同期を $\theta$ とすると、
rotation 比較の直線束 $L_P$ には
$q_P:L_P^{\otimes3}\simeq\ell_t\otimes\ell_s^\vee$ が付随する。
$q_P(\eta^3)=\theta$ を満たす比較は $\mu_3$-torsor であり、
その類を $\gamma(P,\theta)$ とすると実現条件は
$\alpha_P=0$ と $\gamma(P,\theta)=0$ である
（定理\ref{mbks:thm:kernel-criterion}）。
前者の band は周期を標識しても $\mathbf G_m$ のままである。
この判定は Kummer 理論\cite{NFMStacksKummer}の相対的な適用であり、
積 Picard 補正型では、さらに両障害を因子データで記述できる。
$p_i$ の Picard gerbe の類を $\delta_i$ とすると、外積の band は
$(u_i)\mapsto\prod_i u_i$ なので $\alpha_{\rm Pic}=\sum_i\delta_i$ となる
（命題\ref{mbpd:prop:product-gerbe}）。
因子比較の Hom 直線は
$\lambda_i=(f_i)_*\Omega^1_{E_i/S}$ に標準同型であり、
指定周期から $q_i:\lambda_i^{\otimes3}\simeq\OO_S$ を得る。
$j_i\ne0,1728$ の開では $h_i=c_{6,i}/c_{4,i}$ として
$u_i=q_i^{\otimes2}(h_i^{\otimes3})$ と置けば、
因子周期の類は $\gamma_i=[u_i]$ である
（補題\ref{mbpd:lem:hodge-line}、系\ref{mbpd:cor:unit-phase}）。
この三乗写像そのものも、自然な中心・点三次同期に対して計算できる。
Jacobian の Weierstrass 表示で
$\omega=dx/(2y+a_1x+a_3)$、$v_\omega=2y+a_1x+a_3$ と置けば、
\[
 q_i(\omega_i^{\otimes3})=-v_{\omega_i}(p_i)^{-1},\qquad
 \gamma_i=[v_{\omega_i}(p_i)]
\]
である（定理\ref{mbhp:thm:period-formula}）。
Weierstrass 座標と微分の変換則は古典的である\cite{NFMSilvermanArithmetic}。
本稿では実際の中心三次を全 Hesse パラメータ曲線上で正則化し、
原点での MF 周期と graph--diagonal の余法写像を比較して、
この指定周期の係数と符号を固定する。
二捩れ点での局所付値だけから大域式を推測する議論ではない。
中心作用を保つ Picard 型の比較 stack は
$\mathscr G_{\rm Pic}\times D$、
rotation と指定周期まで保持するものは
$\mathscr G_{\rm Pic}\times\prod_iT_i$ となる
（定理\ref{mbpd:thm:coherent-corrections}）。
従って前者では $\sum_i\delta_i=0$ と全 $\gamma_i$ の一致、
後者では $\sum_i\delta_i=0$ と全 $\gamma_i=0$ が必要十分である。
前者の体上の存在次数は $3^{2k-1}$ に改良される
（系\ref{mbpd:cor:field-degree}）。
この分類では中心作用の標識を保持しており、
その標識を忘れた全比較核の分類とはしていない。
自然な同期でこれらの位相が消えない例も構成する。
Hesse 族 $(a,b,c)=(1+d,1,c)$ では、$d=0$ の点曲線と中心元は非退化だが、
自然な target の方程式は $3dF_H$ となる。
最高 Koszul syzygy の直線を保持した一因子 window をこの境界へ延長し、
正則な auxiliary 方程式 $F_H$ の周期と比較すると、
自然な位相の剰余は $1\in\mathbb Z/3$ である
（系\paperref{main}{mbrn:cor:factor-boundary}、命題\ref{mbpd:prop:boundary-phase}）。
独立な $k$ 変数を使えば、$\alpha_{\rm Pic}=0$ でありながら、
上記 Picard 型の中心作用付き比較に必要な最小拡大次数が $3^{k-1}$、
rotation と指定周期も保持する場合は $3^k$ となる。
これは有限型族の函数体 $\mathbb C(d_1,\ldots,d_k)$ でも成立する
（系\ref{mbpd:cor:phase-independence}）。
Orlov のアーベル多様体に対する
Fourier--Mukai 完全列\cite[Propositions~2.15, 3.3, Theorem~4.14]{MBFCOrlovAbelian}
を使うと、Orlov 像が block 対角である自己同値は因子ごとへ分解できる。
非同種因子では全ての核がこの条件を満たす。
さらに $k\ge3$ では、全双対中心作用との可換性だけで非対角成分を消せる。
実際、異なる二因子の間の成分 $u$ に対し、第三因子を使った作用を取れば
$(r-I)u=0$ となり、$r^2+r+I=0$ から $3u=0$、従って $u=0$ となる
（補題\ref{mbfc:lem:marked-block-rigidity}）。
この整数 Hom 群上の証明は同種性や CM に依存しない。
三次周期が因子間の余分な shift を除き、全 underlying 比較型は
偏極自己同型と rotation・shift の外積から得られる。
体上の scalar 障害はなお $\sum_i\delta_i$ であり、
位相条件は $\gamma_i+\nu_{E_i}(\phi_i)$ の一致、
指定 rotation・周期付きではその全ての消滅となる
（定理\ref{mbfc:thm:full-criterion}）。
ここで $g(f_i)=c_gf_i$ なら $\nu_{E_i}(\phi_i)=[c_g]$ であり、
一般の体上では非零になり得る。
固定 Hesse level の偏極自己同型は既知の十八元であり
\cite[Section~4, pp.~246--247]{NFMArtebaniDolgachev}、
三次式を厳密に保つ lift があるため、この追加位相は消える。
従って上の独立境界族の最小次数 $3^{k-1}$、$3^k$ は、
全ての中心同変核型について鋭い（定理\ref{mbfc:thm:sharp-all-types}）。
三ブロック以上の作用の正規化群が誘導する $H$ の自己同型は、因子置換だけである。
指標の共通反転は、因子周期の式 $[2]=[-2+3n]$ と矛盾して実現されない
（命題\ref{mbfc:prop:normalizer}）。
従って相異なる $j$ の三因子以上、または非同種の二因子では、
中心群の生成元の標識を外して部分群として保つ場合にも判定は変わらない。
二つの因子が同種の場合は異なり、任意の isogeny から Pell 方程式を使って
無限位数の混合 Orlov isometry を作り、実際の coherent 中央化核へ持ち上げる
（命題\ref{mbfc:prop:two-isogenous-mixing}、補題\ref{mbfc:lem:cyclic-centralizer-lift}）。
その各型の定義体、scalar gerbe、周期位相は引き続き算術的な課題である。

中心部分群を保つという条件そのものは、通常の残余圏の同値からは従わない。
二ブロックの滑らかな分離和四様体では、二つの flex 接線から三つの平面
$P_0,P_1,P_2$ を作ると、中心生成元がそれらを巡回する。
平面イデアル層は Kuznetsov 圏の二次元球面対象であり、
その数値類の Mukai 行列は対角成分 $-2$、非対角成分 $2$ となる。
Seidel--Thomas の既知の球面捻り\cite[\S2b, Proposition~2.10]{MBUNSeidelThomas}を
この三つの対象へ適用し、実際に中心を正規化しない自己同値を得る。
二つの捻りの積の冪から中心の無限個の相異なる共役部分群が生じる
（定理\ref{mbun:thm:non-normal}、系\ref{mbun:cor:infinite-orbit}）。
既存の中心同変比較により、この障害は基本領域の二ブロック非可換残余にも移る
（系\ref{mbun:cor:nc-non-normal}）。
従って無標識比較では、中心作用の共役軌道とその算術的降下を別に調べる必要がある。
この非正規性だけで無標識比較の最小次数が小さくなるとは結論しない。
また、一般 scalar 障害を含む上界の鋭さも未解決である。

相対積核の構造移送も次数ごとに比較する。
共通 scalar の自然自明化と中心同変対象の構造射が、
$\Gamma_k$-次数 $\boldsymbol a$ 上でともに
$\prod_i\zeta_i^{a_i}$ となることから、核は射影的表示へ降下し、
さらに中心 rigidification 上で非捩れの $\Perf(\prod_iE_i)$ を与える
（定理\ref{mbic:thm:projective-product}、命題\ref{mbic:prop:rigidified-product}）。
Picard gerbe と位相 torsor の積による比較 stack も、
両側の偏極・周期の twist が相殺するため因子モジュライへ降下する。
既存の全構造対象認識開では全 $\OO$ 射に自然となる
（定理\ref{mbic:thm:kernel-stack}、系\ref{mbic:cor:all-o}）。
核の stack が降下することと、その大域切断が存在することは異なる。
個々の rotation の非捩れ降下も、相対的な rotation 比較からは従わない。
既に全データが置換両立している場合には、連結な順序付き被覆上の
周期 torsor の切断は $\Hom(\mathfrak S_k,\mu_3)=0$ により置換不変となる
（命題\ref{mbks:prop:permutation-phase}）。
この限定的な消滅は核自身の $\mathbf G_m$-gerbe へは適用できない。
Sklyanin の幾何、Heisenberg 群、Zhang twist、同変圏と有限核降下の既知理論と、
本稿の中心商・第三 Veronese・同変一般点復元への適用を区別する。
これらの組合せの新規性について網羅的な確認は未完了である。
量子化論文では周囲代数を$A$、中心商を$B$と記すのに対し、
本稿では周囲を$B$、中心商を$A$と記す。
中心元の規格化と中心イデアルを保つかどうかは、引用ごとに区別する。

\subsection{可換復元の入力と先行研究}

可換復元稿の一般中心三次変形の基礎定理と、その一ブロック型への適用は、
論理的にも対象範囲としても異なる。完備noetherian局所底での
有限rank graded MF の同値は定理\paperref{main}{cft:thm:complete}、
$\mathbb C[h]_{(h)}$ 上の次数別有限自由な中心三次変形に対する
全coherent残余圏と可換Kuznetsov成分の一般点同値は
定理\paperref{main}{gdr:thm:generic}による。後者の可換方程式は元の関数体にあるが、
比較核を実現するためには有限拡大を許す。その後で初めて、
良い一ブロック型の$\mathcal O$-marked圏の相対構造と
内在的な残余三次$X_{\mathrm{res}}$の復元・降下を扱う。
一般中心三次変形の定理を、一ブロック表示への存在定理に読み替えない。
以下で用いる可換復元の入力は、可換復元稿の完備MF移送・相対window・
一般点降下・coherent飽和の明示的な定理とその仮定である。
適用する対象の表示、底の正則性、指定周期、飽和性の条件を各主張で保持する。

関連する公刊研究との関係は次の通りである。
Etingof--Ginzburg \cite{NFMEG} は三変数 potential の PBW 変形と中心元を構成し、
Mori--Smith \cite{MoriSmith} は三生成元三二次関係の Calabi--Yau 代数を分類する。
Behrend--Noohi \cite[Sections~2.3, 3.1.2]{NFMBehrendNoohi} は
次数スカラーによる rigidification、重み付き代数の gerbe、
determinant character による中立化、および偏極付き Grothendieck 圏を扱う。
これらの一般的方法を新規とは数えず、本稿では全 $\mathcal O$-同値の復元を用いて
shift の標識を外した局所一ブロック型の圏 stack と同定し、
階数 $(n-1,3)$ に依存する、全ての非空開上の Brauer 障害を決定する。
Berger--Taillefer \cite{NFMBeTa} の potential 解消と PBW--CY 定理も入力とし、
本稿では任意個の中心差の Koszul 解消と有限射影模型によって、
非平坦・非被約な底変更にも両立する相対双対と trace を構成する。
体上の CY 性や PBW 定理自体を新規成果としては数えず、
$n$ に依存する全同型問題や torsor の降下と区別する。
Le Bruyn--Smith--Van den Bergh \cite{LBSV} の一点中心拡大における
楕円曲線と八点、および二次関係の復元に対して、本稿は任意個数の中心変数を扱い、
有限平坦な八重被覆と高次元の二十四次完全交叉を得る。
さらに Rains の中心公式 \cite[\S9, p.~1414]{NFMEG} の境界評価から
中心三次の復元を強め、相対 trace と尺度付き Hessian による
Poisson 対への射を構成する。この射と既存量子化の単純な逆写像、
および残余圏からの可換復元とは区別する。
形式量子化の中心と Poisson 中心の同型は
Etingof--Ginzburg\cite[Proposition~8.1.1, equation~(8.1.3)]{NFMEG}
にも既に現れ、その証明では unimodularity から相対形式 Poisson tensor を
Jacobian potential で表している（Section~8.1, pp.~1407--1408）。
本稿で扱うのは、有限の非零パラメータでの三次部分に対する
trace と Hessian の明示的処方、その任意底変更、全同値に関する降下と逆構成である。
Dito\cite[Section~2, equation~(3); Section~3, Theorem~2]{NFMDito} の一般 Poisson
星積の二次項と Weyl gauge、および Shoikhet\cite[Sections~1.3, 2.3.1]{NFMShoikhet}
の内向き wheel の消滅を用いて、固定入力の残余 potential の二次計算を
全混合係数へ広げる。星積や wheel 消滅自体は既知結果であり、
本稿では trace--Hessian 処方との差が全次元で非零となる開条件を調べる。
二次代数と二次 Poisson 構造の形式対応は既に
Kontsevich\cite[Section~2.3, Proposition~6]{NFMKontsevichAlg}にある。
また点曲線の Hessian 公式は Mori--Smith\cite[Theorems~1.2, 7.2]{MoriSmith}の結果である。
本稿はこれらを既知の入力として、混合係数の trace を加えた代数的補正を
O-marked moduli 上へ降下し、全混合係数での一般次数九と残余比較への適合を調べる。
さらに Lindberg--Pym\cite[Section~1.6; Theorem~6.9 and Remark~6.10]{NFMLindbergPym}
の量子トーラスの交換比 $q=e^\hbar$ を既知の入力として、
生成元固定の逆 tensor の五次項を三変数三次の共変式から決定する。
全混合項の定数評価には
Banks--Panzer--Pym\cite[Corollary~1.3; Section~6]{NFMBanksPanzerPym}の
harmonic graph weights と、Buring\cite[Sections~3.4--3.6]{NFMBuringGraphs}の
公開データを用いる。線型入力と二次 tensor に有効な73項を抽出し、
元の weight と多重度、生成元固定再帰を有理数として検証する。
既知の graph 積分値自体を新規成果とは数えず、
本稿の貢献は混合項の一様な公式、残余三次の四次比較類とその非消滅、
および全同値へ降下する追加補正への適用である。
形式対応の存在、点曲線公式、量子トーラスの交換比そのものを新規成果とは数えない。
次数計算に用いる Hesse 還元は、純 Hessian 写像に対する
Ciliberto--Ottaviani\cite[Theorem~5.7]{NFMCilibertoOttaviani}の議論にもあり、
四次の $S$ は古典的な Aronhold 不変式\cite{NFMOttavianiAronhold}である。
$I,J$ も Hessian 写像と Aronhold 四次式の方向微分で表示できる
（補題\ref{lem:nfm-classical-mixed-covariants}）。
本稿ではこれらを一般混合係数と商 stack の次数、および固定逆 tensor の係数へ適用する。
Mayanskiy \cite{Mayanskiy} の三様体における Schouten 方程式の計算に対して、
本稿は高次元での Jacobi イデアル自体を比較する。
affine 超曲面の Poisson 延長を gradient の Koszul 複体で記述する手法には
Fraenkel \cite[Proposition~4.1, Theorem~4.3]{Fraenkel} の先行研究がある。
また、一般の Lie 代数 law の Jacobi スキームに冪零元が現れることは
Carles \cite[Section~8, Lemma~8.35]{Carles} および
Carles--M\'arquez \cite[Section~3]{CarlesMarquez} に既知である。
本稿の Jacobi 係数写像は Landsberg--Ottaviani の Young flattening
\cite{LandsbergOttavianiYoung}、Gesmundo--Landsberg の Koszul flattening
\cite[Proposition~6.1]{GesmundoLandsbergPartials}に同定される。
従って普遍階数上界自体は既知であり、本稿は座標による再証明を与える。
apolar Gorenstein 代数の Betti 数を計算した Boij \cite[Remark~6]{NFMBoijBetti}
との関係も明示し、$N=9,10$ の二次階数の達成を新規とは数えない。
Jacobi 節で扱う中心的対象は、三次超曲面に付随する Hessian--Jacobi イデアルと
Grassmann 錐のスキーム論的な差、次数三の乗法が決める長さ330の平方零構造、
および特殊三次の非零点の局所構造である。
Hessian と apolar 積、Thom--Sebastiani 分解の関係については
Bricalli--Favale--Pirola \cite[Proposition~2.1, Theorem~3.4]{NFMBFFHessian}を参照し、
本稿の横断切片と Jacobi スキームの retract の主張と区別する。
Lin--Zhang \cite{NFLZSerre}、Pirozhkov \cite{NFPirozhkov}、
Huybrechts--Rennemo \cite{NFHuybrechtsRennemo} の既知の圏論的不変量と Torelli 理論は、
本稿の残余変形の入力である。

残余三次の勾配空間による一意性は、Wang\cite[Theorem~3.3]{NFMWangMilnor}と
Hwang\cite[Theorems~1.5, 1.6]{NFMHwangSymmetrizer}の既知結果に基づく。
本稿では階数一の積分核を相対的に作り、O-同値の線束 twist による重みを追跡して、
Brauer--Severi ambient と Poisson 直線への射を降下する。
perfect な対角・核を有限係数へ広げる一般原理と同値 locus の開性には
To\"en--Vaqui\'e\cite[Lemmas~2.10, 3.25]{NFMToenVaquieModuli}を用いる。
これらの一般原理、体上の三次の存在、既知の量子トーラス交換比を新規とは数えず、
相対的な残余射・中心作用の保存・有界次数の関係 locus による
very general な固定 tensor の超越性への適用を区別する。

点曲線と残余 active 曲線の比較には、既存の split Hesse 同定を
active 再帰の制限と係数空間の支配性によって一般混合係数へ広げる。
相対随伴による尺度線束と Hodge 線束の双対性、および
Fulton--Olsson\cite[Section~1; proof of Corollary~2.4]{NFMFultonOlsson}に記載された
判別式と楕円 moduli の $\mu_2$-gerbe は既知の幾何である。
本稿では点曲線の量子化から定まる微分をこの幾何と組み合わせ、
同じ残余三次を標的とする符号商 Poisson 射の支配性・一般有限性を証明する。
平面三次の十八個の自己同型や二つの Jacobian 同型そのものを新規とは数えない。
また Fisher\cite[Theorem~1.1, Lemmas~2.2, 2.6]{NFMFisherTernary}の
三変数三次の不変式環と微分付き軌道分類を用いる。
本稿では active 不変量への幾何学的一般ファイバーが整であることを混合係数も含めて示し、
九次の中間被覆を取り出す。さらに $h$ の全次数にわたる混合重みの有限性と
trace の多項式性から混合係数の多項式逆を作り、共通の三捩れ安定化群の比較により
残余射の一般次数を十八と決定する。
幾何学的一般ファイバーの整性に基づく体論
\cite{NFMStacksRegularFields}を用い、符号二重被覆の類を active 不変量の体へ降ろす。
古典的不変式論や正則な体拡大の一般論自体とは区別する。
さらに Hesse 係数体への二次類の単射性と、三捩れ対称な混合族を用いて、
符号類の消滅を証明する。Benoist\cite[Theorems~1.6--1.7]{NFMBenoistSeparation}による
滑らかな偏極三次の分離性・DM 性から、Poisson 符号商標的の有限対角を得る。
有限表示の係数降下\cite{NFMStacksStackLimits}と正規環への整性によって
指定した一般点での同型を延ばし、混合三角多項式から得る近傍延長を降下する。
正規な底上の有限 étale 被覆の一般ファイバーへの完全忠実性
\cite{NFMStacksEtaleNormal}と有理三捩れ点による twist の排除を組み合わせる。
これらの一般的な降下原理自体は既知であり、残余射への適用と区別する。
比較核の scalar 障害には、Knudsen--Mumford\cite{NFMKnudsenMumfordDet}の determinant と
Lieblich\cite[Sections~2.2.4--2.2.5]{NFMLieblichTwisted}の twisted sheaf の重みの議論を使い、
可換復元の入力における stabilization の Euler 数から全次元の明示的な index bound を取り出す。
さらに次数一の中心正則列に対する graded purity と active 曲線への制限により、
最高次 cohomology が線束となる型で次元に依存しない階数三の証明書を作る。
三変数の $a=0$ 比較は Orlov\cite[Theorem~2.5(iii)]{NFMOrlovSingularities}と、
正則係数環への拡張である Burke--Stevenson\cite[Theorem~6.4(3)]{NFMBurkeStevenson}を入力とする。
window と内部 shift の規約は Galinat\cite[Section~1.1, Lemma~1.2; proof of Lemma~2.7]{NFMGalinat}
に合わせる。これらの同値や楕円曲線上の spherical twist 自体を新規とは数えない。
determinant による Brauer 障害の消去原理自体、および極大部分体による分裂
\cite{NFMStacksSplitting}は既知結果である。
Picard stack の gerbe 構造\cite{NFMStacksPicardGerbe}と
重み零降下\cite[Proposition~4.10]{NFMBerghSchnurer}も既知である。
本稿ではこれらを最高次線束へ適用し、三倍の Picard 類を表す線束を
determinant で明示的に作り、選択済み核を相対的な三分点被覆上で規格化する。
楕円曲線の三倍写像の次数九\cite{NFMStacksMultiplication}自体は新しい結果ではない。
右 twisted homogeneous coordinate ring の shift は
Stafford--Van den Bergh\cite[\S3, equation~(3.1)]{NFMStaffordVdB}の規約を用いる。
この shift と spherical twist の共役を計算し、
指定した中心三次を持つ active 商に対し、相対 THCR と window1 の比較を構成する。
相対 ample pair は Keeler\cite[Proposition~7.2, Lemma~7.6]{NFMKeelerAmple}の枠組みを用い、
底変更は各次数の perfect cohomology\cite{NFMStacksPerfectBasechange}と
MF の余核図式から検証する。これによって固定入力の下で、
最高次線束の Picard 類を点の平行移動へ同定する型の非空性を得る。
体上の同値の存在だけから相対同値や底変更を推論してはいない。

genus-one 曲線に付随する Jacobian 上の Azumaya 代数と、その点における
Brauer 評価は既知である
\cite[Theorems~1.1--1.2, Corollary~1.4]{NFMFisherGenusOneAlgebras}。
Ciperiani--Krashen\cite[Theorems~3.1, 3.7]{NFMCiperianiKrashen}は
Picard 表現障害を Tate pairing と同定し、有限 torsion の Weil cup 積で表示する。
本稿ではこの既知公式を flex torsor、平面同型 torsor、三分点 torsor へ適用し、
二つの active 曲線間の障害の差を具体化する。
さらに一般の $j$ を持つ Hesse 型曲線上で接線の norm を直接計算し、
非零の index 三 Picard 障害を求める。巡回代数や norm 判定\cite{NFMMilneCFT}自体は既知である。
Matsuno\cite[\S2.4, Lemma~2.16, Theorem~3.3, \S4]{NFMMatsunoElliptic}の
幾何的 Sklyanin 表示と cyclic potential、および
Mori--Smith\cite[Proposition~2.11, Theorem~4.2(1)]{MoriSmith}の
基底変更と正則性判定を用い、この例を元の体上で勾配代数に実現する。
ここで本稿が確認するのは、原点や Picard 類の線束代表を仮定しない potential の降下と、
非零 Brauer 類を保持する支配的な全係数表示族である。
これにより一般 split 表示の点障害の正確な index 三を得る。
さらに Fisher\cite[proof of Lemma~2.3, \S9.5]{NFMFisherTernary}の
平面三次の三捩れ作用と Heisenberg 群を使い、
その全 ambient 上の中央化群を計算する。
非零の自明中心部分を伴う表現では中央化群は split な $\operatorname{GL}_{n-1}$ であり、
active 平面上では恒等に作用する。
有限対角延長と有限 torsor の混合定数降下を組み合わせると、
一般係数体で半古典平面同型 torsor が分裂し、残余側の Picard 障害も正確に index 三となる。
Heisenberg 群自体や Hilbert 90 を新規結果とは数えず、
本稿で証明するのはこの full ambient と残余射への適用である。
形式平面比較だけでは同じ結論が出ないことは、明示的な巡回 twist の相殺例で区別する。

楕円曲線上の Fourier--Mukai 同値の数値作用とその核は
Hille--Van den Bergh\cite[Proposition~4.2, Corollary~4.3, Theorem~5.3]{NFMHilleVdBFM}
に記述されている。本稿ではこの graph 認識の議論を実際の核に適用し、
三次 rotation の数値 centralizer の六元計算と合わせて、
原点を仮定しない体上の全 intertwiner の存在判定を与える。
相対的には、非被約底を許す normalized kernel stack と Picard gerbe を同定する。
underlying kernel の分類と、rotation 同型自体および指定周期を標識した際の
追加の $\mu_3$-torsor を区別する。

三・四次元の full 核の同型類については、
Liu\cite[Theorem~4.13]{NFMLiuCubicThreefoldAut}と
Huybrechts\cite[Theorem~1.2(i)]{NFMHuybrechtsK3Cubic}の既知の自己同値分類を用いる。
前者は滑らかな複素三次三次元多様体、後者は very general な複素三次四次元多様体を扱う。
一般三次の幾何自己同型の自明性\cite[Theorems~1.1, 1.5]{NFMPoonenAutomorphisms}と合わせ、
両次元の幾何学的一般点では $\rho^{-1}[1]$ が全群を生成する。
これらは関手の同型類の分類なので、核への適用には追加の一意性が要る。
本稿では projected point の simple 性と
Canonaco--Stellari\cite[Corollary~4.4]{NFMCanonacoStellariKernels}の cohomology 一意性から
三次多様体の projector 核の一意性を証明する。
その後、定数無限巡回群の連続 Galois cocycle が零であることを用いて全 full 類を降下し、
既存の Picard 障害と接続して最小比較次数三を得る。
自己同値群の既知分類自体を新規成果とはせず、very general な分類を Zariski 開の主張に置き換えない。
指定周期を加えた比較では、同じ整数冪による移送が周期位相も保つことを示す。
中心商の conormal 同定を用いて方程式生成元の相対尺度の変換則を追い、
位相を消す規格化と、独立な標的尺度を付した一般族での正確な最小次数九を得る。
さらに full tensor と半古典平面同型から方程式生成元を三乗倍まで規格化し、
周期類 $\gamma_{\mathrm{sc}}$ を定める。
この類は形式比較体へ移すと零になるが、後述の Hesse 境界計算により
元の係数体上では非自明であることが分かる。
Murfet\cite[Theorem~6.2, equation~(6.6)]{NFMMurfetResidue}の residue trace は
可換 MF 側での計算法を与える一方、source trace と CF 移送の scalar 比較を別に要する。

高次元では、rotation 中央化群 $G_\rho$ の定数中心部分
$\langle T\rangle$ を除いても各 class の Galois 安定化群が変わらないことを示す。
奇数次元の数値格子の計算から、残る群は
$G_{\rho,0}/\langle[2]\rangle$ と同定される。
これは点型 class の選択を具体化する簡約であり、その存在の証明とは区別する。
本稿ではさらに、最高次線束 $M_P$ の右随伴像から原点像を復元する条件を課す。
形式恒等分枝における Koszul 計算と proper な導来 Nakayama によって、この強い型の非空性を示す。
三連続の原点対象の生成性と $HH^1=0$、
To\"en--Vaqui\'e\cite[Proposition~3.20, Corollary~3.24]{NFMToenVaquieModuli}の
bounded な対象 moduli と自己同値 locus の開性から、その class を動かす群 $H_p$ は有限となる。
従って無限な自己同値群の全分類を先に解かずに、有限 torsor と Picard gerbe を分離できる。
さらに本稿では、全 $n\ge3$ の一般一ブロック族で $H_p=1$ を証明する。
中央三次商の双加群自己交叉から通常 Ext と Tate Ext を比較し、
双次数による切断で通常 Ext を回復する。
Lu--Palmieri--Wu--Zhang\cite[Theorem~2.4, Corollary~2.5]{NFMLPWZKoszul}の
Adams Koszul 双対と、Braun--Chuang--Lazarev
\cite[Proposition~3.17, Theorems~3.18, 5.3]{NFMBCLLocalization}の
derived localization の忠実性を組み合わせ、原点・rotation を保つ核を
$\operatorname{Aut}_{\mathrm{gr}}(B,k\Omega)/\mathbf G_m$ へ単射で移す。
Keller\cite[Sections~9.3--9.4]{NFMKellerOrbit}の enhanced orbit の関手性を用い、
内部次数を保つ covering category へ戻すことで、対象像だけでなく核の class を検出する。
一般表示の射影安定化子は自明なので、強化型の class は全次元で一意となる。
Koszul 双対や局所化そのものの新規性は主張せず、ここでの cutoff と
比較核の算術的 class 選択への適用を区別する。

同じ復元から、指定周期の位相を有限な低次積へ帰着する。
Picard gerbe 上の $Z_p=q_Y^R\mathcal M[1]$ の Adams graded orbit は
重み零なので元の体へ降り、その二次積が周囲代数 $B_p$、
指定周期と三次積が $\Omega_p\in(B_p)_3$ を与える。
Lu--Palmieri--Wu--Zhang\cite[Corollary~B]{NFMLPWZExt}の関係復元を用い、
三次 tensor の gauge 差を二次関係 ideal で除いた上で、
$f:(B_p,K\Omega_p)\simeq(B,K\Omega)$ に対する
$\gamma_{\mathrm{sc},n}=[f(\Omega_p)/\Omega]$ を証明する。
元の体上の実際の比較核を仮定せず、標的の方程式の尺度と
source の自己交叉周期の符号を同じ式で保持する点が、この適用に必要である。
良い integral 表示と原始的な三次関係が閉 fibre にもある場合には
非分岐性も従う。一方、以下の境界では二次関係を飽和した格子内で
三次関係が境界パラメータの一倍を含み、この十分条件を満たさない。

具体的な $q=2$ の Hesse 境界では点曲線は split $I_3$ に退化する。
Silverman\cite[Chapter~XI, Proposition~9.2]{NFMSilvermanArithmetic}の
既知の Kummer 評価公式へ規格化した接線比を代入し、
三分点被覆が三つの同じ完全分岐三次体に分かれることを計算する。
一方、点線束 $\mathcal O(P-O)$ の三行三列線型表示を、完備化前の trait 上の
有理式として構成する。点尺度を保つ中心座標の変更で、実際の一般残余三次の
平坦な模型を取り、その上の右射影対象を有限 mutation で構成する。
選択した模型の特殊 fibre は特異でもよい。
Eisenbud\cite[Theorem~7.2]{NFMEisenbudCI}の既知の構成による有限 MF を明示し、
Burke--Stevenson\cite[Theorem~6.4(1)]{NFMBurkeStevenson}の相対窓を指定して
右射影対象と同定する。この窓は特異な射影 fibre にも適用され、
内部次数と同じ方程式の MF 周期を保つ。
特殊 fibre の active MF と中心 Koszul 因子の分離により、
全 $n\ge3$、全ての整数 $i,j$ について
\[
 \operatorname{rank}_R H^i\operatorname{RHom}(Z_R,\rho^jZ_R)
 =\binom{n+2}{3i+2j}
\]
となり、これらの Hom は自由で閉・一般 fibre への底変更と可換である。
範囲外の二項係数は零とする。従って指定 target 周期は primitive であり、
必要な $m_2,m_3$ を integral に取れる。

しかし特殊 fibre の二次積の階数は $\binom{n+2}{2}-3$ に落ち、
三次 class も零となる。実際の行列の一次変形を計算すると、
失われる三つの二次関係の Smith 指数は全て一である。
二次関係を飽和して得る代数は全次数で平坦となり、
その active 特殊 fibre は三角形型 Sklyanin 代数 $S(0,-2,-p)$ となる。
指定周期から回復される三次 class の初項は、境界パラメータ $z=q-2$ と
非零中央三次の積である。これだけで一般 fibre の表示同型の integral 性は従わないが、
cyclic potential の alternating 成分から滑らかな potential 三次を回復できる。
Benoist\cite[Theorems~1.6(ii), 1.7, Lemma~2.2(iii)]{NFMBenoistSeparation}の
既知の分離性をこの三次へ適用すると、指定した一般 active 射影同型が延びる。
残る自己同型の中央三次への作用が三乗倍であることを用いて、
全 $n\ge3$ の完備局所位相を $[D_R]$ と決定する
（命題\ref{prop:nfm-boundary-period-local-kummer}）。
この結果は滑らかな特殊残余三次や full 比較核の境界延長を必要としない。
楕円 Kummer 理論、$I_3$ 成分群、相対 Orlov 窓、Shamash 構成や
滑らかな三次の分離性そのものを新規成果とせず、
ここでの integral 行列・低次積・規格化周期への適用を区別する。

五次元の安定性条件には Peize Liu\cite[Theorem~1.1, Remark~5.6]{NFMPeizeLiuFivefold}の
Serre 不変な条件の構成があるが、同論文はその軌道の一意性を仮定していない。
本稿では同論文の Clifford rotation を普遍平面上の相対核として構成し、
$HH^1=0$ と Fano 曲面の連結性から、その核 class の平面非依存性を示す。
Clifford 比較の入力は Kuznetsov\cite[Theorem~4.2]{NFMKuznetsovQuadric}と
Liu\cite[Theorem~2.8, Corollary~3.10]{NFMPeizeLiuFivefold}、
曲面の連結性は Collino\cite[Propositions~1.8, 1.10]{NFMCollinoFivefold}による。
二捩れの実現 gerbe を、既知の奇数 Pl\"ucker degree
$2835$\cite[Example~4.4]{NFMCilibertoZaidenbergFano}によって中立化し、
元の体上の半 rotation の実核を得る。
さらに spinor 列を相対随伴の unit と同定し、相対・絶対 coevaluation を比較すると、
$O\simeq T=\rho^{-1}[1]$ という perfect 核の自然同型を得る。
従って $\epsilon=OT^{-1}\simeq\operatorname{id}$ である。
この同定は平面が定義された任意の滑らかな五次元三次に成立し、
数値作用や安定性作用の忠実性を前提としない。

非分裂表示では、実際の次数三線束と偏極類の区別が指数に現れる。
theta group による divisor-class pair の降下
\cite[Proposition~2.2, Lemma~2.4]{NFMCremonaFisherDescentAlgebra}と
Tate--Weil 公式\cite[Theorem~3.7]{NFMCiperianiKrashen}を用いて、
次数零 Picard 障害が period 三・index 九となる genus-one 曲線を構成する。
Brauer--Severi 平面と三捩れの Heisenberg 表現の関係は
Saltman\cite[Sections~3--4, Theorem~5]{NFMSaltmanSeveri}でも扱われる。
本稿ではこの既知の構成原理を、混合係数と中心三次 torsor の twist によって
全 $n\equiv1\pmod3$ の良い一ブロック成分の任意の非空開へ実現し、
点型 full 比較の九整除性と普遍 Picard gerbe の最小正階数九へ接続する。
表示安定化子が自明な共通稠密開内では、全該当次元で
随伴復元点型の唯一の class により最小実現次数九を達成する。
さらに非分裂 rotation 自身の実現 gerbe が偏極障害 $\beta$ であることを示し、
$\rho^rP$ の障害を $\alpha+r\beta$ と計算する。
次数三 Picard 類の標準的な押し下げと低次元自己同値分類を組み合わせると、
全 full 比較の最小次数は $\alpha\in\langle\beta\rangle$ に応じて一または三となる。
従って四次元の指数九族でも全型の最小次数は三であり、点型を保つ比較との相違が現れる。
曲線の period--index の相違そのもの、既知の degree $2835$、
Clifford 半 rotation の存在自体を本稿の新規成果として数えていない。

さらに Wu\cite[Theorems~1.1, 1.3, Proposition~7.2]{NFMWuStability}と
Aravena--Liu--Tevelev\cite[Theorem~1.2]{NFMAravenaLiuTevelevGepner}は
超曲面圏の Gepner 安定性条件に関する一般的な構成を述べている。
これらの存在定理と、自己同値が特定の成分を保つこと、
その作用の忠実性および係数体上の Galois 降下は異なる問題である。
本稿の有限群への還元は、これらの安定性条件の存在・一意性を証明入力に使わない。
また Iliev--Manivel\cite[Proposition~3.1, Corollary~5.3]{NFMIlievManivelSevenfold}の
既知の球面対象を用いると、一般三次七様体の full 自己同値群は
$\mathbb Z^2$ を含み、rotation 中央化群の指数は無限になる。
球面対象の構成自体は既知入力である。この帰結は、低次元の全群の巡回性を
全高次元へ広げる方法に対する障害を明示する。
Torelli 型の対象復元や Fourier--Mukai partner 数の有限性
\cite{NFPirozhkov,NFMLinRennemoZhangWeighted,NFMLinZengZhangSingularity,NFMLinZhangPartners}を、
核の同型類の一意性または降下と同一視しない。

全次元の三分点被覆については、普遍 Hesse 曲線の真の一般点から translation 群を得る。
Fisher の三捩れ安定化群を $g$ と $J_g$ で同一視して線型 monodromy を調べ、
標準 split 係数体上の affine monodromy が全 $\operatorname{ASL}_2(\mathbf F_3)$ であることを示す。
三倍写像の次数や古典的 Hesse 群自体と、この係数体上の連結性・分裂次数の決定を区別する。

既知の公式からの帰結と、本稿に証明を与える拡張は以下で区別する。
これらの拡張について、独立な新規性の網羅的確認はなお完了していない。

\subsection{中心復元と構造対象付きモジュライ}
以下 $n\ge3$、$N=n+2$、$r=n-1$ とする。
生成元空間を $V$、中心一次空間を $U\subset V$、$\dim U=r$ と書く。
点の空間は $W=V^*$ であり、$\PP(W)$ 上の三次式は
$\Sym^3W^*=\Sym^3V$ の元である。
周囲の二次代数を $B$、中心三次商を $A=B/(w)$ と書く。

\begin{maintheorem}[中心復元と局所一ブロック型モジュライ]
可換 $\C$-代数上の次数一生成代数 $A,A'$ が、それぞれ
中心多項式環 $C=\Sym U,C'=\Sym U'$ 上平坦、
$\rank U=\rank U'\ge2$ であり、
普遍的に $Z(A)_d=\Sym^dU$ と $Z(A')_d=\Sym^dU'$ を満たすとする。
このとき $d$-Veronese 間の次数付き同型は fpqc 局所で $A\simeq A'$ へ持ち上がり、
持上げ層は $\mu_d$-torsor である。

一ブロック PBW・楕円中心条件を満たす非空開上で、
次数スカラーを除いた次数付き周囲代数の同型問題は、
有限安定化子開に制限すると滑らかな既約有限型 DM stack $\mathfrak B_n$ となる。
中心三次直線の stack $\mathfrak C_n$ からの射
$\tau_n:\mathfrak C_n\to\mathfrak B_n$ は、
階数 $s_n$ のベクトル束の加法群の torsor であり、
\[
 \dim\mathfrak B_n=b_n=\frac{n^2+7n-4}{2},\qquad
 s_n=\binom{n+1}{3},\qquad
 \dim\mathfrak C_n=b_n+s_n=\binom{n+2}{3}+3n-2.
\]
普遍相対双対と全同値比較の良い開に制限すると、
$\mathfrak C_n^{\mathrm{good}}$ は fpqc 局所的にこれらの表示を持つ
$\mathcal O$-marked Grothendieck Abel 圏族の stack
$\mathfrak M_n^{\mathrm{loc}}$ に同値である。
大域の次数付き表示を妨げる gerbe の類 $\beta$ は
$\gcd(n-1,3)\beta=0$ を満たす。
$n\equiv1\pmod3$ では、その制限は任意の非空開上で位数がちょうど三である。
\end{maintheorem}

正確な普遍中心条件と証明は定理\ref{thm:nfm-central-veronese}に、
表示空間と全次数付き同型の比較は命題\ref{prop:nfm-graded-presentation}に、
有限安定化子と torsor の証明は定理\ref{thm:nfm-ambient-dimension}、
\ref{thm:nfm-central-torsor}に置く。
相対 Serre 関手が与える $d=r$ にこの復元を適用するため、
楕円自己同型 $\sigma$ には
$\sigma^j\ne1$ $(1\le j\le\max(3,r))$ を課す。
中心スカラーの重みを考慮した束
$\mathcal Hom(\mathcal L,\Sym^3\mathcal U)$ が商 stack 上に降下する。
$\mathcal U$ や中心三次束そのものの降下を主張するものではない。

\begin{center}
\begin{tabular}{r|r|r|r|r}
$n$&可換三次の次元 $M_n$&周囲 $b_n$&中心三次 $s_n$&対 $c_n$\\\hline
3&10&13&4&17\\
4&20&20&10&30\\
5&35&28&20&48\\
6&56&37&35&72\\
7&84&47&56&103\\
8&120&58&84&142
\end{tabular}
\end{center}
ここで $M_n=\binom{n+4}{3}-(n+2)^2=\binom{n+2}{3}$ である。
$c_n-M_n=3n-2$ という数値一致から、可換モジュライへの射の存在や
そのファイバーの次元を結論してはいない。
良い相対 PBW 表示チャートでは、定理\ref{thm:nfm-relative-ambient-cy}により
真正な双加群双対 $\Theta_{B/R}\simeq B(N)[-N]$ を構成する。
定理\ref{thm:nfm-relative-line-cohomology}と\ref{thm:nfm-relative-serre}は、
任意底変更後の compact generation、相対 properness、相対 Serre 関手
$S=(1-n)[n]$ とその trace を与える。
固定表示模型間の係数作用と全 Hom の fpqc 降下を構成することにより、
定理\ref{thm:nfm-fixed-model-comparison}は
$\underline{\operatorname{Isom}}_{\mathrm{gr}}/\mathbf G_m
\simeq\operatorname{Eq}_{\mathcal O}$ を与える。
大域的な次数付き代表の障害は $\operatorname{Pic}(R)[n-1]$ の線束であり、
次数一加群がともに自由なら $\operatorname{Pic}(R)[\gcd(n-1,3)]$ に属する。
$n\equiv1\pmod3$ では、滑らかな affine 楕円曲線上で非零障害の例も構成する
（命題\ref{prop:nfm-nontrivial-three-torsion}）。
切断飽和性自体は中心正則対によって全整数次数で証明する
（補題\ref{lem:nfm-central-saturation}）。
定理\ref{thm:nfm-local-moduli}では、係数変更を圏内の係数作用として定義し、
表示 gerbe 上の重み付き次数加群の tails により対象の効果化を構成する。
表示 gerbe の $H^2$ の類と、指定同値の上の $H^1$ の線束障害は異なる。
定理\ref{thm:nfm-brauer-dense-open}は、$n\equiv1\pmod3$ では
どの非空開にも $\C(s,t)$ 上で非零 Brauer 障害を持つ対象が存在することを示す。
全ての内在的成分の一般点に局所一ブロック表示がある、という分類は主張しない。

\subsection{相対点対応と Poisson 対への射}
\begin{maintheorem}[相対点スキーム、中心復元、Poisson 対]
一ブロック模型の楕円境界、固定点なしの shift、および中心三次の正則切断条件の下で、
\[
 \operatorname{Pt}(A)=E\amalg D,\qquad
 D\subset\PP^{n+1}\text{ は型 }(2,2,2,3)\text{ の完全交叉},
\]
\[
 \dim D=n-3,\qquad\deg D=24,\qquad\omega_D=\OO_D(7-n).
\]
三つの偏微分で定まる $C$ は $\PP^{n-2}$ への次数八の有限平坦被覆である。
$n=4$ で $D$ が滑らかなら、$D$ は種数 $37$ の連結曲線である。
これらは任意底変更と両立する相対完全交叉と八重被覆を与え、
点商の関手、shift のグラフ、中心射影は
$\mathfrak M_n^{\mathrm{loc}}$ の良い開上へ降下する。
ambient は一般には Brauer--Severi 束であり、上の双対化層の式は
各ファイバーでの式である。

非空開上では、埋込み点対応から $B$ と中心直線 $[w]$、従って $A$ を復元できる。
さらに尺度込みの周囲 Poisson tensor $\Pi$ に対し、
中心三次束と古典三次 Casimir 束の間に、純中心三次を固定する同変同型
\[
 \mathfrak d_\Pi:Z(B)_3\xrightarrow{\sim}
 \mathcal R_\Pi=\C F\oplus\Sym^3U
\]
を構成できる。この同型は任意の基底変更と両立する。
従って滑らかな三次を与える非空開
$\mathfrak M_n^{\mathrm{Pois}}\subset\mathfrak M_n^{\mathrm{loc}}$ 上で
\[
 (\mathcal C,\mathcal O)\longmapsto
 \bigl(V(\mathfrak d_\Pi(w)),\Pi|_{V(\mathfrak d_\Pi(w))}\bigr)
\]
が、偏極類を持つ三次 Poisson 対への射を定める。
\end{maintheorem}

相対点関手と束の正確な降下は定理\ref{thm:relative-point-nfold}、
\ref{thm:point-omarked-descent}に、中心直線の復元は
定理\ref{thm:point-central-recovery}に置く。
Poisson 射の標的は、斉次 quadratic Poisson tensor と滑らかな三次 Casimir の
局所閉スキームの $\PGL(V)$ 商として具体的に構成する
（定理\ref{thm:point-trace-hessian}）。
定理\ref{thm:poisson-strong-lift}により affine lift は実は Poisson 対から一意に回復する。
従って追加の lift 標識なしに、Brauer--Severi 周囲空間と偏極類を持つ対を扱える。
$n\equiv1\pmod3$ で実際の相対線束 $\OO(1)$ を標識に加えるには、
上記の Brauer 障害への対処が必要である。

この射は、最大可換代数商の八重有限平坦構造が与える
正規化 trace と、境界 Jacobian 環の次数三部分への Hessian 同型を使う。
中心元の単純な可換化は古典 Casimir の制限と一致しない
（命題\ref{prop:point-casimir-restriction-obstruction}）。
逆構成は次のように強められる。

\begin{maintheorem}[Poisson 対の開部分との同型]
任意の可換 $\C$-代数上の滑らかな三次 $X_H\subset\PP(W)$ について、
各相対双ベクトルは一意な斉次二次 affine 持上げ
$\Pi=\iota_{dH}P$ を持ち、$\iota_{dH}\Pi=0$ を多項式として満たす。
この同定と Poisson 条件の同値は任意基底変更に両立する。

$P$ が非零分解可能な Poisson 対の stack を $\mathfrak P_n^{\mathrm{dec}}$ とする。
周囲代数単独の有限安定化子を要求せず、良い PBW・中心・相対双対・全同値比較の
入力と滑らかな $H_w$ を課した圏 stack $\mathfrak M_n^{\mathrm{sm}}$ には同型
\[
 \mathfrak M_n^{\mathrm{sm}}\simeq
 \mathfrak P_n^{\mathrm{good}}\subset\mathfrak P_n^{\mathrm{dec}}
\]
がある。右辺は非空開であり、両辺は次元 $\binom{n+2}{3}+3n-2$ の
滑らかな既約有限型 DM stack である。逆写像は $\Pi$ の二次関係と
$[w]=\mathfrak d_\Pi^{-1}([H])$ による。表示 gerbe の位数の区分も同じである。

既存の一般三次の被約 Poisson 分類を入力すると、一般三次の非空開上で
これは全非零 Poisson モジュライへの開埋込みとなる。
全 Poisson モジュライのうち滑らかな三次のモジュライを支配する既約成分は一つで、
本構成はその一般部分の開を同定する。
\end{maintheorem}

証明は第\ref{sec:poisson-inverse}節に置く。
これは全 intrinsic 圏モジュライの成分分類でも、Kuznetsov 成分の同値定理でもない。
後者の比較では、むしろ次の障害がある。

\begin{maintheorem}[残余復元との比較障害]
可換復元稿の完備 MF 移送と一般点復元の仮定を満たすとする。
全ての $n\ge3$ に対し、一ブロック中心三次族で
\[
 H^{\mathrm{TH}}=F_0+hxyz,\qquad
 H^{\mathrm{res}}=F_0+
 \frac{h+h^2/3+h^5/12}{1-h^4/4}xyz,\qquad
 F_0=\frac{x^3+y^3+z^3+\sum_{i=1}^{n-1}u_i^3}{3}
\]
となるものがある。二つの滑らかな一般ファイバーは代数閉拡大後も射影同型でない。
従って非可換残余圏と $\operatorname{Ku}(V(H^{\mathrm{TH}}))$ との間には
rotation を保つ dg Morita 同値がない。
$n\not\equiv1\pmod3$ では rotation を指定しない dg Morita 同値もない。
同じ周囲 tensor の滑らかな三次 Casimir の選択だけでは、この障害を解消できない。

一般の混合 Hesse 表示では、短い中心公式から $h^2ABC/6$ を引いて
Hessian 二次 jet を持つ中心元に規格化する。指定した生成元固定の formality 模型の下で
\[
 W_{\mathrm{res}}-H^{\mathrm{TH}}
 =h^2\left(\frac{\mathscr H_F}{24}
       +\frac{\operatorname{tr}(\mathscr H_F)}{12}\right)+O(h^3),
 \qquad \mathscr H_F=\det\operatorname{Hess}_{x,y,z}F.
\]
括弧内の全 Jacobi 環の class は非空 Zariski 開で非零である。
従って特殊ファイバーを恒等に標識した形式比較は、一般の表示係数においても
既に二次 jet で成立しない。
\end{maintheorem}

厳密 Hesse 係数は固定入力の計算であり、次元に一様な比較障害への適用、
固定 tensor の全 Casimir に対する障害、混合係数上の二次 class を本稿で証明する
（第\ref{sec:nfm-residual-th-obstruction}節）。
一般混合係数における有限 jet の障害から、標識を外した全ての一般ファイバーの
圏非同値までを結論してはいない。


\subsection{残余三次に適合する補正}

\begin{maintheorem}[代数的補正と形式的な残余比較]
分解可能な Poisson 対を $(H,P)$ と表し、$K=\operatorname{supp}P$、
$U=\operatorname{ann}K$ とする。$H$ と $H|_{\PP(K)}$ が滑らかな開で、
active 臨界代数の正規化 trace を $\tau_{H,K}$ と置く。
座標不変な三次式
\[
 \mathcal D_P(H)=\langle\det(D_K^2H),P^{\otimes2}\rangle,
 \qquad H^\sharp=H+\frac{\mathcal D_P(H)}{24}
                      +\frac{\tau_{H,K}(\mathcal D_P(H))}{12}
\]
は、非空開で Poisson 対の代数的射 $\Phi:(H,P)\mapsto(H^\sharp,P)$ を定める。
この射は任意底変更と全 O-同値に両立し、始域と終域を適切な非空開へ縮めると、
表現可能な有限 \emph{\'etale} 射で階数 $9$ となる。
補正後の active 平面切断は量子側の楕円点曲線に等しく、
指定した中心正規化では $W_h\equiv H^\sharp\pmod{h^4}$ となる。

一方、滑らかな先頭三次と滑らかな active 制限を持つ一ブロック族について、
固定 formality の生成元固定模型では、形式 Poisson tensor $\pi_h$ は
元の中心一次空間を厳密に保存し、残余 potential $W_h$ に対して
\[
 \pi_h=a_h\,\iota_{dW_h}P_0,\qquad a_h=h+O(h^2)
\]
と一意に書ける。この tensor は周囲二次代数だけから決まる。
その相対化は被約有限型の普遍 PBW 係数模型とその任意底変更で成立する。

二次関係の係数が $h$ に依存せず、Hessian 正規化中心元を用いると、
$\pi_h$ と $a_h$ は奇級数、$W_h$ は偶級数である。
$g=F|_{U=0}$、$S(xyz)=1$ と規格化した四次不変式を $S$ とすると、
全混合係数について
\[
a_h=h+\frac{S(g)}{240}h^5+O(h^7),\qquad
 \pi_h=h\Pi_F+\frac{h^3}{24}\Pi_{\mathscr H_F}
                 +h^5\Pi_{Q_5}+O(h^7)
\]
を得る。また active 制限は
\[
 [h^5]\pi_h|_{U=0}=\frac{S(g)}{80}\Pi_g.
\]
従って $S(g)\ne0$ の稠密開では、固定した $h$ と生成元に関して
代数的 tensor 補正と形式逆 tensor は五次から異なる。
五次 tensor の全混合係数は、明示的な相対三次式 $I(F),J(F)$ を用いて
\[
 \pi_5(F)=\Pi_{Q_5(F)},\qquad
 Q_5(F)=\frac{4I(F)+J(F)-2S(g)F}{240}
\]
と書ける。指定した中心元
$\operatorname{sym}(F-h^2\mathscr H_F/12+h^4\Psi_\lambda(F))$ の残余三次は
\[
 W_h=F+\frac{h^2}{24}\mathscr H_F
       +h^4\left(\Psi_\lambda(F)+\frac{I(F)}{24}-\frac{J(F)}{80}\right)+O(h^6)
\]
である。全 Jacobi 環での有限補正三次との差は
\[
 [h^{-4}(W_h-H^\sharp)]_{h=0}
 =\frac{[I(F)]}{60}+\frac{[\tau_F(I(F))]}{40}
                      +\frac{S(g)[\tau_F(F)]}{80}.
\]
この類は各 $n\ge3$ の混合係数の非空開で非零であり、
$\C[h]/(h^5)$ 上で特殊ファイバーの同定を保つ三次の同型を妨げる。
固定 MF・window 入力の下では、同じ標識を保つ残余圏の同値も妨げる。

一方、式\eqref{eq:nfm-fourth-invariant-correction}の座標不変な三次式
$\mathcal E_P(H)$ による
\[
 (H,P)\longmapsto
 \left(H^\sharp+\mathcal E_P(H),
       \left(1+\frac{S_P(H)}{240}\right)P\right)
\]
は、適切な非空開で任意底変更と全 O-同値へ降下する。
指定中心族に沿い、この出力三次は $W_h$ と $h^6$ を法として、
出力 tensor は $\pi_h$ と $h^7$ を法として一致する。


分離型 $g_\lambda+G(U)$、
$g_\lambda=(x^3+y^3+z^3)/3+\lambda xyz$、$1+\lambda^3\ne0$ では、
$G$ の偏微分の共通零点が原点だけであると仮定する。
量子化尺度 $s$ に対する補正後三次は
\[
 H^{\mathrm{corr}}=g_\lambda+G+
               \frac{s^2}{24}\det\operatorname{Hess}(g_\lambda).
\]
可換復元稿の完備 MF・相対 window・固定標的への有限降下と
coherent 残余圏の飽和性を仮定すると、
全 $n\ge3$ において、ある有限拡大 $K'/\C(s)$ 上で
$(\mathcal K_{B/(w)},\rho)\simeq
(\operatorname{Ku}(V(H^{\mathrm{corr}})),\rho_{\mathrm{corr}})$
が指定三次周期を保って成立する。
\end{maintheorem}

代数的射と次数は定理\ref{thm:nfm-algebraic-hessian-correction}、
\ref{thm:nfm-correction-degree-nine}、形式構成は定理\ref{thm:nfm-residual-formal-poisson}、
jet と五次の相違は系\ref{cor:nfm-mixed-even-jet}、
\ref{cor:nfm-fifth-restriction-obstruction}、混合項の還元は
命題\ref{prop:nfm-mixed-fifth-reduction}、\ref{prop:nfm-mixed-fifth-explicit}、
四次 potential と比較障害は定理\ref{thm:nfm-mixed-fourth-potential}、
系\ref{cor:nfm-corrected-fourth-obstruction}、
追加の代数的射は定理\ref{thm:nfm-fourth-algebraic-correction}、
split 比較は
定理\ref{thm:nfm-split-corrected-residual}に置く。
これらの有限補正式そのものの一般混合係数での全次数比較は主張しない。
内在的な残余三次への全次数比較は、次の定理で別に構成する。
有限 jet の標識付き障害から、一般ファイバーで標識を外した Ku の非同値は結論しない。
従来の $\Phi$ の一般次数九と点曲線の厳密な一致を、
追加補正の射にも自動的に移すことはしない。

\subsection{一般点の代数的残余射と very general な超越性}

\begin{maintheorem}[残余三次と射影 Poisson 対の一般次元での復元]
$n\ge3$ とし、可換復元稿の完備 MF 移送、正則底の window、
一般点生成性と coherent 飽和性を用いる。
良い一ブロック O-marked stack の非空稠密開上で、低 rotation 積
\[
 V_\rho=\operatorname{Nat}^0(\mathrm{id},\rho),\qquad
 Q_\rho=\ker\bigl(\Sym^2V_\rho\to
              \operatorname{Nat}^0(\mathrm{id},\rho^2)\bigr)
\]
から、任意底変更に両立する方程式直線
\[
 \mathcal L=\ker\bigl(\Sym^3V_\rho\to
       V_\rho\otimes(\Sym^2V_\rho/Q_\rho)\bigr)
\]
を復元できる。中心一次作用 $c:U\hookrightarrow V_\rho$ と合わせた
$(\mathcal L,\operatorname{ann}(cU))$ は、尺度を忘れた分解可能 Poisson 対への
全 O-同値に両立する代数的射 $\Psi_{\mathrm{res}}$ を与える。
この射は支配的で、さらに開を縮めれば相対次元一で滑らかである。
有限 \'etale 被覆上で、対応する同じ $X_{\mathrm{res}}=V(\mathcal L)$ に対し
\[
 \mathcal K\simeq\operatorname{Ku}(X_{\mathrm{res}})
\]
が成立し、被覆上で選んだ rotation と指定三次周期も保存できる。

生成元固定の harmonic-angle 模型からの相対 flow を選べば、
一般形式像は $(V(W_h),[\pi_h])$ と同定される。
他方、その実際の $\pi_h$ の固定座標での係数は、
各 $n\ge3$ の very general な滑らかな係数で $\C(h)$ 上代数的でない。
この非代数性は有限代数拡大では解消せず、普遍係数の関数体上でも成立する。
\end{maintheorem}

代数的な族と全同値への降下は定理\ref{thm:nfm-projective-residual-morphism}、
支配性は命題\ref{prop:nfm-projective-residual-dominance}、
比較被覆は定理\ref{thm:nfm-residual-comparison-cover}、
形式像の同定は系\ref{cor:nfm-projective-formal-compatibility}で証明する。
非代数性は定理\ref{thm:nfm-normalized-tensor-transcendence}による。
後者は指定された固定 frame の主張であり、形式座標変更を許した
射影三次の代数化や、任意の intrinsic な尺度の選択を否定しない。
尺度を持つ lift は $\mathcal L\otimes\det K_{\mathrm{res}}$ の非零切断に対応する。
次の定理は、この線束の平方に点曲線から標準切断を与える。

方程式の一意性は比較核の一意性を含まない。
核全体は自己同値 stack の torsor であり、同型類を整合させた後に
スカラー二 cocycle の障害が残る（命題\ref{prop:nfm-residual-kernel-descent}）。
一変数 trait では、形式比較を選んだ滑らかな三次模型の存在も示すが、
その正準性や核の $h=0$ への延長は主張しない
（系\ref{cor:nfm-residual-good-reduction}）。

\subsection{点曲線が定める尺度と一般有限な残余射}

\begin{maintheorem}[楕円点曲線による残余 Poisson 射の規格化]
上と同じ固定入力の下で、良い一ブロック O-marked stack の
非空稠密開 $\mathfrak U$ を取る。
楕円点曲線を $E_{\mathrm{pt}}$、残余 active 曲線を
$C_{\mathrm{res}}=X_{\mathrm{res}}\cap\mathbf P(K_{\mathrm{res}})$ とすると、
全 $n\ge3$ について
\[
 j(E_{\mathrm{pt}})=j(C_{\mathrm{res}}),\qquad
 \mathscr T=\mathcal L\otimes\det K_{\mathrm{res}}
     \simeq\lambda_{\mathrm{res}}^*.
\]
$j=0,1728$ を除けば、点曲線の標準微分は
$\mathscr T^{\otimes2}$ の至る所非零な標準切断を定め、
全 O-同値と任意基底変更に両立する代数的射
\[
 \Psi_{\mathrm{res}}^\pm:\mathfrak U
    \longrightarrow[\mathfrak P_n^{\mathrm{dec}}/\mu_2]
\]
を与える。$\mu_2$ は Poisson tensor の符号を変える。
この射は支配的かつ一般有限で、適切な非空開へ制限すると有限 étale 全射となる。
さらに反対代数の自由な作用 $\iota$ と整な中間被覆 $Z$ により
\[
 \mathfrak U\xrightarrow{\ 2\ }\mathfrak U/\langle\iota\rangle
       \xrightarrow{\ \sim\ }Z\xrightarrow{\ 9\ }Y,
 \qquad \deg\Psi_{\mathrm{res}}^\pm=18
\]
と因子化する。$Y$ は符号商標的の稠密開であり、矢印上は有限 étale 階数を表す。

標準的な有限 étale 階数二の被覆
\[
 \mathfrak D=
 \underline{\operatorname{Isom}}_{\mathrm{gp}}
 (\operatorname{Jac}(E_{\mathrm{pt}}),\operatorname{Jac}(C_{\mathrm{res}}))
 \longrightarrow\mathfrak U
\]
上では尺度を含む射
$\widetilde\Psi_{\mathrm{res}}:\mathfrak D\to\mathfrak P_n^{\mathrm{dec}}$ を得る。
両曲線の平面同型の sheaf $\mathfrak I$ は有限 étale 階数十八であり、
$\mathfrak I\to\mathfrak D$ は $\operatorname{Jac}(E_{\mathrm{pt}})[3]$-torsor である。
点 shift は $\mathfrak D$ 上で残余 active 曲線の平行移動へ自然に移り、
deck 反転はその逆と Poisson tensor の負を同時に与える。
また $L_0=\C(c_4,-c_6)$ を補正前の active 不変量の体とすると、
$\mathfrak D$ の一般点の Kummer 類は、一意な
$\chi_{\mathrm{act}}\in L_0^\times/(L_0^\times)^2$ の引戻しである。
この類は零であり、適切な非空開上で $\mathfrak D$ は分裂する。
二つの切断のうち半古典先頭項で規格化したものは、
全 O-同値に自然な full Poisson 射
\[
 \Psi_{\mathrm{res}}^{\mathrm{sc}}:
 \mathfrak U\longrightarrow\mathfrak P_n^{\mathrm{dec}},
 \qquad\deg\Psi_{\mathrm{res}}^{\mathrm{sc}}=9
\]
を定める。適切な非空開への制限後に有限 \'etale 階数九となり、
$\Psi_{\mathrm{res}}^{\mathrm{sc}}\circ\iota
 =-\Psi_{\mathrm{res}}^{\mathrm{sc}}$ を満たす。
\end{maintheorem}

証明は\S\ref{sec:nfm-residual-point-scale}に置く。
核心となる全次数の恒等式は、生成元固定模型における
\[
 W_h|_{U=0}=b_h\left(g+\frac{h^2}{24}\det\operatorname{Hess}(g)\right),
 \qquad g=F|_{U=0},\quad b_h\in1+hR[[h]]
\]
である。この式は MF 圏の $U=0$ への基底変更を仮定せず、
active 再帰と既存の同座標 Hesse 同定から証明する。
点曲線による tensor は形式 frame で $(h/c_h)\pi_h$ となり、
固定 harmonic tensor と尺度が同じとは主張しない。
また $\mathfrak D$ 上に両 genus-one torsor の同型があるとは限らず、
平面同型にはさらに $\mathfrak I\to\mathfrak D$ を用いる。
正確な次数の証明は、混合係数の全次数三角形と共通の Hesse 安定化群を用いる
（命題\ref{prop:nfm-residual-mixed-triangular}、定理\ref{thm:nfm-residual-exact-degree-eighteen}）。
反対被覆は full Poisson 符号被覆の active な二次 twist であり、その一般点の類は
$\chi_{\mathrm{op}}=\chi_{\mathrm{sgn}}+\chi_{\mathrm{act}}$ を満たす
（系\ref{cor:nfm-opposite-cover-active-twist}）。
Hesse 係数体への二次類の写像は単射である
（補題\ref{lem:nfm-hesse-kummer-injective}）。
混合係数に関する三角多項式は、三捩れ対称点でも滑らかな三次を与える。
標的の有限対角と、指定された一般点同型の正規基底への延長により、
その近傍の Poisson 対を代数的に降下する
（補題\ref{lem:nfm-poisson-target-finite-diagonal}、
\ref{lem:nfm-finite-diagonal-rational-extension}）。
延長上の三捩れ自己同型から有理三捩れ点を得て Jacobian の二次 twist を排除し、
$\chi_{\mathrm{act}}=0$ と半古典規格化された一般次数九の full 射を得る
（定理\ref{thm:nfm-active-sign-vanishing}、系\ref{cor:nfm-semiclassical-full-residual}）。
従って $\chi_{\mathrm{op}}=\chi_{\mathrm{sgn}}$ である。
この延長の特殊ファイバーを低 rotation 積の特殊化から再構成できることや、
比較核がそこへ延びることまでは主張しない。
一般の反対代数の対は O-marked 圏として異なるが、同じ残余三次と、
共通の有限被覆上で同値な残余 dg 圏を持つ
（系\ref{cor:nfm-opposite-same-residual-category}）。

比較核については、split な表示 chart で stabilization の混合 Euler 数を
可換復元の入力における三次 Euler 数に合わせる開かつ閉な数値型を選べる。
その中で核の同型類が整合的に選ばれた後の scalar gerbe の類 $\alpha$ は
\[
 d_n\alpha=0,\qquad
 \operatorname{ind}(\alpha_{\mathrm{gen}})\mid d_n,
 \qquad d_n=3^{\lfloor(n+1)/2\rfloor}
\]
を満たす（命題\ref{prop:nfm-kernel-three-primary}）。
試験対象を置くための表示 gerbe を分裂させる必要がある体上では、
$n\equiv1\pmod3$ のときに限り上界にさらに因子三を加えられる
（系\ref{cor:nfm-generic-kernel-bound}）。
この bound は比較核の同型類を選ぶ前段階を解決せず、
標準二重被覆だけに核が降りることも意味しない。
むしろ既に選んだ同型類の scalar index は、三に素な体拡大では不変である
（系\ref{cor:nfm-kernel-prime-to-three-invariance}）。
全 rotation 軌道の Euler pairing の整数イデアルも正確に $d_n\mathbf Z$ であり、
その反復だけでは bound を改善できない。
選んだ厳密相対 flow は完備中心環の斉次係数商に降り、active MF 比較を与える
（命題\ref{prop:nfm-completed-central-restriction}）。
有限の代数的比較核にこの coherent な中心線型性を降ろす段階は未完成である。

\begin{maintheorem}[active 制限による比較核の三捩れ評価]
残余復元と比較被覆に用いた固定入力を仮定する。
split 表示の一般点 $\Spec K$ 上で、原点剰余体の stabilization $\kappa_{\mathrm{nc}}$ と、
可換側の係数商・window $1$ からなる関手
\[
 q_Y:\operatorname{Ku}(X_{\mathrm{res}})
       \longrightarrow\operatorname{Perf}(C_{\mathrm{res}}),
 \qquad C_{\mathrm{res}}=X_{\mathrm{res}}\cap\PP(K_{\mathrm{res}})
\]
を固定する。一般点の有限分離拡大上に、rotation と指定周期を保つ同値核 $P$ で、
\[
 R_P=q_Y\Phi_P(\kappa_{\mathrm{nc}}),\qquad
 H^{i}(R_P)=0\ (i>-1),\qquad
 H^{-1}(R_P)\text{ が次数零の線束}
\]
を満たすものが存在する。

この型に属する比較核の同型類が $K$ 上で整合的に選ばれ、その実現の
$\mathbf G_m$-gerbe を $\mathcal G$、類を $\alpha$ とする。
すると $\mathcal G$ 上には重み一・階数三の束があり、
\[
 3\alpha=0,\qquad \operatorname{ind}(\alpha)\mid3
\]
である。その endomorphism 代数は $K$ 上の次数三の Azumaya 代数へ降下し、
符号を除いて $\alpha$ を表す。
\end{maintheorem}

証明では、次数一の中心正則列を殺した剰余体の導来像を、内部次数の purity により
非正準に分裂させる。形式比較の後に負の rotation を計算すると、
最高次 cohomology に線束が一つだけ残る。
この cohomology 自体は正準であり、選んだ full kernel の gerbe 上に重み一で降下する。
これを実際の次数三線束 $\OO_{C_{\mathrm{res}}}(1)$ と tensor して押し下げると、階数三束を得る。
型の非空性は命題\ref{prop:nfm-kernel-active-line-nonempty}、
index の証明は定理\ref{thm:nfm-kernel-active-line-bound}に置く。
さらに、この最高次線束の相対 Picard 類を $m_P$ とすると、核の gerbe は
Picard gerbe の $m_P$ による引戻しに同値である
（系\ref{cor:nfm-kernel-picard-obstruction}）。
より強く、scalar 障害は $J_C(S)/3J_C(S)$ を通り、$m_P=3q$ なら消える。
一般にも三分点被覆
\[
 T_P^{(3)}=S\mathbin{\times}_{m_P,J_C,[3]}J_C\longrightarrow S
\]
は階数九の有限 \'etale 全射であり、その上に最高次線束を規格化した
full 比較核が一意な同型を除いて一意に存在する
（定理\ref{thm:nfm-kernel-tripling-cover}、系\ref{cor:nfm-kernel-picard-mod-three}）。
これは選択済みの full 同型類と、固定した表示・window に関する構成である。
また active 曲線上では実際の比較核が存在し、体上では一般直線の切片が
次数三の有限 étale な scalar 分裂被覆を与える
（系\ref{cor:nfm-kernel-active-curve-splitting}）。
表示が非分裂な体では、表示分裂後にこの型となるという仮定の下で、
$n\not\equiv1\pmod3$ なら三、$n\equiv1\pmod3$ なら九が index の上界となる
（系\ref{cor:nfm-kernel-active-line-unsplit}）。
この強い型の評価を、従来の Euler 数だけを指定した型の全ての核へは拡張しない。
また核の同型類の選択や gerbe の中立性は含まない。
active 同値核そのものの選択や coherent な中心線型性の代数的降下は仮定せず、
full kernel から得る一つの対象を用いている。

さらに支配的な完備族上で相対 THCR/window1 比較と底変更を構成する
（命題\ref{prop:nfm-relative-active-thcr}、\ref{prop:nfm-relative-active-window-basechange}）。
従って同じ固定入力の下で、少なくとも一つの支配的な有限 étale 被覆上に、
平面同型 $\varphi:E_{\mathrm{pt}}\simeq C_{\mathrm{res}}$ と比較核 $P$ で
\[
 \begin{gathered}
 \operatorname{Jac}(\varphi)=\theta_{\mathrm{sc}},\qquad
 m_P=\theta_{\mathrm{sc}}(p_{\mathrm{pt}})=p_{\mathrm{res}},\\
 \alpha=\delta_{C_{\mathrm{res}}}(p_{\mathrm{res}})
 \end{gathered}
\]
となる型が存在する。最後の式はこの型の full 同型類を選んだ後の主張である
（命題\ref{prop:nfm-kernel-point-line-nonempty}、
系\ref{cor:nfm-kernel-point-picard-obstruction}）。
$\theta_{\mathrm{sc}}$ は半古典規格化による Jacobian 同型であり、
$\varphi$ はその上の階数九の $J_E[3]$-torsor $\mathfrak I_{\mathrm{sc}}$ から選ぶ。
証明の核は、楕円曲線上の $U_{a,b}=t_a^*\circ(-\otimes P_b)$ に対する
rotation の共役式であり、右 point shift を織り込むと $3a=0$、$b=p_{\mathrm{pt}}$ となる。
同じ部分群の rotation 保存自己同値は Picard 類を変えないため、
これだけで $m_P$ を任意に零へ調整することはできない。
この非空性は、全ての比較核への適用や full 同型類の選択を含まない。

この障害は実際の active 点対応において一般に非零となる。

\begin{maintheorem}[一般点の平面同型と比較核の Picard 障害]
良い split 一ブロック表示の標準一般係数体を $k_n$ とする。
残余復元と完備 MF・window 等の固定入力の下で、任意の $n\ge3$ に対して
\[
 \mathfrak I_{\mathrm{sc}}(k_n)\ne\varnothing.
\]
従って半古典同型 $\theta_{\mathrm{sc}}$ を誘導する偏極曲線同型
$E_{\mathrm{pt}}\simeq C_{\mathrm{res}}$ が元の係数体上に存在する。
生成元固定形式模型の恒等分枝に一致する同型を選ぶと、これは表示変更に共変であり、
良い一ブロック O-marked stack の共通な滑らかな稠密開上の切断
$\varphi_{\mathrm{sc}}\in\mathfrak I_{\mathrm{sc}}(\mathfrak U)$ へ降下する。
この切断は全 O-保存同値と任意の底変更に両立する。

一般 split 係数体の actual active 点曲線とその shift 点について
\[
 \begin{gathered}
 \delta_{C_{\mathrm{res}}}(p_{\mathrm{res}})
   =\delta_{E_{\mathrm{pt}}}(p_{\mathrm{pt}}),\\
 \operatorname{per}\delta_{E_{\mathrm{pt}}}(p_{\mathrm{pt}})
   =\operatorname{ind}\delta_{E_{\mathrm{pt}}}(p_{\mathrm{pt}})=3
 \end{gathered}
\]
となる。$p_{\mathrm{pt}}$ と $p_{\mathrm{res}}$ はそれぞれの Jacobian の
$k_n$-有理点群で三倍ではない。

前述の点型の実際の full 比較核 $P$ が有限拡大 $L/k_n$ 上に存在すれば、
$3\mid[L:k_n]$ である。
従ってこの型の核は元の体および次数が三に素な有限拡大上には存在しない。
この不存在には、元の体での full 同型類の事前選択は不要である。

点型 full 核の同型類を $k_n$ 上で選択済みの場合、その scalar 類は
\[
 \alpha=\delta_{C_{\mathrm{res}}}(p_{\mathrm{res}})
       =\delta_{E_{\mathrm{pt}}}(p_{\mathrm{pt}})
\]
であり、$\operatorname{per}\alpha=\operatorname{ind}\alpha=3$ である。
この scalar 段階には三次分離分裂体があり、最小次数は三である。
\end{maintheorem}

平面同型の分裂は定理\ref{thm:nfm-generic-plane-isom-splitting}で証明する。
まず全次数模型の混合次数が有限であることから、平面同型 torsor の類を
active 係数体 $K=\C(g_\alpha,h)$ へ降ろす。
次に $H_g=\operatorname{Stab}_{\operatorname{SL}_3}(g)$ の全 ambient 上の作用を、
混合項が零の滑らかな対称点まで延長する。
\[
 Z_{\operatorname{PGL}(X\oplus U)}(H_g)\simeq\operatorname{GL}(U)
\]
によりその点の active 平面同型が元の体上に存在し、残った定数 torsor が消える。
最後に全係数を $\widetilde F=hF$ と変数変換すると体は $k_n(h)$ であり、
有限 \'etale scheme の有理点が純超越拡大で増えないことから $k_n$ へ戻せる。
特殊点で比較核を延長することや、形式完備化が三捩れ類を検出することを仮定しない。
表示変更への共変性と stack 上への降下は
命題\ref{prop:nfm-semiclassical-plane-isom-descent}による。
純中心項も含めた TH potential の parabolic 共変性を用い、
せん断による項を落とさずに形式恒等分枝を比較する。
点型の表示非依存性は系\ref{cor:nfm-point-type-O-descent}で別に証明する。

点側の非零性は命題\ref{prop:nfm-active-point-sklyanin-realization}、
\ref{prop:nfm-active-point-obstruction-dominant}と
系\ref{cor:nfm-generic-point-picard-index-three}で示す。
出発点は $\C(a,b,\lambda)$ 上の
$X^3-aY^3-bZ^3+\lambda XYZ=0$ と三捩れ平行移動の障害
$\bigl(a,b+\lambda^3/(27a)\bigr)_3$ の直接計算である。
Jacobian の関数体上で無限位数の点に替えた後、
Sklyanin potential を元の体上へ降下する。
支配性を検査する有限 Hesse 拡大と、非零類を保持する実際の純超越拡大とを区別する。
表示係数体におけるこの障害は、複素数体上の各点で非零という意味ではない。

残余側への移送と full 核への帰結は
系\ref{cor:nfm-generic-residual-picard-index-three}、
\ref{cor:nfm-point-comparison-cover-three-divisibility}による。
この全次元の系は点型を扱う。$n=3,4$ では、以下の自己同値分類を介して
全 full 核へ障害を拡張し、同型類選択を含めて最小次数三を決定する。
一般の二曲線については式\eqref{eq:nfm-picard-plane-transfer}の cup 補正があり、
それは二曲線の genus-one torsor の差だけに依存する
（補題\ref{lem:nfm-picard-torsor-quotient}）。
形式的には恒等に合同な平面同型を持ちながら、この補正が非零となる例もある
（命題\ref{prop:nfm-cyclic-plane-twist}）。
今回の一般点では $[\mathfrak I_{\mathrm{sc}}]=0$ なので、その相殺は起こらない。

active 曲線間の比較については、さらに完全な存在判定がある。
実次数三偏極 $(E,L_E),(C,L_C)$ 上で
$F\rho_{E,p}\simeq\rho_CF$ となる Fourier--Mukai 同値が存在することは、
偏極曲線同型の存在と $\delta_E(p)=0$ に同値である
（系\paperref{main}{cor:nfm-elliptic-any-intertwiner-criterion}）。
数値作用を制限しないこの判定から、一般点でそのような active 同値は存在しない。
同型類の相対 sheaf は $\mathfrak I_\theta$、実現 gerbe は
$\delta_C(\theta(p))$ の引戻しである
（命題\ref{prop:nfm-elliptic-intertwiner-gerbe}）。
full 核から active 核を作る代数的 assignment は未構成なので、
この非存在を full 核そのものの非存在と同一視しない。
この active 分類からの推論と、以下で与える full 核の独立した降下証明を区別する。

\subsection{点型の内在性、低次元の full 降下と三分点被覆}
\begin{maintheorem}[三・四次元の full 比較と周期標識の最小次数]
\label{main:nfm-low-dimensional-full-descent}
固定入力の下で、良い標準 split 一ブロック表示の一般係数体 $k_n$ を考える。
$n=3,4$ なら、幾何的な full 比較核の全ての同型類を $k_n$ 上で選択できる。
各類の実現 gerbe は band を保って互いに同値であり、その共通の Brauer 類は
\[
 \alpha=\delta_{C_{\mathrm{res}}}(p_{\mathrm{res}}),
 \qquad\operatorname{per}(\alpha)=\operatorname{ind}(\alpha)=3
\]
である。従って、点型かどうかにかかわらず実際の full 核は $k_n$ 上にも
三に素な次数の有限拡大上にも存在しない。
一方、ある三次分離体上には実際の点型 full 核が存在する。
これは rotation を intertwine する性質を持ち、
指定周期の追加標識を忘れた full 比較の最小体次数は正確に三である。

両側の三次方程式の生成元から定まる指定周期を加えると、
全 full 同型類に共通な $\mu_3$-torsor が追加される。
source の生成元を保ち、target の生成元を適切な $k_n$-単元倍に取り替えれば、
この torsor を分裂させて標識付き最小次数を三にできる。
一方、target の生成元だけを独立変数 $t$ 倍した $k_n(t)$ 上の一般族では、
標識付き full 比較核の任意の有限定義体の次数は九の倍数となり、最小次数は九である。
\end{maintheorem}

詳細は補題\ref{lem:nfm-projector-kernel-uniqueness}、
\ref{lem:nfm-cyclic-kernel-class-descent}と
定理\ref{thm:nfm-low-dimensional-full-class-index-three}にある。
幾何的自己同値群は実際の $k_n$-線型核 $T=\rho^{-1}[1]$ が生成する $\mathbb Z$ で、
$T^2=\rho$、$T^3=[1]$ となる。
比較核の類の Galois cocycle は連続準同型 $\Gamma_{k_n}\to\mathbb Z$ なので零である。
$T$ との合成が scalar gerbe の band を保つため、全類の障害が一致する。
指定周期に適合する rotation 同型の torsor の Kummer 類を
$\gamma\in k_n^\times/k_n^{\times3}$ とする。
命題\ref{prop:nfm-period-phase-class-invariance}により、$\gamma$ は全 full class に共通である。
source と target の生成元をそれぞれ $u,v$ 倍すると
$\gamma\mapsto[v/u]\gamma$ となる
（命題\ref{prop:nfm-period-equation-frame-change}）。
独立尺度 $t$ の場合、$\gamma=[d]$ と書けばその成分体は
$k_n(t)[s]/(s^3-dt)=k_n(s)$ である。
この純超越体上では $\alpha$ の index 三が保たれるため、
全体の最小次数は九となる（定理\ref{thm:nfm-free-period-frame-degree-nine}）。
両側の同時規格化では位相は変わらない。
元の復元が与える target のデータは方程式直線であるが、
命題\ref{prop:nfm-period-equation-cube-frame}により full tensor と
指定平面同型から生成元を三乗倍まで定められる。
一般係数体ではこの平面同型自体も唯一である
（系\ref{cor:nfm-generic-plane-isom-unique}）。
全 $n\ge3$ の正準強化 class についても、
命題\ref{prop:nfm-adjoint-period-all-dimensions}により
$\gamma_{\mathrm{sc},n}$ が一意に定義され、形式比較体への像は自明である。
混合有理 CF 比較が同じ完備核へ移り、忠実な体拡大から随伴復元条件が反映されるので、
全 class の分類を使わずにこの正準 class を追跡できる。
さらに定理\ref{thm:nfm-period-phase-active-descent}と
命題\ref{prop:nfm-period-phase-invariant-descent}により、これは一意な類
\[
 \overline\gamma_n\in\C(s,t)^\times/\C(s,t)^{\times3},
 \qquad s=c_4(g),\quad t=-c_6(g),\qquad n\ge3
\]
から来る。
有限 MF の内部次数と正の混合重みにより無限 arity の挿入後も混合単項式数を抑え、
$\kappa(m)$ 上に実際の周期付き full Morita 比較を作ることが最初の段階である。
次に線型持上げと周期解を $(f,\eta)\mapsto(\lambda f,\lambda^{-1}\eta)$ で商にし、
周期 torsor 自身の $\operatorname{SL}_3$ 同変構造を得る。
一般 Heisenberg 安定化群の character が消えるため、不変量体への降下が成立する。
これは形式体への Kummer 引戻しが単射だと仮定した結果ではない。

$\Delta=s^3-t^2$、$D=1-s/24-t/216-s^2/6912$ と置く。
Hesse 係数 $q=2+z$ の境界における計算は、次の値を与える。

\begin{maintheorem}[正準周期位相の Hesse 境界値]
\label{main:nfm-period-hesse-boundary}
全 $n\ge3$ の標準 split 一般係数体 $k_n$ 上で、上の固定入力と
三乗まで規格化した方程式を用いる。
$m$ を混合係数の組とし、$R_H=\C(p,m)[[z]]$、
$K_H=\C(p,m)((z))$ を $q=2+z$ に対応する完備 Hesse trait とする。
このとき
\[
 \left.\gamma_{\mathrm{sc},n}\right|_{K_H}=[D_R],\qquad
 D_R=-\frac{(p^3+8)^2}{27}z+O(z^2).
\]
従って $v_D(\overline\gamma_n)=1\pmod3$ であり、
位相被覆はこの境界で分岐指数三を持つ。
全ての正準位相 $\gamma_{\mathrm{sc},n}$ は非自明である。
$\gamma_{\mathrm{sc},n}=[a_n]$ に対し
$M_n=k_n(\sqrt[3]{a_n})$ と置けば $[M_n:k_n]=3$ で、
正準強化 class の指定周期を標識した最小実現次数は
\[
 3\operatorname{ind}(\alpha_{M_n})\in\{3,9\},
 \qquad \alpha=\delta_{C_{\mathrm{res}}}(p_{\mathrm{res}})
\]
となる。$n=3,4$ では、この値は全 full class に共通する。
\end{maintheorem}

局所 Kummer 類の証明は命題\ref{prop:nfm-boundary-period-local-kummer}、
最小次数の公式は系\ref{cor:nfm-adjoint-period-marked-minimum}による。
Hom の平坦性だけでなく、二次関係の飽和、三次積の最初の非零項、
指定した一般 active 表示同型の延長まで確認することで、
付値と剰余単元の双方を比較している。

この torsor が $\Spec\C[s,t,1/(\Delta D)]$ 上で有限 \'etale に延びるという
追加仮定の下では、補題\ref{lem:nfm-period-two-boundary-criterion}の
三候補 $1,[D],[D]^2$ から上の付値が $[D]$ を選ぶ。
従って大域式 $\overline\gamma_n=[D]$ はこの非分岐仮定の下で成立する。
その仮定と異なる次元の大域的な $\overline\gamma_n$ の一致は未証明である。
残る位相の計算は、$\overline\gamma_n[D]^{-1}$ の他の因子での分岐と、
位相体 $M_n$ 上の $\alpha$ の index を対象とする。
後者が正準規格化での標識付き最小次数三／九を決める。

全次元では、原点対象と active 制限に内在的な表示
\[
 \kappa_{\mathrm{nc}}\simeq j_{\mathrm{nc}}^!(\mathcal O)[1],
 \qquad q_Y\simeq\mathbf Li_C^*j_Y
\]
を与える。従って点型は
$Q_P=\mathbf Li_C^*j_Y\Phi_Pj_{\mathrm{nc}}^!(\mathcal O)$ の
最高次層 $H^0(Q_P)$ が $p_{\mathrm{res}}$ を表すという条件となり、
系\ref{cor:nfm-point-type-O-descent}により全 O-同値に適合する。
点型と $\alpha=\delta_C(p_{\mathrm{res}})$ は共通な良い稠密開上で表示非依存である。
最高次線束型の選択済み類には、非分裂 ambient でも
$\alpha=\delta_C(m_P)$ と $3[C]=0$ から $3\alpha=0$ が従うが、
index は既存の三または九の上界と区別する
（系\ref{cor:nfm-unsplit-point-scalar-period-three}）。

\begin{maintheorem}[随伴復元点型の正準 class と最小次数]
\label{main:nfm-finite-adjoint-point-classes}
全 $n\ge3$ の標準 split 一般係数体 $K=k_n$ で、固定入力の仮定を保つ。
$Z=j_{\mathrm{nc}}^!\mathcal O$、$q=\mathbf Li_C^*j_Y$、
$q^R=j_Y^!i_{C,*}$ と置く。
点型核の最高次線束 $M_P=H^0(q\Phi_PZ)$ に対する自然な随伴射
\[
 \Phi_PZ\longrightarrow q^RM_P
\]
が同型となる型は非空である。この型の幾何的 full class は、
有限 \'etale 群 scheme $H_p$ の torsor $S_p$ をなし、その実現 stack は
\[
 S_p\times_K\mathscr P_{p_{\mathrm{res}}}
\]
に同値である。ここで $\mathscr P_{p_{\mathrm{res}}}$ は残余曲線上の
$p_{\mathrm{res}}$ を表す線束の Picard gerbe である。
一般一ブロック族では全 $n\ge3$ で
\[
 H_p=1,\qquad S_p=\operatorname{Spec}K.
\]
従って class は一意である。その scalar gerbe はなお
$\alpha=\delta_{C_{\mathrm{res}}}(p_{\mathrm{res}})$ を持ち、
$\operatorname{per}(\alpha)=\operatorname{ind}(\alpha)=3$ である。
周期同型そのものを標識しないこの型の最小有限実現次数は正確に三となる。
\end{maintheorem}

非空性と有限 torsor の記述は
定理\ref{thm:nfm-adjoint-point-finite-class-torsor}、
最小次数の一般公式は系\ref{cor:nfm-adjoint-point-realization-degree}にある。
全次元の自明性は定理\ref{thm:nfm-origin-rotation-faithful}と
系\ref{cor:nfm-generic-adjoint-class-unique}による。
原点と内部 shift を保つ kernel の作用を全 $A_\infty$ 構造と homotopy まで回復するので、
三連続の生成対象の対象像だけから忠実性を推論してはいない。
元の最高次線束条件だけの全点型や、任意の full 核の分類は主張しない。

\begin{maintheorem}[非分裂点型の指数九と普遍階数の最小性]
\label{main:nfm-unsplit-picard-index-nine}
$n\ge4$、$n\equiv1\pmod3$ とし、固定入力の下で点型と半古典平面同型が
全 $\mathcal O$-降下する共通な良い一ブロック開を取る。
その任意の非空開に、体値点であって
\[
 \operatorname{per}\delta_{C_{\mathrm{res}}}(p_{\mathrm{res}})=3,
 \qquad\operatorname{ind}\delta_{C_{\mathrm{res}}}(p_{\mathrm{res}})=9
\]
となるものが存在する。この点で実際の点型 full 比較核を持つ有限体拡大の次数は、
九の倍数である。表示安定化子が自明な共通稠密開内では、
全該当次元で随伴復元点型の最小実現次数は九である。

普遍 Picard gerbe $\mathscr P_{p_{\mathrm{res}}}$ 上には重み一・階数九の束があり、
底をどの非空開へ縮めても重み一の正の階数は九の倍数である。
表示 gerbe の index 三と、この点の Picard index 九は異なる。
全ての体値点で index 九となること、
共通な良い開の外での class 選択は主張しない。
型を指定しない全比較核については次の定理が別の最小次数を与える。
\end{maintheorem}

証明は命題\ref{prop:nfm-unsplit-genus-one-picard-nine}、
\ref{prop:nfm-unsplit-point-index-nine-family}、
系\ref{cor:nfm-universal-unsplit-picard-rank-nine}、
\ref{cor:nfm-unsplit-point-minimum-nine}にある。
独立な二つの三次 Kummer 座標から、
次数九の division algebra を二重付値で明示する。
この九は周期標識を追加しない点型の scalar 障害から来る。

\begin{maintheorem}[非分裂 rotation と低次元の全型の最小次数]
\label{main:nfm-unsplit-full-minimum}
最高次線束型の full class $P$ を選択し、その障害を
$\alpha=\delta_C(m_P)$、次数三偏極類の障害を $\beta$ とする。
幾何的 rotation の class は Galois 不変で、その実現 gerbe は $\beta$ である。
$\rho^rP$ の障害は $\alpha+r\beta$ であり、$r=\pm1$ なら index は三を割る。

$n=3,4$ で幾何学的自己同値群が $\mathbf Z$ となる範囲では、
$\alpha=\delta_C(p_{\mathrm{res}})$ として全型の最小有限実現次数は
\[
 d_{\mathrm{full}}(K)=\min_{r=0,1,-1}\operatorname{ind}(\alpha+r\beta)
 =\begin{cases}1,&\alpha\in\langle\beta\rangle,\\
 3,&\alpha\notin\langle\beta\rangle.
 \end{cases}
\]
特に前定理の四次元の指数九族では、強化点型の最小次数は九、全型では三である。
周期同型そのものを標識する場合は含まない。
\end{maintheorem}

証明は補題\ref{lem:nfm-unsplit-rotation-gerbe}、
定理\ref{thm:nfm-unsplit-rotation-index-three}、
\ref{thm:nfm-unsplit-lowdim-full-degree}、
系\ref{cor:nfm-unsplit-nine-versus-three}にある。
高次元にも class 選択後の三次上界は有効だが、
rotation で型が変わり得るので、点型自身の指数九は保持される。

五次元では、Clifford 表示の半 rotation は
\[
 O_\Pi\simeq T=\rho^{-1}[1],\qquad O_\Pi^3\simeq[1]
\]
と perfect 核として同定できる
（定理\ref{thm:nfm-fivefold-half-rotation-identification}）。
従って既存の実核降下で得た $O$ も元の体上で $T$ に等しく、
候補 $\epsilon=OT^{-1}$ は恒等である。
この具体的同定と、上記の全次元の安定化群の忠実性は独立の論証である。

\begin{maintheorem}[標準三分点被覆の正確な次数と Galois 群]
\label{main:nfm-third-division-monodromy}
全 $n\ge3$ の同じ一般係数体で
$[3]^{-1}(p_{\mathrm{pt}})\simeq[3]^{-1}(p_{\mathrm{res}})$ は連結な次数九の体である。
一つの三分点を定義する最小次数は九、
Galois 閉包の群は $\mathbf F_3^2\rtimes\operatorname{SL}_2(\mathbf F_3)$、次数は二百十六となる。
この九次体は非 Galois で、真の中間体を持たない。
一つの三分点を持つ最小の Galois 拡大の次数も二百十六である。
\end{maintheorem}

証明は定理\ref{thm:nfm-generic-point-third-division-monodromy}による。
最小三分点次数九は、上の最小 scalar 分裂次数三と異なる。
従って三次分裂体上でも点は三で割れず、九次体自身に三次の中間体はない。
任意の full 核が三分点を選ぶという主張は含まない。
$n\ge5$ でも強化点型の class は一意に選べるが、
全型への非存在障害の拡張は未解決である。

\subsection{Jacobi スキームの次元境界}
\begin{maintheorem}[Jacobi スキームの境界と平方零構造]
滑らかな可換三次 $X_F\subset\PP^{n+1}$ の自然な Jacobi スキームについて、
$3\le n\le7$ の一般三次では
\[
 \operatorname{Pois}(X_F)\simeq\widehat{\operatorname{Gr}}(3,n+2)
\]
である。
$n\ge8$ では任意の滑らかな $F$ についてこの等式は成立せず、
頂点完備化も抽象局所環として同型でない。
Poisson 幾何の研究稿\cite{NCCubicA}の一般三次に対する被約分類を入力すると、この場合の一般の
$\operatorname{Pois}(X_F)$ は非被約であり、
冪零根基の次数二部分の次元は少なくとも
\[
 \epsilon_N=5\binom{N+1}{6}-N\binom N4+\binom N3>0\quad(N=n+2\ge10)
\]
である。$n=8$ ではさらに
\[
 J_{10}/I_F\simeq\C(-2)^{330},\qquad
 H_{R_{10}/I_F}(t)=H_{R_{10}/J_{10}}(t)+330t^2
\]
となる。冪零根基は極大イデアルで消える平方零イデアルであり、
この拡大は完備局所環としても分裂しない。
$n=8$ のこの結論は、既存の被約分類を用いない独立な計算証明による。
さらに全ての $n\ge3$ の一般三次について
\[
 \operatorname{Pois}(X_F)\setminus\{0\}
 =\widehat{\operatorname{Gr}}(3,N)\setminus\{0\},\qquad
 \operatorname{Proj}(R_N/I_F)=\operatorname{Gr}(3,N)
\]
がスキームとして成り立つ。
$n\ge8$ では冪零根基は頂点に支持される非零の有限長加群であり、
頂点の局所環は深さ零である。
一方、全ての $n\ge11$ で、非零 Poisson 点の局所環に
330次元の平方零部分空間を持つ滑らかな特殊三次が存在し、射影化しても非被約である。
\end{maintheorem}

Jacobi スキームの主張は\S\ref{sec:nfold-jacobi-scheme}で証明する。
Jacobi イデアル $I_F$ の $N\binom N4$ 個の生成二次式の間には
$\binom N3$ 個の普遍的な独立従属関係がある。
Pl\"ucker イデアル $J_N$ の二次部分の次元 $5\binom{N+1}{6}$ との比較が上の下界を与える。
$N\le10$ の最大階数には、明示的な整数小行列が法 $32003$ で非零である証明書を付す。
$n=8$ では Pl\"ucker の線形 syzygy を330次元の商へ射影し、
$39600\times39600$ 小行列を法二で検証することにより、次数三以降の差を消す
（定理\ref{thm:nfold-eightfold-square-zero}）。
標数二の三次幾何を使うのではなく、元の整数係数行列の階数を特性零へ移す。
非零分解可能 $P$ の support を $A$ とし、
$K_{F,A}=\{w\in W:T_F(A,w,W)=0\}$ と置くと、法方向の接空間は
$\Lambda^3K_{F,A}$ となる（定理\ref{thm:nfold-normal-tangent}）。
これが一般三次での非零部分の一致と頂点支持を与える。
さらに、support を固定した横断切片そのものを $\Lambda^3K_{F,A}$ 上の
polar Jacobi スキームとして記述する（定理\ref{thm:nfold-transverse-slice}）。
Thom--Sebastiani 直和にはスキームの retract があり、
これを用いて特殊三次の非零点へ冪零構造を移す。
高次数の Hilbert 関数の比較から、冪零部分の最高次数と
頂点の極大イデアルによる消滅指数が次元によらず一様に有界ではないことも示す
（命題\ref{prop:nfold-higher-nilradical}）。
非被約性の結論に用いる被約分類の依存と、方程式数の比較だけで従う結論を分ける。
また、これらの Poisson スキームを全 Abel 圏の Artin 変形空間と同一視したわけではない。


% !TEX root = ../nc-cubic.tex
\section{満支持二ブロック商の中心と次数付き自己同型}
\label{tba:sec:algebra}

本節は $n=4$, $N=6=3+3$ の二ブロック表示を扱う。
点対応によるブロックの認識と、次数付き代数からの中心商の認識を組み合わせる。
得られる自己同型群は\emph{次数付き表示}の群であり、
次数一 shift を忘れた全 $\OO$-保存 Abel 同値の分類とは区別する。

\subsection{設定と低次数中心}

$k$ を標数零の代数閉体とし、$i=1,2$ に対し
\[
 S_i=T(W_i)/(R_i),\qquad \dim W_i=3,
 \qquad T_i=S_i/(\Omega_i)
       =B(E_i,\mathcal L_i,\sigma_i)
\]
とする。ここで $S_i$ は二次の楕円 Sklyanin 代数、
$\Omega_i\in Z(S_i)_3$ は非零な正則中心元、
$E_i\subset\PP(W_i^*)$ は滑らかな平面三次曲線、
$\mathcal L_i=\OO_{E_i}(1)$、$\sigma_i=t_{p_i}$ は平行移動である。
二次関係スキームは $\sigma_i$ のグラフとする。
これらの Sklyanin 入力と
$H_{S_i}(t)=(1-t)^{-3}$ は既知の結果である
\cite[Theorems~6.6, 6.8]{GRSklyaninATV}。
以下の有限個の位数条件を課す：
\begin{equation}\label{tba:eq:finite-order-open}
 \sigma_i^j\ne\operatorname{id}\qquad (i=1,2,\ 1\leq j\leq3).
\end{equation}
これは無限位数を要求する条件ではない。
パラメータ族では有限個の捩れ条件を除く開条件として扱える。

生成元に全次数一を与え、
\begin{equation}\label{tba:eq:two-block-pair}
 B=S_1\otimes_kS_2,
 \qquad w_{a,b}=a\Omega_1+b\Omega_2,
 \qquad A_{a,b}=B/(w_{a,b}),\qquad ab\ne0
\end{equation}
と置く。$B$ の異なるブロックの生成元は可換である。
$\Omega_1$ は $\Omega_1\otimes1$、$\Omega_2$ は
$1\otimes\Omega_2$ の略記である。

\begin{lemma}[テンソル積と楕円商の中心]
\label{tba:lem:tensor-center}
体上の二つの代数 $D_1,D_2$ に対し
$Z(D_1\otimes_kD_2)=Z(D_1)\otimes_kZ(D_2)$ である。
また \eqref{tba:eq:finite-order-open} の下で
\[
 Z(T_i)_j=0\quad(1\leq j\leq3),\qquad
 Z(S_i)_1=Z(S_i)_2=0,\qquad Z(S_i)_3=k\Omega_i.
\]
\end{lemma}
\begin{proof}
テンソル積の元を、第二因子が線型独立となる有限和として書く。
$D_1\otimes1$ の全元との可換性から、第一因子の各係数は
$Z(D_1)$ に属する。その係数を線型独立に取り直して
$1\otimes D_2$ との可換性を使えば第二因子も中心に属する。

$T_i$ を捻った Laurent 関数環
$k(E_i)[t,t^{-1};\sigma_i]$ へ埋め込む。
次数 $j>0$ の非零中心元 $ct^j$ があれば、次数一の二つの切断との
可換性を割り算することにより、$\sigma_i^j$ はそれらの比を固定する。
$\mathcal L_i$ は非常に豊富なので、これらの比は $k(E_i)$ を生成する。
従って $\sigma_i^j=\operatorname{id}$ となり、
\eqref{tba:eq:finite-order-open} に反する。
$S_i\to T_i$ の核は $(\Omega_i)$ だから、次数一・二では中心元が消え、
次数三の中心元は $k\Omega_i$ に属する。逆の包含は中心性による。
\end{proof}

\begin{proposition}[二ブロックの低次数中心と正則性]
\label{tba:prop:low-centers}
\eqref{tba:eq:finite-order-open}--\eqref{tba:eq:two-block-pair} の下で
\begin{align}
 Z(B)_1=Z(B)_2&=0,& Z(B)_3&=k\Omega_1\oplus k\Omega_2,
 \label{tba:eq:ambient-center}\\
 Z(A_{a,b})_1=Z(A_{a,b})_2&=0,&
 Z(A_{a,b})_3&=k\overline\Omega_2
              =k\overline\Omega_1.
 \label{tba:eq:quotient-center}
\end{align}
$w_{a,b}$ は左右正則であり、
\begin{equation}\label{tba:eq:hilbert}
 H_{A_{a,b}}(t)=\frac{1-t^3}{(1-t)^6}.
\end{equation}
特に二次までの積は $B$ と $A_{a,b}$ で同じであるが、
三次中心の次元は二から一へ下がる。
\end{proposition}
\begin{proof}
$B$ の中心の等式は補題\ref{tba:lem:tensor-center}から従う。
$B$ を第一ブロックの次数で増大濾過する。
$w_{a,b}$ の最高記号は $a\Omega_1$ であり、$B$ 上の左右非零因子である。
従って最高記号を比較すれば $w_{a,b}$ も左右非零因子となる。
中心性によりイデアル $(w_{a,b})$ の全元は $w_{a,b}b$ の形に書け、
その最高記号は $a\Omega_1\operatorname{in}(b)$ である。
よって新たな初期関係を生じず、
\begin{equation}\label{tba:eq:active-associated-graded}
 \operatorname{gr}_{W_1}A_{a,b}\simeq T_1\otimes_kS_2.
\end{equation}
全次数ごとの濾過は有限であり、この議論に完備化は要らない。
正則中心元による完全列から \eqref{tba:eq:hilbert} を得る。

全次数 $1\leq d\leq3$ の中心元 $z\in A_{a,b}$ を取る。
第一ブロック次数 $j$ に関する最高記号も
\eqref{tba:eq:active-associated-graded} の中心に属する。
$j>0$ ならば、その成分は
$Z(T_1)_j\otimes Z(S_2)_{d-j}=0$ に属し矛盾する。
従って $z$ は濾過零部分、すなわち自然に埋め込まれた $S_2$ に属する。
$S_2$ との可換性から $z\in Z(S_2)_d$ であり、補題を適用できる。
$\overline\Omega_2$ は非零で中心的、
$a\overline\Omega_1+b\overline\Omega_2=0$ なので
\eqref{tba:eq:quotient-center} を得る。
\end{proof}

\begin{remark}[中心計算の適用範囲]
\label{tba:rem:center-scope}
この命題は三次までの中心を計算する。
$Z(A^{(3)})_1$ を $Z(A)_3$ と自動的に同一視してはいけない。
Veronese の中心は増えることがあり、ここではその別の計算を行っていない。
また一ブロック型の中心 Veronese 復元定理は、階数二以上の中心一次空間を
用いるため、この満支持型にはそのまま適用できない。
相対的な中心核の任意底変更を主張する場合にも、
関係する有限行列の核と余核の局所自由性を別途確認する必要がある。
\end{remark}

\subsection{二次包絡と中心三次直線}

\begin{proposition}[次数付き商からの二次包絡の復元]
\label{tba:prop:quadratic-envelope}
$B=T(V)/(R)$ が次数一生成の二次代数、
$0\ne w\in B_3$ が左右正則な中心元、$A=B/(w)$ であるとする。
次数付き代数 $A$ から
\begin{equation}\label{tba:eq:envelope-recovery}
 B\simeq T(A_1)/\bigl(\ker(A_1\otimes A_1\to A_2)\bigr),
 \qquad kw=\ker(B_3\to A_3)
\end{equation}
を自然に復元できる。
従ってこの形の二つの表示間の全次数付き同型は、
その周囲二次代数間の、中心三次直線を保つ同型へ一意に持ち上がる。
とくに
\[
 \Aut_{\mathrm{gr}}(A)
 \simeq \Aut_{\mathrm{gr}}(B,kw).
\]
\end{proposition}
\begin{proof}
次数三の元 $w$ が生成するイデアルには二次以下の成分がない。
従って $A_1=B_1=V$, $A_2=B_2$ であり、二次乗法核はちょうど $R$ である。
三次の核は、連結次数付き性 $B_0=k$ により $kw$ に等しい。
次数付き同型は次数一部分で決まり、二次乗法核と三次核を保つから、
\eqref{tba:eq:envelope-recovery} によって一意に持ち上がる。
逆に、そのような周囲同型は商へ降りる。
\end{proof}

点対応だけから三次直線を読めるという主張は、この命題に含まれない。
二ブロックでは定理\ref{tbp:thm:point-graph}の点対応が
$[a:b]$ に依存しない一方、\eqref{tba:eq:envelope-recovery} は
次数三までの\emph{乗法}を使って、その表示内の直線を回復する。

\begin{proposition}[中心直線の尺度商]
\label{tba:prop:central-pencil-gerbe}
二因子 $S_i$ と中心元 $\Omega_i$ を固定する。
両係数が非零な三次中心直線の空間は
\[
 J=\{[a:b]\in\PP(k\Omega_1\oplus k\Omega_2):ab\ne0\}
   \simeq\mathbf G_m,
 \qquad r=b/a
\]
である。二つのブロックの grading scalar による作用は
\begin{equation}\label{tba:eq:central-pencil-action}
 (t_1,t_2)\cdot r=(t_2/t_1)^3r.
\end{equation}
従って共通の grading scalar を除いた相対尺度群
$T=(\mathbf G_m^2)/\mathbf G_m^{\mathrm{diag}}$ に対し、
\begin{equation}\label{tba:eq:pencil-gerbe}
 [J/T]\simeq B\mu_3.
\end{equation}
この同値は表示した尺度群による商についてのものであり、
全 $\OO$-marked moduli の同定ではない。
\end{proposition}
\begin{proof}
各ブロックの三次中心元は $\Omega_i\mapsto t_i^3\Omega_i$ と変換される。
従って作用は \eqref{tba:eq:central-pencil-action} となる。
$q=t_2/t_1$ により $T=\mathbf G_m$ と同定すると、
$r=1$ の軌道写像は $q\mapsto q^3$ である。
これは有限 \'etale 全射で核は $\mu_3$ なので、fpqc 商は $B\mu_3$ である。
\end{proof}

\begin{remark}
固定周囲を\emph{同定付き}で扱うと $J$ は一次元である。
しかし代数閉体上で因子の尺度変更を許せば、全ての $ab\ne0$ は同じ軌道になる。
従ってこの射影方向を、そのまま粗モジュライの新しい一次元と数えてはいけない。
一般の体では三乗根の選択が必要であり、尺度作用だけについては
$k^\times/(k^\times)^3$ がその障害を記述する。
これは、他の因子自己同型や非分裂表示まで含む全ての形の分類ではない。
\end{remark}

\subsection{一般二ブロック表示に残る有限慣性}

\begin{theorem}[一般の次数付き慣性の位数 $243$]
\label{tba:thm:generic-inertia}
\eqref{tba:eq:finite-order-open} に加えて
\begin{equation}\label{tba:eq:generic-j}
 j(E_i)\notin\{0,1728\},\qquad j(E_1)\ne j(E_2)
\end{equation}
と仮定する。$G_i=\Aut_{\mathrm{gr}}(S_i)$、
$\chi_i:G_i\to\mathbf G_m$ を $\Omega_i$ への作用の指標とする。
このとき $\mathcal H_i=\ker\chi_i$ は位数 $27$ の有限 Heisenberg 群であり、
任意の $ab\ne0$ について
\begin{align}
 \Aut_{\mathrm{gr}}(A_{a,b})
  &\simeq\{(g_1,g_2)\in G_1\times G_2:
                         \chi_1(g_1)=\chi_2(g_2)\},
 \label{tba:eq:aut-fiber-product}\\
 \Aut_{\mathrm{gr}}(A_{a,b})/\mathbf G_m^{\mathrm{diag}}
  &\simeq(\mathcal H_1\times\mathcal H_2)/\mu_3^{\mathrm{diag}}.
 \label{tba:eq:generic-aut-group}
\end{align}
右辺の位数は $3^5=243$ であり、中心拡大
\begin{equation}\label{tba:eq:generic-aut-extension}
 1\longrightarrow\mu_3\longrightarrow
 \Aut_{\mathrm{gr}}(A_{a,b})/\mathbf G_m^{\mathrm{diag}}
 \longrightarrow E_1[3]\times E_2[3]\longrightarrow1
\end{equation}
を持つ。左の $\mu_3$ は相対ブロック尺度であり、点曲線の各成分上には自明に作用する。
\end{theorem}
\begin{proof}
まず一因子について考える。
幾何的な二次表示 $R_i=I(\Gamma_{\sigma_i})_{1,1}$ により、
次数付き自己同型を全体の grading scalar で割った群は
\[
 G_i/\mathbf G_m
   =\Aut(E_i,\mathcal L_i,\sigma_i)
\]
である。実際、自己同型は関係グラフを保ち、逆に埋込み曲線の同型で
$\sigma_i$ と可換なものは関係空間を保つ。
曲線はその平面を張るので、埋込み曲線への制限の核はない。

flex を原点として $\mathcal L_i=\OO_{E_i}(3\cdot0)$ と書く。
\eqref{tba:eq:generic-j} より、曲線の自己同型は $t_u\varepsilon$、
$\varepsilon\in\{1,-1\}$ の形である。
$t_{p_i}$ と可換である条件は $\varepsilon(p_i)=p_i$ であり、
$2p_i\ne0$ により $\varepsilon=1$ となる。
$t_u^*\mathcal L_i\simeq\mathcal L_i$ は $3u=0$ と同値なので
$G_i/\mathbf G_m=E_i[3]$ を得る。

grading scalar $t$ は $\Omega_i$ に $t^3$ で作用する。
従って $\ker\chi_i$ は $E_i[3]$ の $\mu_3$ による拡大で、位数は $27$ である。
この群が標準 Heisenberg 群となることも具体的に確認できる。
Hesse 座標では、原始三乗根 $\zeta$ に対する
\[
 (x,y,z)\longmapsto(y,z,x),\qquad
 (x,y,z)\longmapsto(x,\zeta y,\zeta^2z)
\]
は二次 Sklyanin 関係を保つ。両者は位数三で、交換子は
非自明な三乗根による grading scalar であるため、位数 $27$ の Heisenberg 群を生成する。
この群が $\Omega_i$ を固定することは既知であり、
有限位数を仮定しない Smith--Tate の証明
\cite[p.~36, Remark~(2) after Lemma~3.4]{GRSklyaninST}
を適用できる。
従って生成した群は $\ker\chi_i$ 全体である。

次に二ブロックを扱う。
系\ref{tbp:cor:block-spans}により、次数付き自己同型は
二つの生成元部分空間 $W_i$ を置換する。
$j(E_1)\ne j(E_2)$ なので置換は恒等であり、各 $W_i$ を保つ。
命題\ref{tba:prop:quadratic-envelope}により周囲代数へ一意に持ち上がるから、
これは $(g_1,g_2)\in G_1\times G_2$ の作用である。
その中心直線の像は
$k(a\chi_1(g_1)\Omega_1+b\chi_2(g_2)\Omega_2)$ である。
$ab\ne0$ より、元の直線を保つことと
$\chi_1(g_1)=\chi_2(g_2)$ は同値である。
これで \eqref{tba:eq:aut-fiber-product} を得た。

共通の指標値を $c$ とし、$t^3=c$ となる $t$ を選ぶと、
$(t^{-1}g_1,t^{-1}g_2)\in\mathcal H_1\times\mathcal H_2$ となる。
従って \eqref{tba:eq:generic-aut-group} の右辺から左辺への射は全射である。
核は $\mathcal H_1\times\mathcal H_2$ に属する共通の grading scalar、
すなわち $\mu_3^{\mathrm{diag}}$ である。
これで位数と拡大 \eqref{tba:eq:generic-aut-extension} が従う。
相対尺度は各 $\PP(W_i^*)$ に自明に作用するので、最後の主張も従う。
\end{proof}

\begin{corollary}[満支持 $\ell$ ブロックに一様な中心と慣性]
\label{tba:cor:k-block-center-inertia}
$k$ を標数零の代数閉体、$\ell\geq2$ とし、
$S_i$ $(1\leq i\leq\ell)$ は上と同じ楕円 Sklyanin 入力を満たすとする。
$\sigma_i^j\ne\operatorname{id}$ $(1\leq j\leq3)$、
$j(E_i)\notin\{0,1728\}$、$j(E_i)$ は互いに異なる、と仮定する。
\[
 B=\bigotimes_{i=1}^{\ell} S_i,\qquad
 w=\sum_{i=1}^{\ell} a_i\Omega_i\quad(a_i\ne0),\qquad A=B/(w)
\]
と置くと、$w$ は左右正則であり、
\begin{gather}
 Z(B)_1=Z(B)_2=Z(A)_1=Z(A)_2=0,\qquad
 Z(B)_3=\bigoplus_{i=1}^{\ell} k\Omega_i,
 \label{tba:eq:k-block-centers}\\
 Z(A)_3=\bigoplus_{i=2}^{\ell} k\overline\Omega_i,
 \qquad H_A(t)=\frac{1-t^3}{(1-t)^{3\ell}}.
 \label{tba:eq:k-block-quotient}
\end{gather}
非零係数だけを持つ中心直線のトーラス $J_\ell$ と相対尺度トーラス
$T_\ell=\mathbf G_m^\ell/\mathbf G_m^{\mathrm{diag}}$ に対し
\begin{equation}\label{tba:eq:k-block-gerbe}
 [J_\ell/T_\ell]\simeq\mathrm B(\mu_3^{\ell-1}).
\end{equation}
さらに $\mathcal H_i=\ker(\Aut_{\mathrm{gr}}(S_i)\to\mathbf G_m)$ を
次数三中心元への作用の核とすると
\begin{equation}\label{tba:eq:k-block-inertia}
 \Aut_{\mathrm{gr}}(A)/\mathbf G_m^{\mathrm{diag}}
 \simeq \left(\prod_{i=1}^{\ell}\mathcal H_i\right)/\mu_3^{\mathrm{diag}},
 \qquad
 \left|\Aut_{\mathrm{gr}}(A)/\mathbf G_m^{\mathrm{diag}}\right|
 =3^{3\ell-1}.
\end{equation}
これは $n=3\ell-2$ の満支持表示についての結果である。
\end{corollary}
\begin{proof}
第一ブロック次数による濾過で
$\operatorname{gr}A=T_1\otimes S_2\otimes\cdots\otimes S_\ell$ となる。
命題\ref{tba:prop:low-centers}の証明をそのまま適用し、
補題\ref{tba:lem:tensor-center}を反復すれば
\eqref{tba:eq:k-block-centers}--\eqref{tba:eq:k-block-quotient}を得る。
係数比 $a_i/a_1$ $(i\geq2)$ には尺度比 $(t_i/t_1)^3$ が作用するので、
トーラスの三倍写像の核から \eqref{tba:eq:k-block-gerbe} が従う。

ブロックの認識にも、二ブロックの点グラフの証明を反復できる。
第一座標に異なる二ブロックの非零成分があれば、交叉可換関係により
第二座標は第一座標に比例する。各非零ブロックで
$\sigma_i$ の固定点を生じるため、これは不可能である。
第一座標が一ブロック $i$ だけにあれば、交叉関係により第二座標も同じブロックにある。
さらに $\sigma_i$ に固定点がないので、二座標の非零な小行列式を用いる
定理\ref{tbp:thm:point-graph}の局所計算が他の全ブロック成分を零にする。
従って非被約な横断成分もなく、点対応は $\ell$ 個のグラフの直和である。
中心元 $\Omega_i$ はそれぞれの点加群を消すから、商 $A$ の点対応も同じである。
各曲線がその平面を張り、$j$ が互いに異なることから、
全次数付き自己同型は各 $W_i$ を保つ。
二次包絡の復元と $a_i\ne0$ により、中心三次直線を保つ条件は
全 $\chi_i(g_i)$ が等しいことである。
定理\ref{tba:thm:generic-inertia}と同じ共通三乗根の議論で
\eqref{tba:eq:k-block-inertia}を得る。
\end{proof}

\begin{remark}[認識問題への帰結]
\label{tba:rem:inertia-scope}
Heisenberg 作用そのものと一因子の中心元は既知の入力である。
ここで組み合わせて得たのは、二次包絡・中心商・ブロックの認識に基づく
二ブロック表示の全次数付き自己同型群である。
この有限慣性は一般開をさらに縮めても消えない。
従って一ブロック一般点の慣性自明性を用いた議論を、そのまま適用することはできない。
一方、次数付き商 stack が構成できることだけでは、その stack が
全 $\OO$-保存 Abel 同値を記述することは証明されない。
その認識と次数一 shift の選択の問題は、別の課題として残る。
\end{remark}

% !TEX root = ../nc-cubic.tex
\section{満支持ブロック型の次数付き表示スタック}
\label{mbs:sec:presentations}

本節では、前二節の認識を族としてまとめる。
構成するのは、全体の grading scalar を除いた\emph{次数付き表示}のスタックである。
次数一 shift を忘れた全 $\OO$-保存 Abel 同値との比較は、
命題\ref{tbp:prop:zhang-ambiguity}の障害を踏まえて別問題とする。
二ブロックの場合だけでなく、中心一次変数のない任意のブロック数を扱う。
以下、基礎体は $\C$、ブロック数は $k\geq2$、$N=3k$, $n=3k-2$ とする。

\subsection{一因子の族と中心直線束}

$W=\C^3$ とし、$\mathcal H\subset\PP(\Sym^3W)$ を
滑らかな平面三次曲線で $j\ne0,1728$ となる開部分とする。
普遍曲線を $\pi:\mathcal E\to\mathcal H$、
その実際の次数三線束を $\mathcal L=\OO_{\mathcal E}(1)$ と書く。
相対 Jacobian $\mathcal J=\Pic^0(\mathcal E/\mathcal H)$ の中で
\begin{equation}\label{mbs:eq:factor-chart}
 \mathcal T=\{(E,p):p\ne0,\ 2p\ne0,\ 3p\ne0\}
       \subset\mathcal J
\end{equation}
を取る。曲線 $E$ 自体に原点を選ぶ必要はない。
genus-one 曲線はその Jacobian の torsor なので、$t_p:E\to E$ は自然に定まる。
$\mathcal H$ は滑らかな既約九次元スキーム、$\mathcal T$ は
滑らかな既約十次元スキームであり、$\mathcal T\to\mathcal H$ は全射である。

\begin{lemma}[普遍二次代数と普遍三次中心]
\label{mbs:lem:factor-family}
$\mathcal T$ 上に、各幾何点で Sklyanin 代数
$S(E,\mathcal L,t_p)$ を与える二次代数 $\mathcal S$ が存在する。
各 $\mathcal S_d$ は階数 $\binom{d+2}{2}$ の有限局所自由加群であり、
任意の底変更と両立する。三次中心
\[
 \mathcal C=\ker\bigl(\mathcal S_3
      \longrightarrow\Hom(W\otimes\OO_{\mathcal T},\mathcal S_4)\bigr),
 \qquad c\longmapsto(v\longmapsto cv-vc)
\]
は直和因子となる直線束で、その形成も任意の底変更と両立する。
$\GL(W)$ はこの構成全体に作用し、grading scalar $t$ は
$\mathcal T$ 上では自明、$\mathcal C$ 上では重み三で作用する。
\end{lemma}
\begin{proof}
普遍曲線を $\mathcal T$ に引き戻し、添字を省略する。
二次関係束を
\begin{equation}\label{mbs:eq:relative-relations}
 \mathcal R=\ker\left(
 W\otimes W\otimes\OO_{\mathcal T}
 \longrightarrow\pi_*(\mathcal L\otimes t_p^*\mathcal L)\right)
\end{equation}
で定める。第二因子の切断は $t_p^*$ で引き戻してから乗じる。
$3p\ne0$ により、各幾何ファイバーで
$t_p^*\mathcal L\not\simeq\mathcal L$ である。
命題\ref{tbp:prop:quadratic-saturation}の基点のない pencil の議論で
この乗法は全射となる。次数六の線束の $H^1$ は消えるので、
右辺は底変更に両立する階数六束であり、$\mathcal R$ は階数三の直和因子である。
$\mathcal S=T(W\otimes\OO_{\mathcal T})/(\mathcal R)$ と置く。

各ファイバーの AS 正則性・Hilbert 級数・点グラフ・正則中心三次は
Sklyanin 代数の既知の幾何的構成から従う
\cite[Theorems~6.6, 6.8]{GRSklyaninATV}
\cite[\S1, pp.~20--21]{GRSklyaninST}。
相対的な局所自由性は、ファイバーの次元だけを理由に仮定しない。
固定した次数 $d$ では、関係の全ての挿入を用いて
$\mathcal S_d$ を有限自由加群間の行列の余核として局所的に表示できる。
その行列の階数は全幾何点で一定である。
基底 $\mathcal T$ は被約なので、一つ上の大きさの全小行列式は恒等的に零となる。
最大非零小行列式を可逆にした開集合では行列は一定階数の標準形となり、
余核は所定の階数で局所自由、底変更にも両立する。
この議論を各次数に適用する。次数ごとに別の局所自明化でよく、
無限個の開集合の交叉を非空と仮定していない。

中心の交換子行列も同様に扱える。
補題\ref{tba:lem:tensor-center}により、各幾何点の三次中心は一次元である。
従ってその行列は階数九であり、同じ小行列式の議論で
核 $\mathcal C$ は直和因子、像と余核も局所自由となる。
そのため非平坦・非被約な試験基底への変更でも中心核の等式が保たれる。
座標変換への共変性は \eqref{mbs:eq:relative-relations} と交換子の定義から従い、
三次元への scalar 作用は $t^3$ である。
\end{proof}

\subsection{中心商のパラメータと全次数付き同型}

$\mathcal U_k\subset\mathcal T^k$ を $j(E_i)\ne j(E_j)$ $(i\ne j)$ の開部分とし、
そこへ引き戻した中心直線束を $\mathcal C_i$ と書く。
本節の $\PP(\bigoplus_i\mathcal C_i)$ は直線部分束のパラメータ空間とする。
各成分への射が可逆となる開部分
\begin{equation}\label{mbs:eq:center-chart}
 \mathcal X_k=\left\{\ell\subset\bigoplus_{i=1}^k\mathcal C_i:
       \ell\longrightarrow\mathcal C_i\text{ が全て同型}\right\}
       \subset\PP\left(\bigoplus_{i=1}^k\mathcal C_i\right)
\end{equation}
を取る。各ファイバーは $(\mathbf G_m)^{k-1}$ である。
周囲代数 $\mathcal B=\bigotimes_i\mathcal S_i$ と
三次イデアル $\ell$ による商 $\mathcal A=\mathcal B/(\ell)$ が存在する。
局所的な中心生成元で書けば $w=\sum_i a_i\Omega_i$, $a_i\in R^\times$ となる。

\begin{theorem}[満支持型の滑らかな表示 DM stack]
\label{mbs:thm:presentation-stack}
共通の grading scalar を除いた群を
\[
 \overline G_k=\GL_3^k/\mathbf G_m^{\mathrm{diag}},
 \qquad G_k=\overline G_k\rtimes\mathfrak S_k
\]
と置く。このとき
\begin{equation}\label{mbs:eq:stack}
 \mathfrak P_k=[\mathcal X_k/G_k]
\end{equation}
は有限型の滑らかな既約 Deligne--Mumford stack であり、
\begin{equation}\label{mbs:eq:dimension}
 \dim\mathfrak P_k=2k.
\end{equation}
これは、fpqc 局所的に上の満支持 tensor 商表示を持つ対象と、
全次数付き同型を共通の grading scalar で割った射のスタックである。
ブロックの順序は標識に含めない。

各幾何点の慣性群は
\begin{equation}\label{mbs:eq:inertia}
 \left(\prod_{i=1}^k\mathcal H_i\right)/\mu_3^{\mathrm{diag}},
 \qquad |\mathcal H_i|=27,
 \qquad |\operatorname{Inertia}|=3^{3k-1}
\end{equation}
であり、$\mathcal H_i$ は各因子の中心三次を固定する Heisenberg 群である。
とくに $k=2$ は次元四で、慣性群の位数は $243$ である。
$\mathcal A$ は次数別有限局所自由で、任意底変更に両立し、
Hilbert 級数は $(1-t^3)/(1-t)^{3k}$ となる。
\end{theorem}
\begin{proof}
$\mathcal X_k$ は滑らかな既約 $(11k-1)$ 次元スキームである。
共通 scalar は各 $\mathcal C_i$ に同じ三乗で作用するため、
その射影直線 $\ell$ を固定する。従って作用は $G_k$ を経由する。
$G_k$ は滑らかで次元 $9k-1$ なので、商は滑らかで表示した次元を持ち、既約である。

各ファイバーで $w$ の第一ブロックの最高記号は
$a_1\Omega_1$ であり、左右非零因子である。
命題\ref{tba:prop:low-centers}と同じ濾過の議論により $w$ は正則である。
固定次数の乗法
$\ell\otimes\mathcal B_{d-3}\to\mathcal B_d$ は全幾何ファイバーで単射なので、
単元となる最大小行列式の開集合で分裂単射となる。
従って余核は局所自由で、底変更と Hilbert 級数の主張を得る。

次に全ての次数付き同型を認識する。
定理\ref{tbp:thm:point-graph}の証明を各二ブロックに適用すると、
相対点対応は $\coprod_i\Gamma_{\sigma_i}$、点スキームは $\coprod_iE_i$ である。
一つの純ブロック上の可逆小行列式は、他の全ブロックの二組の座標を消す。
従ってこの等式は任意の試験基底でも成立する。
曲線の族は proper smooth で幾何学的に連結なファイバーを持つので、
その成分の相対スキームは
$\operatorname{Spec}(\pi_*\OO_{\coprod E_i})$、すなわち階数 $k$ の有限 \'etale スキームである。
全同型はこの成分スキームを保ち、\'etale 局所的にはブロックを置換する。
各曲線が張る平面は、次数一切断の制限核から底変更と両立して回復される。
このため、非被約な試験基底を含め、余分なブロック混合の同型は現れない。

ブロックの順序を合わせると、同型は各 $W_i$ の可逆線形写像からなり、
二次関係により各幾何的因子を同定する。
命題\ref{tba:prop:quadratic-envelope}の積の核による議論は局所自由な族でも同じであり、
さらに三次核 $\ell$ を保つことが商の同型の必要十分条件となる。
従って全次数付き同型の層は上の $\GL_3^k\rtimes\mathfrak S_k$ 作用で与えられ、
共通 scalar を除くと \eqref{mbs:eq:stack} の射となる。
非分裂表示はこの局所的記述の fpqc stack 化で扱う。
scalar を除いた商では、対象が大域的な次数一ベクトル束や実際の次数付き代数を
持つとは限らない。それらを選ぶには
$\GL_3^k\rtimes\mathfrak S_k$ への持上げが別に必要である。
ここでの分類はこの意味での射影的な次数付き表示である。

幾何点では $j$ が相異なるため安定化子に非自明なブロック置換はない。
定理\ref{tba:thm:generic-inertia}の証明を $k$ 因子で繰り返すと、
全ての中心指標 $\chi_i(g_i)$ が共通の値となることが必要十分である。
その共通三乗根による scalar を除けば
$\prod_i\mathcal H_i$ となり、残る共通核は $\mu_3^{\mathrm{diag}}$ である。
これで \eqref{mbs:eq:inertia} を得る。
標数零で各安定化子は有限 \'etale、従って商の対角は不分岐であり、
有限型のこの商 stack は Deligne--Mumford である。
\end{proof}

\begin{remark}[この表示スタックと内在的モジュライの間]
\label{mbs:rem:intrinsic-boundary}
上の結果は非空開上の次数付き表示の同型問題を解くが、
$\mathfrak P_k$ を全 $\OO$-marked moduli と同定するものではない。
相対 Serre 関手と切断飽和は定理\ref{mbsd:thm:serre}と系\ref{mbsd:cor:saturation}で、
全 Abel 同値の第三 Veronese までの認識は定理\ref{mbv:thm:all-o-third}で得る。
次数一の全ての根の分類、全同値関係の平坦性と対象の効果化は別途必要である。
また、全 intrinsic 成分の被覆、Poisson 射、Ku 比較の全 $\OO$-降下は結論しない。
Heisenberg 作用、Sklyanin 因子の幾何的分類、quotient stack の一般論は既知の入力であり、
ここではそれらと中心商・ブロック認識を組み合わせた表示を与えた。
この組合せの新規性について網羅的な文献確認は未完了である。
\end{remark}

% !TEX root = ../nc-cubic.tex
\section{満支持型の第三 Veronese と全構造対象保存同値}
\label{mbv:sec:veronese}

本節では、次数一の表示そのものではなく、第三 Veronese までを内在的に回復する。
第\ref{mbs:sec:presentations}節の満支持表示を用い、$k\geq2$、
$N=3k$, $n=N-2$, $r=N-3=3(k-1)$ とする。
以下の $C=A^{(3)}$ の次数は三で割った次数である。
次数一の中心がない $A$ に中心 Veronese 復元を直接適用することはできないが、
$C$ の中心一次部分の階数は $k-1$ となる。
この中心の計算を済ませてから、定理\ref{thm:nfm-central-veronese}を適用する。

\subsection{有限個の位数条件による Veronese 中心の計算}

まず基礎体を標数零の代数閉体とする。正整数 $d$ を固定し、
各因子の平行移動について
\begin{equation}\label{mbv:eq:order-bound}
 \sigma_i^j\ne\operatorname{id}
 \qquad(1\leq i\leq k,\ 1\leq j\leq3d)
\end{equation}
を仮定する。これは有限個の捩れ条件を除く非空開条件である。
$V=Z(A)_3=(\bigoplus_i\mathbf C\Omega_i)/\mathbf Cw$ と書く。

\begin{proposition}[第三 Veronese の有限次数中心]
\label{mbv:prop:center}
式\eqref{mbv:eq:order-bound}の下で、$\operatorname{Sym}V\to C$ は単射であり、
\begin{equation}\label{mbv:eq:center}
 Z(C)_e=\operatorname{Sym}^eV\qquad(0\leq e\leq d).
\end{equation}
特に $Z(A^{(3)})_1=Z(A)_3$ が成り立つ。
\end{proposition}
\begin{proof}
$T_i=S_i/(\Omega_i)=B(E_i,\mathcal L_i,\sigma_i)$ と置く。
補題\ref{tba:lem:tensor-center}の関数体による議論は、任意の正整数 $j$ について
$\sigma_i^j\ne1$ なら $Z(T_i)_j=0$ を与える。
中心元 $z\in(S_i)_j$ の像が零なら $z=\Omega_i z'$ であり、
$\Omega_i$ の正則性から $z'$ も中心的である。
従って $j\leq3d$ の中心は、$3\mid j$ なら $\mathbf C\Omega_i^{j/3}$、
そうでなければ零である。

$D=\bigotimes_{i=2}^kS_i$ と置き、第一ブロックの次数で濾過する。
命題\ref{tba:prop:low-centers}と同じ最高記号の計算で
\[
 \operatorname{gr}_{W_1}A=T_1\otimes D
\]
を得る。$D$ は濾過零部分として $A$ に埋め込まれる。
$z\in Z(C)_e$ の最高第一ブロック次数を $j$ とする。
$j>0$ なら、その最高記号は $T_1$ の純粋な三次部分の全元と可換する。
$D$ 側の係数を線型独立に書けば、非零な各係数
$u\in(T_1)_j$ が $(T_1)_3$ を中心化することになる。
捻った関数環内で $u=ft^j$、$a t^3,b t^3\in(T_1)_3$ と書き、
二つの可換式を割ると
$\sigma_1^j(a/b)=a/b$ を得る。
次数九の線束の切断比は $\mathbf C(E_1)$ を生成するので、
$\sigma_1^j=1$ となる。$j\leq3e\leq3d$ に反する。

従って $j=0$、すなわち $z\in D$ である。
$0\ne\bar u\in(T_1)_2$ の持上げ $u\in(S_1)_2$ と
$v\in D_1$ に対して $uv\in A_3=C_1$ なので、$u[z,v]=0$ である。
最高記号は $\bar u\otimes[z,v]$ だから $[z,v]=0$ となる。
よって $z\in Z(D)$ であり、テンソル積の中心と最初の次数別計算から
$z\in\mathbf C[\Omega_2,\ldots,\Omega_k]_{3e}$ を得る。
逆の包含は中心性による。
各 $S_i$ は正則な $\Omega_i$ に関して
$\mathbf C[\Omega_i]$ 上自由である
（$T_i$ の各次数の基底を持ち上げ、次数について帰納する）。
従って $D$ 内で $\Omega_2,\ldots,\Omega_k$ は代数的に独立であり、
$a_1\ne0$ による $\Omega_1$ の消去は $V$ とその張る多項式環を同定する。
\end{proof}

\subsection{中心平坦性と非被約底への変更}

$\mathcal X_{k,d}$ を、第\ref{mbs:sec:presentations}節の $\mathcal X_k$ から
式\eqref{mbv:eq:order-bound}の捩れ点を除いた開部分とする。
中心関係直線束を $\ell$ とし、
$\mathcal V=(\bigoplus_i\mathcal C_i)/\ell$ と置く。

\begin{lemma}[相対中心多項式環上の平坦性]
\label{mbv:lem:relative-center}
$\mathcal C=\mathcal A^{(3)}$ は
$\operatorname{Sym}\mathcal V$ 上平坦である。
任意の試験スキーム $T\to\mathcal X_{k,d}$ に対して
\[
 Z(\mathcal C_T)_e=\operatorname{Sym}^e\mathcal V_T
 \qquad(0\leq e\leq d)
\]
が成り立ち、これらの中心は次数成分の直和因子である。
\end{lemma}
\begin{proof}
各中心直線束を自明化した affine 開 $\operatorname{Spec}R$ 上で考える。
$\mathcal T_i=\mathcal S_i/(\Omega_i)$ の各次数は有限射影 $R$-加群である。
実際、中心元の乗法は全幾何ファイバーで単射なので分裂単射となり、
その余核が所定の Hilbert 関数を持つ有限局所自由加群となる。
各次数で余核への射の $R$-線型切断を選ぶ。
これらの持上げと $\Omega_i$ の冪を用いる次数帰納により、
\[
 \mathcal S_i\cong
 \bigoplus_{a\geq0}(\mathcal T_i)_a\otimes_RR[z_i],
 \qquad \deg z_i=3
\]
を中心多項式環上の次数付き加群として得る。
全ての加群を同時に自由化する開集合は要求していない。
各次数の射影性による切断を個別に選べば十分である。
テンソル積を取り、単元 $a_1$ を使って $\sum a_i z_i=0$ から $z_1$ を消去すると、
$\mathcal A$ は $R[z_2,\ldots,z_k]$ 上、有限射影 $R$-加群を係数とする直和となる。
第三 Veronese は、この係数の元の次数が三の倍数となる直和成分である。
従って $\mathcal C$ は $\operatorname{Sym}\mathcal V$ 上射影、特に平坦である。
この加群記述を底変更しても同じ結論が成り立つ。

中心は有限局所自由加群間の交換子行列
\[
 \mathcal C_e\longrightarrow
 \mathcal Hom(\mathcal C_1,\mathcal C_{e+1})
\]
の核である。$\mathcal C$ は次数一生成なので、この核が実際の中心を計算する。
命題\ref{mbv:prop:center}により、全幾何ファイバーで核の次元は
$\operatorname{Sym}^e\mathcal V$ の階数に等しい。
$\mathcal X_{k,d}$ は滑らかで被約なので、上位小行列式は恒等的に零となる。
最大非零小行列式で局所化すれば行列は一定階数の標準形となる。
従って核・像・余核は局所自由であり、中心核の形成は任意底変更に両立する。
中心に入る $\operatorname{Sym}^e\mathcal V$ は各幾何ファイバーでこの核と一致するので、
束としても一致する。非被約な $T$ に対する主張は、この分裂行列を引き戻して得る。
\end{proof}

\begin{theorem}[第三 Veronese までの一意な持上げ]
\label{mbv:thm:third-lift}
$k\geq3$, $d=k-1$ とし、$A,A'$ は同じ基底上の
$\mathcal X_{k,d}$ から引き戻された二つの表示とする。
このとき制限写像
\[
 \underline{\operatorname{Isom}}_{\rm gr}(A^{(3)},(A')^{(3)})
 \longrightarrow
 \underline{\operatorname{Isom}}_{\rm gr}(A^{(r)},(A')^{(r)})
\]
は fpqc 層として $\mu_{k-1}$-torsor である。
核は第三 Veronese の次数スカラーである。
$k=2$ では $r=3$ なので、この写像は恒等である。
\end{theorem}
\begin{proof}
$C=A^{(3)}$ は次数一生成で、中心一次束 $\mathcal V$ の階数は $k-1\geq2$ である。
補題\ref{mbv:lem:relative-center}は中心多項式環上の平坦性と、
任意底変更後の次数 $d$ の中心条件を与える。
従って $C,C'$ に定理\ref{thm:nfm-central-veronese}を適用できる。
$C^{(d)}=A^{(3d)}=A^{(r)}$ である。
\end{proof}

\subsection{全構造対象保存同値の第三 Veronese による認識}

以下では $d=k-1$ とし、第\ref{mbs:sec:presentations}節の表示を
$\mathcal X_{k,d}$ 上に制限する。
相対 Serre 関手の基底直線束を
\[
 \mathcal D=\det\!\left(\bigoplus_iW_i\right)\otimes\ell^\vee
\]
と書く。定理\ref{mbsd:thm:serre}で確認した式は
$\mathsf S=(-r)[n]\otimes\mathcal D$ である。
従って大域的に内在的なのは、まず
$\bigoplus_{q\geq0}A_{rq}\otimes\mathcal D^{-q}$ であり、
基底直線束を忘れて $A^{(r)}$ と同一視しない。

\begin{theorem}[全 $\mathcal O$-保存同値の有限次数表示への還元]
\label{mbv:thm:all-o-third}
上の範囲の二つの表示 $A,A'$ に対して、自然な fpqc 層の同型
\begin{equation}\label{mbv:eq:all-o-third}
 \frac{\underline{\operatorname{Isom}}_{\rm gr}(A^{(3)},(A')^{(3)})}
 {\mathbf G_m}
 \ \xrightarrow{\sim}\
 \operatorname{Eq}_{\mathcal O}
 (\operatorname{QGr}A,\operatorname{QGr}A')
\end{equation}
がある。右辺は線型完全同値と構造対象の同定の対を、
その同定を保つ自然同型で割ったものである。
これは自己同型を持たない groupoid、すなわち集合値 fpqc 層である。
左辺の $\mathbf G_m$ は第三 Veronese の次数スカラーである。
\end{theorem}
\begin{proof}
一般 Veronese の Gabriel 同値
（補題\ref{lem:nfm-gabriel}）により、左辺の次数付き同型は右辺の同値を与える。
次数スカラーは正規化された自然同型で吸収される。

逆に $\mathcal O$-保存同値 $F$ を取る。
線対象の compact 性、任意底変更に両立する切断飽和と相対 Serre 双対を用いる。
これらは定理\ref{mbsd:thm:serre}と系\ref{mbsd:cor:saturation}で、各因子の有限 Koszul 解消と正則中心商から得られる。
fpqc 局所で $\mathcal D,\mathcal D'$ の frame を選ぶと、
$H=\mathsf S^{-1}[n]$ は $(r)$ となる。
Serre 双対の表現対象の一意性から得た比較 $FH\simeq H'F$ を正負に反復し、
$F(\mathcal O)\simeq\mathcal O'$ と合わせる。
全整数の切断飽和により、これは合成を保つ次数付き同型
$A^{(r)}\simeq(A')^{(r)}$ を与える。
定理\ref{mbv:thm:third-lift}により、それはさらに fpqc 局所で
$C=A^{(3)}\simeq C'=(A')^{(3)}$ へ持ち上がる。
$k=2$ ではこの持上げ操作は不要である。
選んだ比較は全 $r$-orbit の切断と合成を同定しているので、
Gabriel 切断の完全忠実性により、その構造移送は元の $F$ と自然同型になる。

核も確認する。$C$ の次数付き自己同型 $\psi$ の構造移送が恒等なら、
$\mathcal O_C(j)$ 上の自然同型は各々一つの単元 $b_j$ である。
隣接 Hom は $C_1=A_3$ であり、正の階数の有限局所自由加群なので忠実である。
自然性から $\psi_1$ は $b_{j+1}/b_j$ によるスカラーとなり、
その比は $j$ に依存しない。
$C$ は次数一生成なので、$\psi$ 全体が次数スカラーとなる。
同様に恒等関手の自然自己準同型は、全 $\mathcal O_C(j)$ 上で同じ基底スカラーとなる。
Gabriel 切断の忠実性から全対象でもそうであり、$\mathcal O$ 上を固定すれば恒等である。

同値と自然同型の fpqc 降下は
命題\ref{prop:nfm-equivalence-descent}の係数作用付き Gabriel 商の議論に従う。
この議論は中心一次変数の存在に依存しない。
従って右辺は fpqc 層であり、上の局所全射性と核計算は
式\eqref{mbv:eq:all-o-third}を与える。
基底直線束の frame と Serre 比較の変更は局所的な構成にだけ用いており、
自然な左辺から右辺への写像自体にはそれらの選択は残らない。
\end{proof}

\subsection{第三 Veronese の二次表示と残る分類問題}

\begin{lemma}[有限な二次表示]
\label{mbv:lem:quadratic}
上の第三 Veronese は
\[
 C=T(A_3)\big/\bigl(\ker(A_3\otimes A_3\longrightarrow A_6)\bigr)
\]
という二次表示を持つ。関係束は直和因子であり、任意底変更に両立する。
\end{lemma}
\begin{proof}
$A$ は次数一で生成され、定義関係の次数は二と三だけである。
自由代数の長さ $3m$ の語を、長さ三の語に区切る。
一つの二次または三次関係を挿入した語は、区切りをまたいでも
高々二つの隣接した長さ三ブロックに収まる。
一ブロック内の関係は生成元空間を $A_3$ とすることで既に除かれ、
二ブロックにまたがる関係は表示した次数二の核から生成される。
関係イデアルの全次数成分はこのような挿入の和なので、他の高次関係は不要である。
$A_3\otimes A_3\to A_6$ は全射で、$A_6$ は有限局所自由だから分裂し、
その核の形成も任意底変更に両立する。
\end{proof}

従って式\eqref{mbv:eq:all-o-third}の左辺は、有限次元の生成元束の
射影線型同型がこの二次関係束を保存するという有限個の代数方程式で表される。
固定した二つの族の間の全 $\mathcal O$-同値層は、この意味で有限表示のスキームとして表現される。
二ブロックでは $\dim A_3=55$, $\dim A_6=406$ なので、二次関係の階数は $2619$ である。

この認識は次数一の $A$ を復元するものではない。
同じ $C$ を持つ Heisenberg twist は別の次数一表示を与える。
可換 tensor 表示の幾何点に限れば、中心商の余法加群によって
元の次数付き同型類を回復できる（系\ref{mbiq:cor:commuting-o-rigidity}）。
全 $\mathcal O$-局所型の非空開は定理\ref{mbg:thm:generic-stack}で構成する。
全ての第三根表示の分類、$C$ の射影的自己同型群の具体的決定、
全境界、および全 intrinsic 変形モジュライの成分性は別の課題である。
有限表示の Isom 層を得たことだけから、元の次元 $2k$ の表示スタックを
全 $\mathcal O$-marked モジュライと同定しない。

% !TEX root = ../nc-cubic.tex
\section{満支持型の次数一表示の根と Zhang twist}
\label{mbr:sec:roots}

本節では、同じ第三 Veronese と同じ $\OO$-marked 圏を持つ表示を、
まず一つの次数付き自己同型から作る twist について分類する。
さらに、両端が可換 tensor 積表示である場合には、
一般の twisting system も新しい次数付き同型類を生まないことを示す。
ここでの直接の証明は、全 shift を対象として保つ同値に適用される。
可換 tensor チャート内の任意の全 $\OO$-保存同値についても、
定理\ref{mbca:thm:full-inertia}と系\ref{mbiq:cor:commuting-o-rigidity}により
同値関手そのものの次数付き認識が成立する。
また系\ref{mbca:cor:marked-roots}は、構造対象を $\OO(1)$ へ送る根について、
本節の共役 groupoid が全ての根を尽くすことを示す。

基礎体 $k$ は標数零の代数閉体とする。
ブロック数を $\ell\ge2$ とし、
\begin{equation}\label{mbr:eq:setup}
 B=\bigotimes_{i=1}^{\ell}S_i,\qquad
 S_i=S(E_i,\mathcal L_i,t_{p_i}),\qquad
 A=B/(w),\quad w=\sum_{i=1}^{\ell}a_i\Omega_i,\quad a_i\ne0
\end{equation}
を固定する。前節までの楕円 Sklyanin の仮定に加えて、
\begin{equation}\label{mbr:eq:good-open}
 6p_i\ne0,\qquad
 j(E_i)\notin\{0,1728\},\qquad
 j(E_i)\ne j(E_j)\quad(i\ne j)
\end{equation}
を課す。これは非空な Zariski 開条件である。
$6p_i\ne0$ は、後で三捩れ点を足しても符号反転による同定が生じないように使う。
$\mathcal H_i$ を $\Omega_i$ を固定する位数 $27$ の Heisenberg 群とし、
\begin{equation}\label{mbr:eq:finite-group}
 \mathcal H=\prod_i\mathcal H_i,\qquad
 \mathcal K_\ell=\mathcal H/\mu_3^{\mathrm{diag}}
\end{equation}
と置く。系\ref{tba:cor:k-block-center-inertia}により
$\Aut_{\mathrm{gr}}(A)/\mathbf G_m=\mathcal K_\ell$ である。
$\mathcal H_i$ は指数三、すなわち全ての元の三乗が一である。
これは Hesse 座標の二生成元とその中心交換子の表示からも直接従う。

Zhang twist の基本性質と一般の twisting system の規格化は既知である
\cite[Proposition~2.4, Theorems~3.1 and~3.7]{TBPZhangTwist}。
Cooney--Grabowski はこの方法を点幾何による groupoid として研究した
\cite[Proposition~10, Definition~11 and Theorem~28]{MBRCooneyGrabowski}。
ただし本稿の周囲代数では、系\ref{tbp:cor:all-blocks}により
$R_B\ne H^0(\mathcal I_{\Gamma_B}(1,1))$ である。
従って同論文の geometric algebra の仮定は満たされず、
その分類定理を本節へ直接適用することはできない。
以下では混合二次関係と中心三次関係も使って分類する。

\subsection{三捩れ平行移動を伴う代数的 twist}

$\tau=(\tau_i)\in\mathcal H$ に対し、斉次元の積を
\begin{equation}\label{mbr:eq:twist-product}
 u*_{\tau}v=u\tau^{\deg u}(v)
\end{equation}
とする代数を $B^\tau$、その中心商を $A^\tau$ と書く。
本節で\emph{代数的 twist} というときは、
一つの次数付き自己同型の冪から来る twisting system を意味し、
一般の twisting system と区別する。

\begin{lemma}[三捩れ twist の中心商と点 shift]
\label{mbr:lem:torsion-twist}
$\tau\in\mathcal H$ に対して $w$ は $B^\tau$ でも正則な中心元で、
\begin{equation}\label{mbr:eq:common-veronese}
 A^\tau=B^\tau/(w),\qquad
 (A^\tau)^{(3)}=A^{(3)},\qquad
 (A^\tau)^{(3\ell-3)}=A^{(3\ell-3)}
\end{equation}
が同じ次数付きベクトル空間の上で成立する。
これらは全て同じ $\OO$-marked 圏を与える。

$\tau_i^*$ が $E_i$ 上で与える三捩れ平行移動を $t_{q_i}$ とする。
ここで $\tau_i^*$ は点の座標への双対線型写像であり、
$q_i$ の符号はこの規約で固定する。
$A^\tau$ の点スキームは $\coprod_iE_i$ で、各成分上の shift は
\begin{equation}\label{mbr:eq:point-shift}
 \sigma_i^\tau=t_{p_i+q_i},\qquad q_i\in E_i[3]
\end{equation}
となる。
\end{lemma}
\begin{proof}
$\tau^3=1$ かつ $\tau(w)=w$ なので、
$w*_{\tau}u=wu$、$u*_{\tau}w=uw$ である。
これで中心性・正則性と商の等式を得る。
次数が三の倍数の元から始まる積は元の積だから、Veronese の等式も従う。
右加群の作用を $m*_{\tau}u=m\tau^{\deg m}(u)$ と変える関手は
torsion と正則右加群を保つので、$\OO$-marked 圏の同値を与える。

$R$ を $B$ の二次関係空間とすると
\begin{equation}\label{mbr:eq:twisted-relations}
 R^\tau=(1\otimes\tau^{-1})R.
\end{equation}
この関係は $B^\tau$ を二次代数として表示する。
実際、各 tensor 次数で
$v_1\otimes\cdots\otimes v_m\mapsto
v_1\otimes\tau(v_2)\otimes\cdots\otimes\tau^{m-1}(v_m)$
を使うと、twist 後の関係イデアルと元の関係イデアルが対応する。
従って双線型関係スキームは
$\Gamma_{B^\tau}=(1\times\tau^*)\Gamma_B$ であり、
系\ref{tbp:cor:all-blocks}から \eqref{mbr:eq:point-shift} を得る。
各 $\Omega_i$ は普遍点列上で零であり、$\tau$ に固定されるので、
中心商を課しても点スキームは変わらない。
\end{proof}

\begin{theorem}[代数的 twist の共役 groupoid]
\label{mbr:thm:algebraic-twist-groupoid}
$A$ の全ての代数的 Zhang twist を対象とし、
射を次数付き同型の共通 grading scalar による商とする groupoid を
$\mathfrak T_{\mathrm{alg}}(A)$ と書く。
この groupoid は有限群 $\mathcal K_\ell$ の共役作用 groupoid
\begin{equation}\label{mbr:eq:conjugacy-groupoid}
 \mathfrak T_{\mathrm{alg}}(A)\simeq
 [\mathcal K_\ell/\mathcal K_\ell]_{\mathrm{conj}}
\end{equation}
に同値である。
具体的に、$\tau,\upsilon\in\mathcal H$ に対し
\begin{equation}\label{mbr:eq:twist-isom}
 \operatorname{Isom}_{\mathrm{gr}}(A^\tau,A^\upsilon)/k^\times
 \simeq
 \{h\in\mathcal K_\ell:
            h\overline\tau h^{-1}=\overline\upsilon\}.
\end{equation}
従って対象の同型類は $\mathcal K_\ell$ の共役類であり、
$A^\tau$ の射影的な次数付き自己同型群は
$C_{\mathcal K_\ell}(\overline\tau)$ である。
ここでは固定した代数閉体上の対象と射の groupoid を述べており、
任意基底上の twist の族を分類する相対 stack の同値は主張しない。
\end{theorem}

\begin{proof}
まず任意の $f\in\Aut_{\mathrm{gr}}(A)$ を考える。
二次包絡とブロックの認識により、$f$ は各因子の自己同型 $f_i$ からなり、
その中心指標 $\chi_i(f_i)$ は共通の値 $c$ となる。
$t^3=c$ を選び、共通 grading scalar $t$ を除けば
$\tau=t^{-1}f\in\mathcal H$ となる。
共通 scalar による twist は、次数 $m$ の線型写像を
$t^{-\binom m2}$ 倍することで取り除けるので、
$A^f\simeq A^\tau$ である。従って $\mathcal H$ の元による twist を分類すればよい。

次数付き同型 $F:A^\tau\to A^\upsilon$ を取り、その次数一部分を $h$ とする。
点成分の線型包と $j$ の相異性から $h=\bigoplus_i h_i$ である。
$h_i$ が曲線に誘導する同型は、原点を flex に選べば
$t_{u_i}\epsilon_i$、$u_i\in E_i[3]$、$\epsilon_i\in\{1,-1\}$ の形である。
点 shift の intertwining 条件と \eqref{mbr:eq:point-shift} により、
もし $\epsilon_i=-1$ なら
$2p_i=-(q_i+q'_i)$ となる。ただし $q_i,q'_i\in E_i[3]$ は
$\tau_i,\upsilon_i$ に対応する点である。
これは $6p_i=0$ を意味するので除外される。
従って $\epsilon_i=1$ かつ $q_i=q'_i$ であり、
$h_i$ は元の $t_{p_i}$ とも可換な射影自己同型を与える。
その線型持上げ $h_i$ は元の因子 $S_i$ の次数付き自己同型である。
この段階では各 $\chi_i(h_i)$ の一致はまだ使わない。

異なる二ブロック $i,j$ の twist 後の関係は
\begin{equation}\label{mbr:eq:mixed-graph}
 \{x\otimes\tau_j^{-1}y-y\otimes\tau_i^{-1}x:
                         x\in W_i,\ y\in W_j\}
\end{equation}
である。この空間を
$(W_i\otimes W_j)\oplus(W_j\otimes W_i)$ の第一成分上のグラフと見て、
$h\otimes h$ による移送を比較すると
\[
 (h_j\tau_jh_j^{-1})\otimes(h_i\tau_i^{-1}h_i^{-1})
       =\upsilon_j\otimes\upsilon_i^{-1}
\]
を得る。
二つの非零な単純 tensor が等しければ、それぞれの因子は互いに逆の
スカラー倍である。全てのブロック対についてこれを使うと、
一つの共通スカラー $c$ が存在して
\begin{equation}\label{mbr:eq:conjugate-scalar}
 h\tau h^{-1}=\lambda_c\upsilon
\end{equation}
となる。$\lambda_c$ は次数 $m$ に $c^m$ を掛ける grading scalar である。
$\tau_i,\upsilon_i$ は各中心元を固定するので $c^3=1$ である。

逆に \eqref{mbr:eq:conjugate-scalar} を満たす因子自己同型 $h$ は
\begin{equation}\label{mbr:eq:isom-formula}
 F|_{B_m}=c^{-\binom m2}h|_{B_m}
\end{equation}
によって $B^\tau\to B^\upsilon$ の次数付き同型を与える。
実際、次数 $m,n$ の元について
$h(u*_{\tau}v)=c^{mn}h(u)*_{\upsilon}h(v)$ であり、
$\binom{m+n}{2}=\binom m2+\binom n2+mn$ が積の両立性を与える。
二次包絡の普遍性により、与えられた次数一写像を延ばす同型はこれに限る。
次数三では $c^{-\binom32}=c^{-3}=1$ なので、
$F(\Omega_i)=\chi_i(h_i)\Omega_i$ である。
従って $F$ が中心商へ降りる必要十分条件は
全ての $\chi_i(h_i)$ が等しいことであり、
これは $h\in\Aut_{\mathrm{gr}}(A)$ と同値である。

共通 grading scalar を除くと $h$ は $\mathcal K_\ell$ の元となり、
\eqref{mbr:eq:conjugate-scalar} はそこでの共役条件に等しい。
これで \eqref{mbr:eq:twist-isom} を得る。
$\mathcal H$ の二つの代表が共通 $\mu_3$ だけ異なる場合も
\eqref{mbr:eq:isom-formula}が次数一で恒等な同型を与える。
各 $\mathcal K_\ell$ の元の代表を一つ選べば、
射の対応は次数一の線型写像の合成を保つので groupoid の同値となる。
\end{proof}

\subsection{同型類の数とブロックスカラーの部分群}

\begin{corollary}[全代数的 twist の個数]
\label{mbr:cor:twist-count}
定理\ref{mbr:thm:algebraic-twist-groupoid}の仮定の下で、
代数的 Zhang twist の次数付き同型類の個数は
\begin{equation}\label{mbr:eq:count}
 c_\ell=\frac{11^\ell+2\cdot8^\ell}{3}.
\end{equation}
従って二ブロックでは $83$、三ブロックでは $785$ である。
全ての同型類は、補題\ref{mbr:lem:torsion-twist}の意味で
同じ第三 Veronese と同じ $\OO$-marked 圏を持つ表示によって代表される。

共役類を $q=(q_i)\in\prod_iE_i[3]$ ごとに分け、
$s=|\{i:q_i\ne0\}|$ とすると、その $q$ の上の同型類数と自己同型群の位数は
\begin{equation}\label{mbr:eq:support-count}
 \#\{\text{同型類 over }q\}
       =3^{\max(\ell-1-s,0)},\qquad
 |\Aut_{\mathrm{gr}}(A^\tau)/\mathbf G_m|
       =3^{3\ell-1-\min(s,\ell-1)}.
\end{equation}
\end{corollary}

\begin{proof}
$\mathcal K_\ell$ の中心とその商は
\[
 Z=\mu_3^\ell/\mu_3^{\mathrm{diag}},\qquad
 Q=\prod_i E_i[3]
\]
である。各因子の Heisenberg 交換子は非退化な三捩れ pairing であり、
$q_i\ne0$ なら、その因子との交換子は $\mu_3$ 全体を動く。
$q_i=0$ なら、その成分の交換子は零である。
従って support が $I$ である元の交換子像は
\[
 \mathbf F_3^I\longrightarrow
 \mathbf F_3^\ell/\mathbf F_3(1,\ldots,1)
\]
の像である。その次元は $s<\ell$ なら $s$、$s=\ell$ なら $\ell-1$ である。
なお $q\ne0$ のときこの像は非零なので、上記の $Z$ が中心全体である。

$q$ の上には $3^{\ell-1}$ 個の元があり、
各共役類の大きさは $3^{\min(s,\ell-1)}$ である。
これで \eqref{mbr:eq:support-count} を得る。
support の大きさ $s$ の $q$ は $\binom\ell s8^s$ 個あるので、総数は
\[
 \sum_{s=0}^{\ell-1}\binom\ell s8^s3^{\ell-1-s}+8^\ell
   =\frac{11^\ell+2\cdot8^\ell}{3}
\]
となる。
\end{proof}

\begin{corollary}[ブロックスカラー twist の完全な分類]
\label{mbr:cor:scalar-roots}
$\tau|_{W_i}=\varepsilon_i\operatorname{id}$、$\varepsilon_i\in\mu_3$ とする。
二つのこのような twist が次数付き同型となる必要十分条件は、
その比が全ブロックで同じスカラーとなることである。
従ってこの部分の同型類は正確に
$\mu_3^\ell/\mu_3^{\mathrm{diag}}$、個数 $3^{\ell-1}$ である。
各対象の射影的な次数付き自己同型群は $\mathcal K_\ell$ である。
\end{corollary}
\begin{proof}
これらの twist は $\mathcal K_\ell$ の中心の元に対応し、
その共役類は一点である。
定理\ref{mbr:thm:algebraic-twist-groupoid}を適用する。
\end{proof}

\begin{remark}[有限検算と証明の区別]\label{mbr:rem:verification}
\texttt{computations/verify\_multiblock\_twists.py} は
$\mathbf F_3$ 上の上三角 Heisenberg 群の直積を共通中心で割った群を、
明示的な積と逆元によって実装する。
$\ell=2,3$ について全群要素を列挙し、生成元による共役軌道を求めて
$83,785$ 個の共役類と \eqref{mbr:eq:support-count} を厳密に検算する。
この検算は群の共役類を確認するものであり、
代数の同型認識、圏同値、一般 $\ell$ の公式の証明を代用しない。
\end{remark}

\subsection{可換 tensor 積表示どうしの twist 剛性}

\begin{theorem}[可換 tensor 積チャート内の一般 twisting system]
\label{mbr:thm:commuting-twist-rigidity}
$A,A'$ を \eqref{mbr:eq:setup}--\eqref{mbr:eq:good-open} の
可換 tensor 積中心商とする。
両者のブロック数は同じとし、因子や中心係数は異なってよい。
一般の Zhang twisting system $\boldsymbol\tau$ により
$A'\simeq A^{\boldsymbol\tau}$ であるならば、
$A'$ と $A$ は次数付き同型である。
従ってこのチャートの二つの対象について
\begin{equation}\label{mbr:eq:graded-category-rigidity}
 \operatorname{Gr}A\simeq\operatorname{Gr}A'
 \quad\Longleftrightarrow\quad
 A\simeq_{\mathrm{gr}} A'.
\end{equation}
この定理の直接の結論は同値な\emph{対象}の認識である。
$\operatorname{Gr}$ 圏の任意の同値には次数零生成対象の保存を仮定していない。
全 $\OO$-保存 tails 同値の関手そのものの認識は、
系\ref{mbiq:cor:commuting-o-rigidity}で別に得られる。
\end{theorem}

\begin{proof}
Zhang の規格化 \cite[Proposition~2.4]{TBPZhangTwist}により
$\tau_0=1$ としてよい。$V=A_1$ とし、
$g=\tau_1|_V\in\GL(V)$ と置く。
$A'$ をこの次数一空間へ同定すると、その二次関係は
\[
 R'=(1\otimes g^{-1})R_B
\]
であり、二次関係スキームは
$(1\times g^*)\Gamma_B$ である。
このスキームの第一射影の像は $E=\coprod_iE_i$、第二射影の像は $g^*E$ である。
一方、$A'$ も可換 tensor 積表示なので、その二次関係スキームは
その点スキームの自己同型のグラフであり、両射影の像は一致する。
従って $g^*E=E$ となる。
$j(E_i)$ が相異なるため $g^*$ は各 $E_i$ とその線型包を保ち、
$g=\bigoplus_i g_i$ である。
また第一射影により認識される $A'$ のブロック分解も、この同じ分解である。

各混合ブロックの関係は
$x\otimes g_j^{-1}y-y\otimes g_i^{-1}x$ である。
これが可換 tensor 積の関係空間
$\langle x\otimes y-y\otimes x\rangle$ に等しいので、
二つの混合直和成分間のグラフを比較して
\[
 g_j\otimes g_i^{-1}=1_{W_j}\otimes1_{W_i}
 \qquad(i\ne j)
\]
を得る。従って全ての $g_i$ は同じスカラーである。
ゆえに $R'=R_B$、すなわち両者の二次包絡は同じ $B$ である。

$A'$ の三次核は $B$ の中心三次直線であり、
全非零係数という仮定から $\sum_i a_i'\Omega_i$ と書ける。
系\ref{tba:cor:k-block-center-inertia}のブロック別尺度作用により、
代数閉体上ではこの直線を $kw$ へ移せる。
従って $A'\simeq_{\mathrm{gr}}A$ を得る。
最後に、連結次数付き代数の graded module 圏の同値と twist の存在の同値は
Zhang の定理 \cite[Theorem~3.5]{TBPZhangTwist}による。
\end{proof}

\begin{corollary}[全 shift を対象として保つ同値への帰結]
\label{mbr:cor:sss-recognition}
前定理の $A,A'$ に対し、自然な切断写像が
\[
 A_j\simeq\Hom(\OO,\OO(j)),\qquad
 A'_j\simeq\Hom(\OO',\OO'(j))\qquad(j\in\mathbf Z)
\]
を与えると仮定する。
同値 $F:\operatorname{QGr}A\to\operatorname{QGr}A'$ が
\begin{equation}\label{mbr:eq:sss}
 F(\OO(j))\simeq\OO'(j)\qquad(j\in\mathbf Z)
\end{equation}
を満たすならば、$A\simeq_{\mathrm{gr}}A'$ である。
さらに定理\ref{mbva:thm:sss}により、$F$ 自体が次数付き同型の構造移送と自然同型である。
\end{corollary}
\begin{proof}
全ての $\OO(j)$ 上の同型を選び、$F$ をそれらの間の Hom と合成に適用する。
これは二つの section $\mathbf Z$-代数を同定する。
Zhang の定理 \cite[Theorems~3.3 and~3.7(2)]{TBPZhangTwist}により、
二つの section 代数は一般の twisting system によって結ばれる。
切断飽和性でそれらを $A,A'$ と同定し、
定理\ref{mbr:thm:commuting-twist-rigidity}を適用する。
関手そのものについては、定理\ref{mbva:thm:sss}の全整数の frame 調整、
混合二次関係と三次関係の認識、および Gabriel 生成の議論を適用する。
\end{proof}

本稿の $\C$ 上の満支持表示族では、上の切断飽和性は
系\ref{mbsd:cor:saturation}によって任意の通常の底変更後にも満たされる。

\begin{remark}[全自己同値の認識と未指定の第三根]
\label{mbr:rem:remaining-root-problem}
条件 \eqref{mbr:eq:sss} は、各 shift を\emph{対象として}保つ条件であり、
定義の段階で関手と次数一 shift の可換同型を指定するものではない。
本稿の可換 tensor チャートでは、定理\ref{mbva:thm:sss}により
この条件から関手そのものの次数付き持上げを得る。
さらに、定理\ref{mbv:thm:all-o-third}の有限位数条件を満たす範囲では、
定理\ref{mbca:thm:full-inertia}が全 $\OO$-自己同値群を $\mathcal K_\ell$ と同定する。
従って系\ref{mbiq:cor:commuting-o-rigidity}と合わせると、
可換 tensor チャートどうしでは $F(\OO)\simeq\OO'$ だけから
与えられた関手の次数付き持上げと \eqref{mbr:eq:sss} が従う。
任意の非被約底を含む相対同型層についても、
系\ref{mbrr:cor:stack-isomorphism}がさらに選んだ非空開上で認識を与える。

この結論は、同じ第三 Veronese を持つ任意の次数一表示が
可換 tensor 積表示であることを意味しない。
本節は代数的 Zhang twist の全体と、一般 twist の可換チャート内の剛性を決定した。
根の関手としては、系\ref{mbca:cor:marked-roots}により
$\mathsf r\OO\simeq\OO(1)$ を満たす全ての $\mathsf r$ が
$\mathsf s\Phi_h$、$h\in\mathcal K_\ell$ の形であり、
自動的に $\mathsf r^3\simeq\mathsf s^3$ を満たす。
従ってこの条件付きの全根の共役 groupoid も
$[\mathcal K_\ell/\mathcal K_\ell]_{\mathrm{conj}}$ である。
二ブロックでは根の自然同型類が $243$、共役類が $83$ であり、
一般の共役類数は式\eqref{mbr:eq:count}の $c_\ell$ である。
三乗同型 $\mathsf r^3\simeq\mathsf s^3$ の選択や、その coherence は追加標識に含めない。

$\mathsf r^3\simeq\mathsf s^3$ だけを仮定し、$\mathsf r\OO$ を指定しない
全ての第三根の分類、および一般の非代数的 twisting system の全分類は未決定である。
$c_\ell$ は全 intrinsic moduli の成分数や、像を指定しない全ての根の個数ではない。
Zhang の同値、Heisenberg 表現、既存の一因子幾何は既知であり、
本節で組み合わせた分類の新規性の網羅的確認は未完了である。
\end{remark}

% !TEX root = ../nc-cubic.tex
\section{第三 Veronese の点と余法線による因子の回復}
\label{mbiq:sec:quotients}

本節は代数閉体上の幾何学的対象を扱う。
第三 Veronese の中心一次部分で割った代数から、まず純ブロック商を内在的に認識する。
その点対応が直接与えるのは三乗の平行移動と次数九の偏極である。
さらに商の余法線双加群を用いると、元の次数一 shift そのものを回復できる。
この結果は、可換 tensor 積表示どうしの全 $\OO$-同値に関する対象剛性を与える。
さらに定理\ref{mbca:thm:full-inertia}と組み合わせると、
与えられた全同値関手そのものも次数付き同型として認識できる。
非被約底上の同型層については、系\ref{mbrr:cor:stack-isomorphism}が
別の相対議論により非空開上で認識を与える。
点加群関手の非被約方向の分類をこれらから逆に結論するものではない。

\subsection{中心商と純三次部分}

基礎体を標数零の代数閉体 $K$、ブロック数を $k\geq2$ とする。
各 $i$ に対し、滑らかな楕円曲線 $E_i$、次数三の線束 $L_i$、
平行移動 $\sigma_i=t_{p_i}$ と
\begin{equation}\label{mbiq:eq:order}
 p_i\ne0,\qquad 2p_i\ne0,\qquad 3p_i\ne0
\end{equation}
を与える。
$S_i=S(E_i,L_i,\sigma_i)$ を既存の二次 Sklyanin 代数、
$\Omega_i$ をその正則中心三次元とし、
\[
 T_i=S_i/(\Omega_i)=B(E_i,L_i,\sigma_i),\qquad W_i=(T_i)_1
\]
と置く。中心係数が全て非零な
$A=(\bigotimes_iS_i)/(\sum_i a_i\Omega_i)$ に対して、
\begin{equation}\label{mbiq:eq:central-quotient}
 C=A^{(3)},\qquad
 D=C/(Z(C)_1)\simeq
       \left(\bigotimes_iT_i\right)^{(3)}
\end{equation}
である。
ここでは命題\ref{mbv:prop:center}の $d=1$ を用いた。
括弧 $(Z(C)_1)$ は中心一次空間が生成するイデアルを表す。
有限位数の場合に現れ得る高次数中心まで全て零にする操作ではない。
実際、三次中心の全像を $A$ で零にすると各 $\Omega_i$ が零になり、
この中心商と第三 Veronese は可換する。

以下では $D$ 自体の表示だけを使う。
\[
 U_i=T_i^{(3)},\qquad
 F=\bigotimes_iU_i\subset D,
 \qquad L_{i,3}=L_i\otimes\sigma_i^*L_i\otimes\sigma_i^{2*}L_i
\]
と置く。$U_i=B(E_i,L_{i,3},\sigma_i^3)$ であり、$\deg L_{i,3}=9$ である。
全ての Veronese の次数は三で割る。
多重次数 $\boldsymbol a=(a_1,\ldots,a_k)$ には
$T_{\boldsymbol a}=\bigotimes_i(T_i)_{a_i}$ と書く。

各 $T_i$ と $U_i$ は次数一生成である。
また THCR の標準同値と本稿の右加群規約では
\begin{equation}\label{mbiq:eq:thcr-shift}
 \qgr T_i\simeq\operatorname{coh}E_i,\qquad
 (1)\longleftrightarrow s_i=\sigma_{i,*}\circ(-\otimes L_i),
 \qquad s_i^{-1}=(-\otimes L_i^{-1})\circ\sigma_i^*.
\end{equation}
曲線上の正次数線束は任意の自己同型について ample sequence を与えるので、
ここに無限位数の仮定は要らない
\cite[\S3, equations~(3.1), (3.3), Theorem~3.2]{NFMStaffordVdB}。
次数一生成は、次数が少なくとも三の楕円線束の切断乗法の全射性からも従う。

\begin{lemma}[純三次部分上の有限生成性]
\label{mbiq:lem:finite-pure}
$D$ は $F$ 上、左右とも有限生成な次数付き加群である。
生成元の次数は $\lfloor2k/3\rfloor$ 以下に取れる。
\end{lemma}
\begin{proof}
$|\boldsymbol a|=3n$ を満たす多重次数を
$a_i=3q_i+e_i$, $0\leq e_i\leq2$ と書く。
各因子が次数一生成なので、積は
\[
 \left(\bigotimes_i(T_i)_{e_i}\right)
 \otimes\left(\bigotimes_i(T_i)_{3q_i}\right)
 \longrightarrow T_{\boldsymbol a}
\]
で全射であり、左右の順序を逆にしても同じである。
残余多重次数 $\boldsymbol e$ は有限個で、
$|\boldsymbol e|$ は三の倍数、$|\boldsymbol e|\leq2k$ である。
対応する有限次元空間の基底が求める生成元となる。
\end{proof}

\subsection{全点加群の純ブロック性}

点加群とは、次数零で生成され、各非負次数が一次元の次数付き右加群をいう。
THCR $B(E,L,\sigma)$ の点加群は $E$ の点により一意に決まる。
ここで使う範囲を確認しておく。
点加群の非零部分加群は tail 全体を含むので、その tails 像は単純対象である。
THCR 同値の下でこれは skyscraper sheaf であり、切断加群の各次数は一次元である。
点加群からその飽和切断加群への標準写像は torsion-free 性により単射なので、
非負次数で同型となる。
従って点加群はその点から来る標準切断加群であり、shift は $\sigma$ に対応する。
この議論にも平行移動の無限位数性は不要である。

\begin{theorem}[中心商の点加群]
\label{mbiq:thm:point-locus}
$D$ の全ての点加群は、ただ一つの純ブロック商
\begin{equation}\label{mbiq:eq:pure-quotient}
 q_i:D\longrightarrow U_i
\end{equation}
を通じて得られる $U_i$ の点加群である。
従って、第一の作用座標がなす被約閉集合と、その二項の点対応は
\begin{equation}\label{mbiq:eq:reduced-point-locus}
 \coprod_{i=1}^k(E_i,|L_{i,3}|)
       \subset\PP(D_1^*),\qquad
 \coprod_{i=1}^k\Gamma_{\sigma_i^3}
\end{equation}
である。
この主張は全ての幾何学的点加群を分類する。
点加群関手の非被約方向や、有限長の切断点スキームが既にこのグラフと等しいことは主張しない。
\end{theorem}
\begin{proof}
$M$ を $D$ の点加群とする。
補題\ref{mbiq:lem:finite-pure}により、$M$ は $F$ 上も有限生成である。
従ってある $b$ 以上で $M_nF_1=M_{n+1}$ となり、
$M_{\geq b}(b)$ は $F$ の点加群となる。
その次数 $n$ における各純部分 $(U_i)_1$ の作用を
$\alpha_{i,n}\in(U_i)_1^*$ で表し、非零となる添字集合を $J_n$ とする。
$J_n$ は空ではない。
異なる二ブロック間の交換関係は
\[
 \alpha_{i,n}(u)\alpha_{j,n+1}(v)
 =\alpha_{j,n}(v)\alpha_{i,n+1}(u)
 \qquad(i\ne j)
\]
を与える。もし $i\in J_n\setminus J_{n+1}$ なら
$j\in J_{n+1}$ を選んで矛盾を得る。逆向きも同じなので $J_n=J_{n+1}$ である。

この集合が二つ以上の添字を含むと、非零な単純 tensor の等式により、
各 $[\alpha_{i,n}]\in\PP((U_i)_1^*)$ は $n$ によらなくなる。
一方、$U_i$ への制限は点加群なので、その点列は $\sigma_i^3$ で移る。
これは $\sigma_i^3$ の固定点を与え、$3p_i\ne0$ に反する。
従って tail 上で active な純ブロックはただ一つ、$i$ である。

混合部分の作用も除く。
$|\boldsymbol e|=3$ に対し $Q_{\boldsymbol e}=T_{\boldsymbol e}\subset D_1$ と置く。
各 $T_i$ の次数一生成性により、$D_2$ の部分空間として
\begin{equation}\label{mbiq:eq:mixed-propagation}
 (U_i)_1Q_{\boldsymbol e}=Q_{\boldsymbol e}(U_i)_1
\end{equation}
が成り立つ。
tail のある次数で $Q_{\boldsymbol e}$ の作用が非零なら、
次の次数での $(U_i)_1$ の作用と合成して非零な積を取れる。
\eqref{mbiq:eq:mixed-propagation}を用いると、
その次の次数でも $Q_{\boldsymbol e}$ の作用が非零でなければならない。
従って非零性は全ての高次数へ伝播する。

$\boldsymbol e\ne3\boldsymbol e_i$ なら inactive なブロック $j$ に $e_j>0$ がある。
連続する三次数で $Q_{\boldsymbol e}$ の非零作用を選んで合成すると、
点加群上の非零作用となる。
しかしこの積の多重次数は $3\boldsymbol e$ であり、$F$ に属する。
さらに inactive な $U_j$ の正次数を含むので、tail 上の作用は零である。
これは矛盾である。よって全ての混合部分と inactive な純部分は tail 上で零に作用する。

最後に初期次数へ戻る。
ある次数以上で $Q_{\boldsymbol e}$ の作用が零、$(U_i)_1$ の作用が非零であるとする。
直前の次数で $Q_{\boldsymbol e}$ の作用が非零なら、
\eqref{mbiq:eq:mixed-propagation}によって次の次数でも非零になり、矛盾する。
従って直前の次数でも全ての $\boldsymbol e\ne3\boldsymbol e_i$ が零に作用する。
$D_1$ は次の一次元成分を生成するので、残る $(U_i)_1$ の作用は非零である。
逆帰納により全次数で純一ブロックの点加群となる。

逆に $q_i$ を通じた $U_i$ の点加群は $D$ の点加群である。
$U_i$ の標準的な点加群の記述から、第一座標は $|L_{i,3}|$ による埋込み、
shift は $\sigma_i^3$ となる。
これらの曲線とグラフは既に閉で被約なので、表示した被約閉集合が得られる。
\end{proof}

\begin{corollary}[純ブロック商の内在的認識]
\label{mbiq:cor:intrinsic-quotients}
$I_i=\ker(q_i)$ とする。
次数付き代数 $D$ は、順序を除いてイデアル $I_i$ と商 $D/I_i=U_i$ を決定する。
具体的には、\eqref{mbiq:eq:reduced-point-locus}の各既約成分の線型包の annihilator が
$(I_i)_1$ であり、$I_i=((I_i)_1)$ である。
従ってこの形の代数間の次数付き同型は純ブロック商を置換し、
$(E_i,L_{i,3},\sigma_i^3)$ を対応する triple と同定する。
\end{corollary}
\begin{proof}
各 $|L_{i,3}|$ は完全線型系であり、その像は $\PP((U_i)_1^*)$ を張る。
従ってその線型包の annihilator は $D_1\to(U_i)_1$ の核である。
多重次数 $\boldsymbol a$ が $i$ 以外の正成分を持ち、$|\boldsymbol a|=3n$ なら、
その成分を含む $|\boldsymbol e|=3$, $\boldsymbol e\leq\boldsymbol a$ を選べる。
積 $T_{\boldsymbol e}T_{\boldsymbol a-\boldsymbol e}=T_{\boldsymbol a}$ は全射なので、
この成分は $(I_i)_1D_{n-1}$ に含まれる。
よって $I_i$ は次数一で生成される。
次数付き同型は点加群とその shift を保ち、第一座標に射影線型に作用するから、
曲線・偏極・点対応、およびこれらのイデアルを置換する。
\end{proof}

\begin{remark}[Segre 対角は商ではない]
\label{mbiq:rem:segre}
$|\boldsymbol e|=3$ の混合多重次数について、
$\bigoplus_{n\geq0}T_{n\boldsymbol e}$ は積曲線上の THCR 型の対角部分代数である。
しかし他の $D_1$ 成分を全て零にすることで、この代数を商として取り出すことはできない。
実際、その操作は全ての純部分 $(U_j)_1$ を零にする。
一方 $Q_{\boldsymbol e}^3=T_{3\boldsymbol e}$ はこれらの純部分の積で張られるので、
その商では $Q_{\boldsymbol e}^3=0$ となる。
この商は次数二以下にしか成分を持たず、点加群を持たない。
従って Segre 対角の積曲線を $D$ の追加の点成分とみなしてはいけない。
\end{remark}

\subsection{余法線双加群が回復する元の shift}

点のグラフだけでは $\sigma_i$ の三乗までしか見えない。
この不足は、内在的な商 $D\to U_i$ の余法線
\[
 N_i=I_i/I_i^2
\]
を左右の $U_i$ 作用付きで保持することで解消する。
以下の tensor は次数付き tensor であり、商圏へ移してから関手として解釈する。

\begin{proposition}[余法線による次数一の負 shift]
\label{mbiq:prop:conormal-root}
$r_i=3(k-1)>0$ とし、$U_i$-双加群
\[
 M_i=(T_i(-1))^{(3)},\qquad (M_i)_n=(T_i)_{3n-1}
\]
を置く。
次数一にのみ集中し、左右の $(U_i)_{>0}$ が零に作用する有限次元双加群 $Q_i$ があり、
\begin{equation}\label{mbiq:eq:conormal}
 N_i\simeq M_i^{\oplus r_i}\oplus Q_i
\end{equation}
である。
$-\otimes_{U_i}N_i$ が tails 上で誘導する完全関手を $\mathcal N_i$ と書くと、
THCR 同値の下で
\begin{equation}\label{mbiq:eq:conormal-functor}
 \mathcal N_i(\mathcal F)
 \simeq\bigl(\sigma_i^*\mathcal F\otimes L_i^{-1}\bigr)^{\oplus r_i}
 \qquad(\mathcal F\in\operatorname{coh}E_i)
\end{equation}
となる。
\end{proposition}
\begin{proof}
多重次数 $\boldsymbol a$ に対して
$b=\sum_{j\ne i}a_j$ を outside 次数と呼ぶ。
$I_i$ は $b\geq1$ の全成分からなる。
新しい次数 $n\geq2$ では、$I_i^2$ はちょうど $b\geq2$ の全成分である。
一方の包含は明らかである。
逆に $|\boldsymbol a|=3n$, $b\geq2$ なら、
\[
 1\leq t\leq\min(3,b-1),\qquad 3-t\leq a_i
\]
となる整数 $t$ を選べる。
outside 成分から合計 $t$ 個、$i$ 成分から $3-t$ 個を取って
$|\boldsymbol e|=3$, $\boldsymbol e\leq\boldsymbol a$ とする。
$\boldsymbol e$ と $\boldsymbol a-\boldsymbol e$ は両方 outside 次数が正なので、
切断乗法の全射性により $T_{\boldsymbol a}\subset I_i^2$ となる。

従って $n\geq2$ で生き残るのは $b=1$ の成分だけであり、
\[
 (N_i)_n=\bigoplus_{j\ne i}(T_i)_{3n-1}\otimes W_j.
\]
この部分は $n=1$ から存在し、左右の純 $i$-ブロック作用を保持して
$M_i\otimes(\bigoplus_{j\ne i}W_j)$ に等しい。
$n=1$, $b\geq2$ の残りを $Q_i$ とすると、
正次数の $U_i$ を左右どちらから掛けても $n\geq2$, $b\geq2$ に移るため零になる。
これで \eqref{mbiq:eq:conormal} を双加群として得る。

次に $M_i$ の tensor 関手を調べる。
$P$ を次数付き右 $U_i$-加群とし、$P^{[3]}$ は次数を三倍する操作とする。
次数成分を直接取ると、自然な同型
\begin{equation}\label{mbiq:eq:veronese-root-tensor}
 P\otimes_{U_i}M_i
 \simeq
 \left(\left(P^{[3]}\otimes_{U_i^{[3]}}T_i\right)(-1)\right)^{(3)}
\end{equation}
を得る。
右辺の induction と Veronese は、補題\ref{lem:nfm-gabriel}の互いに逆な tails 同値である。
従って $-\otimes_{U_i}M_i$ は、$\qgr T_i$ の $(-1)$ を
$\qgr U_i$ へ移した完全自己同値を誘導する。
有限次元項 $Q_i$ との tensor は正次数作用が零なので tails で消える。
よって $\mathcal N_i$ はこの自己同値の $r_i$ 重複である。
最後に \eqref{mbiq:eq:thcr-shift} の逆を使えば
\eqref{mbiq:eq:conormal-functor}となる。
\end{proof}

\begin{theorem}[中心商からの元の楕円 triple の回復]
\label{mbiq:thm:triple-recovery}
上の形の二つの中心商 $D,D'$ の間に次数付き同型があれば、
ブロック数は等しく、添字の置換後に曲線同型 $f_i:E_i\xrightarrow{\sim}E_i'$ が存在して
\begin{equation}\label{mbiq:eq:recovered-triple}
 f_i^*L_i'\simeq L_i,\qquad
 f_i\sigma_i=\sigma_i'f_i
\end{equation}
を満たす。
特に回復されるのは $\sigma_i^3$ だけでなく元の $\sigma_i$ である。
\end{theorem}
\begin{proof}
系\ref{mbiq:cor:intrinsic-quotients}により、同型は $I_i$、$U_i$ を置換し、
$f_i^*L_{i,3}'\simeq L_{i,3}$、$f_i\sigma_i^3=(\sigma_i')^3f_i$ を与える。
商の次数付き同型 $U_i\simeq U_i'$ が THCR 同値の下で誘導する関手は $f_{i,*}$ である。
実際、次数一写像はこの曲線同型と偏極同型の切断写像に grading scalar を掛けたものであり、
両代数は次数一生成なので全同型がそこで決まる。
THCR 切断関手の自然性と、grading scalar が tails 上で恒等に自然同型であることから従う。
基底線束との tensor の余分な曖昧さはない。

元の同型は $I_i/I_i^2$ を左右の商作用ごと同定する。
従って命題\ref{mbiq:prop:conormal-root}の関手を intertwine し、
\[
 f_{i,*}\bigl(\sigma_i^*\mathcal F\otimes L_i^{-1}\bigr)^{\oplus r_i}
 \simeq
 \bigl((\sigma_i')^*f_{i,*}\mathcal F\otimes(L_i')^{-1}\bigr)^{\oplus r_i}
\]
を得る。
$\mathcal F=K(x)$ とすると、左辺の支持は $f_i(\sigma_i^{-1}x)$、
右辺の支持は $(\sigma_i')^{-1}f_i(x)$ である。
$r_i>0$ なのでこれらは一致する。
全閉点での等式と曲線の被約性により $f_i\sigma_i=\sigma_i'f_i$ となる。

さらに $\mathcal F=\OO_{E_i}$ を代入すると
\[
 (f_{i,*}L_i^{-1})^{\oplus r_i}
       \simeq((L_i')^{-1})^{\oplus r_i}
\]
を得る。
行列の少なくとも一成分は零でないので、二つの次数 $-3$ の線束の間に非零な射がある。
同次数の線束間の非零射は同型であり、$f_i^*L_i'\simeq L_i$ が従う。
\end{proof}

\subsection{可換 tensor 表示に対する全構造対象保存同値}

\begin{corollary}[可換 tensor チャート内の全 $\OO$-同値と対象剛性]
\label{mbiq:cor:commuting-o-rigidity}
$A,A'$ を第\ref{mbs:sec:presentations}節の満支持可換 tensor 積表示とし、
代数閉体上で、定理\ref{mbv:thm:all-o-third}の有限個の位数除外条件
$jp_i\ne0$, $jp_i'\ne0$ $(1\leq j\leq3k-3)$ を満たすとする。
このとき
\begin{equation}\label{mbiq:eq:object-rigidity}
 (\QGr A,\OO_A)\simeq(\QGr A',\OO_{A'})
 \quad\Longleftrightarrow\quad A\simeq_{\mathrm{gr}}A'.
\end{equation}
さらに、与えられた任意の $\OO$-保存同値関手は、
次数付き同型 $A\to A'$ の構造移送と自然同型である。
この次数付き同型は共通 grading scalar を除いて一意である。
この関手の認識には、対象剛性に加えて定理\ref{mbca:thm:full-inertia}を用いる。
\end{corollary}
\begin{proof}
全 $\OO$-保存同値から、定理\ref{mbv:thm:all-o-third}により
$C=A^{(3)}$ と $C'=(A')^{(3)}$ の次数付き同型を得る。
ここでは代数閉体上で、grading scalar の torsor を自明化してよい。
中心一次生成イデアルで割ると、\eqref{mbiq:eq:central-quotient}により $D\simeq D'$ となる。
定理\ref{mbiq:thm:triple-recovery}から、因子の元の楕円 triple が置換を除いて一致する。
二次関係は
\[
 \ker\left(H^0(E_i,L_i)\otimes H^0(E_i,L_i)
 \longrightarrow H^0(E_i,L_i\otimes\sigma_i^*L_i)\right)
\]
なので、この triple 同型は Sklyanin 因子の次数付き同型を与える。
従って周囲代数も同型である。
各三次中心は直線であり、全ての中心係数は非零なので、
ブロックごとの grading scalar の三乗を選んで商の中心直線を一致させられる。
このとき周囲同型は $A\simeq A'$ へ降りる。
逆向きは次数付き同型の構造移送による。

最後に、与えられた同値を $F$ とし、上で得た次数付き同型の構造移送を $F_0$ とする。
$F_0^{-1}F$ は $\QGr A$ の $\OO$-保存自己同値である。
定理\ref{mbca:thm:full-inertia}により、これは $A$ の次数付き自己同型から来る。
従って $F$ 自体も次数付き同型の構造移送である。
共通 grading scalar を除いた一意性は、命題\ref{mbva:prop:faithful}の忠実性による。
ここで追加した関手の認識は、前半の対象剛性や定理\ref{mbiq:thm:triple-recovery}の証明には用いていない。
\end{proof}

\begin{remark}[全自己同型群・相対認識・第三根の区別]
一般の $j$ が相異なり、$j_i\ne0,1728$ であるとき、
全 $\OO$-保存自己同値が各点曲線に誘導する自己同型は、
次数三偏極を保って $\sigma_i$ と可換する $E_i[3]$ の平行移動に限られる。
次数九の偏極だけから得られる $E_i[9]$ への拡大は、余法線によって排除される。
全曲線上で恒等な自己同値の核も、定理\ref{mbca:thm:full-inertia}により
$\mu_3^k/\mu_3^{\mathrm{diag}}$ と決定される。

非被約底上の同型層の一致は幾何点の認識だけからは従わないが、
系\ref{mbrr:cor:stack-isomorphism}によって、さらに選んだ非空開上で成立する。
また、補題\ref{mbca:lem:saturated-powers}は $I_i$ の飽和冪から
全多重次数を回復し、純因子 $U_i\subset D$ の埋込み自体も内在的に認識する。

これらは任意の第三根表示を可換 tensor 積として認識する主張ではない。
第\ref{mbr:sec:roots}節の非可換な混合関係を持つ Zhang twist では、
余法線の左右作用を本節の可換 tensor 公式に置き換えられない。
従って同じ第三 Veronese を持つ非同型の twist の存在とは矛盾しない。
系\ref{mbca:cor:marked-roots}は $\mathsf r\OO\simeq\OO(1)$ を満たす根を尽くすが、
$\mathsf r^3\simeq(3)$ だけを仮定した全ての根の分類は残る。
\end{remark}

% !TEX root = ../nc-cubic.tex
\section{第三 Veronese の表示作用と次数一軌道の認識}
\label{mbva:sec:automorphisms}

本節では、表示の自己同型群が第三 Veronese へ忠実に作用することと、
その像が構造対象の全次数一軌道を保つ同値群に一致することを示す。
二ブロックでは次数三の Heisenberg 表現も明示する。
これらは定理\ref{mbca:thm:full-inertia}による全自己同値群の決定の入力となる。
一般二次代数に対する中心還元核の上界も記録するが、
満支持型では定理\ref{mbcr:thm:reduction-rigidity}が核の自明性を与える。

基礎体は標数零の代数閉体とする。
第\ref{mbr:sec:roots}節の記号により、ブロック数を $\ell\ge2$、
\[
 B=\bigotimes_{i=1}^{\ell}S_i,\qquad
 A=B/(w),\qquad w=\sum_i a_i\Omega_i,\quad a_i\ne0,\qquad C=A^{(3)}
\]
とする。因子の $j$ 不変量は相異なり $0,1728$ を避け、
$p_i,2p_i,3p_i\ne0$ とする。全 $\OO$-同値を扱う箇所では
定理\ref{mbv:thm:all-o-third}の有限捩れ条件も課す。
第\ref{mbr:sec:roots}節の直接的な Zhang twist 表示の分類を引用する場合には、
その箇所で追加条件 $6p_i\ne0$ を課す。
以下の次数一軌道の議論は、全整数の切断飽和性が成立する基礎体上で行う。
本稿の $\C$ 上の表示族では系\ref{mbsd:cor:saturation}がこれを与える。
記号を
\[
 \mathscr C=\operatorname{QGr}A,\qquad
 \mathsf s=(1),\qquad \OO(j)=\mathsf s^j\OO,\qquad
 G=\operatorname{Aut}_{\OO}(\mathscr C)
     \cong\operatorname{Aut}_{\mathrm{gr}}(C)/\mathbf G_m
\]
とする。$G$ の元は構造対象の同定を持つ同値を、
その同定を保つ自然同型で割ったものである。
表示自己同型群は
\[
 \mathcal K_\ell=
 \operatorname{Aut}_{\mathrm{gr}}(A)/\mathbf G_m
 =\Bigl(\prod_i\mathcal H_i\Bigr)/\mu_3^{\mathrm{diag}}
\]
である。ここで各 $\mathcal H_i$ は中心三次元を固定する指数三・位数 $27$ の
Heisenberg 群である。

\subsection{全次数の第三 Veronese への忠実性}

\begin{proposition}[表示慣性は第三 Veronese 上でも忠実である]
\label{mbva:prop:faithful}
制限による準同型
\[
 \mathcal K_\ell\longrightarrow
 \operatorname{Aut}_{\mathrm{gr}}(C)/\mathbf G_m
\]
は単射である。二ブロックでは位数 $243$ の $\mathcal K_2$ が忠実に作用する。
定理\ref{mbca:thm:full-inertia}は、これが全 $\OO$-自己同値群に等しいことを示す。
\end{proposition}
\begin{proof}
$h=(h_i)\in\prod_i\mathcal H_i$ の作用が $C$ 上の次数スカラーになるとし、
$C_1=A_3$ 上の値を $\lambda$ とする。
$h$ は全ての $\Omega_i$ を固定し、その像の張る $Z(A)_3$ は
次元 $\ell-1>0$ なので $\lambda=1$ である。

異なる $i,j$ について、混合部分
\[
 (S_i)_2\otimes W_j\subset B_3
\]
は $A_3$ へ単射に入る。実際、三次関係 $w$ は純粋なブロック三次部分にのみ成分を持つ。
この部分で $h$ は $(h_i)_2\otimes(h_j)_1$ として作用する。
非零な有限次元空間上で二つの可逆写像の tensor 積が恒等なら、
各写像は互いに逆のスカラーである。
従って各 $(h_j)_1=\varepsilon_j\operatorname{id}$ であり、
$h_j\in\mathcal H_j$ より $\varepsilon_j^3=1$ となる。
混合部分上の恒等式は
$\varepsilon_i^2\varepsilon_j=1$、すなわち $\varepsilon_j=\varepsilon_i$ を与える。
よって $h$ は共通 $\mu_3$ に属し、$\mathcal K_\ell$ では恒等である。
\end{proof}

\begin{remark}[twist の積と自己同型の作用]
\label{mbva:rem:scalar-action}
原始三乗根 $\zeta$ を固定する。
二ブロックで $h=(\zeta\operatorname{id}_{W_1},\operatorname{id}_{W_2})$ とすると、
$A_3$ の三つの固有空間は
\[
 \begin{array}{c|c|c}
 \text{部分空間}&\text{固有値}&\text{次元}\\ \hline
 ((S_1)_3\oplus(S_2)_3)/kw&1&19\\
 (S_1)_2\otimes W_2&\zeta^2&18\\
 W_1\otimes(S_2)_2&\zeta&18
 \end{array}
\]
である。したがってブロックスカラー $\mu_3$ は第三 Veronese 上で消えない。
$(A^h)^{(3)}=A^{(3)}$ という等式は、次数三の倍数から始まる
Zhang twist の積が元の積に等しいという意味であり、
$h$ 自身のこの自己同型作用が恒等になるという意味ではない。
\end{remark}

\subsection{全次数一軌道を保つ同値関手の認識}

\begin{theorem}[SSS 部分群の関手レベルの認識]
\label{mbva:thm:sss}
$A,A'$ を上の仮定を満たす二つの可換 tensor 積中心商とする。
全整数について
$F(\OO(j))\simeq\OO'(j)$ を満たす $\OO$-保存同値の集合は、
自然に
\[
 \operatorname{Isom}_{\mathrm{gr}}(A,A')/k^\times
\]
と同定される。特に $G$ の SSS 部分群
\[
 G_{\mathrm{SSS}}=
 \{F\in G:F(\OO(j))\simeq\OO(j)\ \text{for all }j\in\mathbf Z\}
\]
は $\mathcal K_\ell$ に正確に一致する。
\end{theorem}
\begin{proof}
各整数 $j$ について同型 $F(\OO(j))\simeq\OO'(j)$ を選ぶ。
切断飽和と $F$ の完全忠実性により、隣接 Hom 上に可逆線型写像
\[
 h_j:V=A_1\longrightarrow V'=A'_1
\]
を得る。合成を長さ二へ適用すると
\begin{equation}\label{mbva:eq:adjacent-relations}
 (h_{j+1}\otimes h_j)R_B=R_{B'}
\end{equation}
である。右加群では切断 $x$ の写像は左乗法であり、
二つの射の合成は $f_y\circ f_x=f_{yx}$ となるため、この tensor 順序を用いる。

式\eqref{mbva:eq:adjacent-relations}の双線型零点スキームの第一射影と第二射影を取る。
系\ref{tbp:cor:all-blocks}により、それらは各々の点曲線の非交和であり、
$h_j,h_{j+1}$ は両側の点曲線の線型包を対応させる。
$j$ 不変量が相異なるため、成分の対応は $j$ 不変量の一致から一意である。
したがって、このブロックの置換は整数 $j$ に依存しない。
標的のブロックを一度並べ直すと、
$h_j=\bigoplus_i h_{j,i}$ となる。

異なる二ブロック $i,t$ の混合関係は可換関係である。
式\eqref{mbva:eq:adjacent-relations}を
$(W'_i\otimes W'_t)\oplus(W'_t\otimes W'_i)$ に制限し、
第一成分上のグラフとして比較すると
\[
 (h_{j,t}h_{j+1,t}^{-1})
       \otimes(h_{j+1,i}h_{j,i}^{-1})
   =1_{W'_t}\otimes1_{W'_i}.
\]
従って一つの $c_j\in k^\times$ が存在して、全ブロックについて
$h_{j+1,i}=c_jh_{j,i}$ となる。
二ブロックの場合にも一つの混合対が両方のスカラーを一致させ、
三ブロック以上では各対の共通成分を通して同じ結論を得る。
ゆえに $h_j=d_jh_0$ となる $d_j\in k^\times$ が全整数で存在する。

各対象の同型の frame を $b_j\in k^\times$ で変更すると、
隣接 Hom の写像は $(b_{j+1}/b_j)h_j$ になる。
$b_0=1$ とし、$b_{j+1}/b_j=d_j^{-1}$ を正負両方向へ再帰して、
全ての $h_j$ を同じ写像 $h=h_0$ に揃えられる。
次数二の合成から $(h\otimes h)R_B=R_{B'}$ なので、
$h$ は二次包絡 $B\simeq B'$ を与える。
長さ三の合成は、その二次包絡内の三次核 $kw$ を $kw'$ へ移す。
従って $h$ は $A\simeq A'$ の次数付き同型である。

$A,A'$ は次数一生成なので、全ての
$\Hom(\OO(j),\OO(j+m))$ への作用は隣接 Hom と合成で決まる。
上の frame の下で $F$ と $h$ の構造移送はこれら全てに一致する。
Gabriel 切断による完全忠実性と $\{\OO(j)\}_{j\in\mathbf Z}$ の Abel 生成性から、
両関手は構造対象の同定を保つ自然同型を持つ。
最後に、次数付き同型の構造移送が同じになる核は共通 grading scalar である。
これは隣接 Hom 上の自然性、または命題\ref{mbva:prop:faithful}と
定理\ref{mbv:thm:all-o-third}から従う。
\end{proof}

\begin{corollary}[二つの対象と次数一自己同値の安定化群]
\label{mbva:cor:pair-stabilizer}
全 $\OO$-同値の第三 Veronese 認識が成立する範囲では、
\[
 \mathcal K_\ell
  =\operatorname{Stab}_{G}([\OO(1)],[\OO(2)])
  =\{F\in G:F\mathsf s\simeq\mathsf s F\}.
\]
さらに定理\ref{mbca:thm:full-inertia}と合わせると、
本稿の満支持型ではこれらは全て $G$ に等しく、
$\operatorname{Stab}_G([\OO(1)])=G=\mathcal K_\ell$ でもある。
\end{corollary}
\begin{proof}
$G$ の各元は $C$ の次数付き自己同型から来るので、
$\mathsf s^3$ と可換する。
従って $\OO,\OO(1),\OO(2)$ を保てば全整数の $\OO(j)$ を保つ。
逆は明らかで、定理\ref{mbva:thm:sss}を適用する。
また $\mathsf s$ と可換する $\OO$-保存同値は SSS である。
逆に次数付き自己同型の構造移送は次数一 shift と可換する。
最後の等式は定理\ref{mbca:thm:full-inertia}の $G=\mathcal K_\ell$ による。
\end{proof}

\subsection{既知の三乗根の対象軌道と全同値による共役}

右加群の構造移送の規約として、$\Phi_h$ は
作用を $m\cdot a=m h^{-1}(a)$ と変える関手とする。
これにより $h\mapsto\Phi_h$ は準同型となる。
共通 grading scalar は自然同型で恒等になるので、
$h\in\mathcal K_\ell$ に対して $\Phi_h$ が定まる。

\begin{proposition}[代数的 twist の根は同じ対象軌道を持つ]
\label{mbva:prop:root-functors}
$\tau\in\prod_i\mathcal H_i$ に対し、共通の第三 Veronese を使って
$\operatorname{QGr}A^\tau$ を $\mathscr C$ へ同定する。
その次数一 shift を移した自己同値は
\begin{equation}\label{mbva:eq:root-functor}
 \mathsf s_\tau=\Phi_{\tau^{-1}}\mathsf s,\qquad
 \mathsf s_\tau^3\simeq\mathsf s^3
\end{equation}
であり、全整数 $j$ について
$\mathsf s_\tau^j\OO\simeq\OO(j)$ となる。
したがって $\mathcal K_\ell$ の各元は、固定した圏の上で互いに異なる
$\mathsf s^3$ の三乗根の自然同型類を与える。
\end{proposition}
\begin{proof}
Zhang 関手を $\mathsf T_\tau:\operatorname{Gr}A\to\operatorname{Gr}A^\tau$ と書く。
次数 $d$ の元の右作用は $m*a=m\tau^d(a)$ である。
第三 Veronese では $d=3n$ なので、この関手の制限は恒等である。
従って共通 Veronese による比較は $\mathsf T_\tau^{-1}$ に等しい。

$\mathsf T_\tau^{-1}(1)\mathsf T_\tau$ を計算すると、
shift 後の次数 $d$ の元の作用は
$m\tau^{d+1}(\tau^{-d}(a))=m\tau(a)$ である。
これが式\eqref{mbva:eq:root-functor}を与える。
$\Phi_\tau$ は $\mathsf s$ と可換し、$\tau^3=1$ なので三乗の式も従う。
$A^\tau(j)$ を元の圏へ移した右加群の作用は
$m\cdot a=m\tau^j(a)$ であり、$m\mapsto\tau^{-j}(m)$ が
元の $A(j)$ との同型を与える。これは負の整数にも成立する。
最後に $\mathsf s_\tau\simeq\mathsf s_\upsilon$ なら
$\Phi_{\tau^{-1}}\simeq\Phi_{\upsilon^{-1}}$ であり、
命題\ref{mbva:prop:faithful}から $\bar\tau=\bar\upsilon$ となる。
\end{proof}

\begin{corollary}[全 $\OO$-同値は既知の共役類を融合しない]
\label{mbva:cor:root-orbits}
$F\in G$ が $\mathsf s_\tau$ を $\mathsf s_\upsilon$ へ共役するなら、
$F\in\mathcal K_\ell$ である。さらに
\[
 \operatorname{Stab}_{G}(\mathsf s_\tau)
       =C_{\mathcal K_\ell}(\bar\tau).
\]
従って有限群 $\mathcal K_\ell$ の
$(11^\ell+2\cdot8^\ell)/3$ 個の共役類は
（第\ref{mbr:sec:roots}節の有限群計算）、
全自己同型群 $G$ による共役を許しても互いに異なる。
二ブロックでは、固定した圏上の $243$ 個の既知の三乗根が
全 $G$ の下でも正確に $83$ 個の異なる共役類に属する。
系\ref{mbca:cor:marked-roots}は、これらが
$\mathsf r\OO\simeq\OO(1)$ を満たす全ての根を尽くすことを示す。
構造対象の像を指定しない全ての根の分類とは区別する。
また定理\ref{mbr:thm:algebraic-twist-groupoid}によって
直接 Zhang twist 代数の同型 groupoid と比較する場合には、$6p_i\ne0$ を追加する。
\end{corollary}
\begin{proof}
$F\OO\simeq\OO$ と共役条件から、全整数 $j$ について
$F(\mathsf s_\tau^j\OO)\simeq\mathsf s_\upsilon^j\OO$ である。
命題\ref{mbva:prop:root-functors}によりこれは SSS 条件そのものである。
従って $F\in\mathcal K_\ell$ となる。
この群は $\mathsf s$ と可換し、その $\mathsf s_\tau$ への共役作用は
$\bar\tau$ への通常の共役作用に一致する。
\end{proof}

定理\ref{mbca:thm:full-inertia}により $|G|=3^{3\ell-1}$ であり、
順序付き対象対 $([\OO(1)],[\OO(2)])$ の $G$-軌道は一点である。
その全群の決定は、中心商の多重次数の回復と中心方向の持上げ剛性に基づく。

\subsection{一般二次代数の中心還元における有限核}

以下は一般二次代数に使える有限群の判定である。
満支持二ブロックでは $Z(C)_1=kg$ となるが、
その還元核は定理\ref{mbcr:thm:reduction-rigidity}により、
下記の有限性や変形類の条件を仮定せずに自明となる。

\begin{proposition}[二次代数の中心還元の有限核]
\label{mbva:prop:central-reduction-kernel}
$C$ を次数一生成の連結二次代数、$g\in Z(C)_1$ を正則元とし、
$Z(C)_1=kg$、$D=C/(g)$、$D_1\ne0$ とする。
$G=\operatorname{Aut}_{\mathrm{gr}}(C)/k^\times$ が有限群であると仮定する。
このとき制限写像
\[
 G\longrightarrow\operatorname{Aut}_{\mathrm{gr}}(D)/k^\times
\]
の核は自明群または位数二の中心部分群である。

さらに、$C/(g^2)$ が定める一次変形類
$\kappa_C\in HH^2(D,D)_{-1}$ が非零なら、この核は自明である。
後半は一般的な十分条件である。
満支持二ブロックの一次変形類の非消滅は、系\ref{mbcr:cor:nonzero-class}で別途証明する。
\end{proposition}
\begin{proof}
核の元には、$D$ 上の作用が文字通り恒等となる
唯一の次数付き自己同型の代表 $\alpha$ がある。
実際、共通 grading scalar によってそのように規格化でき、
$D_1\ne0$ が一意性を与える。この規格化は積と可換するので、
これらの代表自身も有限群をなす。
$\alpha(g)=\lambda g$ と書くと、$\alpha\mapsto\lambda$ は核の指標となる。
$\lambda=1$ なら $C_1$ 上で
$\alpha=1+g\otimes l$、$l(g)=0$ であり、これは unipotent である。
有限群の元なので標数零では恒等となり、$C$ の次数一生成性から $\alpha=1$ である。
従ってこの指標は単射で、核は循環群である。

その位数を $m\ge3$ と仮定し、生成元 $\alpha$ を取る。
有限群表現の半単純性により
$C_1=kg\oplus U$ と分解でき、$\alpha$ は $g$ に原始 $m$ 乗根 $\lambda$、
$U\simeq D_1$ に恒等として作用する。
二次 tensor の固有値は $1,\lambda,\lambda^2$ の三つで、互いに異なる。
二次関係空間をそれらの固有空間へ分解する。
固有値 $\lambda$ の関係は $g\otimes u-u\otimes g$ だけである。
それ以外があれば $g$ と $U$ の非零な元との積が零となり、正則性に反する。
固有値 $\lambda^2$ の関係は $g^2\ne0$ により存在しない。
従って全ての残りの関係は $U\otimes U$ に属し、
\[
 C\simeq D[g]
\]
となる。しかしこの場合は $U$ を固定し $g$ だけを尺度変更する
$\mathbf G_m$ が $G$ に単射に入るため、有限性に反する。
よって $m\le2$ である。位数二の核は正規部分群なので $G$ の中心に入る。

最後に $C/(g^2)$ を、次数一の変形変数 $g$ による
$D$ の一次変形と見る。
$D$ 上恒等で $g\mapsto\lambda g$ となる自己同型が存在すれば、
その一次変形類は $\kappa_C=\lambda\kappa_C$ を満たす。
$\kappa_C\ne0$ なら $\lambda=1$ であり、前半の単射性から核は自明となる。
\end{proof}

満支持二ブロックでは $Z(C)_1=kg$、$g$ の正則性と二次表示が成立し、
定理\ref{mbca:thm:full-inertia}によって全群も有限である。
従ってこの一般判定を適用できるが、
実際の還元核は定理\ref{mbcr:thm:reduction-rigidity}により自明である。

ここで $g$ を固定する条件を省いて、単に
$\alpha\bmod g=1$ から unipotent 性を推論してはならない。
例えば二次代数
$k\langle x,y,g\rangle/([g,x],[g,y],xy-yx-g^2)$ では、
$g\mapsto-g$、$x,y$ 固定は、$g$ で割ると恒等となる位数二の自己同型である。
これは上の一般的な議論で二次の可能性が残る理由を示す例であり、
本稿の二ブロック代数にこの自己同型が存在するという主張ではない。

\subsection{二ブロックの次数三表現}

以下は $\ell=2$ とする。$Q_i=\mathcal H_i/\mu_3\simeq\mathbf F_3^2$、
$Q=Q_1\oplus Q_2$ と置く。
$\mathcal K_2$ の中心 $Z$ は位数三で、$Q=\mathcal K_2/Z$ である。
中心の二つの非自明な指標を持つ九次元既約表現を
$U_+=W_1\otimes W_2^\vee$、$U_-=W_1^\vee\otimes W_2$ と書く。
二つの因子の線型指標は、$Q$ の指標へ膨張して使う。

\begin{proposition}[五十五次元表現の完全な分解]
\label{mbva:prop:representation}
$\mathcal K_2$-表現として
\begin{equation}\label{mbva:eq:representation}
 C_1=A_3\simeq
 k^{\oplus3}
 \oplus\bigoplus_{\substack{\chi\in Q_1^\vee\\\chi\ne1}}k_\chi
 \oplus\bigoplus_{\substack{\chi\in Q_2^\vee\\\chi\ne1}}k_\chi
 \oplus U_+^{\oplus2}\oplus U_-^{\oplus2}.
\end{equation}
特に、中心の固有空間の次元は $19,18,18$ である。
$h\in\mathcal K_2$ の $Q$ における像の support を $s$ とすると、
$C_1$ 上の指標は
\[
 \operatorname{tr}(h|C_1)=
 \begin{cases}
 55&h=1,\\
 10&s=1,\\
 1&\text{それ以外}
 \end{cases}
\]
となる。
\end{proposition}
\begin{proof}
各 $\mathcal H_i$ の二次関係空間 $R_i$ は、
非自明中心指標を持つ三次元既約表現なので $W_i^\vee$ と同型である。
従って $(S_i)_2\simeq(W_i^\vee)^{\oplus2}$ となる。
Heisenberg の標準生成元は行列式一であり、
三次 syzygy 直線 $P_i\simeq\det W_i$ は自明表現である。
ここでは補題\ref{mbsd:lem:potential-line}の potential 直線の同定を使う。
Koszul 解消の次数三部分と半単純性から
\[
 (S_i)_3\simeq(W_i\otimes W_i^\vee)\oplus k
             \simeq k[Q_i]\oplus k
\]
を得る。基底 $e_j$（$j\in\mathbf Z/3$）上の
$Xe_j=e_{j-1}$、$Ye_j=\zeta^je_j$ を用いると、
最後の同型は九つの行列 $X^aY^b$ が
$\operatorname{End}(W_i)$ において $Q_i$ の全指標を各一回ずつ与えることによる。

純三次部分 $(S_1)_3\oplus(S_2)_3$ から不変直線 $kw$ を除くと、
自明表現が三回、二つの $Q_i$ の非自明指標が各一回残る。
二つの混合部分は、それぞれ $U_-^{\oplus2},U_+^{\oplus2}$ である。
これで式\eqref{mbva:eq:representation}を得る。

各 $W_i$ の指標は、中心元でのみ非零である。
$(S_i)_3$ の指標は $Q_i$ における像が零なら $10$、そうでなければ $1$ である。
この値と混合部分の tensor 積の指標を足せば、表示した trace の公式を得る。
\end{proof}

\begin{corollary}[線型空間内の中心化群]
\label{mbva:cor:linear-centralizer}
この作用の射影線型中心化群は
\[
 Z_{\operatorname{PGL}(C_1)}(\mathcal K_2)
 \cong
 \bigl(\operatorname{GL}_3\times\mathbf G_m^{16}
                  \times\operatorname{GL}_2\times\operatorname{GL}_2\bigr)/\mathbf G_m
\]
であり、次元は $32$ である。
これは $C_1$ という線型空間における中心化群であって、
代数 $C$ の自己同型群における中心化群ではない。
\end{corollary}
\begin{proof}
線型中心化群の表示は式\eqref{mbva:eq:representation}と Schur の補題による。
射影的に中心化する写像は、線型持上げにおいてある一次指標 $\chi$ を用いた
$C_1\simeq C_1\otimes\chi$ を与える。
しかし一次指標のうち重複度三を持つものは自明指標だけなので、$\chi=1$ である。
従って射影中心化群は線型中心化群の共通スカラー商に等しい。
\end{proof}

\begin{proposition}[線型空間内の正規化群の有限商]
\label{mbva:prop:linear-normalizer}
$N= N_{\operatorname{PGL}(C_1)}(\mathcal K_2)$、
$Z_{\mathrm{lin}}= Z_{\operatorname{PGL}(C_1)}(\mathcal K_2)$ と置く。
$N/Z_{\mathrm{lin}}$ は位数 $81\cdot2304=186624$ の有限群であり、
\[
 1\longrightarrow Q
 \longrightarrow N/Z_{\mathrm{lin}}
 \longrightarrow\Gamma\longrightarrow1,\qquad
 |\Gamma|=2304
\]
となる。$\Gamma$ は $Q_1,Q_2$ の順序なし分解を保つ symplectic similitude 群である。
ブロックを保つ部分は
\[
 \{(g_1,g_2)\in\operatorname{GL}_2(\mathbf F_3)^2:
                    \det g_1=\det g_2\}
\]
であり、これにブロック交換を加えた群が $\Gamma$ である。
\end{proposition}
\begin{proof}
$\mathcal K_2$ は指数三の extraspecial 群であり、その交換子は
$Q$ 上の非退化交代形式 $\omega=\omega_1-\omega_2$ を定める。
半分の交換子を積に使う Heisenberg 座標では、
積は $(q,t)(q',t')=(q+q',t+t'+\tfrac12\omega(q,q'))$ であり、
multiplier が $\varepsilon$ の similitude $u$ は
$(q,t)\mapsto(uq,\varepsilon t)$ に持ち上がる。
自己同型群から $\operatorname{GSp}(Q,\omega)$ への写像は全射である。
その核は $\operatorname{Hom}(Q,\mathbf F_3)$ であり、
非退化 pairing によって内部自己同型群 $Q$ と同一視される。

射影正規化群の元も、ある一次指標を掛けた表現の intertwiner である。
前系と同じ重複度三の一意性により、その指標は自明である。
従って正規化群の有限商は
式\eqref{mbva:eq:representation}の同型類を保つ $\mathcal K_2$ の自己同型群に等しい。
一次指標の support は
$(Q_1^\vee\cup Q_2^\vee)\setminus\{0\}$ なので、
許される similitude は二つの二次元平面を置換しなければならない。
実際、両平面の和集合に含まれる二次元部分空間は、この二平面のいずれかである。
逆にそれらを置換する similitude は一次指標の重複度を保ち、
中心指標の符号変更も $U_+,U_-$ の等しい重複度を保つので、
必ず線型 intertwiner を持つ。

各二次元ブロックでは similitude multiplier は行列式である。
二ブロックの行列式が等しい組は
$2\cdot24^2=1152$ 個あり、交換を許せば $2304$ 個となる。
内部自己同型群の位数 $|Q|=81$ を掛けて主張を得る。
\end{proof}

この線型正規化群を実際の $G$ の部分群と取り違えない。
$N$ の中で代数自己同型を求める際にも、
さらに $C$ の $2619$ 次元の二次関係空間を保存する条件を課す必要がある。
定理\ref{mbca:thm:full-inertia}により $G=\mathcal K_2$ なので、
関係空間を保つ実際の射影的代数自己同型は、この有限部分群に正確に一致する。

\begin{remark}[検算と残る認識]
\label{mbva:rem:verification}
\texttt{computations/verify\_veronese\_representation.py} は、
既存の Heisenberg 群の座標実装と Koszul 指標式を入力とし、
全 $243$ 元での次数三の指標、
既約分解の重複度、中心の三つの固有空間次元、および
$\Gamma$ の位数 $2304$ を有限体・整数演算で検算する。
これは有限群表現の検算であって、
全 $G$ の計算や次数一の全根の分類ではない。
本節の認識は、表示作用の忠実性、SSS 関手そのものの認識、
および表示から来る三乗根の共役作用である。
全群 $G=\mathcal K_\ell$ の決定は
定理\ref{mbcr:thm:reduction-rigidity}と定理\ref{mbca:thm:center-automorphisms}を
組み合わせる定理\ref{mbca:thm:full-inertia}による。
$\OO(1)$ の $G$-軌道は一点であり、その安定化群は全 $G$ である。
構造対象の像を指定しない第三根の分類と、任意の非被約底上の相対認識は別問題である。
これらを組み合わせた主張の新規性の網羅的確認は未完了である。
\end{remark}

% !TEX root = ../nc-cubic.tex
\section{中心商上の線対象と次数付き持上げの剛性}
\label{mbcr:sec:central-rigidity}

本節では、中心商上で次数一軌道の対象を保つ自己同型は、
元の第三 Veronese 上でもそれらを保つことを示す。
二ブロックで残っていた中心還元核を、自己同値群の有限性を仮定せずに除く。
証明は中心変数が複数あっても同じなので、全ての満支持 $k\geq2$ を扱う。
必要なのは中心上の平坦性、負次数の $\Ext^1$ と $\Hom$ の消滅、
および正次数の中心イデアルによる商の次数別安定性である。

基礎体を標数零の代数閉体 $K$ とし、
第\ref{mbiq:sec:quotients}節の仮定の下で
\[
 B=\bigotimes_{i=1}^kS_i,\qquad
 A=B/\left(\sum_i a_i\Omega_i\right),\quad a_i\ne0,\qquad
 T_i=S_i/(\Omega_i),\qquad E=\bigotimes_iT_i
\]
と置く。本節の $E$ は次数付き代数であり、楕円曲線 $E_i$ と区別する。
さらに
\begin{equation}\label{mbcr:eq:setup}
 C=A^{(3)},\quad Z_1=Z(C)_1,\quad
 R=\Sym_K Z_1,\quad \mathfrak m=R_{>0},\quad I=\mathfrak m C,\quad
 D=C/I=E^{(3)}
\end{equation}
とする。$R$ の生成元の次数は一、$\dim_KZ_1=k-1$ である。
中心一次空間の同定には命題\ref{mbv:prop:center}の $d=1$ を使う。
各整数 $j$ について
\[
 M_j=(A(j))^{(3)},\qquad \overline M_j=M_j/IM_j=(E(j))^{(3)}
\]
と書く。内部次数の規約は $M(j)_a=M_{a+j}$ である。
以下の $\Ext_{\Gr}$ と $\Hom_{\Gr}$ は次数零の射について取る。
内部次数を全て集める場合には下線を付ける。

\subsection{中心上の自由性と完全交叉の局所コホモロジー}

\begin{lemma}[中心多項式環上の自由性]
\label{mbcr:lem:central-free}
$R\to C$ は単射であり、$C$ と全ての $M_j$ は
次数付き $R$-加群として自由である。
自由基底は無限でもよいが、各内部次数は有限次元で、次数は下に有界である。
\end{lemma}
\begin{proof}
$T_i$ の斉次基底を $S_i$ へ持ち上げ、その張る空間を $H_i$ とする。
$\Omega_i$ の正則性と
$S_i/\Omega_iS_i=T_i$ により、乗法は加群同型
\[
 K[z_i]\otimes_KH_i\xrightarrow{\sim}S_i,\qquad z_i\longmapsto\Omega_i,
 \quad\deg z_i=3
\]
を与える。全射性は次数に関する帰納で従う。
線型関係があれば、最小の $z_i$ 冪を正則性によって消去し、
$z_i$ で割った商を取ることで、その最小係数が零となる。
これを繰り返すと単射性も得られる。

従って $H=\bigotimes_iH_i$ とすると、加群として
\[
 B\simeq K[z_1,\ldots,z_k]\otimes_KH,\qquad
 A\simeq
 \left(K[z_1,\ldots,z_k]/\left(\sum_i a_iz_i\right)\right)\otimes_KH .
\]
中心多項式環の次数を三で割ると、括弧内は $R$ となる。
$H$ の次数 $d$ の基底元は、$d\equiv j\pmod3$ のときに限り
$M_j$ の自由基底を与え、その新しい次数は $(d-j)/3$ である。
特に $1\in H_0$ の成分が $R\to C$ の単射性を与える。
剰余を取れば、表示した $\overline M_j$ の同定も従う。
\end{proof}

\begin{lemma}[中心完全交叉の双対と低次切断]
\label{mbcr:lem:ci-cohomology}
有限次元の次数付き右 $E$-加群 $F$ に対し、
係数の一次元空間を自明化して
\begin{equation}\label{mbcr:eq:finite-duality}
 \mathbf R\underline{\Hom}_E(F,E)\simeq F^\vee[-2k]
\end{equation}
を得る。従って、irrelevant torsion の導来関手を
$\mathbf R\Gamma_E$ と書けば
\begin{equation}\label{mbcr:eq:ci-local-cohomology}
 \mathbf R\Gamma_E(E)\simeq E^\vee[-2k].
\end{equation}
特に $E$ は局所コホモロジーが次元 $2k$ に集中する
AS--Cohen--Macaulay 代数である。
$\mathcal O_E=\pi_EE$ とすると、全整数 $b$ について
\begin{equation}\label{mbcr:eq:ci-low-sections}
 \Hom_{\QGr E}(\mathcal O_E,\mathcal O_E(b))=E_b,\qquad
 \Ext^q_{\QGr E}(\mathcal O_E,\mathcal O_E(b))=0
 \quad(1\leq q\leq2k-2)
\end{equation}
となる。
\end{lemma}
\begin{proof}
前補題の tensor 積表示により、$\Omega_1,\ldots,\Omega_k$ は
$B$ 上の中心正則列である。
定理\ref{mbsd:thm:ambient-duality}の有限 Koszul 双対から
\[
 \mathbf R\underline{\Hom}_B(F,B)\simeq F^\vee(3k)[-3k].
\]
一方、この中心正則列の Koszul 解消を双対化すると
\[
 \mathbf R\underline{\Hom}_B(E,B)\simeq E(3k)[-k].
\]
change of rings の導来随伴を適用して二式を比較すると
\eqref{mbcr:eq:finite-duality}を得る。
これは定理\ref{mbsd:thm:gorenstein}の二項解消による計算を
$k$ 回行ったものである。中心性により normalizing twist は生じない。
同じ議論は左右を交換しても成立する。
また任意の $E$-加群に対する change of rings と
$B$ の大域次元 $3k$ から、$E$ の左右単射次元は高々 $2k$ である。
$F=K$ の非零な第 $2k$ Ext により等号となる。

$F_l=E/E_{\geq l}$ は有限次元である。
torsion は
$\Gamma_E(-)=\varinjlim_l\underline{\Hom}_E(F_l,-)$ で計算され、
filtered colimit は完全なので
\[
 R^q\Gamma_E(E)
   =\varinjlim_l\underline{\Ext}^q_E(F_l,E).
\]
\eqref{mbcr:eq:finite-duality}を代入して
\eqref{mbcr:eq:ci-local-cohomology}を得る。
Hilbert 級数は
$(1+t+t^2)^k/(1-t)^{2k}$ なので、ここでの次元は $2k$ である。

切断に移るための局所化三角の入力も確認する。
定理\ref{thm:nfm-relative-line-cohomology}の証明は、
有限 Koszul 解消を持つ $B$ 上で有界幅の derived torsion 模型を構成している。
その中心商への比較は、今回の正則列の各段階で反復できる。
実際、$Q'=Q/(\Omega)$、$\deg\Omega=3$ とし、
$J$ を injective な次数付き右 $Q'$-加群とすると
\[
 \underline{\Ext}^{\,q}_Q(Q/Q_{\geq l},J)
 \simeq
 \underline{\Hom}_{Q'}\!
 \left(\operatorname{Tor}^{Q}_q(Q/Q_{\geq l},Q'),J\right).
\]
正則元の二項解消により $q>1$ の Tor は零である。
第 $1$ Tor は切断の上端にある幅三の次数に限られ、
$l'\geq l+3$ から $l$ への遷移は零となる。
従って忘却した $J$ は $\Gamma_Q$-acyclic であり、
restriction は derived torsion の比較を与える。
これを全段階で反復すれば、$E$ 上でも torsion 複体上で恒等となる
有界幅の derived torsion と Gabriel 局所化三角
\begin{equation}\label{mbcr:eq:localization}
 \mathbf R\Gamma_E(N)\longrightarrow N
 \longrightarrow\mathbf R\omega_E\pi_EN\longrightarrow
\end{equation}
を得る。標準的な局所化三角と切断右随伴の記述については
\cite[\S4, Proposition~4.5 and equation~(4.15)]{NFMBurkeStevenson}も参照されたい。
ここでは上記の有限解消と正則中心列から、その必要な入力を直接確認した。
$N=E$ の内部次数 $b$ に
\eqref{mbcr:eq:ci-local-cohomology}を代入すれば
\eqref{mbcr:eq:ci-low-sections}が従う。
この計算は $\QGr E$ の有限大域次元や smoothness を用いない。
\end{proof}

\subsection{Veronese の根加群に対する負次数 Ext}

\begin{proposition}[根加群の負次数消滅]
\label{mbcr:prop:negative-ext}
任意の整数 $j$ と $n>0$ について
\begin{equation}\label{mbcr:eq:negative-ext}
 \Hom_{\Gr D}(\overline M_j,\overline M_j(-n))=0,\qquad
 \Ext^1_{\Gr D}(\overline M_j,\overline M_j(-n))=0 .
\end{equation}
\end{proposition}
\begin{proof}
補題\ref{lem:nfm-gabriel}の完全同値を
$\mathsf V:\QGr E\xrightarrow{\sim}\QGr D$ と書く。
このとき
\[
 \mathsf V(\mathcal O_E(j))=\pi_D\overline M_j,\qquad
 \mathsf V(\mathcal O_E(-3a))=\mathcal O_D(-a).
\]
切断右随伴を次数 $a$ で評価することは
$\Hom(\mathcal O_D(-a),-)$ に等しい。
従って、右随伴の各次数とその積作用を比較すると、
加群上の Veronese を $\mathsf V_{\Gr}$ として
$\mathsf V_{\Gr}\omega_E\simeq\omega_D\mathsf V$ を得る。
$\mathsf V_{\Gr}$ は完全であり、$\mathsf V$ は完全同値なので injective を保つ。
この関手の同型を導来化すれば、単位射に両立する同型
\begin{equation}\label{mbcr:eq:veronese-sections}
 \mathbf R\omega_D\pi_D\overline M_j
 \simeq
 \left(\mathbf R\omega_E\pi_E E(j)\right)^{(3)}
\end{equation}
を得る。これは切断右随伴の同定であり、
通常の加群 $\Ext$ が Veronese で自動的に保存されるという主張ではない。

式\eqref{mbcr:eq:localization}と
\eqref{mbcr:eq:ci-local-cohomology}により、
$E(j)\to\mathbf R\omega_E\pi_EE(j)$ の cone は
コホモロジー次数 $2k-1$ にのみ成分を持つ。
Veronese は次数付き加群上で完全なので、
\[
 \operatorname{Cone}\!
 \left(\overline M_j\longrightarrow\mathbf R\omega_D\pi_D\overline M_j\right)
 \in D^{\geq2k-1}(\Gr D).
\]
同じことは $\overline M_j(-n)$ にも成立する。
次数零の複体 $\overline M_j$ からの導来 Hom は、この下界を下げない。
従って $q=0,1$ について、導来随伴と上の cone は同型
\begin{align*}
 \Ext^q_{\Gr D}(\overline M_j,\overline M_j(-n))
 &\simeq
 \Ext^q_{\QGr D}(\pi_D\overline M_j,\pi_D\overline M_j(-n))\\
 &\simeq
 \Ext^q_{\QGr E}(\mathcal O_E(j),\mathcal O_E(j-3n))\\
 &\simeq
 \Ext^q_{\QGr E}(\mathcal O_E,\mathcal O_E(-3n))
\end{align*}
を与える。最後の同型は、次数 shift が $\QGr E$ の完全自己同値であることだけを使う。
$\mathcal O_E(j)$ を可換多様体上の線束と仮定していない。
式\eqref{mbcr:eq:ci-low-sections}と $E_{-3n}=0$ が
\eqref{mbcr:eq:negative-ext}を与える。
\end{proof}

\subsection{中心 thickening 上の同型の一意持上げ}

\begin{lemma}[正次数中心に沿った次数付き持上げ]
\label{mbcr:lem:graded-lifting}
$V$ を有限次元 $K$-空間、$R=\Sym_K V$ を $V$ を次数一に置く多項式環とし、
$C$ は中心写像 $R\to C$ を持つ非負次数付き代数で、
$R$ 上平坦であるとする。
$I=R_{>0}C$、$D=C/I$ と置く。
$P,Q$ は下に有界な次数付き右 $C$-加群で、
$R$ 上平坦とし、$\overline P=P/IP$, $\overline Q=Q/IQ$ と書く。
全ての $n>0$ に対して
\[
 \Hom_{\Gr D}(\overline P,\overline Q(-n))
 =\Ext^1_{\Gr D}(\overline P,\overline Q(-n))=0
\]
なら、指定された次数付き同型
$\overline P\xrightarrow{\sim}\overline Q$ は、
ただ一つの次数付き同型 $P\xrightarrow{\sim}Q$ に持ち上がる。
\end{lemma}
\begin{proof}
$\mathfrak m=R_{>0}$、$R_n=R/\mathfrak m^n$、
$C_n=C/I^n$、$P_n=P/I^nP$、$Q_n=Q/I^nQ$ とする。
平坦性により、$C_n$ と $P_n,Q_n$ はそれぞれ $R_n$ 上平坦である。
従って
\begin{equation}\label{mbcr:eq:derived-base-change}
 P_{n+1}\otimes^{\mathbf L}_{C_{n+1}}D
 \simeq
 P_{n+1}\otimes^{\mathbf L}_{C_{n+1}}
       \left(C_{n+1}\otimes^{\mathbf L}_{R_{n+1}}K\right)
 \simeq P_{n+1}\otimes^{\mathbf L}_{R_{n+1}}K
 \simeq\overline P .
\end{equation}
最初の同定には $C_{n+1}$ の平坦性、最後の同定には
$P_{n+1}$ の平坦性を使う。
中心の正則性だけから導来 tensor を通常の tensor に置き換えていない。
導来随伴により、$D$-加群 $L$ について
\[
 \mathbf R\Hom_{\Gr C_{n+1}}(P_{n+1},L)
 \simeq\mathbf R\Hom_{\Gr D}(\overline P,L)
\]
となる。

$Q$ の平坦性は短完全列
\begin{equation}\label{mbcr:eq:thickening-extension}
 0\longrightarrow
 \overline Q(-n)\otimes_K\Sym^nV
 \longrightarrow Q_{n+1}\longrightarrow Q_n\longrightarrow0
\end{equation}
を与える。ここで $\Sym^nV$ は次数零に置いた係数空間であり、
中心単項式の次数 $n$ は $(-n)$ に含めた。
$P_n\simeq Q_n$ が既に与えられたとする。
$P_{n+1}\to P_n\to Q_n$ を持ち上げる障害と持上げの差は、
\eqref{mbcr:eq:thickening-extension}への Hom 完全列により、それぞれ
\[
 \Ext^1_{\Gr C_{n+1}}
  (P_{n+1},\overline Q(-n)\otimes_K\Sym^nV),\qquad
 \Hom_{\Gr C_{n+1}}
  (P_{n+1},\overline Q(-n)\otimes_K\Sym^nV)
\]
に属する。\eqref{mbcr:eq:derived-base-change}と仮定により両者は零である。
従って同型の候補は一意に持ち上がる。

この持上げは同型である。
余核は $\mathfrak m$ で割ると零になり、
$\mathfrak m^{n+1}=0$ により余核そのものが零となる。
全射の核については、$Q_{n+1}$ の $R_{n+1}$-平坦性から
剰余後も左端の完全性が保たれる。
剰余写像が同型なので核も $\mathfrak m$ で割ると零であり、
同じ冪零性により零となる。
これで互いに両立する同型 $P_n\simeq Q_n$ を全ての $n$ で得る。

最後に、形式的同型を実際の次数付き同型へ戻す。
$P,Q$ の次数の下界を固定すると、各内部次数 $a$ について
$n\gg0$ で $(I^nP)_a=(I^nQ)_a=0$ となる。
従ってその次数の同型は実際の $P_a\simeq Q_a$ に安定する。
任意の斉次元による作用の可換性は、始域と終域の両次数で十分大きな
$n$ を取れば有限商上で確認できる。
よって安定した写像は $C$-線型であり、逆も同じ方法で構成される。
一意性も各有限商での一意性から従う。
完備化した加群の存在だけで結論を止めてはいない。
\end{proof}

\subsection{自己同型への適用と中心還元の忠実性}

自己同型の群については、さらに $j_i$ が相異なり $0,1728$ を避けるとする。
位数の仮定は本節冒頭の $p_i,2p_i,3p_i\ne0$ のままでよい。
全 $\mathcal O$-同値群と同定する場合には
定理\ref{mbv:thm:all-o-third}の有限個の捩れ条件も課す。
$\Phi_\alpha$ は右作用を
$m\cdot c=m\alpha^{-1}(c)$ と変更する関手とする。

\begin{theorem}[中心還元からの次数一軌道の認識]
\label{mbcr:thm:reduction-rigidity}
$\alpha\in\Aut_{\mathrm{gr}}(C)$ とし、その中心商上の作用を
$\overline\alpha\in\Aut_{\mathrm{gr}}(D)$ とする。
$j=1,2$ について
\[
 \Phi_{\overline\alpha}(\overline M_j)\simeq\overline M_j
 \quad\text{in }\Gr D
\]
なら
\[
 \Phi_\alpha(M_j)\simeq M_j\quad\text{in }\Gr C .
\]
各剰余同型を指定すると、その持上げは一意である。
従って $\alpha$ の射影的同型類は
$\mathcal K_k=\Aut_{\mathrm{gr}}(A)/K^\times$ に属する。
さらに、中心還元による準同型
\begin{equation}\label{mbcr:eq:faithful-reduction}
 \Aut_{\mathrm{gr}}(C)/K^\times
 \longrightarrow \Aut_{\mathrm{gr}}(D)/K^\times
\end{equation}
は単射である。
これらの結論に $\Aut_{\mathrm{gr}}(C)/K^\times$ の有限性は不要である。
\end{theorem}
\begin{proof}
$\alpha$ は $Z_1$ を保ち、中心多項式環 $R$ の次数付き自己同型を誘導する。
従って補題\ref{mbcr:lem:central-free}により、
$M_j$ と $\Phi_\alpha(M_j)$ はともに $R$ 上平坦で、下に有界である。
後者の剰余は $\Phi_{\overline\alpha}(\overline M_j)$ となる。
指定した剰余同型で二つを同定し、
命題\ref{mbcr:prop:negative-ext}と補題\ref{mbcr:lem:graded-lifting}を適用する。
これで最初の主張を得る。

Veronese 同値で $M_j$ は $\mathcal O_A(j)$ に対応する。
$C$ の次数付き自己同型から来る関手は $(3)$ と可換するので、
$j=1,2$ の対象を保てば全整数の $\mathcal O_A(j)$ を保つ。
定理\ref{mbva:thm:sss}の証明により、その関手は元の次数付き表示の自己同型から来る。
その証明で使うのは点グラフの $p_i,2p_i,3p_i\ne0$ と $j$ の相異性、
および混合二次関係であり、同節の twist 分類用の $6p_i\ne0$ は使わない。
第三 Veronese 上の次数付き自己同型が同じ構造移送を与える核は
grading scalar であるから、$\alpha$ の射影的同型類は $\mathcal K_k$ に属する。

単射性のため、まず $\mathcal K_k$ の $D$ への作用が射影的に忠実であることを確認する。
$h=(h_i)\in\prod_i\mathcal H_i$ の作用が $D_1$ 上のスカラー $\lambda$ だとする。
$i\ne j$ に対し、混合成分
$(T_i)_2\otimes(T_j)_1\subset D_1$ 上では
$(h_i)_2\otimes(h_j)_1=\lambda\operatorname{id}$ である。
二つの可逆写像の tensor 積がスカラーなら各因子もスカラーなので、
全ての $(h_j)_1=\varepsilon_j\operatorname{id}$ となる。
$h_j$ は $\Omega_j$ を固定するため $\varepsilon_j^3=1$ であり、
純成分 $(T_i)_3\subset D_1$ を見ると $\lambda=1$ である。
混合成分の等式は
$\varepsilon_i^2\varepsilon_j=1$、従って $\varepsilon_j=\varepsilon_i$ を与える。
よって $h$ は共通 $\mu_3$ に属し、$\mathcal K_k$ では恒等である。

次に\eqref{mbcr:eq:faithful-reduction}の核の元を取る。
次数スカラーを除くと、$\overline\alpha=\operatorname{id}_D$ となる代表を選べる。
その代表は全ての $\overline M_j$ を保つので、既に示した部分から
$\mathcal K_k$ に属する。上の忠実性によって、この元は恒等である。
\end{proof}

\begin{corollary}[次数一代数への持上げを持つ還元]
\label{mbcr:cor:reduction-lifts}
前定理の $\overline\alpha$ が、$E=\bigotimes_iT_i$ の
次数付き自己同型 $\beta$ の第三 Veronese への制限であるなら、
$\alpha$ の射影的同型類は $\mathcal K_k$ に属する。
\end{corollary}
\begin{proof}
$\Phi_\beta(E(j))\simeq E(j)$ は
$m\mapsto\beta(m)$ による次数付き加群同型である。
第三 Veronese を取ると
$\Phi_{\overline\alpha}(\overline M_j)\simeq\overline M_j$ となるので、
定理\ref{mbcr:thm:reduction-rigidity}を適用する。
\end{proof}

\begin{corollary}[二ブロックの一次中心変形の非消滅]
\label{mbcr:cor:nonzero-class}
$k=2$ とし、$Z(C)_1=Kg$ の生成元を選ぶ。
$C/(g^2)$ が定める次数付き一次変形類は
\[
 \kappa_C\ne0\quad\text{in }HH^2(D,D)_{-1}
\]
である。従って $C$ は中心次数一変数に関して、一次から既に自明でない変形である。
\end{corollary}
\begin{proof}
次の一般的な事実を使う。
標数が二でない体上で、$C$ が連結二次代数、
$g\in Z(C)_1$ が正則、$D=C/(g)$、$D_1\ne0$ とする。
もし $\kappa_C=0$ なら、$D$ 上恒等で
$g\mapsto-g$ となる非自明な射影的自己同型が存在する。

実際、$C/(g^2)$ の積を斉次線型切断で
$\mu_0+g\mu_1$ と書くと、$\mu_1$ は内部次数 $-1$ の
Hochschild cocycle である。
$[\mu_1]=0$ なら、それを与える内部次数 $-1$ の
一 cochain を用いて切断を変更することで、単位と $g$ を保ち、
$D$ 上恒等となる次数付き同型
$D[g]/(g^2)\simeq C/(g^2)$ を得る。
その $D_1$ の像を $C_1$ へ持ち上げ、$C_1=U\oplus Kg$ と書く。
$R_D\subset U\otimes U$ を $D$ の二次関係空間とすると、
各 $r\in R_D$ の $C_2$ における値は $c_rg^2$、$c_r\in K$ となる。
従って $r-c_rg\otimes g$ は $C$ の二次関係である。

これらと $g\otimes u-u\otimes g$ が全二次関係を張る。
実際、関係空間の $U\otimes U$ への射影は $R_D$ へ全射である。
その核に入る関係は
$g\otimes u+v\otimes g+c\,g\otimes g$ と書け、
積が零という条件は $g(u+v+cg)=0$ である。
$g$ の正則性と直和分解から $v=-u$, $c=0$ となる。
よって $U$ を固定し $g$ を $-g$ に送る写像は全二次関係を保ち、
$C$ の次数付き自己同型を与える。
$U\simeq D_1\ne0$ なので、この自己同型は共通 grading scalar ではない。

本節の $C$ は補題\ref{mbv:lem:quadratic}により二次であり、
$g$ の正則性は補題\ref{mbcr:lem:central-free}から従う。
上の自己同型は定理\ref{mbcr:thm:reduction-rigidity}の中心還元の忠実性に反する。
従って $\kappa_C\ne0$ である。
\end{proof}

\begin{remark}[二ブロックと主張の範囲]
$k=2$ では $R=K[g]$、$I=(g)$ であり、
\eqref{mbcr:eq:ci-local-cohomology}の集中次数は四である。
従って中心還元核は有限開に限らず自明である。
一般二次代数に対する高々 $\mu_2$ という有限群の上界や、
一次中心変形類の非消滅をこの結論の仮定として使う必要はない。
後者の非消滅は、逆に系\ref{mbcr:cor:nonzero-class}で結論として得られる。
この系は非消滅を証明するもので、$\kappa_C$ の具体的な Hochschild cocycle や値は計算していない。

ここでの定理は代数閉体上の固定された表示に関する主張である。
$I$-進の非被約商は元の次数付き代数の内部の thickening であり、
任意の非被約パラメータ底上の自己同型層の認識と同一視しない。
また本節だけでは $\Aut_{\mathrm{gr}}(D)$ の全ての元が $E$ へ持ち上がるとは主張しない。
その分類を別途与えれば、系\ref{mbcr:cor:reduction-lifts}が全自己同値群の決定へ接続する。
\end{remark}

% !TEX root = ../cubic-moduli-descent.tex
\section{中心商の自己同型と全構造対象保存群}
\label{mbca:sec:automorphisms}

本節では、中心商の純ブロック商だけでなく、元の多重次数空間を内在的に回復する。
そのために純商の核の冪を右飽和する。
この認識と楕円曲線上の線束の自己準同型環を使うと、
中心商の全次数付き自己同型が元の因子自己同型から来ることが分かる。
中心方向の加群の持上げ剛性と合わせて、
全 $\OO$-自己同値群を既知の有限 Heisenberg 商群へ同定する。

基礎体は標数零の代数閉体 $K$、ブロック数は $k\ge2$ とする。
第\ref{mbiq:sec:quotients}節の記号により
\[
 T_i=B(E_i,L_i,\sigma_i),\qquad
 E=\bigotimes_{i=1}^kT_i,\qquad
 D=E^{(3)},\qquad U_i=T_i^{(3)},\qquad F=\bigotimes_iU_i
\]
と置く。$L_i$ の次数は三、
$\sigma_i=t_{p_i}$ は $p_i,2p_i,3p_i\ne0$ を満たす平行移動とし、
$j(E_i)$ は相異なり $0,1728$ を避ける。
各 $T_i$ は次数一生成の domain である。
本節の自己同型の計算には $\sigma_i$ の無限位数性を要求しない。
群の同定はまず固定した代数閉体上の対象について述べる。

多重次数 $\boldsymbol a=(a_1,\ldots,a_k)$ に対して
$T_{\boldsymbol a}=\bigotimes_i(T_i)_{a_i}$、
$|\boldsymbol a|=\sum_i a_i$ と書く。
$D_n=\bigoplus_{|\boldsymbol a|=3n}T_{\boldsymbol a}$ である。
純商 $q_i:D\to U_i$ の核を $I_i$ とする。
これらは系\ref{mbiq:cor:intrinsic-quotients}によって内在的であり、
$j$ 不変量の相異性から全自己同型が各 $I_i$ を保つ。

\subsection{飽和冪による多重次数の回復}

次数付き右イデアル $J\subset D$ の右飽和を
\[
 J^{\mathrm{sat}}=
 \{x\in D:xD_{\ge r}\subset J\text{ for some }r\ge0\}
\]
と定める。以下では $I_i^0=D$ とする。
全次数付き自己同型は、イデアルの積とこの右飽和を保つ。

\begin{lemma}[純商の核の飽和冪]
\label{mbca:lem:saturated-powers}
任意の $i$、$m\ge0$ について
\begin{equation}\label{mbca:eq:saturated-powers}
 (I_i^m)^{\mathrm{sat}}\cap D_n
   =\bigoplus_{\substack{|\boldsymbol a|=3n\\
                        \sum_{j\ne i}a_j\ge m}}T_{\boldsymbol a}.
\end{equation}
従って、$|\boldsymbol a|=3n$ に対し
\begin{equation}\label{mbca:eq:multidegree-recovery}
 T_{\boldsymbol a}
  =D_n\cap\bigcap_{i=1}^k
                (I_i^{\,3n-a_i})^{\mathrm{sat}}.
\end{equation}
特に、全次数付き自己同型は各多重次数空間と純部分代数
$U_i\subset D$ を保つ。
\end{lemma}
\begin{proof}
$m=0$ は明らかなので $m>0$ とする。
まず $n\ge m$ において
\[
 (I_i^m)_n
  =\bigoplus_{\substack{|\boldsymbol a|=3n\\
                       \sum_{j\ne i}a_j\ge m}}T_{\boldsymbol a}
\]
を示す。左辺の元は $i$ 以外の次数が少なくとも $m$ なので一方の包含は従う。
逆にその outside 次数が $m$ 以上なら、
多重次数を表す $3n$ 個の文字から outside の文字を $m$ 個選び、
それぞれを異なる三文字の組へ入れる。
$n\ge m$ なので残りの文字でこれらの組を満たし、余りをさらに三文字ずつに分けられる。
最初の $m$ 組に対応する多重次数は全て $(I_i)_1$ に属する。
各 $T_j$ の次数一生成性から、指定した次数分割に対する切断乗法は全射である。
従って $T_{\boldsymbol a}\subset(I_i^m)_n$ となる。

outside 次数が $m$ 以上の成分は、十分高い次数を右から掛けると
この等式の適用範囲へ入るため、右飽和に含まれる。
逆に outside 次数が $m$ 未満の非零成分を持つ斉次元 $x$ を取る。
任意に高い次数の非零な純元 $u\in(U_i)_r$ を右から掛ける。
$T_i$ は domain なので、その純因子における右乗法は単射である。
他の因子との tensor 積でも単射であり、
$x$ の異なる多重次数成分どうしは掛けた後も異なる。
従って $xu$ は outside 次数が $m$ 未満の非零成分を持ち、
$I_i^m$ に入らない。
これは $x\notin(I_i^m)^{\mathrm{sat}}$ を示す。

式\eqref{mbca:eq:saturated-powers}で $m=3n-a_i$ を取ると、
許される多重次数 $\boldsymbol b$ の条件は $b_i\le a_i$ となる。
全ての $i$ についてこの条件を課し、$|\boldsymbol b|=|\boldsymbol a|$ を使えば、
$\boldsymbol b=\boldsymbol a$ である。
これで式\eqref{mbca:eq:multidegree-recovery}を得る。
\end{proof}

ここでは有限長の切断点スキームの非被約構造を使っていない。
幾何学的点加群が与える $I_i$ と、代数内のイデアル演算・右飽和だけから
多重次数を認識した。

\subsection{剰余多重次数部分の自己準同型}

$0\le e\le2$ に対し、次数付き右 $U_i$-加群を
\[
 P_i(e)=(T_i(e))^{(3)},\qquad
 P_i(e)_n=(T_i)_{3n+e}\quad(n\ge0)
\]
と置く。これは次数零で生成される。

\begin{lemma}[剰余次数加群の次数零自己準同型]
\label{mbca:lem:residue-end}
$\operatorname{End}_{\mathrm{gr}\,U_i}(P_i(e))=K$ である。
さらに任意の $\boldsymbol e\in\{0,1,2\}^k$ について、
$\bigotimes_iP_i(e_i)$ の多重次数を保つ右 $F$-自己準同型はスカラーに限られる。
\end{lemma}
\begin{proof}
Veronese 同値と THCR 同値の下で、$P_i(e)$ の tails 像は
$\operatorname{coh}E_i$ の線束となる。
具体的には $T_i(e)$ の像であり、
右加群の shift は式\eqref{mbiq:eq:thcr-shift}の自己同値である。
この既知の THCR 同値は、曲線上の正次数線束の ampleness により
有限位数の場合にも使える
\cite[\S3, equations~(3.1), (3.3), Theorem~3.2]{NFMStaffordVdB}。
滑らかな幾何学的連結射影曲線上の線束の自己準同型環は $K$ である。
また $P_i(e)$ は $T_i$ の成分からなるため右 torsion を持たない。
従って graded 自己準同型から tails 自己準同型への写像は単射であり、
スカラー写像が既に全てを与える。

tensor 積については二因子の場合を示して帰納する。
$P,Q$ を二つの加群とし、多重次数を保つ $P\otimes Q$ の自己準同型を考える。
$Q_0$ の基底を選び、$P\otimes Q_0$ へ制限すると、
各行列成分は $P$ の次数零自己準同型である。
従ってこの制限は $1_P\otimes b_0$ の形となる。
$Q$ が $Q_0$ から生成されるので、右 $U_Q$-線型性は
$b_0$ が $Q$ 全体の次数零自己準同型 $b$ へ延びることを要求する。
実際、$Q_0$ の関係へ一つの非零な $p\in P$ を tensor すれば、
$b_0$ が同じ関係を保つことが分かる。
$b$ もスカラーであるから全体がスカラーとなる。
同じ議論を残りの因子へ繰り返す。
\end{proof}

剰余多重次数群を
\[
 \Lambda=\left\{\bar{\boldsymbol e}\in(\mathbf Z/3)^k:
                    \sum_i\bar e_i=0\right\}
\]
とする。各元の代表を $0\le e_i\le2$ に取れば、
$F$-双加群として
\[
 D=\bigoplus_{\bar{\boldsymbol e}\in\Lambda}D_{\bar{\boldsymbol e}},
 \qquad
 D_{\bar{\boldsymbol e}}=
 \bigoplus_{\boldsymbol q\in\mathbf N^k}
       \bigotimes_i(T_i)_{3q_i+e_i}.
\]
多重次数を $3\boldsymbol q+\boldsymbol e$ で標識すると、
右 $F$-加群 $D_{\bar{\boldsymbol e}}$ は $\bigotimes_iP_i(e_i)$ である。
特に $D_{\bar0}=F$ である。

\begin{proposition}[純部分を点ごと固定する自己同型]
\label{mbca:prop:pure-kernel}
各 $U_i\subset D$ を点ごと固定する次数付き自己同型は、
剰余次数部分 $D_{\bar{\boldsymbol e}}$ 上で
$\lambda(\bar{\boldsymbol e})$ 倍として作用する。
ここで
\[
 \lambda\in\operatorname{Hom}(\Lambda,\mu_3)
       \cong\mu_3^k/\mu_3^{\mathrm{diag}}.
\]
従ってこれらの自己同型は、全てブロック別の三乗根スカラーから来る。
\end{proposition}
\begin{proof}
補題\ref{mbca:lem:saturated-powers}により自己同型は多重次数を保ち、
$F$ を点ごと固定する。
従って各 $D_{\bar{\boldsymbol e}}$ 上の作用は
多重次数を保つ $F$-加群自己同型である。
補題\ref{mbca:lem:residue-end}により、その作用は一つの非零スカラー
$\lambda(\bar{\boldsymbol e})$ である。

任意の二つの剰余次数部分には、積が非零となる元がある。
各因子で非零な斉次元を選び、その tensor を取ればよい。
各 $T_i$ の domain 性が、その積の非零性を保証する。
積は和の剰余次数部分に入るので、
\[
 \lambda(\bar{\boldsymbol e}+\bar{\boldsymbol f})
    =\lambda(\bar{\boldsymbol e})\lambda(\bar{\boldsymbol f})
\]
となる。$\Lambda$ の指数が三なので値は $\mu_3$ に属する。
逆に任意のこのような指標は積を保つ。
標準 pairing
$\mu_3^k\times\Lambda\to\mu_3$ の左側の核は共通 $\mu_3$ であり、
全ての指標を与える。これがブロックスカラーによる表示である。
\end{proof}

\subsection{中心商の全自己同型群}

\begin{theorem}[中心商の自己同型は因子へ持ち上がる]
\label{mbca:thm:center-automorphisms}
制限写像は完全列
\begin{equation}\label{mbca:eq:aut-d}
 1\longrightarrow\mu_3^{\mathrm{diag}}
 \longrightarrow\prod_{i=1}^k\operatorname{Aut}_{\mathrm{gr}}(T_i)
 \longrightarrow\operatorname{Aut}_{\mathrm{gr}}(D)
 \longrightarrow1
\end{equation}
を与える。
従って、$D$ の次数スカラーによる商は
\begin{equation}\label{mbca:eq:projective-aut-d}
 H_D:=\operatorname{Aut}_{\mathrm{gr}}(D)/K^\times
   \cong
   \left(\prod_i\operatorname{Aut}_{\mathrm{gr}}(T_i)\right)/
                                         \mathbf G_m^{\mathrm{diag}}.
\end{equation}
特に、$D$ の全次数付き自己同型は $E=\bigotimes_iT_i$ の
次数付き自己同型へ持ち上がる。
\end{theorem}
\begin{proof}
$\alpha\in\operatorname{Aut}_{\mathrm{gr}}(D)$ を取る。
補題\ref{mbca:lem:saturated-powers}により、純部分代数 $U_i$ も保たれる。
定理\ref{mbiq:thm:triple-recovery}から、
その作用は $(E_i,L_i,\sigma_i)$ の自己同型 $f_i$ を与える。
線束同型 $f_i^*L_i\simeq L_i$ を一つ選ぶと、
THCR の構成の自然性により $h_i\in\operatorname{Aut}_{\mathrm{gr}}(T_i)$ が得られる。

$\alpha|_{U_i}$ と $h_i|_{U_i}$ は、点曲線とその次数九偏極に同じ射影作用を与える。
従ってその比は $U_i$ の次数スカラーである。
実際、次数一の射影作用が曲線上で恒等なら、
曲線の既約性と線型非退化性からその線型写像はスカラーであり、
$U_i$ の次数一生成性が全次数の作用を決める。
このスカラーの三乗根を選んで $h_i$ を grading scalar 倍すれば、
$h_i|_{U_i}=\alpha|_{U_i}$ とできる。

$\alpha$ から $\bigotimes_i h_i$ の制限を除くと、
各 $U_i$ を点ごと固定する自己同型となる。
命題\ref{mbca:prop:pure-kernel}により、それもブロックスカラーから来る。
これで式\eqref{mbca:eq:aut-d}の全射性を得る。

核を調べる。各 $h_i|_{U_i}=1$ なら、
点曲線の作用が恒等なので $h_i$ は $T_i$ の次数スカラー
$\varepsilon_i$ であり、$\varepsilon_i^3=1$ となる。
その $D$ 上の作用が恒等となるのは
$\Lambda$ への指標が自明な場合、すなわち全ての $\varepsilon_i$ が同じ場合に限る。
これが表示した核である。
最後に、$T_i$ の共通 grading scalar $t$ は
$D$ の新しい次数に対して $t^3$ の grading scalar を与える。
代数閉体上では三乗写像が全射なので、
この共通スカラーを除けば式\eqref{mbca:eq:projective-aut-d}となる。
\end{proof}

\begin{corollary}[相対ブロックトーラスと有限部分]
\label{mbca:cor:center-torus}
$\mathcal T=\mathbf G_m^k/\mathbf G_m^{\mathrm{diag}}$ は $H_D$ の中心部分群であり、
\[
 1\longrightarrow\mathcal T
 \longrightarrow H_D
 \longrightarrow\prod_iE_i[3]\longrightarrow1
\]
となる。また既知の有限表示群 $\mathcal K_k$ の像について
\[
 H_D=\mathcal T\,\mathcal K_k,\qquad
 \mathcal T\cap\mathcal K_k
       =\mu_3^k/\mu_3^{\mathrm{diag}}.
\]
従って $\mathcal K_k$ は $H_D$ の正規部分群である。
二ブロックでは $H_D$ の次元は一で、連結成分群の位数は $81$ である。
\end{corollary}
\begin{proof}
$j(E_i)\ne0,1728$ と $2p_i\ne0$ より、
次数三偏極を保ち $\sigma_i$ と可換する曲線自己同型は $E_i[3]$ の平行移動である。
従って $\operatorname{Aut}_{\mathrm{gr}}(T_i)$ は
$E_i[3]$ の、grading scalar による中心拡大である。
これらは元の Sklyanin 因子へも持ち上がり、
その中心三次指標を三乗根スカラーで規格化すると $\mathcal H_i$ に入る。
定理\ref{mbca:thm:center-automorphisms}と既存の Heisenberg 群の記述を使えば、
表示した全ての式が従う。
\end{proof}

\subsection{全自己同値群と構造対象の像を固定した根}

\begin{theorem}[満支持型の全構造対象保存群]
\label{mbca:thm:full-inertia}
$A=(\bigotimes_iS_i)/(\sum_i a_i\Omega_i)$、全 $a_i\ne0$ とし、
$C=A^{(3)}$ とする。上記の基本仮定の下で
\begin{equation}\label{mbca:eq:full-inertia}
 \operatorname{Aut}_{\mathrm{gr}}(C)/K^\times
       =\mathcal K_k
       =\Bigl(\prod_i\mathcal H_i\Bigr)/\mu_3^{\mathrm{diag}}
\end{equation}
である。
さらに全 $\OO$-同値の第三 Veronese 認識の有限位数条件
\[
 jp_i\ne0\qquad(1\le i\le k,\ 1\le j\le3k-3)
\]
を課すと、
\[
 \operatorname{Aut}_{\OO}(\operatorname{QGr}A)=\mathcal K_k.
\]
その位数は $3^{3k-1}$、二ブロックでは $243$ である。
ここでは事前に全自己同値群の有限性を仮定しない。
\end{theorem}
\begin{proof}
$\alpha\in\operatorname{Aut}_{\mathrm{gr}}(C)$ を取る。
中心一次生成イデアル $I=(Z(C)_1)$ は $\alpha$ によって保たれ、
$D=C/I$ に自己同型 $\bar\alpha$ を誘導する。
定理\ref{mbca:thm:center-automorphisms}により、
$\bar\alpha$ は $E=\bigotimes_iT_i$ の次数付き自己同型 $\beta$ へ持ち上がる。

$j=1,2$ に対し
$M_j=(A(j))^{(3)}$、$\bar M_j=(E(j))^{(3)}$ と置く。
$\beta$ の構造移送は $E(j)$ と同型なので、
$\bar\alpha$ の構造移送は $\bar M_j$ と同型である。
定理\ref{mbcr:thm:reduction-rigidity}により、
この同型は $M_j$ の同型へ持ち上がる。
従って $\alpha$ の tails 関手は $\OO(1),\OO(2)$ を保つ。
$C$ の次数シフト、すなわち $A$ の $(3)$ とも可換するので、
全整数の $\OO(j)$ を保つ。

定理\ref{mbva:thm:sss}の隣接 Hom と混合二次関係による認識を適用すると、
この関手は $A$ の次数付き自己同型から来る。
その認識の証明は $p_i,2p_i,3p_i\ne0$ と $j$ の相異性だけを使い、
$6p_i\ne0$ は使わない。
$C$ の自己同型の構造移送の核は共通 grading scalar なので、
式\eqref{mbca:eq:full-inertia}を得る。
逆の包含と忠実性は命題\ref{mbva:prop:faithful}で既に示されている。
最後に定理\ref{mbv:thm:all-o-third}を適用する。
\end{proof}

\begin{corollary}[第一対象を保つ三乗根の完全な分類]
\label{mbca:cor:marked-roots}
全 $\OO$-同値の第三 Veronese 認識が成立する範囲で
$\mathsf s=(1)$ とする。
固定した圏 $\operatorname{QGr}A$ の自己同値 $\mathsf r$ が
\[
 \mathsf r\OO\simeq\OO(1)
\]
を満たすなら、一意な $h\in\mathcal K_k$ に対して
$\mathsf r\simeq\mathsf s\Phi_h$ となる。
従って
\[
 \mathsf r^3\simeq\mathsf s^3,\qquad
 \mathsf r^j\OO\simeq\OO(j)\quad(j\in\mathbf Z)
\]
も自動である。
この条件を持つ根の自然同型類は $3^{3k-1}$ 個であり、
全 $\OO$-自己同値による共役 groupoid は
$[\mathcal K_k/\mathcal K_k]_{\mathrm{conj}}$ である。
同型類数は $(11^k+2\cdot8^k)/3$、
二ブロックでは $243$ 個の根・$83$ 個の共役類である。
\end{corollary}
\begin{proof}
$\mathsf h=\mathsf s^{-1}\mathsf r$ は $\OO$ を保つので、
定理\ref{mbca:thm:full-inertia}により一意な $h\in\mathcal K_k$ に対応する。
$\Phi_h$ は $\mathsf s$ と可換し、$\mathcal K_k$ は指数三なので、
表示した三乗と全軌道の式を得る。
共役作用と個数は第\ref{mbr:sec:roots}節の有限群計算による。
\end{proof}

この根の分類は $\mathsf r\OO\simeq\OO(1)$ を仮定している。
$\mathsf r^3\simeq\mathsf s^3$ だけを仮定した全ての根について、
構造対象の像を先に決めたことにはならない。
数えているのは自己同値の自然同型類であり、
三乗同型 $\mathsf r^3\simeq\mathsf s^3$ の選択やその coherence を追加標識としてはいない。
また、本節の証明から任意の非被約基底上の全同型層の一致を無条件に主張しない。
THCR 同値、有限 Heisenberg 表現と中心方向の持上げ剛性を入力として用い、
新たに行ったのは飽和冪による多重次数認識と剰余次数加群からの全自己同型の回収である。
これらの組合せについての新規性の網羅的確認は未完了である。

% !TEX root = ../cubic-moduli-descent.tex
\section{満支持型の全構造対象保存モジュライの非空開}
\label{mbg:sec:generic-moduli}

本節では、全 $\OO$-自己同値群の具体的な群構造を入力とせずに、
その一般有限性と、満支持局所型のモジュライの非空開を構成する。
第三 Veronese の有限表示、中心商の余法加群による幾何点の認識、
可換三次への特殊化を組み合わせる。
構成する対象は、指定した満支持族に fpqc 局所的に同値な
$\OO$-marked Abel 圏族の stack である。
任意の非可換三次の変形がこの局所表示を持つという主張は含めない。
具体的な慣性群は定理\ref{mbca:thm:full-inertia}が決定し、
さらに系\ref{mbrr:cor:stack-isomorphism}が表示スタックとの一般相対認識を与える。
以下の可換特殊化による有限性証明は、これらの群の完全な同定に依存しない。

\subsection{有限な関係行列と自己同値群の接空間}

補題\ref{mbv:lem:quadratic}の状況で、
\[
 C=A^{(3)}=T(V)/(R_2),\qquad V=A_3,\qquad
 R_2=\ker(V\otimes V\longrightarrow A_6)
\]
と書く。係数基底上で $V,R_2,C_2$ は有限局所自由である。
次の有限行列を考える：
\begin{equation}\label{mbg:eq:derivation-matrix}
 \overline\Delta_C:\End(V)/\OO\operatorname{id}
 \longrightarrow\Hom(R_2,C_2),\qquad
 [u]\longmapsto\mu(u\otimes1+1\otimes u)|_{R_2}.
\end{equation}
次数スカラーは関係を保つので、この式は商上で定義される。

\begin{lemma}[自己同値の接空間と開条件]
\label{mbg:lem:tangent-open}
幾何点での $\ker\overline\Delta_C$ は
$\Aut_{\mathrm{gr}}(C)/\mathbf G_m$ の単位元における接空間である。
この核が零となる条件は表示基底上の Zariski 開条件であり、
全 $\OO$-同値によって不変である。
この開上の各幾何点では全 $\OO$-自己同値群は有限 \'etale であり、
任意の二つの族の間の全 $\OO$-同値スキームは有限表示かつ不分岐・準有限である。
最後の二つの性質は、同値スキームから二つの族の共通基底への射について述べる。
\end{lemma}
\begin{proof}
$1+\epsilon u$ が $R_2$ を保つ条件は、
$(u\otimes1+1\otimes u)R_2\subset R_2$ である。
これが \eqref{mbg:eq:derivation-matrix} の核であり、
二次表示によって高次の条件を別に加える必要はない。
有限局所自由束間の写像が各ファイバーで単射となる条件は、
所定の最大小行列式の非消滅として開である。
定理\ref{mbv:thm:all-o-third}により、この条件は全 $\OO$-同値で不変となる。

射影的自己同型群は $\PGL(V)$ の閉部分群スキームで有限型である。
標数零の有限型群スキームは滑らかである
\cite[Tag~047N]{MBGStacks}。
接空間が零なら次元零、従って有限 \'etale となる。
二模型間の射影的 Isom の幾何ファイバーは、空でなければ
一方の自己同型群の torsor なので有限 \'etale である。
その表現可能性と有限表示性は補題\ref{mbv:lem:quadratic}から従う。
ファイバーによる不分岐性の判定と不分岐射の準有限性
\cite[Tag~02G3, Lemmas~29.36.10 and~29.36.12]{MBGStacks}を適用する。
この議論は相対 Isom の平坦性や有限性までは与えない。
\end{proof}

\begin{lemma}[滑らかな可換三次の切断環の導分]
\label{mbg:lem:commutative-tangent}
標数零の体上の滑らかな三次超曲面
$X=V(F)\subset\PP^{N-1}$、$N\ge6$ に対し、
$C_0=(k[x_1,\ldots,x_N]/(F))^{(3)}$ と置く。
このとき
\[
 \operatorname{Der}_0(C_0)=k\cdot\operatorname{Eu},
 \qquad \operatorname{Eu}|_{(C_0)_q}=q\operatorname{id}
\]
である。
\end{lemma}
\begin{proof}
超曲面の完全列と射影空間のコホモロジーから
$(C_0)_q=H^0(X,\OO_X(3q))$、$H^1(X,\OO_X)=0$ を得る。
次数零の導分は $X=\Proj C_0$ 上の vector field を誘導する。
その核を確認する。誘導された vector field が零なら、
非零な次数一切断 $s,t$ の比に対して
$\delta(s)t-s\delta(t)=0$ となる。
基点のない $\OO_X(3)$ の局所 frame ごとに $\delta(s)/s$ は正則で、
重なり上で一致する。従ってこれは $H^0(X,\OO_X)=k$ の元である。
次数一生成性により $\delta$ は Euler 導分のスカラー倍となる。

次に $H^0(X,T_X)=0$ を示す。
制限 Euler 列により、ambient の vector field は
線型形式 $l_i$ による $\sum_i l_i\partial_{x_i}$ をスカラーで割ったものである。
normal 列で $T_X$ に入る条件は
$\sum_i l_iF_{x_i}=cF$ である。
Euler 恒等式を用いて $l_i$ から $(c/3)x_i$ を引けば
\[
 \sum_i\bigl(l_i-(c/3)x_i\bigr)F_{x_i}=0.
\]
滑らかさと Euler 恒等式により、偏微分 $F_{x_i}$ の共通零点は
affine 原点だけである。従ってこれらの二次式は多項式環の正則列をなし、
第一 syzygy は Koszul syzygy で生成される。
後者の係数の次数は二であり、次数一係数の非零 syzygy はない。
よって全ての $l_i-(c/3)x_i$ が零で、射影 vector field も零となる。
以上から主張が従う。
\end{proof}

\begin{proposition}[満支持一般点での有限自己同値]
\label{mbg:prop:finite-open}
第\ref{mbs:sec:presentations}節の既約な表示スタック $\mathfrak P_k$ で、
定理\ref{mbv:thm:all-o-third}の有限位数条件を課した開の中に、
全 $\OO$-自己同値群が有限 \'etale となる非空開 $\mathfrak P_k^{\mathrm{fin}}$ が存在する。
\end{proposition}
\begin{proof}
開性は補題\ref{mbg:lem:tangent-open}によるので、一般点での非空性を示す。
$\alpha_i,\beta_i$ を代数的に独立な $2k$ 個の変数とし、
\[
 k_0=\C(\alpha_1,\beta_1,\ldots,\alpha_k,\beta_k),\qquad
 R=k_0[h]_{(h)}
\]
と置く。各因子を巡回関係
\begin{equation}\label{mbg:eq:degeneration}
 (1+h\alpha_i)yz+(-1+h\alpha_i)zy+h\beta_i x^2=0
\end{equation}
とその巡回置換で定める。
一般ファイバーは一般の楕円 Sklyanin 代数で、特殊ファイバーは三変数多項式環である。
両者の Hilbert 級数が $(1-t)^{-3}$ なので、
命題\paperref{main}{gre:prop:input}と同じ DVR 上の構造定理により、全次数成分は有限自由となる。
中心三次の一般ファイバーの直線と $(S_i)_3$ の交叉を取り、
その primitive な生成元を $\Omega_i$ とする。
交換子の標的 $(S_i)_4^{\oplus3}$ が自由なので、
中心核は飽和部分加群であり、この生成元の剰余は非零である。

一次の Poisson 括弧は、符号を除けば
\[
 F_i=\frac{\beta_i}{3}(x^3+y^3+z^3)+2\alpha_i xyz
\]
の Jacobian 括弧である。$f_i=\Omega_i\bmod h$ はその三次 Casimir だから
$\nabla f_i\wedge\nabla F_i=0$ である。
Euler 恒等式を使うと $F_i\,df_i-f_i\,dF_i=0$、すなわち
$d(f_i/F_i)=0$ となり、標数零では $f_i=\lambda_iF_i$、
$\lambda_i\in k_0^\times$ である。
定数倍を調整して $\Omega_i\bmod h=F_i$ とできる。
各 $F_i$ は滑らかなので、分離した和 $F=\sum_iF_i$ も
$\PP^{3k-1}_{k_0}$ の滑らかな三次超曲面を定める。

$B_R=\bigotimes_iS_i$、$A_R=B_R/(\sum_i\Omega_i)$ と置く。
剰余の $F$ は非零因子であり、次数ごとの自由性から
この中心商も次数別有限自由となる。
具体的に、$w b\in hB_R$ なら剰余の正則性から $b\in hB_R$ なので、
各商次数成分は $h$-torsion を持たない。
正則性自体も同じ議論と有限自由加群の $h$-進分離性から従う。
従って $C_R=A_R^{(3)}$ の二次関係束と乗法核は底変更に両立する。
特殊ファイバーでは補題\ref{mbg:lem:commutative-tangent}により
$\overline\Delta_{C_0}$ が単射である。
その最大小行列式の一つは単元だから、一般ファイバーでも単射となる。

最後に、この一般ファイバーが満支持表示の一般点を検出することを確認する。
Hesse 座標で $[a:b:c]$ を用い、$a-b\ne0$ で $a-b=2$ と規格化すれば、
$a=1+u$, $b=-1+u$, $c=v$ となる。
式\eqref{mbg:eq:degeneration}では $u=h\alpha_i$, $v=h\beta_i$ であり、
$h$ を可逆にすると各因子の二つの独立なパラメータを与える。
楕円 Sklyanin 因子の Hesse 表示と座標変換により、この族は表示の一般点を支配する
\cite[\S1]{GRSklyaninST}。
中心係数の全非零部分はブロックスカラーの一軌道なので、
和の係数を一にしたことは表示モジュライの一般性を失わない。
従って上で示した接空間の消滅は $\mathfrak P_k$ の一般点で成立する。
有限個の捩れ条件を除く既存の非空開と交わらせればよい。
\end{proof}

\subsection{関係 Grassmann 多様体による一般局所型の効果化}

定理\ref{mbca:thm:full-inertia}は、基本開上の全第三 Veronese 自己同型群を
$\mathcal K_k$ と同定するので、全 $\OO$-認識の有限位数条件を課した開では
全 $\OO$-慣性も常に $\mathcal K_k$ である。
従って命題\ref{mbg:prop:finite-open}の有限性は、その結果からも従う。
ここでは群の同定とは独立な導分行列と退化の議論を保持し、
関係 Grassmann 多様体上で開を選ぶために用いる。

以下では $\mathcal X\subset\mathcal X_k$ を、前命題と
系\ref{mbiq:cor:commuting-o-rigidity}の仮定を満たす非空開とする。
元の群 $G_k$ に不変な開を取る。
$m=\rank A_3$、$s=m^2-\rank A_6$ とし、
\[
 \mathcal F=\operatorname{Fr}_{\mathcal X}(A_3),\qquad
 \mathcal G=\Gr(s,\C^m\otimes\C^m),\qquad H=\PGL_m
\]
と置く。第三 Veronese の関係核は $\GL_m$-同変な射
\begin{equation}\label{mbg:eq:relation-map}
 \phi:\mathcal F\longrightarrow\mathcal G
\end{equation}
を定める。標的への作用は $H$ を経由し、frame の共通 scalar は関係空間を変えない。
$Z$ をこの像の被約閉包とする。$Z$ は既約で $H$-不変である。

\begin{theorem}[全 $\OO$-局所型の一般 DM stack]
\label{mbg:thm:generic-stack}
$Z$ に $H$-不変な非空開 $Z^\circ$ が存在し、
\begin{equation}\label{mbg:eq:generic-stack}
 \mathfrak M_k^\circ=[Z^\circ/\PGL_m]
\end{equation}
は有限型・滑らか・既約の Deligne--Mumford stack で、次元は $2k$ である。
この stack は、対応する非空開 $\mathcal X^\circ\subset\mathcal X$ の
満支持表示に fpqc 局所的に同値な $\OO$-marked Abel 圏の局所型を、
全 $\OO$-保存同値を射として分類する。
幾何点での慣性群は
\[
 G_A=\Aut_{\mathrm{gr}}(A^{(3)})/\mathbf G_m
\]
であり、有限 \'etale である。定理\ref{mbca:thm:full-inertia}と合わせると
$G_A=\mathcal K_k$、$|G_A|=3^{3k-1}$ と具体的に決定される。
また、次数一の表示による幾何点の同型類と、この stack の幾何点の同型類は一致する。
この構成だけから任意基底上の二つの stack の同型を結論せず、
その相対認識には系\ref{mbrr:cor:stack-isomorphism}を用いる。
\end{theorem}
\begin{proof}
まず像の効果化を行う。
$\mathcal F\to Z$ は有限型の支配射なので、Chevalley の定理により
その像は非空開を含む。generic flatness によってさらに縮めれば、
$\mathcal F\times_Z Z_0\to Z_0$ は有限表示・忠実平坦となる
\cite[Tags~054K and~052A]{MBGStacks}。
$Z$ の smooth locus、および式\eqref{mbg:eq:derivation-matrix}の
核が零となる開と交わらせる。いずれも一般点を含む。
得られた開の全 $H$-平行移動の和を $Z^\circ$ とすれば、
これは $H$-不変な非空開である。$H$ の元は fpqc 局所的に $\GL_m$ へ持ち上がるので、同変性により
\[
 \mathcal F^\circ:=\mathcal F\times_Z Z^\circ\longrightarrow Z^\circ
\]
は依然として忠実平坦・有限表示である。
Noetherian 性によりこの射は準コンパクトであり、fppf 被覆となる。
$\mathcal F^\circ$ は $\GL_m$-不変な開なので、frame torsor を通じて
非空開 $\mathcal X^\circ\subset\mathcal X$ へ降下する。

Grassmann 多様体上の普遍二次代数を $Z^\circ$ に制限する。
その各次数成分は、上の fppf 被覆後に $A^{(3)}$ の各次数成分となる。
従って全次数の局所自由性と底変更は忠実平坦降下によって従う。
無限個の平坦性開集合を交わしているのではなく、
一つの fppf 被覆上で既知の全次数平坦性を降ろしている。
これにより任意の試験基底、特に非被約底上でも、
$Z^\circ$ の族は fpqc 局所的に指定された満支持表示を持つ。

この局所表示の間の全 $\OO$-同値は、
定理\ref{mbv:thm:all-o-third}によって第三 Veronese の
射影的な次数付き同型に等しい。
二次表示により、その同型はちょうど $\PGL_m$ による関係空間の移送である。
従って局所表示とその全同値の fpqc stack 化は
式\eqref{mbg:eq:generic-stack}と同値である。
次数スカラーは tails 上で標準的に恒等関手と同型であり、
その同型は積に両立する。従って $\GL_m$ の frame を選ぶ必要はなく、
射影的 frame の降下で構造対象も降下する。
実際の大域次数付き代数の存在を別途仮定してはいない。
対象は局所 tails 圏と同値の descent datum として構成でき、
その射・合成・係数変更の降下は命題\ref{prop:nfm-equivalence-descent}に従う。

$Z^\circ$ は滑らかで、$H$ も滑らかなので、商 stack は滑らかな有限型代数 stack である。
関係空間の安定化子は $G_A$ そのもので、補題\ref{mbg:lem:tangent-open}により
有限 \'etale である。従ってこの商は Deligne--Mumford となる
\cite[Tag~0DSL]{MBGStacks}。既約性は $Z$ の既約性から従う。

次元を計算する。$\dim\mathcal X=11k-1$ なので
$\dim\mathcal F=11k-1+m^2$ である。
幾何点 $z\in Z^\circ$ を固定する。
系\ref{mbiq:cor:commuting-o-rigidity}により、$z$ と同じ第三 Veronese を持つ
$\mathcal X$ の表示は一つの次数付き同型類をなし、そのパラメータは
$G_k$-軌道である。表示の幾何慣性が有限なので、この軌道の次元は $9k-1$。
各表示上で $z$ を与える線型 frame の空間は、
$\Aut_{\mathrm{gr}}(C)$ の torsor であり、次元は一である。
従って $\phi$ の幾何ファイバーの次元は $9k$ で、
\[
 \dim Z^\circ=(11k-1+m^2)-9k=m^2+2k-1.
\]
$\dim\PGL_m=m^2-1$ を引くと $\dim\mathfrak M_k^\circ=2k$ となる。
幾何点の同型類についての最後の主張も同じ対象剛性から従う。
\end{proof}

\begin{remark}[得られた開と残る問題]
\label{mbg:rem:scope}
この構成は非空開の存在を与えるが、$Z$ の全境界や全満支持表示を覆うとは主張しない。
開を切り出す最大小行列式・flat locus を明示的に計算したわけでもない。
全 $\OO$-慣性は定理\ref{mbca:thm:full-inertia}により表示慣性と一致し、
二ブロックでは位数 $243$ の $\mathcal K_2$ である。
さらに定理\ref{mbrr:thm:finite-etale}は表示を忘れる射を
非空開上で有限 \'etale とし、この慣性の一致から
系\ref{mbrr:cor:stack-isomorphism}で二つの stack の同型を得る。
これは幾何点の一致だけを根拠にした相対認識ではない。

系\ref{mbca:cor:marked-roots}により、$\mathsf r\OO\simeq\OO(1)$ を満たす
次数一根の分類も完了している。
構造対象の像を指定しない全ての第三根の分類、残余比較核の全 $\OO$-降下、
この局所型が全非可換三次の変形モジュライの既約成分をなすことは別問題である。
有限表示の認識と一般開の商構成を、任意の抽象 Abel 圏変形の認識へ拡張していない。
\end{remark}

% !TEX root = ../cubic-moduli-descent.tex
\section{表示スタックと全構造対象スタックの一般相対認識}
\label{mbrr:sec:relative-recognition}

幾何点の対象剛性と自己同型群の一致だけでは、非被約底上の同型層の一致は従わない。
本節では、定理\ref{mbg:thm:generic-stack}の関係 Grassmann 模型を用いて、
表示を忘れる射がさらに非空開上で有限 \'etale となることを先に示す。
その後で幾何点の慣性の一致を用いる。
この順序により、全ての非被約底へ余法線の認識を直接拡張することなく、
一般開上の相対認識を得る。

以下では定理\ref{mbg:thm:generic-stack}の
$\mathcal X^\circ$, $Z^\circ$, $H=\PGL_m$ を固定し、
\[
 \mathfrak P^\circ=[\mathcal X^\circ/G_k],
 \qquad
 \mathfrak M^\circ=[Z^\circ/H]
\]
と書く。ここで $\mathfrak P^\circ$ は
定理\ref{mbs:thm:presentation-stack}の表示スタックの対応する開部分である。
第三 Veronese と tails を取る操作は自然な射
\begin{equation}\label{mbrr:eq:forgetful}
 \nu:\mathfrak P^\circ\longrightarrow\mathfrak M^\circ
\end{equation}
を定める。両者の対象は射影的な表示の降下を含み、
大域的な実次数付き代数を持つとは限らない。

\subsection{射影的な枠による相対ファイバー}

\begin{lemma}[表示を忘れる射の忠実性]
\label{mbrr:lem:faithful}
射\eqref{mbrr:eq:forgetful}は任意の試験基底上で射に忠実であり、
代数空間によって表現可能である。
\end{lemma}
\begin{proof}
次数付き自己同型 $\alpha$ の構造移送が、構造対象の同定を保って
恒等関手と自然同型になる場合を調べればよい。
fpqc 局所で実際の次数付き代数 $A$ と $\alpha$ を選ぶ。
各 $\OO(j)$ は $\alpha$ の構造移送で保たれる。
系\paperref{main}{mbsd:cor:saturation}は任意底変更後にも
\[
 \End(\OO(j))=\OO_T,\qquad
 \Hom(\OO(j),\OO(j+1))=A_1
\]
を与える。従って選んだ自然同型の各成分は単元 $b_j$ であり、
隣接 Hom 上の自然性から
$\alpha_1=(b_{j+1}/b_j)\operatorname{id}$ となる。
$A_1$ は正の階数の有限局所自由加群なので、この比は $j$ に依存しない。
次数一生成性により $\alpha$ は共通 grading scalar である。
よって射影的な表示の射としては恒等である。
これは非被約底でも成立し、fpqc 局所で確認した忠実性は降下する。
忠実な algebraic stack の射の表現可能性
\cite[Tag~04Y5]{MBGStacks}を用いる。
\end{proof}

線型枠では共通 grading scalar が $A_3$ に重み三で作用するため、
ここでは射影的な枠
\begin{equation}\label{mbrr:eq:projective-frames}
 \mathcal P
 =\operatorname{Fr}_{\mathcal X^\circ}(A_3)/\mathbf G_m
 \longrightarrow\mathcal X^\circ
\end{equation}
を使う。これは $H$-torsor である。
$G_k$ は $\mathcal P$ に作用し、その作用は $H$ と可換する。
関係空間の射は $G_k$-不変な $H$-同変射
$\bar\phi:\mathcal P\to Z^\circ$ を与える。

\begin{proposition}[忘却射の平坦性と準有限性]
\label{mbrr:prop:forgetful}
$\nu$ は表現可能、有限表示、忠実平坦かつ準有限である。
幾何点 $y$ を表す満支持模型 $A$ を一つ選ぶと、
その幾何ファイバーの点の集合は
\begin{equation}\label{mbrr:eq:coset-fiber}
 |\mathfrak P^\circ_y|
 \simeq G_A/\mathcal K_k,\qquad
 G_A=\Aut_{\OO}(\QGr A)
\end{equation}
と同定される。ここではファイバーの被約性までは主張しない。
\end{proposition}
\begin{proof}
射影的な枠を用いた二つの商を取ると
\begin{equation}\label{mbrr:eq:atlas-basechange}
 \mathfrak P^\circ\times_{\mathfrak M^\circ} Z^\circ
 \simeq Q:=[\mathcal P/G_k]
\end{equation}
となる。具体的に、右辺の $H$-商は
$[\mathcal X^\circ/G_k]$ であり、その $Z^\circ$ への射は
$\bar\phi$ から誘導される。
補題\ref{mbrr:lem:faithful}により $Q$ は代数空間である。
その atlas $\mathcal P\to Q$ は滑らかな全射である。

定理\ref{mbg:thm:generic-stack}で構成した
$\operatorname{Fr}(A_3)\to Z^\circ$ は忠実平坦・有限表示である。
$\operatorname{Fr}(A_3)\to\mathcal P$ は $\mathbf G_m$-torsor なので、
平坦性の降下により $\bar\phi$ も忠実平坦・有限表示である。
さらに $\mathcal P\to Q$ による降下と
式\eqref{mbrr:eq:atlas-basechange}から $\nu$ の忠実平坦性を得る。
有限表示性はこの有限型の商表示から従う。

代数閉体上では系\ref{mbiq:cor:commuting-o-rigidity}により、
$y$ の上の表示は一つの次数付き同型類である。
その構造対象付き比較の選択は $G_A$ の元であり、
表示の自己同型による変更が $\mathcal K_k$ の作用となる。
補題\ref{mbrr:lem:faithful}により相対自己同型はないので、
ファイバーの点は式\eqref{mbrr:eq:coset-fiber}の剰余類である。
$G_A$ はこの開上有限であるから各幾何ファイバーの点数は有限である。
有限型代数空間のファイバーに対する準有限性の判定により結論を得る。
\end{proof}

\subsection{一般有限 \'etale 比較と stack の同型}

\begin{theorem}[一般開上の有限 \'etale な表示射]
\label{mbrr:thm:finite-etale}
$Z^\circ$ をさらに非空な $H$-不変開へ縮め、
両スタックをその逆像へ制限すると、$\nu$ は有限 \'etale 全射になる。
その階数は一定の正整数 $d$ であり、全幾何点で
\begin{equation}\label{mbrr:eq:degree}
 d=[G_A:\mathcal K_k]
   =\frac{|G_A|}{3^{3k-1}}
\end{equation}
となる。
\end{theorem}
\begin{proof}
$\mathcal P$ は滑らかな $\mathcal X^\circ$ 上の $H$-torsor なので、
$\mathbb C$ 上滑らかである。
$Z^\circ$ の一般点を $\eta$ とすると、
$\mathcal P_\eta$ の局所環は正則環の局所化である。
標数零の体上では正則性から幾何学的正則性、従って滑らかさが従う
\cite[Tags~038U, 038X]{MBGStacks}。
したがって $\mathcal P_\eta$ は $\kappa(\eta)$ 上滑らかである。
滑らかな全射 $\mathcal P_\eta\to Q_\eta$ により、
$Q_\eta$ もこの体上滑らかとなる。
前命題により $Q_\eta$ は準有限かつ準コンパクトなので、有限 \'etale である。

一般点は、一つの affine 開集合内の主開集合の逆極限である。
$Q$ と $Z^\circ$ は準コンパクト・準分離な有限型代数空間である。
前者の準分離性は、有限型の $G_k$ による有限型スキーム $\mathcal P$ の
商表示から従う。従って有限性と \'etale 性の有限表示降下
\cite[Tags~084Z, 07SL]{MBGStacks}を適用できる。
これにより非空開 $Z_1\subset Z^\circ$ 上で $Q\to Z^\circ$ は有限 \'etale となる。
この議論では $\mathfrak P^\circ$ の分離性を仮定していない。

$Q\to Z^\circ$ は $H$-同変なので、その性質は各 $H$-平行移動 $hZ_1$ 上でも成立する。
全ての平行移動の和へ置き換えても、有限 \'etale 性の標的上の局所性により保たれる。
こうして $H$-不変な非空開を得る。
式\eqref{mbrr:eq:atlas-basechange}から $\nu$ も有限 \'etale となり、
全射性は前命題から従う。
既約な標的上では階数は一定であり、
今は幾何ファイバーが被約なので
式\eqref{mbrr:eq:coset-fiber}がその階数を与える。
\end{proof}

\begin{corollary}[満支持一般開の全相対認識]
\label{mbrr:cor:stack-isomorphism}
本節の満支持局所型では、定理\ref{mbrr:thm:finite-etale}で
さらに選んだ非空開上で
\begin{equation}\label{mbrr:eq:stack-isomorphism}
 \mathfrak P^\circ\xrightarrow{\sim}\mathfrak M^\circ
\end{equation}
である。特に、この開に属する二つの表示族について、任意の非被約底を含む
試験スキーム上で自然な fpqc 層の同型
\begin{equation}\label{mbrr:eq:relative-isom}
 \underline{\operatorname{Isom}}_{\rm gr}(A,A')/\mathbf G_m
 \xrightarrow{\sim}
 \operatorname{Eq}_{\OO}(\QGr A,\QGr A')
\end{equation}
を得る。
\end{corollary}
\begin{proof}
定理\ref{mbca:thm:full-inertia}により、チャートの全幾何点で
$G_A=\mathcal K_k$ である。
従って式\eqref{mbrr:eq:degree}の一定階数は一である。
階数一の有限 \'etale 射は同型である。
同じ結論は、表現可能な \'etale 射の幾何ファイバーが一つの被約点となることと、
普遍単射な \'etale 射が開埋込みであること
\cite[Tags~05W5, 0DUD]{MBGStacks}からも従う。
スタックの同型は全ての基底変更と両立するので、
表示スタックの同型層を取れば式\eqref{mbrr:eq:relative-isom}となる。
ここで与えられた同値の次数付き持上げは fpqc 局所での主張であり、
大域的な実代数の frame や持上げを追加で指定したわけではない。
\end{proof}

\begin{remark}[必要な相対議論]
一般に、前系の議論には幾何一般点での慣性の一致だけでも十分である。
有限 \'etale 射の階数が一般点で一なら、既約な標的上で恒等的に一となるためである。
この系は、幾何点の全単射から直ちに導いたものではない。
例えば $\operatorname{Spec}\mathbb C\to
\operatorname{Spec}\mathbb C[\epsilon]/(\epsilon^2)$ は、
全ての代数閉体で同じ点と自明な慣性を持つが同型ではない。
ここでは射影枠による表現可能性、有限表示、および一般ファイバーからの
有限 \'etale 性の降下がこの不足を埋める。
縮小前の全ての表示パラメータ上で式\eqref{mbrr:eq:relative-isom}が成立することや、
抽象 Abel 圏の全変形がこの局所型を持つことは結論しない。
\end{remark}

\subsection{全慣性の中心の内在的な認識}

\begin{corollary}[中心三乗根群の相対的な回復]
\label{mbrr:cor:central-band}
式\eqref{mbrr:eq:stack-isomorphism}の非空開上で、
$\mathfrak M^\circ$ の慣性の中心は、階数 $3^{k-1}$ の有限 \'etale な中央部分群である。
ブロックを順序付ける有限 \'etale 被覆上では、その群は自然に
\begin{equation}\label{mbrr:eq:central-band}
 Z_0=(\mu_3^k)/\mu_3^{\rm diag}
\end{equation}
となる。一つのブロックを基準に取れば $Z_0\simeq(\mu_3)^{k-1}$ であり、
これはそのブロックを固定して他の因子を独立に三乗根倍する標準作用である。
順序を忘れた底上の descent は、$\mathfrak S_k$ が
式\eqref{mbrr:eq:central-band}の因子を置換する作用による。
\end{corollary}
\begin{proof}
まず幾何点の
$\mathcal K_k=(\prod_i\mathcal H_i)/\mu_3^{\rm diag}$ を考える。
その元が中心なら、第 $i$ 因子だけに属する全ての元との交換子が共通対角群に入る。
$k\ge2$ なので、その交換子の他の成分が一であることから共通対角値も一である。
Heisenberg pairing の非退化性により、各因子の元は中心 $\mu_3$ に属する。
従って $Z(\mathcal K_k)=Z_0$ である。

順序付き表示チャートでは、$Z_0$ はブロックスカラーとして
$G_k$ のブロック保存部分に入り、$\mathcal X^\circ$ に自明に作用する。
次数付き同型の相対認識
（定理\ref{mbs:thm:presentation-stack}）により、
全ての試験基底でブロック保存自己同型と可換する。
従って慣性への有限 \'etale な中央部分群を与える。
慣性は不分岐であるから、この部分群の単射は開埋込みである。
幾何点で中心を尽くすので、任意基底上の中心の切断もこの開部分へ因子化する。
これで中心そのものを表すことが分かる。
ブロックの置換はこの群の因子を置換するため、
順序付き被覆からの降下は記載した $\mathfrak S_k$ 作用となる。
最後に式\eqref{mbrr:eq:stack-isomorphism}を通して全 $\OO$-stack へ移す。
\end{proof}

ここで回復した中心部分群は、定理\paperref{main}{mbza:thm:central-action}で
coherent な圏作用へ移される。compact 残余圏と projector の降下は
定理\ref{mbsd2:thm:projector-descent}による。
共通 scalar の自然自明化は次数一 shift に重み一を残すため、
rotation の線型化、関係直線を含む周期、比較核の実現はそれぞれ区別する。
第\paperref{main}{mbrd:sec:residual-descent}節は同変積表示と核 stack の降下を与えるが、
指定した一般点の核の大域切断を、この中心群の認識だけから結論しない。

% !TEX root = ../cubic-moduli-descent.tex
\section{共通次数スカラーと残余構造の降下}
\label{mbsd2:sec:scalar-descent}

次数付き代数の共通スカラーは tails 上で恒等関手と自然同型であるが、
その自然同型は次数 shift と厳密には可換しない。
本節ではこの差を計算し、表示 gerbe、例外的成分、rotation の降下を区別する。
相対 Serre 双対や残余 projector の構成には compact 部分を用いる。
三次周期を残余圏へ移す箇所では、別途 shift と両立する MF/window 比較を仮定する。
全表示上の tails の滑らかさや、compact と全 coherent 対象の一致は仮定しない。

\subsection{表示 gerbe とスカラーの自然同型}

第\ref{mbs:thm:presentation-stack}定理の満支持楕円開を用い、
ブロック数を $k\geq2$ とする。記法を
\[
 H_k=\GL_3^k\rtimes\mathfrak S_k,
 \qquad G_k=H_k/\mathbf G_m^{\mathrm{diag}},
 \qquad
 q:\widetilde{\mathfrak P}_k=[\mathcal X_k/H_k]
       \longrightarrow\mathfrak P_k=[\mathcal X_k/G_k]
\]
と固定する。これは共通スカラーを band とする $\mathbf G_m$-gerbe であり、
その類を $\beta\in H^2_{\mathrm{fppf}}(\mathfrak P_k,\mathbf G_m)$ と書く。
順序付き被覆では、それぞれ $H_k$ と $G_k$ から $\mathfrak S_k$ を除く。
実際の代数 $A=B/(\ell)$、$B=\bigotimes_iS_i$ と関係直線 $\ell\subset B_3$ は
$\widetilde{\mathfrak P}_k$ 上に存在する。
共通スカラー $\lambda$ は $A_d$ に $\lambda^d$、$\ell$ に $\lambda^3$ として作用する。

右加群と内部 shift の規約を
\begin{equation}\label{mbsd2:eq:convention}
 \Phi_h(M)=M^{h^{-1}},\qquad m*a=m\,h^{-1}(a),
 \qquad M(j)_n=M_{n+j}
\end{equation}
とする。このとき $h\mapsto\Phi_h$ は群作用である。
以下で重みというとき、係数線束については通常の gerbe の重み、
対象・関手については次の自然同型で共通スカラーを消した後の重みをいう。

\begin{proposition}[次数 shift のスカラー重み]
\label{mbsd2:prop:scalar-weight}
共通 grading scalar $s_\lambda$ に対して
\begin{equation}\label{mbsd2:eq:tau}
 \tau_{\lambda,M}:\Phi_{s_\lambda}(M)\xrightarrow{\sim}M,
 \qquad m\in M_n\longmapsto\lambda^n m
\end{equation}
は graded 加群、tails、導来 enhancement 上の自然同型である。
これは積 $\lambda\mu$ と全 $H_k$ 作用に整合し、$\OO$ 上の同定を保つ。
従って $\OO$ 付き Abel 圏とその導来 compact 部分は
$\mathfrak P_k$ 上の線型圏の stack へ降下する。

全 $j\in\mathbf Z$ に対して
\begin{equation}\label{mbsd2:eq:shift-commutator}
 \tau_{\lambda,M(j)}
   =\lambda^{-j}\bigl(\tau_{\lambda,M}\bigr)(j).
\end{equation}
この規約で、標準の同変構造を持つ $\OO(j)$ と関手 $(j)$ はともに重み $j$ である。
重み $a$ の係数線束 $L$ を tensor すると重みは $j+a$ になる。
特に
\begin{equation}\label{mbsd2:eq:corrected-shift}
 (3m)\otimes\ell^{-m},\qquad
 \OO(3m)\otimes\ell^{-m}\quad(m\in\mathbf Z)
\end{equation}
は $\mathfrak P_k$ へ降下する。
\end{proposition}

\begin{proof}
$a\in A_d$ に対して
$\tau_\lambda(m*a)=\lambda^{n+d}\lambda^{-d}ma
=\tau_\lambda(m)a$ なので \eqref{mbsd2:eq:tau} は加群同型である。
全写像が次数を保つことから自然性が従い、
$\lambda^n\mu^n=(\lambda\mu)^n$ が積の coherence を与える。
次数付き同型 $h$ とも可換する。これらの恒等式は複体の微分と両立し、
torsion 加群を保つので、局所化と enhancement に移る。
局所的な $H_k$ への持上げの差は共通スカラーである。
その差を $\tau$ で同定すると、三重重なりの条件は上の積の恒等式となる。
加群の fpqc 降下を適用し、得た線型圏の導来圏と compact 部分を取ればよい。
$\Phi_h(A)\to A$ の同定は $a\mapsto h(a)$ であり、
共通スカラーでは $\tau_{\lambda,A}$ そのものなので構造対象も降下する。

$M(j)_n=M_{n+j}$ 上で二つの辺は、それぞれ $\lambda^n$ と $\lambda^{n+j}$ を掛ける。
従って \eqref{mbsd2:eq:shift-commutator} を得る。
符号を明確にするため、同変関手 $F$ の比較を
$\gamma_\lambda:\Phi_{s_\lambda}F\to F\Phi_{s_\lambda}$ と書くと、
誘導されるスカラー作用は
\begin{equation}\label{mbsd2:eq:functor-weight}
 F\xrightarrow{\tau_{\lambda,F}^{-1}}\Phi_{s_\lambda}F
 \xrightarrow{\gamma_\lambda}F\Phi_{s_\lambda}
 \xrightarrow{F(\tau_\lambda)}F
\end{equation}
である。$F=(j)$ の標準比較は恒等なので、この合成は $\lambda^j$ となる。
また $\Phi_h(A(j))\to A(j)$ の標準同定は古い次数で $h$ を作用させる。
次数 $n$ の成分では $\lambda^{n+j}$ であり、$\tau$ との差は同じ $\lambda^j$ である。
tensor 積の重みは加算される。
重み零の加群・複体が gerbe の底からの引戻しであることは
\cite[Propositions~4.10, 5.7]{NFMBerghSchnurer}、
tensor と Hom の重みは同文献 Corollaries~4.12, 5.9 による。
ここでは対応する加群・双加群の作用写像も重み零なので、関手と合成も降下する。
\end{proof}

\subsection{表示 gerbe の正確な周期}

\begin{theorem}[全非空開で残る三次障害]
\label{mbsd2:thm:gerbe-period}
$\mathfrak P_k$ の任意の非空開 $\mathfrak U$ に対し、
$\beta|_{\mathfrak U}$ の位数は正確に三である。
対応する表示 gerbe 上の線束が純重み $j$ を持つための必要十分条件は $3\mid j$ である。
従って正の shift の冪を係数線束の tensor によって重み零にできる最小の冪は三である。

順序付き被覆上では、純重み一のベクトル束の最小の正階数は三である。
順序を忘れた底でも純重み一のベクトル束の階数は三の倍数であり、
$V=\bigoplus_iW_i$ は階数 $3k$ の例を与える。
\end{theorem}

\begin{proof}
重み三の線束 $\ell$ が存在するので $3\beta=0$ である。
これは局所表示の持上げの三重積を $c_{abc}\in\mathbf G_m$ と書けば、
$\ell$ の遷移写像により $c_{abc}^3$ が coboundary になることを意味する。
同様に重み $j$ の線束の存在は $j\beta=0$ と同値である。

非空開の任意の幾何点を取る。その射影的慣性は
$K_k=(\prod_i\mathcal H_i)/\Delta\mu_3$ である。
実表示の慣性は、定理\ref{tba:thm:generic-inertia}の証明により
\begin{equation}\label{mbsd2:eq:central-extension}
 \widetilde K_k=
 \bigl(\mathbf G_m\times\prod_i\mathcal H_i\bigr)
 \big/\{(\zeta,\zeta^{-1},\ldots,\zeta^{-1}):\zeta\in\mu_3\}.
\end{equation}
各 Heisenberg 群の commutator はその中心 $\mu_3$ 全体である。
従って $\prod_i\mathcal H_i$ の対角 $\mu_3$ も commutator に含まれ、
\eqref{mbsd2:eq:central-extension}ではこれが共通 grading scalar の $\mu_3$ となる。
一次元表現は commutator を殺すため、$\widetilde K_k$ の任意の指標が
共通 $\mathbf G_m$ 上で $\lambda^j$ なら $3\mid j$ でなければならない。
gerbe 上の線束をこの幾何点の慣性へ制限すれば同じ条件が得られる。
よって重み一・二の線束はなく、$\beta$ は非零である。
これは任意の非空開で通る。逆に $j=3m$ なら $\ell^m$ が存在する。
ここで用いたのは残余 gerbe $BK_k$ への制限であり、
代数閉体の点そのものへの引戻しが非中立という主張ではない。

重み一・階数 $d$ のベクトル束の行列式は重み $d$ の線束なので $3\mid d$ である。
順序付き被覆には重み一・階数三の $W_i$ があり、下界を達成する。
\end{proof}

\begin{corollary}[関係方程式を標識する $\mu_3$-gerbe]
\label{mbsd2:cor:relation-frame}
$\ell$ の非零 frame を選ぶ stack
\[
 \widetilde{\mathfrak P}_k^{\ell=1}
   :=\operatorname{Fr}(\ell)\longrightarrow\widetilde{\mathfrak P}_k
       \longrightarrow\mathfrak P_k
\]
は $\mu_3$-gerbe である。その band を
$\mu_3\hookrightarrow\mathbf G_m$ で拡張すると元の表示 gerbe を得る。
特にこれも任意の非空開で中立にはならない。
\end{corollary}
\begin{proof}
fpqc 局所的に実表示と frame を取れる。
共通スカラーは frame に三乗で作用するので、frame を保つ相対慣性は $\mu_3$ である。
三乗写像 $\mathbf G_m\to\mathbf G_m$ は fppf 全射なので、任意の二つの持上げは
局所的に同型となる。band 拡張の記述はこの局所模型から従う。
もしこの $\mu_3$-gerbe が中立なら、その band 拡張も中立となり前定理に反する。
\end{proof}

\subsection{例外的成分と残余 projector}

$N=3k$, $r=N-3$, $n=N-2$ と置く。
実表示上の compact 圏を $\mathscr D_A^c=D(\QGr A)^c$ とし、
\[
 \mathscr K_A=
 \langle\OO,\OO(1),\ldots,\OO(r-1)\rangle^\perp
 \subset\mathscr D_A^c
\]
と定める。右直交はこれらの対象からの全 RHom の消滅を意味する。

\begin{theorem}[残余圏と mutation の降下]
\label{mbsd2:thm:projector-descent}
各 $\OO(j)$ の生成する admissible 成分、その評価 projector
\[
 E_j(M)=\mathbf R\Hom(\OO(j),M)\otimes\OO(j),
\]
評価射 $E_j\to\operatorname{id}$、left mutation $L_{\OO(j)}$ は重み零である。
従って $\mathscr K_A$ とその projector は $\mathfrak P_k$ へ降下する。
この主張と構成は、普遍表示からの任意の底変更に両立する。

これは $\OO(j)$ 自体の降下を意味しない。
重み $a$ の層を $a\beta$-twisted と呼ぶ規約では、降下した第 $j$ 成分は
\begin{equation}\label{mbsd2:eq:twisted-block}
 \operatorname{Perf}(\mathfrak P_k,-j\beta)
\end{equation}
から、局所的な $F\mapsto F\otimes\OO(j)$ によって与えられる。
$3\nmid j$ ならこの block の局所的生成対象を大域的な $\OO(j)$ として選ぶことはできない。
\end{theorem}

\begin{proof}
定理\paperref{main}{mbsd:thm:serre}より $\OO(j)$ は相対的に exceptional であり、
compact 対象間の RHom は係数 perfect である。
従って表示した評価とその cone は compact 部分上で定義され、
左右の随伴を持つ admissible 成分を与える。
$0\leq i<j<r$ なら
$\mathbf R\Hom(\OO(j),\OO(i))=0$ なので、表示した有限列は半直交である。
右直交への projector は、右端の例外成分から順に left mutation を施して得る。
全評価写像とその cone は dg enhancement で取り、合成の coherence を保持する。

$\OO(j)$ の重みは $j$、その第一変数から来る RHom の重みは $-j$ である。
よって $E_j$ と評価射の重みは零となり、その cone と有限合成も同じである。
実際、$\OO(j)$ を係数線束 $L$ だけ変更すると
\[
 \mathbf R\Hom(\OO(j)\otimes L,M)\otimes(\OO(j)\otimes L)
 \simeq\mathbf R\Hom(\OO(j),M)\otimes\OO(j)
\]
が評価射と両立する。これが三重重なりでも成り立つ降下同型である。
RHom、有限 cone、tensor の底変更は定理\paperref{main}{mbsd:thm:serre}による。

係数が重み $-j$ なら $F\otimes\OO(j)$ の全重みは零になる。
逆は $\mathbf R\Hom(\OO(j),-)$ で与えられるため、
\eqref{mbsd2:eq:twisted-block} を得る。
さらに $\End(\OO(j))=\OO$ で負の自己 Ext は零なので、
この局所対象の実現 groupoid は $j\beta$ を類とする $\mathbf G_m$-gerbe である。
定理\ref{mbsd2:thm:gerbe-period}により $3\nmid j$ の場合は中立ではない。
\end{proof}

\subsection{関係直線を保持した rotation と三次周期}

標準 rotation 候補を、実表示上で
\[
 \rho=L_{\OO}\circ(1)|_{\mathscr K_A}
\]
と書く。まず自己同値性と相対 Serre 公式は有限 mutation だけから従う。

\begin{proposition}[compact 残余の rotation と Serre]
\label{mbsd2:prop:rotation-serre}
普遍表示からの任意の底変更後に、$\rho$ は $\mathscr K_A$ の自己同値であり、
その相対 Serre 関手は
\begin{equation}\label{mbsd2:eq:rotation-serre}
 \mathsf S_{\mathscr K_A}
    \simeq\rho^{-r}\otimes\mathscr D[n],
 \qquad\mathscr D=\det V\otimes\ell^\vee
\end{equation}
である。この結論には MF 比較を必要としない。
\end{proposition}

\begin{proof}
$s=(1)$ と置く。$M\in\mathscr K_A$ に対する評価三角
\[
 \mathbf R\Hom(\OO,sM)\otimes\OO\longrightarrow sM\longrightarrow\rho M
\]
を使う。$\mathbf R\Hom(\OO,\rho M)=0$ は mutation の定義から従う。
$1\leq a\leq r-1$ では
$\mathbf R\Hom(\OO(a),\OO)=0$ および
$\mathbf R\Hom(\OO(a),sM)
=\mathbf R\Hom(\OO(a-1),M)=0$ なので、$\rho M\in\mathscr K_A$ である。

包含を $i:\mathscr K_A\hookrightarrow\mathscr D_A^c$、その左随伴を $P$ と書くと、
半直交列から
\[
 P=L_{\OO}L_{\OO(1)}\cdots L_{\OO(r-1)}
\]
となる。右端から適用する合成である。
$sL_{\OO}=L_{\OO(1)}s$ を反復すれば
\[
 \rho^r=P s^r i
\]
を得る。相対的に proper な圏の admissible 部分圏の逆 Serre は
$\mathsf S_{\mathscr K_A}^{-1}=P\mathsf S_A^{-1}i$ である。
実際、左随伴と Serre pairing を適用すると、両対象 $M,N\in\mathscr K_A$ について
\[
 \mathbf R\Hom(P\mathsf S_A^{-1}iM,N)
 \simeq\mathbf R\Hom(\mathsf S_A^{-1}iM,iN)
 \simeq\mathbf R\Hom(N,M)^\vee
\]
となる。右随伴による Serre とこの逆は互いに逆である。
定理\paperref{main}{mbsd:thm:serre}を代入すると
\[
 \mathsf S_{\mathscr K_A}^{-1}
   =P s^r i\otimes\mathscr D^\vee[-n]
   =\rho^r\otimes\mathscr D^\vee[-n].
\]
従って $\rho^r$ は自己同値であり、$\rho$ 自身も自己同値となる。
逆を取れば \eqref{mbsd2:eq:rotation-serre}を得る。
全比較は有限 mutation と既存の相対 Serre pairing から作られ、任意底変更に両立する。
\end{proof}

以下の周期に関する主張では、扱う残余 dg 圏と graded MF 圏の間に、
内部 shift をこの $\rho$ へ移す比較があることを入力とする。
例えば第\paperref{main}{tw:thm:window}定理およびその底変更が与える圏では、この条件を満たす。
同定が構成された compact 残余と、全 coherent 残余への拡張の範囲は区別する。

\begin{proposition}[係数直線を含む三次周期]
\label{mbsd2:prop:line-period}
残余 projector と両立する標準 rotation の重みは一である。
MF 圏では関係直線を保持した正準的な周期は
\begin{equation}\label{mbsd2:eq:mf-line-period}
 [2]\simeq(3)\otimes\ell^\vee
\end{equation}
となる。上記の window 比較がある場合、その指定周期は
\begin{equation}\label{mbsd2:eq:rotation-line-period}
 \rho^3\otimes\ell^\vee\xrightarrow{\sim}[2],
 \qquad\text{すなわち}\qquad
 \rho^3\xrightarrow{\sim}[2]\otimes\ell
\end{equation}
である。左式は重み零の関手間の同型として $\mathfrak P_k$ へ降下する。
$\ell$ の frame を選ぶ $\mu_3$-gerbe 上では、通常の指定式 $\rho^3\simeq[2]$ となる。

一般に $\rho^j$ の標準同変構造は重み $j$ を持つ。
係数線束による補正で重み零にできる正の最小冪は三であり、
その補正は $\rho^3\otimes\ell^\vee$ で達成される。
\end{proposition}

\begin{proof}
mutation の重みは零、shift の重みは一なので、最初の主張が従う。
関係の包含 $\ell\to B_3$ は、中心自然変換
\[
 w:\operatorname{id}\longrightarrow T,
 \qquad T=(3)\otimes\ell^\vee
\]
を与える。対応する line-valued MF は
$P_1\xrightarrow{d_1}P_0\xrightarrow{d_0}T(P_1)$ と
$d_0d_1=w$, $T(d_1)d_0=w$ からなる。
翻訳は $(P_0,P_1)\mapsto(T(P_1),P_0)$ と両微分の符号反転であり、
二回繰り返すと $T$ そのものになる。
従って \eqref{mbsd2:eq:mf-line-period} は定義から従う。
この line-valued 規約は
\cite[Definitions~1.1, 1.3]{MBS2PolishchukVaintrob}と同じであり、
中心元による非可換 MF でも同じ二つの合成式だけを使う。

window の shift 比較を三回合成すると、\eqref{mbsd2:eq:rotation-line-period}を得る。
これは局所 frame での中心元の二拡張を用いる
式\papereqref{main}{tw:eq:specified-period}、\papereqref{main}{tw:eq:rotation-period-square}を
関係直線を保持して書いたものである。
frame の変更は $w$ とその二拡張を同時に変えるので、係数 $\ell^\vee$ がその差を消す。
全構成は関係の包含と評価から作られ、同変性と指定周期の coherence を保つ。
最小性は定理\ref{mbsd2:thm:gerbe-period}による。
\end{proof}

\begin{corollary}[scalar-simple な rotation 核の実現障害]
\label{mbsd2:cor:rotation-obstruction}
ある開部分で上の rotation が可逆な dg 核によって実現され、
その核の実現 stack が fpqc 降下を満たすとする。
さらに局所的な核について
\[
 \mathcal H^a\mathbf R\operatorname{Hom}(\rho,\rho)=0\ (a<0),
 \qquad
 \mathcal H^0\mathbf R\operatorname{Hom}(\rho,\rho)=\OO
\]
が任意底変更と両立すると仮定する。
このとき局所 rotation の同型類を実現する $\mathbf G_m$-gerbe の類は $\beta$ である。
従って任意の非空開で、この同型類の実際の大域的 rotation 核は存在しない。
これは残余圏とその projector の降下とは両立する。
\end{corollary}

\begin{proof}
局所 rotation は表示同型によって移されるので、その同型類の実現 stack は
局所的に非空である。仮定した Ext 条件により、その同型類の部分は
通常の $\mathbf G_m$-gerbe であり、高次の負の自己 Ext による追加の慣性はない。
実表示をその標準 rotation 核へ送ると表示 gerbe からこの gerbe への射を得る。
命題\ref{mbsd2:prop:scalar-weight}と projector の重み零性により、
この射の band 上の写像は $\lambda\mapsto\lambda$ である。
同じ band を持つ gerbe 間の band を保つ射は同値なので、その類は $\beta$ となる。
前定理の非中立性を適用する。
\end{proof}

\begin{remark}[条件付き非降下と標準同変構造の区別]
直前の系の scalar-simple 条件は、残余圏の $HH^{<0}=0$, $HH^0=\OO$ から従う。
滑らかな可換三次との比較が与えられた場合には
命題\paperref{main}{prop:nf-hh}が幾何ファイバーでの入力となるが、
相対的な適用には上に明記した底変更も確認する。
この条件を確かめていない全表示開で、本節が無条件に示すのは、
標準 rotation の重み一、係数線束による補正の最小冪三、および
projector の重み零である。特定の同変構造の失敗だけから、
全ての可能な核の実現を否定してはいない。
同様に $\rho^3$ の標準同変構造自体の重みは三であり、
降下するのは \eqref{mbsd2:eq:rotation-line-period}の補正付き関手である。
\end{remark}

\begin{corollary}[残余 Serre 式の重み零部分]
\label{mbsd2:cor:residual-serre}
上の window 比較により三次周期が得られた範囲では
\begin{equation}\label{mbsd2:eq:residual-serre}
 \mathsf S_{\mathscr K}\simeq[k]\otimes
       \bigl(\det V\otimes\ell^{-k}\bigr).
\end{equation}
係数直線 $\det V\otimes\ell^{-k}$ は重み零であり、$\mathfrak P_k$ へ降下する。
\end{corollary}
\begin{proof}
$r=3(k-1)$ なので、\eqref{mbsd2:eq:rotation-line-period}から
$\rho^{-r}=[-2(k-1)]\otimes\ell^{-(k-1)}$ を得る。
$n=3k-2$ と $\mathscr D=\det V\otimes\ell^{-1}$ を
命題\ref{mbsd2:prop:rotation-serre}に代入すれば
\eqref{mbsd2:eq:residual-serre}となる。
$\det V$ の重みは $3k$、$\ell^k$ の重みも $3k$ である。
\end{proof}

\begin{remark}[二つの表示のスカラーを同期する条件]
\label{mbsd2:rem:kernel-synchronization}
二つの表示 gerbe を添字 $s,t$ で区別し、その積の band を
$\mathbf G_{m,s}\times\mathbf G_{m,t}$ とする。
比較核 $P$ 自身の重みを固定すると、$P\rho_s$ と $\rho_tP$ の重みには
それぞれ $(1,0)$ と $(0,1)$ が加わる。
従って rotation 比較
$\eta:P\rho_s\to\rho_tP$ の Hom は相対重み $(-1,1)$、
すなわち $\lambda_t/\lambda_s$ で変換される。
この差を消す共通スカラーは対角部分群である。
比の band へ押し出した gerbe の類は $\beta_t-\beta_s$ となる。
三重比較の周期図式では、関係直線の比較
$\ell_s\to\ell_t$ に重み $(-3,3)$ が現れる。
これらは比較\emph{データ}の重みの計算であり、比の gerbe の中立化、
直線の同定、同変な $\eta$、あるいは実際の Morita 核の存在を与えるものではない。
\end{remark}

\begin{remark}[全構造対象モジュライへの適用範囲]
定理\ref{mbrr:thm:finite-etale}と系\ref{mbrr:cor:stack-isomorphism}の非空開では、
表示スタックと全 $\OO$-局所型スタックは同型である。
従って本節の圏、projector、twisted block、および重み零の構造はその全 $\OO$-開へ移る。
これは比較核の大域的 section、全表示上の残余 smoothness、
あるいは別々の表示 gerbe 上の rotation を無補正で同期できることを主張しない。
共通スカラーのこの障害は、因子を置換する際の有限中心群の splitting の障害とも区別する。
\end{remark}

% !TEX root = ../nc-cubic.tex
\section{因子モジュライへの射と点曲線積の相対模型}
\label{mbfm:sec:factor-moduli}

本節では、中心商の係数を忘れる射と、各Sklyanin因子の平行移動を忘れる射を分ける。
点曲線の積は中心三次の生成元を選ばずに全$\OO$開上へ降下する。
一方、その積との残余圏比較には別の核が必要である。
一般のrigidificationについては\cite[Theorem~5.1]{MBFMRomagny}を参照するが、
ここで必要な商とbandは表示群の作用から直接求める。

\subsection{中心係数を忘れる有限gerbe}

第\ref{mbs:sec:presentations}節の$\mathcal H,\mathcal T,\mathcal U_k,
\mathcal X_k$を使い、$\mathcal H_k^{\ne}\subset\mathcal H^k$を
$j$が対ごとに異なる開部分とする。次の三つのstackを置く。
\begin{equation}\label{mbfm:eq:three-stacks}
\begin{aligned}
 \mathfrak P_k&=[\mathcal X_k/
 ((\GL_3^k/\mathbb G_m^{\rm diag})\rtimes\mathfrak S_k)],\\
 \mathfrak B_k&=[\mathcal U_k/(\PGL_3^k\rtimes\mathfrak S_k)],\\
 \mathfrak E_k&=[\mathcal H_k^{\ne}/(\PGL_3^k\rtimes\mathfrak S_k)].
\end{aligned}
\end{equation}
$\mathfrak B_k$は射影的なSklyanin因子の組、$\mathfrak E_k$は
順序を忘れた平面三次の組を分類する。
後者の偏極はBrauer--Severi平面上の次数三偏極類として扱い、
全ての底で実際の$\mathcal O(1)$が存在するとはしない。

\begin{proposition}[中心慣性によるrigidification]
\label{mbfm:prop:rigidification}
中心関係直線を忘れる射
\[
 q_k:\mathfrak P_k\longrightarrow\mathfrak B_k
\]
は、順序付き被覆上で$Z_0=\mu_3^k/\Delta\mu_3$をbandとするgerbeである。
順序を忘れた底のbandは因子置換で降下する有限\'etale群である。
さらに$q_k$は$Z_0$を相対慣性から除くrigidificationであり、
\begin{equation}\label{mbfm:eq:rigidification}
 \mathfrak P_k\mathbin{/\mkern-6mu/}Z_0\simeq\mathfrak B_k.
\end{equation}
幾何点の慣性群の射は
\[
 K_k=(\mathcal H_{27}^k)/\Delta\mu_3
 \longrightarrow\prod_i E_i[3]
\]
で、核は$Z_0$である。
\end{proposition}
\begin{proof}
まず順序を固定する。群の完全列
\[
 1\longrightarrow T_k=\mathbb G_m^k/\mathbb G_m^{\rm diag}
 \longrightarrow\GL_3^k/\mathbb G_m^{\rm diag}
 \longrightarrow\PGL_3^k\longrightarrow1
\]
を用いる。中心直線束を局所的に自明化すると、
$\mathcal X_k\to\mathcal U_k$のファイバーは全係数が非零な
射影$(k-1)$次元トーラスとなる。$T_k$の作用は
\[
 [a_1:\cdots:a_k]\longmapsto
 [\lambda_1^3a_1:\cdots:\lambda_k^3a_k].
\]
従ってこの作用はfppf局所に推移的で、全ての相対安定化子は
$T_k[3]=\mu_3^k/\Delta\mu_3$である。
三乗写像が標数零で有限\'etale全射であることも使える。

$\mathfrak B_k$上の対象をfppf局所で座標付き表示へ戻せば、
中心直線束のframeと全係数非零な関係直線を選べる。
二つのそのような選択は、係数比の三乗根を取った後に同型となる。
同型の相対核は上記$Z_0$である。これでgerbeの二条件を直接確認した
（商groupoidによる判定は\cite[Tag~06PD]{MBFMStacksGerbes}）。
同型層からこの核を割り、stack化すると、因子の全射影同型の層になる。
対象も局所的に全て持ち上がるので、そのstackは$\mathfrak B_k$と同型である。
これが式\eqref{mbfm:eq:rigidification}の普遍性も与える。

各因子の射影自己同型は、点曲線の次数三偏極と非二捩れの平行移動を保つ。
$j\ne0,1728$では原点を保つ群は$\{\pm1\}$であり、
$2p_i\ne0$が負符号を除く。そのため残るのは$E_i[3]$である。
Heisenberg群の中心拡大を因子ごとに射影化すると表示した慣性の射を得る。
因子置換は上の群完全列と三乗写像を保つので、全構成が
$\mathfrak S_k$に同変であり、順序を忘れた結論も従う。
\end{proof}

\begin{remark}[二種類のgerbe]
\label{mbfm:rem:two-gerbes}
$q_k$のbandは有限群$Z_0$であり、
定理\ref{mbsd2:thm:gerbe-period}の実表示
$\widetilde{\mathfrak P}_k\to\mathfrak P_k$のband$\mathbb G_m$とは異なる。
中心係数を忘れたことから、実表示の障害$\beta$の消滅は従わない。
また$q_k$の存在は、その有限gerbeの中立性を主張するものではない。
\end{remark}

\begin{corollary}[中心同変化の因子モジュライへの降下]
\label{mbfm:cor:equivariant-rigidification}
$\mathfrak P_k^\circ$上のcompact残余圏の族を
$\mathscr K$とし、$Z_0$はその内在的慣性として作用するものとする。
$\mathfrak B_k$の対応する開部分$\mathfrak B_k^\circ$上に
圏の族$\mathscr D$があり、
\[
 q_k^*\mathscr D\simeq\mathscr K^{hZ_0}
\]
となる。右辺はcoherenceを含む同変化である。
これは$\mathscr K$自身の降下や、$\mathscr D$と非捩れの曲線積の圏との
同値を主張しない。
本稿の具体的な表示族について、この追加の非捩れ同定は
命題\ref{mbic:prop:rigidified-product}で実積核の慣性を計算して得る。
\end{corollary}
\begin{proof}
gerbeの上の圏の切断を取り、$\mathscr D=(q_k)_*\mathscr K$と定める。
fpqc局所でgerbeの対象$x$を選ぶと、その表示は$BZ_0$となり、
切断の圏は$\mathscr K_x$のhomotopy fixed point圏に一致する。
二つの中立化を比較する同型の持上げは$Z_0$だけ異なるが、
同変対象の構造写像$\theta_z$がその差を自然に同定する。
三重の重なりでの等式は、$\mathscr K$の内在的作用のcoherenceと
$\theta$の群作用条件である。順序の変更は$Z_0$の群の輸送も含めて扱う。
これで局所の同変圏は$\mathfrak B_k^\circ$上へ貼り合わさる。
有限群の位数は可逆であり、有限軌道和と平均射影によってcompact対象の
同変化もcompactな模型内で行える。
一般にはこの降下結果にgerbe由来の捩れが残り得る。
同値核の選択やその捩れの消滅は、この降下の証明には含まれない。
\end{proof}

\subsection{平行移動を忘れる滑らかな射}

\begin{proposition}[点曲線のモジュライへの射]
\label{mbfm:prop:point-map}
自然な射
\begin{equation}\label{mbfm:eq:point-map}
 \operatorname{pt}_k:\mathfrak P_k
 \xrightarrow{q_k}\mathfrak B_k
 \xrightarrow{v_k}\mathfrak E_k
\end{equation}
がある。$v_k$は表現可能な滑らかな全射で相対次元$k$を持つ。
$\dim\mathfrak E_k=k$、$\dim\mathfrak B_k=2k$であり、
順序付き平面三次の組を固定した$v_k$の幾何ファイバーは
\[
 \prod_i\bigl(\operatorname{Jac}(E_i)
       \setminus(E_i[2]\cup E_i[3])\bigr)
\]
である。ここで捩れ部分群はJacobianのものを意味する。
$\operatorname{pt}_k$の対応するファイバーは、この積上の$Z_0$-gerbeである。
全$\OO$開$\mathfrak M_k^\circ\simeq\mathfrak P_k^\circ$へ制限しても
自然な射を得て、像は$\mathfrak E_k$の非空開となる。
\end{proposition}
\begin{proof}
$\mathcal T\to\mathcal H$は普遍Jacobianから有限捩れ閉部分を除いた射であり、
滑らか・全射・相対次元一である。従って
$\mathcal U_k\to\mathcal H_k^{\ne}$は滑らか・全射・相対次元$k$である。
両者を同じ群$\PGL_3^k\rtimes\mathfrak S_k$で割ると$v_k$を得る。
その基底変更は、点曲線を固定したJacobianの開部分の積であるため表現可能で、
滑らかさとファイバーの記述が従う。
商の次元は$9k-8k=k$および$10k-8k=2k$である。
有限gerbe$q_k$の局所模型は$BZ_0$なので、
$\operatorname{pt}_k$もstackの射として滑らかで相対次元$k$である。
その非空開への制限の像は滑らかな射の開性から非空開となる。
全$\OO$射に対する自然性は系\ref{mbrr:cor:stack-isomorphism}による。
\end{proof}

\begin{proposition}[点曲線積の全$\OO$相対模型]
\label{mbfm:prop:product-family}
$\mathfrak E_k$上に、順序付きチャートでは$E_1\times\cdots\times E_k$となる
相対次元$k$の滑らかな射影的族$\mathscr Y_k\to\mathfrak E_k$がある。
従って$\operatorname{pt}_k$による引戻しは$\mathfrak P_k$上、
さらに$\mathfrak M_k^\circ$上の点曲線積の族と
smooth properな$\operatorname{Perf}$の族を与える。
この構成は任意底変更と両立し、幾何的な族への中心慣性$Z_0$の作用は自明である。
\end{proposition}
\begin{proof}
$\mathcal H_k^{\ne}$上の普遍曲線の積には、各因子の射影座標変換と
因子置換が自然に作用する。これを基底と同時に割ると族$\mathscr Y_k$を得る。
滑らかさ・proper性と相対次元は表示チャート上で検査できる。
射影性を確かめるには、各平面上の
\[
 \mathcal O_{\mathbb P(W_i^*)}(3)\otimes(\det W_i)^{-1}
\]
を用いる。この線束は共通scalarの重みが零なので$\PGL(W_i)$に同変であり、
曲線への制限は次数九である。その外積には因子置換の自然な線型化があり、
相対ample線束として降下する。
従って族は射影的であり、$\operatorname{Perf}$は相対smooth properとなる。
普遍族の引戻しとして定義したため、底変更も保たれる。
中心$Z_0$は各$W_i$のscalarによって作用し、射影曲線上では恒等である。
\end{proof}

\begin{remark}[点曲線積と残余比較]
\label{mbfm:rem:comparison-boundary}
命題\ref{mbfm:prop:product-family}は、点曲線積という可換標的の族を
全$\OO$開上で構成する。非可換同変残余圏からこの族への比較核を自動的には与えない。
その一般点比較には系\ref{mbrd:cor:nc-product}の入力が必要であり、
相対比較では核の実現とrotation・周期を別に確認する。
また、点曲線の組のモジュライへの射と、
中心関係の係数を含む一つの分離和三次超曲面への射を区別する。
後者に必要な中心三次線と点三次線の自然な比較は、
命題\ref{mbpl:prop:line-comparison}で最大可換商を介して非空開上に構成する。
従って定理\ref{mbpl:thm:commutative-family}の可換三次族と、
本節の点曲線積の族はともに得られるが、非可換残余圏との相対核はなお別に必要である。
\end{remark}

% !TEX root = ../nc-cubic.tex
\section{最大可換商による中心直線と点三次直線の比較}
\label{mbpl:sec:point-line}

中心三次と点曲線の方程式はともに共通 scalar の重み三を持つが、
この事実だけから両直線を自然には同定できない。
本節では、同じ代数の最大可換商へ両直線を写して比較する。
二つの像が可逆になる非空開上では、係数の選択なしに中心直線を点三次直線へ移せる。
その結果、満支持の分離和を全 $\OO$-開上の滑らかな可換三次族として降下し、
両側の表示 gerbe と関係直線を同期できる。
この可換族と非可換残余圏の同値は、既存の一般点比較がある族に限って述べる。

\subsection{共通の一次元商と自然な線同型}

第\ref{mbs:thm:presentation-stack}定理の一因子の楕円開を用いる。
普遍因子を $S=T(W)/(R_2)$、$\operatorname{rank}W=
\operatorname{rank}R_2=3$ と書き、中心三次直線を
$\mathcal C=Z(S)_3$ とする。
点曲線 $E\subset\mathbb P(W^\vee)$ の方程式直線を
\[
 \mathcal J=\ker\bigl(\operatorname{Sym}^3W
                  \longrightarrow\pi_*\OO_E(3)\bigr)
       \subset\operatorname{Sym}^3W
\]
と書く。これは階数一の直和因子であり、任意底変更に両立する。
ここで $\mathbb P(W^\vee)$ の斉次座標環を $\operatorname{Sym}W$ とする。

最大可換代数商を
\begin{equation}\label{mbpl:eq:abelianization}
 Q=S_{\mathrm{ab}}
   :=S/\langle ab-ba:a,b\in S\rangle
   =\operatorname{Sym}W/(\overline R_2)
\end{equation}
と置く。これは交換子の\emph{線型包}による商とは異なる。
次数三は有限表示
\begin{equation}\label{mbpl:eq:finite-presentation}
 R_2\otimes W\longrightarrow\operatorname{Sym}^3W
       \longrightarrow Q_3\longrightarrow0
\end{equation}
を持つ。左辺の階数は九、中央の階数は十である。
従って各幾何ファイバーで $\dim Q_3\geq1$ であり、
左の射の階数が九となる条件は開条件である。
この開を $\mathcal T^{\mathrm{ab}}$ と書く。

二つの自然な射
\begin{equation}\label{mbpl:eq:two-maps}
 c:\mathcal C\longrightarrow S_3\longrightarrow Q_3,
 \qquad
 j:\mathcal J\longrightarrow\operatorname{Sym}^3W\longrightarrow Q_3
\end{equation}
を取る。$\mathcal T^{\mathrm{ab}}$ 上で両射が可逆となる開を
$\mathcal T^\natural$ と定める。

\begin{proposition}[中心三次と点三次の自然な比較]
\label{mbpl:prop:line-comparison}
$\mathcal T^\natural$ 上には一意な線同型
\begin{equation}\label{mbpl:eq:canonical-comparison}
 \vartheta:\mathcal C\xrightarrow{\sim}\mathcal J,
 \qquad j\vartheta=c
\end{equation}
がある。この同型は $\GL(W)$-同変であり、次数付き代数の全同型と任意底変更に両立する。
$\mathcal T^\natural$ 自身も $\GL(W)$-不変な開である。
\end{proposition}
\begin{proof}
階数九の小行列式を逆にした開では
式\eqref{mbpl:eq:finite-presentation}の左の射は直和因子への同型である。
従って $Q_3$ は可逆加群で、その形成は非平坦・非被約な底変更にも両立する。
その上で $c,j$ は線束間の射なので、可逆となる条件も開条件である。
$\vartheta=j^{-1}c$ と定めれば存在と一意性が従う。
最大可換商、点曲線の方程式直線、中心および二つの射は全て代数同型に自然である。
従って全構成は同変である。
\end{proof}

\begin{remark}[Jacobian chart における意味]
\label{mbpl:rem:jacobian-chart}
体上で $\varepsilon\ne0$ とし、関係を
$[y,z]=\varepsilon\operatorname{sym}(g_x)$ とその巡回版で書くと、
\[
 Q=K[x,y,z]/(g_x,g_y,g_z).
\]
$V(g)$ が滑らかなら三つの偏微分は二次正則列であり、
$\operatorname{Hilb}Q=(1+t)^3$、特に $\dim Q_3=1$ である。
点三次の代表 $J=g+\varepsilon^2\det\operatorname{Hess}(g)/24$ は
\[
 [J]=\frac{\varepsilon^2}{24}[\det\operatorname{Hess}(g)]\in Q_3
\]
へ写る。Euler 恒等式により $[g]=0$ を使った。
ここではさらに中心の像 $[\Omega]$ との比を取る必要があり、
$J$ の方程式だけからその比を省略してはならない。
点三次の Hessian 公式は既知の
\cite[Theorem~7.2]{MoriSmith}、本稿の規約では
式\eqref{mpc:eq:point-determinant}による。
\end{remark}

\subsection{非空性を与える明示的な中心恒等式}

\begin{lemma}[Fermat 型の厳密計算]
\label{mbpl:lem:fermat}
$K=\mathbb C(t)$ とし、
\[
 S_t=K\langle x,y,z\rangle/
 ([y,z]-t x^2,\ [z,x]-t y^2,\ [x,y]-t z^2)
\]
と置く。$C=x^3+y^3+z^3$、$U=xyz+yzx+zxy$ と書くと
\begin{equation}\label{mbpl:eq:fermat-center}
 \Omega_t=(3+t^3)C-2t^2U
\end{equation}
は非零の中心三次である。点三次の代表は
\begin{equation}\label{mbpl:eq:fermat-point}
 F_t=x^3+y^3+z^3+t^2xyz
\end{equation}
であり、最大可換商では
\begin{equation}\label{mbpl:eq:fermat-images}
 Q=K[x,y,z]/(x^2,y^2,z^2),\qquad
 [\Omega_t]=-6t^2xyz,\quad [F_t]=t^2xyz.
\end{equation}
従ってこの因子は $\mathcal T^\natural$ に属し、
$\vartheta(\Omega_t)=-6F_t$ である。
同じ表示を $R=\mathbb C[t]_{(t)}$ 上で取ると、各次数は有限自由であり、
$\Omega_t$ は $R$ 上 primitive な中心三次である。
\end{lemma}

\begin{proof}
以下の交換子計算は三つの表示関係だけを使う。
$D_x=[yz,x^2]$ と置くと
\[
 [C,x]=t(zy^2z-yz^2y)=t^2D_x.
\]
また
\[
 [U,x]=D_x+t(y^3x-xz^3).
\]
関係を二回代入すれば
\[
 D_x=t(y^3x+xy^3-z^3x-xz^3)
     =2t(y^3x-xz^3)-t[C,x]
\]
となるので、
\[
 [U,x]=\frac{3+t^3}{2}D_x.
\]
従って $[\Omega_t,x]=0$ であり、巡回置換から $y,z$ とも可換する。
最大可換商の表示と式\eqref{mbpl:eq:fermat-images}は直接従う。
$xyz\ne0$ なので、$\Omega_t$ は非零である。

$g=C/3$, $\varepsilon=t$ を点三次の公式へ代入する。
$\det\operatorname{Hess}(g)=8xyz$ だから、
$J=(C+t^2xyz)/3$ となり、式\eqref{mbpl:eq:fermat-point}を得る。
これは Hesse パラメータ $(a:b:c)=(1:-1:-t)$ の通常の Sklyanin 表示である。
点曲線が滑らかな一般パラメータでは三次元二次 AS 正則代数となる
\cite[Section~1, Theorems~6.6, 6.8]{GRSklyaninATV}。
式\eqref{mbpl:eq:fermat-point}は $t=0$ で滑らかな Fermat 三次に特殊化するので、
一般点で滑らかである。中心三次の次元が一となる基本楕円開では
$\Omega_t$ がその直線を生成し、最後の結論を得る。

$R$ 上の各次数は有限生成加群である。
一般ファイバーの Hilbert 級数は $(1-z)^{-3}$、特殊ファイバーは三変数多項式環なので、
各次数の一般階数と特殊ファイバーの次元は一致する。
DVR 上の有限加群の構造定理により torsion 部分はなく、各次数は自由である。
中心恒等式は $R$ 上でも成立し、$\Omega_t\bmod t=3C\ne0$ なので primitive である。
\end{proof}

\begin{proposition}[良い楕円開および全構造対象開との交わり]
\label{mbpl:prop:nonempty-open}
$\mathcal T^\natural$ は非空であり、任意の有限個の条件
$mp\ne0$ を加えた良い楕円開と交わる。
各因子を $\mathcal T^\natural$ に制限した満支持表示の開
$\mathfrak P_k^\natural$ も非空であり、
系\ref{mbrr:cor:stack-isomorphism}の全 $\OO$-認識開との交わりは非空である。
\end{proposition}

\begin{proof}
前補題の族を $t=0$ の trait 上で見る。
点三次 $F_t$ の族はこの trait 上で滑らかである。
点対応の行列は $t=0$ で外積行列となり、曲線上で常に階数二を持つ。
従ってその核による点自己同型は曲線族上へ延長され、特殊ファイバーでは恒等になる。
Sklyanin の楕円表示ではこれは平行移動である。
$t\ne0$ で恒等なら点対応の対角上の関係から
$x^2=y^2=z^2=0$ となり、射影曲線に点がなく矛盾する。
よって平行移動切断は非零で、特殊化は零である。
標数零の楕円曲線の形式群に非零有限位数点はない。
従って一般点の平行移動は無限位数であり、全ての有限個の捩れ条件を避ける。
Hesse の $j$ 関数\eqref{gre:eq:hesse-j}へ $\mu=t^2/3$ を代入すると非定数なので、
$j=0,1728$ も一般点では避けられる。

各因子について独立な $t_i$ を取れば $j_i$ は相異なり、満支持係数も取れる。
従って $\mathfrak P_k^\natural$ は非空である。
表示スタックは既約なので、既存の非空な全 $\OO$-認識開との交わりも非空になる。
\end{proof}

\subsection{可換分離和の自然な族と表示障害の同期}

$\mathfrak P_k^\natural$ の実表示 gerbe 上で、
中心直線 $\ell\subset\bigoplus_i\mathcal C_i$ に
\begin{equation}\label{mbpl:eq:point-relation-line}
 \ell_{\mathrm{pt}}=
 \left(\bigoplus_i\vartheta_i\right)(\ell)
 \subset\bigoplus_i\mathcal J_i\subset\operatorname{Sym}^3V,
 \qquad V=\bigoplus_iW_i
\end{equation}
を対応させる。

\begin{theorem}[全構造対象開上の可換三次族]
\label{mbpl:thm:commutative-family}
式\eqref{mbpl:eq:point-relation-line}は
$H_k=\GL_3^k\rtimes\mathfrak S_k$ に同変であり、関係直線の自然な同型
\begin{equation}\label{mbpl:eq:relation-synchronization}
 \ell\xrightarrow{\sim}\ell_{\mathrm{pt}}
\end{equation}
を伴う。これによって
\[
 Y_{\mathrm{pt}}=V(\ell_{\mathrm{pt}})\subset\mathbb P(V^\vee)
\]
は $\mathfrak P_k^\natural$ 上へ降下する滑らかな三次 $3k-2$ 次元族を定める。
降下した ambient は相対次元 $3k-1$ の Brauer--Severi scheme であり、
大域的な次数一ベクトル束の存在は要求しない。

滑らかな三次超曲面の射影的表示スタックを
$\mathfrak C_{3k}^{\mathrm{sm}}$ とすると自然な射
\begin{equation}\label{mbpl:eq:moduli-map}
 \Phi_{\mathrm{pt}}:\mathfrak P_k^\natural
       \longrightarrow\mathfrak C_{3k}^{\mathrm{sm}}
\end{equation}
がある。全 $\OO$-認識開との交わりでは、これは全 $\OO$-局所型モジュライ上の射となる。
構成はこの開の任意の非被約試験底にも両立する。

両側の実表示 gerbe の類を $\beta,\beta_{\mathrm{cub}}$ とすると
\begin{equation}\label{mbpl:eq:brauer-synchronization}
 \beta=\Phi_{\mathrm{pt}}^*\beta_{\mathrm{cub}}.
\end{equation}
この等式は band を保つ gerbe の同値で実現され、
式\eqref{mbpl:eq:relation-synchronization}も同時に保持される。
\end{theorem}

\begin{proof}
命題\ref{mbpl:prop:line-comparison}の自然性により各因子の同型は線型座標変更に両立する。
同じ構成を全因子に用いるので置換にも両立する。
満支持条件により $\ell_{\mathrm{pt}}$ の各 $\mathcal J_i$ への射は可逆である。
局所的に方程式を $\sum_i b_iJ_i$、$b_i\in R^\times$ と書くと、
各 $V(J_i)$ は滑らかな平面三次である。
勾配はブロックごとに分かれ、全偏微分の共通アフィン零点は原点だけなので、
分離和の射影超曲面は全幾何ファイバーで滑らかである。
相対 Cartier 方程式の平坦性と Jacobian 判定により、族も滑らかである。

共通スカラーは $\mathbb P(V^\vee)$ に自明に作用し、
$\ell_{\mathrm{pt}}$ と $\operatorname{Sym}^3V$ に同じ三乗で作用する。
従って射影 ambient と超曲面は $G_k$ へ降下する。
projective frame を局所的に取れば滑らかな三次の係数空間への射を得て、
その $\PGL_{3k}$ 商が式\eqref{mbpl:eq:moduli-map}である。
全 $\OO$-射についての主張は、表示スタックとの同型
系\ref{mbrr:cor:stack-isomorphism}を使う。

実表示からは $V$ と実際の方程式直線 $\ell_{\mathrm{pt}}$ が得られるため、
元の表示 gerbe から、標的の実表示 gerbe を引き戻したものへの射がある。
$H_k\to\GL(V)$ は共通スカラー $\lambda$ を同じ $\lambda$ に送るので、
band 上の写像は恒等である。従ってこの二つの $\mathbf G_m$-gerbe は同値であり、
式\eqref{mbpl:eq:brauer-synchronization}を得る。
関係直線の比較は最初から式\eqref{mbpl:eq:point-relation-line}に含まれている。
以上は有限局所自由な加群の射と同変降下から作られ、任意底変更に両立する。
\end{proof}

\begin{remark}[局所方程式の係数を省略しない]
\label{mbpl:rem:scale}
中心生成元と点方程式を独立に $\Omega_i,J_i$ と選び、
$w=\sum_i a_i\Omega_i$ と書いた場合には、
\[
 \vartheta_i(\Omega_i)=u_iJ_i,
 \qquad
 F_{\mathrm{pt}}=\sum_i a_iu_iJ_i,
 \qquad
 u_i=[\Omega_i]/[J_i]\in R^\times
\]
である。$u_i$ は共通の一次元商 $Q_{i,3}$ 内の比であり、任意には選ばない。
例えば前補題の frame では $u_i=-6$ である。
$F=\sum_i a_iJ_i$ と書けるのは、点方程式を
$J_i=\vartheta_i(\Omega_i)$ と規格化した場合である。
共通 scalar 重み三だけを理由に $u_i=1$ と置く操作は用いない。
\end{remark}

\subsection{既存の一般点比較との接続と境界}

\begin{corollary}[Jacobian tensor 一変数族内での残余比較]
\label{mbpl:cor:generic-comparison}
第\ref{mpc:sec:point-comparison}節の入力と PBW 仮定を満たす一変数族を取り、
その一般点が $\mathfrak P_k^\natural$ に属すると仮定する。
するとある有限代数拡大後に、同節の中心群同変・rotation・指定周期保存の比較の標的を、
定理\ref{mbpl:thm:commutative-family}の可換族の一般ファイバーへ取れる。
これは同節の一変数族での主張であり、全モジュライの一般対象についての
残余比較を新たに証明するものではない。
\end{corollary}

\begin{proof}
$K=\mathbb C(h)$ 上で、同節の点方程式の frame を $J_i$ と書く。
式\eqref{mpc:eq:algebraic-potential}の既存標的は
$H_{\mathrm{pt}}=\sum_i c_i\lambda_iJ_i$ であり、
本節の標的は $F_{\mathrm{pt}}=\sum_i c_iu_iJ_i$ である。
各 $u_i,\lambda_i$ は $K^\times$ に属する。
有限拡大で $s_i^3=\lambda_i/u_i$ を満たす $s_i$ を選べば
\[
 F_{\mathrm{pt}}(s_1x_1,\ldots,s_kx_k)=H_{\mathrm{pt}}.
\]
このブロック対角座標変換は中心三乗根群と可換し、potential を正確に移す。
従って定理\ref{mpc:thm:point-model}の比較と合成でき、
内部 shift と指定周期も保たれる。
\end{proof}

\begin{remark}[開の非空性と未主張の境界]
例えば $t_i=d_i h$、$d_i\in\mathbb C^\times$、$d_i^6$ 相異なるとして
補題\ref{mbpl:lem:fermat}の因子を tensor すれば、一般点で各比較射は可逆となる。
$j(E_{t_i})=512d_i^6h^6+O(h^{12})$ なので各 $j_i$ も相異なる。
この例は $g_i=(x_i^3+y_i^3+z_i^3)/3$ と
$\varepsilon_i=d_i h$ の Jacobian tensor 族である。
補題\ref{mbpl:lem:fermat}の DVR 上の自由性が PBW 入力を与え、
$\Omega_{d_ih}$ は primitive、満支持の特殊三次は Fermat 型の滑らかな分離和になる。
従って実際に第\ref{mpc:sec:point-comparison}節と直前の系の適用範囲に入る。
この族での交わりの非空性と、全モジュライの既約性から得る
全 $\OO$-認識開との交わりの非空性は、異なる主張として用いている。

$\mathcal T^\natural$ は、有限表示\eqref{mbpl:eq:finite-presentation}と
二つの射\eqref{mbpl:eq:two-maps}をともに可逆にするための開である。
その外で別の自然な線比較が存在しないとは主張しない。
$Q_3$ の階数が増える場合、または二つの射の一方が零になる場合には、
本節の構成だけから延長は得られない。
また表示 gerbe と関係直線の同期は、比較核や rotation 比較の大域的 section の存在、
全点での smooth proper 残余比較、指定形式核の全 jet の代数化を与えない。
\end{remark}

% !TEX root = ../nc-cubic.tex
\section{比較核・関係直線・三次周期の相対同期}
\label{mbks:sec:synchronization}

本節は、比較核の stack の降下と、その大域切断の存在を区別する
命題\ref{mbrd:prop:kernel-stack}を、rotation と指定周期を付加した場合に精密化する。
一つの underlying 核の型を固定すると、残る障害は核の scalar gerbe と、
比較同型を入れる直線束の三次 Kummer 類に分かれる。
表示 gerbe の同期を固定しない場合には、この後者も固定された不変量ではない。

体上の周期位相と、核の scalar 慣性が周期標識後も残ることは、
式\eqref{eq:nfm-period-phase-convention}--\eqref{eq:nfm-period-phase-kummer}および
注意\ref{rem:nfm-full-kernel-period-nine}ですでに扱った。
ここでは係数直線、相対 Picard 群、全構造対象局所型への降下を保持する。
用いる Kummer 完全列と、三乗を自明化した直線束による
$H^1(\mu_3)$ の記述は既知である\cite{NFMStacksKummer}。
以下の相対比較への適用について、初出の網羅的な確認はしていない。

\subsection{固定した比較核に対する直線と三次形式}

$S$ は $\C$ 上の scheme または代数 stack とする。
以下の層とコホモロジーは fppf 位相で扱う。
$3$ は可逆なので、現れる $\mu_3$-torsor は有限 \'etale である。
$\mathcal C_s,\mathcal C_t$ を $S$-線型 dg 圏の族とし、
fpqc 局所に smooth proper な dg 代数で表されると仮定する。
両側に実際の可逆核 $\rho_s,\rho_t$ と直線束 $\ell_s,\ell_t$、
指定された中心的周期
\begin{equation}\label{mbks:eq:periods}
 \pi_i:\rho_i^3\xrightarrow{\sim}[2]\otimes\ell_i,\qquad i=s,t
\end{equation}
を与える。ここで中心的とは、rotation、shift、底の直線束との標準交換同型に関し、
周期を左右どちらから適用しても一致するという意味であり、
その dg coherence もデータに含める。
本稿の残余圏へ適用する際の周期の存在には、
命題\ref{mbsd2:prop:line-period}の MF/window 比較を引き続き要求する。

$P:\mathcal C_s\simeq\mathcal C_t$ を実際の Morita 核とする。
任意底変更後に
\begin{equation}\label{mbks:eq:simple}
 \mathcal H^a\mathbf R\operatorname{Hom}(P,P)=0\quad(a<0),
 \qquad
 \mathcal H^0\mathbf R\operatorname{Hom}(P,P)=\OO_S
\end{equation}
が成立し、これらの同定が底変更と両立すると仮定する。
また $P\rho_s$ と $\rho_tP$ が fpqc 局所で同型であるとする。
有限 \'etale 群の coherent 作用を同時に扱う場合には、
核、局所同型、周期をすべて同変なものとし、
以下の Hom も同変核の圏で取る。
標数零では不変部分を取ることが完全であり、
underlying な式\eqref{mbks:eq:simple}から同変版も従う。

\begin{proposition}[比較直線と関係直線の三乗則]
\label{mbks:prop:comparison-line}
以上の仮定の下で
\[
 L_P=\mathcal H^0\mathbf R\operatorname{Hom}(P\rho_s,\rho_tP)
\]
は $S$ 上の直線束であり、同型の層はその frame torsor
$\mathit{Isom}(P\rho_s,\rho_tP)=L_P^\times$ に等しい。
周期の三重図式は標準同型
\begin{equation}\label{mbks:eq:cubic-line}
 q_P:L_P^{\otimes3}\xrightarrow{\sim}
 M_{st},\qquad M_{st}=\ell_t\otimes\ell_s^\vee
\end{equation}
を与える。すべての構成は任意底変更と両立する。

指定した関係直線同型 $\theta:\ell_s\xrightarrow{\sim}\ell_t$ に対し、
周期を保つ比較同型の層は
\begin{equation}\label{mbks:eq:period-torsor}
 T_{P,\theta}
 =\{\eta\in L_P^\times:q_P(\eta^{\otimes3})=\theta\}
\end{equation}
という有限 \'etale $\mu_3$-torsor である。
従って大域比較同型が存在する必要十分条件は
\begin{equation}\label{mbks:eq:kummer-class}
 \gamma(P,\theta):=[T_{P,\theta}]=0
       \quad\text{in }H^1(S,\mu_3)
\end{equation}
である。
\end{proposition}
\begin{proof}
fpqc 局所で同型 $\eta_0:P\rho_s\to\rho_tP$ を取る。
可逆核との合成は導来 Hom を保つので、
式\eqref{mbks:eq:simple}により全ての次数零射は一意に $u\eta_0$ と書ける。
これが同型である条件は $u\in\OO^\times$ である。
従って局所 Hom は階数一自由で、重なりの移行は可逆スカラーであり、
$L_P$ と $L_P^\times$ が得られる。
負の Ext 消滅は、同型の mapping space にさらに高次の自己同型がないことを保証する。

合成の順序を右から左として
\[
 \eta^{(3)}
  =(\rho_t^2\eta)(\rho_t\eta\rho_s)(\eta\rho_s^2)
  :P\rho_s^3\longrightarrow\rho_t^3P
\]
と置く。周期で始域と終域を移すと、
\[
 (\pi_tP)\eta^{(3)}(P\pi_s)^{-1}
 :P[2]\otimes\ell_s\longrightarrow [2]P\otimes\ell_t
\]
を得る。$P$ の scalar-simple 性により、この Hom は $M_{st}$ と同定される。
この写像を
\begin{equation}\label{mbks:eq:cubic-convention}
 (\pi_tP)\eta^{(3)}
   =\bigl([2]P\otimes q_P(\eta^{\otimes3})\bigr)(P\pi_s)
\end{equation}
と記述する。$\eta$ を $u\eta$ に替えると右辺の係数は $u^3$ 倍になるので、
局所的な式は $L_P^{\otimes3}\to M_{st}$ へ貼り合わさる。
frame $\eta_0$ の三重合成は同型であり、$q_P$ も同型である。

局所 frame $e$ を選び、
$\theta^{-1}q_P(e^{\otimes3})=c\in\OO^\times$ と書けば、
式\eqref{mbks:eq:period-torsor}は $u^3=c^{-1}$ である。
これは階数三の有限 \'etale 被覆で、$\mu_3$ が $u$ の乗法で自由推移的に作用する。
異なる frame の式はこの作用を保って貼り合わさる。
torsor に大域切断があることとその類の消滅は同値である。
各段階の底変更両立性は仮定した Hom の底変更と、合成・周期の自然性から従う。
\end{proof}

\begin{corollary}[相対 Picard 成分と尺度の成分]
\label{mbks:cor:picard-phase}
関係直線同型 $\theta$ がある場合、$\gamma(P,\theta)$ の
$\operatorname{Pic}(S)[3]$ への像は $[L_P]$ である。
従って次の二つの条件を順に検査すればよい：
\begin{enumerate}
 \item $L_P$ が自明である。
 \item その frame $e$ について
 $\theta/q_P(e^{\otimes3})\in\Gamma(S,\OO_S^\times)$ が三乗である。
\end{enumerate}
第2の条件は frame の選択によらない。
$\theta$ を $u\theta$ に替えると
\begin{equation}\label{mbks:eq:theta-change}
 \gamma(P,u\theta)=\gamma(P,\theta)+[u]
\end{equation}
となる。ここで $[u]$ は $z^3=u$ の Kummer 類である。

関係直線同型をまだ指定していない場合にも、
$[M_{st}]=3[L_P]$ は必ず成立する。
$\theta$ と $\eta$ の双方を自由に選べるなら、両者の存在は $L_P$ の自明性と同値である。
固定した $\theta$ に対する周期障害と、この自由な選択問題は異なる。
\end{corollary}
\begin{proof}
Kummer 完全列\cite{NFMStacksKummer}の部分
\[
 0\longrightarrow
 \Gamma(S,\OO_S^\times)/\Gamma(S,\OO_S^\times)^3
 \longrightarrow H^1(S,\mu_3)
 \longrightarrow\operatorname{Pic}(S)[3]\longrightarrow0
\]
を、対 $(L_P,\theta^{-1}q_P)$ に適用する。
残りは式\eqref{mbks:eq:cubic-line}と局所式 $u^3=c^{-1}$ から従う。
scheme に対する通常の証明は、直線束とその frame の fppf 降下により
代数 stack にも同じ形で適用できる。
\end{proof}

\subsection{核の型を固定した stack と二つの障害}

次に $P$ 自体は大域的に選ばず、
underlying Morita 核の局所同型型を $S$ 上で一つ固定する。
その forms の stack を $\mathscr G_P$ とする。
局所核に式\eqref{mbks:eq:simple}を仮定すれば、
注意\ref{mbrd:rem:kernel-section}により $\mathscr G_P$ は
$\mathbf G_m$-gerbe となる。その類を $\alpha_P\in H^2(S,\mathbf G_m)$ と書く。
以下では、この型の局所核が rotation を局所的に intertwine することも仮定する。

\begin{theorem}[指定周期付き核の相対実現判定]
\label{mbks:thm:kernel-criterion}
上記の仮定と、指定した $\theta:\ell_s\simeq\ell_t$ の下で、
普遍核から作る $L_P,q_P,T_{P,\theta}$ は $\mathscr G_P$ から $S$ へ降下する。
周期を保つ比較同型を持つ、この型の核の stack は
\begin{equation}\label{mbks:eq:product-stack}
 \mathscr G_{P,\theta}^{\rm per}
       \simeq\mathscr G_P\times_S T_{P,\theta}
\end{equation}
である。従って大域的な核と比較同型の対が存在する必要十分条件は
\begin{equation}\label{mbks:eq:two-obstructions}
 \alpha_P=0,\qquad \gamma(P,\theta)=0.
\end{equation}
一つの実際の $P$ とそれへの同定を固定した fiber は $T_{P,\theta}$ である。
一方、核を動かす stack の自己同型群は周期標識後も $\mathbf G_m$ であり、
$\mu_3$ に縮まるわけではない。
\end{theorem}
\begin{proof}
局所核の scalar 自己同型 $v\operatorname{id}_P$ は、
$P\rho_s$ と $\rho_tP$ の双方に同じ $v$ を掛ける。
従って Hom による $L_P$ への作用は自明である。
局所核の同型による移送も、この同型を scalar 倍しても変わらない。
そのため $L_P$、合成による $q_P$、および $T_{P,\theta}$ の降下 cocycle は
実際の等式として成立する。

降下後の torsor の $U$ 上の切断を与えることは、
$U$ 上の任意の局所核に周期を保つ比較同型を与え、
核の全ての同型による移送と両立させることに等しい。
これが式\eqref{mbks:eq:product-stack}を与える。
大域切断の存在は、二つの因子のそれぞれに大域切断があることと同値である。
また任意の scalar 自己同型は
\[
 (\rho_t(v\operatorname{id}_P))\eta
      =\eta((v\operatorname{id}_P)\rho_s)
\]
を満たすので、$\eta$ を保つ。逆に核の自己同型は全て scalar であり、
自己同型群は $\mathbf G_m$ のままである。
負の自己 Ext 消滅により、ここで省略された高次慣性はない。
\end{proof}

\begin{remark}[実現型と比較の合成]
\label{mbks:rem:composition}
同じ核の型の別の form を選んでも、$L_P,q_P,\gamma(P,\theta)$ は変わらない。
特に $P$ を底の直線束で tensor しても、核の重みは Hom の両側で相殺する。
これは後述の rotation の同期自体を直線束で変更する操作とは異なる。

二つの比較 $P:s\to t$, $Q:t\to u$ について同じ scalar-simple 条件を満たすなら、
比較同型の合成は
$L_{QP}\simeq L_Q\otimes L_P$ を与える。
関係直線同型を $\theta_{su}=\theta_{tu}\theta_{st}$ と選んだ場合、
三重図式の合成により
$\gamma(QP,\theta_{su})=\gamma(Q,\theta_{tu})+\gamma(P,\theta_{st})$
となる。これは比較核の合成の coherence を用いる恒等式であり、
同型類の不変性だけから降下を推論しているのではない。
\end{remark}

\subsection{表示 gerbe の同期と全構造対象族への適用}

$\mathscr G_s,\mathscr G_t\to S$ を二つの表示 $\mathbf G_m$-gerbe とし、
類を $\beta_s,\beta_t$ とする。
rotation はそれぞれの gerbe 上で重み一、$\ell_s,\ell_t$ は重み三とする。
この段階では rotation 自体が $S$ 上の実核であることは要求しない。

\begin{proposition}[比の gerbe と同期選択の変換則]
\label{mbks:prop:gerbe-synchronization}
band を保つ同値 $b:\mathscr G_s\simeq\mathscr G_t$ が存在する必要十分条件は
$\beta_t-\beta_s=0$ である。
その一つを固定し、両側を共通 gerbe $\mathscr G_s$ 上で比較する。
このとき $M_{st}=\ell_t\ell_s^\vee$ は重み零であり、$S$ 上の直線束となる。
核の型の gerbe $\mathscr G_P$ をさらに取った場合にも、
普遍 Hom 直線の二つの scalar 重みはともに零であるから、
命題\ref{mbks:prop:comparison-line}と
定理\ref{mbks:thm:kernel-criterion}は $S$ 上に降下して成立する。
ここでの $\eta$ は、共通 gerbe 上の二つの重み一 rotation を比較するデータであり、
個々の rotation の $S$ 上への実現を意味しない。

同期 $b$ を底の直線束 $N$ による target 側の twist で変更すると
\begin{equation}\label{mbks:eq:synchronization-change}
 \rho_t\longmapsto\rho_t\otimes N,\qquad
 \ell_t\longmapsto\ell_t\otimes N^{\otimes3},\qquad
 L_P\longmapsto L_P\otimes N,\qquad
 M_{st}\longmapsto M_{st}\otimes N^{\otimes3}.
\end{equation}
従って $\gamma(P,\theta)$ は、固定した同期と関係直線同型に対する障害である。

逆に、一つの核の型が $S$ 上で選ばれており、
$\mathscr G_P\times_S\mathscr G_s\times_S\mathscr G_t$ 上の普遍核が
rotation を fpqc 局所で intertwine するなら、
その Hom 直線の存在だけで $\beta_t-\beta_s=0$ が必要となる。
\end{proposition}
\begin{proof}
最初の同値は、同じ band の gerbe の同型類が $H^2(S,\mathbf G_m)$ で
分類されることの適用である。局所表示の二コサイクルで書けば、
両 cocycle の比が coboundary となることが、band を保つ同期の条件である。
同期後は rotation の重みが双方で一となり、Hom では相殺する。
核の scalar も両側で相殺する。
重み零の直線束とその写像は gerbe から底へ降下するため、
前二小節の証明をそのまま適用できる。

二つの同期の差は $\mathbf G_m$-gerbe の band を保つ自己同値であり、
その同型類は底の $\mathbf G_m$-torsor、すなわち直線束 $N$ で記述される。
重み $j$ の対象はこの変更で $N^{\otimes j}$ 倍されるので、
式\eqref{mbks:eq:synchronization-change}を得る。

最後の主張では、同期を選ぶ前の Hom 直線の重みは
核・source・target の順に $(0,-1,1)$ である。
まず核の gerbe の重み零を降ろす。
残った直線束を局所表示で trivialize すると、その遷移の三重積は
target の表示 cocycle と source の逆 cocycle の積である。
従ってこの比は coboundary となり、$\beta_t-\beta_s=0$ を得る。
これは任意の underlying 同値に対する結論ではなく、
一つの大域的な核の型が局所 rotation 型を保つという仮定の必要条件である。
\end{proof}

\begin{corollary}[全構造対象開上の条件付き判定]
\label{mbks:cor:all-o-criterion}
系\ref{mbrr:cor:stack-isomorphism}の全構造対象局所型の開から、
二つの残余族と中心作用を $S$ 上へ引き戻す。
比較する両族が命題\ref{mbrd:prop:kernel-stack}の smooth proper 仮定を満たし、
指定周期の MF/window 比較、核の型、普遍 scalar-simple 条件、
局所 rotation 比較を与えたとする。
さらに表示 gerbe の同期 $b$ と関係直線同型 $\theta$ を固定する。
このとき、同変な比較核・rotation 比較・指定周期を持つ stack は
全 $\OO$-表示変更に関して fpqc 降下し、その実際の大域切断の存在は
式\eqref{mbks:eq:two-obstructions}と同値である。

この判定は、任意の全 $\OO$-底上で $\alpha_P$ や $\gamma(P,\theta)$ が消えること、
または非正則な底で全 coherent 残余と compact 残余が一致することを述べない。
本稿の自然な可換標的は定理\ref{mbpl:thm:commutative-family}で構成し、
順序付き実表示上の相対模型と Picard 補正型の具体的判定には
定理\ref{mbrn:thm:relative-product}と系\ref{mbpd:cor:descent-criterion}を用いる。
\end{corollary}
\begin{proof}
残余族・中心作用の降下は系\ref{mbrd:cor:residual-family}による。
比較核と追加の同型の stack の降下は命題\ref{mbrd:prop:kernel-stack}を用いる。
命題\ref{mbks:prop:gerbe-synchronization}により Hom と周期のデータは重み零である。
あとは定理\ref{mbks:thm:kernel-criterion}を適用する。
\end{proof}

\begin{remark}[同期を自由に変える問題との違い]
\label{mbks:rem:free-synchronization}
式\eqref{mbks:eq:synchronization-change}で $N=L_P^\vee$ を選び、
関係直線同型も $q_P$ に従って選び直せば、
比較直線とその三乗はともに自明化できる。
従って同期と指定周期を自由に変更してよい問題には、
固定した $(b,\theta)$ の非零な $\gamma$ をそのまま障害として使えない。
定理\ref{mbpl:thm:commutative-family}の中心関係・点三次の比較は、
基本楕円開全体で自然な $b$ と $\theta$ を与える。
この同期を保持した問題では上の変更は許されず、
Picard 補正型の二つの障害は系\ref{mbpd:cor:descent-criterion}で記述される。
自然な同期に対する非零位相は命題\ref{mbpd:prop:boundary-phase}で検出する。
全ての核の型や周期三乗写像の全係数の計算は、この指定型の結果に含めない。
\end{remark}

\subsection{非消滅例と因子置換}

\begin{example}[相対 Picard 障害と単数の Kummer 障害]
\label{mbks:ex:picard-unit}
滑らかな平面三次 $E$ 上の
$R=T_{\OO_E}\circ(-\otimes\OO_E(1))$ は $R^3\simeq[2]$ を持つ
（命題\ref{mbrd:prop:factor-rotations}）。
滑らかな楕円曲線 $S$ 上で、非自明な $N\in\operatorname{Pic}^0(S)[3]$ と
$v:N^{\otimes3}\simeq\OO_S$ を取る。
定数族 $\operatorname{Perf}(E\times S/S)$ で
\[
 P=\operatorname{id},\qquad \rho_s=R,\qquad \rho_t=R\otimes N,\qquad
 \ell_s=\ell_t=\OO_S
\]
とし、target の周期に $v$ を用いる。
このとき $L_P=N$, $q_P=v$ である。
両表示 gerbe は中立、実際の $P$ と関係直線同型も存在するが、
$L_P^\times(S)=\varnothing$ なので大域 rotation 比較は存在しない。
smooth proper 性、負の自己 Ext 消滅、$\operatorname{End}(P)=\OO_S$ と
その底変更は定数族の Hochschild 複体から従う。

一方、$S=\operatorname{Spec}\C(t)$ とし、同じ $R$ を両側に取る。
source の周期を $\pi$、target の周期を $t^{-1}\pi$ とし、
$\theta=\operatorname{id}$ とする。
この場合は $L_P=\OO_S$ であるが、
式\eqref{mbks:eq:cubic-convention}から周期条件は $\eta^3=t$ となる。
$t=0$ における付値が一なので $t$ は $\C(t)$ の三乗ではなく、
$T_{P,\theta}$ は非自明な三次 torsor である。
同じ例は $\operatorname{Spec}\C[t,t^{-1}]$ 上の有限 \'etale 被覆 $u^3=t$ として表せ、
どの非空 Zariski 開へ縮めても切断を持たない。

これらは $E$ の代わりに二曲線の積を取り、
系\ref{mbrd:cor:symmetric-rotation}の
$R_{\rm sym}=(R_1R_2)^{-1}[2]$ を用いても成立する。
従って $3\nmid k$ による対称 split の存在だけでは、相対 Picard 障害も
周期の Kummer 障害も消えない。
ここで変更した rotation や周期は意図的に別のデータであり、
本稿の標準量子化族の自然な同期でこれらの類が非零であると主張する例ではない。
自然な Sklyanin 同期を保持した非零位相と独立な位相の例は、
命題\ref{mbpd:prop:boundary-phase}と系\ref{mbpd:cor:phase-independence}で与える。
\end{example}

\begin{proposition}[整合した置換作用の下での位相切断]
\label{mbks:prop:permutation-phase}
$U$ を連結な $\C$-scheme、$H=\mathfrak S_k$、$k\ge2$ とし、
$[U/H]$ 上で前節の全てのデータがすでに coherent に与えられているとする。
特に核の型、rotation、関係直線比較、指定周期も置換で保たれ、
$H$ は係数 $\mu_3$ に自明に作用する。
$T_{P,\theta}|_U$ に切断があれば、その切断は自動的に $H$-不変であり、
$[U/H]$ 上へ降下する。
同値に、この状況の引戻し
\[
 H^1([U/H],\mu_3)\longrightarrow H^1(U,\mu_3)
\]
は単射である。
\end{proposition}
\begin{proof}
切断 $e$ とその移送 $g(e)$ の比は
$a_g\in\Gamma(U,\mu_3)=\mu_3(\C)$ である。
作用の coherence と係数への自明作用により
$a_{gh}=a_ga_h$ となる。
任意の置換は位数二の隣接互換で生成され、$\mu_3$ には位数二の非自明元がない。
従って $\operatorname{Hom}(\mathfrak S_k,\mu_3)=0$ であり、全ての $a_g$ は一である。
よって $e$ は不変で、torsor の fpqc 降下により商 stack の切断を与える。
\end{proof}

\begin{remark}[置換位相の消滅が使える範囲]
\label{mbks:rem:permutation-scope}
命題\ref{mbks:prop:permutation-phase}は、ordered 被覆上での位相の消滅そのものを
与えるものではない。
Picard 補正型の位相とその非消滅は第\ref{mbpd:sec:picard-descent}節で扱う。
また、$3\mid k$ における rotation の標準線型化の対称選択の不存在
（系\ref{mbza:cor:permutation-obstruction}）を取り除かない。
その場合は、命題の前提となる rotation データの置換両立性を先に与える必要がある。
自然同型は次数零で、作用は $S$-線型であり、周期と関係直線比較まで
coherent に保存されるという条件を、単なる同型類の置換不変性に弱めてはいけない。

さらに核の scalar 慣性は $\mathbf G_m$ であるため、
この $\mu_3$ の議論は核自体の置換降下には適用できない。
例えば $K=\C(t)$、$Y=E\times E$ とし、因子交換の関手を $\sigma^*$ とする。
source には通常の $\mathfrak S_2$ 作用、target には同じ関手を用いるが、
$\sigma^*\sigma^*\simeq\operatorname{id}$ の乗法比較を $t$ 倍した弱作用を与える。
正規化二コサイクルの唯一の非自明な値を $c(\sigma,\sigma)=t$ と置けば、
結合則は $t$ が scalar であることから成立する。
恒等核の型は交換で不変であるが、その同変構造は
$u^2t=1$ を満たす $u\in K^\times$ を要求し、存在しない。
両族は $B_K\mathfrak S_2$ 上へ降下するが、
この型の実際の比較核の gerbe は非中立である。
対称 rotation とその三次周期は scalar cocycle と可換するため両側で保たれる。
これは整合した dg 作用から作る一般的な反例であり、
標準のブロック置換にこの余分な cocycle が発生すると断定するものではない。
具体的な外積核の Koszul symmetry が全ての置換関係を満たすと証明できれば、
その構成についてはこの別の障害を排除できる。
\end{remark}

% !TEX root = ../nc-cubic.tex
\section{積 Picard 補正型の核 gerbe と相対周期位相}
\label{mbpd:sec:picard-descent}

本節では、点曲線の各 $p_i$ を次数零線束で表すという
定理\ref{mbrc:thm:finite-point-comparison}の構成を、族の上で分類する。
分類するのは積表示上で対角線に支持される Picard 補正型であり、
任意の残余比較核の分類ではない。
この型について、scalar gerbe は因子の Picard 障害の和となる。
双対群作用の比較には因子周期位相の差が現れ、
rotation 自体の比較にはさらに共通位相が残る。
これらは Picard gerbe、Kummer 理論、有限双対同変化という既知の一般論の適用であり、
本節の相対比較への適用について初出の網羅的な確認はしていない。

\subsection{相対模型と分類する局所型}

$S$ を滑らかな準射影有限型 $\C$-scheme とし、$k\geq2$ とする。
順序付き実表示の点データ
\[
 f_i:E_i\longrightarrow S,\qquad
 L_i=\OO_{E_i}(1),\qquad p_i\in J_i(S),\qquad
 J_i=\operatorname{Pic}^0_{E_i/S}
\]
を固定する。各 $f_i$ は滑らかで固有、幾何学的に連結な genus-one 曲線族であり、
$L_i$ の相対次数は三である。
基本楕円開の条件、とくに $p_i,2p_i,3p_i\ne0$ を課す。
さらに
\[
 Y=\prod_{i=1}^k{}_S E_i,\qquad
 G=Z_0=(\mu_3^k)/\Delta\mu_3,
 \qquad H=\widehat G
\]
と置く。以下の核は相対 perfect 核または可逆 perfect dg 双加群であり、
すべての同型・作用には dg coherence を含める。

相対模型への入力を明示しておく。
定理\ref{mbrn:thm:relative-product}と
定理\ref{mbrp:thm:relative-product}の構成により、実際の $S$-線型圏
$\mathcal C_p,\mathcal C_0$ と、作用を保持する積表示
\begin{equation}\label{mbpd:eq:relative-models}
 \mathcal C_p^G\simeq\Perf(Y/S)\simeq\mathcal C_0^G
\end{equation}
を固定する。
非可換側は compact 残余圏、可換側は自然な点三次族の perfect 残余圏である。
命題\ref{mbrn:prop:rotations}の因子 rotation と可換側の相対 window は
\begin{equation}\label{mbpd:eq:factor-rotations}
 R_{i,p}=T_{\OO_{E_i}}\circ(t_{p_i})_*\circ(-\otimes L_i),
 \qquad R_{i,0}=T_{\OO_{E_i}}\circ(-\otimes L_i)
\end{equation}
を与える。ここで twist は相対評価の cone である。
定理\ref{mbpl:thm:global-line-comparison}の関係直線同期により、
両側に共通の直線束 $\ell$ と指定された中心周期
\begin{equation}\label{mbpd:eq:periods}
 \pi_{i,p}:R_{i,p}^3\xrightarrow{\sim}[2]\otimes\ell,
 \qquad
 \pi_{i,0}:R_{i,0}^3\xrightarrow{\sim}[2]\otimes\ell
\end{equation}
を固定する。以下ではこれらを勝手に scalar 倍しない。
基準因子 $b$ に対して、$H$ の生成元は
\[
 Q_{i,p}=R_{i,p}^{-1}R_{b,p},\qquad
 Q_{i,0}=R_{i,0}^{-1}R_{b,0},\qquad i\ne b
\]
であり、同じ $\ell$ を用いるため三乗の関係では係数直線が相殺する。

式\eqref{mbpd:eq:relative-models}を入力しない場合にも、以下の積曲線上の
Picard gerbe と因子位相の主張は成立する。
通常の残余圏上の核への移送には、この相対模型とその作用・周期を要求する。
任意の noetherian 底変更についての主張は、
系\ref{mbrn:cor:base-change}と系\ref{mbrp:cor:base-change}の
導来引戻しによる perfect/compact 模型の範囲で述べる。
非正則な底上の全 $D^b\operatorname{coh}$ へは読み替えない。

\subsection{外積における scalar 障害の和}

$p_i$ を表す実際の線束の gerbe を
\[
 \mathscr G_i=\mathcal{Pic}^{p_i}_{E_i/S},\qquad
 \delta_i=[\mathscr G_i]\in H^2(S,\mathbb G_m)
\]
と書く。Picard stack が Picard functor 上の gerbe であることは
\cite{NFMStacksPicardGerbe}による。
射影 $\operatorname{pr}_i:Y\to E_i$ に対して
\[
 m=\sum_i\operatorname{pr}_i^*p_i\in\operatorname{Pic}_{Y/S}(S),
 \qquad \mathscr G=\mathcal{Pic}^{m}_{Y/S}
\]
と置く。$\mathscr G$ の対象 $M$ は、fpqc 局所に
$\bigotimes_i\operatorname{pr}_i^*M_i$、$[M_i]=p_i$ と書ける。
この $M$ に対応する自己同値
$F_M=-\otimes M$ を本節の \emph{Picard 補正型} と呼ぶ。
核は $\Delta_{Y*}M$ である。

\begin{proposition}[外積 gerbe の band と Brauer 類]
\label{mbpd:prop:product-gerbe}
外積は
\[
 \prod_i\mathscr G_i\longrightarrow\mathscr G,
 \qquad (M_i)_i\longmapsto\bigotimes_i\operatorname{pr}_i^*M_i
\]
を定め、その band の写像は
$\mathbb G_m^k\to\mathbb G_m$、$(u_i)\mapsto\prod_i u_i$ である。
この band の核で rigidify した stack は $\mathscr G$ と同値であり、
\begin{equation}\label{mbpd:eq:alpha}
 \alpha_{\mathrm{Pic}}=[\mathscr G]
       =\sum_{i=1}^k\delta_i\quad\text{in }H^2(S,\mathbb G_m),
 \qquad 3\alpha_{\mathrm{Pic}}=0.
\end{equation}
従って実際の積補正 $F_M$ が存在する必要十分条件は
$\sum_i\delta_i=0$ である。
各 $\delta_i$ の個別の消滅は必要条件ではない。
\end{proposition}
\begin{proof}
各 $E_i$ および $Y$ の構造層の相対零次直像は、任意底変更後も $\OO_S$ である。
従って線束の自己同型は底の単元であり、外積の scalar はその積となる。
fppf 局所にすべての $M_i$ を選べば、写像は
$B\mathbb G_m^k\to B\mathbb G_m$ の積写像に一致する。
相対 Picard 類が $m$ である線束はこの外積と fppf 局所で同型なので、
band の核を除いた射は局所で完全忠実かつ本質的全射となり、主張した同値を得る。

同じ計算を descent cocycle で書けば、$i$ 番目の局所線束の三重重なりの
scalar $c_{i,abc}$ に対し、外積の scalar は
$\prod_i c_{i,abc}$ である。これが式\eqref{mbpd:eq:alpha}の加法則である。
必要な細分を許した fppf descent の計算なので、
単一の通常 \v{C}ech 被覆で全ての $H^2$ を表せるとは仮定していない。
補題\ref{lem:nfm-picard-tripling-line}は、$3p_i$ を表す実際の線束を
$S$ 上で与える。Picard gerbe の tensor 積に関する加法則から
$3\delta_i=\delta_{E_i}(3p_i)=0$ となる。
gerbe が対象を持つこととその類の消滅は同値であり、最後の主張が従う。
\end{proof}

\subsection{因子比較直線の Hodge 表示}

\begin{lemma}[因子の Hom 直線と周期 torsor]
\label{mbpd:lem:factor-lines}
$\mathscr G_i$ 上の普遍線束に対応する $F_i=-\otimes M_i$ について、
\[
 V_i=\mathcal H^0\mathbf R\operatorname{Hom}
       (F_iR_{i,p},R_{i,0}F_i)
\]
は直線束である。これは $S$ 上へ標準的に降下し、同じ記号で表す。
式\eqref{mbpd:eq:periods}の三重図式は
\[
 q_i:V_i^{\otimes3}\xrightarrow{\sim}\OO_S
\]
を与える。従って
\begin{equation}\label{mbpd:eq:factor-torsors}
 T_i=\{\eta_i\in V_i^\times:q_i(\eta_i^{\otimes3})=1\},
 \qquad \gamma_i=[T_i]\in H^1(S,\mu_3)
\end{equation}
は、$S$ 上の有限 \'etale な $\mu_3$-torsor とその類を定める。
これらは実線束 $M_i$ の選択によらず、任意底変更と両立する。
\end{lemma}
\begin{proof}
命題\ref{prop:nfm-elliptic-intertwiner-gerbe}を
$C=E_i$、graph $=\mathrm{id}$ に適用すれば、
$F_iR_{i,p}\simeq R_{i,0}F_i$ は fppf 局所で成立する。
この相対命題の graph/Picard 計算は非被約な底にも適用される。
対角核 $\Delta_*M_i$ は層であり、負の相対 Ext は零である。
その零次相対自己準同型は $(f_i)_*\OO_{E_i}=\OO_S$ で、任意底変更と両立する。
可逆核との合成はこれらの性質を保つ。
従って命題\ref{mbks:prop:comparison-line}を適用できる。

普遍線束の scalar は比較射の始域と終域に同じ重みで作用するため、
$V_i$、$q_i$ および $T_i$ 上では相殺する。
重み零降下、または $M_i\mapsto M_i\otimes f_i^*N_i$ における
底の線束 $N_i$ の直接の相殺から、これらは $S$ に降下する。
式\eqref{mbpd:eq:periods}では両側の関係直線を同定済みなので、
三乗写像の値域は正確に $\OO_S$ である。
\end{proof}

\begin{lemma}[比較直線は Hodge 線である]
\label{mbpd:lem:hodge-line}
$\lambda_i=(f_i)_*\Omega^1_{E_i/S}$ と置くと、
$p_i\ne0$ の開上に、任意底変更と両立する標準同型
\begin{equation}\label{mbpd:eq:hodge-line}
 V_i\simeq\lambda_i
\end{equation}
がある。特に指定周期は $q_i:\lambda_i^{\otimes3}\simeq\OO_S$ を与え、
$\gamma_i$ の $\operatorname{Pic}(S)[3]$ への像は $[\lambda_i]$ である。
\end{lemma}
\begin{proof}
添字を省き、$E\times_SE$ の対角を $\Delta$、$t_p$ の graph を $\Gamma_p$ とする。
入力座標を第一因子とする Fourier--Mukai 規約では、
相対評価の完全列
$0\to\OO(-\Delta)\to\OO\to\OO_\Delta\to0$ により、
$R_p$、$R_0$ の核はそれぞれ
\[
 K_p=\operatorname{pr}_1^*L\otimes\OO(-\Gamma_p)[1],
 \qquad K_0=\operatorname{pr}_1^*L\otimes\OO(-\Delta)[1]
\]
である。従って $F_MR_p$ から $R_0F_M$ への核の比線束は
\begin{equation}\label{mbpd:eq:kernel-ratio}
 N=\OO(\Gamma_p-\Delta)\otimes
       \operatorname{pr}_1^*M\otimes\operatorname{pr}_2^*M^{-1}.
\end{equation}
前補題の相対 intertwining により、$N$ は fppf 局所に自明であり、
射影 $g:E\times_SE\to S$ に対して
$N\simeq g^*V$ となる。
ここでは各幾何ファイバーでの自明性だけから非被約底上の自明性を推論していない。

$p\ne0$ なので $\Gamma_p\cap\Delta=\varnothing$ であり、
$\OO(\Gamma_p)|_\Delta$ は graph の標準切断で自明化される。
二つの $M$ の因子は $\Delta$ 上で標準的に相殺し、
\[
 N|_\Delta\simeq\OO(-\Delta)|_\Delta
       \simeq\Omega^1_{E/S}\simeq f^*\lambda
\]
となる。$f_*\OO_E=\OO_S$ により底の線束の引戻しは完全忠実なので、
$V\simeq\lambda$ を得る。
この構成は普遍 $M$ の scalar に不変であり、
式\eqref{mbpd:eq:hodge-line}は $S$ に降下する。
最後は Kummer 完全列と系\ref{mbks:cor:picard-phase}による。
\end{proof}

\begin{corollary}[基本 $j$-開における位相の単元表示]
\label{mbpd:cor:unit-phase}
さらに $j_i\ne0,1728$ とする。
Jacobian の Weierstrass 不変量を
$c_{4,i}\in\lambda_i^{\otimes4}$、$c_{6,i}\in\lambda_i^{\otimes6}$ と書き、
\[
 h_i=c_{6,i}/c_{4,i}\in\Gamma(S,\lambda_i^{\otimes2})^\times,
 \quad e_i=(q_i\otimes\mathrm{id}_{\lambda_i})(h_i^{\otimes2}),
 \quad u_i=q_i^{\otimes2}(h_i^{\otimes3})\in\Gamma(S,\OO_S^\times)
\]
と置く。このとき $e_i$ は $\lambda_i$ の大域 frame であり、
\begin{equation}\label{mbpd:eq:unit-phase}
 q_i(e_i^{\otimes3})=u_i^2,
 \qquad \gamma_i=[u_i]\in H^1(S,\mu_3).
\end{equation}
特にこの開では位相の Picard 成分は零であり、
$T_i$ は大域方程式 $z_i^3=u_i$ で表示できる。
\end{corollary}
\begin{proof}
局所 Weierstrass 方程式の変換則
\cite[Chapter~III, \S1, Table~3.1]{NFMSilvermanArithmetic}により、
$c_4,c_6$ はそれぞれ Hodge 線の四乗、六乗の切断である。
この変換則は任意の $\C$-代数上の座標変換恒等式であり、
Jacobian の局所 Weierstrass 表示から相対的にも貼り合わさる。
$c_4^3-c_6^2=1728\Delta$、$j=c_4^3/\Delta$ と $\Delta\ne0$ により、
仮定した開では両者は零点を持たない。
よって $h_i$ は二乗の frame であり、$q_i$ を併用すると $e_i$ は一乗の frame となる。
tensor 冪を取れば
$q_i(e_i^{\otimes3})=q_i^{\otimes4}(h_i^{\otimes6})=u_i^2$ である。
$\eta_i=s_i e_i$ の周期条件は $s_i^3=u_i^{-2}$ なので、
$z_i=u_i s_i$ と置けば $z_i^3=u_i$ を得る。
従って Kummer 類は $[u_i^{-2}]=[u_i]$ となる。
\end{proof}

ここで $q_i$ は指定された MF/window 周期の三重合成で定義された実際の写像である。
式\eqref{mbpd:eq:unit-phase}はその写像から位相を単元へ還元する公式であって、
自然な Sklyanin 周期に対する $u_i$ が全基底上で三乗になることの証明ではない。
$q_i$ の全係数に関する数値式や、$\gamma_i$ の一般点での非消滅も、
この公式だけからは結論しない。

\subsection{双対作用の位相と通常の比較核}

因子 torsor の積とその対角商を
\begin{equation}\label{mbpd:eq:phase-quotient}
 T=\prod_i T_i,
 \qquad D=T/\Delta\mu_3
\end{equation}
と置く。$T$ は $\mu_3^k$-torsor、$D$ は $G$-torsor である。
基準因子 $b$ を選ぶと、後者の $k-1$ 個の座標は
$\eta_i^{-1}\eta_b$ に対応し、
\begin{equation}\label{mbpd:eq:difference-classes}
 [D]_b=(\gamma_b-\gamma_i)_{i\ne b}
       \in H^1(S,\mu_3)^{k-1},
 \qquad T\xrightarrow{\sim}T_b\times_SD.
\end{equation}
ここで最後の射は $\eta\mapsto(\eta_b,[\eta])$ である。
$G$ と $H$ の pairing の符号を明記すると、
$(\zeta_i)\in\mu_3^k$ による変更は $Q_i$ の比較射を
$\zeta_i^{-1}\zeta_b$ 倍するという規約である。

\begin{theorem}[Picard 型の coherent 比較 stack]
\label{mbpd:thm:coherent-corrections}
$\mathscr K_{\mathrm{Pic}}^H$ を、$F_M$、$M\in\mathscr G$ と、
二つの指定された $H$ 作用を intertwine する coherent 構造の stack とする。
また $\mathscr K_{\mathrm{Pic}}^{H,\rho,\pi}$ では、
$FR_{b,p}\simeq R_{b,0}F$ という $H$-同変な rotation 比較と
式\eqref{mbpd:eq:periods}の周期条件も標識する。
このとき $S$ 上に標準同値
\begin{equation}\label{mbpd:eq:kernel-stacks}
 \mathscr K_{\mathrm{Pic}}^H\simeq\mathscr G\times_SD,
 \qquad
 \mathscr K_{\mathrm{Pic}}^{H,\rho,\pi}\simeq\mathscr G\times_ST
\end{equation}
がある。両 stack の全対象の自己同型群は $\mathbb G_m$ である。
有限双対同変化により、これらはそれぞれ
$\mathcal C_p\simeq\mathcal C_0$ の \emph{$G$-同変な Picard 補正型}の比較核、
およびその rotation・指定周期付き比較核の stack に移る。
この移送は band を保ち、任意の許された底変更と両立する。
\end{theorem}
\begin{proof}
まず fppf 局所にすべての $M_i$ と周期適合 $\eta_i$ を選ぶ。
外積 $F=\boxtimes_iF_i$ に対し、$\eta_i$ と $\eta_b$ から
$FQ_{i,p}\simeq Q_{i,0}F$ を作る。
異なる因子の比較射の次数は零であり、外積の交換則と可換する。
$Q_i^3$ を消す周期の図式は $\eta_i,\eta_b$ の周期条件により保たれる。
従って $\prod_{i\ne b}Q_i^{d_i}$、$0\leq d_i<3$ という正規形を用いれば、
交換と三乗の消去を生成関係とする全 coherent 作用が比較される。
これは定理\ref{mbrc:thm:finite-point-comparison}の coherence 証明と同じ核の計算であり、
中心周期の係数直線が共通なので相対的にも成立する。

ここで全 coherent 構造を分類する。
任意底変更後に $\Delta_{Y*}M$ の負の相対 Ext は零、
零次相対自己準同型は $\OO_S$ である。
従って一つの coherent $H$ 構造に対する別の構造は、
各 $h\in H$ について一意の scalar $a_h$ を掛けることで得られる。
積の coherence は正確に $a_{h+h'}=a_ha_{h'}$、$a_0=1$ であり、
その差の層は $\operatorname{Hom}(H,\mathbb G_m)=G$ となる。
この計算は高次の自己同型を除いているので、同型類だけの計算ではない。
一方、$\eta_i\mapsto\zeta_i\eta_i$ の差は
$Q_i$ 上で $\zeta_i^{-1}\zeta_b$ となり、
核が正確に $\Delta\mu_3$、像が全 character 群となる。
よって $T\to$ coherent 構造の層は $D$ を介して同型となる。

$M_i$ を底の線束で変更しても比較 Hom ではそれらが相殺し、
外積の scalar も比較構造を変えない。
従ってこの同定は $\prod_i\mathscr G_i$ から
命題\ref{mbpd:prop:product-gerbe}の $\mathscr G$ へ降下する。
一般の $M\in\mathscr G$ が大域的には因子線束の外積で表せない場合にも、
局所因子分解間の変更が相殺するため、$D$ から coherent 構造を作れる。
これで第一の stack 同値を得る。

次に $d\in D$ を固定する。
前の因子計算では $FR_{b,p}\to R_{b,0}F$ の比較射自体も $H$-同変である。
その Hom 線全体は $V_b$ であり、局所に一つの非零同変射があることと
scalar-simple 性により、同変 Hom も正確に $V_b$ である。
周期条件は $T_b$ となる。
式\eqref{mbpd:eq:difference-classes}により
$D\times_ST_b\simeq T$ なので第二の同値を得る。
核の scalar 自己同型はすべての比較射の共役で相殺する。
従って周期を標識しても自己同型群は $\mathbb G_m$ のままであり、
$\mu_3$ に縮まるわけではない。

最後に、$\OO_S$ は $1/|G|$ と全三乗根を含む。
Elagin の二重同変化\cite[Proposition~4.1, Theorem~4.2, \S8]{MEDElagin}の
証明で用いる平均化、有限直和、Fourier 射影はこの底上でも定義される。
作用付き dg 模型上の実際の同変 perfect 双加群へ適用すると、
同変関手および自然同型に自然な bidual の同値を与える。
逆向きはもう一度同変化することで得られるため、これは指定作用を保つ核の
groupoid を同定する。scalar 射はそのまま scalar 射へ移る。
実相対模型の有限核・projector を用いる構成なので、
導来底変更とも両立する。
\end{proof}

\begin{corollary}[二段階の降下条件]
\label{mbpd:cor:descent-criterion}
本節で指定した $G$-同変 Picard 補正型について、次が成り立つ。
\begin{enumerate}
 \item rotation 比較を標識しない実際の比較核が $S$ 上に存在する必要十分条件は
 \[
   \sum_i\delta_i=0,\qquad
   \gamma_i=\gamma_b\quad\text{for all }i
 \]
 である。
 \item $d\in D(S)$ により局所核の型を一つ選んだ後、
 その forms gerbe の類は $\alpha_{\mathrm{Pic}}=\sum_i\delta_i$ である。
 この型の rotation 比較直線は $\lambda_b$、指定周期 torsor の類は $\gamma_b$ である。
 \item rotation と指定周期まで保つ実際の比較が存在する必要十分条件は
 \[
   \sum_i\delta_i=0,\qquad \gamma_i=0\quad\text{for all }i
 \]
 である。
\end{enumerate}
基本 $j$-開では、第一の位相条件は各 $u_b/u_i$ が大域単元の三乗であること、
第三の位相条件は各 $u_i$ が大域単元の三乗であることに等しい。
\end{corollary}
\begin{proof}
式\eqref{mbpd:eq:kernel-stacks}の大域対象の存在条件を取る。
列 $1\to\Delta\mu_3\to\mu_3^k\to G\to1$ は基準因子 $b$ で分裂し、
式\eqref{mbpd:eq:difference-classes}が成り立つ。
したがって $D(S)\ne\varnothing$ と各 $\gamma_b-\gamma_i=0$ は同値である。
固定した $d$ の上で $T\to D$ のファイバーを $T_b$ へ射影すると、
同じ band の $\mu_3$-torsor の同型となる。
残りは命題\ref{mbpd:prop:product-gerbe}、
定理\ref{mbks:thm:kernel-criterion}および式\eqref{mbpd:eq:unit-phase}による。
\end{proof}

この系では、局所型を選ぶ $D$、選んだ型を実現する $\mathscr G$、
選んだ実核の周期を合わせる $T_b$ を区別している。
全因子に共通の非零位相 $\gamma_i=\gamma$ は $D$ では消えるが、
rotation の指定周期には残る。
一般の周期付き圏の族でこの現象を実現するには、定数族の全 target 周期を
$S=\Spec\C[t,t^{-1}]$ 上で同じ単元 $t^{-1}$ 倍すればよい。
その場合 $\alpha_{\mathrm{Pic}}=0$、$D$ は自明でも、残りは $[t]\ne0$ である。
これは周期標識を変更した場合の例であり、自然な Sklyanin 周期の一般値を
計算した例とはしない。

\subsection{有限 \'etale 被覆と次数上界}

\begin{proposition}[三分点による標準被覆と横断切片による被覆]
\label{mbpd:prop:finite-covers}
各 $i$ の三分点 torsor とその積を
\[
 P_i=[3]^{-1}(p_i)\longrightarrow S,
 \qquad P=\prod_i P_i
\]
と置く。これは階数 $9^k$ の有限 \'etale 被覆である。
このとき次の被覆上には、実際の比較核を構成できる：
\[
 \begin{array}{c|c|c}
  \text{被覆} & \text{階数} & \text{保持する標識}\\ \hline
  P\times_SD & 3^{3k-1} & G\text{ 作用}\\
  P\times_ST & 3^{3k} & G\text{ 作用、rotation、指定周期}.
 \end{array}
\]
これらの被覆は、固定した相対模型・周期データに対して標準的であり、
因子の置換と底変更に両立する。

$S$ が既約である場合、ある非空開 $U\subset S$ では
各曲線の相対横断直線切片 $Z_i\to U$ を階数三の有限 \'etale 被覆として選べる。
$Z=\prod_i Z_i$ とすれば、上の $P$ を $Z$ で置き換えた被覆の階数は、
それぞれ $3^{2k-1}$ と $3^{2k}$ になる。
この後半の被覆には直線の選択と開への制限を要する。
\end{proposition}
\begin{proof}
楕円曲線上の $[3]$ は有限 \'etale 階数九である
\cite{NFMStacksMultiplication}。
相対 Jacobian 上でも微分は三倍で可逆だから $[3]$ は \'etale であり、
固有かつ各ファイバーが有限なので有限となる。階数はファイバーの九に等しい。
$P_i$ 上の普遍 $q_i'$、$3q_i'=p_i$ に対して、
補題\ref{lem:nfm-picard-tripling-line}の標準線束
$\mathscr N_{L_i}$ を $q_i'$ に制限する。
これで $p_i$ を表す実際の $M_i$ を標準的に得る。
従って $P$ 上では全因子の gerbe が中立化され、外積 $M$ を選べる。
さらに $D$ または $T$ の普遍切断を用い、
定理\ref{mbpd:thm:coherent-corrections}を適用すればよい。
階数はそれぞれ $9^k\cdot3^{k-1}$ と $9^k\cdot3^k$ である。

後半について、まず $S$ の一般点の無限体上で各平面三次に横断する直線を選ぶ。
係数をある非空開へ広げ、必要ならさらに縮めれば、
交わり $Z_i$ は有限 \'etale 階数三となる。
底変更 $E_i\times_UZ_i\to Z_i$ は普遍交点 $o_i$ を切断に持つので、
\[
 M_i=\OO_{E_i\times_UZ_i}\bigl(t_{p_i}(o_i)-o_i\bigr)
\]
が $p_i$ の実線束を与える。
あとは同じ構成を適用する。これらの階数は上界であり、最小性を主張しない。
\end{proof}

\begin{corollary}[体上で rotation を標識しない場合の改良上界]
\label{mbpd:cor:field-degree}
体 $K\supset\C$ 上の順序付き実表示に対し、
ある $[K':K]=3^e$、$e\leq2k-1$ の有限分離拡大上で
\[
 \Ku_{\mathrm{coh}}(A_{K'})\simeq\Ku(X_{\mathrm{pt},K'})
\]
を与える $G$-同変な実際の核が存在する。
rotation と指定周期まで標識する場合は、
定理\ref{mbrc:thm:finite-point-comparison}の $e\leq2k$ が得られる。
\end{corollary}
\begin{proof}
補題\ref{mbrc:lem:degree-three-point}を各因子へ順次適用すれば、
高々 $k$ 回の三次拡大により全 $M_i$ を選べる。
rotation を標識しない場合には、$D$ の $k-1$ 個の $\mu_3$-torsor を
順次分裂させれば十分であり、さらに高々 $k-1$ 回の三次拡大で足りる。
体は全三乗根を含むため、非自明な Kummer 拡大は次数三である。
体上の積表示と bidual は定理\ref{mbnp:thm:product}、
系\ref{mbnp:cor:coherent}、系\ref{mbrp:cor:comparison-field}から得られる。
従ってこの系に相対模型の仮定を追加する必要はない。
指定周期を保つには全 $T_i$ を分裂させるので、追加回数は高々 $k$ となる。
\end{proof}

\subsection{二捩れ境界から得られる非零位相と鋭い位相次数}

前節までの公式は位相の自動消滅を意味しない。
実際、自然な関係直線同期のまま、非零の位相を次のように検出できる。
用いる形式基底は有限型 Hesse 族からの底変更であり、
全 $\OO$-認識を二捩れ境界へ拡張するものではない。

\begin{proposition}[自然な周期位相の二捩れ境界における剰余]
\label{mbpd:prop:boundary-phase}
$c\in\C^\times$ を、Hesse 三次
\[
 F_0=c(x^3+y^3+z^3)-(2+c^3)xyz
\]
が滑らかで $j(F_0)\ne0,1728$、かつ $c^3\ne1$ となるように選ぶ。
$R=\C[[d]]$、$K=\C((d))$ とし、
式\eqref{mbpl:eq:hesse-relations}に
\[
 (a,b,c)=(1+d,1,c)
\]
を代入する。source の中心元を $\Omega_H$ と固定し、
target の方程式を自然な比較が与える
$\vartheta(\Omega_H)=3dF_H$ とする。
一般ファイバー上の Picard 補正型について
\begin{equation}\label{mbpd:eq:boundary-phase}
 \delta_E(p)=0,\qquad
 \gamma=[d]\in K^\times/K^{\times3}.
\end{equation}
特に離散付値による Kummer 剰余は $1\in\mathbb Z/3$ である。
\end{proposition}
\begin{proof}
まず、ここで必要な因子模型が閉ファイバーにも延びることを確認する。
Hesse の点曲線は $R$ 上滑らかであり、
$o=(1:-1:0)$ は全族の flex 切断である。
$d=0$ では二次関係が左右交換で不変なので、
点グラフの自己同型 $\sigma$ は $\sigma^2=1$ を満たす。
もし $\sigma=1$ なら二次関係は $\wedge^2W$ になるが、
実際の関係は $cx^2+yz+zy$ とその巡回式であり、その可能性はない。
従って閉ファイバーの $p$ は非零二捩れである。
特に $3p\ne0$ であり、$p$ と $3p$ の零切断はこの形式近傍に現れない。

滑らかな楕円データの Sklyanin 代数が AS 正則であること、および
中心正則三次による商が THCR となることには、
二捩れの除外は要らない
\cite[Theorems~6.6, 6.8]{GRSklyaninATV}
\cite[\S1, pp.~19--21]{GRSklyaninST}。
閉ファイバーでは
\[
 \Omega_H=(1-c^3)(U-V_0),
 \qquad
 (S_3)^{H_{27}}
 =\langle C,U,V_0\rangle/\langle cC+U+V_0\rangle
\]
であるから、$\Omega_H$ は非零である。
この計算は補題\ref{mbpl:lem:hesse-center}の不変部分の計算であり、
$a-b$ を可逆にする必要がない。
Sklyanin 代数は domain なのでこの中心元は正則である。
また $3p\ne0$ により THCR の三次中心は零なので、
$\Omega_H$ は THCR への三次 kernel の生成元となる。

相対的にも、各次数の Hilbert 次元が一般・特殊ファイバーで等しいため、
DVR 上の有限生成次数成分には torsion がなく、有限自由となる。
三段 Koszul 分解は各次数で有限自由加群の複体であり、
ファイバーでの完全性と Nakayama により相対的にも完全である。
同じ議論で $S_R/(\Omega_H)$ の次数成分と THCR の次数成分を同定する。
従って、系\ref{mbrn:cor:factor-boundary}の一因子の相対
parameter-zero window と有限 MF 模型をこの正則基底上で適用できる。
その構成は実際の最高 Koszul syzygy の直線を保持して行い、
二捩れ境界で退化し得る $\det W$ との交代化による同定は使わない。
source の $\Omega_H$ と auxiliary target の $F_H$ はともに正則な生成元なので、
両者による指定周期は $R$ 上の同型として定義される。
以上のデータと構成は、完備化前の正則局所環
$R_0=\C[d]_{(d)}$ 上でも定義される。

flex を原点とし、$M=\OO_E(t_p(o)-o)$ を取る。
これは $R$ 上の実際の線束なので Picard 障害は零である。
$p\ne0$ は閉ファイバーでも成り立つため、
補題\ref{mbpd:lem:hodge-line}の比較直線 $\lambda$ も全 $R$ 上で定義される。
先ほどの正則な二つの周期を使った三乗写像を
$q':\lambda^{\otimes3}\simeq\OO_R$ と書く。
$\lambda$ の $R$-frame で表した $q'$ の値は単元である。

一方、自然な同期の target は $F_H$ から $3dF_H$ へ変わる。
命題\ref{prop:nfm-period-equation-frame-change}の規約により
\[
 q_{\mathrm{nat}}=(3d)^{-1}q',\qquad
 \gamma_{\mathrm{nat}}=[3d]\gamma'
\]
となる。$R^\times$ の全ての元は $R$ 内で三乗根を持つ。
これは剰余体が $\C$ であることと Hensel の補題による。
従って $\gamma'=0$、$[3]=0$ であり、
$\gamma_{\mathrm{nat}}=[d]$ を得る。
付値 $v_d(d)=1$ からこの類は非零である。
\end{proof}

\begin{corollary}[独立な境界位相と最小分裂次数]
\label{mbpd:cor:phase-independence}
任意の $k\geq2$ に対し、定数 $c_1,\ldots,c_k\in\C$ を
前命題の条件および閉ファイバーの $j$ 値の相異なる条件を満たすように選べる。
\[
 R=\C[[d_1,\ldots,d_k]],\qquad K=\operatorname{Frac}(R),
 \qquad (a_i,b_i,c_i)=(1+d_i,1,c_i)
\]
とし、source の関係元を $\sum_i\Omega_{H,i}$、
自然な可換 target を $\sum_i3d_iF_{H,i}$ とする。
一般ファイバーは基本楕円領域の順序付き実表示であり、
\[
 \alpha_{\mathrm{Pic}}=0,\qquad
 \gamma_i=[d_i],\qquad
 [D]_b=([d_b/d_i])_{i\ne b}
\]
となる。
この表示の Picard 補正型において、$G$ 作用を保つ実際の比較核を定義する
有限拡大の最小次数は $3^{k-1}$ であり、
rotation と指定周期も保つ場合の最小次数は $3^k$ である。
それぞれの分裂体として
\begin{equation}\label{mbpd:eq:sharp-phase-fields}
 K_D=K\bigl(\sqrt[3]{d_b/d_i}:i\ne b\bigr),
 \qquad K_T=K\bigl(\sqrt[3]{d_i}:1\leq i\leq k\bigr)
\end{equation}
を取れる。
同じ二つの最小次数は、元の有限型 Hesse 族の函数体
$K_0=\C(d_1,\ldots,d_k)$ 上でも成り立つ。
この場合、分裂体の表示には各 $\gamma_i$ の実際の有理関数代表を用いる。
\end{corollary}
\begin{proof}
閉ファイバーの Hesse パラメータ
$\mu=(2+c^3)/(3c)$ は $c$ の非定数関数であり、
$j$ も非定数である。従って有限個の除外条件を避けて
所要の定数を選べる。
各因子は前命題の一変数族の引戻しであり、
その flex と実線束 $M_i$ も引き戻せるので全 $\delta_i$ は零である。
正則な三乗写像の値は $R^\times$ に属する。
完全局所環 $R$ の単元も Hensel により三乗根を持つので、
前命題と同じ frame 変換で $\gamma_i=[d_i]$ となる。

各高さ一素因子 $(d_j)$ の離散付値を $v_j$ とする。
もし $\prod_i d_i^{e_i}$、$e_i\in\{0,1,2\}$ が $K$ で三乗なら、
全 $j$ に対して $e_j=v_j(\prod_i d_i^{e_i})$ は三で割れる。
従って $[d_1],\ldots,[d_k]$ は $K^\times/K^{\times3}$ 上独立である。
同じ付値で $[d_b/d_i]$、$i\ne b$ の独立性も得る。
Kummer 理論\cite{NFMStacksKummer}により
式\eqref{mbpd:eq:sharp-phase-fields}の次数はそれぞれ
$3^{k-1}$、$3^k$ である。
具体的には、$r$ 個の三乗根を加えた体の Galois 群は $\mu_3^r$ の部分群である。
もし真の部分群なら、その上で恒等となる非自明な character があり、
対応する三乗根の単項式は元の体に属する。
三乗すると、元の $r$ 個の Kummer 類の非自明な関係となり、独立性に反する。

Picard 障害は既に零なので、
系\ref{mbpd:cor:descent-criterion}はこの二つの体上の核の存在を与える。
逆に所要の型の核が有限拡大 $K'/K$ 上にあれば、
$D$ または $T$ が $K'$-点を持つため、
$K'$ は対応する三乗根を全て含む。
$K$ が $\mu_3$ を含むので、これらの体は Galois であり、
所要の最小次数の主張が従う。

最後に $K_0$ 上で、前命題の完備化前の比較直線の frame を各
$\C[d_i]_{(d_i)}$ 上で取る。
regular な周期の三乗写像を $a_i(d_i)\in\C[d_i]_{(d_i)}^\times$ と書けば、
自然な位相は
$\gamma_i=[3d_i/a_i(d_i)]$ である。
従って $d_j=0$ に沿う Kummer 剰余は $\delta_{ij}$ となる。
同じ付値の議論により $\gamma_i$ と $\gamma_b-\gamma_i$ の独立性が従い、
同じ Kummer 拡大次数と最小性を得る。
形式完備化は、その単元因子を三乗として消し、
式\eqref{mbpd:eq:sharp-phase-fields}の単純な表示を得るために使っただけである。
\end{proof}

この系が鋭くするのは、実線束を既に持つ場合の \emph{位相段階} の次数である。
系\ref{mbpd:cor:field-degree}の全上界 $3^{2k-1}$ や
定理\ref{mbrc:thm:finite-point-comparison}の $3^{2k}$ の最小性ではない。
また Picard 補正型以外の比較核の不存在は結論しない。
一方、元の有限型 Hesse 開上にも Picard 補正型の大域比較核は存在しない。
もし存在すれば、その函数体上に必要な $D$ の切断を与えるためである。
これは各因子の Picard gerbe がすでに中立である族での、
自然な関係直線同期に対する実際の位相障害である。

\begin{remark}[全比較核、順序降下、表示 gerbe との区別]
\label{mbpd:rem:scope}
式\eqref{mbpd:eq:alpha}は本節の対角 Picard 補正型についての計算である。
因子の平行移動を graph に持つ核、数値作用の異なる核、因子を混合する核まで
同じ gerbe で分類するとはしていない。
また式\eqref{mbpd:eq:kernel-stacks}では中心 $G$ 作用を保つ構造を標識している。
この標識を忘れた全比較核に対し、異なる $D$ の点が依然として相異なることや、
追加の自己同型がないことは主張しない。

因子の置換は $\delta_i$、$T_i$、$u_i$ を対応する添字へ移し、
和 $\sum_i\delta_i$ と商 $T/\Delta\mu_3$ を保つ。
従って相対模型と作用が置換に coherent に両立する範囲で、
ここで構成した stack と有限被覆も同じ descent datum を持つ。
基準 $b$ は差の座標表示と rotation の線型化を選ぶために用いており、
商 $D$ の定義には現れない。
$3\nmid k$ における対称な中心作用の分裂は、
$\alpha_{\mathrm{Pic}}$ や $\gamma_i$ の消滅を意味しない。
$3\mid k$ における分裂選択の障害も、
本節の scalar gerbe および Kummer 類とは別である。

最後に、本節は順序付きの実表示上で計算している。
全 $\OO$-開へ移す場合には、
命題\ref{mbks:prop:gerbe-synchronization}の表示 gerbe と関係直線の同期、
および実際の相対模型の descent datum を維持する必要がある。
表示 gerbe の差が零であることだけから、
$D$ の切断、$\sum_i\delta_i=0$、または指定周期の位相の消滅は従わない。
\end{remark}

% !TEX root = ../cubic-moduli-descent.tex
\section{Hesse 中心の正則化と一因子周期の全係数}
\label{mbhp:sec:hesse-phase}

本節では、補題\ref{mbpd:lem:hodge-line}の比較直線と、
定理\ref{mbpl:thm:global-line-comparison}の自然な関係直線同期を用いる。
その指定周期の三乗写像を、Jacobian 上の Weierstrass 奇函数によって計算する。
中心元の既知の Hesse 式を使うが、中心直線の flex における正則化、
比較の原点での規格化、および周期 torsor の全係数の決定が本節の内容である。
以下の境界への延長は計算のための一因子模型であり、
全 $\OO$ 認識の領域を境界へ広げる主張ではない。

\subsection{固定した Hesse 曲線上の中心直線}

$K\supset\C$、$\mu^3\ne1$ とし、
\[
 E=V(F_\mu),\qquad
 F_\mu=x^3+y^3+z^3-3\mu xyz,\qquad
 o=(1:-1:0),\qquad L=\OO_E(1)=\OO_E(3o)
\]
を固定する。平行移動点のパラメータ曲線を $P=E$ と書く。
$p=(a:b:c)\in P$ に対し、族 $\mathcal S$ の関係は
\[
 cx^2+ayz+bzy,\qquad cy^2+azx+bxz,\qquad cz^2+axy+byx
\]
である。関係の共通 scalar を除いたこの表示は $P$ 上の関係部分束を定める。
以下、$x,y,z$ は代数の生成元、$a,b,c$ はパラメータである。
$C=x^3+y^3+z^3$、$U=xyz+yzx+zxy$、
$V_0=xzy+zyx+yxz$ と書く。

\begin{lemma}[九つの flex を含む中心直線の正則化]
\label{mbhp:lem:regularized-center}
中心式\eqref{mbpl:eq:hesse-center}について、
$abc\ne0$ 上の
\begin{equation}\label{mbhp:eq:regularized-center}
 \widetilde\Omega=\frac{\Omega_H}{abc}
\end{equation}
は、$P$ 全体で非零な切断
$\mathcal S_3\otimes\OO_P(1)$ に延長する。
その像は直線部分束
\[
 \mathcal C\simeq\OO_P(-1)=\OO_P(-3o)\longrightarrow\mathcal S_3
\]
であり、全幾何ファイバーで中心正則三次を与える。
その商は $(E,L,t_p)$ の相対 THCR である。
点方程式の固定直線を
$\mathcal J=KF_\mu\otimes\OO_P$ とすると、基本開の自然な比較は
束の射
\begin{equation}\label{mbhp:eq:extended-theta}
 \vartheta:\mathcal C\longrightarrow\mathcal J,\qquad
 \vartheta(\widetilde\Omega)=3(a-b)F_\mu
\end{equation}
へ延長する。その零因子は
$E[2]\setminus\{o\}$ の三つの点の和である。
\end{lemma}
\begin{proof}
各ファイバーは、非 flex では通常の滑らかな楕円 Sklyanin 代数である。
通常の AS-regular 性の例外十二点は、三つの座標頂点と
$a^3=b^3=c^3$ の九点であり、$\mu^3\ne1$ の $E_\mu$ 上にはない。
flex では $p\in E[3]$ なので $t_p^*L\simeq L$ であり、
$t_p$ は $H^0(E,L)$ 上の線型持上げを持つ。
この持上げによる $\operatorname{Sym}H^0(E,L)$ の Zhang twist を取る。
平面三次の次数二イデアルは零なので、その二次関係は
$(E,L,t_p)$ の幾何二次関係そのものである。
従って $a=0$、$b=0$、$c=0$ の全九 flex で、
$\mathcal S_p$ はこの多項式環の Zhang twist に一致する。
特に全ファイバーは三次元 quadratic AS-regular domain であり、
Hilbert 級数は $(1-t)^{-3}$ である
\cite[Theorems~6.6, 6.8]{GRSklyaninATV}
\cite[Section~1]{GRSklyaninST}。
有限関係行列の階数は各次数で一定なので、
滑らかな被約曲線 $P$ 上の $\mathcal S_n$ は局所自由である。
この議論は中心商の幾何を先に仮定していない。

次数三の不変関係は
\[
 cC+aU+bV_0=0
\]
である。$b\ne0$ で $V_0$ を消去し、
$a^3+b^3+c^3=3\mu abc$ を用いると
\begin{equation}\label{mbhp:eq:division-certificate}
 \Omega_H
   =3ac\bigl((a^2-\mu bc)C+(\mu ab-c^2)U\bigr)
       \quad\text{in }\mathcal S_3
\end{equation}
を得る。従って
\begin{equation}\label{mbhp:eq:regular-chart}
 \widetilde\Omega
   =\frac3b\bigl((a^2-\mu bc)C+(\mu ab-c^2)U\bigr).
\end{equation}
$a\ne0$ では $U$ を消去して
\[
 \widetilde\Omega
   =-\frac3a\bigl((b^2-\mu ac)C+(\mu ab-c^2)V_0\bigr)
\]
を得る。$\{a\ne0\}$ と $\{b\ne0\}$ は $P$ を覆うので、
この二式は九つの flex の全てを含む。
これらは稠密開で一致するため、局所自由な $\mathcal S_3$ 上で貼り合わさる。
係数は $(a,b,c)$ に関して斉次次数一であり、
正確には $\mathcal S_3\otimes\OO_P(1)$ の切断を得ている。
ここでの除算は $\mathcal S_3$ 内の恒等式である。
パラメータ torsion を持ち得る最大可換商の次数三で
$abc$ を消去しているのではない。

$c=0$、$b=-\zeta a$、$\zeta^3=1$ では
\[
 \widetilde\Omega=-3a\zeta^{-1}C+3\mu aU.
\]
補題\ref{mbpl:lem:hesse-center}の不変部分の計算はこの点でも適用でき、
唯一の不変関係は $aU+bV_0$ である。
よって $C,U$ はその商で独立であり、上の元は非零である。
非 flex では同補題の非零性の証明を用いる。
$a=0$ では
$\widetilde\Omega=-3\mu cC-3c^2U/b$ であり、
唯一の不変関係は $cC+bV_0$ なので非零である。
$b=0$ では
$\widetilde\Omega=3\mu cC+3c^2V_0/a$、
唯一の不変関係は $cC+aU$ であり、やはり非零である。
既知の中心式
\cite[Section~2, p.~3]{GRSklyaninDLL}
\cite[p.~35, Remark~(1)]{GRSklyaninST}
と $\mathcal S_4$ の局所自由性により、
中心性は稠密開から全体に延長する。
各ファイバーが domain であるため、非零中心元は正則である。

幾何的関係の評価は $\mathcal S$ から相対 THCR への射を与える。
THCR の各次数は局所自由であり、$\widetilde\Omega$ の像は稠密開で零なので
全体で零である。次数三の核は階数一で、
上で得た非零中心元がこれを生成する。
中心正則性による商の Hilbert 級数と THCR の Hilbert 級数は一致する。
評価の次数一生成性から全射であり、従って全次数で同型となる。

$F_H=abcF_\mu$ なので、
式\eqref{mbpl:eq:hesse-comparison}は
式\eqref{mbhp:eq:extended-theta}を与える。
この射は $\OO_P(-1)\to\OO_P$ という斉次性とも整合する。
$o$ を原点とする反転は $(a:b:c)\mapsto(b:a:c)$ である。
その固定点は $o$ と直線 $a=b$ の三交点であり、
標数零の滑らかな楕円曲線の二捩れ点は被約である。
従って主張した零因子を得る。
\end{proof}

この族への相対 window の適用には
系\paperref{main}{mbrn:cor:factor-boundary}を用いる。
ここで $\mathcal C$ は指定した THCR 商の核の直線であり、
中心の次数三全体の次元が境界でも一であるとは主張しない。
特に最上 Koszul syzygy 直線は実際の
$(W\otimes R_2)\cap(R_2\otimes W)$ のまま保持する。
二捩れ点で退化することのある $\det W$ への交代化は使わない。
中心商、有限 Koszul 模型、指定周期はこの補題の全パラメータ上で存在する。

\subsection{比較直線の原点を含む延長}

以下の符号規約を固定する。
対角の余法束の同型は
\begin{equation}\label{mbhp:eq:conormal-sign}
 I_\Delta/I_\Delta^2\longrightarrow\Omega^1_E,\qquad
 \overline{\operatorname{pr}_2^*f-\operatorname{pr}_1^*f}\longmapsto df
\end{equation}
とする。これが補題\ref{mbpd:lem:hodge-line}の同型の符号を定める。

\begin{lemma}[比較直線と同期前の三乗写像]
\label{mbhp:lem:global-comparison-line}
$E_x\times P$ 上の線束
$M=\OO(\{x=p\}-\{x=o\})$ を用い、$F_M=-\otimes M$ とする。
$R_p=T_{\OO_E}(t_p)_*(-\otimes L)$、
$R_0=T_{\OO_E}(-\otimes L)$ に対する比較射の直線 $\mathcal V$ は、
$p=o$ を含めて
\begin{equation}\label{mbhp:eq:global-lines}
 \mathcal V\simeq\lambda\otimes\OO_P(o),\qquad
 \widetilde q:\mathcal V^{\otimes3}
       \xrightarrow{\sim}\mathcal J\otimes\mathcal C^\vee,
 \qquad \lambda=H^0(E,\Omega_E^1)
\end{equation}
を満たす。
$p\ne o$ 上の $\mathcal V\simeq\lambda$ は
補題\ref{mbpd:lem:hodge-line}のものと一致する。
$\vartheta$ で同期し、標準切断による
$\lambda\to\lambda\otimes\OO_P(o)$ を三乗して得る
有理写像 $q:\lambda^{\otimes3}\dashrightarrow\OO_P$ の因子は
\begin{equation}\label{mbhp:eq:period-divisor}
 \operatorname{div}(q)=3[o]
       -\sum_{\tau\in E[2]\setminus\{o\}}[\tau].
\end{equation}
\end{lemma}
\begin{proof}
$E_x\times E_y\times P$ 上で $\Delta=(y=x)$、
$\Gamma=(y=x+p)$ と置く。
式\eqref{mbpd:eq:kernel-ratio}と同じ計算により、核の比は
\[
 N=\OO(\Gamma-\Delta)\otimes
       \operatorname{pr}_{x,p}^*M\otimes
       \operatorname{pr}_{y,p}^*M^{-1}
\]
である。$p\ne o$ では相対 Picard intertwiner の計算を使い、
$p=o$ では二つの核と $M_o=\OO_E$ が一致する。
従って全パラメータで $N$ のファイバーは自明な線束である。
proper cohomology と底変更、または同じ相対 intertwiner の構成により、
$g:E\times E\times P\to P$ に対して
$\mathcal V=g_*N$ は直線束、$g^*\mathcal V\to N$ は同型となる。
ここでは滑らかな被約のパラメータ曲線上で証明しており、
非被約底上の結論を幾何ファイバーの一致だけから出していない。

対角上では二つの $M$ は標準的に相殺する。
$\Gamma$ と $\Delta$ の交わりは正確に $p=o$ なので、
\[
 N|_{\Delta\times P}
     \simeq\Omega_E^1\boxtimes\OO_P(o).
\]
これが $\mathcal V$ の主張を与える。
$p\ne o$ では $\OO_P(o)$ の標準切断で自明化するため、
既存の Hodge 同型とも一致する。

中心商の相対 window は、まだ同期していない二つの指定周期
\[
 R_p^3\simeq[2]\otimes\mathcal C,\qquad
 R_0^3\simeq[2]\otimes\mathcal J
\]
を与える。比較射を三回合成してこれらを適用すると、
$\mathcal V^3\to\mathcal J\otimes\mathcal C^\vee$ を得る。
各ファイバーで比較射と周期はいずれも可逆なので、この写像は同型である。
これは系\paperref{main}{mbrn:cor:factor-boundary}の実際の window 周期であり、
途中で scalar を取り替えたものではない。
$\vartheta$ は $\mathcal J\otimes\mathcal C^\vee$ の切断であるから、
同期後の $q$ は $\widetilde q$ をこの切断で割った写像になる。
$\lambda\to\mathcal V$ の零因子は $[o]$、
$\vartheta$ の零因子は前補題の三つの二捩れ点である。
これから式\eqref{mbhp:eq:period-divisor}を得る。
\end{proof}

\subsection{Weierstrass 奇函数による全係数}

楕円曲線の原点付き Weierstrass 表示と微分を
\[
 y^2+a_1xy+a_3y=x^3+a_2x^2+a_4x+a_6,\qquad
 v_\omega=2y+a_1x+a_3,\qquad \omega=dx/v_\omega
\]
と書く。この表示とその座標変換則には
\cite[Chapter~III, \S1, Table~3.1; \S5]{NFMSilvermanArithmetic}
を用いる。$v_\omega$ は反転で符号を変え、因子は
\[
 \operatorname{div}(v_\omega)
   =\sum_{\tau\in E[2]\setminus\{o\}}[\tau]-3[o]
\]
である。

\begin{theorem}[自然な同期に対する一因子周期の全係数]
\label{mbhp:thm:period-formula}
式\eqref{mbhp:eq:conormal-sign}の規約を用いる。
定理\ref{mbpl:thm:global-line-comparison}の自然な同期に対し、
補題\ref{mbpd:lem:hodge-line}の三乗写像は
\begin{equation}\label{mbhp:eq:period-formula}
 q(\omega^{\otimes3})=-\frac1{v_\omega(p)}.
\end{equation}
これは $p,2p,3p\ne0$ の全基本楕円開上の式であり、
普遍族からの任意の許された底変更に両立する。
より内在的には、
$\mathfrak v(p)=v_\omega(p)\omega^{\otimes3}$
という $\lambda^{\otimes3}$ の可逆切断について
$q(\mathfrak v(p))=-1$ である。
\end{theorem}
\begin{proof}
まず固定した $E_\mu$ 上で証明する。
前補題の因子公式により、
$q(\omega^{\otimes3})v_\omega$ は $p$ に依存しない非零定数である。
その値を原点で求める。
$a=1,b=-1,c=0$ では代数は可換多項式環であり、
式\eqref{mbhp:eq:regular-chart}から
\begin{equation}\label{mbhp:eq:origin-curvatures}
 \widetilde\Omega_o=-3C+3\mu U=-3F_\mu,\qquad
 \vartheta(\widetilde\Omega)_o=6F_\mu
\end{equation}
となる。ここでは可換代数の中で $U=3xyz$ である。
曲率を $f$ から $cf$ へ変えると指定周期は
$\pi_{cf}=c^{-1}\pi_f$ となる
（命題\ref{prop:nfm-period-equation-frame-change}）。
原点では両側の商代数と window の section 切断が同一であり、
$M_o=\OO_E$ の標準自明化により比較射は恒等射 $\eta_o$ である。
従って、同期済みだが $\mathcal V_o$ 上で評価した三乗写像は
\begin{equation}\label{mbhp:eq:origin-cube}
 q_{\mathcal V,o}(\eta_o^{\otimes3})
        =\frac{1/6}{-1/3}=-\frac12.
\end{equation}
同一の quotient/window に二つの方程式 frame を入れた比較なので、
ここに未指定の window scalar はない。

Weierstrass の局所径数を $t=-x/y$ とすれば
\[
 \omega=dt+O(t)\,dt,\qquad
 v_\omega=-2t^{-3}+O(t^{-2}).
\]
$\lambda\to\mathcal V$ の一次項も必要である。
曲線上の原点とは異なる点の局所座標を $u$、出力座標を $z$ とし、
平行移動を
$z=u+A(u)t+O(t^2)$、$\omega=du/A(u)$ と書く。
この近傍では $M$ の標準有理切断を用い、
核の比線束の局所 frame を
\[
 e=\frac{z-u}{z-u-A(u)t+O(t^2)}
\]
と取れる。$t=0$ でこれは恒等射 $\eta_o$ に一致する。
対角上で分子の余法類を式\eqref{mbhp:eq:conormal-sign}により $du$ へ送り、
分母を $-A(u)t+O(t^2)$ へ制限すると、
$\mathcal V\to\lambda$ の有理同型は
$e\mapsto-\omega/t+O(1)$ となる。
従って逆向きの正則射は
\[
 \omega\longmapsto-t\eta_o+O(t^2).
\]
式\eqref{mbhp:eq:origin-cube}と合わせて
$q(\omega^3)=t^3/2+O(t^4)$ を得る。
$v_\omega$ の主部を掛ければ定数は $-1$ となり、
式\eqref{mbhp:eq:period-formula}が従う。

この議論をまず $K=\C(\mu)$ に適用する。
滑らかな普遍 Hesse 被覆上では両辺は正則な線束の射であり、
稠密な一般点での一致は全基本開上の一致を与える。
中心線比較と window は表示の同型に自然である。
また Weierstrass 座標変更
$x=u^2x'+r$、$y=u^3y'+su^2x'+t_0$ に対して
$v_\omega=u^3v_{\omega'}$、$\omega=u^{-1}\omega'$ なので、
$v_\omega\omega^3$ は不変である。
従って普遍 Hesse 被覆の重なりでも式は一致し、
補題\ref{mbpl:lem:hesse-cover}の fpqc 降下により
全基本楕円族の式となる。
最後にこの普遍族上の射の等式を引き戻す。
これにより非被約な底変更にも対応し、
任意底上で幾何点だけを比較する議論を避けている。
\end{proof}

\begin{corollary}[全基本開上の Kummer torsor]
\label{mbhp:cor:global-kummer}
原点を選ばない genus-one 曲線族についても、
Jacobian $J=\Pic^0(E/S)$ とその点 $p$ を用いて
$\mathfrak v(p)\in\Gamma(S,\lambda^{\otimes3})^\times$ が内在的に定まる。
自然な同期に対する因子周期 torsor は
\begin{equation}\label{mbhp:eq:cube-torsor}
 T\simeq\{\xi\in\lambda^\times:
              \xi^{\otimes3}=\mathfrak v(p)\}.
\end{equation}
特に Hodge frame $\omega$ を選んだ底上で
\[
 \gamma=[v_\omega(p)]\in H^1(S,\mu_3).
\]
この記述は因子の同型、共通 grading scalar、
中心生成元とその自然な像の同時変更、および任意底変更に不変である。
\end{corollary}
\begin{proof}
$E$ と $J$ の Hodge 線は標準的に一致する。
Weierstrass 変換則により $\mathfrak v$ は Jacobian 上で定義され、
原点と非零二捩れ点を除けば可逆である。
定理の式より、周期適合条件 $q(\eta^3)=1$ は
$\eta^3=-\mathfrak v(p)$ に等しい。
$\xi=-\eta$ とすれば式\eqref{mbhp:eq:cube-torsor}を得る。
これは $\mu_3$-同変である。
Hodge frame の変更は函数を三乗倍するだけであり、
中心方程式の同時変更は二つの周期を同じ比で変更する。
その他の自然性は普遍族上の構成と前定理の降下による。
\end{proof}

\begin{corollary}[Hesse 座標での有理式]
\label{mbhp:cor:hesse-expression}
$a=1$ の chart 上で
$\omega_H=dc/(3b^2-3\mu c)$ と置く。
これは $E_\mu$ の不変微分であり、
$l=a+b+\mu c$ と書けば
\begin{equation}\label{mbhp:eq:hesse-phase}
 \begin{split}
 v_{\omega_H}(p)&=
       9(1-\mu^3)\frac{a-b}{a+b+\mu c},\\
 q(\omega_H^{\otimes3})&=
       \frac{a+b+\mu c}{9(\mu^3-1)(a-b)},\\
 \gamma&=\left[
       9(1-\mu^3)\frac{a-b}{a+b+\mu c}\right].
 \end{split}
\end{equation}
\end{corollary}
\begin{proof}
直線 $l=0$ は $o$ における接線であり、$E$ 上の因子は $3[o]$ である。
したがって $(a-b)/l$ は $v_{\omega_H}$ と同じ因子を持つ。
原点近傍で Hesse 方程式を解くと
\[
 b=-1-\mu c+\frac{\mu^3-1}{3}c^3+O(c^4),\qquad
 l=\frac{\mu^3-1}{3}c^3+O(c^4),\qquad
 \omega_H=(1/3+O(c))\,dc.
\]
$t=c/3+O(c^2)$ に対する主部
$v_{\omega_H}=-54c^{-3}+O(c^{-2})$ から係数が決まる。

別の直接検算として
\[
 X=\frac{3(\mu^3-1)c}{l},\qquad
 V=\frac{9(1-\mu^3)(a-b)}{l}
\]
を用いれば、Hesse 方程式の剰余として
\[
 \frac{dX}{\omega_H}=V,\qquad
 V^2=4X^3-27\mu^2X^2+
       54\mu(\mu^3-1)X-27(\mu^3-1)^2
\]
が成立する。従って $Y=V/2$ は
$\omega_H=dX/(2Y)$ の Weierstrass 表示を与える。
残りの式は前定理と前系による。
\end{proof}

\begin{corollary}[既存の単元表示と境界剰余]
\label{mbhp:cor:unit-and-residue}
$j\ne0,1728$ 上で $h=c_6/c_4\in\lambda^2$ と置くと、
系\ref{mbpd:cor:unit-phase}の単元と frame は
\[
 u=\frac{h^3}{\mathfrak v(p)^2},\qquad
 e=-\frac{h^2}{\mathfrak v(p)}
\]
である。Hodge frame に関する $h=H\omega^2$、
$\mathfrak v=V\omega^3$ では
$[u]=[H^3/V^2]=[V]=\gamma$ となる。
また良い還元を持つ一径数族で $p$ が非零二捩れ因子と横断的に交われば、
自然な位相の Kummer 剰余は $+1\in\mathbb Z/3$ である。
\end{corollary}
\begin{proof}
$q(\mathfrak v)=-1$ を
$e=(q\otimes\mathrm{id})(h^2)$、$u=q^{\otimes2}(h^3)$ へ代入する。
剰余の主張は $v_\omega$ が非零二捩れ点で単純零を持つことによる。
従って命題\ref{mbpd:prop:boundary-phase}の非零剰余と一致する。
\end{proof}

本節が決定したのは、自然な関係直線同期に対する一因子の
Picard 補正周期である。
これを積に適用する際の scalar gerbe、対角位相商、
および選んだ比較核型の条件は定理\ref{mbpd:thm:coherent-corrections}に従う。
この有理式だけから全核型を同定したり、
表示 gerbe や比較核の大域実現障害が消えると結論したりはしない。

中心の除算恒等式、Hesse--Weierstrass 変換、原点での形式展開、および
関連する period multiplier の例の有限多項式計算は
\texttt{computations/verify\_multiblock\_hesse\_phase.py} で再検算できる。
この計算は係数と恒等式を確認するものであり、
window による圏比較、Hom 直線、または fpqc 降下の機械的証明ではない。

% !TEX root = ../nc-cubic.tex
\section{相対積比較と Picard 型核の全構造対象降下}
\label{mbic:sec:intrinsic-comparison}

第\ref{mbrn:sec:relative-nc}節の相対積関手を、実表示を選ばない基底へ移す。
中心同変残余圏と曲線積の比較核は降下する。
通常の残余圏どうしについては、Picard 補正型の核の
\emph{stack} が降下し、その切断の存在はなお Picard 障害と周期位相に依存する。
この二つの結論を区別する。

重み零の複体の gerbe 降下と tensor--Hom の重み則は
\cite[Propositions~4.10, 5.7; Corollaries~4.12, 5.9]{NFMBerghSchnurer}、
rigidification と群作用の両立は
\cite[Theorem~5.1]{MBFMRomagny}の既知の一般論である。
以下では、既存の実関手がその降下条件を満たすことを次数ごとに確認する。
任意の圏の変形や、全ての残余比較核の型の分類は主張しない。

\subsection{三つの表示基底と関手の自然性}

まずブロックの順序を固定し、基本楕円開上で
\[
 \widetilde{\mathfrak P}
   =[\mathcal X_k/\GL_3^k],\qquad
 \mathfrak P=[\mathcal X_k/(\GL_3^k/\mathbb G_m^{\mathrm{diag}})],
 \qquad
 \mathfrak B=[\mathcal U_k/\PGL_3^k]
\]
と置く。射
\[
 a:\widetilde{\mathfrak P}\longrightarrow\mathfrak P,\qquad
 b:\mathfrak P\longrightarrow\mathfrak B
\]
はそれぞれ共通次数スカラーの $\mathbb G_m$-gerbe と、
$G=Z_0=\mu_3^k/\Delta\mu_3$-gerbe である
（定理\ref{mbsd2:thm:gerbe-period}、命題\ref{mbfm:prop:rigidification}）。
各 $\mathfrak B$ 上の曲線を $E_i$、Jacobian の切断を $p_i$、
積を $Y=\prod_iE_i$ とする。
ここで $L_i=\OO_{E_i}(1)$ は $\widetilde{\mathfrak P}$ 上では実線束であるが、
$\mathfrak B$ 上では一般に偏極類である。

非可換側と自然な点三次側の compact 残余圏をそれぞれ
$\mathscr K_p,\mathscr K_0$ と書く。
両者は定理\ref{mbsd2:thm:projector-descent}により $\mathfrak P$ 上にある。
定理\ref{mbpl:thm:commutative-family}の band を保つ表示 gerbe の同値と
$\ell\simeq\ell_{\mathrm{pt}}$ を用い、以後はこの\emph{自然な同期}を固定する。
二つの表示 gerbe を独立に自明化したり、周期を独立に scalar 倍したりしない。

\begin{lemma}[実 window と外積の構造移送]
\label{mbic:lem:naturality}
定理\ref{mbrn:thm:relative-product}および
定理\ref{mbrp:thm:relative-product}の積関手は、順序付き実表示の
次数付き同型と関係直線の同型に、coherence を含めて自然である。
共通次数スカラーに対する関手の重みは零である。
これは核、逆核、単位、余単位、および指定された三次周期の図式について成立する。
\end{lemma}
\begin{proof}
一因子の THCR の切断関手は、局所的に
\[
 F\longmapsto
 \bigoplus_{n\geq0}R(f_i)_*(F\otimes L_{i,n})
\]
で書ける。ここで $L_{i,n}$ は $\sigma_i=t_{p_i}$ による積線束である。
因子の次数スカラー $\lambda_i$ は $L_{i,n}$ に $\lambda_i^n$ として作用する。
曲線を恒等に保つこのスカラーに対し、幾何側の $F$ には恒等の移送を用いると、
切断の次数 $n$ に生じる比較は
\[
 m\longmapsto\lambda_i^n m
\]
である。これは命題\ref{mbsd2:prop:scalar-weight}の $\tau_{\lambda_i}$ と一致する。
切断三角で用いた $\kappa_i\leftrightarrow\OO_{E_i}[1]$ についても、
augmentation の次数零の作用は恒等なので、追加の character は入らない。
内部 shift $(j)$ の重み $j$ は、幾何側の $j$ 回の因子 rotation の重みと一致する。

余核、導来切断、次数による切断、特異 dg 商、相対評価 cone は、
いずれも次数付き同型と可換する。
有限 MF の余核比較の逆および window の逆は、Morita $\infty$-圏での
同値の逆として取る。従って逆の選択に新しい scalar cocycle はなく、
同値と単位・余単位の自然性を同時に保てる。
有限 Koszul--Shamash higher homotopy の選択は生成性の証明にのみ使い、
関手の構成には使っていないことがここで重要である。

外積は因子の次数付き構造移送を tensor する操作であり、比較は乗法的である。
共通関係直線の frame の変更は、第\ref{mbrn:sec:relative-nc}節で用いた
twisted MF の一方の矢印および odd--odd 項の係数直線で補正される。
この変更は source と target の同じ関係包含の変更である。
余核二拡張の合成は指定した中心関係の乗法であり、
その周期図式も構造移送に自然となる。
以上を合成すると、共通スカラーを $\tau$ で消した後の積関手の作用は恒等になる。
これは関手の同型類の一致だけでなく、各段階の自然同型の乗法性による主張である。
\end{proof}

\begin{theorem}[射影的表示上の実際の積核]
\label{mbic:thm:projective-product}
$\mathfrak P$ 上に相対 perfect Morita 核による同値
\begin{equation}\label{mbic:eq:projective-product}
 \Perf(Y/\mathfrak P)
   \xrightarrow{\ \sim\ }\mathscr K_p^{hG},
 \qquad
 \Perf(Y/\mathfrak P)
   \xrightarrow{\ \sim\ }\mathscr K_0^{hG}
\end{equation}
がある。両側は $\mathfrak P$ 上 smooth proper である。
実表示 gerbe の中立化は必要ない。
構成は許された任意の noetherian 底変更に両立する。
この最後の主張は構成済み compact/Perf 模型についてのものであり、
非正則底上の全 coherent 圏との同一視ではない。
\end{theorem}
\begin{proof}
滑らかな表示チャートとその重なりに相対積定理を適用する。
補題\ref{mbic:lem:naturality}は $\GL_3^k$ に対する比較を与え、
共通 $\mathbb G_m$ の比較は $\tau$ によって自明になる。
この $\tau$ は積に整合するので、$\GL_3^k/\mathbb G_m^{\mathrm{diag}}$ に対する
実際の降下 datum となる。
命題\ref{mbrd:prop:kernel-stack}の perfect 双加群とその作用写像の降下を用いれば、
核・逆核・単位・余単位が降下する。
ここで一つの大域 dg 代数や大域 generator の存在を仮定しているのではない。
局所の perfect Morita 双加群が相対 Morita 圏の核を与えるという意味である。

局所生成対象を用いて複体として見る場合にも、比較のスカラー重みは零である。
重み零複体の降下と、その tensor・Hom の両立性
\cite[Proposition~5.7, Corollary~5.9]{NFMBerghSchnurer}
は双加群の作用写像・合成にも適用できる。
smoothness、properness、および核の可逆性はこの fpqc 局所表示で検査できる。
任意底変更との両立は、系\ref{mbrn:cor:base-change}、
系\ref{mbrp:cor:base-change}と同じ降下 datum を引き戻して得る。
\end{proof}

\begin{proposition}[中心 rigidification 上の非捩れの積表示]
\label{mbic:prop:rigidified-product}
系\ref{mbfm:cor:equivariant-rigidification}の構成で得る
$\mathfrak B$ 上の圏を $\mathscr D_p=b_*\mathscr K_p$、
$\mathscr D_0=b_*\mathscr K_0$ とする。
本稿の表示族については、実際の perfect Morita 核による同値
\begin{equation}\label{mbic:eq:rigidified-product}
 \mathscr D_p\simeq\Perf(Y/\mathfrak B)\simeq\mathscr D_0
\end{equation}
がある。従ってこの中心同変化の降下には、曲線積に対する追加の捩れは残らない。
通常の残余圏 $\mathscr K_p,\mathscr K_0$ 自身が $\mathfrak B$ へ降りる、
あるいはその間の通常の比較核が大域的に存在する、という主張ではない。
\end{proposition}
\begin{proof}
前定理の核について、相対慣性 $G$ の比較を確認すればよい。
$\Gamma_k$ の character 格子を
\[
 \Lambda_k=\mathbb Z^k/\langle3e_i-3e_j\rangle
\]
とする。$z\in G$ の持上げ $\widetilde z=(\zeta_i)\in\mu_3^k$ に対し、
$\Gamma_k$-次数付き加群の次数 $\alpha$ の部分で
\begin{equation}\label{mbic:eq:central-trivialization}
 \tau_{\widetilde z}(m_\alpha)
       =\left(\prod_i\zeta_i^{\alpha_i}\right)m_\alpha
\end{equation}
と置く。$\zeta_i^3=1$ によりこれは $\alpha$ の表示によらず、
右作用を $\widetilde z^{-1}$ で移した加群から元の加群への同型である。
共通曲率 character も $\widetilde z$ 上で一なので、MF の両微分と可換する。

命題\ref{mbza:prop:gamma-action}の同定では、
式\eqref{mbic:eq:central-trivialization}は同変対象の構造写像
$\theta_z:\alpha_zM\to M$ そのものである。
一方、各因子の切断に対する $\tau_{\zeta_i}$ を外積すると、
正確に同じ式となる。幾何側で $\widetilde z$ は各 $E_i$ を恒等に保つので、
これが積関手と幾何側の恒等作用の比較である。
持上げを対角 $\zeta$ 倍すると比較の差は、全内部次数 $n$ 上の $\zeta^n$ となる。
これは既に除いた共通スカラーの $\tau_\zeta$ であり、新しい選択を残さない。
積 $\widetilde z\widetilde z'$ との整合性も上式から直接従う。

gerbe $b$ を局所的に中立化すると、
$b_*\mathscr K$ の切断は $\mathscr K^{hG}$ である。
中立化の比較の持上げが $z$ だけ変わる際の同定は $\theta_z$ であり、
上で積関手の比較と一致することを確認した。
従って前定理の核は $\mathfrak B$ 上の式\eqref{mbic:eq:rigidified-product}へ降下する。
これは一般の gerbe 上の圏の降下が常に非捩れになるという主張ではなく、
この実関手に対する慣性の計算である。
\end{proof}

\subsection{偏極を実現しない底上の Picard 障害と周期}

\begin{proposition}[因子モジュライ上の Picard gerbe]
\label{mbic:prop:intrinsic-picard}
$\mathfrak B$ 上に
\[
 \mathscr G_i=\mathcal{Pic}^{p_i}_{E_i/\mathfrak B},
 \qquad
 \mathscr G=\mathcal{Pic}^{m}_{Y/\mathfrak B},
 \qquad m=\sum_i\operatorname{pr}_i^*p_i
\]
が定まり、
\begin{equation}\label{mbic:eq:intrinsic-alpha}
 \alpha_{\mathfrak B}=[\mathscr G]
      =\sum_i\delta_{E_i/\mathfrak B}(p_i),\qquad
 3\delta_{E_i/\mathfrak B}(p_i)=0
\end{equation}
となる。実表示上の $\alpha_{\mathrm{Pic}}$ はこの類の引戻しである。
この主張には、$\mathfrak B$ 上で $L_i$ 自体を実線束として選ぶ必要はない。
\end{proposition}
\begin{proof}
曲線族と Jacobian の切断 $p_i$ は射影的な因子データであり、
命題\ref{mbfm:prop:product-family}の構成で $\mathfrak B$ 上にある。
Picard stack の構成と gerbe 性は fppf 局所であり
\cite[Tag~0DME]{NFMStacksPicardGerbe}、
外積の band は $\mathbb G_m^k\to\mathbb G_m$ の積写像である。
従って命題\ref{mbpd:prop:product-gerbe}と同じ cocycle の積で和の式を得る。

三捩れ性には偏極の実現に注意する。
実表示を選んだチャートでは
\[
 N_i=(\sigma_i)_*L_i\otimes L_i^{-1}
\]
は相対 Picard 類 $3p_i$ を持つ。実際、
$M_i=\OO(t_{p_i}(o)-o)$ を $p_i$ とする右加群の規約では、
$t_{-p_i}^*L_i\otimes L_i^{-1}$ の類が $3p_i$ である。
表示の lift を scalar 倍しても二つの $L_i$ の重みは相殺し、
$N_i$ は $\PGL_3$ に対して自然である。
重なりで $L_i$ が $L_i\otimes f_i^*Q_i$ に変わっても $Q_i$ は相殺する。
従って $N_i$ は $\mathfrak B$ 上の実線束へ降下し、
$\delta(3p_i)=3\delta(p_i)=0$ を与える。
偏極の Brauer 障害とこの Picard 障害を同一視してはいない。
\end{proof}

\begin{lemma}[周期データの射影的自然性]
\label{mbic:lem:phase-descent}
$\lambda_i=(f_i)_*\Omega^1_{E_i/\mathfrak B}$ とすると、
第\ref{mbpd:sec:picard-descent}節の比較直線、三乗写像、位相 torsor は
$\mathfrak B$ 上のデータ
\[
 V_i=\lambda_i,\qquad q_i:\lambda_i^{\otimes3}\xrightarrow{\sim}\OO_{\mathfrak B},
 \qquad T_i,\qquad \gamma_i=[T_i]
\]
の引戻しである。これらは中心係数 $a_i$ の選択によらない。
幾何点の射影 Heisenberg 慣性 $E_i[3]$ は $\lambda_i,q_i,T_i$ に自明に作用する。
\end{lemma}
\begin{proof}
一因子の実表示 gerbe 上で中心直線 $\mathcal C_i$ と点三次直線 $\mathcal J_i$ を用い、
自然な $\vartheta_i:\mathcal C_i\simeq\mathcal J_i$ で二つの周期を同期する。
係数 $a_i$ を持つ多因子表示は、これを
$\ell\xrightarrow{\sim}\mathcal C_i$ で引き戻したものである。
両側の周期に同じ直線同型が入るため、
三重比較の比 $q_i$ ではその同型が相殺する。

$L_i$ と因子 rotation は presentation scalar に対して重み一、
中心線と点三次線は重み三である。
一方、幾何曲線上の Picard 線束 $M_i$ と Hodge 線束は重み零である。
従って $F_iR_{i,p}$ と $R_{i,0}F_i$ の重みは等しく、
その Hom、自然な $\vartheta_i$ による周期の比、および $T_i$ は重み零となる。
普遍 $M_i$ の \emph{Picard gerbe 自身の} scalar は別の重みであるが、
これも Hom の両側で相殺する。
この二つの相殺を混同しない。

補題\ref{mbpd:lem:hodge-line}の graph--diagonal の構成は、
曲線、$p_i$、および線束の pullback に自然である。
従ってその $V_i\simeq\lambda_i$ と $q_i$ は射影的な同型に対して貼り合う。
定理\ref{mbpl:thm:global-line-comparison}の $\vartheta_i$ も
全次数付き同型に自然なので、period の追加 character は入らない。
さらに定理\ref{mbhp:thm:period-formula}と系\ref{mbhp:cor:global-kummer}は、
これを Jacobian の Hodge 線に値を持つ奇 Weierstrass 座標 $v_i(p_i)$ による
\[
 q_i=-v_i(p_i)^{-1},\qquad
 T_i\simeq\operatorname{Cube}(v_i(p_i))
\]
として内在的に記述する。

幾何点の残る射影慣性は $E_i[3]$ の平行移動である。
これは Jacobian、$p_i$、および不変微分を恒等に保つため、
上の三つのデータに恒等に作用する。
しかし $t_q^*M_i\simeq M_i$ という同型を全 $q\in E_i[3]$ に対して
coherent に選べるとは結論しない。
そのような線束の実現の不足は、前命題の Picard gerbe に残る。
\end{proof}

\subsection{通常の比較核の stack と全構造対象開}

$T=\prod_iT_i$、$D=T/\Delta\mu_3$ と置く。
これは $\mathfrak B$ 上の $\mu_3^k$-torsor と $G$-torsor である。
関係直線を保持した rotation 比較というときは、
共通の表示 gerbe 上の二つの重み一 rotation を比較するデータを意味する。
個々の rotation を $\mathfrak P$ 上の実関手として降ろすことは要求しない。

\begin{theorem}[Picard 型比較 stack の内在化]
\label{mbic:thm:kernel-stack}
$\mathfrak P$ 上で、式\eqref{mbic:eq:projective-product}による積表示の下で
対角 Picard 補正となる、$G$-同変な通常残余比較核の stack を
$\mathscr K_{\mathrm{Pic}}$ とする。
順序付き基準因子を一つ選び、相対 rotation 比較と指定周期も加えた stack を
$\mathscr K_{\mathrm{Pic}}^{\rho,\pi}$ とする。
このとき
\begin{equation}\label{mbic:eq:kernel-stacks}
 \mathscr K_{\mathrm{Pic}}
     \simeq b^*(\mathscr G\times_{\mathfrak B}D),\qquad
 \mathscr K_{\mathrm{Pic}}^{\rho,\pi}
     \simeq b^*(\mathscr G\times_{\mathfrak B}T).
\end{equation}
従ってこれらの核の stack は因子モジュライ上の stack から来る。
全対象の自己同型群は、周期標識後も $\mathbb G_m$ である。
この同定は許された任意の底変更と両立する。
\end{theorem}
\begin{proof}
実表示上では定理\ref{mbpd:thm:coherent-corrections}がこの同定を与える。
積関手は定理\ref{mbic:thm:projective-product}で降下済みなので、
Picard 補正型という条件は $\mathfrak P$ 上に定義できる。
共通 scalar は Picard 核に重み零で作用する。
rotation 比較の source と target はともに重み一であり、その Hom と周期の比は零である。
命題\ref{mbks:prop:gerbe-synchronization}の意味で、
比較データは個々の rotation の実現を待たずに降下する。

$\mathfrak B$ まで降ろす際には、各因子の偏極の lift が独立に変わる。
重なりで $L_i'=L_i\otimes f_i^*N_i$ と書けば、
\[
 R_{i,p}'=R_{i,p}\otimes N_i,\qquad
 R_{i,0}'=R_{i,0}\otimes N_i
\]
であり、両方の三次周期の係数直線は $N_i^3$ だけ変わる。
したがって因子比較射とその周期条件は変わらない。
基準 $b$ を用いた双対作用の生成元 $R_i^{-1}R_b$ は
$N_i^{-1}N_b$ だけ変わり得るが、この係数も source と target に共通である。
両作用の intertwining の合成・交換・三乗の図式では、すべて同時に相殺する。
三重重なりの scalar も同じ理由で相殺し、実際の cocycle 条件を満たす。

Picard 核の因子分解を変える操作については
定理\ref{mbpd:thm:coherent-corrections}の band の積と
負の自己 Ext 消滅をそのまま用いる。
coherent な双対作用の比較の選択は $D$、
さらに相対 rotation と周期の比較を選ぶことは $T$ に一致する。
これらは前補題によって $\mathfrak B$ 上にあり、
普遍核の scalar は比較射の共役で相殺する。
従って式\eqref{mbic:eq:kernel-stacks}の同定自体が降下する。
有限双対同変化の平均化と Fourier 射影は構造移送に自然なので、
積上の記述から通常の残余比較核へ戻す際にも同じ降下 datum を保つ。
\end{proof}

\begin{corollary}[全構造対象開上の比較と二つの障害]
\label{mbic:cor:all-o}
系\ref{mbrr:cor:stack-isomorphism}の開
$\mathfrak M_k^\circ\simeq\mathfrak P_k^\circ$ 上で、
中心同変残余圏と点曲線積の perfect 圏を比較する実際の核が存在する。
Picard 型の通常比較核の stack もこの開へ移る。

この開への任意の許された射 $S\to\mathfrak M_k^\circ$ に対し、
その型の $G$-同変な通常比較核が存在する必要十分条件は
\[
 \alpha_S=0,\qquad D_S(S)\ne\varnothing,
 \qquad
 \alpha_S=\sum_i\delta_{E_i/S}(p_i)
\]
である。相対 rotation と指定周期も保持する条件は
$\alpha_S=0$、$T_S(S)\ne\varnothing$ である。
ここでは順序が局所的な場合、和と torsor は次の置換降下で解釈する。
\end{corollary}
\begin{proof}
相対認識はこの開上で全 $\OO$ 同値と射影的な次数付き構造移送を同定する。
従って前の核と stack の降下は全 $\OO$ 射に対して自然である。
判定は式\eqref{mbic:eq:kernel-stacks}の二因子がそれぞれ大域対象を持つ条件である。
類は $\mathfrak B$ の類をまず $\mathfrak P^\circ$ へ、次いで $S$ へ引き戻す。
$H^2(\mathfrak B,\mathbb G_m)\to H^2(S,\mathbb G_m)$ の単射性は仮定せず、
因子モジュライ上の gerbe が中立であることを必要条件にしているのではない。
\end{proof}

\subsection{順序を忘れる場合と残る標識}

\begin{proposition}[因子置換と相対 rotation 比較]
\label{mbic:prop:permutation}
上の積核、Picard gerbe、および核 stack の記述は
$\mathfrak S_k$ に coherent に同変である。
順序を忘れた底では $\mu_3^k$ と $G$ を因子置換で降下した有限群とし、
\[
 \mathscr G\times D,\qquad \mathscr G\times T
\]
もその意味で降下する。
第二の stack は、両側の線型化を同時に変更することを許した
\emph{相対的な} rotation・周期比較を分類する。
一つの標準 $G$-線型化付き rotation を順序なしに選ぶ主張ではない。
\end{proposition}
\begin{proof}
曲線積、$m=\sum_i\operatorname{pr}_i^*p_i$、外積 Picard gerbe は因子の置換に自然である。
各 $T_i$ は対応する曲線・$p_i$・Hodge データとともに置換される。
従って $T$ の置換作用は対角 $\mu_3$ を保ち、$D$ にも作用する。
核の外積と finite MF の外積は Koszul 符号を含む標準 braiding を用いる。
その braid 関係と二乗関係が三重重なりの coherence を与える。
共通 charge と全有限指標を保持する window もこの作用に自然である。
ここで「符号を無視できる」とはしていない。

rotation の比較では基準 $b$ の変更を明示する。
定理\ref{mbpd:thm:coherent-corrections}の
$D\times T_b\simeq T$ は、すべての因子の周期適合比較射を同時に与える記述である。
したがって別の基準 $c$ に移ると、その同じ組の $c$ 成分を取ればよい。
関手としては、命題\ref{mbza:prop:shift-character}の同じ character twist が
source と target の両方に現れる。
$G$-同変比較核の同変化は双対指標 twist を intertwine するため、
この共通 twist を用いて二つの相対比較を同定できる。
三つの基準の比較は既存の character cocycle と核の同変 coherence により一致する。
period 図式も同じ三乗関係を用いるので保たれる。
これにより $\mathscr G\times T$ は基準選択に依存せず置換で降下する。

この議論は source と target の線型化を\emph{同時に}変えている。
個々の rotation の線型化を一つ選ぶ torsor
$\{\boldsymbol a\in\mathbf F_3^k:\sum a_i=1\}$ は別に残る。
系\ref{mbza:cor:permutation-obstruction}により、その置換不変な切断は
$3\nmid k$ のときに限り存在する。
共通表示 $\mathbb G_m$-gerbe の中立化も、この相対比較からは従わない。
\end{proof}

\begin{remark}[核 stack と選択した核の降下の境界]
\label{mbic:rem:scope}
式\eqref{mbic:eq:rigidified-product}は中心同変化についての大域核を与えるが、
双対群の作用を忘れた通常残余圏の比較核は別である。
また $\mathfrak B$ 上では $L_i$ の独立な twist が残るため、
各 $R_i$ や $R_i^{-1}R_b$ 自体を非捩れの大域関手として降ろしたとはしていない。
降りるのは、比較の両側でこれらの twist を同期した核の stack である。

この stack が $\mathscr G\times D$ と書けても、
$\mathscr G$ や $D$ の大域切断は自動ではない。
さらに、別の残余自己同値と合成して Picard 型の外へ出る核の存在・非存在を、
この型の scalar 類や位相から結論することはできない。
指定された形式核の全 jet の代数化も、本節の降下とは別問題である。
\end{remark}

% !TEX root = ../cubic-moduli-descent.tex
\section{点曲線の積における全中心同変核と混合型}
\label{mbfc:sec:full-comparison}

本節では、三ブロック以上、または二つの Jacobian が幾何学的に非同種である場合に、
第\ref{mbpd:sec:picard-descent}節の Picard 補正型を全ての中心同変比較核へ拡張する。
三ブロック以上では、指定した双対中心作用が因子を分離するため、同種性を除外しない。
二つの因子が同種の場合には、実際に外積でない混合核が生じる。
一般の体では、偏極曲線自己同型による追加の Kummer 補正を許す必要がある。
一方、固定 Hesse level を持つ族ではこの追加補正が消えるため、
系\ref{mbpd:cor:phase-independence}の最小次数は全ての中心同変比較核について成り立つ。

以下では $K\supset\C$ を体、$k\ge2$ とし、
第\paperref{main}{mbrc:sec:rotation-comparison}節の順序付き実表示と自然な関係直線同期を固定する。
$E_i/K$ は滑らかな genus-one 点曲線、$L_i$ は実次数三線束、
$p_i\in J_{E_i}(K)$ は点の平行移動であり、
\[
 Y=\prod_{i=1}^k E_i,\qquad
 R_{i,p}=T_{\OO_{E_i}}(t_{p_i})_*(-\otimes L_i),\qquad
 R_{i,0}=T_{\OO_{E_i}}(-\otimes L_i)
\]
と書く。関係直線は体上で一度自明化し、指定周期を
$\pi_{i,p}:R_{i,p}^3\simeq[2]$、
$\pi_{i,0}:R_{i,0}^3\simeq[2]$ と表示する。
この自明化を全因子で同時に変更しても以下の障害類は変わらない。

基準因子 $b$ を固定し、
\[
 G=Z_0=(\mu_3^k)/\Delta\mu_3,\qquad H=\widehat G,
 \qquad Q_{i,p}=R_{i,p}^{-1}R_{b,p},\quad
 Q_{i,0}=R_{i,0}^{-1}R_{b,0}\quad(i\ne b)
\]
とする。積上の関手には他の因子の恒等関手を外積するものとする。
$H$ の作用は指定周期を含む coherent な作用であり、
その同定を任意の $\operatorname{Aut}(H)$ で変更することはしない。
本節の全核という語は実際の perfect Fourier--Mukai 核を指す。
有限群の二重同変化により、これは元の残余圏の
$G$-同変 dg Morita 比較核の問題に等しい
\cite[Proposition~4.1, Theorem~4.2]{MEDElagin}。
体上の適用と rotation の同変性は
定理\paperref{main}{mbrc:thm:finite-point-comparison}の証明で確認済みである。

\subsection{外積分解と shift の曖昧さ}

\begin{lemma}[Orlov 像が block 対角である核の外積分解]
\label{mbfc:lem:factorization}
$\kappa$ を標数零の代数閉体とし、$E_1,\ldots,E_k$ を楕円曲線とする。
$\operatorname{FM}(E_i)$ と $\operatorname{FM}(Y)$ を
Fourier--Mukai 自己同値の同型類の群とし、
$\operatorname{FM}_{\mathrm{block}}(Y)$ を Orlov 像が
$Y\times\widehat Y=\prod_i(E_i\times\widehat E_i)$ に関して
block 対角である部分群とする。このとき外積は同型
\begin{equation}\label{mbfc:eq:factorization}
 \operatorname{FM}_{\mathrm{block}}(Y)
 \simeq
 \left(\prod_i\operatorname{FM}(E_i)\right)
 \big/
 \left\{([n_1],\ldots,[n_k]):\sum_i n_i=0\right\}.
\end{equation}
を与える。外積が恒等関手となる因子関手は、それぞれ shift である。
因子が互いに非同種なら $\operatorname{FM}_{\mathrm{block}}(Y)=\operatorname{FM}(Y)$ である。
\end{lemma}
\begin{proof}
Orlov の写像と完全列
\cite[Propositions~2.15, 3.3, Theorem~4.14]{MBFCOrlovAbelian}を用いる。
原点を選んで $Y$ をアーベル多様体とみなすと、
\[
 0\longrightarrow \mathbf Z\oplus(Y\times\widehat Y)(\kappa)
 \longrightarrow\operatorname{FM}(Y)
 \longrightarrow U(Y\times\widehat Y)\longrightarrow1
\]
がある。
isometric 自己同型の block 対角部分群は
\[
 U(Y\times\widehat Y)_{\mathrm{block}}=\prod_i U(E_i\times\widehat E_i)
\]
である。各成分を因子の Fourier--Mukai 同値へ持ち上げ、その外積を取る。
与えられた同値との差は上の完全列の核に入り、
平行移動・次数零線束の tensor・shift である。
平行移動と $\Pic^0(Y)=\prod_i\Pic^0(E_i)$ の線束も外積であるから全射性を得る。
非同種なら異なる因子とその双対の間の群準同型は零なので、全 Orlov 像が block 対角である。

核を確認する。因子同値の核 $P_i$ の外積が $\OO_{\Delta_Y}$ と同型なら、
Orlov の層性\cite[Proposition~3.2]{MBFCOrlovAbelian}により
$P_i=F_i[n_i]$ と書ける。
支持の外積が $\Delta_Y=\prod_i\Delta_{E_i}$ なので、各 $F_i$ の支持は
$\Delta_{E_i}$ である。対角のイデアルによる積も外積後に零となり、
非零な他因子との体上の tensor は忠実なので、各 $F_i$ は対角に scheme-theoretically 支持される。
外積が対角上の線束であることを、他の因子の点へ制限すると、
各 $F_i$ 自体が対角上の線束となる。
従って $P_i=\Delta_{E_i*}M_i[n_i]$ である。
同じ制限により各 $M_i$ は自明であり、cohomology の次数から
$\sum_i n_i=0$ を得る。逆の包含は外積の shift 則による。
\end{proof}

\subsection{指定された中心作用による三因子以上の分離}

ここでは代数閉体 $\kappa$、$\operatorname{char}\kappa=0$ 上で、
楕円曲線 $E_i$ と次数三線束 $L_i$ を考える。
因子間の同種性、CM の有無、または $j$ の相異なることは仮定しない。
原点 $o_i$ を選び、主偏極
\[
 \varpi_i:E_i\xrightarrow{\sim}\widehat E_i,\qquad
 x\longmapsto t_x^*\OO_{E_i}(o_i)\otimes\OO_{E_i}(o_i)^{-1}
\]
によって $V_i=E_i\times\widehat E_i$ を $E_i^2$ と表示する。
Orlov の写像を $\mathfrak u$ と書く。
その規約は
\[
 \Phi_{(x,\xi)}=t_{x*}(-)\otimes P_\xi,\qquad
 F\Phi_vF^{-1}\simeq\Phi_{\mathfrak u(F)v}
\]
である \cite[equation~(2.1), Corollary~2.13,
Propositions~2.15, 3.3]{MBFCOrlovAbelian}。

\begin{lemma}[因子 rotation の Orlov 行列]
\label{mbfc:lem:orlov-rotation}
上の $V_i=E_i^2$ の座標において、
$R_i=T_{\OO_{E_i}}\circ(-\otimes L_i)$ の Orlov 像は
\[
 r_i=r=
 \begin{pmatrix}-2&-1\\3&1\end{pmatrix},\qquad
 r^2+r+I=0.
\]
従って $r-I$ は isogeny であり、
$(r-I)(r+2I)=-3I$ である。
これは rank--degree 格子上の作用を代用したものではなく、
$E_i\times\widehat E_i$ 上の実際の群準同型の行列である。
\end{lemma}
\begin{proof}
まず $-\otimes L_i$ と $\Phi_{(x,\xi)}$ の共役を取る。
$\varpi_{L_i}=3\varpi_i$ なので、その作用は
$(x,\xi)\mapsto(x,\xi+3\varpi_i(x))$、従って行列は
\[
 t_3=\begin{pmatrix}1&0\\3&1\end{pmatrix}
\]
である。$\Pic^0$ 成分はこの共役行列に影響しない。
主偏極で元の曲線へ戻した normalized Poincare Fourier--Mukai 変換を
$\mathsf S$ とすると、skyscraper と次数零線束の像から
\[
 \mathfrak u(\mathsf S)=
   s=\begin{pmatrix}0&-1\\1&0\end{pmatrix}.
\]
実際、$\mathsf S(k_x)=P_{\varpi_i(x)}$、
$\mathsf S(P_{\varpi_i(y)})=k_{-y}[-1]$ である。
点の spherical twist は $T_{k_{o_i}}=-\otimes\OO(o_i)$、
$\mathsf S(k_{o_i})=\OO_{E_i}$ なので、
\[
 \mathfrak u(T_{\OO_{E_i}})
   =s\begin{pmatrix}1&0\\1&1\end{pmatrix}s^{-1}
   =\begin{pmatrix}1&-1\\0&1\end{pmatrix}.
\]
Orlov 写像の関手性から、この行列と $t_3$ の積が $r$ となる。
残りは整数行列の積で確認できる。
特に shift は Orlov 写像の核に入るので、
rank--degree 格子での shift の作用との同一視はしていない。
\end{proof}

$Y=\prod_{i=1}^kE_i$、$V=Y\times\widehat Y=\prod_iV_i$ とする。
基準 $b$ を選び、$Q_i=R_i^{-1}R_b$ とする。
$H=(\mathbf Z/3)^{k-1}$ の Orlov 像は
\[
 q_{\mathbf a}=\operatorname{diag}(r_1^{a_1},\ldots,r_k^{a_k}),
 \qquad
 \mathbf a\in(\mathbf Z/3)^k,\qquad \sum_i a_i=0
\]
である。

\begin{lemma}[三因子以上の整数 block rigidity]
\label{mbfc:lem:marked-block-rigidity}
$k\ge3$ とする。
$g\in U(Y\times\widehat Y)$ が全ての $q_{\mathbf a}$ と可換するなら、
$g$ は因子 $V_i$ に関して block 対角である。
従って
\[
 \operatorname{Cent}_{U(V)}(\mathfrak u(H))
   =\prod_i\operatorname{Cent}_{U(V_i)}(r_i).
\]
この主張は同種な因子や CM 係数を含む整数 Hom 群上で成立する。
\end{lemma}
\begin{proof}
$g_{ij}:V_j\to V_i$ を block 成分とする。
$i\ne j$ を固定し、これらと異なる $\ell$ を選ぶ。
$a_i=1$、$a_\ell=-1$、他の成分を零とした元について可換性を取れば
\[
 (r_i-I)g_{ij}=0.
\]
前補題の恒等式を左から掛けると $3g_{ij}=0$ となる。
アーベル多様体間の Hom 群は torsion-free なので $g_{ij}=0$ である。
具体的には、$3g_{ij}=0$ なら連結な $V_j$ の像は有限群 $V_i[3]$ に入り、
群準同型の像は零となる。
全ての非対角成分が零になった後、
isometry 条件と可逆性は因子ごとの条件となる。
逆包含は明らかである。
\end{proof}

補題\ref{mbfc:lem:factorization}により、ここで得た block 対角性は
実際の Fourier--Mukai 核の外積分解を与える。
同種因子間の混合線束は Neron--Severi 部分に属し、
Orlov 写像の核である $\Pic^0(Y)$ へ紛れ込むことはない。

ここで用いた Orlov の完全列と楕円曲線上の数値分類は既知の結果である。
本節で必要なのは、それらを指定された双対 rotation 作用へ適用する次の議論である。

\begin{proposition}[双対 rotation 作用の全 underlying centralizer]
\label{mbfc:prop:centralizer}
代数閉体上で、次のいずれかを仮定する：
\begin{equation}\label{mbfc:eq:block-hypothesis}
 k\ge3,\qquad\text{または}\qquad k=2\text{ かつ }E_1,E_2\text{ が非同種}.
\end{equation}
各 $L_i$ を次数三とし、$R_i=R_{i,0}$、
$\varepsilon_i=R_i^{-1}[1]$ と置く。
$\varepsilon_i^3=[1]$、$\varepsilon_i^2=R_i$ である。
全ての $Q_i=R_i^{-1}R_b$ と関手の同型類で可換する自己同値の群は
\begin{equation}\label{mbfc:eq:centralizer}
 \operatorname{Cent}_{\operatorname{FM}(Y)}\{Q_i:i\ne b\}
 \simeq
 \left(\prod_i\operatorname{Aut}(E_i,L_i)\right)\times\Lambda_k,
 \qquad
 \Lambda_k=\mathbf Z^k/
       \{(3n_i)_i:\sum_i n_i=0\}.
\end{equation}
ここで $(a_i)$ は外積 $\boxtimes_i\varepsilon_i^{a_i}$ を表す。
基準因子を選ぶと $\Lambda_k\simeq\mathbf Z\oplus(\mathbf Z/3)^{k-1}$ となる。
これは underlying 関手の同型類の分類であり、coherent な同変構造はまだ標識していない。
\end{proposition}
\begin{proof}
三ブロック以上では補題\ref{mbfc:lem:marked-block-rigidity}、
二つの非同種因子では補題\ref{mbfc:lem:factorization}により、
自己同値を $\boxtimes_i F_i$ と書く。
$FQ_iF^{-1}\simeq Q_i$ を右辺の逆関手と合成すると、
第 $i$、第 $b$ 因子だけを除いて恒等となる外積を得る。
同補題の核の記述から、ある整数 $n_i$ に対し
\[
 F_iR_iF_i^{-1}\simeq R_i[n_i]
\]
となり、基準因子にも同様の式が成り立つ。
三乗を取ると左辺は $[2]$、右辺は $[2+3n_i]$ である。
非零有界 coherent 複体は非零 shift と同型にならないため $n_i=0$ である。
従って全ての $F_i$ が $R_i$ と可換する。

系\ref{cor:nfm-elliptic-full-centralizer}により、各因子で
$F_i=\phi_{i*}\varepsilon_i^{a_i}$、
$\phi_i\in\operatorname{Aut}(E_i,L_i)$ と一意に書ける。
補題\ref{mbfc:lem:factorization}の shift の核は、
$\varepsilon_i^3=[1]$ の下で式\eqref{mbfc:eq:centralizer}の部分格子に一致する。
偏極自己同型は $R_i$ と shift に可換するので直積である。
\end{proof}

同じ議論を source の $R_{i,p}$ と target の $R_{i,0}$ に適用すると、
幾何的な underlying $H$-intertwiner は、各因子の全一因子 intertwiner の外積となる。
周期を三乗した上の議論は、因子間で余分な shift を相殺する可能性も除いている。

\subsection{双対群の抽象自己同型を許した正規化群}

この小節では $H$ の標識を $\operatorname{Aut}(H)$ によって変更することを許す。
まず関手の同型類の群における normalizer を扱い、
coherent な作用構造についての主張は明記する。
代数閉体 $\kappa$、$\operatorname{char}\kappa=0$ 上で
$Y=\prod_{i=1}^kE_i$、$R_i=T_{\OO_{E_i}}(-\otimes L_i)$ とし、
各 $L_i$ の次数は三、指定周期は $R_i^3\simeq[2]$ とする。
\[
 H=\left\{(t_i)\in\mathbf F_3^k:\sum_i t_i=0\right\}
\]
と表示する。整数 lift $\widetilde t_i$ に対し、$H$ の作用関手の同型類は
\begin{equation}\label{mbfc:eq:normalizer-action}
 Q_{\mathbf t}=
 \left(\mathop{\boxtimes}_{i=1}^k R_i^{\widetilde t_i}\right)
 \left[-\frac{2}{3}\sum_i\widetilde t_i\right]
\end{equation}
である。三次周期により lift の変更に依存しない。
特に $Q_{e_i-e_j}=R_iR_j^{-1}$ であり、
coherence は既存の $\Gamma$-同変 MF 模型からのものを用いる。
この作用は忠実である。

\begin{lemma}[三因子以上の指標配置]
\label{mbfc:lem:normalizer-characters}
$k\ge3$ とし、$\chi_i\in H^\vee$ を第 $i$ 座標の制限とする。
$\{\pm\chi_i:1\le i\le k\}$ を保つ $\operatorname{GL}(H^\vee)$ の部分群は
\[
 \{\pm1\}\times S_k
\]
である。すなわち、その任意の元は一つの置換と全因子共通の符号から成る。
$k=3$ にも例外はない。
\end{lemma}
\begin{proof}
$H^\vee=\mathbf F_3^k/\langle(1,\ldots,1)\rangle$ において
$\chi_i$ は標準基底の像である。
$k\ge3$ なら $\chi_i=\pm\chi_j$ は $i=j$ で符号が正の場合に限る。
実際、異なる $i,j$ に対して第三の座標を比較すれば、
$e_i\mp e_j$ が $(1,\ldots,1)$ の倍となることはない。
また、$\chi_i$ の間の関係全体は
\[
 \sum_i\chi_i=0
\]
の一次元空間である。
従って配置を保つ線型同型は $\chi_i\mapsto\epsilon_i\chi_{\sigma(i)}$
と書け、上の関係を移すと全ての $\epsilon_i\in\{\pm1\}$ が等しい。
逆に共通符号と任意の置換はこの関係を保ち、線型同型を定める。
これらの作用は忠実である。
特に $k=3$ では、四点からなる $\mathbf P^1(\mathbf F_3)$ のうち
選ばれた三点の集合を保つ射影作用は $S_3$ であり、
残る共通 scalar が上の $\{\pm1\}$ である。
\end{proof}

\begin{proposition}[全双対作用の Fourier--Mukai normalizer]
\label{mbfc:prop:normalizer}
$k\ge3$ とする。因子の非同種性、CM の有無、$j$ の相異なることは仮定しない。
Fourier--Mukai 自己同値の同型類の群において
\[
 N_H=\left\{F:\ \exists\alpha\in\operatorname{Aut}(H),\quad
 FQ_{\mathbf t}F^{-1}\simeq Q_{\alpha(\mathbf t)}\quad
 (\mathbf t\in H)\right\}
\]
と置く。
このとき $F\in N_H$ が誘導する $\alpha$ は因子置換だけであり、
共通反転は生じない。
許される置換は正確に
\[
 \Pi=\left\{\sigma\in S_k:
       (E_i,L_i)\simeq(E_{\sigma(i)},L_{\sigma(i)})\text{ for all }i\right\}
\]
であり、完全列
\begin{equation}\label{mbfc:eq:normalizer-sequence}
 1\longrightarrow\operatorname{Cent}_{\operatorname{FM}(Y)}(H)
 \longrightarrow N_H\longrightarrow\Pi\longrightarrow1
\end{equation}
を得る。
各 $F$ は置換後に因子 Fourier--Mukai 同値の外積となり、
その因子同値は $R_i$ を対応する $R_{\sigma(i)}$ へ shift なしで移す。

代数閉体上では、$\Pi$ は $j$ の値を保つ置換群でもある。
従って $j_i$ が相異なれば
\begin{equation}\label{mbfc:eq:normalizer-equals-centralizer}
 N_H=\operatorname{Cent}_{\operatorname{FM}(Y)}(H).
\end{equation}
同じ $j$ の因子ごとに偏極同型を選ぶと、
式\eqref{mbfc:eq:normalizer-sequence}は分裂し、
命題\ref{mbfc:prop:centralizer}の群に $\Pi$ が因子置換で作用する半直積となる。
この分裂は選択に依存する。
\end{proposition}
\begin{proof}
原点を選んで $V_i=E_i\times\widehat E_i$、$V=\prod_iV_i$ とする。
補題\ref{mbfc:lem:orlov-rotation}により、Orlov 写像での作用は
\[
 \mathfrak u(Q_{\mathbf t})
       =\operatorname{diag}(r^{t_1},\ldots,r^{t_k}),\qquad
 r=\begin{pmatrix}-2&-1\\3&1\end{pmatrix}.
\]
これは数値 Grothendieck 群の代用ではなく、実際の $V$ 上の作用である。
一つの素数 $\ell\ne3$ を固定し、Tate 加群を
$\mathbf Q_\ell(\zeta_3)$ へ拡大する。
$r$ の固有値は $\zeta_3,\zeta_3^{-1}$ なので、
$H$ の指標の集合は正確に $\{\pm\chi_i\}$ である。
各指標の重複度は二であり、因子の同種性や CM によって変わらない。

Orlov 写像の関手性
\cite[Proposition~2.15]{MBFCOrlovAbelian}から、
$u=\mathfrak u(F)$ は $\alpha$ に沿ってこの表現を移す。
従って前補題より、ある $\sigma\in S_k$ と共通符号
$\epsilon\in\{\pm1\}$ があり、
\[
 (\alpha\mathbf t)_{\sigma(i)}=\epsilon t_i
 \qquad(\mathbf t\in H)
\]
となる。

ここで有理的な指標分解から実際の整数 block 分解へ戻る。
各 $V_i$ の Tate 加群は共役指標対 $\{\chi_i,-\chi_i\}$ の成分である。
従って $u$ の $V_i\to V_j$ 成分は、$j\ne\sigma(i)$ なら
Tate 加群上で零となる。
アーベル多様体の準同型は Tate 加群への作用で忠実に検出されるので、
その成分は実際に零である。
これは Tate 同種定理を用いた全射性の主張ではなく、
非零な像を持つ準同型は Tate 加群上でも非零になるという忠実性だけを用いる。
同じ議論を $u^{-1}$ に適用すると、残る
\[
 u_i:V_i\xrightarrow{\sim}V_{\sigma(i)}
\]
は isogeny だけでなく実際の同型である。
Orlov の isometry 条件
\cite[Definition~2.17, Proposition~2.18]{MBFCOrlovAbelian}は
この block 分解に沿って各 $u_i$ の isometry 条件となる。

標数零の代数閉体上なので、各 $u_i$ は
$F_i:D^b(E_i)\xrightarrow{\sim}D^b(E_{\sigma(i)})$ へ持ち上がる。
具体的には、まず isometry の存在から一つの同値 $F_i^0$ を得て
\cite[Theorem~4.13]{MBFCOrlovAbelian}、その像を $u_i^0$ とする。
次に $(u_i^0)^{-1}u_i\in U(V_i)$ を同曲線の自己同値へ持ち上げ
\cite[Theorem~4.14]{MBFCOrlovAbelian}、$F_i^0$ と合成すれば、
指定された isometry $u_i$ 自身を持つ $F_i$ が得られる。
その外積と置換の合成を $F$ から除くと Orlov 写像の核に入り、
それは shift、平行移動、$\Pic^0(Y)$ の線束の tensor から成る
\cite[Proposition~3.3]{MBFCOrlovAbelian}。
$\Pic^0(Y)=\prod_i\Pic^0(E_i)$ なので、これらも因子の外積である。
従って $F$ 自身を、置換を伴う $F_i$ の外積としてよい。
ここでは異なる因子間の isometry の持上げを代数閉体上でだけ用いており、
任意の定義体上の全射性は仮定していない。

$\mathbf t=e_i-e_j$ の intertwining を取り、
外積が恒等になる核が総 shift 零だけであるという
補題\ref{mbfc:lem:factorization}の記述を適用する。
すると各因子について整数 $n_i$ があり、
\[
 F_iR_iF_i^{-1}\simeq R_{\sigma(i)}^{\epsilon}[n_i]
\]
となる。三乗を取れば
\[
 [2]\simeq[2\epsilon+3n_i]
\]
である。
非零有界 coherent 複体は非零 shift と同型にならない。
$\epsilon=-1$ なら $3n_i=4$ となって不可能であり、
$\epsilon=1$ なら $n_i=0$ である。
従って共通反転は生じず、各 $F_i$ は shift なしの rotation intertwiner となる。

系\paperref{main}{cor:nfm-elliptic-any-intertwiner-criterion}を $p=0$ に適用すると、
$F_i$ の存在から $(E_i,L_i)\simeq(E_{\sigma(i)},L_{\sigma(i)})$ が従う。
逆にこの偏極同型の pushforward の外積は所要の置換を実現する。
従って像は正確に $\Pi$ であり、核は定義から centralizer である。
代数閉体上では同じ $j$ の genus-one 曲線は同型であり、
次数三線束の差は平行移動で吸収できる。
実際、平行移動から $\Pic^3$ への写像の線型部は $[3]$ で、これは全射である。
同型類ごとの代表への偏極同型を固定すれば置換が厳密な群作用となり、
最後の分裂と半直積を得る。
\end{proof}

\begin{remark}[coherent 構造と符号の別検算]
\label{mbfc:rem:normalizer-coherence}
前命題は underlying 関手の同型類を分類する。
従って、coherent な $H$-normalizer に対しても誘導される自己同型の必要条件は同じである。
代数閉体上では、上で構成した各偏極置換の因子比較を三乗根で規格化して
指定周期に合わせ、その外積を取れば coherent な置換を実現できる。
一般の体上では period multiplier や比較核の降下を別に確認する必要がある。

共通反転の排除は CM の有無に依存しない関手の周期から得られた。
数値格子上でも、$Xr=r^{-1}X$ を満たす実行列は
\[
 X=\begin{pmatrix}a&b\\3b-3a&-a\end{pmatrix},\qquad
 \det X=-(a^2-3ab+3b^2)
\]
なので、非零なら行列式は負である。
従って $r$ と $r^{-1}$ は $\operatorname{SL}_2(\mathbf Z)$ では共役でない。
一方、上の証明では因子の Orlov isometry をこの二次元数値行列と同一視せず、
実際の Fourier--Mukai 持上げと周期を用いた。
\end{remark}

\begin{corollary}[二つの非同種因子と体上の位相下界]
\label{mbfc:cor:normalizer-bound}
$k=2$ で $E_1,E_2$ が幾何学的に非同種なら、
代数閉包上でも $N_H=\operatorname{Cent}_{\operatorname{FM}(Y)}(H)$ である。

さらに $K\supset\C$ 上の順序付き実表示を考え、次のいずれかを仮定する：
\begin{enumerate}
 \item $k\ge3$ かつ $j(E_i)$ が相異なる。
 \item $k=2$ かつ $E_1,E_2$ が幾何学的に非同種である。
\end{enumerate}
source の $H$ 作用を target の $H$ 作用へ
ある抽象自己同型 $\alpha\in\operatorname{Aut}(H)$ に沿って移す
coherent Fourier--Mukai 比較核があれば、必ず $\alpha=\mathrm{id}$ である。
従ってこれらの範囲では、中心 band の基底を忘れ、抽象自己同型を許しても、
指定された中心作用についての全核型判定は変わらない。
特に定理\ref{mbfc:thm:sharp-all-types}の独立な Hesse 族では、
この意味で標識を弱めても最小次数は $3^{k-1}$ と $3^k$ のままである。
\end{corollary}
\begin{proof}
二つの非同種因子の場合、任意の Fourier--Mukai 自己同値は外積となる。
$\operatorname{Aut}(H)=\{\pm1\}$ であり、反転を誘導する外積があれば、
前命題と同じく各因子で
$F_iR_iF_i^{-1}\simeq R_i^{-1}[n_i]$、従って $3n_i=4$ となる。
これで第一の主張を得る。

一般の $K$ 上では、まず代数閉包へ移す。
各 $p_i$ はそこで実 Picard 線束を持つので、既存の Picard 比較により
source の $H$ 作用を target の作用へ $\alpha=\mathrm{id}$ で移す基準核がある。
与えられた比較とその逆を合成すれば、target の $H$-normalizer になる。
前命題または第一の主張から $\alpha=\mathrm{id}$ である。
$H$ は定数群なので、この結論は元の $K$ 上でも同じである。
有限体拡大後にも仮定は保たれる。
従って全型判定とその Kummer 分裂次数に新しい核型は加わらない。
独立な Hesse 族の因子は幾何学的に非同種であることを
定理\ref{mbfc:thm:sharp-all-types}で既に確認している。
\end{proof}

この系は中心作用を比較データから完全に除く主張ではない。
normalizer に入らない残余圏自己同値の可能性は除いておらず、
完全に無標識な Kuznetsov 圏の全比較核や、その最小定義次数を分類していない。

\subsection{偏極自己同型が加える Kummer 位相}

\begin{proposition}[平面自己同型の period multiplier]
\label{mbfc:prop:multiplier}
$E/K$ を滑らかな平面三次、$L=\OO_E(1)$ とし、
$R_0=T_{\OO_E}(-\otimes L)$ とその三次方程式から来る指定周期
$\pi:R_0^3\simeq[2]$ を固定する。
$\phi\in\operatorname{Aut}_K(E,L)$ に対し、
$\eta:\phi_*R_0\simeq R_0\phi_*$ の指定周期への適合障害を
\[
 \nu_E(\phi)\in H^1(K,\mu_3)=K^\times/K^{\times3}
\]
と置く。これは群準同型である。

$W=H^0(E,L)$ とし、$\phi^*L\simeq L$ の選択から得る
pullback $g:W\xrightarrow{\sim}W$ を取る。
$f\in\operatorname{Sym}^3W$ を $E$ の方程式とし、
\begin{equation}\label{mbfc:eq:multiplier}
 g(f)=c_g f
\end{equation}
と書けば、pushforward の規約において
\begin{equation}\label{mbfc:eq:multiplier-class}
 \nu_E(\phi)=[c_g].
\end{equation}
$g$ の scalar 変更は $c_g$ を三乗倍するので右辺は選択に依存しない。
$R_0$ の冪、$[1]$ および $\varepsilon=R_0^{-1}[1]$ の multiplier は零である。
\end{proposition}
\begin{proof}
比較射の Hom は一次元であり、比較射を scalar 倍すると三重図式はその三乗倍になる。
二つの自己同型の比較射を合成すると三重図式の scalar は積となるため、
$\nu_E$ は群準同型である。

符号を確認する。$\phi_*$ に対応する graded 加群の制限は、
自由加群の行列に $g^{-1}$ を施す構成である。
従って MF の曲率は $f$ から $g^{-1}(f)=c_g^{-1}f$ へ移る。
この構成は内部 shift と MF の指定周期を保つ。
一方、命題\ref{prop:nfm-period-equation-frame-change}により
\[
 \pi_{c_g^{-1}f}=c_g\pi_f.
\]
従って同じ $\pi_f$ を source と target に使って三重図式を表示すると、
比較射の三乗の scalar は $c_g^{-1}$ である。
period-normalized 比較射の方程式は $v^3=c_g$ となるから、
式\eqref{mbfc:eq:multiplier-class}を得る。
この MF/window の比較が $\phi_*$ であることは、商の graded 加群上での制限と
section・evaluation による window の自然性から従う。

方程式を $u f$ に替えても両側の period が同じ $u^{-1}$ 倍を受けるので
$\nu_E$ は変わらない。
rotation の冪と shift については、MF の指標 shift と三次周期の中心 coherence を用いる。
それら自身の標準交換射は三乗周期を保つ。周期の cohomological degree は二なので、
$[1]$ との交換にも余分な符号は残らない。
\end{proof}

\begin{example}[基本楕円領域でも非零となる multiplier]
\label{mbfc:ex:nonzero-multiplier}
$K=\C(a)$ とし、$\mu\in\C$ を、Hesse 三次の判別式が非零かつ
$j\ne0,1728$ となるように選ぶ。
\[
 f=x^3+a y^3+a^2z^3-3\mu a xyz,
 \qquad g(x)=az,\quad g(y)=x,\quad g(z)=y
\]
と置くと $g^3=aI$、$g(f)=a f$ である。
代数閉包上で変数を三乗根により変更すれば、これは滑らかな Hesse 三次となる。
従って $g$ は $E=V(f)$ の実際の偏極自己同型を定め、
\[
 \nu_E(\phi)=[a]\ne0
\]
となる。非零性は $a$-進付値で分かる。
幾何的な $\operatorname{Aut}(E,0)$ は $\{\pm1\}$ であるから、
この位数三の自己同型は幾何的には非零三捩れ点の平行移動である。
一般の体上で全ての $\nu_E$ を零とみなすことはできない。
\end{example}

\subsection{全核型のスカラー障害と位相判定}

各因子の Picard gerbe と位相を第\ref{mbpd:sec:picard-descent}節の通り
\[
 \delta_i=\delta_{E_i}(p_i)\in H^2(K,\mathbb G_m),\qquad
 \gamma_i\in H^1(K,\mu_3),\qquad
 \alpha=\sum_i\delta_i
\]
と書く。$\gamma_i$ は、個々の $p_i$ が実線束で表されない場合にも
factor Picard gerbe 上の scalar 重みが零であることにより $K$ 上で定義される。
本節では $H^1(K,\mu_3)$ を加法記法で書く。

\begin{theorem}[因子分離の下での全中心同変核の存在判定]
\label{mbfc:thm:full-criterion}
上の順序付き実表示について、$k\ge3$、または
$k=2$ かつ $J_{E_1,\overline K},J_{E_2,\overline K}$ が非同種であるとする。
以下の条件は、$K$ の任意の体拡大後にも同じ形で成立する。

\begin{enumerate}
 \item 全ての underlying $H$-intertwiner の幾何的同型類は、
       $\boldsymbol\phi=(\phi_i)_i$、
       $\phi_i\in\operatorname{Aut}(E_{i,\overline K},L_i)$ と
       $\lambda\in\Lambda_k$ により一意に表示される。
       その型が $K$ 上に実際の核を持つためには $\phi_i$ が $K$ 上に定義される必要があり、
       このとき実現 gerbe の類は正確に $\alpha$ である。
 \item 中心 $G$ の作用を保つ通常の残余圏の比較核が $K$ 上に存在するための必要十分条件は
       \begin{equation}\label{mbfc:eq:ordinary-criterion}
        \alpha=0,\qquad
        \exists\,\phi_i\in\operatorname{Aut}_K(E_i,L_i):\quad
        \gamma_i+\nu_{E_i}(\phi_i)
           \text{ が全ての }i\text{ で等しい}
       \end{equation}
       である。
 \item さらに rotation 比較と指定三次周期を保つ核が存在するための必要十分条件は
       \begin{equation}\label{mbfc:eq:strong-criterion}
        \alpha=0,\qquad
        \exists\,\phi_i\in\operatorname{Aut}_K(E_i,L_i):\quad
        \gamma_i+\nu_{E_i}(\phi_i)=0\quad\text{for all }i
       \end{equation}
       である。
\end{enumerate}
ここで全ての shift と数値型を許している。
二つの同種因子の全算術的核型、中心作用そのものを忘れた核については主張していない。
\end{theorem}
\begin{proof}
まず代数閉包上で $[M_i]=p_i$ を満たす線束を選び、
$F_0=\boxtimes_i(-\otimes M_i)$ とする。
一因子比較により $F_0R_{i,p}\simeq R_{i,0}F_0$ がある。
従って任意の underlying $H$-intertwiner $F$ について、
$FF_0^{-1}$ は全ての $Q_{i,0}$ と可換する。
命題\ref{mbfc:prop:centralizer}により
\begin{equation}\label{mbfc:eq:normal-form}
 F\simeq C_{\boldsymbol\phi,\lambda}F_0,
 \qquad
 C_{\boldsymbol\phi,\lambda}
    =\boxtimes_i\bigl(\phi_{i*}\varepsilon_i^{a_i}\bigr),
 \qquad[(a_i)]=\lambda
\end{equation}
となり、$(\boldsymbol\phi,\lambda)$ は一意である。
$F_0$ の幾何的同型類は各 $p_i$ だけで決まり、Galois 不変である。
また $\varepsilon_i$ は $K$ 上の関手なので、
$F$ が $K$ 上の核なら一意性により各 $\phi_i$ は Galois 不変、従って $K$ 上へ降下する。
その偏極保存性も降下する。
実際、$\phi_i^*L_i\otimes L_i^{-1}$ は幾何的に自明な実線束であり、
固有基底変換により $K$ 上でも自明である。

従って $C_{\boldsymbol\phi,\lambda}$ は $K$ 上の実際の Fourier--Mukai 自己同値となる。
これの逆を $F$ と合成した核は、代数閉包上で対角上の線束である。
cohomology の集中、対角への支持、可逆層であることを忠実平坦降下すると、
その核は $\Delta_{Y*}M$、
\[
 [M]=\sum_i\operatorname{pr}_i^*p_i\quad\text{in }\Pic^0(Y)(K)
\]
となる。
逆にこの類の実線束があれば、式\eqref{mbfc:eq:normal-form}により全ての
$K$-有理な型の underlying 核が得られる。
命題\ref{mbpd:prop:product-gerbe}により、その実現 gerbe は
$\mathfrak G_{\mathrm{Pic}}$ と band を保って同値で、類は $\alpha$ である。
個々の $\delta_i$ の消滅は必要ではなく、ここではその和の相殺を保持している。

次に coherent な構造を調べる。
各 $\phi_i$ の period-normalized commutator の torsor を $V_i(\phi_i)$ と書く。
その類は $\nu_{E_i}(\phi_i)$ である。
$\varepsilon_i$ の標準比較は周期を保つので、
型 $(\boldsymbol\phi,\lambda)$ の各因子の位相 torsor は
\[
 T_i(\phi_i)=T_i\mathbin{\wedge^{\mu_3}}V_i(\phi_i),
 \qquad [T_i(\phi_i)]=\gamma_i+\nu_{E_i}(\phi_i)
\]
である。
これは因子線束が局所的にしか存在しない場合にも、Picard gerbe 上で計算して降下できる。
核の scalar 自己同型は比較射の両端で相殺し、追加の scalar 重みは生じない。

定理\ref{mbpd:thm:coherent-corrections}の証明をこの合成比較射に適用すると、
固定型に対する全 coherent $H$-構造の torsor は
\[
 D_{\boldsymbol\phi}
   =\left(\prod_i T_i(\phi_i)\right)/\Delta\mu_3,
 \qquad
 [D_{\boldsymbol\phi}]_b
   =\bigl(\gamma_b+\nu_{E_b}(\phi_b)
          -\gamma_i-\nu_{E_i}(\phi_i)\bigr)_{i\ne b}
\]
である。
これは構成した比較射だけの集合ではない。
一つの coherent 構造があれば、任意の別の構造との差は
$\operatorname{Hom}(H,\mathbb G_m)=G$ の指標となる。
実際、同値核の degree-zero endomorphism は
$\mathrm{HH}^0(Y)=H^0(Y,\OO_Y)=K$、負の自己 Ext は零であり、
生成元の比較の差は scalar、積の coherence はちょうど指標条件である。
従って上の torsor は全ての構造を尽くす。
rotation と指定周期まで保持すると $\prod_iT_i(\phi_i)$ が残る。
核の scalar 自己同型群はどちらの場合も $\mathbb G_m$ である。

以上の実現 gerbe と有限 torsor の両方が中立となる条件が、
それぞれ式\eqref{mbfc:eq:ordinary-criterion}、
式\eqref{mbfc:eq:strong-criterion}である。
二重同変化の自然性により元の $G$-同変残余比較へ移せる。
この議論は全ての体拡大上で同じなので、指定した任意の拡大体にも適用できる。
\end{proof}

\begin{remark}[全型という範囲]
\label{mbfc:rem:scope}
前定理は幾何的核の型、その各体拡大上の実現 gerbe、および coherent 比較の存在を扱う。
任意の非被約基底上の全核 stack を外積型へ分解する相対定理は主張していない。
また、例\ref{mbfc:ex:nonzero-multiplier}により、
Picard 補正型の $D$ が非自明であることだけから、
一般の $K$ 上で全ての中心同変比較核の不存在を推論することはできない。
全型では上の $\nu_{E_i}$ による補正を先に調べる必要がある。
\end{remark}

\subsection{Hesse 族における全型の最小次数}

\begin{lemma}[固定 Hesse level の全偏極自己同型]
\label{mbfc:lem:hesse-lifts}
$K\supset\C$ 上の滑らかな Hesse 三次
\[
 f=u(x^3+y^3+z^3)+vxyz,\qquad j(E)\ne0,1728
\]
を考える。
任意の体拡大 $K'/K$ 上の $\operatorname{Aut}_{K'}(E,L)$ の全ての元は、
$\zeta\in\C$ を原始三乗根として
\[
 A(x,y,z)=(x,\zeta y,\zeta^2z),\qquad
 B(x,y,z)=(y,z,x),\qquad
 S(x,y,z)=(y,x,z)
\]
の積による $\operatorname{GL}_3(\C)$ の lift を持つ。
それらは $f$ を厳密に保つので、全ての period multiplier は零である。
\end{lemma}
\begin{proof}
$o=(1:-1:0)$ を flex 原点とする。
この level の行列表示は
\cite[\S4, pp.~246--247]{NFMArtebaniDolgachev}の標準表示である。
$A,B$ は射影的に可換する位数三の元であり、相異なる九つの射影変換を生成する。
$j\ne0,1728$ より原点保存自己同型は $\{\pm1\}$ なので、
これらの位数三の元は三捩れ平行移動であり、$E[3]$ を尽くす。
$S$ は $o$ を固定する非恒等な involution なので inversion である。
次数三偏極を保存する曲線自己同型は
$E[3]\rtimes\{\pm1\}$ だから、この十八元が幾何的な群全体である。
全ての行列が $\C$ 上で定義されるため、任意の $K'$ 上でも追加の元はない。
$A,B,S$ は $x^3+y^3+z^3$ と $xyz$ をそれぞれ保つ。
命題\ref{mbfc:prop:multiplier}を適用すれば最後の主張を得る。
\end{proof}

\begin{theorem}[独立境界位相の全比較核に対する鋭い次数]
\label{mbfc:thm:sharp-all-types}
系\ref{mbpd:cor:phase-independence}の独立な $k$ 個の Hesse 族を取る。
すなわち Sklyanin パラメータを
\[
 (a_i,b_i,c_i)=(1+d_i,1,c_i),\qquad
 K_0=\C(d_1,\ldots,d_k)
\]
とし、$c_i\in\C$ は同系の滑らかさと基本楕円条件を満たす一般の定数とする。
source と target には自然な中心・点三次直線比較を使い、指定周期を固定する。
このとき
\begin{enumerate}
 \item 中心 $G$ の作用を保つ通常の残余比較核を得る有限体拡大の最小次数は
       $3^{k-1}$ である。
 \item さらに rotation 比較と指定三次周期を保つ核を得る有限体拡大の最小次数は
       $3^k$ である。
\end{enumerate}
ここでは Picard 補正型への制限を置かず、全ての dg Morita 比較核を許している。
同じ結論は $\operatorname{Frac}\C[[d_1,\ldots,d_k]]$ 上でも成り立つ。
中心作用を忘れた比較についての最小性は主張しない。
\end{theorem}
\begin{proof}
Hesse パラメータの三乗
\[
 \kappa_i=\frac{(a_i^3+b_i^3+c_i^3)^3}{a_i^3b_i^3c_i^3}
\]
は $a_i=1+d_i,b_i=1$ の下で $d_i$ の非定数有理関数である。
Hesse 三次の $j$ は $\kappa_i$ の非定数有理関数である
\cite[\S2, equations~(4)--(5)]{NFMArtebaniDolgachev}。
従って各因子の $j_i$ は $d_i$ だけの非定数有理関数である。
従って $j_1,\ldots,j_k$ は $\C$ 上代数的に独立である。
同種な二つの楕円曲線の $j$ は代数関係を満たすため、
因子は幾何学的に互いに非同種である。
後者の事実は、固定した $j$ の曲線の有限位数部分群とその商が
$\overline{\C(j)}$ 上に定義されることからも分かる。
完備局所体へ移しても、同種性があれば同じ代数関係が生じるので非同種性は保たれる。

Hesse flex により各 $p_i$ の実 Picard 線束があるので $\alpha=0$ である。
補題\ref{mbfc:lem:hesse-lifts}により、任意の有限体拡大上でも全ての
$\nu_{E_i}$ が零となる。
従って定理\ref{mbfc:thm:full-criterion}の全型判定は、
Picard 補正型の判定と一致する。

命題\ref{mbpd:prop:boundary-phase}と系\ref{mbpd:cor:phase-independence}により、
$K_0$ 上ではある $a_i(d_i)\in\C[d_i]_{(d_i)}^\times$ に対し
\[
 \gamma_i=[u_i],\qquad u_i=3d_i/a_i(d_i),\qquad
 v_{d_j}(u_i)=\delta_{ij}
\]
と書ける。
従って $\gamma_i$ の張る Kummer 部分空間の次元は $k$、
$\gamma_b-\gamma_i$ $(i\ne b)$ の張る部分空間の次元は $k-1$ である。
$K_0$ は $\mu_3$ を含むので、それぞれの全類を消す最小分裂体は
\[
 K_0\bigl(\sqrt[3]{u_b/u_i}:i\ne b\bigr),\qquad
 K_0\bigl(\sqrt[3]{u_i}:1\le i\le k\bigr)
\]
であり、その次数は $3^{k-1}$ と $3^k$ である。
全ての類を消す任意の有限体拡大はこれらの分裂体を含むので、この次数は下界でもある。
前定理がその拡大上の全型比較の存在を与えるため、最小性を得る。

完備局所環では全単元が三乗根を持つので $\gamma_i=[d_i]$ となる。
同じ付値計算と Kummer 理論を適用すれば、そこでの主張も従う。
ここで周期が境界まで延びる根拠は
系\paperref{main}{mbrn:cor:factor-boundary}の実 Koszul 最上 syzygy 直線を用いた一因子 window であり、
二捩れ境界で退化し得る $\det W_i$ への交代化同定は用いていない。
\end{proof}

\begin{corollary}[全中心同変核の大域切断の不存在]
\label{mbfc:cor:no-global-kernel}
前定理の有限型 Hesse 積族の基本楕円非空開上では、
中心 $G$ の作用を保つ通常の比較核の大域切断は存在しない。
任意の非空開へさらに縮小しても同じである。
\end{corollary}
\begin{proof}
大域核があれば一般点 $K_0$ へ引き戻せるが、
前定理の最小次数は $3^{k-1}>1$ である。
\end{proof}

\begin{remark}[既知の分類入力と本節の帰結]
\label{mbfc:rem:novelty}
アーベル多様体の Fourier--Mukai 完全列、楕円曲線の自己同値分類、
Hesse level と Kummer 理論は既知の入力である。
本節ではそれらを、満支持比較の指定された双対作用へ適用し、
因子 shift の相殺を除き、全核型の scalar 和と自己同型 multiplier による判定を得た。
Hesse 境界例では追加 multiplier が消えるため、
Picard 型で得た最小次数を全中心同変核へ拡張できる。
この適用・結論と同一の先行結果があるかどうかの新規性調査は未完了である。
\end{remark}

\subsection{二つの同種因子における混合中央化群}

以下では $k=2$、$Q=R_1^{-1}R_2$、
$q=\operatorname{diag}(r^{-1},r)$ とする。
主偏極を用いた Hom 行列の Rosati 転置を $X\mapsto X^\dagger$ とし、
\[
 J_2=\begin{pmatrix}0&1\\-1&0\end{pmatrix},\qquad
 X^\sharp=J_2^{-1}X^\dagger J_2
\]
と書く。全体では $J=\operatorname{diag}(J_2,J_2)$ を用いる。
Orlov の isometry 条件は $g^\sharp g=I$ である。
この式は $\operatorname{End}(E_i)=\mathbf Z$ を仮定していない。

\begin{proposition}[同種対に対する整数係数の混合元]
\label{mbfc:prop:two-isogenous-mixing}
$E_1,E_2$ が同種であれば、
$\operatorname{Cent}_{U(Y\times\widehat Y)}(q)$ は
非対角成分を持つ無限位数の元を含む。
この主張は $E_1\simeq E_2$ を必要とせず、CM の有無にも依存しない。
\end{proposition}
\begin{proof}
isogeny $u:E_2\to E_1$ を選び、その次数を $d$、
主偏極に関する dual isogeny を $u^\dagger$ とする。
$uu^\dagger=[d]$、$u^\dagger u=[d]$ である。
整数行列
\[
 A_0=\begin{pmatrix}-1&-1\\0&1\end{pmatrix},\qquad
 B_1=A_0,\qquad B_3=A_0(r-r^{-1})
     =\begin{pmatrix}-3&-1\\6&3\end{pmatrix}
\]
は、$s=1,3$ について
\[
 r^{-1}B_s=B_sr,\qquad
 rB_s=B_sr^{-1},\qquad
 B_s^\sharp=-B_s,\qquad B_s^2=sI
\]
を満たす。
$d$ と $3d$ の双方が平方数になることはないので、
$D=ds$ が非平方となる $s\in\{1,3\}$ を選べる。

$u,u^\dagger$ を二つの座標へ成分ごとに適用して
\[
 N=\begin{pmatrix}0&uB_s\\u^\dagger B_s&0\end{pmatrix}
\]
と置くと
\[
 Nq=qN,\qquad N^\sharp=-N,\qquad N^2=DI.
\]
Pell 方程式 $a^2-Db^2=1$ の $a>1,b>0$ となる整数解を取る。
この解の存在は非平方正整数に対する通常の Pell の定理による。
すると
\[
 g=aI+bN,\qquad g^{-1}=aI-bN
\]
は整数 Hom 成分を持つ実際のアーベル多様体自己同型であり、
$g^\sharp g=(a^2-Db^2)I=I$、$gq=qg$ を満たす。
二つの非対角成分は非零である。
$D$ は非平方なので $\mathbf Q[N]\simeq\mathbf Q(\sqrt D)$ であり、
この実二次体の埋込みの下で $g$ は
$a+b\sqrt D>1$ とその逆へ移る。従って $g$ は無限位数である。
\end{proof}

\begin{remark}[数値群の整数 Hom 表示]
\label{mbfc:rem:two-integral-centralizer}
二因子の数値中央化群は、$g=(g_{ij})$ に対して
\[
 rg_{ii}=g_{ii}r,\qquad
 r^{-1}g_{ij}=g_{ij}r\ (i\ne j),\qquad
 g^\sharp g=I
\]
を課した整数 Hom 行列の群である。
ここで $r^{-1}g_{ij}=g_{ij}r$ と $rg_{ij}=g_{ij}r^{-1}$ は同値である。
可換成分は $uI+vr$、反転 intertwiner は
$A_0(uI+vr)$ と一意に書ける。
係数 $u,v$ は各成分に対応する $\operatorname{Hom}(E_j,E_i)$ に属する。
これは $r$ の $(1,2)$ 成分が $-1$ であることを使って
四つの成分方程式を解けば分かる。
従って有理化で失われる整除条件を隠しておらず、
CM を含む Hom 群のままでの記述となっている。
\end{remark}

\begin{lemma}[数値中央化元の実際の coherent 核への持上げ]
\label{mbfc:lem:cyclic-centralizer-lift}
代数閉体 $\kappa$ 上で、任意の
$g\in\operatorname{Cent}_{U(Y\times\widehat Y)}(q)$ は
指定された $\langle Q\rangle\simeq\mathbf Z/3$ 作用と
coherently 可換する実際の Fourier--Mukai 自己同値へ持ち上がる。
特に前命題の混合元は、単なる数値的候補ではない。

underlying 自己同値の同型類については正確列
\[
 1\longrightarrow
   \mathbf Z\times E_1[3](\kappa)\times E_2[3](\kappa)
 \longrightarrow \operatorname{Cent}_{\operatorname{FM}(Y)}(Q)
 \xrightarrow{\mathfrak u}
   \operatorname{Cent}_{U(Y\times\widehat Y)}(q)
 \longrightarrow1
\]
がある。固定した underlying 中央化関手上の coherent 構造の同型類は
$\mu_3$-torsor である。
\end{lemma}
\begin{proof}
Orlov の全射 \cite[Theorem~4.14]{MBFCOrlovAbelian}で $g$ を $F$ へ持ち上げる。
可換子の Orlov 像は恒等なので、ある $v\in Y\times\widehat Y$ と整数 $n$ に対し
\[
 FQF^{-1}\simeq\Phi_v[n]\,Q
\]
と書ける。$Q^3\simeq\mathrm{id}$ の指定された関係を三乗に用いると、
$I+q+q^2=0$ により右辺の三乗は $[3n]$ となる。
有界 coherent 圏上の非零 shift は恒等関手ではないので $n=0$ である。

$(1-q)(2+q)=3I$ から $1-q$ は isogeny であり、
代数閉体上で全射である。
$(1-q)w=-v$ となる $w$ を取り、$F'=\Phi_wF$ と置けば
\[
 F'Q(F')^{-1}
   \simeq \Phi_{w+v-qw}\,Q
   \simeq Q.
\]
従って underlying 中央化関手が得られる。
この補正の符号は、$\Phi_w$ を $F$ の左から合成する規約による。

同型 $\eta:F'Q\xrightarrow{\sim}QF'$ を一つ選ぶ。
$Q^3\simeq\mathrm{id}$ の指定された coherence を使って
$\eta$ を三回合成した図式は、$F'$ の scalar 自己同型 $c$ となる。
実際、可逆核の零次自己準同型は
$\mathrm{HH}^0(Y)=H^0(Y,\OO_Y)=\kappa$ であり、負の自己 Ext は零である。
$\eta$ を $t\eta$ に替えるとこの scalar は $t^3c$ になる。
$t^3=c^{-1}$ を選べば、巡回群の三乗関係を保つ比較を得る。
負の自己 Ext の消滅により比較の mapping space に更なる高次自己同型はなく、
指定された群作用の coherence とこの三乗関係が実際の coherent intertwiner を与える。
二つの選択の差は
$\operatorname{Hom}(\mathbf Z/3,\kappa^\times)=\mu_3$ である。

最後に Orlov 写像の核の $\Phi_w[n]$ が $Q$ と可換する条件は $qw=w$ である。
$r$ と $r^{-1}$ の固定部分は、上の主偏極座標でともに
$\{(x,0):x\in E_i[3]\}$ である。
これと shift の独立性から正確列を得る。
\end{proof}

ここまでの二因子の持上げは代数閉体上の主張である。
一般の定義体上で数値中央化元の全てが実核を持つことや、
その各混合型の scalar gerbe・周期位相を分類したとはしていない。
一方、同種な二因子について外積型だけで全核を尽くすことはできない、
という障害は実際の coherent 中央化核によって証明されている。

整数行列と Pell 元の有限検算は
\texttt{computations/verify\_multiblock\_isogenous.py} に収める。
このコードは Orlov の存在定理、忠実平坦降下、または coherent 核の持上げを
機械的に証明するものではない。

\begin{example}[二次同種で生じる混合反転核]
\label{mbfc:ex:two-isogenous-inversion}
二因子の場合、上の非同種性を単に $j_1\ne j_2$ へ弱めることはできない。
代数閉体 $\kappa$、$\operatorname{char}\kappa=0$ 上で
$\phi:E_1\to E_2$ を次数二の isogeny とし、主偏極による dual を $\phi^\vee$ とする。
各因子の任意の次数三線束から $R_i$ を作り、
$Q=R_1^{-1}R_2$、$D=\operatorname{diag}(r^{-1},r)$ と置く。
整数行列
\[
 s=\begin{pmatrix}1&0\\-3&-1\end{pmatrix},\qquad
 J_2=\begin{pmatrix}0&1\\-1&0\end{pmatrix}
\]
は
$s^2=I$、$sr=r^{-1}s$、$s^{\mathsf t}J_2s=-J_2$ を満たす。
$\phi,\phi^\vee$ を二つの座標へ成分ごとに作用させ、
\begin{equation}\label{mbfc:eq:two-inversion-matrix}
 U=\begin{pmatrix}s&\phi^\vee I_2\\\phi I_2&s\end{pmatrix}
   :E_1^2\times E_2^2\longrightarrow E_1^2\times E_2^2
\end{equation}
と置く。このとき
\[
 U^{-1}=\begin{pmatrix}-s&\phi^\vee I_2\\\phi I_2&-s\end{pmatrix},
 \qquad U^\dagger\operatorname{diag}(J_2,J_2)U
          =\operatorname{diag}(J_2,J_2),\qquad UD=D^{-1}U.
\]
ここで $\dagger$ は主偏極に関する Rosati 転置を含む転置である。
これらは $\phi^\vee\phi=[2]$、$\phi\phi^\vee=[2]$ と
$s^{\mathsf t}J_2+J_2s=0$ から直接従う。
従って $U$ は整数 Hom 成分を持つ実際の Orlov isometry であり、
二つの非対角成分は非零である。

これは実際の coherent 反転核へ持ち上がる。
Orlov の定理\cite[Theorem~4.14]{MBFCOrlovAbelian}で $U$ を $F$ へ持ち上げると、
ある $z\in Y\times\widehat Y$、整数 $n$ に対し
\[
 FQF^{-1}\simeq\Phi_z Q^{-1}[n]
\]
と書ける。$I+D+D^2=0$ と $Q^3\simeq\mathrm{id}$ を用いて三乗すると
$[3n]\simeq\mathrm{id}$、従って $n=0$ である。
$1-D^{-1}$ は isogeny なので、代数閉体上で
$(1-D^{-1})w=-z$ となる $w$ が存在する。
$F'=\Phi_wF$ と置けば
\[
 F'Q(F')^{-1}\simeq\Phi_{z+w-D^{-1}w}Q^{-1}\simeq Q^{-1}.
\]
同型 $F'Q\simeq Q^{-1}F'$ の三乗図式のずれは scalar であり、
その逆数の三乗根で比較射を規格化すれば
$\mathbf Z/3$ の指定された関係を保つ。
可逆核の自己準同型は $\kappa$、負の自己 Ext は零なので、
これは実際の coherent $H$-normalizer を与える。
この補正は Orlov 像を変えないため、得られた核は因子の外積ではない。

一般の二次同種対では $j(E_1)\ne j(E_2)$ となるので、
異なる $j$ だけでは二因子の反転を排除できない。
この例では代数閉体上の存在を述べている。
一般の定義体上への核・coherence の降下や、
個々の因子 rotation・指定周期まで保つ比較、
混合型の scalar gerbe と Kummer 位相は結論していない。
\end{example}

% !TEX root = ../cubic-moduli-descent.tex
\section{二ブロック残余圏における中心標識の非正規性}
\label{mbun:sec:unmarked}

二ブロックの中心 $\mu_3$ 作用は、全 $\OO$-局所型からは内在的に得られるが、
通常の残余圏の全自己同値が保つ部分群ではない。
本節では、滑らかな分離和三次四様体の三つの平面から、
この中心群を生成元の逆転を許しても保存しない球面 twist を構成する。
さらに、同じ共役軌道に属する相異なる中心部分群が無限個生じることを示す。
証明には三つの平面の交点数だけを用い、
K3 曲面への幾何化や三次四様体の全自己同値群の分類を仮定しない。

以下、$\kappa\supset\C$ を代数閉体とし、
\[
 E_i=V(f_i)\subset\mathbb P(W_i),\qquad \dim W_i=3,\qquad i=1,2
\]
を滑らかな平面三次曲線とする。滑らかな三次四様体
\[
 X=V(f_1+f_2)\subset\mathbb P(W_1\oplus W_2)
\]
と、その標準 dg enhancement を備えた残余圏 $\mathcal C=\Ku(X)$ を考える。
原始三乗根 $\zeta\in\kappa$ を固定し、
\begin{equation}\label{mbun:eq:block-action}
 g[x:y]=[x:\zeta y],\qquad
 G=\langle g_*|_{\mathcal C}\rangle
\end{equation}
と置く。$g$ は $\OO_X(1)$ を保つので $\mathcal C$ を保つ。
自己同値の群では、dg 自己同値を自然同型で割った同型類を考える。
以下の自己同値は、実際の perfect Fourier--Mukai 核の制限として得られる。

\subsection{平面のイデアル層と三つの球面類}

\begin{lemma}[平面のイデアル層の球面性]
\label{mbun:lem:plane-spherical}
任意の平面 $P\simeq\mathbb P^2\subset X$ に対し、
イデアル層 $I_P$ は $\mathcal C$ の $2$-球面対象である。
従って
\begin{equation}\label{mbun:eq:spherical-twist}
 T_P(F)=
 \operatorname{Cone}\!\left(
   \operatorname{RHom}_X(I_P,F)\otimes_\kappa I_P\longrightarrow F
 \right)
\end{equation}
は $\mathcal C$ の dg 自己同値である。
\end{lemma}
\begin{proof}
残余圏の定義と Serre 関手は
\[
 \mathcal C=
 \{F:H^\bullet(X,F(-j))=0\ (j=0,1,2)\},
 \qquad S_{\mathcal C}=[2]
\]
である \cite[p.~220, (2)--(3)]{NFKuznetsovCubicFourfolds}。
完全列
\[
 0\longrightarrow I_P(-j)\longrightarrow\OO_X(-j)
       \longrightarrow\OO_P(-j)\longrightarrow0
\]
を用いる。$j=0$ では定数切断の写像が同型で、それ以外のコホモロジーは零である。
$j=1,2$ では両方の線束のコホモロジーが全て零である。
従って $I_P\in\mathcal C$ となる。

法束の完全列は
\begin{equation}\label{mbun:eq:plane-normal}
 0\longrightarrow N_{P/X}\longrightarrow\OO_P(1)^{\oplus3}
       \longrightarrow\OO_P(3)\longrightarrow0
\end{equation}
である。$H=c_1(\OO_P(1))$ とすると
\[
 c_1(N_{P/X})=0,\qquad c_2(N_{P/X})=3H^2,\qquad
 \det N_{P/X}\simeq\OO_P .
\]
Riemann--Roch より $\chi(P,N_{P/X})=-1$ であり、
正則埋込みの局所 Ext と局所大域スペクトル系列から
\[
 \chi_X(\OO_P,\OO_P)
 =\sum_{j=0}^2(-1)^j\chi(P,\mathop{\bigwedge}\nolimits^jN_{P/X})
 =1-(-1)+1=3 .
\]
また
\[
 \chi_X(\OO_X,\OO_X)=\chi_X(\OO_X,\OO_P)=1,\qquad
 \chi_X(\OO_P,\OO_X)=\chi(P,\det N_{P/X}[-2])=1 .
\]
従って $[I_P]=[\OO_X]-[\OO_P]$ から
\begin{equation}\label{mbun:eq:plane-self-euler}
 \chi_X(I_P,I_P)=2
\end{equation}
を得る。

$I_P$ は階数一の torsion-free 層である。
その自己準同型は $X\setminus P$ 上で正則関数による乗法となり、
$X$ の正規性と $\operatorname{codim}_X P=2$ から $X$ 全体へ延びる。
従って $\operatorname{Hom}(I_P,I_P)=\kappa$ である。
層の負次数 Ext は零であり、$S_{\mathcal C}=[2]$ により
$\operatorname{Ext}^2(I_P,I_P)=\kappa$ および
$\operatorname{Ext}^q(I_P,I_P)=0$（$q>2$）を得る。
式\eqref{mbun:eq:plane-self-euler}から $\operatorname{Ext}^1(I_P,I_P)=0$ である。
Serre 双対性の完全対と合わせ、これは $\mathcal C$ 内での $2$-球面性である。

球面 twist とその双対評価による逆関手は
\cite[Definition~2.9, Proposition~2.10]{MBUNSeidelThomas}の構成で得られる。
この構成は dg Hom、評価および上記の完全対を用いるので、
admissible な dg 部分圏 $\mathcal C$ に適用できる。
具体的には、$I_P^\vee=\operatorname{RHom}_X(I_P,\OO_X)$ として
\[
 \operatorname{Cone}(I_P^\vee\boxtimes I_P\longrightarrow\OO_{\Delta_X})
       \in\operatorname{Perf}(X\times X)
\]
が式\eqref{mbun:eq:spherical-twist}を表す。
これは $\mathcal C$ を保ち、その制限が自己同値である。
$I_P$ の球面性を $\mathrm D^b(X)$ 全体で主張しているわけではない。
\end{proof}

\begin{lemma}[flex 接線から得られる三平面]
\label{mbun:lem:three-planes}
$g(P_r)=P_{r+1}$（添字は $\mathbf Z/3$）を満たす三つの平面
$P_0,P_1,P_2\subset X$ が存在する。
相異なる二平面の scheme-theoretic な共通部分は同じ直線であり、
\begin{equation}\label{mbun:eq:plane-intersections}
 P_r^2=3,\qquad P_r\cdot P_s=-1\quad(r\ne s)
\end{equation}
である。
\end{lemma}
\begin{proof}
各 $E_i$ の flex $o_i$ とその接線
$\ell_i=\mathbb P(U_i)$（$\dim U_i=2$）を選ぶ。
ある $c_i\in\kappa^\times$ と非零線型式 $u_i\in U_i^\vee$ により
\[
 f_i|_{U_i}=c_i u_i^3
\]
と書ける。ここで $\ker u_i$ が flex を表す。
$\beta^3=-c_1/c_2$ を選ぶと、$\mathbb P(U_1\oplus U_2)\simeq\mathbb P^3$ 上で
\[
 c_1u_1^3+c_2u_2^3
   =c_2\prod_{r=0}^2(u_2-\beta\zeta^r u_1)
\]
となる。その三つの線型因子で定まる平面を $P_r$ とする。
式\eqref{mbun:eq:block-action}はこれらを巡回させ、
任意の二つの共通部分は
\[
 L=V(u_1,u_2)\subset\mathbb P(U_1\oplus U_2)
\]
という直線である。

自己交点は式\eqref{mbun:eq:plane-normal}の $c_2$ から $3$ となる。
相異なる $P=P_r,Q=P_s$ は直線 $L$ に沿って clean に交わる。
実際、線型平面の接空間の共通部分は $T_L$ である。
余剰束 $\mathcal E$ は
\[
 0\longrightarrow N_{L/Q}\simeq\OO_L(1)
       \longrightarrow N_{P/X}|_L\longrightarrow\mathcal E\longrightarrow0
\]
により与えられる階数一束である。
$\det N_{P/X}\simeq\OO_P$ より $\mathcal E\simeq\OO_L(-1)$ であり、
余剰交差公式から $P\cdot Q=\deg c_1(\mathcal E)=-1$ を得る。
\end{proof}

$K_{\mathrm{num}}(\mathcal C)_\mathbf Q$ に Euler 形式の符号を変えた
$(x,y)=-\chi(x,y)$ を入れ、
$\delta_r=[I_{P_r}]$ と書く。
相異なる平面 $P,Q$ について、Riemann--Roch の次数四成分から
$\chi_X(\OO_P,\OO_Q)=P\cdot Q$ である。
これは $\operatorname{ch}(\OO_P^\vee)$ と
$\operatorname{ch}(\OO_Q)$ の最初の非零成分がいずれも余次元二の基本類であり、
その積が既に最高次数であるためである。
従って
\[
 \chi_X(I_P,I_Q)=P\cdot Q-1
\]
となり、三つの類の Gram 行列は
\begin{equation}\label{mbun:eq:gram}
 B=\bigl((\delta_r,\delta_s)\bigr)_{r,s}
  =\begin{pmatrix}
    -2&2&2\\2&-2&2\\2&2&-2
   \end{pmatrix}
  =2J_3-4I_3,\qquad \det B=32 .
\end{equation}
特に $\delta_0,\delta_1,\delta_2$ は独立である。
$W=\mathbf Q\delta_0\oplus\mathbf Q\delta_1\oplus\mathbf Q\delta_2$ と置くと、
$g_*$ はこれらを巡回させる。従って $g_*|_{\mathcal C}$ の位数は正確に三である。

\subsection{中心群の非正規性と無限共役軌道}

\begin{theorem}[通常の残余圏で中心群は正規でない]
\label{mbun:thm:non-normal}
上の任意の滑らかな分離和三次四様体に対して、
$G\subset\operatorname{Aut}_{\mathrm{dg}}(\mathcal C)$ は正規部分群でない。
具体的には、$T_0=T_{P_0}$ とすると
\[
 T_0g_*T_0^{-1}\not\simeq g_*,
 \qquad T_0g_*T_0^{-1}\not\simeq g_*^{-1}.
\]
従って中心作用を生成元の逆転まで保存する自己同値だけでは、
通常の $\Ku(X)$ の全自己同値を覆わない。
\end{theorem}
\begin{proof}
式\eqref{mbun:eq:spherical-twist}を Grothendieck 群へ写すと、
$T_i=T_{P_i}$ の作用は反射
\begin{equation}\label{mbun:eq:reflection}
 r_i(x)=x-\chi(I_{P_i},x)\delta_i
       =x+(x,\delta_i)\delta_i
\end{equation}
である。特に $r_i^2=1$ であり、逆 twist も同じ数値作用を持つ。
式\eqref{mbun:eq:gram}を用いて
\[
 r_0g_*r_0(\delta_0)
  =-r_0(\delta_1)
  =-\delta_1-2\delta_0
\]
を得る。一方、$g_*\delta_0=\delta_1$ および
$g_*^{-1}\delta_0=\delta_2$ である。
三つの類は独立だから、これらの数値作用はどちらも一致しない。
自然同型な関手の数値作用は等しいため、主張が従う。
ここで自己同値群から数値格子への作用が忠実であることは必要ない。
\end{proof}

\begin{corollary}[同じ共役軌道内の無限個の中心部分群]
\label{mbun:cor:infinite-orbit}
$\mathsf U=T_0T_1$ と置くと、部分群
\[
 G_n=\mathsf U^nG\mathsf U^{-n}\qquad(n\in\mathbf Z)
\]
は互いに相異なる。これらは全て同じ共役軌道に属する位数三の部分群であり、
その coherent な作用の同変化は全て
\[
 \mathcal C^{G_n}\simeq\operatorname{Perf}(E_1\times E_2)
\]
と同値である。
\end{corollary}
\begin{proof}
$W$ の順序付き基底 $(\delta_0,\delta_1,\delta_2)$ に関し、
\[
 r_0=\begin{pmatrix}-1&2&2\\0&1&0\\0&0&1\end{pmatrix},
 \qquad
 r_1=\begin{pmatrix}1&0&0\\2&-1&2\\0&0&1\end{pmatrix},
 \qquad
 h=r_0r_1=\begin{pmatrix}3&-2&6\\2&-1&2\\0&0&1\end{pmatrix}.
\]
$N=h-I$ とすると $N^3=0$ なので、全ての $m\in\mathbf Z$ に対し
\[
 h^m=I+mN+\binom m2N^2 .
\]
$g_*$ の $W$ 内の固定部分は
$s=\delta_0+\delta_1+\delta_2$ が張る直線であり、
\begin{equation}\label{mbun:eq:unipotent-orbit}
 h^ms=
 (1+4m^2+2m)\delta_0+
 (1+4m^2-2m)\delta_1+\delta_2 .
\end{equation}
$\mathsf U^m$ が $G$ を正規化するなら、$g_*$ を逆元へ移す場合を含めて、
$h^m$ は固定直線 $\mathbf Qs$ を保つ。
式\eqref{mbun:eq:unipotent-orbit}の第三座標から比例係数は一であり、
$4m^2+2m=4m^2-2m=0$ より $m=0$ である。
従って $G_n=G_{n'}$ なら $n=n'$ となる。

元の幾何学的 $G$ 作用を $\mathsf U^n$ とその逆で共役すれば、
関手の合成同型も含む coherent な $G_n$ 作用を得る。
この共役は同変化の同値
$\mathcal C^G\simeq\mathcal C^{G_n}$ を与える。
また定理\paperref{main}{mbrd:thm:split-product}、二ブロックでは
\cite[Theorems~0.1 and~5.8]{TBHu}により
$\mathcal C^G\simeq\operatorname{Perf}(E_1\times E_2)$ である。
従って、同変化がこの積圏になるという性質だけでも、
指定された部分群 $G$ は一意に選び出せない。
\end{proof}

\begin{corollary}[非可換二ブロック幾何点への移送]
\label{mbun:cor:nc-non-normal}
代数閉体 $\kappa\supset\C$ 上の、定理\ref{mbs:thm:presentation-stack}の
基本楕円領域に属する二ブロック満支持実表示 $A$ を取る。
その中心 $Z_0=(\mu_3^2)/\Delta\mu_3\simeq\mu_3$ の残余作用は、
$\Ku_{\mathrm{coh}}(A)$ の全 dg Morita 自己同値群において正規でない。
さらに、この自己同値群内の共役軌道は無限個の相異なる部分群を含む。
\end{corollary}
\begin{proof}
定理\paperref{main}{mbrc:thm:finite-point-comparison}は、代数閉体上では拡大なしで
実際の $Z_0$-同変 dg Morita 同値
\[
 \Ku_{\mathrm{coh}}(A)\simeq\Ku(X_{\mathrm{pt}})
\]
を与える。右辺は本節の滑らかな分離和三次四様体の残余圏である。
この同値と逆を用いて $T_0,T_1$ および上の中心群作用を移せば、
定理\ref{mbun:thm:non-normal}と系\ref{mbun:cor:infinite-orbit}がそのまま移る。
因子の非同種性や一般点という追加条件は使わない。
\end{proof}

\begin{remark}[標識の忘却について得られたことと残ること]
\label{mbun:rem:scope}
本節は、通常の残余圏の全自己同値に対して不変な
指定中心部分群を回復する方針を否定する。
中心作用の共役類を内在的に認識する可能性や、
比較の前後へ適切な自己同値を合成して中心を一致させる可能性は否定しない。
また、非代数閉体 $K$ 上で新たな無標識比較核が存在すること、
あるいは中心同変比較の最小定義体次数が標識の忘却によって下がることは、
この反例だけからは従わない。
構成に用いた flex、平面および twist の元の体への降下は別途必要である。
とりわけ、系\ref{mbun:cor:infinite-orbit}は
「非共役な作用が無限個ある」という主張ではない。

平面の交差理論、球面 twist、および同変化の積表示は既知の入力である。
本節の具体的な非正規性と無限共役軌道の先行研究との一致については、
現時点で網羅的な新規性確認を終えていない。
式\eqref{mbun:eq:gram}と式\eqref{mbun:eq:unipotent-orbit}の有限行列計算は
\texttt{computations/verify\_multiblock\_isogenous.py} で独立に検算できる。
この検算は球面性や幾何学的な実現の証明を代替しない。
\end{remark}

% !TEX root = ../cubic-moduli-descent.tex
\section{中心 Veronese 錐による次数一の復元}
\label{sec:nfm-veronese}

本節では $n\geq3$, $N=n+2$, $r=n-1$ とする。
周囲の二次代数を $B$、中心三次商を $A=B/(w)$ と書く。
次数一の生成元空間を $V$、その中心部分を $U\subset V$ とする。
一般次元では相対 Serre 関手が与えるのは $r$ 次の twist である。
これから次数一を回復するためには、$r$-Veronese の中心に含まれる
Veronese 錐を用いる。
以下の代数的復元は、Serre 双対性の構成とは独立である。

\subsection{中心 Veronese の同型の持上げ}

\begin{lemma}[Veronese 錐の同型層]
\label{lem:nfm-veronese-cone}
$R$ を可換 $\mathbf C$-代数、$U,U'$ を同じ階数 $s\geq2$ の
有限局所自由 $R$-加群、$d\geq1$ とする。
$C=\operatorname{Sym}_R U$, $C'=\operatorname{Sym}_R U'$ と置く。
次数を $d$ で割った Veronese 代数について、自然な写像
\[
 \underline{\operatorname{Isom}}_R(U,U')
 \longrightarrow
 \underline{\operatorname{Isom}}_{\mathrm{gr},R}(C^{(d)},(C')^{(d)})
\]
は fpqc 層として $\mu_d$-torsor である。
\end{lemma}

\begin{proof}
基底上で局所的に $U,U'$ を自由としてよい。
Veronese 環の Proj はそれぞれ $\mathbf P^{s-1}_R$ であり、
その新しい次数一の線束は通常の $\mathcal O(d)$ である。
次数付き同型はこの線束を保つ射影空間の同型 $f$ を誘導する。
各幾何学的ファイバーでは $f^*\mathcal O(1)=\mathcal O(1)$ なので、
$M=f^*\mathcal O(1)\otimes\mathcal O(-1)$ は各ファイバーで自明である。
射影空間の $H^0(\mathcal O)=R$, $H^i(\mathcal O)=0$ $(i>0)$ と
cohomology and base change により、$M=\pi^*L$ となる基底の可逆層 $L$ がある。
この議論は有限表示のデータを noetherian 部分環へ降ろしてもよいので、
$R$ の noetherian 性を必要としない。
基底を被覆して $L$ を自明化すると、次数一切断への作用から
$f$ は線型同型 $g:U\to U'$ に持ち上がる。
ここでは projective bundle の切断の同定を用いた
\cite[Section~27.21]{NFMStacksProjective}。

与えた Veronese 同型と $\operatorname{Sym}^{(d)}g$ の差は Proj 上で恒等である。
$\mathcal O(d)$ 上の自己同型は基底の単元 $a\in R^\times$ であり、
$C^{(d)}$ は新次数一で生成されるので、この差は新次数 $j$ に $a^j$ を掛ける。
忠実平坦な有限 \'etale 被覆
$R\to R[t]/(t^d-a)$ 上で $t$ を選び、$g$ を $tg$ に置き換えれば差を消せる。
二つの持上げの比は射影空間上で恒等なので $\lambda\operatorname{id}$ であり、
$\lambda^d=1$ が必要十分である。
従って各同型の持上げ層は $\mu_d$ により単純推移的に作用され、局所的に非空である。
\end{proof}

\begin{theorem}[中心を用いた任意次数の Veronese 復元]
\label{thm:nfm-central-veronese}
$R$ を可換 $\mathbf C$-代数、$d\geq1$ とする。
$A,A'$ は非負次数付きで次数一生成の $R$-代数、$A_0=A'_0=R$ とする。
階数 $s\geq2$ の有限局所自由加群 $U,U'$ に対し、中心的な単射
\[
 C=\operatorname{Sym}_R U\longrightarrow A,
 \qquad C'=\operatorname{Sym}_R U'\longrightarrow A'
\]
があり、$A,A'$ はそれぞれ $C,C'$ 上平坦とする。
さらに任意の可換 $R$-代数 $R_1$ への変更後に
\[
 Z(A_{R_1})_d=\operatorname{Sym}^d_{R_1}U_{R_1},
 \qquad
 Z(A'_{R_1})_d=\operatorname{Sym}^d_{R_1}U'_{R_1}
\]
が成り立つと仮定する。このとき制限写像
\[
 \underline{\operatorname{Isom}}_{\mathrm{gr},R}(A,A')
 \longrightarrow
 \underline{\operatorname{Isom}}_{\mathrm{gr},R}(A^{(d)},(A')^{(d)})
\]
は $\mu_d$-torsor である。
その核の $\lambda\in\mu_d$ は、元の次数 $j$ に $\lambda^j$ を掛ける。
\end{theorem}

\begin{proof}
$U$ を自由にし、基底を $u_1,\ldots,u_s$ とする。
$Q=A^{(d)}$ と置く。
$z\in Z(Q)$ と次数一の元 $a\in A_1$ について
$u_1^{d-1}a$ は $Q$ に属し、
\[
 u_1^{d-1}[z,a]=[z,u_1^{d-1}a]=0
\]
である。$A$ の $C$-平坦性により $u_1$ は $A$ 上の非零因子なので、
$[z,a]=0$。次数一生成性から
$Z(Q)=Z(A)^{(d)}$ を得る。
従って $Z(Q)_1=\operatorname{Sym}^dU$ であり、
その生成する中心部分代数は内在的に $D=C^{(d)}$ と同定される。
また $Z(A)_1=U$ も従う。実際、$z\in Z(A)_1$ なら
$u_i^{d-1}z\in C_d$ なので、$i=1,2$ を用いて
$z\in C[u_1^{-1}]\cap C[u_2^{-1}]=C$ と書ける。
局所化した $C$ は局所化した $A$ に単射であり、Laurent 多項式の比較でこの交叉を取れる。
次数一部分を取れば $z\in U$ である。従って全ての次数付き同型は
指定した中心多項式環を保つ。

$Q\simeq Q'$ は $D\simeq D'$ を誘導する。
補題\ref{lem:nfm-veronese-cone}により、fpqc 局所でこの同型を
$C\simeq C'$ へ持ち上げる。
以下では Veronese 側の次数を $d$ 倍して元の次数に戻す。
各 $u_i$ について積写像
\[
 Q[(u_i^d)^{-1}]
 \otimes_{D[(u_i^d)^{-1}]}C[u_i^{-1}]
 \longrightarrow A[u_i^{-1}]
\]
は同型である。実際、左辺の第二因子は
\[
 C[u_i^{-1}]
 =\bigoplus_{j=0}^{d-1}u_i^jD[(u_i^d)^{-1}]
\]
と分解し、右辺の次数が $j$ modulo $d$ の部分も
$u_i^jQ[(u_i^d)^{-1}]$ と一致する。
よって $Q\simeq Q'$ と選んだ $C\simeq C'$ は、全てのこの局所化を同定する。

$C$ の二変数 $u_1,u_2$ に関する列
\[
 0\longrightarrow C\longrightarrow
 C[u_1^{-1}]\oplus C[u_2^{-1}]
 \longrightarrow C[(u_1u_2)^{-1}]
\]
は最初の二箇所で完全である。
これは多項式の Laurent 展開から任意の係数環上で直接確かめられる。
$A$ は $C$ 上平坦なので、これを tensor して
\[
 A=A[u_1^{-1}]\cap A[u_2^{-1}]
 \quad\text{in }A[(u_1u_2)^{-1}]
\]
を得る。従って局所化の同型は一意に $A\simeq A'$ を与える。
構成は局所化の重なりで一致し、$U$ の自明化を変えても一意性により降下する。

$C$ への持上げを固定すると $A$ の持上げは一意であり、
$C$ の二つの持上げの差 $u_i\mapsto\lambda u_i$ は
上の modulo $d$ 分解で次数 $j$ に $\lambda^j$ を掛ける。
よって持上げ層はちょうど $\mu_d$-torsor である。
全ての段階が係数変更と両立する。
\end{proof}

\begin{remark}
この定理は低次数中心の仮定を省略した Veronese 復元を主張しない。
また $s\geq2$ は最後の二変数交叉に使われる。
同型持上げの $\mu_d$ は次数スカラー $\mathbf G_m$ の部分群であり、
構造対象付き Abel 圏の新しい慣性ではない。
一つの同型の大域持上げ障害は $H^1_{\mathrm{fppf}}(-,\mu_d)$ にあり、
対象の有効降下に現れる $H^2$ の障害とは区別する。
\end{remark}

\subsection{一ブロック族での低次数中心}

\begin{proposition}[有限個の楕円位数条件]
\label{prop:nfm-center-bound}
$k$ を代数閉体、$\operatorname{char}k=0$、$U$ を階数 $r\geq2$ の
次数一次元からなる空間とし、$C=\operatorname{Sym}_kU$ とする。
$B$ は三つの active 生成元を持つ濾過付き二次代数で
\[
 \operatorname{gr}_{\mathrm{act}}B=S\otimes_k C,
 \qquad S/(\Omega)=H(E,\mathcal L,\sigma)
\]
とする。ここで $S$ は三次元楕円 AS 正則代数、
$\Omega\in Z(S)_3$ は正則元、$E$ は滑らかな楕円曲線である。
$G\in Z(B)_3$ の最高 active 記号が $\Omega$ であるとし、
$\sigma^j\ne\operatorname{id}$ が
$1\leq j\leq q:=\max(3,d)$ について成り立つとする。
このとき
\[
 Z(B)_1=U,\qquad Z(B)_3=kG\oplus\operatorname{Sym}^3U.
\]
任意の $g\in\operatorname{Sym}^3U$ に対して
$A=B/(G+g)$ は $C$ 上自由であり、
\[
 \operatorname{gr}_{\mathrm{act}}A=H(E,\mathcal L,\sigma)\otimes_k C,
 \qquad Z(A)_d=\operatorname{Sym}^dU.
\]
従って $d=r=n-1$ の場合に定理\ref{thm:nfm-central-veronese}を適用できる。
\end{proposition}

\begin{proof}
$H=H(E,\mathcal L,\sigma)$ は捻った Laurent 関数環
$k(E)[t,t^{-1};\sigma]$ に埋め込まれる。
次数 $j>0$ の非零中心元 $a t^j$ があれば、
次数一切断との可換性から $\sigma^j$ はその全ての比を固定する。
$\deg\mathcal L=3$ の完全線型系は $E$ を埋め込むので、
これらの比は $k(E)$ を生成する。従って $\sigma^j=\operatorname{id}$ となる。
ゆえに $Z(H)_j=0$ $(1\leq j\leq q)$ である。
$S/(\Omega)=H$ から
$Z(S)_1=Z(S)_2=0$, $Z(S)_3=k\Omega$ を得る。

$B$ の中心元の最高 active 記号も中心である。
次数一では正の active 成分が消え、$Z(B)_1=U$。
次数三では active 三次の成分は $k\Omega$ に属するので $G$ の倍を引き、
次いで active 次数二、一の成分を消せば純中心三次だけが残る。

$G+g$ の最高記号は正則元 $\Omega$ なので、
最高次数に関する帰納法により $G+g$ は左右正則で、表示した商の associated graded を得る。
$H$ の各次数の $k$-基底を斉次に持ち上げれば、
それらは $A$ の $C$-基底を与える。
最後に $a\in Z(A)_d$ の正の最高 active 次数 $j\leq d$ は
$Z(H)_j\otimes C_{d-j}=0$ に属するため不可能である。
よって $a\in C_d$ である。
\end{proof}

\begin{remark}[相対版で必要な普遍性]
\label{rem:nfm-relative-center}
有限型被約 $\mathbf C$-基底上で、上記の PBW 分裂と中心三次の
最高記号の直和因子性が係数変更と両立しているとする。
このとき $U\subset B_1$、$RG\oplus\operatorname{Sym}^3U\subset B_3$、
$\operatorname{Sym}^dU\subset A_d$ を除いた商は有限局所自由である。
各商から三つの active 生成元との交換子を取る写像は全閉点で単射なので、
最大小行列式の可逆開上で分裂単射になる。
従って中心の三つの等式は任意の、非平坦・非被約な底変更後にも成立する。
この分裂単射の確認をせずに、幾何ファイバーの中心計算だけを普遍的な等式へ拡張しない。

固定した $n$ では除外する条件は有限個である。
楕円曲線の平行移動を動かす空間で $\sigma^j=1$ は閉条件であり、
無限位数の平行移動は全てのこれらの閉集合を避ける。
従って所定の楕円パラメータ空間における共通開は非空である。
三様体の低次数中心条件だけで任意の $n$ を扱うことはできない。
\end{remark}

\section{一ブロックの次数付きモジュライと中心三次 torsor}
\label{sec:nfm-graded-stack}

\subsection{表示空間と全ての次数付き同型}

$V=X\oplus U$, $\dim X=3$, $\dim U=r$ とし、
$X$ の基底を $x,y,z$ とする。
$F\in\operatorname{Sym}^3V/\operatorname{Sym}^3U$ に対して
\[
 B_F=T(V)/\bigl([U,V],\ [y,z]-\operatorname{sym}(F_x),\
 [z,x]-\operatorname{sym}(F_y),\ [x,y]-\operatorname{sym}(F_z)\bigr)
\]
と置く。
$F$ は生成元を可換化した多項式として扱う規約であり、
純中心三次の変更は関係式を変えない。
$P\subset\operatorname{GL}(V)$ を $U$ を保つ parabolic、
$\overline P=P/\mathbf G_m$ とする。
$g\in P$ の $V/U$ 上の作用を $M$ とすると、作用は
\[
 g\cdot F=(\det M)^{-1}g(F)\pmod{\operatorname{Sym}^3U}
\]
である。中心スカラーはこの作用に自明に作用する。

一ブロックの PBW・楕円・中心三次の入力を満たす $P$-不変な非空開を
$T\subset\operatorname{Sym}^3V/\operatorname{Sym}^3U$ とする。
ここでいう入力は、命題\ref{prop:nfm-center-bound}の associated graded と中心元が
相対的に存在し、注意\ref{rem:nfm-relative-center}の有限局所自由性を持つことを含む。
この開の一ブロック量子化入力は既存の一般次元量子化構成から得られる
\cite{CubicQuantization}。
具体的には、純 active 三次を滑らかな Hesse 開に取り、
active 次数を下げる混合 potential の PBW 変形を用いる。
中心変数の個数はこの三変数 PBW 計算の係数環を変えるだけである。
Hesse 座標と中心平行移動で三つの混合二次交叉項を消した開では、
Rains の有限中心公式の係数に $u$ の一次式・二次式を代入できる
\cite[Theorem~3.3.2 and Section~9]{NFMEG}。
有限個の分母と最高 active 記号の係数を反転し、
$\sigma^j\ne1$ $(1\leq j\leq\max(3,r))$ を課す。
この開を $P$ で飽和すれば、所要の非空不変開を得る。
相対 PBW 分裂と中心核の普遍性は注意\ref{rem:nfm-relative-center}の有限行列で確認する。
この一ブロック入力自体を本節で新しく証明した結果とは数えない。

\begin{proposition}[表示空間の同型問題の忠実性]
\label{prop:nfm-graded-presentation}
$F,F'\in T$ に対して、$B_F\simeq B_{F'}$ の次数付き同型は、
$g(F)=(\det M)F'$ modulo $\operatorname{Sym}^3U$ を満たす $g\in P$ に一致する。
従って $[T/\overline P]$ は、この族の次数付き周囲代数を
次数スカラーで割った同型問題を表す。
\end{proposition}

\begin{proof}
次数一生成性から代数同型は $V$ 上の線型同型で決まる。
中心一次部分が $U$ なので、この同型は $P$ に属する。
標数零での分解 $V^{\otimes2}=\bigwedge^2V\oplus\operatorname{Sym}^2V$ に関して、
$B_F$ の関係空間は二次双ベクトル
$\pi_F:\bigwedge^2V\to\operatorname{Sym}^2V$ の graph である。
従って関係空間から $\pi_F$ は一意に復元される。
$\pi_F$ は三つの active 偏微分によって決まり、
それらが全て零となる三次式は $\operatorname{Sym}^3U$ に限る。
active 三重ベクトルは $M$ により $(\det M)^{-1}$ 倍されるので、
関係空間を移す条件はちょうど表示した式である。
逆向きは関係式への代入で従う。
この議論は有限局所自由関係束を用いれば係数環上でも同じである。
\end{proof}

\begin{lemma}[有限安定化子を持つ点]
\label{lem:nfm-finite-stabilizer}
$T$ の中で $\overline P$-安定化子が有限となる部分 $T^{\mathrm{fin}}$ は
非空な不変開である。
\end{lemma}

\begin{proof}
滑らかな純 active 三次式 $F_3(x,y,z)$ を所定の楕円開に取り、
\[
 F=F_3+xQ_1(u)+yQ_2(u)+zQ_3(u),\qquad
 Q_1=\sum_{i=1}^r u_i^2,\quad Q_2=\sum_{i=1}^r\lambda_i u_i^2
\]
とする。ここで $\lambda_i$ は互いに異なり、$Q_3$ は任意である。
ここではまず、全表示アフィン空間上の $\overline P$-安定化子を調べる。
この計算自体に PBW 条件は必要ない。

$\mathfrak p$ の元を active 部分 $M$、中心部分 $D$、
active から中心への項 $T_x,T_y,T_z\in U$ に分ける。
安定化の式の純 active 三次部分は
$M(F_3)=\operatorname{tr}(M)F_3$ である。
滑らかな平面三次の射影自己同型群は有限なので $M=aI_3$ である。
実際、その楕円曲線上の平行移動のうち次数三の埋込みを保つものは三捩れ点に限られ、
原点を保つ自己同型群も標数零では有限である。
active 次数二の成分は
\[
 T_x(F_3)_x+T_y(F_3)_y+T_z(F_3)_z=0
\]
である。滑らかな $F_3$ の三つの偏微分は線型独立なので全ての $T_i$ は零となる。
active 次数一の成分から $D(Q_i)=2aQ_i$ $(i=1,2)$ を得る。
$D-aI_r=H$ と置けば
\[
 H^{\mathsf t}+H=0,\qquad
 H^{\mathsf t}\Lambda+\Lambda H=0,
 \qquad \Lambda=\operatorname{diag}(\lambda_1,\ldots,\lambda_r).
\]
よって $H$ は $\Lambda$ と可換する反対称行列であり、固有値が相異なるので $H=0$。
残る安定化子の Lie 代数は次数スカラーのみである。
従って $\overline P$ の安定化子の Lie 代数は零となる。
標数零の有限型代数群は滑らかなので、この安定化子は有限である。

一般の点でも無限小安定化子は係数に多項式的に依存する有限行列の核である。
従って核が零となる条件は全表示空間の非空開である。
表示アフィン空間は既約で、$T$ も非空開なので、両者の交わりは非空である。
\end{proof}

\begin{theorem}[一ブロックの代数的周囲モジュライ]
\label{thm:nfm-ambient-dimension}
$\mathfrak B_n=[T^{\mathrm{fin}}/\overline P]$ は滑らかな既約有限型
Deligne--Mumford stack で、その次元は
\[
 b_n=\binom{n+4}{3}-\binom{n+1}{3}
       -\bigl((n-1)^2+3(n-1)+8\bigr)
     =\frac{n^2+7n-4}{2}
\]
である。これは次数付き代数の stack についての結論であり、
周囲の全ての無偏極 Abel 圏同値を分類するとは主張しない。
\end{theorem}

\begin{proof}
$T^{\mathrm{fin}}$ は次元
$\binom{r+5}{3}-\binom{r+2}{3}$ のアフィン空間の非空開である。
$\overline P$ は次元 $r^2+3r+8$ の滑らかな連結代数群である。
従って商 stack は滑らかで既約、有限型である。
安定化子は有限で無限小部分を持たないので対角は非分岐であり、商は DM である。
商 stack の次元公式から表示した値を得る。
\end{proof}

\subsection{中心三次の射影空間と加法 torsor}

\begin{theorem}[一般次元の中心三次 torsor]
\label{thm:nfm-central-torsor}
$T^{\mathrm{fin}}$ 上で
\[
 \mathcal E=Z(\mathcal B)_3,\qquad
 \mathcal W=\operatorname{Sym}^3\mathcal U,\qquad
 \mathcal L=\mathcal E/\mathcal W
\]
と置く。$\mathcal L$ は可逆層である。
$\mathbf P_{\mathrm{lines}}$ を一次元部分束の規約とし、
\[
 J=\mathbf P_{\mathrm{lines}}(\mathcal E)
      \setminus\mathbf P_{\mathrm{lines}}(\mathcal W)
\]
とする。この $J\to T^{\mathrm{fin}}$ は
$\mathcal{H}om(\mathcal L,\mathcal W)$ の加法群の torsor であり、
商によって
\[
 \tau_n:\mathfrak C_n=[J/\overline P]\longrightarrow\mathfrak B_n
\]
を得る。射 $\tau_n$ はスキームで表現可能、アフィン、滑らか、全射で、相対次元は
\[
 s_n=\binom{n+1}{3}
\]
である。$\mathfrak C_n$ は滑らかな既約有限型 DM stack で、
\[
 c_n=\dim\mathfrak C_n
 =\binom{n+4}{3}-(n+2)^2+3n-2.
\]
$J$ の全ての線は左右正則な中心三次元を与え、
その商は一般次元の三次 Hilbert 級数を持つ。
\end{theorem}

\begin{proof}
$0\to\mathcal W\to\mathcal E\to\mathcal L\to0$ を考える。
$J$ の $S$-点である直線 $M\subset\mathcal E_S$ は
$M\to\mathcal L_S$ が同型であるという条件を満たす。
従って $J$ はこの完全列の分裂の層を表す。
二つの分裂の差は $\mathcal{H}om(\mathcal L,\mathcal W)$ の一意な切断であり、
局所自由層の列は Zariski 局所で分裂する。
よって $J$ は所定の加法群の torsor である。

中心スカラーの重みは $\mathcal U$ 上で一、
$\mathcal E,\mathcal W,\mathcal L$ 上で三である。
従って $\mathcal{H}om(\mathcal L,\mathcal W)$ 上の重みは零で、
$\overline P$ へ降下する。一方 $J$ も直線の空間なので中心スカラーは自明に作用する。
分裂の差による作用と torsor 同型は同変であり、商へ降下する。
任意の $S\to\mathfrak B_n$ を $\overline P$-torsor 上で表示して基底変更すると、
得られる射はアフィン空間束である。
fpqc アフィン降下 \cite[Lemma~35.37.1, Tag~0245]{NFMStacksAffineDescent} により、
元の基底でもアフィンスキームで表現される。
滑らかさ、全射性、相対次元も降下する。

局所的に一つの中心元 $G$ を選べば $J$ の全ての線は $k(G+g)$
$(g\in\operatorname{Sym}^3U)$ と書ける。
命題\ref{prop:nfm-center-bound}により最高 active 記号は常に $\Omega$ であり、
正則性と商の PBW 基底が得られる。
特に次数 $j$ の Hilbert 関数は
$\binom{j+N-1}{N-1}-\binom{j+N-4}{N-1}$ である。
最後に滑らかなアフィン空間束と前定理から stack の性質と次元公式を得る。
\end{proof}

\begin{remark}[ファイバーの意味]
周囲代数とその同定を固定した幾何学的 $2$-ファイバーは
$\mathbf A^{s_n}$ の torsor である。
周囲の同型類だけを固定した同型類の集合には有限群
$\operatorname{Aut}_{\mathrm{gr}}(B)/\mathbf G_m$ の作用が残る。
また $\mathcal U,\mathcal E$ 自体を $\mathfrak B_n$ 上の通常の束として降下させたわけではない。
原点を選ばずに torsor を零切断付きベクトル束と同一視しない。
\end{remark}

\begin{proposition}[複数ブロックでの正しい相対形]
\label{prop:nfm-multiblock-torsor}
一般に $0\to W\to E\to L\to0$ が有限局所自由層の完全列で
$\operatorname{rank}L=k$ なら、
$\mathbf P_{\mathrm{lines}}(E)\setminus\mathbf P_{\mathrm{lines}}(W)$ は
$\mathbf P_{\mathrm{lines}}(L)$ 上の
$\mathcal{H}om(\mathcal O(-1),W)$ の加法 torsor である。
従って $k\geq2$ では固定周囲上の中心線の空間にも
射影的な $k-1$ 次元の方向が残る。
\end{proposition}
\begin{proof}
$\mathbf P_{\mathrm{lines}}(L)$ 上の普遍直線
$\mathcal O(-1)\subset L$ に沿って元の完全列を引き戻す。
$\mathbf P(E)\setminus\mathbf P(W)$ の点は、その引戻し列の分裂と一対一に対応する。
前定理と同じ分裂の差の議論を適用する。
\end{proof}


% !TEX root = ../nc-cubic.tex
\section{一ブロック族の普遍相対双対性}
\label{sec:nfm-relative-duality}

本節は、次数付きモジュライから構造対象付き Abel 圏へ移る際の
相対双対入力を構成する。
$N=n+2\geq5$、$r=N-3=n-1$ とする。
三様体稿\cite{Threefold}の potential 解消と導来局所化の方法を、
中心変数が $r$ 個の場合に展開する。
有限中心性、楕円平行移動の有限位数性、あるいは tails の dg smoothness は使わない。

\subsection{中心差の Koszul 解消と真正な双加群双対}

次の仮定は、既に用いた一ブロックの相対 PBW 入力を、
双加群解消の構成に必要な形で明記したものである。
$R$ は被約な有限型 $\mathbf C$-代数、
$Q=R[u_1,\ldots,u_r]$ とし、全ての $u_\alpha,x,y,z$ の次数を一とする。
$Q$ 上の cyclic potential
\[
 \Phi=xyz-yxz-\operatorname{Symm}(F),\qquad
 B=Q\langle x,y,z\rangle/(\partial_x\Phi,\partial_y\Phi,\partial_z\Phi)
\]
を考える。純中心三次は cyclic 微分で消えるので任意に除いてよい。
active 次数の最高三次部分を $\Phi_3$、その Jacobian 代数を
$S=R\langle x,y,z\rangle/(\partial\Phi_3)$ とする。
以下では次を仮定する。
\begin{enumerate}
\item 各 $S_a$ は有限射影 $R$-加群であり、全ての幾何学的ファイバーは
三次元二次 Koszul CY 代数で Hilbert 級数 $(1-t)^{-3}$ を持つ。
\item active PBW 同型 $\operatorname{gr}_{\mathrm{act}}B=S\otimes_RQ$ があり、
その濾過は $Q$-加群として分裂する。
この同型と分裂は任意の係数変更後にも PBW 性を与える。
\end{enumerate}
滑らかな純 active 三次を持つ Jacobian 族では、体上の入力は
Mori--Smith\cite[Theorem~4.2]{MoriSmith}と
Berger--Taillefer\cite[Theorems~3.1, 3.6]{NFMBeTa}による。
相対 PBW 分裂は入力論文\cite{NCCubicA}の表示チャートで構成される。
以下の結論はこれらの体上の定理を任意係数環へ直接引用したものではなく、
有限射影複体と普遍恒等式から導く。

\begin{theorem}[任意個の中心変数を持つ相対 CY 解消]
\label{thm:nfm-relative-ambient-cy}
上の仮定の下で、$B$ は有限な相対二次 Koszul 双加群解消
\[
 K_B^{-i}=B\otimes_RW_i\otimes_RB\quad(0\leq i\leq N),
 \qquad K_B\xrightarrow{\sim}B
\]
を持つ。$W_i$ は内部次数 $i$、階数 $\binom Ni$ の有限射影加群であり、
$W_0=R$、$W_N$ は可逆である。
この解消とその片側 augmentation は任意の可換 $R$-代数への係数変更後にも完全である。
また active 変数と中心変数の順序を固定すると、真正な次数付き双加群として
\begin{equation}
\label{eq:nfm-relative-ambient-cy}
 \Theta_{B/R}:=
 \mathbf R\underline{\operatorname{Hom}}_{B\otimes_RB^{\mathrm{op}}}
       (B,B\otimes_RB^{\mathrm{op}})
 \simeq B(N)[-N]
\end{equation}
を得る。この同型も任意の、非平坦・非被約な係数変更と両立する。
\end{theorem}

\begin{proof}
まず $Q$ に相対的な長さ三の potential 複体を作る。
$x_1=x,x_2=y,x_3=z$、$r_i=\partial_{x_i}\Phi$ とし、
$D=B\otimes_QB^{\mathrm{op}}$ と置く。
自由基底 $e_i,f_i,e$ の内部次数をそれぞれ $1,2,3$ とする。
\[
 K_Q:\quad 0\longrightarrow D(-3)\longrightarrow D(-2)^3
 \longrightarrow D(-1)^3\longrightarrow D\longrightarrow0
\]
の微分を
\[
 d_1(e_i)=\delta_i,\qquad
 d_2(f_i)=\sum_jH_{ji}e_j,\qquad
 d_3(e)=\sum_i\delta_i f_i
\]
で定める。ここで
\[
 \delta_i=x_i\otimes1-1\otimes x_i^{\mathrm{op}},\qquad
 H_{ji}=\sum_{r_i=a x_j b}a\otimes b^{\mathrm{op}}
\]
とし、最後の和は $r_i$ を自由 $Q$-代数の語で展開して係数付きで取る。
普遍微分の telescoping と cyclic 恒等式 $\sum_i[x_i,r_i]=0$ により
\[
 \sum_jH_{ji}\delta_j=r_i\otimes1-1\otimes r_i^{\mathrm{op}},\qquad
 \sum_i\delta_iH_{ji}=r_j\otimes1-1\otimes r_j^{\mathrm{op}}.
\]
従って $K_Q$ は複体である。
active 濾過の随伴複体は $S$ の potential 解消を $Q$ へ延長したものである。
$S$ の解消の各内部次数は有限射影 $R$-加群の有界複体であり、
増大複体は全幾何学的ファイバーで完全である。
有限射影複体を右端から局所的に分裂することにより、
この完全性は任意係数変更後にも保たれる。
active 濾過は各全次数で有限なので、元の $K_Q\to B$ も普遍的に完全である。

自己双対性は fibre ごとの CY 性から推測しない。
反自己同型 $\tau(a\otimes b^{\mathrm{op}})=b\otimes a^{\mathrm{op}}$ により
右 $D$-加群を通常の双加群へ戻す規約を取ると、
cyclic 語の二文字の位置を交換する恒等式から
\[
 H_{ji}=\tau(H_{ij}),\qquad \tau(\delta_i)=-\delta_i
\]
を得る。
双対複体の微分を $d^j(f)=(-1)^{j+1}f d$ とすれば、
\[
 1^\vee\mapsto e,\quad e_i^\vee\mapsto-f_i,\quad
 f_i^\vee\mapsto e_i,\quad e^\vee\mapsto-1
\]
は複体の同型 $K_Q^\vee\simeq K_Q(3)[-3]$ を与える。
従って $\Theta_{B/Q}\simeq B(3)[-3]$ である。
これは Berger--Taillefer の potential 解消の自己双対性を
同じ符号規約で係数環 $Q$ 上に書いた計算である。

次に $C=B\otimes_RB^{\mathrm{op}}$ と置く。
PBW 分裂により $B$ は $Q$ 上射影で、$C$ は
$Q\otimes_RQ=R[u_L,u_R]$ 上平坦である。
従って
\[
 z_\alpha=u_{\alpha,L}-u_{\alpha,R}\quad(1\leq\alpha\leq r)
\]
は $C$ 上の中心正則列で、$D=C/(z_1,\ldots,z_r)$ である。
順序を指定したこの正則列の Koszul 解消を双対化すると、
真正な $D$--$C$ 双加群として
\[
 \mathbf R\underline{\operatorname{Hom}}_C(D,C)\simeq D(r)[-r]
\]
を得る。左 $D$-作用は定義域の右掛算による前合成から生じ、
唯一の cohomology 上では通常の左掛算となる。
中心列であるため normalizing automorphism は現れない。
change of rings の随伴により
\[
 \mathbf R\underline{\operatorname{Hom}}_C(B,C)
 \simeq \mathbf R\underline{\operatorname{Hom}}_D
       (B,D(r)[-r])
 \simeq B(r+3)[-r-3]
\]
となる。全ての右作用を保持したままの随伴であり、
反転規約によって \eqref{eq:nfm-relative-ambient-cy} を得る。
$D$ は $C$-perfect、$B$ は $D$-perfect なので $B$ は $C$-perfect でもある。

最後に二次 Koszul 模型を指定する。
$V=B_1$、$R_2=\ker(V^{\otimes2}\to B_2)$ とし、
$W_i\subset V^{\otimes i}$ を全ての位置の $R_2$ の交わりで定める。
各体値では、全次数を固定した有限の active 濾過の bar spectral sequence により、
$\operatorname{gr}_{\mathrm{act}}B=S\otimes k[u]$ の Koszul 性が $B$ へ移る。
Hilbert 級数は $(1-t)^{-N}$ なので
$\dim W_i=\binom Ni$、$W_{N+1}=0$ である。
$W_i$ を定める有限行列に対し、被約な有限型基底上の一定階数の判定を
$2\leq i\leq N+1$ に適用する。
これにより $W_i$ は有限射影で、その形成は任意係数変更と可換する。
標準 Koszul 双加群複体の各全次数は有限射影複体で、全ファイバーで完全であるから、
同じ局所分裂の議論で普遍的に完全である。
片側で見れば $B$ の有界な射影解消なので分裂し、augmentation との tensor も完全となる。
先に構成した双対同型をこの有限解消で表すと、各内部次数の cone は有界な
flat 加群の完全複体であるから、任意係数変更後にも同型である。
\end{proof}

\subsection{中心三次商の全 twist と compact 対象}

以下では可換環 $R$ 上の $B$ について、定理\ref{thm:nfm-relative-ambient-cy}の
有限 Koszul 解消と双対同型が任意係数変更に両立して与えられているとする。
また $w\in Z(B)_3$、$A=B/(w)$ について、全整数 $a$ の列
\begin{equation}
\label{eq:nfm-relative-central-split}
 0\longrightarrow B_{a-3}\xrightarrow{w}B_a\longrightarrow A_a\longrightarrow0
\end{equation}
が完全で、$A_a$ は有限射影 $R$-加群であると仮定する。
従ってこの列は係数線型に分裂し、任意底変更後にも正則である。
一ブロックの中心三次 $G+g$ では、PBW 最高記号の正則性と
各次数の一定階数によりこの入力が得られる。

\begin{theorem}[普遍的線コホモロジーと相対 properness]
\label{thm:nfm-relative-line-cohomology}
任意の可換 $R$-代数 $T$ と整数 $i,j$ に対し、$a=j-i$ と置くと
\begin{equation}
\label{eq:nfm-relative-line-cohomology}
 \mathbf R\operatorname{Hom}_{\operatorname{QGr}A_T}
   (\mathcal O(i),\mathcal O(j))\simeq
 \begin{cases}
  (A_T)_a[0],&a\geq0,\\
  0,&-r<a<0,\\
  (A_T)_{-a-r}^{\vee}[-n],&a\leq-r.
 \end{cases}
\end{equation}
ここで双対は $T$ 上で取る。
この同型は任意の係数変更と両立する。
全ての $\mathcal O(i)$ は compact であり、
$\bigoplus_{i=0}^{N-1}\mathcal O(i)$ は
$D(\operatorname{QGr}A_T)$ の compact generator である。
compact 対象間の RHom は $T$-perfect で、その形成は任意の導来係数変更と可換する。
\end{theorem}

\begin{proof}
基底を $T$ に変更して $R=T$ と書く。
$M_{<q}=B/B_{\geq q}$ は有限射影係数加群であり、
$K_q=M_{<q}\otimes_BK_B$ は有限生成射影右 $B$-加群の長さ $N$ の解消である。
従って
\[
 C_B(M)=\varinjlim_q\underline{\operatorname{Hom}}_B(K_q,M)
\]
は通常の derived torsion $\mathbf R\Gamma_B(M)$ を計算する。
各項は exact な関手で filtered colimit と直和を保ち、横方向の長さが $N$ なので、
非有界複体にも有限幅の全複体として適用できる。

この模型が torsion 複体上で恒等となることも確認する。
内部次数が上に有界な加群 $L$ では、各固定内部次数について
$\underline{\operatorname{Hom}}_B(K_q,L)$ は $q\gg0$ で
$\underline{\operatorname{Hom}}_B(K_B,L)$ に安定し、その cohomology は $L$ のみである。
任意の torsion 加群は有限個の斉次元で生成される、内部次数が上に有界な
部分加群の filtered colimit である。
従って $\mathbf R\Gamma_B(L)=L$ となり、有界振幅による truncation 比較から
同じ結論が torsion cohomology を持つ非有界複体にも成立する。
torsion 部分圏の包含は完全で、その右随伴が $\Gamma_B$ である。
上で計算した導来随伴の単位が torsion 複体上で同型となるので、
$D(\operatorname{Tors}B)$ は $D(\operatorname{Gr}B)$ に完全忠実に埋め込まれ、
その像は torsion cohomology を持つ複体全体である。
また $\pi_B$ は完全なので、その右随伴 $\omega_B$ は K-injective 複体を保つ。
$\pi_B\omega_B=\operatorname{id}$ から
$\pi_B\mathbf R\omega_B=\operatorname{id}$ が従う。
従って非有界導来 Gabriel 商は
$D(\operatorname{Gr}B)/D_{\mathrm{tors}}(\operatorname{Gr}B)
\simeq D(\operatorname{QGr}B)$ であり、局所化には三角
\[
 \mathbf R\Gamma_B(M)\longrightarrow M\longrightarrow
 \mathbf R\omega_B\pi_BM\longrightarrow
\]
があり、$\mathbf R\omega_B$ は直和を保つ。
従って全ての $\mathcal O_B(i)$ は compact generators となる。
augmentation Koszul 解消を tails に送ると $N+1$ 個の twist の完全列を得る。
その両端の係数が可逆なので、任意の連続する $N$ 個の twist が全 twist を thick に生成する。

有限射影係数加群 $M$ に対し、項別の tensor--Hom 随伴は
\[
 \mathbf R\underline{\operatorname{Hom}}_B(M,B)
 \simeq \Theta_{B/R}\otimes_B^{\mathbf L}M^\vee
 \simeq M^\vee(N)[-N]
\]
を与える。追加の左作用も保持される。
$M=M_{<q}$ として colimit を取ると、真正な双加群として
\[
 \mathbf R\Gamma_B(B)\simeq B^\vee(N)[-N]
\]
となる。
局所化三角の内部次数 $a$ により、周囲の線コホモロジーは
$B_a$ $(a\geq0)$、零 $(-N<a<0)$、
$B_{-a-N}^{\vee}[-N+1]$ $(a\leq-N)$ である。

中心商の導来圏を扱う前に、通常の derived torsion の比較
\begin{equation}
\label{eq:nfm-relative-torsion-comparison}
 i_*\mathbf R\Gamma_A(M)\simeq\mathbf R\Gamma_B(i_*M)
\end{equation}
を確認する。
$J$ を injective $A$-加群とする。induction と restriction の随伴により
\[
 \operatorname{Ext}^q_B(B/B_{\geq l},i_*J)
 \simeq\operatorname{Hom}_A(\operatorname{Tor}^B_q(B/B_{\geq l},A),J)
\]
である。
$q>1$ の Tor は零である。
$q=1$ の kernel は、元の次数が $l-3,l-2,l-1$ の部分にのみ存在するので、
$l'\geq l+3$ に対する遷移は零である。
colimit を取ると $i_*J$ は $\Gamma_B$-acyclic となる。
先の有限幅模型を injective 解消へ適用すれば、非有界複体にも比較が成立する。
従って $\mathbf R\Gamma_A$ も有界振幅を持ち、直和を保ち、torsion 複体上では恒等である。
$B$ の場合と同じ導来随伴の議論により、
$D(\operatorname{Tors}A)\simeq D_{\mathrm{tors}}(\operatorname{Gr}A)$ である。
また $\pi_A$ の完全性と $\pi_A\mathbf R\omega_A=\operatorname{id}$ から
$D(\operatorname{Gr}A)/D_{\mathrm{tors}}(\operatorname{Gr}A)
\simeq D(\operatorname{QGr}A)$ を得る。
この局所化の fibre は $\mathbf R\Gamma_A$ なので、$\mathbf R\omega_A$ も直和を保つ。

中心商への induction と restriction は、この tails の非有界導来圏に随伴
$\mathbf Lj^*\dashv j_*$ を誘導する。
restriction は exact、保守的、直和保存である。
ここで $A$-torsion と忘却後の $B$-torsion は一致し、
$-\otimes_B^{\mathbf L}A$ は $w:M(-3)\to M$ の cone なので torsion 複体を保つ。
従ってこれらの随伴を導来 Gabriel 商へ降ろせる。
また
\[
 j_*\mathcal O_A(a)=\operatorname{Cone}
 \bigl(\mathcal O_B(a-3)\xrightarrow{w}\mathcal O_B(a)\bigr).
\]
周囲の線コホモロジーをこの cone に適用する。
$a\geq0$ では \eqref{eq:nfm-relative-central-split} が $A_a$ を与え、
$-r<a<0$ では両項が零である。
$a\leq-r$ では高次 cohomology 上の写像が
$B_{-a-r}^{\vee}\to B_{-a-N}^{\vee}$ となる。
これは双加群としての自然性により $w$ の双対であり、全射で kernel は
$A_{-a-r}^{\vee}$ である。
cone の次数は $N-2=n$ なので \eqref{eq:nfm-relative-line-cohomology} を得る。

$j_*$ が直和を保つので $\mathbf Lj^*\mathcal O_B(i)=\mathcal O_A(i)$ は compact である。
$j_*$ の保守性から、周囲の compact generator の像も generator となる。
compact 部分はその thick 閉包である\cite[Theorems~3.4, 3.8]{NFMKellerDg}。
線対象間の RHom は表示式により係数 perfect で、全ての係数変更と両立する。
両変数について有限 cone、shift、直和因子で閉じれば、全 compact 対象に結論が延びる。
\end{proof}

\subsection{両側局所双対から相対 Serre trace へ}

以下の双対は内部次数を反転する係数双対
$(M^\vee)_a=\mathbf R\operatorname{Hom}_R(M_{-a},R)$ である。
双加群を双対化すると左右の作用も交換する。

\begin{theorem}[中心三次商の普遍相対 Serre 双対]
\label{thm:nfm-relative-serre}
前定理の入力の下で、任意の可換 $R$-代数 $T$ に対し
$D(\operatorname{QGr}A_T)^c$ の相対 Serre 自己同値は
\[
 \mathsf S=(-r)[n]=(1-n)[n]
\]
である。すなわち任意の compact $K,L$ について、両変数に自然な同型
\[
 \mathbf R\operatorname{Hom}_T
   \bigl(\mathbf R\operatorname{Hom}(K,L),T\bigr)
 \simeq\mathbf R\operatorname{Hom}(L,K(-r)[n])
\]
がある。この同型と合成 pairing・trace は任意の係数変更と両立する。
\end{theorem}

\begin{proof}
$T=R$ と書く。前定理の証明で
\eqref{eq:nfm-relative-torsion-comparison}と非有界 Gabriel 局所化を構成した。
その $\mathbf R\Gamma_A$ は torsion cohomology の導来部分圏への右随伴である。

中心正則元の二項解消と coinduction の随伴から、有限射影係数の $A$-加群 $M$ について
\[
 \mathbf R\underline{\operatorname{Hom}}_B(M,B)
 \simeq\mathbf R\underline{\operatorname{Hom}}_A(M,A)(3)[-1].
\]
前定理で得た $B$ 上の双対を代入すると
$\mathbf R\underline{\operatorname{Hom}}_A(M,A)
 \simeq M^\vee(r)[-N+1]$ である。
$M=A/A_{\geq l}$ のとき、最初は左 $B$ と右 $A$ の作用を保持して計算するが、
唯一の cohomology 上の左 $B$-作用は全射 $B\to A$ を通る。
従ってこれは真正な $A$-双加群の同型である。
colimit と次数ごとの係数双対を取ることにより
\begin{equation}
\label{eq:nfm-relative-local-dual}
 \Gamma:=\mathbf R\Gamma_A(A)\simeq A^\vee(r)[-N+1],\qquad
 \Gamma^\vee\simeq A(-r)[N-1]
\end{equation}
を真正な双加群として得る。

ここから trace へ進むため、導来局所化を双加群のレベルで記述する。
$E=\pi A$ に左 $A$-作用を残した injective 解消 $E\to I$ を取り、
\[
 Q_a=\operatorname{Hom}^{\bullet}(I,I(a))
\]
とする。左作用付き対象の圏は、完全な monad
$M\mapsto\bigoplus_{a\geq0}A_a\otimes_RM(-a)$ の加群圏である。
その自由関手が完全なので、作用を忘れた $I(a)$ も K-injective である。
合成により $Q=\bigoplus_aQ_a$ は内部次数付き dg 代数となり、
左作用は strict な写像 $A\to Q$ を与える。
線対象族の dg Morita 表示\cite[Theorem~3.8(c)]{NFMKellerDg}により
\[
 D(\operatorname{QGr}A)\simeq D(\operatorname{Gr}Q).
\]
この表示で $\mathbf R\omega$ は $A$ への restriction、$\pi$ は extension である。
従って $Q\otimes_A^{\mathbf L}Q\simeq Q$ であり、
基礎 $A$-双加群として $Q\simeq\mathbf R\omega\pi A$ となる。

局所化三角 $\Gamma\to A\to Q\to$ を係数双対化する。
両側局所化 $\mathcal L^e(M)=Q\otimes_A^{\mathbf L}M\otimes_A^{\mathbf L}Q$ は
$Q$-双加群上では恒等である。
実際、任意の $Q$-双加群 $M$ に対し、結合則と homological epimorphism の恒等式から
\[
 Q\otimes_A^{\mathbf L}M
 \simeq(Q\otimes_A^{\mathbf L}Q)\otimes_Q^{\mathbf L}M
 \simeq M
\]
となり、右側の tensor にも同じ計算を適用する。
特に $Q^\vee$ は積を双対化した真正な $Q$-双加群なので
$\mathcal L^e(Q^\vee)\simeq Q^\vee$ である。
一方 $A^\vee$ は内部次数が非正に有界な torsion 加群なので
$\mathcal L^e(A^\vee)=0$ である。
従って \eqref{eq:nfm-relative-local-dual} を使うと、真正な $Q$-双加群の三角
\[
 Q^\vee\longrightarrow0\longrightarrow Q(-r)[N-1]
 \longrightarrow Q^\vee[1]
\]
を得る。その境界写像から
\[
 Q^\vee\simeq Q(-r)[N-2]=Q(-r)[n]
\]
が従う。必要なら係数環上 K-flat な模型を取り、$Q^e$ を導来 tensor で定める。
これにより高次 cohomology から来る両側作用も保持される。

perfect な次数付き $Q$-加群 $K,L$ に対し、評価から来る標準の dg 双対性は
\[
 \mathbf R\operatorname{Hom}_R(\mathbf R\operatorname{Hom}_Q(K,L),R)
 \simeq\mathbf R\operatorname{Hom}_Q
       (L,K\otimes_Q^{\mathbf L}Q^\vee)
\]
である。有限自由 shift では双対の定義であり、両変数について thick に延びる。
先の双加群同型を代入すれば主張を得る。
trace は $L=K$ の双対同型と恒等射での評価によって定める。
この構成は両側作用を保つ dg 写像から来るので、合成 pairing の両変数自然性も従う。

最後に係数変更を確認する。
前定理の有限射影模型は、線対象間の標準 Hom 比較が任意係数変更で同型であることを与える。
導来係数延長は torsion 加群を torsion cohomology の複体へ送るので、
dg 導来商へ降り、合成を保つ比較
$Q_R\otimes_R^{\mathbf L}T\simeq Q_T$ を与える。
局所化三角の fibre も、$B$ 上の有限幅の torsion 模型から標準写像ごと係数変更と可換する。
各内部次数の $Q$ は係数 perfect なので係数双対も可換する。
従って上の境界写像、評価、trace の全構成が任意の係数変更と両立する。
\end{proof}

\begin{remark}[供給された入力と残るモジュライ上の問題]
良い相対 PBW 表示チャートと中心三次の正則性の下で、
各次数の有限射影性、全 twist の compact 性、compact 部分の相対 properness、
真正な相対 Serre 双対、および compact Hom の任意底変更は、本節で供給された。
CY の規格化に用いた座標の順序は双対の構成のためのもので、
Abel 圏の対象に新しい標識を追加したものではない。
固定表示模型間の全同値の係数変更と fpqc 降下は、
次節の定理\ref{thm:nfm-fixed-model-comparison}で構成する。
一般の Abel 圏族がこの表示族に局所的に属することと、
表示チャートから一般の圏モジュライへの対象の効果化は別の問題として残る。
本節は tails の dg smoothness、全 Abel モジュライの代数性、
あるいは大域残余射の存在を結論しない。
\end{remark}

% !TEX root = ../nc-cubic-n-fold.tex
\section{構造対象だけを保つ Abel 同値への条件付き移行}
\label{sec:nfm-abel-conditional}

\begin{lemma}[中心正則対による全次数の切断飽和性]
\label{lem:nfm-central-saturation}
$R$ を可換環、$A=\bigoplus_{j\geq0}A_j$ を
$A_0=R$ を満たす次数一生成の $R$-代数とし、$A_1$ は有限生成とする。
階数 $s\geq2$ の有限局所自由 $R$-加群 $U$ と中心的な次数付き写像
\[
 C=\operatorname{Sym}_R U\longrightarrow A
\]
があり、$A$ は $C$-加群として平坦であると仮定する。
$\operatorname{QGr}A$ は、次数付き右 $A$-加群の圏を
\[
 \operatorname{Tors}A
 =\{M\mid \text{各斉次元 }m\in M\text{ は、ある }q
       \text{ について }mA_{\geq q}=0\}
\]
で割った Gabriel 商とし、$\mathcal O_A=\pi A$ と置く。
このとき全整数 $j$ について、左乗法による標準写像
\begin{equation}
\label{eq:nfm-central-saturation}
 A_j\xrightarrow{\ \sim\ }
 \operatorname{Hom}_{\operatorname{QGr}A}
     (\mathcal O_A,\mathcal O_A(j))
\end{equation}
は同型である。ここで $j<0$ なら $A_j=0$ とする。
特に $\operatorname{End}(\mathcal O_A)=R$ である。
結論は任意の可換 $R$-代数への係数変更後にも成立し、
標準写像は係数変更と両立する。
\end{lemma}

\begin{proof}
$A_1$ の有限生成性と次数一生成性により各 $A_q$ は有限生成である。
従って $\operatorname{Tors}A$ は部分対象・商・拡大・任意直和に閉じた
localizing 部分圏であり、表示した Gabriel 商が存在する。
拡大に閉じることは、$mA_{\geq q}$ が torsion 部分に入ったとき、
有限生成な $A_q$ の像の有限個の生成元を同時に消す上界を取り、
次数一生成性を用いれば従う。

まず、$A(j)$ に非零 torsion 部分加群は存在しない。
これは $U$ を局所自明化し、その基底元 $u$ を一つ取れば分かる。
$u$ は $C$ 上非零因子であり、$C$-平坦性によって $A$ 上も非零因子となる。
torsion 元はある $u$ の冪で零になるので、元自身が零である。
この消滅は基底の局所被覆から従う。

Gabriel 商の射の表示
\cite[Chapitre~III, \S1, p.~365]{NFMGabrielQuotient} では、
射は $A/I$ が torsion である斉次右イデアル $I\subset A$ と
標的の torsion 部分対象についての filtered colimit で与えられる。
今は標的 $A(j)$ の torsion 部分は零である。
また $A/I$ は $1$ の像で生成されるので、これが torsion なら
$A_{\geq m}\subset I$ となる $m\geq1$ が存在する。
従って
\begin{equation}
\label{eq:nfm-gabriel-colimit}
 \operatorname{Hom}_{\operatorname{QGr}A}
       (\mathcal O_A,\mathcal O_A(j))
 \simeq\varinjlim_{m\geq1}
       \operatorname{Hom}_{\operatorname{Gr}A}
       (A_{\geq m},A(j)).
\end{equation}
この議論には $R$ や $A$ の noetherian 性を必要としない。

右辺の射を、次数零の右 $A$-線型写像
$f:A_{\geq m}\to A(j)$ で表す。
$U$ が自由な基底開上で、二つの異なる基底元を $u,v$ とする。
多項式環の二変数 \v{C}ech 列を $C$-平坦な $A$ と tensor すれば
\[
 A=A[u^{-1}]\cap A[v^{-1}]
       \quad\text{in } A[(uv)^{-1}]
\]
を得る。各局所化への写像は単射である。
$\alpha=f(u^m)$、$\beta=f(v^m)$ と置くと、中心性と右線型性から
\[
 \alpha v^m=f(u^mv^m)=f(v^mu^m)=\beta u^m.
\]
よって
\[
 a=\alpha u^{-m}=\beta v^{-m}
 \in A[u^{-1}]\cap A[v^{-1}]=A.
\]
$f$ の次数から $a\in A_j$ である。
任意の斉次元 $b\in A_{\geq m}$ に対して
\[
 f(b)u^m=f(bu^m)=f(u^mb)=\alpha b=a u^m b=a b u^m
\]
となり、$u$ の非零因子性から $f(b)=ab$ となる。
従って $f$ は左乗法による写像 $A\to A(j)$ に一意に延長する。
一意性は、延長の差が $u^m$ を零に送れば差の係数も零になることによる。
局所自明化上で得られる元 $a$ はこの一意性により重なりで一致し、
$R$-加群 $A_j$ の切断として貼り合わさる。
ここでは一般の Hom に局所化が可換することを仮定していない。
既に与えた $f$ を局所化し、回復した係数 $a$ を貼り合わせている。
以上を \eqref{eq:nfm-gabriel-colimit} に適用すれば
\eqref{eq:nfm-central-saturation} が従う。

最後に、任意の $R\to R_1$ に対して
$C_{R_1}=\operatorname{Sym}_{R_1}U_{R_1}$ であり、
$A_{R_1}\simeq A\otimes_C C_{R_1}$ は $C_{R_1}$ 上平坦である。
次数一の有限生成性と非負次数性も保たれる。
従って同じ証明を $R_1$ 上で繰り返せる。
同型は左乗法という標準写像から作ったものであるから、
係数変更の比較もこの写像の可換図式によって与えられる。
\end{proof}


\begin{lemma}[一般 Veronese の Gabriel 復元]
\label{lem:nfm-gabriel}
$A$ は可換環 $R$ 上の非負次数付き、次数一生成代数で、
$A_1$ は有限生成とする。任意の $d\geq1$ に対して
\[
 \operatorname{QGr}A\simeq\operatorname{QGr}A^{(d)}
\]
である。全整数 $i$ を用いた切断関手
\[
 N\longmapsto\bigoplus_{i\in\mathbf Z}
 \operatorname{Hom}(\mathcal O_A(-di),N)
\]
はこの同値の下で Gabriel 右随伴に一致し、完全忠実である。
\end{lemma}
\begin{proof}
$Q=A^{(d)}$、$F(M)=M^{(d)}$ とし、
$G(N)=N^{\mathrm{rescale}}\otimes_{Q^{\mathrm{rescale}}}A$ とする。
rescale は次数を $d$ 倍する操作である。
$FG=1$ であり、counit $GFM\to M$ の核と余核は
$d$ の倍数の次数で零となる。
次数一生成性から、その斉次元は高々 $d-1$ 個の生成元を掛ければ
$d$ の倍数の次数へ達して零になるので、これらは tails-torsion である。

$F$ が torsion を反映することも同様である。
斉次元 $m$ に次数を $d$ の倍数へ合わせる長さ $<d$ の生成語を掛け、
有限個の生成語の像を同時に消す次数上界を取ればよい。
よって $F,G$ は torsion 同値を保ち、商上で互いに逆となる。
$F(\mathcal O_A(-di))=\mathcal O_Q(-i)$ であるから、
表示した切断関手は $\operatorname{Gr}Q\to\operatorname{QGr}Q$ の Gabriel 右随伴である。
その完全忠実性が最後の主張を与える。
\end{proof}

\begin{theorem}[相対全 $\mathcal O$-保存同値の比較]
\label{thm:nfm-full-hom-conditional}
$R$ を可換 $\mathbf C$-代数、$n\geq3$, $d=n-1$ とする。
$A,A'$ は定理\ref{thm:nfm-central-veronese}の普遍中心条件と平坦性を満たすとする。
さらに任意の可換係数変更後に次の入力を仮定する。
\begin{enumerate}
\item 各 $A_j,A'_j$ は有限局所自由である。
中心平坦性と補題\ref{lem:nfm-central-saturation}により、
全整数 $j$ で標準切断写像は同型となり、負次数の Hom は零である。
特に $\operatorname{End}(\mathcal O_A)=R$ であり、
$C_d\subset A_d$ と $A'_d$ の対応する単射により $A_d,A'_d$ は忠実である。
\item 全ての $\mathcal O_A(i)$ は $D(\operatorname{QGr}A)$ で compact であり、
compact 部分は係数環上 proper な相対 Serre 双対を持ち、
その Serre 関手は $(-d)[n]$ に同型である。$A'$ も同様である。
\item 係数作用の延長、固定模型間の Hom の fpqc 降下、
標準の Gabriel 局所化が与えられた族の係数変更と両立する。
\end{enumerate}
このとき自然な fpqc 層の同型
\[
 \underline{\operatorname{Isom}}_{\mathrm{gr},R}(A,A')/\mathbf G_m
 \simeq
 \underline{\operatorname{Eq}}_{\mathcal O,R}
 (\operatorname{QGr}A,\operatorname{QGr}A')
\]
がある。右辺は構造対象を保つ全ての完全線型同値を自然同型で割った層である。
構造対象上で恒等である自然自己同型は自明である。
\end{theorem}

\begin{proof}
完全線型同値 $F$ は複体に項別に作用し、非有界導来圏と compact 部分を同定する。
相対 Hom と係数双対を保つので、表現対象の一意性から Serre 関手と可換する。
$H=\mathsf S^{-1}[n]=(d)$ とすると、$F(\mathcal O_A)=\mathcal O_{A'}$ と
一つの可換同型 $FH\simeq H'F$ の正負の反復によって、
全ての $\mathcal O_A(di)$ とその間の Hom・合成が対応する。
補題\ref{lem:nfm-central-saturation}の切断飽和性により、次数付き同型
$\phi:A^{(d)}\simeq(A')^{(d)}$ を得る。
一つの可換同型を一貫して反復しているので、
この同型は orbit 列全体の section $\mathbf Z$-代数の作用を記述する。

補題\ref{lem:nfm-gabriel}により、その全 section 代数上の作用は
元の関手 $F$ を自然同型まで決定する。
定理\ref{thm:nfm-central-veronese}で $\phi$ を fpqc 局所的に
$A\simeq A'$ へ持ち上げれば、その構造移送と $F$ は自然同型となる。
これで層としての全射性を得る。

構造移送が標識付き恒等関手となる次数付き自己同型を考える。
各 $\mathcal O_A(di)$ 上の自然同型の成分は基底の単元であり、
隣接 Hom $A_d$ の忠実性と合成の整合性から、その Veronese 上の作用は
新次数 $j$ に $a^j$ を掛ける一つの次数スカラーとなる。
fpqc 局所で $a$ の $d$ 乗根を取り、それを元の次数スカラーで除くと、
残る作用は定理\ref{thm:nfm-central-veronese}の $\mu_d$ に属する。
これも元の次数スカラーなので、核はちょうど $\mathbf G_m$ である。

自然自己同型の各 orbit 線上の成分も、隣接 Hom の忠実性から全て等しい。
$\mathcal O_A$ 上で恒等なら全て一であり、Gabriel 切断の完全忠実性から
関手全体で恒等である。最後に仮定した Hom 降下と係数変更の整合性により、
局所的に構成した比較は元の基底へ降下する。
\end{proof}

\begin{remark}[証明の状態と適用の境界]
定理\ref{thm:nfm-relative-ambient-cy}、\ref{thm:nfm-relative-line-cohomology}、
\ref{thm:nfm-relative-serre}により、良い相対 PBW 表示チャートでは、
有限 Koszul 対角、真正な相対双加群双対、compact 性、相対 properness、
Serre trace とその非平坦係数変更が供給される。
残る入力は、一般の Abel 圏族の係数作用、全同値の fpqc 降下、
およびそれらと表示模型の係数変更との整合性である。
これらを備えたモジュライ問題に対して、$\mathfrak C_n$ は
全 $\mathcal O$-保存 Abel 問題の全部分 stack へ移る。
この条件付き結論から全 Abel モジュライの代数性・成分性や
Kuznetsov 残余射の幾何化を推論しない。
\end{remark}

% !TEX root = ../cubic-moduli-descent.tex
\section{局所一ブロック型の構造対象付き圏モジュライ}
\label{sec:nfm-local-moduli}

本節では、固定表示模型間の全同値比較を対象の効果化へ進める。
扱う族の範囲を、良い一ブロック模型を fpqc 局所的に持つ族に限定する。
この局所表示条件を満たすこと自体を、任意の非可換三次多様体に対して
証明するわけではない。
一方、この範囲では次数一 shift を標識として要求せず、
構造対象を持つ Grothendieck 圏の族を代数的 stack として分類する。

次数スカラーを除く商と、その上の次数別に重みを持つ twisted algebra は
Behrend--Noohi\cite[Section~2.3]{NFMBehrendNoohi}の一般的構成に既に現れる。
同論文の偏極付き Grothendieck 圏は自己同値もデータに含む
\cite[Section~3.1.2]{NFMBehrendNoohi}。
ここで必要な追加点は、定理\ref{thm:nfm-fixed-model-comparison}により
\emph{全ての} $\mathcal O$-保存同値を表示の同型へ戻すことと、
twisted algebra から大域の $\mathcal O$-marked 圏を実際に構成することである。
後半では、一ブロックの中心階数と active 階数から
持上げ障害の正確な位数を調べる。

\subsection{係数変更を備えた圏族}

$R$ を可換 $\mathbf C$-代数とする。
$R$-線型 Grothendieck Abel 圏 $\mathcal C$ と可換 $R$-代数 $T$ に対し、
\[
 \mathcal C_T=\operatorname{Mod}_T(\mathcal C)
\]
は $T$-作用付き対象の圏を表す。
係数延長は copower $X\mapsto T\otimes_RX$ であり、
反復係数変更には標準の結合同型を用いる。
したがって $(\mathcal C,\mathcal O)$ は、各 affine $R$-scheme 上の
圏と構造対象 $(\mathcal C_T,T\otimes_R\mathcal O)$ を定める。
この係数変更構造を持つ族だけを以下で扱い、
幾何学的ファイバーの圏を単に並べたものとは区別する。

\begin{lemma}[Grothendieck 圏内の忠実平坦降下]
\label{lem:nfm-grothendieck-fpqc}
$\mathcal C$ を $R$-線型 Grothendieck Abel 圏とする。
任意の忠実平坦 $R\to T$ に対して、
\[
 T\otimes_R-:\mathcal C\longrightarrow\mathcal C_T
\]
は完全かつ保守的であり、$\mathcal C$ は $\mathcal C_T$ の
通常の fpqc descent datum の圏に同値である。
また、二つのこのような係数圏族の間の線型完全同値、自然同型、
および指定対象上の同定は fpqc 降下を満たす。
\end{lemma}

\begin{proof}
加群 $M$ による copower $M\otimes_R X$ は、自由加群による表示を用いて
余核として構成でき、$M$ と $X$ のそれぞれについて colimit を保つ。
Lazard の定理\cite[Theorem~10.81.4, Tag~058G]{NFMStacksLazard}により、
平坦 $T$ は有限自由加群の filtered colimit である。
Grothendieck 圏の filtered colimit は完全なので、$T\otimes_R-$ は完全である。

保守性のために、任意の $X$ に対して $X\to T\otimes_RX$ が単射であることを示す。
忠実平坦性により $R\to T$ は純単射で、$Q=T/R$ は平坦である。
実際、任意の $R$-加群 $M$ に対する $M\to T\otimes_RM$ は、
$T$ を tensor した後には積写像による左逆を持つため単射であり、
$\operatorname{Tor}_1^R(Q,M)=0$ が従う。
$Q=\varinjlim Q_i$ を有限自由加群による表示とし、
$T_i=T\times_QQ_i$ と置くと
\[
 0\longrightarrow R\longrightarrow T_i\longrightarrow Q_i\longrightarrow0
\]
は分裂完全で、$\varinjlim T_i=T$ である。
これらに $-\otimes_RX$ を施し、完全な filtered colimit を取ると
$X\to T\otimes_RX$ は単射となる。
完全性と合わせて保守性を得る。

対象と射の降下は、補題\ref{lem:nfm-tails-fpqc}の等化子による証明を適用する。
必要な入力は、今証明した完全性・保守性と Abel 圏内の核だけである。
具体的に、descent datum を持つ $Y\in\mathcal C_T$ の coaction を
$\theta:UY\to T\otimes_RUY$ と書けば、
\[
 X=\ker\bigl(\theta-\eta:UY\rightrightarrows T\otimes_RUY\bigr)
\]
が降下対象である。cocycle 恒等式と $T$-作用の counit 恒等式から
$T\otimes_RX\simeq Y$ を得る。
同じ等化子は射にも適用できる。

関手と自然同型の降下には命題\ref{prop:nfm-equivalence-descent}の証明を用いる。
局所同値を降下対象ごとに適用して関手を作り、逆同値も降下させる。
得られた関手の係数変更が元の関手と一致することは、自由 $T$-作用付き対象と
標準の reflexive coequalizer
\[
 T\otimes_RT\otimes_RUY\rightrightarrows T\otimes_RUY\longrightarrow Y
\]
上で確認する。同値は colimit を保つので、この確認で全対象を扱える。
指定対象上の同定は降下対象間の射として降下する。
\end{proof}

以後 $r=n-1\ge2$, $N=n+2=r+3$ とする。
第\ref{sec:nfm-graded-stack}節の $T^{\mathrm{fin}}$ を、
必要ならさらに $P$-不変な非空開 $T^{\mathrm{good}}$ へ制限する。
この開上の一ブロック模型は、普遍中心条件、中心上の平坦性、
次数別有限局所自由性、および
定理\ref{thm:nfm-fixed-model-comparison}に必要な相対双対性を
全ての係数変更後に満たすものとする。
相対 PBW 分裂と純 active Koszul CY 入力は
第\ref{sec:nfm-relative-duality}節の仮定である。
これらは本節で追加の分類定理として証明するものではない。
中心三次の torsor をこの開へ制限して $J^{\mathrm{good}}$ と書き、
\[
 \mathfrak C_n^{\mathrm{good}}=[J^{\mathrm{good}}/\overline P]
\]
と置く。以下の議論はこの stack の任意の非空開部分にも制限できる。

\begin{definition}[局所一ブロック型の圏族]
\label{def:nfm-local-category}
$\mathfrak M_n^{\mathrm{loc}}(R)$ の対象は、
$R$-線型 Grothendieck Abel 圏と指定対象の対 $(\mathcal C,\mathcal O)$ であって、
ある忠実平坦 $R$-代数 $T$ と $j\in J^{\mathrm{good}}(T)$ に対し
\[
 (\mathcal C_T,\mathcal O_T)
 \simeq (\operatorname{QGr} A_j,\mathcal O_{A_j})
\]
となるものとする。
射は $R$-線型完全同値 $F$ と同定
$\alpha:F(\mathcal O)\xrightarrow{\sim}\mathcal O'$ の対、
二射は $\alpha$ を保つ自然同型とする。
圏族の係数変更は上で定めた $\operatorname{Mod}_T$ によるものとし、
次数一 shift、次数一生成元の束、および局所表示は対象の標識に含めない。
\end{definition}

\subsection{twisted graded module による対象の効果化}

$x\in\mathfrak C_n^{\mathrm{good}}(R)$ とし、
対応する $\overline P$-torsor を $Q\to\operatorname{Spec}R$ とする。
中心拡大
\[
 1\longrightarrow\mathbf G_m\longrightarrow P
 \longrightarrow\overline P\longrightarrow1
\]
に関する $Q$ の持上げ gerbe を $\mathcal G_x$ と書く。
これは $x$ による $[J^{\mathrm{good}}/P]$ の引戻しである。
普遍次数付き代数は $\mathcal G_x$ 上の代数
$\mathcal A=\bigoplus_{j\ge0}\mathcal A_j$ を与え、
inertia の $\lambda\in\mathbf G_m$ は $\mathcal A_j$ 上で $\lambda^j$ として作用する。
この規約を「重み $j$」と呼ぶ。

\begin{lemma}[重み付き tails の大域模型]
\label{lem:nfm-weighted-tails}
$\operatorname{Gr}^{\mathrm{wt}}_{\mathcal G_x}(\mathcal A)$ を、
各 $M_j$ が重み $j$ の準連接層であるような次数付き右
$\mathcal A$-加群 $M=\bigoplus_{j\in\mathbf Z}M_j$ の圏とする。
fpqc 局所的に通常の graded torsion となる加群全体を
$\operatorname{Tors}^{\mathrm{wt}}$ と書けば、
\[
 \mathcal C_x=
 \operatorname{Gr}^{\mathrm{wt}}_{\mathcal G_x}(\mathcal A)
 /\operatorname{Tors}^{\mathrm{wt}},\qquad
 \mathcal O_x=\pi\mathcal A
\]
は $R$-線型 Grothendieck Abel 圏とその指定対象を定める。
任意の $R$-代数 $T$ について標準同値
\[
 (\mathcal C_x)_T\simeq\mathcal C_{x_T}
\]
があり、構造対象と反復係数変更を保つ。
gerbe を中立化する fpqc 被覆上では、これは通常の
$(\operatorname{QGr}A,\mathcal O_A)$ に同値である。
\end{lemma}

\begin{proof}
代数的 stack 上の準連接層の圏は Grothendieck 圏である
\cite[Proposition~96.15.2, Tag~0781]{NFMStacksQCohGrothendieck}。
$\mathbf G_m$-gerbe 上の準連接層には重み分解があり、
各重みへの射影は完全で colimit を保つ。
この主張は gerbe を中立化した後の $\mathbf G_m$-表現の重み分解で確認でき、
fpqc 降下により大域化する。
従って各重みの圏、さらにその全整数にわたる積も Grothendieck 圏である。
ここで generator は各重みへ射影した generator の直和から得られる。

重み付き次数付き層から自由右 $\mathcal A$-加群を作る関手は
\[
 (E_i)_i\longmapsto
 \left(\bigoplus_{i\le j}E_i\otimes\mathcal A_{j-i}\right)_j
\]
である。加群の核、余核、filtered colimit は各次数で計算でき、
この自由加群関手を generator に適用すると加群圏の generator が得られる。
したがって $\operatorname{Gr}^{\mathrm{wt}}(\mathcal A)$ は Grothendieck 圏である。
次数一の有限生成性により torsion は部分対象、商、拡大、直和で閉じる。
このことは fpqc 局所で通常の次数付き加群の証明に帰着する。
従ってこれは localizing 部分圏で、表示した Gabriel 商も Grothendieck 圏である。

次に係数変更を調べる。
$\mathcal G_{x_T}\to\mathcal G_x$ は affine base change であるから、
その上の準連接層は $\mathcal G_x$ 上の $T$-作用付き準連接層に等しい。
この同一視は affine atlas 上の加群の同一視と、その fpqc 降下から従い、
inertia の重みおよび $\mathcal A$-作用を保つ。
従って
\[
 \operatorname{Gr}^{\mathrm{wt}}_{\mathcal G_{x_T}}(\mathcal A_T)
 \simeq\operatorname{Mod}_T
 \bigl(\operatorname{Gr}^{\mathrm{wt}}_{\mathcal G_x}(\mathcal A)\bigr).
\]
Gabriel 商の完全左随伴と完全忠実右随伴に $T$-作用を持たせる。
補題\ref{lem:nfm-coefficient-objects}と同様に、その核はちょうど
$\mathcal A_T$-torsion である。
これで任意の、非平坦なものも含む係数変更との整合を得る。

最後に $\mathcal G_x$ を中立化すると、各重み $j$ の準連接層は通常の
$R$-加群に一意に対応し、$\mathcal A$ は通常の次数付き代数になる。
重み付き加群・torsion・構造対象はそれぞれ通常のものへ移る。
\end{proof}

\begin{theorem}[局所一ブロック型の $\mathcal O$-marked モジュライ]
\label{thm:nfm-local-moduli}
上記の良い表示の入力の下で、
\[
 \mathfrak C_n^{\mathrm{good}}
 \xrightarrow{\ \sim\ }
 \mathfrak M_n^{\mathrm{loc}},\qquad
 x\longmapsto(\mathcal C_x,\mathcal O_x)
\]
は fpqc stack の同値である。
右辺の $2$-groupoid は通常の groupoid と同値であり、
指定対象上で恒等となる自然自己同型は全て恒等である。
特に $\mathfrak M_n^{\mathrm{loc}}$ は滑らかな既約有限型 DM stack であり、
\[
 \dim\mathfrak M_n^{\mathrm{loc}}
 =\binom{n+4}{3}-(n+2)^2+3n-2.
\]
この結論は、任意の局所模型の descent datum に対し、
大域の Grothendieck Abel 圏とその構造対象が有効に存在することを含む。
\end{theorem}

\begin{proof}
まず射を比較する。二つの表示 $A=B/(w)$ と $A'=B'/(w')$ の間の
次数付き同型は、次数二の関係空間から $B\simeq B'$ を誘導する。
実際、三次関係は次数二の関係を変えず、$B$ はその二次関係で定義される。
得られた同型は三次の核直線を互いに移すので、
命題\ref{prop:nfm-graded-presentation}により $J^{\mathrm{good}}$ の $P$-作用の射に一致する。
定理\ref{thm:nfm-fixed-model-comparison}により、次数スカラーを除いたこの射の層は
全ての $\mathcal O$-保存同値の層に一致する。
したがって主張する射は fpqc 局所で完全忠実である。
補題\ref{lem:nfm-grothendieck-fpqc}により、この完全忠実性は大域化する。

局所模型の同値の自然自己同型はスカラーであり、
その $\mathcal O$ 上の成分を恒等に指定するとスカラーは一となる
\textup{(}補題\ref{lem:nfm-natural-scalars}\textup{)}。
一般の対象でも、自然自己同型を良い被覆へ延長してこの結論を適用し、
忠実平坦降下で恒等性を確認できる。
従って各同値の groupoid は離散 groupoid に同値である。
ここで構造対象の標識を忘れた二自己同型の $\mathbf G_m$ と、
表示の次数スカラーを除いた通常の stack とを混同しない。

次に $(\mathcal C,\mathcal O)\in\mathfrak M_n^{\mathrm{loc}}(R)$ を取り、
定義の被覆上で $(\mathcal C_T,\mathcal O_T)$ を模型 $A_j$ と同定する。
重なりで得られる marked equivalence の類は、前段により
商 stack の射に一意に対応する。
元の係数圏族の結合同型は cocycle を与え、
marked natural automorphism の自明性により三重重なりでの曖昧さはない。
商 stack の対象降下から $x\in[J^{\mathrm{good}}/\overline P](R)$ を得る。
補題\ref{lem:nfm-weighted-tails}で構成した $(\mathcal C_x,\mathcal O_x)$ は
同じ局所模型と同じ同値の cocycle を持つ。
両者間の局所同値は補題\ref{lem:nfm-grothendieck-fpqc}で降下するので、
$(\mathcal C,\mathcal O)\simeq(\mathcal C_x,\mathcal O_x)$ である。

同じ議論は、最初から大域の $(\mathcal C,\mathcal O)$ が与えられていない
局所対象と marked equivalence の descent datum にも適用できる。
商 stack 上の対象を先に降下し、重み付き tails を作ることで、
大域の Grothendieck 圏を得る。
したがって対象の効果化を定義に含めて仮定したのではない。
滑らかさ、既約性、有限型性、DM 性と次元は
定理\ref{thm:nfm-central-torsor}から従う。
\end{proof}

\begin{remark}[証明した範囲]
局所的に良い一ブロック表示を持つという条件の下で、
標識は $(\mathcal C,\mathcal O)$ だけで十分である。
ただし、この条件が全ての内在的モジュライ成分の一般点で成立すること、
この stack が未構成の全モジュライの開かつ閉な成分をなすこと、
任意の Abel 圏の変形が局所表示条件を保つことは、この定理には含まれない。
これらには独立の局所表示・変形・成分の定理が必要である。
\end{remark}

\subsection{大域次数付き表示の gerbe と三捩れ障害}

\begin{proposition}[表示の持上げ gerbe]
\label{prop:nfm-presentation-gerbe}
写像
\[
 \rho_n:[J^{\mathrm{good}}/P]\longrightarrow
 \mathfrak M_n^{\mathrm{loc}}
\]
は $\mathbf G_m$-gerbe である。
その class を $\beta$ とし、$e=\gcd(r,3)=\gcd(n-1,3)$ と置くと
\[
 r\beta=3\beta=0,\qquad e\beta=0.
\]
任意の affine $S$ と $x\in\mathfrak M_n^{\mathrm{loc}}(S)$ に対し、
$\beta_x=0$ であることと、$x$ が同じ良い一ブロック型の
大域次数付き代数の tails として表示されることは同値である。
ここで次数一は有限局所自由でよく、自由な基底や中心三次直線の生成元は要求しない。

特に $n\not\equiv1\pmod3$ ならこの gerbe は全体で中立である。
$ar+3b=1$ となる整数 $a,b$ を一つ固定すると、
\[
 \chi(g)=\det(g|_U)^a\det(g|_{V/U})^b
\]
がその中立化を与える。
\end{proposition}

\begin{proof}
次数スカラーは $J^{\mathrm{good}}$ に自明に作用するため、
$P$ から $\overline P$ への商の fibre は
その torsor の $P$-持上げの gerbe である。
$[J^{\mathrm{good}}/P]$ 上には、重み一の階数 $r$ の束 $\mathcal U$ と
階数三の束 $\mathcal V/\mathcal U$ がある。
その determinant はそれぞれ重み $r$ と三の可逆層である。
gerbe 上の重み $m$ の可逆層は $m\beta=0$ を与える。
これは局所表示の三重重なりの scalar cocycle $c_{ijk}$ に対し、
可逆層の遷移写像が $c_{ijk}^{m}$ を coboundary として消すことから分かる。
従って表示した annihilator を得る。

$\beta_x=0$ は持上げ gerbe の中立性、すなわち $P$-torsor への持上げの存在と同値である。
持上げを取れば $P$-同変な普遍代数が $S$ 上へ降下し、
次数一は有限局所自由、中心一次部分は階数 $r$ の部分束になる。
三次関係は中心三次の可逆部分束で生成される。
逆にこの種の大域次数付き表示の、中心一次部分を保つ frame torsor は
$P$-torsor となり、$x$ を $[J^{\mathrm{good}}/P]$ へ持ち上げる。

$e=1$ のとき、$\chi(\lambda\operatorname{id}_V)=\lambda$ である。
写像 $p\mapsto\chi(p)^{-1}p$ は $\overline P$ の $\ker\chi$ への同型を与え、
\[
 P\simeq\mathbf G_m\times\ker\chi,\qquad
 \overline P\simeq\ker\chi
\]
と中心拡大を分裂する。
従って全ての $\overline P$-torsor はこの分裂により持ち上がる。
同じことは、gerbe 上の普遍次数 $j$ の代数束を
$\mathcal A_j\otimes\chi^{-j}$ に置き換えると重みが零になることからも分かる。
これは同じ局所代数を表すが、選んだ整数 $a,b$ に依存する中立化である。
\end{proof}

\begin{remark}[対象の障害と同値の障害]
$\beta_x\in H^2_{\mathrm{fpqc}}(S,\mathbf G_m)$ は
\emph{対象}が大域次数付き表示を持つかどうかの障害である。
一つの持上げが存在すれば、他の持上げとの差は基底の可逆層 $L$ であり、
代数は $\mathcal A_j\mapsto\mathcal A_j\otimes L^{\otimes j}$ と変わる。
一方、命題\ref{prop:nfm-lifting-line}の $L_F$ は、
\emph{既に二つの大域模型を固定した後}の指定同値の持上げ障害である。
後者の $H^1$ の現象から前者の $H^2$ の効果化を推測してはならない。
また $\beta_x=0$ は次数一が自由であること、中心三次直線が自明であること、
あるいは次数一 shift の標準的な選択を意味しない。
\end{remark}

\begin{theorem}[三捩れ障害はどの非空開でも消えない]
\label{thm:nfm-brauer-dense-open}
$n\equiv1\pmod3$、すなわち $r=3m$ とする。
任意の非空開部分 stack
$\mathfrak U\subset\mathfrak M_n^{\mathrm{loc}}$ に対し、
\[
 K=\mathbf C(s,t),\qquad x\in\mathfrak U(K)
\]
で $\beta_x\ne0$ となるものが存在する。
従って $\beta|_{\mathfrak U}$ は位数がちょうど三であり、
$\rho_n$ はどの非空開部分上でも中立化しない。
この主張は非空開全体についてのもので、特殊次元や境界点の例に依存しない。
\end{theorem}

\begin{proof}
まず $K$ 上の非自明な degree 三の中心 division algebra を固定する。
$\zeta$ を原始三乗根とし、
\[
 D=(s,t)_3
 =K\langle a,b\rangle/
 (a^3-s,\ b^3-t,\ ba-\zeta ab)
\]
と置く。$a^i b^j$ $(0\le i,j<3)$ は基底であり、中心は $K$ である。
division 性は $K'=\mathbf C(s)((t))$ へ延長して確認できる。
$L=\mathbf C(s)(a)$、$a^3=s$ とすると $[L:\mathbf C(s)]=3$ で、
$D\otimes_KK'$ は
\[
 f_0+f_1b+f_2b^2,\qquad f_i\in L((t)),\qquad
 bf=\sigma(f)b,\quad\sigma(a)=\zeta a
\]
と書ける。非零元について
\[
 v(f_0+f_1b+f_2b^2)
 =\min_i\bigl(v_t(f_i)+i/3\bigr)
\]
と定める。
三つの候補値は小数部分が異なるので最小項は一意であり、
積の最小項も一意で消えない。従って積の valuation は和に等しく、
この有限次元代数に零因子はない。
有限次元代数の非零元による乗法は単射なら全射なので、これは division algebra である。
よって $D$ 自身も division algebra であり、その Brauer class は非零である。

$U\simeq\mathbf C^3\otimes\mathbf C^m$、$V/U\simeq\mathbf C^3$ を選び、
\[
 \operatorname{GL}_3\longrightarrow P,\qquad
 g\longmapsto\operatorname{diag}(g^{\oplus m},g)
\]
を考える。中心スカラーは $V$ 全体の同じスカラーへ移るから、
$\operatorname{PGL}_3\to\overline P$ を誘導する。
この写像に沿って $P\to\overline P$ を引き戻した中心拡大は
$\operatorname{GL}_3\to\operatorname{PGL}_3$ そのものである。
$D$ に対応する $\operatorname{PGL}_3$-torsor
\cite[Tag~03R6]{NFMStacksBrauerForms}をこの写像で押し出して
$\overline P$-torsor $Q$ を得ると、その持上げ障害は $[D]$ である。
Brauer class の符号規約を逆に取っても、以下の非零性には影響しない。

$\mathfrak U$ の逆像を $J'\subset J^{\mathrm{good}}$ とする。
これは非空な $\overline P$-不変開である。
表示 potential の空間
\[
 W_F=\operatorname{Sym}^3V/\operatorname{Sym}^3U
\]
には $\det(V/U)^{-1}$ で捩じった作用を用いており、
次数スカラーの重みは零である。
従って $W_F$ は $\overline P$ の線型表現であり、
$Q\times^{\overline P}W_F$ は $K$ 上の通常の有限次元ベクトル空間となる。
$T^{\mathrm{good}}_Q$ はその非空開である。
さらに $J^{\mathrm{good}}_Q\to T^{\mathrm{good}}_Q$ は
ベクトル束の加法群による torsor であり、特に滑らかである。

非空開 $J'_Q$ の $T^{\mathrm{good}}_Q$ における像は非空開である。
$K$ は無限体なので、この像には $K$-有理点 $z$ がある。
$z$ 上の fibre $J^{\mathrm{good}}_{Q,z}$ は有限次元ベクトル空間の加法群の torsor であり、
体上では自明である。
ここでは characteristic zero のため、有限 Galois 拡大上の一点の
Galois 軌道の affine barycenter を取ることでも有理点を構成できる。
従ってこの fibre は affine 空間であり、その非空開 $J'_{Q,z}$ も $K$-有理点を持つ。
この一点は $Q\to J'$ の同変写像、すなわち
$x\in[J'/\overline P](K)=\mathfrak U(K)$ を与える。
$x$ の持上げ gerbe の障害は固定した $Q$ の障害 $[D]\ne0$ である。
命題\ref{prop:nfm-presentation-gerbe}の $3\beta=0$ と合わせて、位数はちょうど三となる。
\end{proof}

\begin{remark}[一般点と幾何学的一般点]
この非中立性は、非空な稠密開に制限するだけでは
大域次数付き表示の障害を取り除けないことを意味する。
代数閉体上の個々の幾何学的点では gerbe は中立であり、矛盾はない。
定理は商 stack の幾何学的一般点に非零の体の Brauer class が残るという意味ではなく、
\emph{全ての非空開部分 stack 上で普遍 gerbe が非中立である}という主張である。
\end{remark}

\begin{remark}[Poisson 標的の偏極に対する制約]
後に構成する射がこの表示から得られる ambient 射影空間束を保つとする。
その ambient は $\overline P\to\operatorname{PGL}(V)$ による
Brauer--Severi 束であり、符号を除いて class は $\beta$ である。
もし標的を、実際の相対可逆層 $L$ を持ち、$p_*L$ が階数 $N$ で
基底変更と両立し、$L$ の全切断がその ambient 埋込みを回復する三次族に限定すれば、
ambient は $\mathbf P((p_*L)^\vee)$ となり Brauer class は零でなければならない。
従って $n\equiv1\pmod3$ では、その種の標的への ambient と両立する射を
この stack の非空開全体に作るには、gerbe の中立化という追加データが必要である。
これは Poisson 射そのものの不存在ではない。
Brauer--Severi ambient 内の三次族、または実際の可逆層を選ばない偏極類を
許す標的は、引き続き有効な候補である。
\end{remark}

% !TEX root = ../nc-cubic-n-fold.tex
\section{一ブロック量子化の点スキームと二十四次の定常成分}
\label{sec:points}

この節では $k$ を標数零の代数閉体、$n\geq3$、$r=n-1$、
$V=\langle x,y,z,u_1,\ldots,u_r\rangle$ を生成元空間とする。
$F\in\Sym^3V$ と $h\in k^\times$ に対し、$u_\alpha$ を中心とし
\begin{equation}\label{eq:point-relations}
 [y,z]=h\operatorname{sym}(F_x),\qquad
 [z,x]=h\operatorname{sym}(F_y),\qquad
 [x,y]=h\operatorname{sym}(F_z)
\end{equation}
を課した二次代数を $B$ とする。
$g=F|_{u=0}$ は滑らかな平面三次式と仮定する。
さらに $B/(u)$ の双線型関係の零スキームが、滑らかな平面三次曲線
$E\subset\PP(\langle x,y,z\rangle^*)$ 上の固定点を持たない自己同型
$\sigma$ のグラフであると仮定する。
$w\in Z(B)_3$ の境界への像は、この $E$ の全点加群を消すものとし、
$A=B/(w)$ と置く。
これらは、適切な開にある楕円型一ブロック中心商の仮定である。
古典曲線 $V(g)$ と量子境界 $E$ は区別する。

一点の中心拡大における楕円曲線と八点、および点対応からの
二次関係の復元は Le Bruyn--Smith--Van den Bergh
\cite[Theorem~4.2.2, Proposition~4.2.5 and \S4.3]{LBSV} に先行する。
三様体稿 \cite[\S23]{Threefold} の二中心変数の議論を、ここでは
$r$ 個の中心変数へ拡張する。定常成分は $n>3$ で正次元となる。

\begin{theorem}[点関手の分解]\label{thm:point-nfold}
$w$ の tensor 代表の可換化を $w_{\mathrm{ab}}$ とし、
\begin{equation}\label{eq:point-CD}
 C=V(F_x,F_y,F_z),\qquad
 D=V(F_x,F_y,F_z,w_{\mathrm{ab}})\subset\PP(V^*)
\end{equation}
と置く。$D$ は代表の選択に依存しない。
ここで点加群は、基底 $S$ 上で次数零から生成され、各次数成分が可逆層であり、
$M_0\simeq\OO_S$ の同定を備えるものとする。同型はこの同定を保つ。
この正規化された点加群の関手は
\begin{equation}\label{eq:point-disjoint}
 \operatorname{Pt}(B)=E\amalg C,\qquad
 \operatorname{Pt}(A)=E\amalg D
\end{equation}
で表される。shift は $E$ 上で $\sigma$、定常成分上で恒等である。
$C$ は型 $(2,2,2)$、次元 $n-2$、次数 $8$ の完全交叉であり、
中心座標への射
\begin{equation}\label{eq:point-critical-cover}
 f:C\longrightarrow\PP^{r-1}_{u_1:\cdots:u_r}
\end{equation}
は次数 $8$ の有限平坦射である。
$w_{\mathrm{ab}}$ が $C$ の各既約成分上で非零ならば、
$D$ は型 $(2,2,2,3)$ の完全交叉であり、
\begin{equation}\label{eq:point-invariants}
 \dim D=n-3,\qquad \deg D=24,\qquad
 \omega_D\simeq\OO_D(7-n).
\end{equation}
\end{theorem}

\begin{proof}
二次関係の双線型化の零スキームを $\Gamma_B$ とし、その二座標を $p,q$ と書く。
中心関係は全ての一次生成元 $v_i$ に対し
\[
 u_\alpha(p)v_i(q)-v_i(p)u_\alpha(q)=0
\]
を与える。$u_\alpha(p)$ が可逆な chart では、第二因子の一座標を
可逆に取ってこの式を使えば $u_\alpha(q)$ も可逆である。
両者を一に規格化すると $p=q$ となり、残りの関係は
$F_x=F_y=F_z=0$ である。この計算は非被約な基底上でも成り立つ。

一方、$p,q$ の二座標の小行列式が可逆な chart では、中心関係の
二本から全ての $u_\alpha(p),u_\alpha(q)$ が零となる。
従ってこの部分は $\Gamma_\sigma$ に等しい。
幾何点で中心座標が全て零なら境界の対応上にあり、$\sigma$ が
固定点を持たないため、この小行列式 chart がその全体を覆う。
$C$ と平面 $V(u)$ は交わらない。交点があれば
$g_x=g_y=g_z=0$ となり、Euler の公式と $V(g)$ の滑らかさに反する。
従って二種類の chart は全体を覆い、scheme として
\[
 \Gamma_B=\Gamma_\sigma\amalg\Delta_C
\]
を得る。これは自己同型のグラフなので、二次関係を満たす無限点列は
初期点から一意に定まる。

中心三次関係を加えると、$E$ 上では仮定によって三重線型化が零である。
滑らかな $E$ 上の切断として零なので、任意の基底変更でも零である。
定常列上では三重線型化は $w_{\mathrm{ab}}(p)$ となる。
二つの tensor 代表の差の可換化は $(F_x,F_y,F_z)$ に属するので、
$D$ は代表に依存しない。これで点関手の主張を得た。

有限平坦性を環上で証明する。$S=k[u_1,\ldots,u_r]$ と置き、
\[
 T=S[x,y,z]/(F_x,F_y,F_z)
\]
を active 次数で濾過する。三つの最高記号 $g_x,g_y,g_z$ は
$k[x,y,z]$ の正則列である。従って、濾過 Koszul 複体の最高記号からの
帰納法により
\[
 \operatorname{gr}T\simeq
 S\otimes_k k[x,y,z]/(g_x,g_y,g_z)
\]
となる。具体的には、最高次数の関係を Koszul syzygy で消し、
active 次数を下げれば、初期イデアルに新しい生成元が生じないことが分かる。
右辺の Jacobian 環の Hilbert 級数は $(1+t)^3$ である。
その斉次基底を持ち上げると、$T$ は $S$ 上自由であり、全次数付き加群として
\begin{equation}\label{eq:critical-free-module}
 T\simeq S\oplus S(-1)^3\oplus S(-2)^3\oplus S(-3)
\end{equation}
を得る。この分解自体の標準性は主張しない。
$C\cap V(u)=\varnothing$ と合わせて Proj を取れば、$f$ は有限平坦で
\begin{equation}\label{eq:critical-pushforward}
 f_*\OO_C\simeq\OO\oplus\OO(-1)^3\oplus\OO(-2)^3\oplus\OO(-3)
\end{equation}
となる。

環 $T$ の次元は $r$ なので、三つの二次式の高さは三であり、
多項式環の Cohen--Macaulay 性から正則列となる。
従って $C$ の次元と次数が従う。
完全交叉は埋込み素点を持たないので、各既約成分で非零な
$w_{\mathrm{ab}}$ は非零因子である。
正則列に三次式を加えて次元と次数を得る。
随伴公式により
$\omega_D=\OO_D(2+2+2+3-(n+2))=\OO_D(7-n)$ である。
\end{proof}

\begin{corollary}[中心射影と二十四次の行列式]\label{cor:point-norm}
定理の完全交叉条件の下で、$\mathcal V=f_*\OO_C$ と置くと
\begin{equation}\label{eq:point-norm-sequence}
 0\longrightarrow\mathcal V(-3)
 \xrightarrow{\cdot w_{\mathrm{ab}}}\mathcal V
 \longrightarrow f_*\OO_D\longrightarrow0
\end{equation}
は完全である。この乗法行列の行列式は非零な
$\OO_{\PP^{r-1}}(24)$ の切断であり、$f_*\OO_D$ の
零次 Fitting イデアルを定める。
$n=4$ で $D$ が滑らかならば、$D$ は次数 $24$、種数 $37$ の連結曲線である。
\end{corollary}

\begin{proof}
中心座標が基点を持たないので $\OO_C(1)=f^*\OO(1)$ である。
$f$ は有限、$w_{\mathrm{ab}}$ は非零因子だから、中心三次切断の完全列を
押し出して表示した完全列を得る。階数 $8$ より行列式の次数は $3\cdot8=24$。
一般点で乗法は同型なので行列式は非零である。
これは scheme 論的像の被約方程式や、$D$ の射影先への双有理性を主張するものではない。
正次元の射影完全交叉では Koszul 解消と射影空間の線束コホモロジーから
$H^0(D,\OO_D)=k$ であり、従って連結である。
$n=4$ では $\omega_D=\OO_D(3)$ なので
$2g(D)-2=3\cdot24=72$ となる。
\end{proof}

\begin{remark}[滑らかな定常成分の非空性]\label{rem:point-nonempty}
三つの一般の二次式 $Q_i(u)$ に対して
\[
 F=(x^3+y^3+z^3)/3-xQ_1(u)-yQ_2(u)-zQ_3(u)+G(u)
\]
とすると、$C$ は $x^2=Q_1,y^2=Q_2,z^2=Q_3$ で与えられる。
$\PP^{r-1}$ 上の三つの分岐二次超曲面を滑らかで交わりが横断的になるよう選べば、
その二重被覆のファイバー積は滑らかである。
局所的には分岐因子の方程式を座標の一部に取り、その平方根を新座標にすればよい。
$r=2$ では分岐点を全て異なるものに取る。
純中心三次式を動かす線型系は $C$ 上基点を持たないため、
Bertini により一般の中心三次切断 $D$ は滑らかである。
これは $n=3$ では被約な二十四点を意味する。
境界の滑らかさ・固定点なしの条件と、後の Hessian 条件も、係数空間の
非空開条件として共通に課せる。ここで示すのは該当する開の非空性であり、
全ての分離型や全ての中心三次切断が滑らかだという主張ではない。
\end{remark}

\begin{theorem}[点対応からの包絡復元]\label{thm:point-envelope}
定理\ref{thm:point-nfold}の完全交叉条件を仮定し、
$R_B\subset V\otimes V$ を二次関係空間とする。
全点スキーム上の shift のグラフ
$\Gamma_A=\Gamma_\sigma\amalg\Delta_D$ について
\begin{equation}\label{eq:point-envelope}
 H^0(\mathcal I_{\Gamma_A}(1,1))=R_B
\end{equation}
が成立する。従って、埋込み点スキームと shift から二次包絡 $B$ と、
式\eqref{eq:point-relations}の尺度込みの双ベクトル $h\Pi_F$ を復元できる。
\end{theorem}

\begin{proof}
完全交叉のイデアルは飽和しており、三次生成元は二次部分に寄与しないから
$Q=I(D)_2=\langle F_x,F_y,F_z\rangle$ である。
従って $\Delta_D$ 上で消える双線型形式は
$\bigwedge^2V\oplus\operatorname{sym}(Q)$ で、次元は $\binom N2+3$ である。
$U=\langle u_1,\ldots,u_r\rangle$ を含む反対称部分は
次元 $\binom N2-3$ であり、境界上で零となる。
これで割った六次元の候補空間は
\[
 \bigwedge^2(V/U)\oplus
 \operatorname{sym}\langle g_x,g_y,g_z\rangle
 \subset(V/U)^{\otimes2}
\]
に単射する。

$L=\OO_E(1)$、$M=\sigma^*L$ とすると、乗法
$H^0(L)\otimes H^0(M)\to H^0(LM)$ は全射である。
$L\not\simeq M$ の場合は $L$ の基点のない pencil の完全列と
$H^1(ML^{-1})=0$、$L\simeq M$ の場合は平面三次に二次関係がないことから従う。
従って $\Gamma_\sigma$ の双線型イデアルは三次元であり、これは境界代数の
三つの関係で生成される。候補空間に含まれるこの三次元核を戻せば
全体の核の次元は $\binom N2$ である。
既に $R_B$ を含み、$R_B$ は反対称成分への射影が同型なので、両者は等しい。
関係を $v_i\otimes v_j-v_j\otimes v_i-\operatorname{sym}(q_{ij})$ と一意に書けば
$\sum q_{ij}\partial_{v_i}\wedge\partial_{v_j}$ が復元される。
本模型では定義から $h\Pi_F$ に等しい。
\end{proof}

\begin{proposition}[古典 Casimir 空間]\label{prop:point-casimir}
上の完全交叉条件に加えて、$g$ の Hessian に対する定数行列の条件
\begin{equation}\label{eq:point-hessian-scalar}
 M\operatorname{Hess}(g)\text{ が対称}\quad\Longrightarrow\quad M=cI_3
\end{equation}
を仮定する。このとき埋込み点スキームから
\[
 U=I(E)_1,\qquad Q=I(D)_2,
\]
および
\begin{equation}\label{eq:point-casimir}
 \mathcal R_Q=\{H\in\Sym^3V:\partial_aH\in Q\text{ for all }a\in\operatorname{ann}(U)\}
 =kF\oplus\Sym^3U
\end{equation}
を復元できる。従って $\dim\mathcal R_Q=1+\binom{n+1}{3}$ であり、
古典 Jacobian Poisson 直線 $[\Pi_F]$ が定まる。
\end{proposition}

\begin{proof}
$E$ はその境界平面を張る。
$n=3$ では正次元成分、$n>4$ では一次元成分として $E$ を識別できる。
$n=4$ では $D$ は連結な次数 $24$ の完全交叉曲線なので、
次数 $3$ の連結成分 $E$ と区別できる。
よって $U$ と $Q$ は埋込みデータから得られる。
$H\in\mathcal R_Q$ に対して
$\nabla_{x,y,z}H=M\nabla_{x,y,z}F$ となる定数行列 $M$ がある。
偏微分の可換性と $u=0$ への制限から式\eqref{eq:point-hessian-scalar}が適用でき、
$M=cI$ である。従って $H-cF$ は中心変数だけの三次式となる。
逆包含は直接従う。
$F$ を $cF+G(u)$ に替えることは Jacobian 括弧を $c$ 倍するだけなので、
中心空間と表示した三次空間は $[\Pi_F]$ を定める。

さらにこの空間は affine Poisson 括弧の三次 Casimir 全体でもある。
Casimir $H$ の active 勾配は分数体上で $a/b$ 倍の $\nabla F$ となる。
$F_x,F_y,F_z$ は正則列で共通因子を持たないので、既約表示の分母 $b$ は単元である。
斉次次数を比較して $a/b$ は定数となり、同じ結論を得る。
この最後の Casimir 計算自体には式\eqref{eq:point-hessian-scalar}を必要としない。
\end{proof}

\begin{remark}[明示的な非空条件と限界]
$g=(x^3+y^3+z^3)/3+\lambda xyz$ では
$\lambda(\lambda^3-8)\ne0$ のとき式\eqref{eq:point-hessian-scalar}が成り立つ。
実際、$M\operatorname{Hess}(g)$ の対称性は対角成分を等しくし、非対角成分には
$\lambda m_{13}=2m_{21}$、$\lambda m_{21}=2m_{32}$、$\lambda m_{32}=2m_{13}$
ともう一組の循環式を課すので、全非対角成分が零になる。
滑らかさの条件 $\lambda^3\ne-1$ と合わせても非空である。
この条件による勾配復元の原理と、本稿での任意 $r$ の点関手への適用を区別する。
また $\mathcal R_Q/\Sym^3U$ の直線を得ることは、その中の特定の三次式を
選ぶことではない。量子中心から残余圏が定める可換三次式への補正は別問題であり、
点スキームを有限係数の Poisson 零スキームと同一視してはいない。
\end{remark}

% !TEX root = ../nc-cubic-n-fold.tex
\section{相対点関手と構造対象付きモジュライへの降下}
\label{sec:relative-points}

\subsection{相対点関手と任意の基底変更}
\label{subsec:relative-points}

点スキームを族にする際には、中心三次の生成元ではなくその直線を用いる。
また、各ファイバーでの完全交叉性と、非平坦な基底変更後にも保たれる
相対完全交叉性を区別する。
以下の $S$ は良い一ブロック表示チャートの任意の基底変更とし、
$\mathcal V$ は階数 $N=n+2$ の生成元束、
$\mathcal U\subset\mathcal V$ は階数 $r=n-1$ の中心部分束、
$\mathcal X=\mathcal V/\mathcal U$ とする。
$\mathbb P(\mathcal V^*)$ は $\mathcal V$ の可逆商を表す規約とし、
その標準全射を $\mathcal V\otimes\OO\to\OO(1)$ と書く。

局所的に $\mathcal V=\langle x,y,z,u_1,\ldots,u_r\rangle$ と表示し、
式\eqref{eq:point-relations}の関係を用いる。
$g=F|_{u=0}$ は各幾何学的ファイバーで滑らかな平面三次式とする。
境界の双線型関係スキームは、相対的に滑らかな平面三次曲線
$\mathcal E\to S$ の固定点を持たない $S$-自己同型
$\sigma$ のグラフであると仮定する。
中心三次直線を $\ell\subset Z(\mathcal B)_3$ とし、
その境界への像は $\mathcal E$ 上の普遍点加群を消すものとする。
この最後の条件は切断の恒等式であり、幾何点での消滅だけを仮定するのではない。
$\mathcal A=\mathcal B/(\ell)$ と置く。

三つの偏微分が張る階数三の部分束を
$\mathcal Q\subset\operatorname{Sym}^2\mathcal V$ とする。
中心三次の可換化は、代表に依存しない写像
\[
 \ell\longrightarrow
 \bigl(\operatorname{Sym}\mathcal V/(\mathcal Q)\bigr)_3
\]
を与える。
これにより
\[
 \mathcal C=V(\mathcal Q),\qquad
 \mathcal D=V(\mathcal Q,\ell_{\mathrm{ab}})
 \subset\mathbb P(\mathcal V^*)
\]
を定める。各幾何学的ファイバーの $\mathcal C$ 上で
$\ell_{\mathrm{ab}}$ が非零因子を生成する開を、点スキームの良い開と呼ぶ。
これは以下の行列式によって検出される開条件である。

\begin{theorem}[普遍的な相対点スキーム]
\label{thm:relative-point-nfold}
以上の点スキームの良い開上では、任意の $S$-スキーム $S'$ に対して
正規化された相対点加群の関手が
\[
 \operatorname{Pt}_{S'}(\mathcal B_{S'})
   =\mathcal E_{S'}\amalg\mathcal C_{S'},\qquad
 \operatorname{Pt}_{S'}(\mathcal A_{S'})
   =\mathcal E_{S'}\amalg\mathcal D_{S'}
\]
で表される。ここで各次数成分は $S'$ 上の可逆層で、次数零を
$\OO_{S'}$ と同定し、次数零から生成されることを要求する。
shift は $\mathcal E$ 上で $\sigma$、$\mathcal C,\mathcal D$ 上で恒等である。

中心射影
\[
 f:\mathcal C\longrightarrow\mathbb P(\mathcal U^*)
\]
は次数八の有限局所自由射である。
$\mathcal C\to S$ と $\mathcal D\to S$ は平坦であり、
それぞれ型 $(2,2,2)$ と $(2,2,2,3)$ の相対完全交叉である。
これらの主張、正則埋込み、および点関手の同定は、任意の基底変更と両立する。
各ファイバーの次元と次数は、それぞれ $(n-2,8)$ と $(n-3,24)$ である。

$p_D:\mathcal D\to S$ とすると、相対双対化層は
\begin{equation}\label{eq:relative-point-dualizing}
 \omega_{\mathcal D/S}\simeq
 \OO_{\mathcal D}(7-n)\otimes
 p_D^*\bigl(\det\mathcal V\otimes
              (\det\mathcal Q)^{-1}\otimes\ell^{-1}\bigr).
\end{equation}
従って基底上の determinant 因子を自明化していない族では、
右辺を単に $\OO_{\mathcal D}(7-n)$ と書かない。
\end{theorem}

\begin{proof}
まず有限局所自由性を係数環上で確認する。
局所的に $S=\operatorname{Spec}R$ として
\[
 K=R[x,y,z]/(g_x,g_y,g_z),\qquad K_i=K_i^{\mathrm{gr}}
\]
と置く。三つの偏微分の Koszul 複体の各次数部分は有限自由複体である。
各幾何学的ファイバーでは正次数のホモロジーが零であり、
$K_i$ の階数は $i=0,1,2,3$ に対して $1,3,3,1$、それ以外では零である。
有限自由複体で可逆な小行列式を順に分離すると、
各次数の複体は局所的に可縮な複体と次数零に置いた自由加群の直和になる。
従って $K_i$ は表示した階数の有限局所自由加群であり、
この Koszul 完全性と $K_i$ の形成は任意の係数変更と両立する。
これは被約基底上の階数計算だけを非被約基底へ移す議論ではない。

$R_u=R[u_1,\ldots,u_r]$ および
\[
 T=R_u[x,y,z]/(F_x,F_y,F_z)
\]
を考える。active 次数の最高記号は $g_x,g_y,g_z$ である。
最高記号の Koszul 関係で最高次数を消す帰納法により
\[
 \operatorname{gr}_{\mathrm{act}}T\simeq R_u\otimes_R K
\]
を得る。この帰納法は全ての係数変更後にも同じである。
従って $T$ は長さ四の有限濾過を持ち、その次数付き随伴商は
$R_u(-i)\otimes_R K_i$ $(0\leq i\leq3)$ である。
各 $K_i$ が局所自由なので、局所的に斉次持上げを選べば
\begin{equation}\label{eq:relative-critical-free}
 T\simeq R_u\oplus R_u(-1)^3\oplus
                 R_u(-2)^3\oplus R_u(-3)
\end{equation}
となる。この分解の標準性は主張しない。
とくに $T$ は $R_u$ 上階数八で有限局所自由である。

滑らかな $g$ の偏微分は射影平面上に共通零点を持たないから、
$\mathcal C$ は $V(\mathcal U)$ と交わらない。
よって中心射影は全域で定義され、\eqref{eq:relative-critical-free} の Proj により
有限局所自由性を得る。
$\mathcal C$ の三つの方程式についても、上の Koszul 複体は完全であり、
その商は $S$ 上平坦である。
射影空間は $S$ 上平坦かつ有限表示なので、これは相対正則埋込みであり、
任意の基底変更後にも正則である
\cite[Tags~063P, 063U]{NFMStacksRelativeRegular}。

三次直線の像は各ファイバー上の非零因子である。
平坦かつ有限表示な $\mathcal C\to S$ 上の局所主イデアルに対する
ファイバー判定により、$\mathcal D\subset\mathcal C$ は相対有効 Cartier 因子である
\cite[Tag~062Y]{NFMStacksRelativeCartier}。
従って $\mathcal D\to S$ は平坦で、この Cartier 性は
非平坦なものを含む全ての基底変更で保たれる
\cite[Tag~056Q]{NFMStacksRelativeCartier}。
完全交叉型とファイバーの次元・次数が従う。

点関手については、定理\ref{thm:point-nfold}の二種類の chart の議論を
そのまま係数環上で用いる。
中心座標の一つが可逆なら中心関係から二つの点は等しく、
残りの式は $\mathcal Q=0$ となる。
二点の二次小行列式の一つが可逆なら、全中心座標が零となって
境界の $\sigma$ のグラフを得る。
固定点なしの仮定と $\mathcal C\cap V(\mathcal U)=\varnothing$ により
これらは全体を覆う。従って scheme として
\[
 \Gamma_{\mathcal B}=\Gamma_\sigma\amalg\Delta_{\mathcal C}
\]
である。この対応は自己同型のグラフなので、無限点列は初期点から一意に決まる。
三次直線の多重線型化は境界上で恒等的に零であり、
定常列上では $\ell_{\mathrm{ab}}$ の評価となる。
これで任意の基底上の点関手と shift の主張を得る。

最後に、射影空間の相対標準束は
$\OO(-N)\otimes\det\mathcal V$ である。
二次三本と三次直線の conormal determinant は
$\OO(-9)\otimes\det\mathcal Q\otimes\ell$ なので、
相対随伴公式から\eqref{eq:relative-point-dualizing}を得る。
\end{proof}

\begin{corollary}[相対 norm と良い点スキーム開]
\label{cor:relative-point-norm}
$p_U:\mathbb P(\mathcal U^*)\to S$、
$\mathcal H=f_*\OO_{\mathcal C}$ と置く。
点スキームの良い開上には、任意の基底変更と両立する完全列
\[
 0\longrightarrow p_U^*\ell\otimes\mathcal H(-3)
 \longrightarrow\mathcal H
 \longrightarrow f_*\OO_{\mathcal D}\longrightarrow0
\]
がある。最初の写像の行列式は
\[
 \mathcal N\in
 H^0\bigl(\mathbb P(\mathcal U^*),
       \OO(24)\otimes p_U^*\ell^{-8}\bigr)
\]
であり、$f_*\OO_{\mathcal D}$ の零次 Fitting イデアルを定める。
各ファイバーで $\mathcal N$ が零切断でないという条件は基底の開条件であり、
上の非零因子条件と同値である。
良い一ブロック表示族ではこの開は非空である。
\end{corollary}

\begin{proof}
中心射影から $\OO_{\mathcal C}(1)=f^*\OO(1)$ である。
相対 Cartier 因子の完全列を有限射で押し出せば表示した完全列を得る。
$\mathcal H$ の階数は八なので、その determinant は
$\OO(24)\otimes p_U^*\ell^{-8}$ に値を持つ。
各ファイバーでは $T$ は多項式環 $k[u]$ 上自由であり、
その上の三次乗法の単射性は行列式が非零多項式であることと同値である。
$T$ は $k[u]$ 上自由なので、三次乗法の核に射影化で消える
非零部分加群はない。従ってこれは $\mathcal C$ 上の非零因子条件とも同値である。
行列式の全係数が同時に零になる条件は閉条件なので、所定の条件は開である。

非空性のため、一つの良い周囲代数とその中心三次の持上げを固定する。
純中心三次を加える線型系は $\mathcal C$ 上基点を持たない。
有限個の既約成分のいずれかで恒等的に零となる条件はそれぞれ真の affine
線型条件であるから、その和集合を避けることができる。
これは一ブロック表示成分の非空開についての主張である。
\end{proof}

\begin{proposition}[相対点対応からの二次包絡]
\label{prop:relative-point-envelope}
$\rho=\sigma\amalg\operatorname{id}_{\mathcal D}$ とし、
$\Gamma=\Gamma_\rho\subset
\mathbb P(\mathcal V^*)\times_S\mathbb P(\mathcal V^*)$ と置く。
$p$ をこの積から $S$ への射とすると
\begin{equation}\label{eq:relative-point-envelope}
 p_*\mathcal I_\Gamma(1,1)=\mathcal R_{\mathcal B}
 \subset\mathcal V\otimes\mathcal V
\end{equation}
であり、左辺の形成は任意の基底変更と両立する。
また、埋込みイデアルの低次数部分は任意の基底変更と両立して
\begin{equation}\label{eq:relative-point-ideal-pieces}
 p_*\mathcal I_{\mathcal D}(2)=\mathcal Q,
 \qquad
 p_*\mathcal I_{\mathcal C}(3)=\mathcal Q\otimes\mathcal V,
 \qquad
 p_*\mathcal I_{\mathcal D}(3)\big/
       p_*\mathcal I_{\mathcal C}(3)\simeq\ell
\end{equation}
となる。ここで後二式の $p$ は射影空間から $S$ への射を表す。
\end{proposition}

\begin{proof}
型 $(2,2,2,3)$ の Koszul 解消に $\OO(2)$ または $\OO(3)$ を tensor する。
射影空間の相対次元は $n+1\geq4$、イデアル層のこの解消の長さは三なので、
射影空間の最高次コホモロジーは全次数零の計算に寄与しない。
中間コホモロジーも零であり、残る次数零の列から
\eqref{eq:relative-point-ideal-pieces}を得る。
二次生成元の一次係数による syzygy はないので、
$\mathcal Q\otimes\mathcal V\to\operatorname{Sym}^3\mathcal V$ は単射である。
三次生成元の tensor 代表を替えた差はその像に入るため、最後の同型は代表に依存しない。
この相対 Koszul 計算は全ての基底変更後にも同じである。
とくに $n=3$ で $\mathcal D$ が零次元でも成立するが、
$\operatorname{Sym}^3\mathcal V\to p_*\OO_{\mathcal D}(3)$ の全射性は主張しない。

従って対角成分上で消える双線型形式の束は
$\bigwedge^2\mathcal V\oplus\operatorname{sym}(\mathcal Q)$ である。
その中の $\mathcal U\wedge\mathcal V$ を割ると、境界での候補は
\[
 \bigwedge^2\mathcal X\oplus
 \operatorname{sym}\langle g_x,g_y,g_z\rangle
 \subset\mathcal X^{\otimes2}
\]
となる。境界の $L=\OO_{\mathcal E}(1)$ に対して
$L\otimes\sigma^*L$ は相対次数六であるから、その切断束は階数六で、
任意の基底変更と両立する。
乗法 $p_*L\otimes p_*\sigma^*L\to p_*(L\otimes\sigma^*L)$ は
各幾何学的ファイバーで全射であることを
定理\ref{thm:point-envelope}の pencil の議論で確認できる。
従って相対的にも分裂全射であり、三次元の核は境界の二次関係束そのものである。
これを上の候補束に引き戻すと、核は $\mathcal R_{\mathcal B}$ に等しい。
全てが部分束と分裂全射で計算されるので、この等式も任意の基底変更で保たれる。
\end{proof}

\subsection{$\mathcal O$-marked モジュライへの点データの降下}
\label{subsec:point-descent}

上の条件を満たす $P$-不変な非空開を $J^{\mathrm{pt}}\subset J^{\mathrm{good}}$ とし、
\[
 \mathfrak M_n^{\mathrm{pt}}
  =[J^{\mathrm{pt}}/\overline P]
  \subset\mathfrak M_n^{\mathrm{loc}}
\]
と置く。最後の同定には定理\ref{thm:nfm-local-moduli}を用いる。
これは良い一ブロック表示に fpqc 局所的に同値な圏族の開であり、
全ての内在的成分の一般点を覆うとは主張しない。

\begin{theorem}[点関手、shift、中心射影の $\mathcal O$-marked 降下]
\label{thm:point-omarked-descent}
$\mathfrak M_n^{\mathrm{pt}}$ 上には Brauer--Severi 束
\[
 \mathfrak P_V=[\mathbb P(\mathcal V^*)/\overline P],
 \qquad
 \mathfrak P_U=[\mathbb P(\mathcal U^*)/\overline P]
\]
と、それぞれに付随するデータ
\[
 \mathfrak Y=\mathfrak E\amalg\mathfrak D\hookrightarrow\mathfrak P_V,
 \qquad \rho:\mathfrak Y\xrightarrow{\sim}\mathfrak Y,
 \qquad f:\mathfrak C\longrightarrow\mathfrak P_U
\]
がある。$f$ は次数八で有限局所自由であり、
$\mathfrak C\hookrightarrow\mathfrak P_V$ と
$\mathfrak D\hookrightarrow\mathfrak P_V$ はそれぞれ fpqc 局所的に
型 $(2,2,2)$ と $(2,2,2,3)$ の相対完全交叉である。
これらは全ての $\mathcal O$-保存同値に関して自然であり、基底変更と両立する。

さらに $\mathfrak Y$ は次の標識を増やさない点関手を表す。
対象 $(\mathscr C,\mathcal O)$ に対し、基底 $T$ 上の点とは全射
\[
 q:\mathcal O_T\twoheadrightarrow\mathcal P
\]
であって、fpqc 局所的に良い表示
$(\mathscr C_T,\mathcal O_T)\simeq
 (\operatorname{QGr}\mathcal A_T,\pi\mathcal A_T)$ の下で、
正規化された点加群 $\mathcal A_T\twoheadrightarrow M$ の tails に同型となるものとする。
射は $q$ を保つ同型とする。
この条件は良い表示の選択によらず、当該 groupoid は自己同型を持たない。
従って $\mathfrak Y$ 上にはこの意味での普遍点商が存在する。
次数一 shift の大域的な選択を点関手の定義に加える必要はない。
\end{theorem}

\begin{proof}
次数付き同型は生成元の射影空間、全関係、点加群、およびその正規化 shift を移す。
次数スカラー $\lambda$ は射影空間上で恒等であり、点加群の次数 $j$ には
$\lambda^j$ を掛ける標準同型を与える。
次数零の同定は保たれるので、この作用は正規化点加群の同型類上で自明である。
従って幾何データは $P$ から $\overline P$ へ降下する。
定理\ref{thm:nfm-fixed-model-comparison}により、全ての
$\mathcal O$-保存同値は局所的に次数付き同型から来て、
二つの持上げは次数スカラーだけ異なる。
よってこの降下は表示同型だけでなく全ての $\mathcal O$-保存同値に関して自然である。
有限局所自由性、平坦性、および相対正則埋込みも fpqc 降下する。

点関手と tails の関係を確認する。
点加群は各次数から上へ生成されるので非零 torsion 部分加群を持たない。
二つの点加群の tails が同型なら、その同型と逆は十分大きい次数で互いに逆な
次数付き写像により表せる。これは点加群が次数零の一つの可逆層で生成され、
各打切り $M_{\geq m}$ も一つの有限局所自由次数成分から生成されることによる。
従って十分大きい次数以後の点列が等しい。
点対応 $\rho$ が自己同型なので、その初期点も等しい。
相対点関手の表示から、元の正規化点加群は同型である。
また $q$ は全射なので、$q$ を保つ自己同型は恒等だけである。
これにより点加群と上述の局所点商の同型類を同定できる。

全 $\mathcal O$-同値の局所持上げによって、局所点商という条件は
良い表示から独立である。
圏族と対象の有効 fpqc 降下により、この点商は貼り合わさる。
正規化 shift は局所的に
\[
 M\longmapsto M(1)_{\geq0}\otimes M_1^{-1}
\]
であり、初期点を $\rho$ で移す。
次数スカラーに関してこの式の次数 $j$ の重みは再び $j$ になるので、
これも点商の関手とそのグラフとして降下する。
ここでは圏全体の次数一 shift を大域に選んだわけではない。
\end{proof}

\begin{proposition}[偏極、双対化層、包絡の正確な降下形]
\label{prop:point-polarization-descent}
命題\ref{prop:nfm-presentation-gerbe}の表示 gerbe
\[
 \mathfrak G=[J^{\mathrm{pt}}/P]
   \longrightarrow\mathfrak M_n^{\mathrm{pt}}
\]
の類を $\beta$ とし、$e=\gcd(r,3)$ とする。
整数 $a,b$ を $ar+3b=e$ となるように一組固定し、
$\chi=(\det\mathcal U)^a\otimes(\det\mathcal X)^b$ と置く。
gerbe 上の線束
\begin{equation}\label{eq:point-descended-polarization}
 \OO_{\mathbb P(\mathcal V^*)}(e)\otimes p_V^*\chi^{-1},
 \qquad
 \OO_{\mathbb P(\mathcal U^*)}(e)\otimes p_U^*\chi^{-1}
\end{equation}
はそれぞれ $\mathfrak P_V,\mathfrak P_U$ 上の相対 ample 線束に降下する。
中心射影によるそれらの引戻しは $\mathfrak C$ 上で一致する。
別の整数解を用いた降下線束との差は、モジュライ基底上の線束の引戻しである。

特に ambient Brauer--Severi 束上で、
$n\not\equiv1\pmod3$ なら幾何学的ファイバーで
$\OO(1)$ となる偏極を選べる。
$n\equiv1\pmod3$ では幾何学的に $\OO(3)$ となる偏極は常に選べるが、
$\OO(1)$ の存在は $\beta$ の消滅と同値である。
素の標準 $P$-線型化を持つ $\OO(1)$ は、いずれの場合にも重み一なので
そのまま $\overline P$ へ降りるわけではない。

一方、$\omega_{\mathcal D/S}$、$f_*\OO_{\mathcal C}$、
相対 norm の切断とその Fitting 因子、および
\eqref{eq:relative-point-envelope}が定める二次包絡の
次数スカラーを除いた同型類は、追加の偏極を選ばずに降下する。
通常の生成元束 $\mathcal V$ や二次関係束 $\mathcal R_{\mathcal B}$ 自体の
降下を主張するものではない。
\end{proposition}

\begin{proof}
gerbe の中心スカラーの重みは、
$\mathcal V,\mathcal U,\mathcal X$ 上で一、
$\OO(1)$ 上で一、$\mathcal Q$ 上で二、$\ell$ 上で三である。
$\chi$ の重みは $e$ なので、\eqref{eq:point-descended-polarization} の重みは零となる。
重み零の準連接層の fpqc 降下を用いれば表示した線束を得る。
その相対 ample 性は局所的な $\OO(e)$ の ample 性から従う。
二つの $\chi$ の比は重み零の基底線束なので、選択の差についての主張も従う。
$\OO(1)$ 型線束の存在と Brauer 類の消滅の同値、および
$r\beta=3\beta=0$ は命題\ref{prop:nfm-presentation-gerbe}による。

式\eqref{eq:relative-point-dualizing}の基底因子の重みは $N-9=n-7$ であり、
$\OO(7-n)$ の重みと相殺する。
norm の値を取る線束の重みも $24-8\cdot3=0$ である。
$\mathcal C$ とその構造層には次数スカラーが自明に作用するので、
$f_*\OO_{\mathcal C}$ は通常の有限局所自由層として降下する。

包絡の関係束 $\mathcal R_{\mathcal B}$ は重み二である。
従って一般には、これをモジュライ基底上の通常の部分束と記すことはできない。
しかし埋込みグラフのイデアルそのものは降下し、各表示上での
\eqref{eq:relative-point-envelope}による復元は射影的座標変更と両立する。
よって二次包絡は次数スカラーを除いた同型問題として復元される。
反対称部分への射影を用いて関係から復元する
$\bigwedge^2\mathcal V\to\operatorname{Sym}^2\mathcal V$ は重み零であるため、
対応する尺度込みの相対 Poisson 双ベクトルも ambient Brauer--Severi 束上へ降下する。
この結論だけから中心三次に対応する古典 Casimir を選択したことにはならない。
\end{proof}

\begin{remark}[普遍点加群と普遍点商]
普遍点加群の次数 $j$ の可逆層は gerbe 上で重み $j$ を持つ。
したがって、普遍点商が $\mathcal O$-marked な圏族に降下するという主張と、
全次数成分が通常の線束として降下するという主張は異なる。
前者は上で証明したが、素の線型化を持つ次数成分は通常の線束として降下しない。
モジュライ基底上に通常の次数付き代数族を置き、その点加群として記述するには、
表示 gerbe の中立化を選ぶ必要がある。
これは点スキーム自身への引戻しで障害が消えることとは矛盾しない。
実際、点スキームへの引戻し gerbe 上では普遍点加群の $M_1$ が重み一の線束であり、
これによってその gerbe は中立化される。
そこで次数 $j$ ごとに $\mathcal A_j$ と $M_j$ の双方を
$M_1^{-j}$ で捻れば、点空間上では通常の次数付き代数とその普遍点加群を得る。
またここでの $\beta\in H^2(\mathbf G_m)$ は対象の表示障害であり、
命題\ref{prop:nfm-lifting-line}の一つの同値の持上げを妨げる
$\operatorname{Pic}$ 上の線束とは別の障害である。
\end{remark}

% !TEX root = ../nc-cubic-n-fold.tex
\section{点対応による中心復元と Poisson 三次対への射}
\label{sec:poisson-morphism}

\subsection{中心三次の評価と境界 Jacobi 線}

この節では一ブロックの良い点スキーム開をさらに縮める。
全ての主張は $n\ge3$、$r=n-1\ge2$ に一様である。
次数三中心の束を $\mathcal E$、純中心三次の部分束を
$\mathcal W=\Sym^3U$、商線束を $\mathcal L=\mathcal E/\mathcal W$ とする。
固定ファイバーでは
\[
 S=k[U],\qquad T=\Sym(V)/(F_x,F_y,F_z),\qquad
 J_g=T/(U)=k[x,y,z]/(g_x,g_y,g_z)
\]
と書く。$T$ は $B$ の最大可換\emph{代数}商である。
これは交換子の張る線形部分空間による商 $\mathrm{HH}_0(B)$ とは異なる。
代数の可換化は中心上の評価写像
\[
 \alpha:\mathcal E\longrightarrow T_3
\]
を与える。tensor 代表を可換化し、三つの偏微分のイデアルで割る表示と同じ写像である。
純中心三次は $J_g$ で零となるので、合成は
\begin{equation}\label{eq:point-boundary-jacobi-map}
 \epsilon:\mathcal L=\mathcal E/\mathcal W\longrightarrow (J_g)_3
\end{equation}
を定める。右辺は階数一である。
$C\cap V(U)=\varnothing$ であるから、この商を $C$ と境界平面との
scheme 理論的交叉への制限と解釈してはならない。
これは次数付き切断環の Artin 商であり、
$T_3=H^0(C,\mathcal O_C(3))$ を用いて
\[
 (J_g)_3=T_3/(U\cdot T_2)
\]
と求めるものである。

\begin{proposition}[境界 Jacobi 線と中心評価]\label{prop:point-jacobi-evaluation}
$\epsilon$ が同型である条件は、表示変更群 $P$ に不変な非空開条件である。
この開上で $\alpha:\mathcal E\to T_3$ は余核が局所自由な単射であり、
任意の基底変更後にも単射である。
ここで分裂は局所的な存在の主張であり、標準的な大域分裂は選ばない。
\end{proposition}

\begin{proof}
$g$ の滑らかさより三つの偏微分は正則列であり、
$J_g$ の Hilbert 級数は $(1+t)^3$ である。
従って $(J_g)_3$ は線束となり、\eqref{eq:point-boundary-jacobi-map}
が同型であることは線束間の写像の非消滅条件である。
全ての写像は代数 $B$ とその中心一次空間 $U$ から定まるので $P$-共変である。

非空性を、公刊された Etingof--Ginzburg の Rains 中心公式
\cite[\S9, p.~1414]{NFMEG} で確認する。
\[
 g=(x^3+y^3+z^3)/3+\lambda xyz,\qquad 1+\lambda^3\ne0,
 \qquad \delta=1-h\lambda/2\ne0
\]
とし、$q=(1+h\lambda/2)/\delta$、$t=h/\delta$ と置く。
境界の関係は $xy-qyx-tz^2=0$ およびその巡回形となる。
同文献の式で全ての低次係数を零とした中心元は
\begin{equation}\label{eq:point-rains-homogeneous}
 \Psi=t(q+1)\bigl[
 t(t^3+1)y^3+(q^3-t^3)yzx-q(t^3+1)zyx+t(q^3-t^3)z^3
 \bigr].
\end{equation}
$J_g$ では $x^3=y^3=z^3=-\lambda xyz$ であり、
\begin{equation}\label{eq:point-rains-socle}
 [\Psi_{\mathrm{ab}}]
 =-\frac{4h^4(1+\lambda^3)}{\delta^6}[xyz]\ne0.
\end{equation}
この計算では $[xyz]\ne0$ も必要である。
例えば $x^2y=-\lambda y^2z=\lambda^2xz^2=-\lambda^3x^2y$
から、次数三の二種類の変数を含む単項式は全て零である。
立方は上の式で $xyz$ の倍数になり、$(J_g)_3$ は一次元なので $xyz$ は非零である。
係数恒等式の直接の証明書は
\[
 \Psi_{\mathrm{ab}}-\kappa xyz
 =t^2(q+1)\bigl((t^3+1)y g_y+(q^3-t^3)z g_z\bigr),\qquad
 \kappa=-4h^4(1+\lambda^3)/\delta^6
\]
である。
良い Sklyanin 開では中心三次は一次元であるから、その任意の非零生成元にも
同じ非消滅が成立する。混合中心係数を加えても境界への像は変わらない。
従来の PBW・滑らかな楕円境界・必要な有限範囲の位数除外の各非空開と交わらせられる。
全ての有限位数を除く場合には、これとは別に very general 条件として扱う。

次に $\mathcal W\to T_3$ は単射で、局所的には直和因子である。
実際、active 次数の濾過と $J_g$ の斉次基底の持上げにより
$T\simeq S\oplus S(-1)^3\oplus S(-2)^3\oplus S(-3)$ となり、
最初の項に $\mathcal W=S_3$ が含まれる。
商 $T_3/\mathcal W$ への $\mathcal L$ の写像は、さらに $(J_g)_3$ へ写したとき
同型である。従ってこれは局所的な直和因子であり、
$\mathcal E\to T_3$ も単射、余核も局所自由である。
この議論は有限局所自由な加群と単元を用いるため、任意の基底変更後にも成立する。
\end{proof}

\begin{theorem}[点対応から中心商までの復元]\label{thm:point-central-recovery}
上の開と点スキームの完全交叉条件を仮定する。
埋込み点スキームと shift のグラフから、二次包絡 $B$ に加えて
中心三次直線 $[w]\subset Z(B)_3$、従って $A=B/(w)$ を復元できる。
復元は表示変更と両立する。
相対的には定理\ref{thm:relative-point-nfold}の良い点スキーム開に
$\epsilon$ の同型条件を加えれば、中心直線の復元も任意の基底変更と両立する。
\end{theorem}

\begin{proof}
定理\ref{thm:point-envelope}で $B$ を復元した後、$U=Z(B)_1$、
$Q=\langle F_x,F_y,F_z\rangle$、$C=V(Q)$ と $\alpha$ は全て定まる。
固定点部分は $D$ である。完全交叉イデアルの次数三部分から
\[
 I(D)_3/I(C)_3=k[w_{\mathrm{ab}}]\subset T_3
\]
を得る。命題\ref{prop:point-jacobi-evaluation}の単射 $\alpha$ による
この直線の逆像は、ただ一つの直線 $kw\subset\mathcal E$ である。
従って二次関係にこの三次関係を加えることで $A$ が決まる。

$n=3$ で $D$ が零次元でも、型 $(2,2,2,3)$ の斉次完全交叉イデアルは
飽和している。この場合に $T_{D,3}=H^0(D,\mathcal O_D(3))$ という
誤った全射性を使う必要はなく、多項式からの制限写像の核を取ればよい。
相対版では命題\ref{prop:relative-point-envelope}および
式\eqref{eq:relative-point-ideal-pieces}により、二次包絡と
$p_*\mathcal I_D(3)/(Q\cdot V)$ の直線部分束が任意の基底変更と両立して定まる。
中心評価の余核は局所自由なので、その逆像の形成も基底変更と両立する。
\end{proof}

\begin{remark}[中心変数の個数と復元]
中心三次の復元では $r\ge2$ が本質的である。
比較のため $r=1$ とし、$C$ が八つの被約点である周囲代数を固定する。
純中心三次を加えるパラメータのうち有限個を除けば、
中心三次の可換化は $C$ の全ての点で非零となるので $D=\varnothing$ である。
従って中心商の点対応は同じ $\Gamma_\sigma$ だけになり、中心直線を復元できない。
一点中心拡大の周囲二次関係を復元する
\cite[Proposition~4.2.5]{LBSV}と、本稿の $r\ge2$ における中心商の復元を区別する。
\end{remark}

\begin{proposition}[定常三次切断と古典 Casimir の相違]
\label{prop:point-casimir-restriction-obstruction}
上の開で、古典 Casimir 空間を $\mathcal R=kF\oplus\Sym^3U$ とする。
その $T_3$ への制限像は $(J_g)_3$ への写像の核に入る。
一方、$w\notin\Sym^3U$ ならば $[w_{\mathrm{ab}}]$ はこの核に入らない。
特に、
\[
 \alpha(Z(B)_3)\cap \operatorname{im}(\mathcal R\longrightarrow T_3)
 =\Sym^3U.
\]
従って $w_{\mathrm{ab}}|_C$ を古典 Casimir の制限と同一視する構成、
および純中心三次だけを加えてその同一視を実現する構成は成立しない。
$D$ が上記の Cartier 切断であるとき、古典 Casimir $H$ の切断により
$D=C\cap V(H)$ を scheme として実現することもできない。
\end{proposition}

\begin{proof}
$F$ の $J_g$ への像は $g$ の類であり、
$3g=xg_x+yg_y+zg_z$ なので零である。
純中心三次の像も零である。
一方、中心三次の像は同型 $\epsilon$ を経て零になるのが純中心三次の場合に限る。
さらに $S_3\subset T_3$ は単射なので、交叉についての等式も従う。

幾何学的な切断としての主張のために、$T_3=H^0(C,\mathcal O_C(3))$ を確認する。
$T$ は次元 $r\ge2$ の Cohen--Macaulay な斉次完全交叉環なので、
irrelevant ideal に関する局所コホモロジーの次数零と一次は消える。
切断環との標準完全列から、全ての次数で $T_j=H^0(C,\mathcal O_C(j))$ となる。
また $H^0(C,\mathcal O_C)=k$ である。
二つの非零因子切断が同じ Cartier 零スキームを定めるなら、
その比は $C$ の大域可逆関数なので $k^\times$ に属する。
従って同じ $D$ を与える古典 Casimir が存在すれば、
その制限は $w_{\mathrm{ab}}|_C$ の非零定数倍となり、既に示した障害に反する。
\end{proof}

\begin{remark}[古典制限写像そのものの一般単射性]
上の障害は古典 Casimir の制限写像の核が大きいことによるものではない。
$\mathcal R\to T_3$ 自体は非空開上で単射となる。
実際、$F=g-xu_1^2$ と取れば、Euler の式から $[F]=-(2/3)xu_1^2$ であり、
active 次数の濾過によりこれは純中心三次の像に属さない。
従って $kF\oplus S_3$ の制限は単射であり、単射性は係数空間の開条件である。
三変数制限の Hessian 剛性条件や良い点スキーム条件とも交わらせられる。
一方、$F=g+G(u)$ では $[F]=G(u)$ となるから、この単射性は全係数点の主張ではない。
この相違は補正を用いる Poisson モジュライへの射の非存在を意味しない。
次の trace と Hessian による構成は、この相違を保ったまま中心三次の両空間を同定する。
\end{remark}

\subsection{Trace と Hessian による Poisson 三次対への射}

\begin{lemma}[任意基底上の三次 Poisson 中心]\label{lem:point-relative-casimirs}
定理\ref{thm:relative-point-nfold}の良い周囲族を考える。
任意の、非被約・非平坦なものも含む係数環への基底変更後に、
三次 Poisson Casimir の加群は
\[
 Z_{\mathrm{Pois}}(\Sym(V))_3=R_0F\oplus\Sym^3U
\]
となる。ここで局所表示の係数環を $R_0$ と書いた。
従ってこれは普遍的な局所自由加群であり、その形成は任意の基底変更と両立する。
Hessian のスカラー剛性条件は、この主張には不要である。
\end{lemma}

\begin{proof}
$A_0=\Sym_{R_0}(V)$ とし、$f=(F_x,F_y,F_z)$ と置く。
相対点定理の仮定により、この三元は任意の基底変更後にも正則列である。
三次式 $H$ が Jacobian Poisson 括弧の Casimir であることは
$\nabla_{x,y,z}H\wedge f=0$ と同値である。
正則列の Koszul cochain 複体の $A_0^3$ における完全性から
$\nabla H=a f$ となる $a\in A_0$ が一意に存在する。
両辺の斉次次数を比較し、$A_0\to A_0^3$, $a\mapsto af$ の単射性を使うと
$a\in R_0$ である。
従って $H-aF$ の三つの active 偏微分は零となる。
$R_0$ は標数零の体上の代数なので、$H-aF$ は純中心三次である。
逆包含は定義から明らかである。
滑らかな三変数制限は各点で非零なので、表示した和は局所的に直和となる。
\end{proof}

尺度を含む周囲 Poisson tensor $\Pi=h\Pi_F$ を固定する。
その support は点ベクトル空間 $V^*$ の三次元部分空間
$K=\operatorname{ann}(U)$ である。
ここでは既存の二次包絡復元から得られる $\Pi$ の尺度を捨てない。
$\mathcal R=kF\oplus\Sym^3U$ を補題\ref{lem:point-relative-casimirs}の
三次 Poisson Casimir の空間とし、
$\mathcal L_{\mathrm{cl}}=\mathcal R/\Sym^3U$ と置く。
尺度込みの $\Pi$ を用いているので、命題\ref{prop:point-casimir}で
$Q$ だけから復元するための Hessian 剛性条件を追加する必要はない。

有限自由 $S$-代数 $T$ は階数八なので、正規化した trace
\begin{equation}\label{eq:point-normalized-trace}
 \operatorname{tr}:T\longrightarrow S,
 \qquad a\longmapsto\frac18\operatorname{Tr}_{T/S}(m_a)
\end{equation}
が定まる。これは次数を保ち、$S$ 上で恒等である。
任意の係数環上でも有限局所自由代数の乗法行列の trace で定義されるため、
任意の基底変更と可換する。
以下、$\operatorname{tr}(H)$ は $H$ を $T$ に送ってから trace を取る意味とする。
従って
\begin{equation}\label{eq:point-trace-splittings}
 \begin{aligned}
 \mathcal E&=\mathcal W\oplus\mathcal E^0,
 &\mathcal E^0&=\ker(\operatorname{tr}\circ\alpha),\\
 \mathcal R&=\mathcal W\oplus\mathcal R^0,
 &\mathcal R^0&=\ker(\operatorname{tr}|_{\mathcal R})
 \end{aligned}
\end{equation}
は標準的な分裂であり、後半は標準的な持上げ
$s_{\mathrm{cl}}:\mathcal L_{\mathrm{cl}}\to\mathcal R^0$ を与える。

\begin{lemma}[尺度付き Hessian 同型]\label{lem:point-hessian-line}
尺度込みの $\Pi$ から、表示に依存しない線束同型
\begin{equation}\label{eq:point-hessian-line}
 \nu_\Pi:\mathcal L_{\mathrm{cl}}\xrightarrow{\sim}(J_g)_3
\end{equation}
が定まる。局所的に非零の $\overline F\in\mathcal L_{\mathrm{cl}}$ を選び、
$\Pi=\iota_{dF}P$ となる $P\in\Lambda^3K$ を取れば、
\[
 \nu_\Pi(\overline F)=
 \left[\mathscr H(P,F)\right],\qquad
 \mathscr H(P,F)=\frac16\sum P^{ijk}P^{abc}F_{ia}F_{jb}F_{kc}.
\]
特に $\Pi=h\Pi_F$、$P=h\partial_x\wedge\partial_y\wedge\partial_z$ ならば
$\nu_\Pi(\overline F)=h^2[\det\operatorname{Hess}(g)]$ である。
\end{lemma}

\begin{proof}
この縮約はテンソルとして定義され、基底変更に不変である。
$F$ に純中心三次を加えても $K$ に沿う Hessian は変わらない。
さらに $F$ を $cF$ に替えると、同じ $\Pi$ を与える $P$ は $c^{-1}P$ となり、
縮約は $c$ 倍となる。従って表示した式は商直線上の線形写像を定める。
Hesse 座標では
\[
 [\det\operatorname{Hess}(g)]=8(1+\lambda^3)[xyz]\ne0
\]
なので同型である。滑らかな三変数三次式を Hesse 形にする変換と
縮約の共変性から一般の場合が従う。
族の場合にも、全ての幾何ファイバー上の非消滅により線束同型となる。
\end{proof}

\begin{theorem}[Trace--Hessian による中心三次の同定]
\label{thm:point-trace-hessian}
境界 Jacobi 写像 $\epsilon$ が同型となる良い一ブロック開上で、
\begin{equation}\label{eq:point-trace-hessian-map}
 \mathfrak d_\Pi:\mathcal E\xrightarrow{\sim}\mathcal R,
 \qquad
 w\longmapsto
 -12\,s_{\mathrm{cl}}\bigl(\nu_\Pi^{-1}(\epsilon(\overline w))\bigr)
 +\operatorname{tr}(\alpha(w))
\end{equation}
は線型同型であり、純中心三次上で恒等である。
これは基底変更および次数付き表示変更と両立する。
従って中心三次直線の torsor の同型
\begin{equation}\label{eq:point-central-classical-torsors}
 \mathbb P(\mathcal E)\setminus\mathbb P(\mathcal W)
 \xrightarrow{\sim}
 \mathbb P(\mathcal R)\setminus\mathbb P(\mathcal W)
\end{equation}
を定める。

$H_w=\mathfrak d_\Pi(w)$ が滑らかな射影三次超曲面を定める非空開上で、
\begin{equation}\label{eq:point-poisson-pair}
 [w]\longmapsto (X_w,\pi_w),\qquad
 X_w=V(H_w),\quad \pi_w=\Pi|_{X_w}
\end{equation}
は次数付き表示のスカラー商上の射となる。
より具体的には、$\mathcal Z_N^{\mathrm{sm}}$ を
\[
 \mathbb P(\Sym^3V)\times
 \bigl(\Sym^2V\otimes\Lambda^2V^*\bigr)
\]
内で、三次式の滑らかさおよび $[\Pi,\Pi]=0$、$\Pi(dH,-)=0$ を満たす
$([H],\Pi)$ の局所閉スキームとすれば、標的を代数的商
$[\mathcal Z_N^{\mathrm{sm}}/\mathrm{PGL}(V)]$ と取れる。
斉次 affine lift のデータを忘れることで、Brauer--Severi 周囲空間内の
滑らかな埋込み三次 Poisson 対とその偏極\emph{類}へ射を得る。
大域的な $\mathcal O(1)$ 自体の降下を仮定する必要はない。
定理\ref{thm:point-omarked-descent}と局所的な良い表示のモジュライ比較により、
この射は同じ開の $\mathcal O$-marked な族についても表示に依存しない。
\end{theorem}

\begin{proof}
\eqref{eq:point-trace-splittings} の両分裂の下で、写像は $\mathcal W$ 上の恒等と、
商直線間の同型 $-12\nu_\Pi^{-1}\epsilon$ の直和である。
従って線型同型であり、非純中心の直線を非純中心の直線へ送る。
全ての構成は最大可換代数商、有限局所自由代数の trace、
尺度付き Poisson tensor およびその三次 Casimir から定まる。
これらは表示変更に共変で、任意の基底変更と両立する。
次数付きスカラーは両三次空間で重み三を持つため、
射影化後には自明に作用する。
このことが表示のスカラー商とその fpqc forms 上での降下を与える。

$H_w$ は $\Pi$ の Casimir であるから、その斉次イデアルは Poisson イデアルである。
括弧は次数零を持つので $\operatorname{Proj}$ 上にも Poisson tensor を定める。
一般点では中心座標が非零で Euler 方向が active support の外にあるため、
得られる射影 Poisson tensor は階数二である。

滑らかな三次式の非空性も各固定周囲で確認できる。
$F+G(u)$ を $G\in\Sym^3U$ で動かす線型系は $V(U)$ の外で基点を持たない。
基点部分 $V(U,g)$ では active 勾配が非零なので、各 $F+G$ は滑らかである。
Bertini により一般の $G$ はその外でも滑らかな超曲面を与える。
同型\eqref{eq:point-central-classical-torsors}で引き戻せば、各固定周囲の
中心三次 torsor に非空開を得る。
全パラメータ空間では、既に課した良い点スキーム開や
滑らかな定常成分の非空開とも交わらせられる。
後者を全ての固定周囲で主張しているのではない。
\end{proof}

\begin{remark}[規格化・既存量子化・残余圏との比較]
式\eqref{eq:point-trace-hessian-map}の同型自体は、係数 $-12$ を
任意の固定した非零定数に替えても成立する。
この値は入力研究\cite[Proposition \texttt{prop:pg-ad-intrinsic-hessian}]{NCCubicA}の
二次 jet に合わせたものである。
すなわち、指定した Weyl 表示の二次 jet
\[
 \Omega_{F,h}=\operatorname{sym}
 \left(F-\frac{h^2}{12}\det\operatorname{Hess}_{x,y,z}(F)+O(h^4)\right)
\]
を満たす族に限れば、$\epsilon(\overline\Omega)
=-h^2[\det\operatorname{Hess}(g)]/12+O(h^4)$ なので、
$\mathfrak d_{h\Pi_F}(\Omega_{F,h})=F+O(h^2)$ となる。
ここで指定するのは、同入力命題で純中心二次補正を調整した中心元である。
同入力の短い混合 Hesse 公式には $h^2ABC/6$ という純中心項もあり、
それを引く前の式を上の $O(h^4)$ 表示と同一視しない。
有限の非零 $h$ で常に入力の $F$ と一致するとは主張しない。
またこの同変な構成の存在は、自然な逆量子化の一意性を意味しない。

点スキームの定常切断 $w_{\mathrm{ab}}|_C$ と $H_w|_C$ は異なる。
前者の境界 Jacobi 類を消すのではなく、その非零類を Hessian 同型によって
古典 Casimir 方向の係数へ読み替え、純中心成分を trace で指定している。
これにより Poisson 三次対への射は構成できるが、
$X_w$ と残余圏が回復する可換三次多様体 $X_{\mathrm{res}}$ の同型、
および $\operatorname{Ku}(A)\simeq\operatorname{Ku}(X_w)$ は別の比較問題として残る。
特定の量子化 chart の写像を全ての内在的モジュライ成分の写像に拡張したとは主張しない。
\end{remark}

% !TEX root = ../nc-cubic-n-fold.tex
\section{強い affine 持上げと Poisson 対からの逆構成}
\label{sec:poisson-inverse}

体上の三次超曲面の双ベクトル表示
\eqref{eq:nfold-bivector}は入力論文\cite[命題
\texttt{prop:pg-represent}]{NCCubicA}にある。
ここではその任意基底上の形を証明し、Poisson 射の標的から affine lift の
追加データを外す。さらに良い一ブロック型では逆構成を与え、
前節の射がどの開部分を同定するかを明示する。
この逆構成は、得られる三次の Kuznetsov 成分との同値を意味しない。

\subsection{任意基底上の強い持上げ}

$R$ を可換 $\mathbf C$-代数、$V$ を階数 $N=n+2\ge5$ の有限局所自由
$R$-加群、$W=V^*$ とする。局所的に
$H\in\Sym^3V$ が $R$ 上滑らかな三次超曲面
$X\subset\mathbb P(W)$ を定めるとする。
$E=\sum z_i\partial_{z_i}$ を Euler ベクトル場と書く。
二次係数の斉次 affine 双ベクトル
$\Pi\in\Sym^2V\otimes\Lambda^2W$ が
\[
 \iota_{dH}\Pi=0
\]
を多項式の恒等式として満たすとき、$H$ に対する\emph{強い持上げ}と呼ぶ。
これは $H$ による剰余での接性より強い条件である。

\begin{theorem}[強い affine 持上げの存在と一意性]
\label{thm:poisson-strong-lift}
任意の $\pi\in H^0(X,\Lambda^2T_{X/R})$ は一意な強い持上げを持つ。
その持上げは、一意な定数三重ベクトル $P\in\Lambda^3W$ を用いて
\[
 \Pi=\iota_{dH}P
\]
と書ける。これらの同定は任意の基底変更と両立する。
$H$ の生成元を変えたときには $H\mapsto cH$, $P\mapsto c^{-1}P$ となる。
従って三次方程式の直線束を $\mathcal L_H$ と書けば、全体として
\begin{equation}\label{eq:poisson-relative-bivectors}
 p_*\Lambda^2T_{X/R}\simeq\mathcal L_H\otimes\Lambda^3W
\end{equation}
である。特に左辺は有限局所自由で、任意の基底変更と両立する。
また $\pi$ が Poisson であることと、強い持上げ $\Pi$ が affine Poisson
双ベクトルであることは同値である。この同値も非被約基底を含めて成立する。
\end{theorem}

\begin{proof}
基底を局所化して $V$ と方程式直線を自明化する。
$\mathbb P=\mathbb P^{N-1}_R$ と書く。
Euler 列とその二次外冪の列は
\[
 0\to\OO\to W\otimes\OO(1)\to T_{\mathbb P/R}\to0,\qquad
 0\to T_{\mathbb P/R}\to\Lambda^2W\otimes\OO(2)
      \to\Lambda^2T_{\mathbb P/R}\to0
\]
である。相対次元が $N-1\ge4$ なので、射影空間の線束コホモロジーから
\[
 R^ip_*\bigl(\Lambda^2T_{\mathbb P/R}(-3)\bigr)=0
 \quad(i=0,1),\qquad R^1p_*T_{\mathbb P/R}=0
\]
を得る。これらの消滅は任意の係数環上で成立する。
従って $\mathbb P$ 上の双ベクトルの $X$ への制限は同型であり、
各切断は二次係数の affine 双ベクトル $q$ で表せる。
この affine 表示の核は、一次係数ベクトル場 $L$ に対する $E\wedge L$ である。

$q$ が $X$ 上で $\Lambda^2T_{X/R}$ に値を持つため、
$\iota_{dH}q$ の $T_{\mathbb P/R}(3)|_X$ における像は零である。
制限列と Euler 列の核を用いて
\[
 \iota_{dH}q=HL+aE
\]
と書ける。ここで $L$ は一次係数のベクトル場、$a$ は三次式である。
もう一度 $dH$ と縮約すると $H(L(H)+3a)=0$ を得る。
滑らかな相対三次の方程式 $H$ は、任意基底変更後にも多項式環上の非零因子である。
従って $a=-L(H)/3$ であり、
\[
 \Pi=q-\tfrac13E\wedge L
\]
は強い持上げである。
二つの強い持上げの差は $\Delta=E\wedge L$ と書ける。
その縮約は $3HL-L(H)E=0$ であるから、さらに $E$ と外積を取ると
$3H\Delta=0$ となる。$H$ の非零因子性から $\Delta=0$ であり、一意性を得る。

滑らかな三次の偏微分 $H_1,\ldots,H_N$ の Koszul 複体は、
各次数の有限自由複体をファイバー上の完全性から局所的に分離することで、
任意基底変更後にも完全である。
従って二次係数の $\Pi$ に対する $\iota_{dH}\Pi=0$ は
$\Pi=\iota_{dH}P$ と同値である。
係数次数が零なので $P$ は定数三重ベクトルであり、
次の Koszul 項の係数次数は $-2$ だから、その $P$ は一意である。
構成の一意性によって局所表示は貼り合わさり、全ての基底変更と座標変更に両立する。

最後に Jacobi 条件を確認する。
標準体積形式の divergence を $\delta$ と書くと、
$\Pi=\iota_{dH}P$ は $\delta\Pi=0$ を満たす。
$J=[\Pi,\Pi]$ は $\delta J=0$ および $\iota_{dH}J=0$ を満たす。
$J$ の射影化が $X$ 上で零なら、Euler 列の三次外冪と制限列により
\[
 J=HT+E\wedge B,\qquad
 T\in\Lambda^3W,\quad B\in\Sym^2V\otimes\Lambda^2W
\]
と書ける。同じ Euler コホモロジー計算がこの表示を任意基底上で与える。
$\delta(E\wedge B)=NB-E\wedge\delta B$ より
\[
 E\wedge B=-N^{-1}E\wedge\iota_{dH}T,\qquad
 0=\iota_{dH}J=(1-3/N)H\iota_{dH}T.
\]
$N\ne3$、$H$ は非零因子であり、定数三重ベクトル上の
$\iota_{dH}$ は上述の Koszul 計算で単射なので、$T=0$、従って $J=0$ となる。
逆向きは制限するだけである。
\end{proof}

\begin{corollary}[一ブロック support の回復]
\label{cor:poisson-support-recovery}
定理の $P$ が非零分解可能なら、その support
$K\subset W$ は階数三の部分束として回復される。
$U=\operatorname{ann}(K)$ は $\Pi$ の一次 Poisson 中心に等しい。
従って滑らかな三次の Poisson 対とその超平面偏極類から、
強い持上げ、中心一次部分、および一ブロック support を回復できる。
\end{corollary}
\begin{proof}
相対的な「非零分解可能」とは、非零 Pl\"ucker 錐のスキームを通るという条件とする。
これは、階数三の部分束 $K$ と
$\mathcal L_H\otimes\det K$ の非零切断を与えることに等しい。
従って support の形成は任意基底変更と両立する。
一次式 $u$ が Poisson 中心に属すれば
$\iota_{dH}(\iota_{du}P)=0$ である。
定数二重ベクトル上でも Koszul 微分は単射だから、$\iota_{du}P=0$、
すなわち $u\in\operatorname{ann}(K)$ である。逆向きは直接従う。
\end{proof}

\subsection{分解可能な Poisson 対と良い量子化の開}

$\mathcal H_N^{\mathrm{sm}}\subset\mathbb P(\Sym^3V)$ を滑らかな三次式の開、
$\mathcal L_H\subset\Sym^3V\otimes\OO_{\mathcal H_N^{\mathrm{sm}}}$ を普遍直線とする。
$\mathcal Z_N^{\mathrm{dec}}$ を、$([H],\Theta)$ で
$0\ne\Theta\in\mathcal L_H\otimes\Lambda^3W$ が分解可能なものの相対スキームとする。
具体的には、$\mathcal H_N^{\mathrm{sm}}\times\operatorname{Gr}(3,W)$ 上の
直線束 $\mathcal L_H\otimes\det\mathcal K$ の零切断を除いた空間である。
従ってこれは滑らかで既約である。
次数スカラーの重みは $3-3=0$ なので $\operatorname{PGL}(V)$ が作用する。
\[
 \mathfrak P_n^{\mathrm{dec}}
   =[\mathcal Z_N^{\mathrm{dec}}/\operatorname{PGL}(V)]
\]
と置く。前定理によりこれは、強い lift に対応する三重ベクトルが
非零分解可能な滑らかな三次 Poisson 対と超平面偏極類のモジュライである。
affine lift 自体を追加標識として保持する必要はない。
偏極類を保つ同型は、局所的に $\OO_X(1)$ を同定し全切断を取ることで、
ambient の射影同型へ一意に延長する。
ここでは $p_*\OO_X(1)=V$ および任意基底変更との整合を使う。

有限安定化子を要求する前の良い一ブロック PBW 開を
$T^{\mathrm{alg}}$ と書く。
ここには、相対 PBW・普遍中心・相対 Serre 双対・全
$\mathcal O$-同値比較に必要な入力を全て課すが、
周囲代数 $B$ 単独の安定化子が有限という条件は課さない。
これらの入力が成立することは、対象とする表示チャートごとに必要である。
中心三次 torsor を $J^{\mathrm{alg}}$ とする。
その上で境界 Jacobi 評価 $\epsilon$ の同型条件、相対点スキームの良い条件、
および $H_w=\mathfrak d_\Pi(w)$ の滑らかさを課した非空開を
$J^{\mathrm{sm}}$ と書く。
\[
 \mathfrak M_n^{\mathrm{sm}}=[J^{\mathrm{sm}}/\overline P]
\]
はこれらの良い\emph{対}を局所模型とする $\mathcal O$-marked 圏族の stack である。
必要なら旧 $T^{\mathrm{good}}$ を縮め、旧 Poisson 射の定義域
$\mathfrak M_n^{\mathrm{Pois}}$ がこの stack の有限周囲安定化子部分の
非空開になるようにする。
旧 $\mathfrak M_n^{\mathrm{loc}}$ 全体は $H_w$ の滑らかさ等を要求していないので、
その全体を $\mathfrak M_n^{\mathrm{sm}}$ の開と同一視しない。

\begin{theorem}[Poisson 対からの逆構成と開部分の同定]
\label{thm:poisson-good-open-isomorphism}
$\mathfrak P_n^{\mathrm{dec}}$ には次の条件で定まる非空開
$\mathfrak P_n^{\mathrm{good}}$ があり、trace--Hessian 射は同型
\begin{equation}\label{eq:poisson-good-open-isomorphism}
 \mathfrak M_n^{\mathrm{sm}}\xrightarrow{\sim}
 \mathfrak P_n^{\mathrm{good}}
\end{equation}
を与える。
\begin{enumerate}
\item 強い持上げ $\Pi$ から定める二次代数
\[
 B_\Pi=T(V)/
 \bigl(v\otimes v'-v'\otimes v-\operatorname{sym}\Pi(v,v')\bigr)
\]
が $T^{\mathrm{alg}}$ の良い模型であり、境界評価 $\epsilon$ が同型である。
\item 一意な中心直線
\[
 [w]=\mathfrak d_\Pi^{-1}([H])\subset Z(B_\Pi)_3
\]
が相対点スキームの良い条件を満たす。
\end{enumerate}
同型は任意の基底変更と両立し、全ての自己同型群を同定する。
$\mathfrak M_n^{\mathrm{sm}}$ は滑らかな既約有限型 DM stack で、次元は
\[
 \binom{n+4}{3}-(n+2)^2+3n-2
 =\binom{n+2}{3}+3n-2
\]
である。周囲代数 $B$ の安定化子が正次元であっても、
ここに列挙した良い入力と $H_w$ の滑らかさがあれば適用できる。
\end{theorem}

\begin{proof}
support を局所的に固定し、$K$ の体積元 $P_0$ を一つ選ぶ。
$\Theta=H\otimes tP_0$ と書けば、対応する二次代数の potential は
$F=tH$ modulo $\Sym^3U$ である。
$H\mapsto cH$ に対して $t\mapsto c^{-1}t$ なので、この $F$ は変わらない。
support を保つ frame の変更では、従来と同じ
$\det(g|_{V/U})^{-1}g(F)$ の作用となる。
従って条件(1)は $P$-不変な良い表示開の引戻しであり、開条件である。
その上では trace--Hessian 同型が任意基底変更で定義される。
条件(1)に含まれる active 楕円三次の仮定により、
$H|_{\mathbb P(K)}$ が滑らかであるため $[H]$ は純中心方向でなく、
条件(2)の中心直線は一意に定まる。条件(2)も相対 norm による開条件である。
従来の非空開 $J^{\mathrm{Pois}}$ の像を含むので、得られる開は非空である。

この逆構成から $A=B_\Pi/(w)$ を作る。
次数一生成元束は一般には表示 gerbe 上にあるが、
第\ref{sec:nfm-local-moduli}節の重み付き tails の構成は
周囲の有限安定化子を用いないので、$\mathcal O$-marked 圏族を与える。
全 $\mathcal O$-同値比較も同じ入力だけを使うため、
この圏族の同型問題は $[J^{\mathrm{sm}}/\overline P]$ に一致する。

量子側から出発すると、元の $\Pi$ は $H_w$ に対する強い持上げである。
定理\ref{thm:poisson-strong-lift}の一意性により、Poisson 対から回復する持上げは
元の $\Pi$ と等しい。従って $B$ をその二次関係から回復し、
trace--Hessian 同型の逆で元の中心直線 $[w]$ を得る。
逆に Poisson 側から作った $A$ に trace--Hessian 射を適用すると、
定義により元の $[H]$ と $\Pi|_X$ に戻る。
これらは係数変更と表示変更に両立する、相互に逆な構成である。

射も一致する。Poisson 対の偏極類を保つ同型は ambient の射影同型を与え、
強い持上げの一意性から $\Pi$ を保つ。
従って $K,U,B_\Pi$ を保ち、trace--Hessian の共変性から $[w]$ も保つ。
逆に全 $\mathcal O$-同値は局所的に次数付き同型に持ち上がるため、
Poisson 対の同型を与える。二つの線型持上げの差は次数スカラーである。
これにより幾何点の全単射だけでなく、fpqc stack の同型を得る。

DM 性を直接確認する。
滑らかな三次 $H$ の射影安定化子の Lie 代数に属する一次ベクトル場 $L$ は
$L(H)=cH$ を満たす。Euler の公式を引けば
\[
 \sum_i(L_i-cz_i/3)H_i=0
\]
となる。滑らかな三次の偏微分の正則列には一次係数の syzygy がないため、
$L=cE/3$ であり、射影 Lie 代数では零である。
従って射影安定化子は有限で、標数零では非分岐である。
Poisson 対の安定化子はその部分群なので、同じ結論を持つ。
$J^{\mathrm{sm}}$ は滑らかな既約表示空間上の加法 torsor の非空開であり、
$\overline P$ は滑らかなので、商は滑らかかつ既約である。
表示空間と群の次元を引けば主張の次元を得る。
\end{proof}

\begin{corollary}[拡大した対の stack における Brauer 障害]
\label{cor:poisson-enlarged-brauer}
$[J^{\mathrm{sm}}/P]\to\mathfrak M_n^{\mathrm{sm}}$ の表示 gerbe の類を
$\beta^{\mathrm{sm}}$ とする。
$n\not\equiv1\pmod3$ では中立であり、
$n\equiv1\pmod3$ ではどの非空開上でも位数が正確に三である。
\end{corollary}
\begin{proof}
determinant character による $r\beta=3\beta=0$ と
$\gcd(r,3)=1$ の場合の中立化には、有限安定化子は使われない。
$\gcd(r,3)=3$ の場合、旧 $\mathfrak M_n^{\mathrm{Pois}}$ は
既約な $\mathfrak M_n^{\mathrm{sm}}$ の非空開なので稠密である。
任意の非空開との交わりに定理\ref{thm:nfm-brauer-dense-open}を適用すれば、
障害が非零の体値点を得る。これで位数三が従う。
\end{proof}

\subsection{全 Poisson モジュライの一般部分との関係}

\begin{corollary}[一般三次上の開埋込みと支配的成分]
\label{cor:poisson-generic-open-immersion}
系\ref{cor:nfold-vertex-support}と注意\ref{rem:nfold-relative-punctured}が適用される
滑らかな三次式の非空不変開を $\mathcal H_N^\circ$ とする。
この開上では、非零三次 Poisson 対の全モジュライは
$\mathfrak P_n^{\mathrm{dec}}$ の制限に等しい。
従って\eqref{eq:poisson-good-open-isomorphism}の制限は、
全非零 Poisson モジュライへの非空な開埋込みである。
その像は滑らかな三次のモジュライを支配する。

滑らかな三次 Poisson 対の全モジュライには、三次のモジュライを支配する
既約成分が一つだけある。本構成はその一般部分の非空開を同定する。
特殊な三次の locus 上にのみ存在する他の既約成分を排除する主張ではない。
\end{corollary}
\begin{proof}
相対非零 Jacobi スキームと非零 Pl\"ucker 錐の等式は、引用した結果により
任意の基底変更と両立する。三次方程式の直線束の捻りと射影群の商を戻すと、
表示したモジュライの等式を得る。
$\mathcal Z_N^{\mathrm{dec}}\to\mathcal H_N^{\mathrm{sm}}$ は滑らかな全射であり、
各幾何ファイバーは非零 Grassmann 錐で既約である。
従ってその任意の非空開の像は基底の非空開となり、支配性が従う。
一般の三次上で全 Poisson スキームの台は一つの Grassmann 錐であるため、
一般ファイバーは既約である。有限型の全モジュライで基底を支配する
既約成分は一般ファイバーの既約成分に対応するので、一つだけである。
頂点に支持される非被約構造は、この台の既約性の主張を変えない。
\end{proof}

\begin{remark}[開埋込みの範囲]
分解可能な三重ベクトルの条件は、全ての滑らかな三次を動かす空間では
一般に閉条件であり、非零条件を加えても自動的に開にはならない。
従って定理\ref{thm:poisson-good-open-isomorphism}の無条件の標的は
$\mathfrak P_n^{\mathrm{dec}}$ である。
全非零 Poisson モジュライへの開埋込みには、上の一般三次の開への制限を用いる。
また、一般対象の局所一ブロック表示をこの Poisson 対のモジュライの範囲で得たことと、
未構成の全 intrinsic $\mathcal O$-marked 圏モジュライの全成分を分類することは異なる。
\end{remark}

% !TEX root = ../nc-cubic-n-fold.tex
\section{JacobiスキームとGrassmann錐の次元境界}
\label{sec:nfold-jacobi-scheme}

本節は、Poisson構造の被約な分類と、その未被約スキーム構造を区別する。
$N=n+2\ge5$、$W=\mathbb C^N$とし、$F\in\operatorname{Sym}^3W^*$が
滑らかな三次超曲面$X_F\subset\mathbb P(W)$を定めるとする。
$\widehat{G}_N\subset\Lambda^3W$を$\operatorname{Gr}(3,W)$のPl\"uckerアフィン錐とし、
$J_N\subset R_N:=\operatorname{Sym}((\Lambda^3W)^*)$をその斉次イデアルとする。
双ベクトルのスカラーは割らない。

\subsection{Hessian二次式によるスキームの表示}
座標を$z_1,\ldots,z_N$、$F_i=\partial_iF$、
$F_{ijr}=\partial_i\partial_j\partial_rF$とする。
$P\in\Lambda^3W$に対し
$q_P=\iota_{dF}P$と置く。
既存稿\cite[命題\texttt{prop:pg-represent}]{NCCubicA}の
Euler列とgradientのKoszul複体による表示は、
\begin{equation}\label{eq:nfold-bivector}
 \Lambda^3W\xrightarrow{\sim}H^0(X_F,\Lambda^2T_{X_F}),
 \qquad P\longmapsto q_P|_{X_F}
\end{equation}
を与える。$F$を$aF$へ取り替えると$P$は$a^{-1}P$へ取り替える。
以下のイデアルはこの選択に依存せず同型となる。

$H\in\operatorname{Sym}^2W^*$について、
\begin{align}\label{eq:nfold-hessian-B}
 B_H(P)(\alpha,\beta,\gamma,\eta)
 &=H(P(\alpha,\beta),P(\gamma,\eta))\notag\\
 &\quad+H(P(\beta,\gamma),P(\alpha,\eta))
       +H(P(\gamma,\alpha),P(\beta,\eta))
\end{align}
と定める。ここで$P(\alpha,\beta)\in W$は
$\eta(P(\alpha,\beta))=P(\alpha,\beta,\eta)$で定める。
$B_H(P)$は四つの共変ベクトルに関して交代的であり、$\Lambda^4W$の元である。
$H_r=(F_{ijr})_{ij}$と置き、$I_F\subset R_N$を
\begin{equation}\label{eq:nfold-jacobi-quadrics}
 B_{H_r}(P)(dz_a,dz_b,dz_c,dz_d),
 \qquad 1\le r\le N,\quad a<b<c<d
\end{equation}
で生成される二次斉次イデアルとする。

\begin{lemma}\label{lem:nfold-jacobi-scheme}
Poisson構造の自然なJacobi零スキームは
\begin{equation}\label{eq:nfold-poisson-scheme}
 \operatorname{Pois}(X_F)=\operatorname{Spec}(R_N/I_F)
\end{equation}
である。この同定は任意の可換$\mathbb C$-代数、特に非被約なArtin係数への
基底変更と両立する。また
\begin{equation}\label{eq:nfold-quadratic-inclusion}
 I_F\subset J_N,\qquad
 \dim_{\mathbb C}(I_F)_2\le N\binom N4
\end{equation}
が成り立つ。
\end{lemma}
\begin{proof}
既存稿のJacobi延長の議論を、係数環を含めて確認する。
$E=\sum z_i\partial_i$と標準体積形式のdivergence $\delta$を用いると、
$q_P$は$\delta q_P=0$および$\iota_{dF}q_P=0$を満たす。
$J=[q_P,q_P]$の射影化が$X_F$で零なら、Euler列と制限列から
$J=FT+E\wedge B$、$T\in\Lambda^3W$、
$B\in\operatorname{Sym}^2W^*\otimes\Lambda^2W$と書ける。
$\delta J=0$と
$\delta(E\wedge B)=NB-E\wedge\delta B$により
$E\wedge B=-N^{-1}E\wedge\iota_{dF}T$。
さらに$\iota_{dF}J=0$から
$(1-3/N)F\iota_{dF}T=0$を得る。
$N\ne3$、$F$は零因子でなく、$F_i$は定数係数で独立なので$T=0$、$J=0$となる。
この議論は\cite[命題\texttt{prop:pg-jacobi-extension}]{NCCubicA}と同じものである。

アフィンJacobiatorの$ijk$成分は、非零定数倍と符号を除いて
\begin{equation}\label{eq:nfold-J-linearized}
 \sum_a F_a L_a^{ijk},\qquad
 L_a^{ijk}=
 \sum_{r,b}\bigl(P^{ira}P^{jkb}+P^{jra}P^{kib}+P^{kra}P^{ijb}\bigr)F_{br}
\end{equation}
である。$F_1,\ldots,F_N$は二次式の正則列なので、
一次式を係数とするsyzygyは存在しない。
従って$J=0$と全ての$L_a^{ijk}=0$は同値であり、後者は
\eqref{eq:nfold-jacobi-quadrics}に一致する。
ここで用いた写像
$\bigoplus_a\operatorname{Sym}^1W^*\to\operatorname{Sym}^3W^*$、
$(\ell_a)\mapsto\sum\ell_aF_a$は、有限次元$\mathbb C$-ベクトル空間の
分裂単射である。Euler列から得た有限次元の同型も分裂し、
$F$による乗法も各次数で分裂単射である。
従って任意の係数環へテンソルしても、上の同値は保たれる。
これは複素点の集合の等式だけでなく、零スキームの等式を示す。

$P$が分解可能なら、そのsupportの三次元部分空間に沿う
Jacobian括弧となるため、任意の$F$についてJacobi恒等式を満たす。
従って全ての生成二次式は$\widehat G_N$上で消え、$I_F\subset J_N$となる。
最後の上界は\eqref{eq:nfold-jacobi-quadrics}の本数である。
\end{proof}

\begin{lemma}\label{lem:nfold-plucker-count}
Grassmann錐の二次イデアルの次元は
\begin{equation}\label{eq:nfold-plucker-count}
 q_N:=\dim(J_N)_2
 =5\binom{N+1}{6}=N\binom N5-\binom N6
\end{equation}
である。
\end{lemma}
\begin{proof}
Pl\"ucker環の標準単項式基底を用いる
\cite[\S2.4, Theorem~2.3]{BCMStandardMonomial}。
三元部分集合を増加列として書き、成分ごとの順序を用いると、
次数二の標準単項式は$p_Ap_B$、$A\le B$である。
従って$(J_N)_2$の次元は、辞書式に$A\le_{\mathrm{lex}}B$と並べた
三元部分集合対で、成分ごとには比較できないものの個数に等しい。
二つ以上の添字を共有する対は比較可能である。
和集合の大きさが五の場合、固定した五元集合上には比較不能な対が五個ある。
添字を$1,\ldots,5$へ移した表示は
\[
 (125,134),\ (125,234),\ (135,234),\ (145,234),\ (145,235).
\]
和集合の大きさが六の場合も五個であり、
\[
 (126,345),\ (136,245),\ (145,236),\ (146,235),\ (156,234)
\]
である。従って$q_N=5\binom N5+5\binom N6=5\binom{N+1}{6}$を得る。
\end{proof}

\begin{lemma}[Jacobi 二次式の普遍的な一次従属]
\label{lem:nfold-universal-row-relations}
$N\ge5$とし、$T\in\operatorname{Sym}^3W^*$を任意の対称三次テンソルとする。
式\eqref{eq:nfold-jacobi-quadrics}の係数写像を
\[
 \Phi_T:W\otimes\Lambda^4W^*\longrightarrow (J_N)_2,
 \qquad e_r\otimes\alpha\longmapsto B_{H_r}(P)(\alpha)
\]
と書く。ただし$\alpha$に関する評価は線形に延長する。
座標に依存しない分裂単射
\[
 \iota:\Lambda^3W^*\longrightarrow W\otimes\Lambda^4W^*,
 \qquad\beta\longmapsto\sum_{r=1}^N e_r\otimes(dz_r\wedge\beta)
\]
について$\Phi_T\iota=0$であり、従って
\[
 \dim(I_F)_2\le N\binom N4-\binom N3.
\]
さらに、全ての対称三次テンソルを同時に考えた共通核は
\[
 \bigcap_{T\in\operatorname{Sym}^3W^*}\ker\Phi_T=\operatorname{im}\iota
\]
に等しい。この共通核の等式だけから、固定した一般の$T$で核が
$\operatorname{im}\iota$に等しいとは結論しない。
\end{lemma}
\begin{proof}
任意の$a<b<c$について、Jacobi二次式の間に
\[
 \sum_{r=1}^N B_{H_r}(P)(dz_r,dz_a,dz_b,dz_c)=0
\]
がある。実際、左辺を展開すると
\[
 \sum_{r,i,j}T_{ijr}
 \bigl(p_{rai}p_{bcj}+p_{abi}p_{rcj}+p_{bri}p_{acj}\bigr)
\]
である。第一項・第三項では$(r,i)$、第二項では$(r,j)$を交換すると、
$T$は不変で該当する$p$は符号を変えるので、各和は零である。
これが$\Phi_T\iota=0$を示す。
縮約写像
\[
 \kappa:W\otimes\Lambda^4W^*\longrightarrow\Lambda^3W^*,
 \qquad e_r\otimes\alpha\longmapsto\iota_{e_r}\alpha
\]
に対し$\kappa\iota=(N-3)\operatorname{id}$だから、
$\iota$は分裂単射であり、$\binom N3$個の独立な従属関係を得る。

共通核を調べる。$u=\sum e_r\otimes\alpha_r$をその元とし、
任意の$\ell\in W^*$に対して$T=\ell^3$を代入する。
$H_r=\ell(e_r)\ell^2$であるから、
$\alpha(\ell):=\sum\ell(e_r)\alpha_r$は
全ての$P$に対して$B_{\ell^2}(P)$を消す。
$\ell\ne0$なら$\iota_\ell P$は$\Lambda^2\ker\ell$の任意の元を走り、
$B_{\ell^2}(P)$はその外積二乗の非零定数倍である。
二重ベクトルの外積二乗は$\Lambda^4\ker\ell$を張るので、
$\alpha(\ell)$の$\ker\ell$への制限は零、すなわち
$\ell\wedge\alpha(\ell)=0$である。

最後の条件から$\alpha(\ell)=\ell\wedge\beta$となる
一定の$\beta\in\Lambda^3W^*$が存在することを座標で確認する。
$\ell=dz_i$を代入すると$\alpha_i=dz_i\wedge\beta_i$と書ける。
$\ell=dz_i+dz_j$を代入すると
$dz_i\wedge dz_j\wedge(\beta_j-\beta_i)=0$である。
任意の三元添字集合$K$について、$i\notin K$の場合の
$\beta_i$の$dz_K$係数は$i$によらない。この共通係数を$\beta$の
$dz_K$係数とすれば、全ての$i$について
$\alpha_i=dz_i\wedge\beta$を得る。
従って$u=\iota(\beta)$である。
\end{proof}


\subsection{一般三次における被約化と未被約構造}
ここで既存稿\cite[系\texttt{cor:pg-generic}]{NCCubicA}を明示的な入力とする：
一般の滑らかな三次式$F$について、全ての非零Poisson構造は一般階数二であり、
\eqref{eq:nfold-bivector}の下で分解可能な三重ベクトルに対応する。
これは既存稿の全次元正規形と高階数層の次元評価の帰結である。
したがって、ある非空Zariski開集合上で
\begin{equation}\label{eq:nfold-reduced-cone}
 \sqrt{I_F}=J_N
\end{equation}
となる。以下では、この被約化の主張からは決まらないスキーム構造を調べる。

\begin{theorem}[八次元以上のJacobiスキームの非被約性]
\label{thm:nfold-high-dimensional-nilpotents}
$n\ge8$、すなわち$N\ge10$とする。
任意の滑らかな三次式$F$に対し$I_F\ne J_N$であり、
頂点での完備局所環も$\widehat G_N$の頂点完備化と同型ではない。
さらに\eqref{eq:nfold-reduced-cone}を満たす一般の$F$について、
$\operatorname{Pois}(X_F)$は非被約であり、その被約化は$\widehat G_N$である。
冪零根基の次数二部分は少なくとも
\begin{equation}\label{eq:nfold-nilpotent-lower-bound}
 \epsilon_N:=5\binom{N+1}{6}-N\binom N4+\binom N3
 =\frac{(N+1)N(N-1)(N-2)(N-4)(N-9)}{144}>0
\end{equation}
次元を持つ。
特に$n=8$では$\epsilon_{10}=330$である。
\end{theorem}
\begin{proof}
補題\ref{lem:nfold-universal-row-relations}と\ref{lem:nfold-plucker-count}から
\[
 \dim (J_N)_2-\dim (I_F)_2\ge\epsilon_N>0
\]
となり、$I_F\subsetneq J_N$である。
\eqref{eq:nfold-reduced-cone}の下では$J_N/I_F$が冪零根基であり、
その二次部分は少なくとも$\epsilon_N$次元である。

頂点の極大イデアルを$\mathfrak m$とする。
$I_F$も$J_N$も二次斉次式で生成され、一次部分は零なので、
対応する局所環の$\mathfrak m/\mathfrak m^2$はともに次元$\binom N3$であるが、
$\mathfrak m^2/\mathfrak m^3$の次元には少なくとも$\epsilon_N$の差がある。
この数は局所環の抽象同型および完備化で不変だから、完備局所環は同型でない。
冪零根基の次数二の元が局所化で消えないことも、定数項が非零の多項式を掛けて
$I_F$に入るなら、その二次成分を比較して元の二次式が$I_F$に入ることから分かる。
\end{proof}

\begin{remark}[被約な分類と形式変形]
この定理は、一般の高次元三次でも非零Poisson構造の複素点が
Grassmann錐をなすことと両立する。
相違は非被約係数を許したJacobiスキームに現れる。
ここから全Abel圏変形の形式環との同一視を結論してはいない。
その比較には、相対formality、非可換模型の平坦性、および降下の独立した証明が必要である。
少なくとも、被約なPoisson分類だけから全Artin変形を
Grassmann錐の完備化に置き換えることはできない。
\end{remark}

\subsection{有限次元の最大階数を与える計算証明書}
\begin{theorem}[八次元以下の一般Jacobi行列の階数]
\label{thm:nfold-low-dimensional-generic}
$3\le n\le8$、すなわち$5\le N\le10$とする。
滑らかな三次式のパラメータ空間には非空Zariski開集合が存在し、
その上で
\[
 \dim(I_F)_2=min\left\{5\binom{N+1}{6},\,
                  N\binom N4-\binom N3\right\}
\]
が成り立つ。従って$3\le n\le7$の一般の$X_F$のPoissonスキームは、被約化を取る前から
$\operatorname{Gr}(3,N)$のアフィン錐である。
この範囲では、局所的に三次式の代表を選んだ相対Poissonスキームも
この開集合上でGrassmann錐との積になる。
$n=8$では一般に$\dim(I_F)_2=1980$、
$\dim(J_{10}/I_F)_2=330$である。
また$N=9,10$では一般の$F$について
$\ker\Phi_{T_F}=\operatorname{im}\iota$である。
$n=3$で全ての滑らかな三次に対する等式は既知であり
\cite[Theorem~1, \S2.2]{Mayanskiy}、ここで新たに用いるのは$n=4,5,6,7$の
一般点のイデアル一致と$n=8$の不足次元の厳密な判定である。
\end{theorem}
\begin{proof}
補題\ref{lem:nfold-jacobi-scheme}から$(I_F)_2\subset(J_N)_2$である。
補題\ref{lem:nfold-universal-row-relations}も合わせると、
表示した最小値が階数の上界である。
従ってJacobi二次式の係数行列がこの上界を達成する三次テンソルが一つあれば十分である。
この階数条件は三次式の係数に関する非空Zariski開条件であり、
滑らかな三次式の非空開集合と交わる。
$N\le9$ではPl\"uckerイデアルが二次式で生成されることから、
二次部分の一致はイデアル全体の一致を与える。
$N=10$では$\dim(J_{10})_2=2310$から$1980$を引いて$330$を得る。
$N=9,10$の核の主張は、その次元が$\binom N3$に等しいことと
普遍的な分裂単射$\iota$から従う。

再現コード\path{computations/verify_jacobi_rank.py}と
証明書\path{computations/jacobi_rank_certificate.json}は、各$N$に対して
対称整数テンソル$T_{ijr}$を明示している。
対応する三次式は
$F_T=\frac16\sum_{i,j,r}T_{ijr}z_i z_j z_r$であり、Hessianの$z_r$係数は$T_{ijr}$である。
行は$(r,a<b<c<d)$、列は辞書式に並べたPl\"ucker変数の非順序二次単項式で添字付ける。
行列の各行は
\begin{equation}\label{eq:nfold-certificate-matrix}
 \sum_{i,j}T_{ijr}
 \bigl(p_{abi}p_{cdj}+p_{bci}p_{adj}+p_{cai}p_{bdj}\bigr)
\end{equation}
の係数であり、添字の重複を持つ$p$は零、その他は交代符号で増加順に直す。
証明書は小行列の行番号・列番号と、法$32003$での非零行列式を含む：
\begin{center}
\begin{tabular}{c|r|r|r}
$N$ & Jacobi二次式の本数 & 小行列の次数 & 行列式$\bmod 32003$\\
\hline
5 & 25 & 5 & 15512\\
6 & 90 & 35 & 14374\\
7 & 245 & 140 & 7973\\
8 & 560 & 420 & 508\\
9 & 1134 & 1050 & 13073\\
10 & 2100 & 1980 & 16095
\end{tabular}
\end{center}
法素数で非零の整数行列式は特性零でも非零である。
ここでは有限体での零を特性零の恒等式と見なす推論を使っていない。
証明書の生成と、指定された小行列の行列式の再計算は別の手順で実行する。

最後に$N\le9$で、ある小行列式が可逆な係数開集合では、
Jacobi二次式から$(J_N)_2$の固定基底を係数環上で解き出せる。
従って相対イデアルも$J_N$の係数拡大に等しく、相対積表示を得る。
\end{proof}

\subsection{非零分解可能点の接空間と頂点に集中する冪零構造}

非被約性が非零のPoisson構造の近傍にも残るかどうかは、二次方程式の
総数だけからは分からない。以下では滑らかな三次に対する接空間を計算し、
被約化がGrassmann錐となる場合には、余分なスキーム構造が頂点だけに
支持されることを示す。

$T_F\in\operatorname{Sym}^3W^*$を三階微分テンソル
$T_F(v,w,u)=\partial_v\partial_w\partial_uF$とする。
非零分解可能な$P\in\Lambda^3W$の三次元supportを$A$と書き、
\begin{equation}\label{eq:nfold-transverse-kernel}
 K_{F,A}:=\{w\in W\mid T_F(a,w,u)=0
                 \text{ for all }a\in A,\ u\in W\}
\end{equation}
と置く。これは$W\to A^*\otimes W^*$、
$w\mapsto T_F(-,w,-)|_{A\times W}$の核であり、$P$の尺度に依存しない。

\begin{theorem}[非零分解可能点の法方向]\label{thm:nfold-normal-tangent}
$F$が滑らかな三次超曲面を定めるとする。任意の非零分解可能な$P$について
$A\cap K_{F,A}=0$であり、自然な同型
\begin{equation}\label{eq:nfold-normal-tangent}
 \frac{T_P\operatorname{Pois}(X_F)}{T_P\widehat G_N}
       \simeq\Lambda^3K_{F,A}
\end{equation}
が成り立つ。右辺から左辺への写像は
$\Lambda^3K_{F,A}\subset\Lambda^3W$を接ベクトルとみなすことによる。
従って
\[
 \dim T_P\operatorname{Pois}(X_F)
 =3(N-3)+1+\binom{\dim K_{F,A}}3.
\]
さらに、$P$のあるZariski近傍で$I_F=J_N$となることと
$\dim K_{F,A}\le2$は同値である。
\end{theorem}
\begin{proof}
$0\ne w\in A\cap K_{F,A}$なら$T_F(w,w,u)=0$が全ての$u\in W$について
成り立つ。三次式の同次性から$dF(w)(u)=T_F(w,w,u)/2$であるので、
$[w]\in\mathbb P(W)$は$X_F$の特異点となる。これは滑らかさに反する。
従って$K:=K_{F,A}$を含む補空間$B$を選んで$W=A\oplus B$とできる。
$P=e_1\wedge e_2\wedge e_3$となる$A$の基底を選ぶ。

Grassmann錐の$P$での接空間は
$\Lambda^3A\oplus(\Lambda^2A\otimes B)$である。
$I_F\subset J_N$だから、この接空間は全てのJacobi方程式の微分で消える。
従って法方向の代表は一意に
\[
 Q=\sum_{a=1}^3e_a\wedge\beta_a+\gamma,
 \qquad\beta_a\in\Lambda^2B,\quad\gamma\in\Lambda^3B
\]
と書ける。$B$の基底を$e_4,\ldots,e_N$とし、$b<c$を$B$の添字とする。
$\{a,a',a''\}=\{1,2,3\}$について、Jacobi二次式の
$(a',a'',b,c)$成分を$P$で微分すると、非零の符号を除いて
\begin{equation}\label{eq:nfold-normal-equations}
 T_F(e_a,q_{bc},e_r)=0\quad(1\le r\le N),\qquad
 q_{bc}:=\sum_{j=1}^3\beta_j^{bc}e_j
             +\sum_{j=4}^N\gamma^{bcj}e_j
\end{equation}
を得る。実際、\eqref{eq:nfold-certificate-matrix}の積のうち、
$P$の成分が非零となるのは二つの$A$添字を持つ因子だけである。
従って\eqref{eq:nfold-normal-equations}は$q_{bc}\in K$を意味する。
$K\subset B$なので、まず全ての$\beta_j^{bc}$が零となる。
残る条件は$\gamma$の全ての二重縮約が$K$に属することであり、
これは$\gamma\in\Lambda^3K$と同値である。
最後の同値は$B=K\oplus B'$に沿う外積の分解で、$B'$を含む成分を
二重縮約すれば$B'$方向のベクトルが得られることから従う。

逆に$Q=\gamma\in\Lambda^3K$なら、上の$(2A,2B)$成分の微分は零である。
残りの可能な$(3A,1B)$成分の微分は
$\sum_{a=1}^3\sum_{j\in B}T_F(e_a,e_j,e_r)\beta_a^{bj}$であり、
$\beta=0$なので零である。$A$添字が一つ以下の成分の微分は、
$P$の非零成分を含めないので初めから零である。
これで\eqref{eq:nfold-normal-tangent}を得る。
$\Lambda^3K$の接空間商への写像は補空間$B$を選ばずに定まるので、
同型は自然である。

$\dim K\le2$なら、Jacobi零スキームとGrassmann錐の接空間は一致する。
$P\ne0$ではGrassmann錐は滑らかである。
周囲の正則局所環を$\mathcal R$、その極大イデアルを$\mathfrak m$とし、
$I=I_F\mathcal R\subset J=J_N\mathcal R$と置く。
滑らかな埋込みなので$J/\mathfrak mJ\to\mathfrak m/\mathfrak m^2$は単射である。
接空間の一致から$I$の像が$J/\mathfrak mJ$を生成し、
$J=I+\mathfrak mJ$を得る。有限生成加群$J/I$にNakayamaの補題を適用して
$J=I$となる。有限生成性により、この等式は$P$のあるZariski近傍で成り立つ。
逆に局所イデアルが等しければ接空間も等しいため、
$\Lambda^3K=0$、すなわち$\dim K\le2$である。
\end{proof}

\begin{lemma}[非零層の横断性を与える係数開集合]
\label{lem:nfold-generic-transverse}
$N\ge6$とする。
滑らかな三次式のうち、ある三次元部分空間$A\subset W$に対して
$\dim K_{F,A}\ge3$となるものは、三次式の係数空間内で
余次元少なくとも$3N$の閉部分集合の中に含まれる。
従って、滑らかな三次式の非空Zariski開集合上では、
全ての非零分解可能な$P$の近傍で$I_F=J_N$となる。
この局所的な結論には、一般三次の被約Poisson分類を用いない。
$N=5$では同じ局所的な結論が全ての滑らかな三次に対して成り立つ。
\end{lemma}
\begin{proof}
$N\ge6$とする。交わらない三次元部分空間の対$(A,L)$の空間は
$\operatorname{Gr}(3,W)^2$の開部分であり、次元は$6(N-3)$である。
一つの対を固定して$W=A\oplus L\oplus C$と分解する。
条件$T_F(A,L,W)=0$は、$F$の単項式のうち、
$A$と$L$の座標を少なくとも一つずつ含むものの係数が零であることを意味する。
これらは型$(2A,1L)$、$(1A,2L)$、$(1A,1L,1C)$に分かれ、
独立な線形条件の数は
\[
 \binom42\,3+3\binom42+3\cdot3\,(N-6)=9N-18
\]
である。従って、この条件を満たす$(F,A,L)$のincidenceの次元は
\[
 \dim\operatorname{Sym}^3W^*-(9N-18)+(6N-18)
 =\dim\operatorname{Sym}^3W^*-3N
\]
である。その係数空間への像の閉包の余次元も少なくとも$3N$となる。

$F$が滑らかで$\dim K_{F,A}\ge3$なら、
$K_{F,A}$内に三次元部分空間$L$を選べる。
定理\ref{thm:nfold-normal-tangent}により$A\cap L=0$であり、
$T_F(A,L,W)=0$だから、$F$は上のincidenceの像に属する。
像の閉包の補集合を滑らかな三次式の開集合と交わらせれば、
求める非空開集合を得る。そこでの局所イデアルの一致は
定理\ref{thm:nfold-normal-tangent}による。
$N=5$では$A\cap K_{F,A}=0$だけで$\dim K_{F,A}\le2$となる。
\end{proof}

\begin{remark}
このincidenceの議論は、Grassmann錐上の全非零点での局所構造を制御する。
Grassmann錐から離れたPoisson構造の不在は、この議論だけからは結論しない。
次の系では、その点に既存の被約分類を用いる。
\end{remark}

\begin{corollary}[冪零構造の支持]\label{cor:nfold-vertex-support}
$F$は滑らかで、さらに$\sqrt{I_F}=J_N$と仮定する。
このとき全ての非零分解可能な$P$について$\dim K_{F,A}\le2$である。
従って
\begin{equation}\label{eq:nfold-punctured-equality}
 \operatorname{Pois}(X_F)\setminus\{0\}
       =\widehat G_N\setminus\{0\}
\end{equation}
はスキームとしての等式であり、
\begin{equation}\label{eq:nfold-projective-equality}
 \operatorname{Proj}(R_N/I_F)\simeq\operatorname{Gr}(3,W).
\end{equation}
冪零根基$\mathcal N_F=J_N/I_F$は有限長加群である。
特に、ある整数$s$について
\begin{equation}\label{eq:nfold-nilradical-annihilator}
 \mathfrak m_0^s\mathcal N_F=0,
 \qquad \mathfrak m_0=(p_{ijk})\subset R_N
\end{equation}
となり、$I_F$の斉次極大イデアルに関する飽和は$J_N$に等しい。
これらの結論は全ての$n\ge3$の一般の滑らかな三次に適用できる。
$n\ge8$では定理\ref{thm:nfold-high-dimensional-nilpotents}により
$\mathcal N_F\ne0$なので、その支持は正確に頂点である。
\end{corollary}
\begin{proof}
ある$P$について$\dim K_{F,A}\ge3$と仮定し、
$K_{F,A}$内の三次元部分空間をsupportとする非零分解可能な
$Q\in\Lambda^3K_{F,A}$を選ぶ。
$P$と$Q$はそれぞれ全てのJacobi二次式を満たす。
その混合項は全て$T_F(A,K_{F,A},W)=0$により零となるので、
$P+tQ$も任意の$t\in\mathbb C$についてJacobi方程式を満たす。
$A\cap K_{F,A}=0$であり、$t\ne0$なら$P+tQ$のsupportは六次元である。
従って$P+tQ$は分解可能でなく、$\sqrt{I_F}=J_N$に反する。

定理\ref{thm:nfold-normal-tangent}により、$\widehat G_N$の全ての
非零点の近傍で$I_F=J_N$となる。
被約化の仮定により他の点は存在しないので、
\eqref{eq:nfold-punctured-equality}と\eqref{eq:nfold-projective-equality}を得る。
$\mathcal N_F$は有限生成加群で支持が頂点に含まれるため、有限長であり、
ある$s$について\eqref{eq:nfold-nilradical-annihilator}が成り立つ。
$I_F\subset J_N$、$J_N$の飽和性、およびこの包含から
$(I_F:\mathfrak m_0^\infty)=J_N$が従う。
一般の三次への適用には、既存稿の被約分類
\eqref{eq:nfold-reduced-cone}を用いる。
\end{proof}

\begin{remark}[相対的な非零層]\label{rem:nfold-relative-punctured}
三次式の代表を選び、\eqref{eq:nfold-reduced-cone}が成立する非空係数開集合を
$U$とする。上の結論は各ファイバーだけでなく、
$U$上の相対Jacobiスキームの非零部分にも適用できる。
実際、$J_N$は固定されたGrassmannイデアルの係数拡大であり、
$R_N/J_N$もそのイデアルも係数環上平坦である。
従って相対商加群$J_N/I$のファイバーは$J_N/I_F$である。
各非零閉点でこのファイバーが零であることから、Nakayamaの補題により
相対商加群自体が零となる。従って相対非零Poissonスキームは
$U\times(\widehat G_N\setminus\{0\})$である。
この主張は任意の係数変更で保たれる。
三次式の代表の変更に伴う直線束の捻りを戻せば、
一般三次上の非零Poisson対の全スキームは滑らかであり、
その被約化だけを取り出す必要はない。
\end{remark}

\begin{corollary}[頂点の局所コホモロジー]\label{cor:nfold-embedded-vertex}
系\ref{cor:nfold-vertex-support}の仮定の下で
\[
 H^0_{\mathfrak m_0}(R_N/I_F)=J_N/I_F
\]
が成り立つ。特に$n\ge8$の一般の滑らかな三次について、
$\mathfrak m_0$は$R_N/I_F$の埋め込まれた随伴素イデアルであり、
\[
 \operatorname{Ass}_{R_N}(R_N/I_F)=\{J_N,\mathfrak m_0\},\qquad
 \operatorname{depth}\bigl((R_N/I_F)_{\mathfrak m_0}\bigr)=0.
\]
従ってJacobiスキームは頂点でCohen--Macaulayではない。
\end{corollary}
\begin{proof}
$J_N/I_F$は有限長なので全て$\mathfrak m_0$-torsionである。
一方、$R_N/J_N$は正の次元を持つ整域であり、
$H^0_{\mathfrak m_0}(R_N/J_N)=0$である。
完全列
$0\to J_N/I_F\to R_N/I_F\to R_N/J_N\to0$から最初の等式を得る。
$n\ge8$では$J_N/I_F\ne0$なので、有限長加群の非零socleに属する元は
ちょうど$\mathfrak m_0$で消される。
従って$\mathfrak m_0$は随伴素イデアルで、頂点の深さは零である。
唯一の極小素イデアルは$J_N$であり、
$\dim(R_N/J_N)=3(N-3)+1>0$だから、これは埋め込まれた素イデアルである。
上の完全列から随伴素イデアルは$J_N$と$\mathfrak m_0$以外には存在しない。
同じ次元の正値性からCohen--Macaulay性も成り立たない。
\end{proof}

\subsection{高次数の下界と冪零部分の次数幅}

頂点に支持されることは、冪零部分が次数二だけに集中することを意味しない。
この点はGrassmann環のHilbert関数と生成二次式の数から定量化できる。
$m_N=\binom N3$、$a_N=N\binom N4-\binom N3$と置き、$d\ge2$について
\begin{align}
 h_N(d)&:=\dim(R_N/J_N)_d
 =\frac{2\prod_{i=1}^3\prod_{j=1}^d(N+j-i)}{d!(d+1)!(d+2)!},
 \label{eq:nfold-grassmann-hilbert}\\
 L_{N,d}&:=\binom{m_N+d-1}{d}-h_N(d)
            -a_N\binom{m_N+d-3}{d-2}
 \label{eq:nfold-higher-defect}
\end{align}
と定める。

\begin{proposition}[高次数の差と次数幅の非一様性]
\label{prop:nfold-higher-nilradical}
任意の滑らかな$F$と$d\ge2$について
\[
 \dim(J_N/I_F)_d\ge L_{N,d}
\]
である。特に$N\ge13$では$L_{N,3}>0$であり、$N=13$では
\[
 \dim(J_{13}/I_F)_3\ge179179.
\]
さらに、各固定した$d\ge2$について、十分大きな$N$では$L_{N,d}>0$となる。
従って、一般三次の有限長冪零根基の非零成分が現れる最高次数、
およびそのsocleが現れる最高次数は、$N$によらない定数では抑えられない。
また$\mathfrak m_0^s(J_N/I_F)=0$となる最小の$s$も一様には抑えられない。
この主張は冪零根基の冪零指数が非有界であることを意味しない。
\end{proposition}
\begin{proof}
Grassmann環の標準単項式基底\cite[\S2.4, Theorem~2.3]{BCMStandardMonomial}は、
次数$d$では形$(d,d,d)$の半標準Young表に対応する。
Weylの次元公式をこの長方形に適用すると、その分母は
\[
 \prod_{j=1}^d j(j+1)(j+2)
 =\frac{d!(d+1)!(d+2)!}{2}
\]
であり、分子は式\eqref{eq:nfold-grassmann-hilbert}の積になる。
従って
\[
 \dim(J_N)_d=\binom{m_N+d-1}{d}-h_N(d).
\]
$I_F$は二次式で生成されるので、乗法
$(R_N)_{d-2}\otimes(I_F)_2\twoheadrightarrow(I_F)_d$は全射である。
補題\ref{lem:nfold-universal-row-relations}により
$\dim(I_F)_2\le a_N$であり、これらを引けば下界を得る。

$d=3$では展開により
\[
 L_{N,3}=
 \frac{N(N-4)(N-2)(N-1)(N+1)}{25920}
 \bigl(17N^4-258N^3+631N^2-270N-360\bigr).
\]
括弧内を$N=t+13$で書くと
$17t^4+626t^3+7807t^2+34726t+21480$となる。
従って$N\ge13$で正であり、$N=13$を代入すると$179179$を得る。

一般の固定した$d\ge2$については
\[
 L_{N,d}=
 \left(\frac1{6^dd!}-\frac2{d!(d+1)!(d+2)!}\right)N^{3d}
       +O(N^{3d-1}).
\]
実際、$m_N=N^3/6+O(N^2)$であり、最後の生成式数による項は
$O(N^5N^{3d-6})=O(N^{3d-1})$である。
先頭係数は$(d+1)!(d+2)!>2\cdot6^d$により正である。
最後の不等式は$d=2$で成立し、$d$を一つ増やすと左辺の倍率が
$(d+2)(d+3)>6$となるため全ての$d\ge2$で成立する。
よって十分大きな$N$で$(J_N/I_F)_d\ne0$である。

一般の三次では$J_N/I_F$は系\ref{cor:nfold-vertex-support}の有限長冪零根基である。
その最高非零次数の成分は全ての正次数変数で消えるのでsocleに含まれる。
従って各$d$に対して最高socle次数が$d$以上となる次元が存在する。
さらに$J_N$が二次式で生成されることから、$J_N/I_F$も次数二で生成される
$R_N$-加群である。従って$(J_N/I_F)_d\ne0$なら
$\mathfrak m_0^{d-2}(J_N/I_F)\ne0$であり、極大イデアルによる消滅指数も非有界である。
ここでは座標環の正次数元による作用を調べており、
冪零根基自身の積の長さを決定してはいない。
\end{proof}

\subsection{非零階数二層と特殊な高階数構造}
$M_n=\binom{n+4}{3}-(n+2)^2$を滑らかな三次$n$-foldのモジュライの次元とする。
三次式と非零分解可能な$P$の対のモジュライは
\begin{equation}\label{eq:nfold-rank-two-dimension}
 M_n+\dim\widehat G_N=M_n+3n-2
\end{equation}
次元の滑らかなDeligne--Mumford stackであり、三次モジュライへの忘却射は
相対次元$3n-2$の滑らかな全射である。
実際、三次式の係数開上では、これは
$\mathcal O(-1)$で捻ったGrassmann錐の頂点を除いた束である。
その相対次元は$3(N-3)+1=3n-2$であり、
$\operatorname{PGL}(W)$の商を取ると上式を得る。
滑らかな三次の射影自己同型群は有限なので商はDeligne--Mumfordである。
一般三次の係数開上では注意\ref{rem:nfold-relative-punctured}により、
これは非零Poisson対の全スキームを与え、被約化を取る必要はない。
この次元はPoisson対の底の次元であり、固定した量子化パラメータで得られる
非可換Abel圏の同型類の像の次元を意味しない。
また$n\ge8$では、一般のJacobiスキームの非被約性は零Poisson構造に集中する。

全ての滑らかな三次に対するGrassmann表示は、$n\ge4$では複素点の集合としても偽である。
例えば
\begin{equation}\label{eq:nfold-fermat-rank-four}
 F=z_1^3+\cdots+z_N^3,\qquad
 P=e_1\wedge e_2\wedge e_3+e_4\wedge e_5\wedge e_6
\end{equation}
は、二つの独立なJacobian括弧の和としてPoisson構造を与え、一般階数は四である。
$P$のsupportは六次元なので分解可能でない。
$N=6$ではEuler方向が像に入る条件が
$z_1^3+z_2^3+z_3^3=z_4^3+z_5^3+z_6^3=0$であり、$X_F$の真の閉集合をなす。
$N>6$では一般点のEuler方向が残りの座標成分を持つので、やはり射影化で階数は落ちない。
従って、一般の三次に関する定理と、特殊な三次の全Poisson構造の分類を区別する必要がある。

% !TEX root = ../nc-cubic-n-fold.tex
\section{Trace--Hessian 三次と残余復元三次の比較障害}
\label{sec:nfm-residual-th-obstruction}

本節では、Trace--Hessian 写像が与える三次と、
rotation を備えた残余圏から復元される三次を区別する。
固定入力論文Bの完備移送・一般点復元定理を用いる箇所を明示する。
同入力の全証明の独立再検証を行ったという意味ではない。
特に、可換特殊ファイバーとの一致や同じ一次変形類だけでは、
二つの復元を同一視できないことを、任意の次元の一ブロック族で示す。

\subsection{残余復元が選ぶデータ}

本稿の定理\ref{cft:thm:complete}は、次数別有限自由な完備中心三次対
$(B,\Omega)$ に対し、特殊ファイバー $F$ を持つ形式三次式 $W$ と、
内部 shift および指定三次周期を保つ graded MF 同値を構成する。
同入力の補題\texttt{cft:lem:model}, \texttt{cft:lem:center}で、
formality と星積模型への厳密同型 $\phi$ を選んだ後の $W$ は
\begin{equation}
\label{eq:nfm-formality-central-inverse}
 W=C_\pi^{-1}\phi(\Omega),\qquad
 C_\pi(f)=\sum_{j\ge0}\frac1{j!}
 \mathcal U_{j+1}(f,\pi,\ldots,\pi)
\end{equation}
として定まる。従って $W\bmod h=F$ ではあるが、$W=F$ は結論ではない。
この写像は有限次数内の形式級数であり、
量子関係式から直接回復する尺度付き二次 Poisson tensor と、
formality が構成する formal Poisson tensor とも区別する。

三次周期は $\deg t=(2,-3)$ の周期変数を係数として固定することにより保存される。
したがって、potential の単元倍を行ったときには周期の規格化も追跡する必要がある。
一方、rotation の次数一・二の積から復元するのは勾配空間であり、
指定周期だけから量子中心と古典三次の Trace--Hessian 同定が強制されるわけではない。

一般点での低次数不変量は
\[
 V_\rho=H^0\operatorname{RHom}_{\mathcal K^e}(\Delta,P_\rho),\qquad
 U_\rho=H^0\operatorname{RHom}_{\mathcal K^e}
 (\Delta,P_\rho\otimes_{\mathcal K}^{\mathbf L}P_\rho),
 \qquad
 Q_\rho=\ker\bigl(\operatorname{Sym}^2V_\rho\to U_\rho\bigr)
\]
である。固定入力の定理\texttt{gcd:thm:field-model}は、
\emph{滑らかな可換模型の幾何学的存在を得た後に}この勾配空間を積分し、
方程式の定義体を元の関数体へ降ろす。
有限体拡大が不要なのは方程式であり、同値核とその周期比較には有限拡大が残る。
この降下から、方程式が $h=0$ で正則であることや剰余が $F$ であることを
一般には結論できない。

\begin{proposition}[既知の可換模型がある場合の比較判定]
\label{prop:nfm-residual-low-product-test}
$k$ を標数零の体とし、$(\mathcal K,\rho)$ は smooth proper な dg pair で、
ある代数閉拡大体上で滑らかな三次超曲面の rotation 付き Kuznetsov pair に
同値であると仮定する。
滑らかな三次式 $H\in\operatorname{Sym}^3W$ に対して、次は同値である。
\begin{enumerate}
\item 代数閉包上で $(\mathcal K,\rho)\simeq(\operatorname{Ku}(V(H)),\rho_H)$ である。
\item 代数閉包上の線型同型 $V_\rho\simeq W$ が存在し、
$Q_\rho$ を $\langle\partial_iH\rangle$ に移す。
\end{enumerate}
成立する場合、同値はある有限代数拡大上に取れる。
元の幾何的模型との同値が指定三次周期を保っていれば、
必要な有限拡大と方程式の規格化の後、その周期も保存できる。
\end{proposition}

\begin{proof}
同値から低 rotation 積を比較する向きは、固定入力の
命題\texttt{prop:ej-geometric-low-degree}である。
逆に既存の可換模型 $Y$ を取ると、第二条件の下で $Y$ と $V(H)$ の勾配空間が一致する。
滑らかな三次式の勾配空間による射影的一意性
\textup{(}同入力の命題\texttt{prop:ej-projective-uniqueness}\textup{)}により、
両者は代数閉包上で射影同型となる。
既存の残余同値と合成すればよい。
座標変換によって方程式が $cH$ に移る場合は、$c$ の三乗根で
全座標を一様に調整して potential の厳密同型にする。
これにより内部 shift と指定 MF 周期を保てる。
座標変換と同値核は、固定入力の有限降下補題\texttt{rog:lem:descent}により
有限拡大上へ降下する。
なお幾何的模型の存在の仮定を除き、低積の次元や可積分性だけから
圏同値を結論したわけではない。
\end{proof}

\subsection{次元に一様な split Hesse 族}

$r=n-1\ge2$、$R=\mathbf C[h]_{(h)}$、$K=\mathbf C(h)$ とする。
$D_R$ を
\[
 [y,z]=h x^2+\frac{h^2}{2}(yz+zy),\qquad
 [z,x]=h y^2+\frac{h^2}{2}(zx+xz),\qquad
 [x,y]=h z^2+\frac{h^2}{2}(xy+yx)
\]
で定義し、$B_R=D_R[u_1,\ldots,u_r]$ と置く。
中心変数と混同しない係数を $a_h,b_h$ と書いて
\[
 \begin{split}
 S&=x^3+y^3+z^3,\qquad P=xyz,\qquad
 G=\frac13\sum_{i=1}^r u_i^3,\qquad F_0=S/3+G,\\
 a_h&=1+\frac{h^4}{2}+\frac{h^5}{6}-\frac{h^8}{48},\qquad
 b_h=\frac{(4-h^4)(12h+h^5-8h^2)}{48},\\
 w_h&=\frac{a_h}{3}S+b_h\operatorname{sym}(xyz)+G
 \end{split}
\]
とする。
固定入力論文Bの命題\texttt{gre:prop:input}の Sklyanin 中心公式と
次数別自由性の証明は active な三変数で完結している。
従って中心変数を任意個付加しても、$B_R$ は次数別有限自由で、
$w_h$ は中心三次元、剰余 $F_0$ は滑らかな三次式である。

\begin{proposition}[Trace--Hessian 出力の厳密計算]
\label{prop:nfm-th-hesse-exact}
上の族の一般点において、Trace--Hessian 写像の $-12$ 規格化は
\[
 H_h^{\mathrm{TH}}=\mathfrak d_\Pi(w_h)=F_0+hxyz
\]
を厳密に与える。
従ってこの射の当該 trait に沿う出力は $h=0$ まで正則に延長する。
これは Trace--Hessian 構成そのものを、中心次元が跳ぶ可換特殊点全体へ
延長したという主張ではない。
\end{proposition}

\begin{proof}
尺度付き二次 tensor は $\Pi=h\Pi_{g_h}$、$g_h=S/3+hP$ である。
最大可換代数商は
\[
 T=K[u_1,\ldots,u_r]\otimes_K
 K[x,y,z]/(x^2+hyz,\ y^2+hzx,\ z^2+hxy)
\]
となる。active な正次数元による乗法は active 次数を真に上げるので、
その trace は零である。従って $\operatorname{tr}(w_{h,\mathrm{ab}})=G$、
$s_{\mathrm{cl}}(\overline g_h)=g_h$ である。
Jacobi 線上で $x^3=y^3=z^3=-hP$ だから
\[
 \epsilon(\overline w_h)=(b_h-ha_h)[P]
 =-\frac23h^2(1+h^3)[P].
\]
一方 Hessian 同型は
\[
 \nu_\Pi(\overline g_h)=8h^2(1+h^3)[P]
\]
を与える。よって
$-12\nu_\Pi^{-1}\epsilon(\overline w_h)=\overline g_h$ であり、
trace の純中心成分 $G$ を加えると主張を得る。
\end{proof}

\begin{proposition}[指定周期付き残余復元の Hesse 係数]
\label{prop:nfm-residual-hesse-input}
固定入力論文Bの完備移送定理とその明示 Hesse 計算を用いると、
上の族の残余圏は有限拡大後に
\[
 H_h^{\mathrm{res}}=F_0+\lambda_{\mathrm{res}}(h)xyz,
 \qquad
 \lambda_{\mathrm{res}}(h)
 =\frac{h+h^2/3+h^5/12}{1-h^4/4}
\]
の三次超曲面の Kuznetsov 成分と、rotation および指定三次周期を保って同値になる。
特に
\[
 H_h^{\mathrm{res}}-H_h^{\mathrm{TH}}
 =\frac{h^2}{3}xyz+O(h^3).
\]
\end{proposition}

\begin{proof}
入力の命題\texttt{gre:prop:hesse-normalform}は、formality と中心逆像を
active な三変数だけで構成し、中心多項式環上へ自明に延長する。
この証明では中心変数が二個であることを用いていない。
正規化も active 座標だけを $A(h)^{-1/3}$ 倍するもので、全ての $u_i$、
底の $h$、周期変数 $t$ を固定する。
従って任意の $r$ で同じ active Hesse 係数を得る。

同入力の定理\texttt{gre:thm:exact-potential}は、この active MF 圏を
指定した中心商の相対楕円曲線の $\operatorname{Perf}$ と比較する。
分数体上の楕円曲線の導来同値から $j$ 不変量の一致を得て、
一次項が $h$ である分枝を選ぶことにより、表示した有理式を決定する。
二次係数は、同入力の命題\texttt{gre:prop:jet}でも、
生成元を固定する harmonic-angle formality と中心写像から独立に計算されている。
この二次計算では Weyl gauge の eye 項と中心写像の inner-wheel 項を区別しており、
中心元の可換化だけを用いているのではない。

任意の $N=r+3\ge5$ について、固定三対象への制限・Laurent 局所化・生成・有限降下は
同入力の一般次元定理を適用できる。
period を固定した完備比較と上の exact potential 座標変更を用いるため、
同値は $w_h$ と表示した規格化 $H_h^{\mathrm{res}}$ の三次周期を保つ。
分母を払った方程式を使う場合は、この周期の規格化を別途保持する。
\end{proof}

\subsection{有限 jet と幾何学的一般点の障害}

\begin{lemma}[同じ低次 jet を持つ三次族の最初の差]
\label{lem:nfm-first-difference-jet}
$F$ を滑らかな三次式とし、$H,W\in k[[h]][x_1,\ldots,x_N]_3$ が
$H\bmod h=W\bmod h=F$ および $H\equiv W\pmod{h^m}$ を満たすとする。
既に固定した低次の標識の下で、$h^{m+1}$ を法とする両族の
射影座標同型が存在するための障害は
\[
 \left[\frac{W-H}{h^m}\bmod h\right]\in\operatorname{Jac}(F)_3
\]
である。これが非零なら、剰余を恒等とする形式射影同型は存在しない。
\end{lemma}

\begin{proof}
方程式の単元倍は、その形式三乗根による scalar 座標変換に吸収する。
滑らかな三次超曲面には無限小射影自己同型がない。
従って低次で同じ方程式を保つ、剰余が恒等の座標変換は、
射影 scalar を除いて $h^m$ を法として恒等でなければならない。
$I+h^mM$ による最初の変化は $\sum_{i,j}M_{ij}x_jF_{x_i}$ であり、
これは Jacobian イデアルの次数三部分全体を動く。
Euler の公式によって $F$ の scalar 倍も同じ部分に含まれる。
これで一段の存在の必要十分条件と、非零の場合の障害を得る。
\end{proof}

\begin{theorem}[Trace--Hessian と残余復元の比較障害]
\label{thm:nfm-th-residual-obstruction}
上の split Hesse 族について、次が成立する。
\begin{enumerate}
\item 二つの出力の $h^2$ の差の class は
$\tfrac13[xyz]\ne0$ in $\operatorname{Jac}(F_0)_3$ である。
従って可換特殊ファイバーの標識と底のパラメータ $h$ を保つ形式射影同型は存在しない。
\item 二つの一般ファイバーは、共通の代数閉拡大体上でも射影同型でない。
\item 固定入力の残余復元定理を用いると、非可換側の残余圏と
$\operatorname{Ku}(V(H_h^{\mathrm{TH}}))$ の間に rotation を保つ dg Morita 同値は存在しない。
さらに $n\not\equiv1\pmod3$ では、rotation を指定しない dg Morita 同値も存在しない。
\end{enumerate}
\end{theorem}

\begin{proof}
$\operatorname{Jac}(F_0)=k[x,y,z,u_1,\ldots,u_r]/(x^2,y^2,z^2,u_1^2,\ldots,u_r^2)$
なので $xyz$ の class は非零である。第一項は前補題から従う。

一般点では $g_\lambda=S/3+\lambda xyz$ の symmetrizer は scalar だけである。
実際、Hessian の対称条件は $\lambda(\lambda^3-8)\ne0$ のとき
対角成分を等しくし、非対角成分を零にする。
具体的には各 Hessian 係数行列に対する上三角三成分を並べた
九元一次方程式の係数行列に、行列式 $\lambda^2(\lambda^3-8)^2$ の
八次小行列がある。恒等行列は核に入るので、核はちょうど scalar である。
小行列の指定と厳密検算を付属コードに収録する。
$\lambda=h$ と $\lambda=\lambda_{\mathrm{res}}(h)$ はともにこの条件を満たす。
この場合の分解の一意性は、一般的な分解定理なしに確認できる。
$H(x)=T_H(x,x,x)$ に対し
\[
 E(H)=\{L:T_H(Lu,v,w)=T_H(u,Lv,w)\ \text{for all }u,v,w\}
\]
と置く。この空間は線型座標変換によって共役される。
滑らかな三次の $T_H$ は第一引数について非退化である。
実際、$T_H(a,-,-)=0$ なる $a\ne0$ があれば、その方向の座標に方程式が依存せず、
$[a]$ が特異点となる。
この非退化性の下で $E(H)$ は可換な単位的部分代数になる。
$L,M\in E(H)$ に対し、三つの引数間で作用を移すと
\[
 \begin{aligned}
 T_H(LMu,v,w)&=T_H(u,MLv,w),\\
 T_H(LMu,v,w)&=T_H(Mu,v,Lw)
             =T_H(u,Mv,Lw)=T_H(u,LMv,w).
 \end{aligned}
\]
非退化性から $LM=ML$、さらに同じ等式から $LM\in E(H)$ を得る。
Hessian を $\operatorname{Hess}(H)=\sum_i x_iH_i$ と書けば、条件は
$L^{\mathsf t}H_i=H_iL$ である。
中心変数 $u_i$ の係数行列は $2E_{ii}$ なので、各中心行の非対角成分は零である。
active な係数行列との条件から、中心列の active 成分も零となる。
実際、それらは active な全 Hessian 係数行列の共通核に入るが、
その核は滑らかさにより零である。
active ブロックには前段の scalar 条件が適用できるから
\[
 E(H)=K I_3\oplus\bigoplus_{i=1}^r K E_{u_i u_i}
 \simeq K^{r+1}.
\]
この代数の原始冪等元のうち階数三のものは active ブロックへの射影ただ一つであり、
残りは全て階数一である。この計算は任意の体拡大後にも成立する。
従って全体の射影同型があれば、唯一の階数三の像を保ち、
active 平面三次の同型を与える。

Hesse 平面三次の $j$ 不変量は、
\cite[Section~2, equations~(4)--(5)]{NFMArtebaniDolgachev} の
Weierstrass 形に $(t_0,t_1)=(1,3\lambda)$ を入れると
\[
 j(\lambda)=\frac{27\lambda^3(8-\lambda^3)^3}{(1+\lambda^3)^3}
\]
であり、厳密有理式の展開は
\[
 j(\lambda_{\mathrm{res}}(h))-j(h)
 =13824h^4+4608h^5+512h^6+O(h^7)\ne0
\]
となる。従って任意の体拡大後も active 三次は同型でなく、第二項を得る。

rotation を保つ Kuznetsov pair の同値があれば、低 rotation 積から
二つの滑らかな三次の射影同型を得るので矛盾する。
$n\not\equiv1\pmod3$ では命題\ref{prop:nf-rotation}により rotation は
Serre 関手と cohomological shift の積から回復できる。
任意の dg Morita 同値はそれらを保つため、rotation の標識を外しても同じ矛盾を得る。
$n\equiv1\pmod3$ で rotation を外した同値の不存在は、この議論からは結論しない。
\end{proof}

\begin{corollary}[Casimir の選択だけでは比較を修復できない]
\label{cor:nfm-casimir-selection-obstruction}
上の split 族で、尺度付き周囲 tensor $\Pi=h\Pi_{g_h}$ を固定する。
その三次 Casimir $H$ が滑らかな超曲面を定めるなら、
$V(H)$ は $V(H_h^{\mathrm{res}})$ と幾何学的に射影同型でない。
従って Trace--Hessian 写像の非零の数値規格化や、
純中心三次成分の選択を替えるだけでは、この例の比較を修復できない。
\end{corollary}

\begin{proof}
三次 Casimir は $H=cg_h+G'(u)$ と書ける。
滑らかさから $c\ne0$ であり、$G'$ の勾配の共通零点も原点だけである。
二つの部分の対称三重線形形式はそれぞれ第一引数について非退化なので、
前証明と同じ mixed-block の条件から
$E(H)=E(cg_h)\times E(G')$ となる。
第一因子は体 $K$ なので、$G'$ のさらなる分解の有無にかかわらず、
active 部分への射影 $(1,0)$ は階数三の原始冪等元である。
もし $V(H)\simeq V(H_h^{\mathrm{res}})$ なら、
この冪等元は後者の唯一の階数三の原始冪等元へ移る。
従って $g_h$ と $g_{\lambda_{\mathrm{res}}}$ の平面三次は同型となり、
前定理の $j$ 不一致に反する。
\end{proof}

\begin{remark}[一般性と適用範囲]
この族の ambient は中心変数への連続自己同型を持つため、
周囲代数だけに有限安定化子を要求した $T^{\mathrm{fin}}$ には属さない。
従って上の定理だけから、その特定の開上で一般対象が比較不能であるとは結論しない。
一方、量子中心で切った pair の有限安定化子条件は ambient の条件とは異なる。
Trace--Hessian 出力が滑らかなら、pair の安定化子はその滑らかな三次の
射影自己同型群へ単射するため有限である。
全同値比較と相対双対入力を保ち、ambient の有限性を要求しない
定理\ref{thm:poisson-good-open-isomorphism}の pair モジュライを用いればよい。

split 族でも、$T$ 上の $w_{h,\mathrm{ab}}$ の乗法は
$G$ 倍と active な次数を上げる冪零作用の和である。
従ってその determinant は $G^8\ne0$ で、$w_{h,\mathrm{ab}}$ は非零因子となる。
よって定常点スキームの完全交叉条件は満たす。
定常点スキーム自体の滑らかさは主張せず、上の比較にも必要ない。
PBW 性は $D_R$ の次数別自由性と中心多項式変数の付加から得られる。
点曲線の平行移動切断は
$p=(1-h^2/2:-1-h^2/2:-h)$ で、$h=0$ では原点 $e=(1:-1:0)$ に特殊化するが、
形式近傍では $p\ne e$ である。
標数零の楕円形式群の乗算 $[m]$ は線型項 $m$ を持つ可逆な形式写像なので、
任意の整数 $m>0$ について $[m]p\ne e$ である。
従って一般点の平行移動は無限位数であり、必要な有限個の位数除外を満たす。
Sklyanin の中心入力と既存の相対双対定理を適用する範囲で、
中心条件・相対双対性・$\epsilon$ の非消滅を確認できる。

混合 potential の一般点における、特殊ファイバーの標識を保つ形式比較は
次小節の有限 jet で判定できる。
一般ファイバー上で標識を外した圏同値の不存在は、
その有限 jet の計算だけからは結論しない。
\end{remark}

\subsection{混合係数の一般点における二次 jet 障害}

以下の計算では入力論文Aの Hesse 中心公式と、入力論文Bの
harmonic-angle formality による完備 MF 移送を使う。
後者の存在定理を任意の係数環へ無条件に拡張するのではなく、
各複素点における形式模型の二次 jet を同じ有限式で計算する。
その有限式自体は係数空間上の正則切断となる。

$U=\langle u_1,\ldots,u_r\rangle$、$r=n-1\ge2$ とし、
次の全混合 Hesse 係数空間を考える。
\[
 F=g_\lambda+Axy+Byz+Czx+Dx+Ey+Hz+G,\qquad
 g_\lambda=(x^3+y^3+z^3)/3+\lambda xyz,
\]
ただし $A,B,C\in U$、$D,E,H\in\Sym^2U$、$G\in\Sym^3U$ である。
$F$ と $g_\lambda$ がともに滑らかな三次を定める開を $\mathcal S_n$ とする。
$\mathscr H_F=\det\operatorname{Hess}_{x,y,z}(F)$ と書き、
\[
 T_F=\operatorname{Sym}U[x,y,z]/(F_x,F_y,F_z),\qquad
 \tau_F=\tfrac18\operatorname{Tr}_{T_F/\operatorname{Sym}U}
\]
と置く。$\tau_F$ は商への写像との合成として多項式にも適用する。
これは全混合係数と両立する $\operatorname{Sym}U$-線型写像である。
量子関係は $[a,b]=h\operatorname{sym}\Pi_F(a,b)$ とし、
入力の短い中心元から $h^2ABC/6$ を引いた規格化を
\[
 \widetilde\Omega_{F,h}
 =\operatorname{sym}\left(F-\frac{h^2}{12}\mathscr H_F
                              +h^4\Psi_\lambda(F)\right)
\]
とする。中心性・次数別 PBW 性は固定入力の該当開で用いる。

\begin{proposition}[残余復元と trace 像の普遍二次 jet]
\label{prop:nfm-relative-residual-second-jet}
Kontsevich の harmonic-angle formality と、次数一生成元を固定する
厳密同型 $\phi$ を選ぶ。
入力の中心逆像 $W^{\mathrm{res}}=C_\pi^{-1}\phi(\widetilde\Omega_{F,h})$ は
\begin{equation}\label{eq:nfm-residual-universal-jet}
 W^{\mathrm{res}}=F+\frac{h^2}{24}\mathscr H_F+O(h^3)
\end{equation}
を満たす。同じ中心元の trace--Hessian 像は
\begin{equation}\label{eq:nfm-th-universal-jet}
 H^{\mathrm{TH}}
 =F-\frac{h^2}{12}\tau_F(\mathscr H_F)
      +h^4\tau_F(\Psi_\lambda(F))
\end{equation}
である。この式は $h\ne0$ で定義した射の出力が当該 trait に沿って正則に延びたもので、
可換特殊点での最大可換商を $T_F$ と同一視するものではない。
従って差の最初の候補は
\begin{equation}\label{eq:nfm-relative-comparison-class}
 \Delta_2(F)=
 \left[\frac1{24}\mathscr H_F+
             \frac1{12}\tau_F(\mathscr H_F)\right]
 \in (J_F)_3
\end{equation}
となる。この類は $\mathcal S_n$ 上の Jacobi 束の正則切断であり、
その形成は任意の係数変更と両立する。
各 $n\ge3$ について非消滅開 $\mathcal S_n^{\mathrm{diff}}$ は非空である。
\end{proposition}

\begin{proof}
Weyl 対称化を $\sigma$ と書く。
$h\mapsto-h$ と語順の反転は反対代数を同定するので、
対称化で戻した積の一次項は $\frac12\Pi_F$ である。
入力論文Bの補題\texttt{cft:lem:model}の逐次構成を
$h^2$ を法とする恒等写像から始め、次数一成分を線型座標変更で調整すると、
$\phi\sigma=1+O(h^2)$ を保って全生成元を固定できる。
ここで $P_0^{ij}=\Pi_F(z_i,z_j)$ とし、添字は全 $N=n+2$ 変数を走る。
Dito の一般 Poisson tensor に対する二次積
\cite[Section~2, equation~(3)]{NFMDito}は、$\pi=P_0/2$ という規約で
\[
 \begin{aligned}
 B_2(P_0;f,g)={}&\frac18P_0^{ij}P_0^{kl}f_{ik}g_{jl}\\
 &+\frac1{12}P_0^{ij}(\partial_iP_0^{kl})
                  (f_{jl}g_k+f_kg_{jl})
   -\frac1{24}\Theta^{ik}f_i g_k
 \end{aligned}
\]
である。線型形式 $X=a_i z_i$ の冪と $X$ に適用すると、第一項は二階微分、
第二項は反対称性により消え、二-cycle 項だけが残る。
入力の命題\texttt{gre:prop:jet}と同じ polarization により
\[
 \phi\sigma=1+h^2\rho_2+O(h^3),\qquad
 \rho_2=-\frac1{48}\Theta^{ik}\partial_i\partial_k,
 \qquad
 \Theta^{ik}=(\partial_lP_0^{ij})(\partial_jP_0^{kl})
\]
を得る。
実際、$X=a_i z_i$ の星積冪の二次項では二-cycle 項だけが残り、
$-r(r-1)X^{r-2}\Theta^{ik}a_i a_k/48$ となる。
二次積は対称なので、量子生成元の関係との比較から
formal Poisson tensor の $h^2$ 係数も零である。

$P_0$ の非零成分は active な三添字だけを持つ。
$M=(F_{ab})_{a,b=x,y,z}$ と置けば
$P_0^{ij}=\varepsilon^{ija}F_a$ であり、縮約恒等式
\[
 \Theta^{ik}
 =\varepsilon^{ija}\varepsilon^{klb}M_{al}M_{bj}
 =-2\operatorname{adj}(M)_{ik}
\]
を得る。その他の成分は零である。
従って、$M\operatorname{adj}(M)=(\det M)I_3$ より
\[
 \rho_2F=\tfrac1{24}\operatorname{tr}(\operatorname{adj}(M)M)
        =\tfrac18\det M=\tfrac18\mathscr H_F.
\]
これは混合係数を含む対称三次式についての多項式恒等式である。

中心写像について、$\mathcal U_2(f,q)=0$ は二本の辺の反対称性による。
$\mathcal U_3(f,P_0,P_0)$ に寄与し得るグラフは inner 二輪だけであり、
その harmonic-angle 重みは零である
\cite[Section~1.3, Theorem]{NFMShoikhet}。
ここではグラフの重みの消滅だけを使い、Poisson tensor の線形性を仮定しない。
逆向き wheel の非零重みとは区別する\cite[Section~2.3.1]{NFMShoikhet}。
よって $C_\pi=1+O(h^3)$ であり、
\[
 C_\pi^{-1}\phi(\widetilde\Omega_{F,h})
  =F+h^2(-\mathscr H_F/12+\rho_2F)+O(h^3)
\]
から\eqref{eq:nfm-residual-universal-jet}を得る。
生成元を固定すること、eye 項を含む星積と同じ formality の中心写像を使うことが
この係数の規約を定めている。

trace 側では、入力の有限 Hesse 公式により
\[
 \Psi_\lambda(F)|_{U=0}
       =\lambda(8-\lambda^3)g_\lambda/48
\]
である。右辺は境界 Jacobi 環で零なので、
$\epsilon(\widetilde\Omega_{F,h})=-h^2[\det\operatorname{Hess}(g_\lambda)]/12$
は厳密等式である。
$\nu_{h\Pi_F}$ が $h^2$ を含むことによる高次項の繰り下がりは、
この規格化では生じない。
従って trace--Hessian 式の非純中心方向の係数は厳密に一であり、
\eqref{eq:nfm-th-universal-jet}を得る。
一般の $F-h^2\mathscr H_F/12+O(h^4)$ という jet だけから
この結論を出したわけではない。

有限局所自由な臨界代数の trace は係数空間上正則である。
滑らかな $F$ の勾配 Koszul 複体は次数別に普遍的に完全であり、
$(J_F)_3$ は階数 $\binom{n+2}{3}$ の局所自由束となる。
従って\eqref{eq:nfm-relative-comparison-class}は任意の係数変更と両立する正則切断である。
Fermat 点
\[
 F_*=(x^3+y^3+z^3+u_1^3+\cdots+u_r^3)/3
\]
では $\mathscr H_{F_*}=8xyz$、$\tau_{F_*}(\mathscr H_{F_*})=0$ であり、
\[
 \Delta_2(F_*)=\tfrac13[xyz]\ne0,
 \qquad J_{F_*}=k[x,y,z,u]/(x^2,y^2,z^2,u_1^2,\ldots,u_r^2).
\]
これで非消滅開の非空性を得る。
\end{proof}

\begin{corollary}[一般混合係数での標識付き形式比較の障害]
\label{cor:nfm-generic-marked-comparison-obstruction}
$F\in\mathcal S_n^{\mathrm{diff}}(\C)$ を、前命題の中心量子化・次数別 PBW・
指定 formality 模型の入力が成立する開の複素点とし、特殊ファイバーを $V(F)$ と同定する。
同じ完備 MF 入力が適用できる係数体についても、同じ主張が成立する。
二つの三次族 $V(W^{\mathrm{res}})$ と $V(H^{\mathrm{TH}})$ は、
既に $\C[h]/(h^3)$ 上で、この同定を保つ同型を持たない。
固定入力の完備 MF 移送と相対 window 同定を適用する場合、
量子中心商の残余圏と $\operatorname{Ku}(V(H^{\mathrm{TH}}))$ の間にも、
同じ $h$ を固定し特殊ファイバー上恒等となる
$\C[[h]]$-線型 dg Morita 同値は存在しない。
この第二の主張は $n=4$ を含む全 $n\ge3$ に対する、入力付きの有限 jet 障害である。
\end{corollary}

\begin{proof}
両族は $h^2$ を法として $V(F)$ の定数族であり、
小拡大 $\C[h]/(h^3)\to\C[h]/(h^2)$ 上の持上げの差は
$h^2\Delta_2(F)$ である。
$H^0(V(F),T_{V(F)})=0$ なので、特殊ファイバー上恒等な同型は
$h^2$ を法としても恒等である。
そのような同型があれば差は $(J_F)_3$ で零となり、仮定に反する。

圏については、命題\ref{prop:nf-hh}の実際の幾何的 Kodaira--Spencer 写像
\[
 \eta_F:(J_F)_3\hookrightarrow HH^2(\operatorname{Ku}(V(F)))
\]
が全ての次元で単射であり、$HH^1=0$ である。
従って恒等特殊ファイバーを持つ圏同値の一次の曖昧さはなく、
二つの幾何的残余変形の持上げの差は非零の
$h^2\eta_F(\Delta_2(F))$ である。
入力の完備 MF 同値は剰余で恒等であり、各 Artin 切断と両立するので、
量子残余変形にも同じ結論が移る。
ここで全 Artin 形式認識の全射性は用いず、
四次元でも成立する幾何的接写像の単射性だけを使う。
\end{proof}

\begin{remark}[一般性と依存範囲]
全混合 Hesse 係数空間は既約であり、Hesse 座標と中心平行移動は
一般の一ブロック表示を覆う。
従って上の非消滅開は、既存の良い一ブロック入力が成立する非空開とも交わる。
ただしここで代数的に相対化したのは二次 jet とその Jacobi 類である。
全次数の残余射を係数空間上に構成したこと、任意の非被約係数環上の
完備 MF 移送を新たに証明したこと、あるいは一般ファイバー上の
標識を外した Kuznetsov 圏が非同値であることは主張しない。
特に、有限 jet 障害と先の split 族における有理式・$j$ 不変量による
幾何学的一般点の障害を区別する。
\end{remark}

% !TEX root = ../nc-cubic-n-fold.tex
\section{残余三次の代数的復元と射影 Poisson 射}
\label{sec:nfm-residual-algebraic}

本節では、残余圏からの三次の復元と、固定生成元での形式 tensor の代数性を分ける。
方程式の直線は低 rotation 積から有限行列の核として復元できる。
中心一次空間の作用と合わせると、Poisson tensor の直線も得られる。
この構成を族へ広げる箇所では、固定入力論文Bの完備 MF 移送・window 比較・
一般点の生成性と有限降下を用いる。これらの未公表入力は新しい定理の仮定として保持する。
他方、固定 harmonic-angle 模型の tensor 自身には、
任意の $n\ge3$ の very general な滑らかな係数で超越性の障害がある。
後者は既知の log-canonical 量子化と多項式再帰から導き、MF 比較を用いない。

\subsection{勾配空間を積分する方程式直線}

滑らかな三次の勾配空間による復元は既知である。
Wang\cite[Theorem~3.3]{NFMWangMilnor}は非 Sebastiani--Thom 部分の勾配射の
単射性と immersion を示し、Hwang\cite[Theorems~1.5, 1.6; Proposition~2.6]{NFMHwangSymmetrizer}
はその fibre を symmetrizer 群で記述する。
以下では今回必要な、方程式直線の任意底変更と線束 twist への両立を明示する。
$V$ は座標関数の空間、従って三次式は $\Sym^3V$ に属する。

\begin{lemma}[scalar symmetrizer の局所底上の保存]
\label{lem:nfm-scalar-symmetrizer-local}
$R$ を $\mathbb Q$-代数である局所環とし、
$F_R$ を $N$ 変数の三次式とする。
その剰余 $F_0$ が滑らかで scalar symmetrizer を持つなら、
$F_R$ の symmetrizer は $R\cdot\mathrm{id}$ であり、その形成は任意底変更と両立する。
勾配部分加群 $Q_{F_R}$ は階数 $N$ の局所直和因子で、
\[
 \{G\in\operatorname{Sym}^3_RV:
                  dG\in V\otimes_RQ_{F_R}\}=R F_R
\]
である。これらは $R=k[[h]]$、$F_R=W_h$ にも適用できる。
\end{lemma}

\begin{proof}
symmetrizer 条件は $N^2$ 個の行列成分に関する線形方程式である。
剰余で核の次元が一なので、$N^2-1$ 次の単元小行列式がある。
恒等行列は全ての底上で解であり、自由直和因子を生成する。
この直和因子を除いた $N^2-1$ 個の変数について、その小行列で消去すれば、
他の解がないことと任意底変更との両立が従う。

滑らかな $F_0$ の $N$ 個の偏微分は線形独立なので、
勾配写像の $N$ 次の小行列式の一つは単元である。
従ってその像 $Q_{F_R}$ は自由直和因子となる。
上記の $G$ に対して $\nabla G=A\nabla F_R$ となる $A$ が一意に存在し、
二階微分の対称性により $A$ は symmetrizer である。
従って $A=a\,\mathrm{id}$、Euler 等式から $G=aF_R$ となる。
\end{proof}

\begin{proposition}[勾配データからの相対的な方程式直線]
\label{prop:nfm-relative-gradient-line}
$V$ を階数 $N$ のベクトル束とし、
$Q\subset\operatorname{Sym}^2V$ を階数 $N$ の部分束とする。
$M=\binom{N+2}{3}$ と置き、自然な微分写像から
\[
 D_Q:\operatorname{Sym}^3V
 \longrightarrow V\otimes(\operatorname{Sym}^2V/Q)
\]
を定める。Grassmann 束上で $\operatorname{rank}D_Q=M-1$ となる
局所閉な determinantal 部分スキームを取り、
$\mathcal L=\ker D_Q$ とする。
さらに
\[
 V^\vee\otimes\mathcal L\longrightarrow Q
\]
が同型で、$Q$ の共通射影零点が存在しない開へ制限する。
この局所閉部分上で、$\mathcal L$ は三次方程式の直線束を定め、
その零スキームは滑らかな相対三次超曲面である。
構成は任意底変更と線型同型に自然である。

この部分スキームと、滑らかで scalar symmetrizer を持つ三次直線の開とは、
勾配空間を取る射と $\ker D_Q$ を取る射によって互いに逆に同型となる。
\end{proposition}

\begin{proof}
rank 条件は $M$ 次小行列式を零とし、
$M-1$ 次小行列式の少なくとも一つを可逆とする条件である。
その上では単元小行列による行列消去により
$\ker D_Q$ と $\operatorname{coker}D_Q$ は局所自由であり、
核の階数は一である。そのため核の形成は非平坦なものを含む任意底変更と両立する。
微分写像の同型条件は、復元した三次直線の偏微分がちょうど $Q$ を張ることを意味する。
$Q$ の共通射影零点がないことから、Jacobian 判定で三次族は滑らかである。
共通零点の存在条件が閉であることは、射影空間との incidence の固有射影から従う。

逆に滑らかな三次直線からの勾配写像は階数 $N$ の部分束を与える。
scalar symmetrizer 条件と前補題により、その $D_Q$ の核は元の直線そのものである。
局所的な生成元を替えても同じ直線であり、両合成は恒等となる。
この証明は非被約な試験スキームにもそのまま通る。
\end{proof}

この勾配射の単射性・immersion は既知であり、
Wang\cite[Theorem~3.3]{NFMWangMilnor}とその引用する先行結果に帰属する。
上の記述では、今回必要な核の任意底変更と局所底上の逆構成を有限行列で明示している。

\begin{lemma}[全次元の scalar symmetrizer を持つ滑らかな例]
\label{lem:nfm-uniform-scalar-symmetrizer}
$N=n+2\ge5$ とし、添字を $\mathbb Z/N\mathbb Z$ で読む。
三次式
\[
 F_{\mathrm{cyc}}=\sum_{i=1}^{N}x_i^2x_{i+1}
\]
は滑らかな射影超曲面を定め、その symmetrizer は scalar のみである。
さらに三次元部分空間
\[
 x_4=\cdots=x_N=x_1
\]
への制限も滑らかである。
従って、滑らかな全三次と滑らかな active 制限を持つ係数空間の中に、
symmetrizer が scalar のみとなる非空 Zariski 開がある。
\end{lemma}

\begin{proof}
全偏微分が零なら
$x_{i-1}^2+2x_ix_{i+1}=0$ である。
一つの座標が零なら直前の座標も零となり、全て零となる。
全て非零なら各式を掛けることで $1=(-2)^N$ を得るが、
これは標数零では成立しない。従って三次は滑らかである。

$A\operatorname{Hess}(F_{\mathrm{cyc}})$ が対称となる行列 $A$ を求める。
これは本文の symmetrizer の転置表示であり、scalar 性は同じである。
$\operatorname{Hess}(F_{\mathrm{cyc}})=\sum_kx_kH_k$ と書けば
\[
 H_k/2=E_{k-1,k-1}+E_{k,k+1}+E_{k+1,k}.
\]
$AH_k$ の非零列は $\{k-1,k,k+1\}$ に限られる。
対称性により同じ三つ以外の行も零なので、$A$ の第 $j$ 列の台は
\[
 \{j,j+1,j+2\}\cap\{j-2,j-1,j\}\cap\{j-1,j,j+1\}
 =\{j\}
\]
に含まれる。ここで $N\ge5$ を用いた。
従って $A$ は対角行列であり、$AH_k$ の $(k,k+1)$ 成分の対称性から
$A_{kk}=A_{k+1,k+1}$ となる。よって $A$ は scalar である。

表示した三次元切断への制限は
\[
 g_c=x^2y+y^2z+z^2x+cx^3,\qquad c=N-3.
\]
勾配の共通非零点では $x\ne0$ なので $x=1$ とできる。
後二式から $z=-y^2/2$、$y^3=1$ を得て、
第一式は $3c+9y/4=0$ となる。
従って特異となる条件は $c^3=-27/64$ であり、$c=N-3$ はこれを満たさない。
$i\ge4$ について $x_i=x_1+z_i$ と線型変数変換すれば、
この切断は固定された active 座標平面 $z_4=\cdots=z_N=0$ になる。
最後に、symmetrizer は係数に多項式的に依存する有限行列の核で、
恒等行列を必ず含む。核の次元が一であることは開条件であり、
上の例によってこの開が非空と分かる。
\end{proof}

\begin{remark}
この例は巡回置換という非自明な射影自己同型を持つ。
従って scalar symmetrizer と自明な射影自己同型群は異なる条件である。
以下の方程式直線の復元に後者は必要ない。
また、この係数開の非空性を全 intrinsic 圏成分の被覆とは読み替えない。
\end{remark}


\subsection{入力と相対有限模型}

$N=n+2\ge5$、$r=N-3$ とする。
入力Bは序論で指定した \texttt{nc-cubic} の固定 commit に統一する。
\texttt{generic\_recovery/field\_model.tex},
\texttt{generic\_descent.tex}, \texttt{window.tex},
\texttt{coherent\_saturation.tex} の各結果を区別する。
体上の三次の存在、幾何学的射影同型類の一意性、
有限体拡大後の比較核は既存入力である。
本節で行うのは、積分直線が一意な開における、
有限個の線型代数データによる正準的な相対構成である。
これらの相対構成・応用についての新規性の網羅的確認は未完了である。

表示 chart 上の既存の相対双対性は、
全ての $\mathcal O(i)$ の compact 性、proper Hom、基底変更を与える。
$B$ の Koszul 複体を $A=B/(\Omega)$ と tensor した複体の二つの Tor は
$R$ と $R(-3)$ であり、tails では消える。
従って連続する $N$ 個の twist
\[
 G=\bigoplus_{i=0}^{N-1}\mathcal O(i)
\]
は全ての twist を生成する。
同じ Hom 計算による相対 exceptional block
$\mathcal O,\ldots,\mathcal O(n-2)$ を有限回の mutation で射影して
残余圏の有限生成対象 $G_{\mathrm{res}}$ を得る。
有限回の mutation は proper Hom による評価 cone なので、係数変更と両立する。
以下の相対残余圏は compact 部分を指し、全 coherent 残余圏との一致は
各幾何ファイバーで後述の飽和性入力を使う。
その代数
\[
 \Gamma_{\mathrm{res}}=R\operatorname{End}(G_{\mathrm{res}})
\]
は proper で、係数変更と両立する。
この議論では $\Gamma_{\mathrm{res}}$ の smoothness を仮定しない。

rotation/window との同定には、正則な係数 chart 上で入力Bの (L) を使う。
一般点の線束生成残余圏と全 coherent 残余圏の一致には、
同入力の coherent saturation を使う。
任意の係数環で相対 saturation が証明されたことに置き換えてはならない。
以下の相対構成では、これらの代数的な有限模型と比較がある開を使う。
単に各幾何学的点で模型が存在することだけを仮定してはいない。

\subsection{支配的な中心三次の一変数族}

$V=X\oplus U$、$\dim X=3$、$\dim U=r$、
$T\subset\operatorname{Sym}^3V/\operatorname{Sym}^3U$ を
一ブロック周囲代数の良い非空開とする。
$F$ を $T$ の一般点での代表、$k=\mathbf C(T)$ とし、
$R=k[h]_{(h)}$ 上で
\[
 B_{hF}=T(V)_R/
 ([U,V],[y,z]-h\operatorname{sym}F_x,
 [z,x]-h\operatorname{sym}F_y,[x,y]-h\operatorname{sym}F_z)
\]
を考える。既存の PBW 入力を用い、必要な係数を反転する。

\begin{lemma}[中心行列による支配的な $h$ 族]
\label{lem:nfm-dominant-central-family}
一般 $F$ に対し、
$\Omega_h\in Z(B_{hF})_3$ を $\Omega_h\bmod h=F$ として代数的に選べる。
$G\in\operatorname{Sym}^3U$ の係数を独立に動かすと、
\[
 (F,G,h)\longmapsto(B_{hF},[\Omega_h+G])
\]
は $h\ne0$ 上で既存の中心三次 torsor $J$ を支配する。
\end{lemma}
\begin{proof}
次数 $3,4$ の PBW 加群は有限自由であり、
\[
 \delta_h:B_{hF,3}\longrightarrow B_{hF,4}^{\oplus N},
 \qquad b\longmapsto(h^{-1}[x_i,b])_i
\]
が定義される。剰余は Hamilton 微分である。
active 勾配 $F_x,F_y,F_z$ の正則列性から
\[
 \ker\delta_0=kF\oplus\operatorname{Sym}^3U.
\]
実際、$d_XG\wedge d_XF=0$ と Koszul 完全性から
$d_XG=a\,d_XF$、次数から $a\in k$ となる。

$(F,h)\mapsto hF$ は $h\ne0$ 上で ambient 係数空間を支配する。
従って既存の一般中心公式により、一般ファイバーの中心三次の次元も
$1+\dim\operatorname{Sym}^3U$ である。
$\delta_h$ の一般・特殊の階数は等しい。
剰余で非零の最大 minor を用いた行消去では、残る Schur 部分は
一般点で零、従って $R$ 上で零となる。
よって中心核は有限自由で、剰余に関する基底変更が成り立つ。
$F$ を持ち上げて $\Omega_h$ を得る。

$h\ne0$ では $F\mapsto hF$ は可逆であり、
固定した ambient 代数の中心三次直線は、
非中心商で $\Omega_h$ と同じ像を持つように一意に規格化すると
$\Omega_h+G$ と書ける。
従って $G$ は中心三次 torsor の全ファイバーを動く。
有理的にはこの係数族は $J\times\mathbf G_m$ と同型である。

\end{proof}

$G$ の係数を基礎体に加えると、
$\Omega_h+G\bmod h=F+G$ は一般の全変数三次である。
滑らかさ、scalar symmetrizer、および active 平面三次の滑らかさを同時に課せる。
各 $N\ge5$ で非空であることは、
補題\ref{lem:nfm-uniform-scalar-symmetrizer}の巡回三次の明示例から従う。

ここで入力Bの (L),(CF) を一般係数体 $k/\mathbf C$ 上へ拡張することを明記する。
Euler 恒等式は整数係数であり、有限核の降下は任意の標数零の体で成立する。
有限 Koszul 行列、固定内部次数の Hom、普遍 formality 恒等式は
係数変更と両立し、完備化は $h$ についてのみ行う。
生成性の段階で必要なら代数閉包へ移す。
複素閉点ごとの定理を単に並べて一般係数体上の結果を結論するのではない。

\subsection{低 rotation 積による相対三次}

\begin{theorem}[稠密開上の正準的な残余三次]
\label{thm:nfm-relative-residual-cubic}
$S$ を標数零の整 noetherian scheme とし、
proper dg 模型 $\Gamma$ と、可逆かつ底上 perfect な rotation の双加群 $P$ が
代数的に与えられるとする。
一般点 $\eta$ のある体拡大上で
$(\operatorname{Perf}\Gamma_\eta,\rho)$ が
滑らかな $n$ 次元三次の $(\operatorname{Ku},\rho)$ に同値であり、
その三次が幾何学的に ST 既約であるとする。
このとき稠密開 $S^\circ$ 上で
\[
 V=H^0R\operatorname{Hom}_{\Gamma^e}(\Gamma,P),\quad
 V_2=H^0R\operatorname{Hom}_{\Gamma^e}
       (\Gamma,P\otimes_\Gamma^{\mathbf L}P),\quad
 Q=\ker(\operatorname{Sym}^2V\longrightarrow V_2)
\]
は有限局所自由で、
\[
 \operatorname{rk}V=N,\quad \operatorname{rk}V_2=\binom N2,\quad
 \operatorname{rk}Q=N .
\]
積は対称かつ全射である。
さらに
\[
 D_Q:\operatorname{Sym}^3V
 \xrightarrow{d}V\otimes\operatorname{Sym}^2V
 \longrightarrow V\otimes(\operatorname{Sym}^2V/Q),\qquad
 \mathcal L=\ker D_Q
\]
は可逆部分束で、その形成は任意の基底変更と可換する。
$\mathcal L$ の局所生成元は滑らかな三次族
$Y\subset\mathbf P(V^\vee)$ を定め、その勾配空間は $Q$ である。
構成は rotation 付き Morita 同値と基底変更に関して自然である。
体上の滑らかな三次模型の存在自体は仮定であり、低積の階数だけから同値を結論しない。
\end{theorem}

\begin{proof}
一般点で $\Gamma_\eta$ は smooth である。
一般点での perfect 対角を、有限 cell の retract として有限個の係数で開へ広げる
\cite[Lemma~2.10]{NFMToenVaquieModuli}。
得られた perfect $\Gamma^e$-加群 $E$ から対角への一般点の同型は、
$E$ の compact 性によって局所化後の射 $E\to\Gamma$ へ広がる。
その cone は、$\Gamma$ の proper 性と $E$ の有限 cell 性から $S$-perfect である。
その有限個の cohomology は一般点で零なので、一つの非零元を反転して消せる。従って、
縮小後の対角は perfect となる。
rotation の逆と単位・余単位も同じ有限 cone の議論で広がる。

従って二つの導来 Hom 複体は $S$-perfect である。
開を縮めて全 cohomology を局所自由にすれば、
$H^0$ は非平坦なものも含めた任意の基底変更と可換する。
体上の低次数計算から、積の対称性、全射性、表示した階数を得る。
$N\ge5$ では非恒等 sector が $H^0$ に現れ得る次数は、
偶数 $N$ で $N/2$ であるため、次数 $1,2$ には寄与しない。
ここは固定入力B\cite{NFNCCubicB}の \texttt{gcd:lem:basechange},
\texttt{gcd:thm:field-model} と、その一次資料 Lin--Zhang の低次数計算を使う。

一般点の $Q$ は比較先の勾配空間である。
その積分空間は symmetrizer 代数と同型であり、ST 既約性から一次元である。
従って $D_Q$ の階数は $\binom{N+2}{3}-1$。
整性により全ての最大サイズの minor は零であり、階数 $\binom{N+2}{3}-1$ の minor を反転すると、核は可逆部分束、余核は局所自由となり、
核の形成は任意の基底変更と可換する。
$\mathcal L\otimes V^\vee\to Q$ の同型性と判別式の非消滅も
一般点から開へ広がる。全ての操作は指定同値に関して自然である。
\end{proof}

一ブロック表示への適用では、支配的な $h$ 族で入力Bを使い、
一般点の smoothness と幾何学的模型を得る。
smoothness は体拡大から降下する。
特殊三次 $F+G$ の scalar symmetrizer minor は完成上の単元なので、
形式三次 $W_h$ と幾何学的一般点も ST 既約である。
このことで元の表示空間の稠密開に定理を適用できる。
全 coherent 残余圏との一致には、ambient の noetherian/Koszul/AS 正則性、
中心元の正則性、線束生成部分の smoothness、および coherent saturation を用いる。


\subsection{中心作用と全 O-同値への降下}

\begin{lemma}[中心一次作用の一般点での単射性]
\label{lem:nfm-residual-central-injection}
良い一ブロック表示の中心一次部分 $U$ の乗法は、代数的な写像
\[
 c:U\longrightarrow V_\rho
   =H^0\operatorname{RHom}_{\mathcal K^e}(\Delta,P_\rho)
\]
を定める。固定入力の MF・window 比較の下で、この写像は
表示空間の稠密開上で階数 $r=n-1$ の部分束となる。
その形成はこの開からの任意基底変更と両立する。
\end{lemma}

\begin{proof}
中心元 $u$ の乗法は graded MF の各対象で厳密な自然変換
$\mathrm{id}\to(1)$ となり、微分と可換する。
window 比較により残余圏の rotation への自然変換を得る。
構成は次数付き同型、内部 shift と係数変更に両立する。

単に特殊ファイバーで単射であるだけでは、一般点での単射性は従わない。
そこで補題\ref{lem:nfm-dominant-central-family}の支配的な族を用い、
中心作用を実際の完備同値で可換側の $\operatorname{MF}^{\mathbb Z}(W_h)$ へ移す。
正則 DVR 上の可換 MF・window 比較により、後者は滑らかな三次族の相対 Ku である。
これは smooth proper な圏の有限 mutation による admissible 成分なので、
低 Nat の導来複体 $E$ は $k[[h]]$-perfect となる。
体上の低次数計算では、特殊体と Laurent 体で $H^0$ の次元はともに $N$ である。
DVR の完全列
\[
 0\longrightarrow H^0(E)\otimes k
 \longrightarrow H^0(E\otimes^{\mathbf L}k)
 \longrightarrow\operatorname{Tor}_1(H^1(E),k)
 \longrightarrow0
\]
において、中央の次元は $H^0(E)$ の一般階数に
$H^0(E),H^1(E)$ の torsion の生成元数を加えたものである。
両体での次元が等しいため、この二つの torsion は零である。
従って $H^0(E)$ は自由階数 $N$ で、特殊化と両立する。
完備同値の剰余は恒等なので、移した $c$ の剰余は通常の $U\subset k[x]_1$ である。
その階数 $r$ の minor の一つは単元となる。

固定三対象への制限を先に行い、その後で Laurent 局所化する。
ここでは入力Bの補題 \texttt{rog:lem:localization}を使う。
その像が代数閉拡大後に可換 MF 全体を生成するという入力から、
定義済みの関手は Laurent 体上ですでに同値である。
実際、smoothness を体拡大から降ろし、
随伴核の単位・余単位の cone が忠実平坦な体拡大後に零であることを使えばよい。
従って上の独立性は代数的な一般点の $c$ の独立性を意味する。
ここでは完備底上で固定三対象が生成するという強化は使わない。
最後に全階数の minor を反転すれば、代数的な係数開上の部分束を得る。
\end{proof}

\begin{lemma}[相対 flow による中心一次作用の保存]
\label{lem:nfm-relative-center-transport}
定理\ref{thm:nfm-residual-formal-poisson}の仮定の下で、
生成元を固定する星積模型
\[
 \phi:B_h\xrightarrow{\sim}(S_R[[h]],*_{\pi_h}),\qquad
 \pi_h=a_h\,\iota_{dW_h}P_0,
 \qquad P_0=\partial_x\wedge\partial_y\wedge\partial_z
\]
を取る。固定入力の完備 MF 移送が適用できるものとし、
その明示的な順序付き flow をこの模型に適用して得る同値を
\[
 \Phi_{\mathrm{rel}}:
 \operatorname{MF}^{\mathbf Z}(B_h,\Omega_h)
 \xrightarrow{\sim}
 \operatorname{MF}^{\mathbf Z}(S_R[[h]],W_h)
\]
と書く。この同値は、各内部次数で完備化した
$\operatorname{Sym}_R U_R[[h]]$ に関して相対的に取れる。
特に、任意の $u\in U_R$ に対する中心乗法の自然変換
$m_u:\operatorname{Id}\to(1)$ を、標的の同じ中心乗法へ送る。

従って、内部 shift に対応する rotation の比較と可換側の低次数同定の下で、
中心作用の像は通常の線型式 $U_R[[h]]\subset S_{R,1}[[h]]$ に対応する。
この主張は、上で選んだ相対 flow に関するものであり、
任意の完備 MF 同値が中心作用を同じ座標で保つと主張するものではない。
\end{lemma}

\begin{proof}
補題\ref{lem:nfm-formal-linear-center}により、
$f\in\operatorname{Sym}_R U_R[[h]]$ の星積による乗法は通常の乗法である。
$\phi$ は中心一次生成元を固定するので、
\[
 \phi(fb)=f*_{\pi_h}\phi(b)=f\phi(b)
\]
となる。従って星積模型への最初の厳密同型は、この中心係数環に関して相対的である。

固定入力の補題\texttt{cft:lem:primitive}における定数三重ベクトルは、
一意性と $\pi_h=a_h\iota_{dW_h}P_0$ により $\Xi=a_hP_0$ である。
同入力の曲率付き経路とその generator は
\[
 \alpha_s=-tW_h+(1-s)\pi_h,\qquad
 \eta=\Xi/t,\qquad
 \gamma_s=\mathcal U^{\alpha_s}(\eta)
\]
であり、$t$ は微分しない係数として扱われる。
各グラフの関数頂点 $W_h$ は出る辺を持たず、
$\pi_h$ と $\Xi$ の出る辺の添字は全て active 方向 $x,y,z$ に属する。
よって境界入力に作用する微分も全て active 方向であり、
$\gamma_s$ の各 arity 成分は
$\operatorname{Sym}_R U_R[[h]][t,t^{-1}]$ に関して多重線型である。
係数や $W_h$ が中心変数に依存していても、この議論は変わらない。

補題\texttt{cft:lem:arity}の graph による単位正規化から、
$\gamma_s$ の正の arity 成分は単位入力で零になる。
中心係数に関する多重線型性と合わせると、
単位の代わりに任意の $f(U)$ を入れても零になる。
ここで別の contracting homotopy や非相対な正規化を加える必要はない。
補題\texttt{cft:lem:flow}の順序付き flow とその逆は、
これらの coderivation の合成と積分から作られるので、
同じ相対的多重線型性と単位正規化を保つ。
従って、その Taylor 成分 $F_j$ は
$F_1(1)=1$、$j\ge2$ に対して単位を含む $F_j$ は零、
かつ中心係数に関して多重線型である。

有限 rank の MF 行列 $Q$ に対し、中心乗法の像は挿入公式により
\[
 \bigl(\Phi_{\mathrm{rel}}\bigr)_1(u\operatorname{id})
 =\sum_{p,q\ge0}
 F_{p+q+1}\bigl(Q(1)^{\otimes p},
                 u\operatorname{id},Q^{\otimes q}\bigr)
\]
と書ける。各項から $u$ を取り出すと、$p+q>0$ の項は単位入力の正規化で零であり、
残る項は $F_1(u\operatorname{id})=u\operatorname{id}$ である。
内部 shift は自由基底のラベルを変えるだけなので、この行列計算と両立する。
高次の自然変換成分も、$u$ を取り出した後には単位入力を含む
arity 二以上の項となるため零である。
従って保存されるのは対象ごとの射だけでなく、中心作用の自然変換そのものである。
全ての式は有限 Artin 商で有限和として成立し、
固定内部次数内の完備性により完備底上でも成立する。
最後に window 比較を用いれば、rotation に対する中心作用の主張を得る。
\end{proof}


\begin{theorem}[代数的な残余射影 Poisson 射]
\label{thm:nfm-projective-residual-morphism}
$n\ge3$ とし、固定入力Bの完備 MF 移送、正則底の window 比較、
一般点生成性、coherent 飽和性を上記の支配的な係数族に適用する。
既存の良い一ブロック O-marked stack
$\mathfrak M_n^{\mathrm{sm}}$ には非空稠密開
$\mathfrak M_n^{\mathrm{res}}$ があり、標準的な代数的射
\begin{equation}
\label{eq:nfm-projective-residual-map}
 \Psi_{\mathrm{res}}:
 \mathfrak M_n^{\mathrm{res}}\longrightarrow
 \overline{\mathfrak P}_n^{\mathrm{dec}}
 :=
 \bigl[(\mathcal H_N^{\mathrm{sm}}\times
              \Gr(3,V_0^*))/\PGL(V_0)\bigr]
\end{equation}
が定まる。ここで $V_0$ は固定した $N$ 次元空間、
$\mathcal H_N^{\mathrm{sm}}\subset\PP(\Sym^3V_0)$ は滑らかな三次の開である。
標的は Brauer--Severi ambient を許し、
非零分解可能 Poisson tensor の尺度を忘れた対の stack である。

局所表示でその値は、低 rotation 積から得る
$\mathcal L=\ker D_{Q_\rho}$ と
\[
 K_{\mathrm{res}}=\operatorname{ann}(c(U))\subset V_\rho^*
\]
で与えられる。特に underlying cubic は
$X_{\mathrm{res}}=V(\mathcal L)\subset\PP(V_\rho^*)$ である。
構成は全 O-同値に降下し、任意基底変更と両立する。
大域的な次数一 shift、方程式の生成元、split ambient は仮定しない。
\end{theorem}

\begin{proof}
定理\ref{thm:nfm-relative-residual-cubic}と前補題で、
$\mathcal L$ と階数三の $K_{\mathrm{res}}$ を代数的に作る。
必要な perfect 性、低次 cohomology の局所自由性、積と微分の階数、
全三次と active 制限の滑らかさは開条件である。
これらは Morita 同値と基底線束による twist に不変であるため、
その成立 locus は表示変更群に不変な開となる。
支配的な $h$ 族と補題\ref{lem:nfm-uniform-scalar-symmetrizer}により非空である。
既存の非空開 $\mathfrak M_n^{\mathrm{sm}}$ との交差も非空である。

線束 $\mathcal L$ と $\det K_{\mathrm{res}}$ の局所生成元をそれぞれ $G,P$ と取ると
$\iota_{dG}P$ は三変数の Jacobian bracket であり、Jacobi 恒等式を満たす。
また $G$ は Casimir である。
定理\ref{thm:poisson-strong-lift}の単射性から、この tensor は非零である。
$G,P$ の変更は tensor の非零スカラー倍だけを変える。
従って $(\mathcal L,K_{\mathrm{res}})$ は所要の射影 Poisson 対を与える。

降下を確認する。全 O-同値は局所的に次数付き同型に持ち上がる。
持上げの線束を $M$ とし、その向きを $V_\rho'=V_\rho\otimes M$ と取ると
\begin{equation}
\label{eq:nfm-residual-line-weights}
 \begin{aligned}
 V_\rho'&=V_\rho\otimes M,&
 Q_\rho'&=Q_\rho\otimes M^{\otimes2},&
 \mathcal L'&=\mathcal L\otimes M^{\otimes3},\\
 c(U)'&=c(U)\otimes M,&
 K_{\mathrm{res}}'&=K_{\mathrm{res}}\otimes M^{-1}.
 \end{aligned}
\end{equation}
局所的な $\mathcal O(i)$ を基底線束で捩じっても、
それらが生成する exceptional block は変わらない。
従って残余圏は全 O-同値に保たれ、rotation だけが $M$ で捩じれる。
式\eqref{eq:nfm-residual-line-weights}はこの変換と積の自然性から従う。

二つの次数付き持上げの差は次数スカラーである。
その $V_\rho$ 上の作用はスカラーなので、ambient、三次、
$K_{\mathrm{res}}$ の射影同型は持上げに依存しない。
三重重なり上の差も次数スカラーであり、射影同型は cocycle 条件を満たす。
これで幾何点の同型類だけでなく、族と全射の fpqc 降下を得る。
核と部分束の形成が任意基底変更と両立することから、射にも同じ性質がある。
重み零の降下と Brauer--Severi ambient の一般的背景は
\cite[Proposition~4.10; Remark~6.1]{NFMBerghSchnurer}に対応する。
\end{proof}

\begin{proposition}[尺度の線束]
\label{prop:nfm-residual-scale-torsor}
表示 gerbe 上で
\[
 \mathscr T=\mathcal L\otimes\det K_{\mathrm{res}}
\]
は重み零であり、$\mathfrak M_n^{\mathrm{res}}$ 上の通常の線束へ降下する。
尺度を持つ Poisson stack
$\mathfrak P_n^{\mathrm{dec}}\to\overline{\mathfrak P}_n^{\mathrm{dec}}$
の $\Psi_{\mathrm{res}}$ に沿う引戻しは $\mathscr T^\times$ である。
従ってその lift は $\mathscr T$ の至る所非零な切断の選択に等しい。
\end{proposition}

\begin{proof}
式\eqref{eq:nfm-residual-line-weights}で重みは $3+(-3)=0$ となる。
$G\otimes P\mapsto\iota_{dG}P$ は単射であり、
その非零値が固定した Poisson 直線の全尺度をちょうど一度ずつ与える。
従ってこの線束の零切断を除いたものが引戻し torsor である。
\end{proof}

ここで $\mathscr T$ が非自明とは主張しない。
必要な区別は、Poisson 直線、尺度の非標準的な代数的選択、
および固定 formality によって指定された形式的尺度の三者である。

\begin{proposition}[支配性と一般相対次元]
\label{prop:nfm-projective-residual-dominance}
定理\ref{thm:nfm-projective-residual-morphism}の射は支配的である。
source をさらに非空開へ制限すれば、相対次元一の滑らかな射となる。
\end{proposition}

\begin{proof}
補題\ref{lem:nfm-dominant-central-family}で、$F+G$ を一般の全変数三次とする。
active support を動かす frame も許せば、特殊点 $(F+G,\operatorname{ann}U)$ は
標的の一般点となる。
実際の完備 MF 比較による低積の同定と積分直線の一意性から、
一般 $h$ での $\Psi_{\mathrm{res}}$ の三次は $V(W_h)$ と射影同型である。
移送した中心作用は前補題の自由な低 Nat 加群の部分束を与え、
その剰余は $U$ である。
従ってこの標的値は $k[[h]]$ 上の族へ延長し、特殊点は所要の一般対である。
像の閉包が真の閉部分なら、その一般形式点だけでなく特殊点も含まねばならず、
標的の一般点を含むことに矛盾する。

標的は滑らかな有限型 DM stack であり、
\[
 \dim\overline{\mathfrak P}_n^{\mathrm{dec}}
 =\binom{N+2}{3}-N^2+3(N-3)
 =\binom{n+2}{3}+3n-3
 =\dim\mathfrak M_n^{\mathrm{sm}}-1.
\]
標数零の generic smoothness を étale chart 上で適用すれば結論を得る。
表現可能性、幾何学的連結 fibre、有理商の記述までは主張していない。
\end{proof}


\begin{corollary}[形式 Poisson 対との全次数比較]
\label{cor:nfm-projective-formal-compatibility}
上の支配的な中心 $h$ 族に沿い、生成元固定模型と
補題\ref{lem:nfm-relative-center-transport}の相対 flow を選ぶ。
$k((h))$ 上の $\Psi_{\mathrm{res}}$ の値は、形式線型同定の下で
\[
 \bigl(V(W_h),[\pi_h]\bigr)
\]
である。この値は $k[[h]]$ 上の
$\bigl(V(W_h),[\pi_h/h]\bigr)$ に延長し、特殊値は
$\bigl(V(F+G),[\Pi_{F+G}]\bigr)$ となる。
ここで $F+G$ は選んだ中心元 $\Omega_h+G$ の剰余であり、
$G$ は純中心三次なので $\Pi_{F+G}=\Pi_F$ である。
\end{corollary}

\begin{proof}
低 rotation 積の同定から、$Q_\rho$ は $\nabla W_h$ に移る。
scalar symmetrizer 条件の下では積分直線が一意なので、
$\mathcal L$ は $k((h))W_h$ に移る。
相対 flow は中心作用を厳密に保つため、$K_{\mathrm{res}}$ は
$\operatorname{ann}U$ に移る。
定理\ref{thm:nfm-residual-formal-poisson}の
$\pi_h=a_h\iota_{dW_h}P_0$ から、両 Poisson 直線は一致する。
$a_h/h$ は定数項一の形式単元であるから、$[\pi_h/h]$ は特殊点でも定義される。
ここで用いる次数一空間の同定は形式行列であり、その行列の代数性は主張しない。
\end{proof}

\subsection{同じ三次への比較核と残る降下障害}

\begin{theorem}[有限 étale 被覆上の残余比較]
\label{thm:nfm-residual-comparison-cover}
定理\ref{thm:nfm-projective-residual-morphism}の固定入力の下で、
$\mathfrak M_n^{\mathrm{res}}$ の非空開
$\mathfrak M_n^{\mathrm{cmp}}$ と有限 étale 全射
$S'\to\mathfrak M_n^{\mathrm{cmp}}$ を取れ、
その上で次数付き表示を選ぶと
\[
 \mathcal K_{S'}\simeq
 \operatorname{Ku}(X_{\mathrm{res},S'})
\]
という dg Morita 同値が perfect kernel で実現する。
局所 rotation と三次方程式の生成元を選んだ後、
rotation と指定三次周期も保存できる。
比較先は $\Psi_{\mathrm{res}}$ が与えた同じ $X_{\mathrm{res}}$ である。
\end{theorem}

\begin{proof}
まず整な有限型表示 chart $S$ とその関数体 $k(S)$ 上で考える。
低積から得た三次は、既存の幾何的模型と代数閉包上で射影同型である。
固定入力Bの \texttt{gcd:thm:field-model}と
\texttt{rog:lem:descent}により、有限分離拡大 $k'/k(S)$ 上で
この固定した三次への比較核を得る。
方程式の単元倍と指定周期を合わせるために必要な三乗根も有限拡大に含める。

$S$ を正規にし、$k'$ での正規化を取る。
分岐 locus を除けば有限 étale 全射となる。
perfect kernel、逆 kernel、単位・余単位、rotation 比較、
周期に関する有限個の図式を局所化後へ広げる
\cite[Lemma~2.10]{NFMToenVaquieModuli}。
同値性を判定する単位・余単位の cone と rotation 比較の cone は底上 perfect であり、
一般点で零なので、その cohomology の台を除けば零になる
\cite[Lemma~3.25 and its proof]{NFMToenVaquieModuli}。
有限個の図式の差も局所化後に零にできる。
被覆上で除いた閉集合の像は有限性により閉なので、
底も縮めれば被覆の全射性を保てる。

stack については、一般点に達する affine étale chart の支配成分を取る。
source は有限型 DM stack で affine diagonal を持つので、
この chart は表現可能・準有限・分離的である。
Zariski の主定理で有限化し、その中で支配成分の scheme 理論的閉包を取る。
この閉包の境界の有限閉像は一般点を含まない。
その像を除けば
非空開上の有限 étale 全射を得る。
その関数体上で表示 gerbe の中立化と必要な核を有限体拡大で選び、
同じ有限表示の議論を適用する。
これである $S'$ と所要の比較を得る。
\end{proof}

この定理は、最初に選んだ形式同値の全 Taylor 係数を代数化する主張ではない。
また裸の O-marked 対は、大域的な rotation や周期の自明化を標識していない。
Brauer--Severi ambient では個々の $\mathcal O_X(i)$ を大域線束と仮定せず、
局所的にそれらが生成する admissible block を用いる。
基底線束による twist ではその評価 projector は変わらないので、
block とその直交は降下する。
一般に block は Brauer 類の冪で捩じれた基底圏となる
\cite[Theorem~6.2; Remark~6.3]{NFMBerghSchnurer}。
上の比較被覆では通常の split な Ku を使える。

\begin{proposition}[同値核の降下の二段階]
\label{prop:nfm-residual-kernel-descent}
同じ残余三次と局所 rotation を固定しても、
比較核全体は rotation に両立する全自己同値 stack の torsor であり、
一般に一つの $\mathbf G_m$-gerbe ではない。
局所核の同型類が重なり上で整合するように選べた後で、
その核の実際の降下はスカラー二 cocycle の消滅に帰着する。
\end{proposition}

\begin{proof}
同値 $P,Q:\mathcal K\to\operatorname{Ku}(X_{\mathrm{res}})$ の差は
$QP^{-1}$ という自己同値であり、自己同値の合成は全ての比較核に単純推移的に作用する。
一つの核 $P$ が存在すれば $P[2m]$、$m\in\mathbb Z$ は
同じ rotation・周期に両立し、次数零の低 Nat 積にも同じ作用を与える。
これらは互いに非同型である。
実際、非零 proper 圏で非零の偶 shift が恒等なら、
非零対象の恒等射から無限個の cohomological 次数の自己 Hom が生じ、
proper 性に矛盾する。
従って低積と方程式直線の一意性は、核の同型類の一意性を意味しない。

一方、幾何的な Ku の $\mathrm{HH}^{<0}=0$、$\mathrm{HH}^0=k$ と
Morita 不変性から、固定した同値核は simple である。
相対的にも必要な開でこの性質と基底変更を保てる。
整合した同型類を選んだ後に重なりの同型 $\alpha_{ij}$ を選ぶと、
三重重なりで
\[
 \alpha_{ij}\alpha_{jk}\alpha_{ki}
    =c_{ijk}\,\mathrm{id},\qquad c_{ijk}\in\mathcal O^\times
\]
となる。$c_{ijk}$ が coboundary であることが、これらを cocycle 条件へ修正して
実際の kernel を降ろすための必要十分条件である。
これは simple-object stack の $\mathbf G_m$-gerbe
\cite[Corollary~3.22]{NFMToenVaquieModuli}であり、
全自己同値 sheaf \cite[Corollary~3.24]{NFMToenVaquieModuli}とは区別される。
\end{proof}

$\mathrm{HH}^1=0$ による infinitesimal rigidity は、
核の同型類の選択やこの二 cocycle の消滅を含まない。
表示 gerbe の Brauer 類 $\beta$ と比較核の gerbe は異なる降下問題から生じ、
両者を同一視する構成はまだない。

\subsection{一変数 trait に沿う良い還元}

以下は稠密開の定理とは別である。
相対 Hochschild 複体を $h=0$ まで perfect と証明しなくても、
形式比較を選んだ後の三次の良い還元の存在は得られる。

\begin{lemma}[完成上の格子の代数的降下]
\label{lem:nfm-completed-lattice-descent}
$R=k[h]_{(h)}$、$K=k(h)$、$\widehat R=k[[h]]$、
$\widehat K=k((h))$ とする。
有限次元 $K$-空間 $V_K$ の
$V_K\otimes_K\widehat K$ 内の任意の $\widehat R$-格子 $\widehat\Lambda$ は
\[
 \Lambda=V_K\cap\widehat\Lambda
\]
という有限自由 $R$-格子の完成である。
\end{lemma}
\begin{proof}
$V_K\simeq K^N$ を選び、
$h^a\widehat R^N\subset\widehat\Lambda\subset h^{-b}\widehat R^N$
とする。両端の商は有限 jet 加群であり、
$R/(h^{a+b})=\widehat R/(h^{a+b})$ により両基底で同じである。
その中の $\widehat\Lambda/h^a\widehat R^N$ の逆像を
$h^{-b}R^N$ 内で取れば $\Lambda$ となる。
これは有限生成 torsion-free、従って自由であり、完成は所要の格子である。
\end{proof}

\begin{corollary}[形式比較を選んだ滑らかな trait 模型]
\label{cor:nfm-residual-good-reduction}
一変数量子化に入力Bの (L),(CF) を仮定し、
$W_0$ が滑らかで scalar symmetrizer を持つとする。
一般点の正準的な積分直線から得る三次 $Y_K$ は、
滑らかな射影三次族 $Y_R/R$ に延長できる。
その完成は $V(W_h)$ と射影同型である。
\end{corollary}
\begin{proof}
固定三対象から得た完全忠実な関手
\[
 \operatorname{Perf}(\Gamma_{\widehat K})
 \longrightarrow\operatorname{Ku}(W_{\widehat K})
\]
は、$\widehat K$ の代数閉包上で生成的である。
この関手は既に $\widehat K$ 上で定義されている。
smoothness を体拡大から降下し、随伴核の単位・余単位の cone を調べれば、
代数閉包上の同値は元の $\widehat K$ 上のこの関手の同値を意味する。
従って有限体拡大への閉点特殊化をする前に
\[
 V_K\otimes\widehat K\simeq
 \widehat K[x_1,\ldots,x_N]_1,\qquad
 Q_K\otimes\widehat K\simeq\nabla W_h
\]
という実際の同定を持つ。
scalar symmetrizer 条件により正準的な積分直線も
$\widehat K W_h$ と同定される。

右辺の標準 $\widehat R$-格子を引き戻し、
補題 \ref{lem:nfm-completed-lattice-descent} により $\Lambda\subset V_K$ へ降ろす。
$Q_R=Q_K\cap\operatorname{Sym}^2\Lambda$ は saturated であり、
DVR 上で商も有限自由である。
その完成は $\nabla W_h$ に一致する。
実際、滑らかな特殊三次の勾配は独立なので
$\nabla W_h$ は $\operatorname{Sym}^2\widehat\Lambda$ の direct summand であり、
同じ一般部分空間を持つ saturated 格子は一意である。
$D_{Q_R}$ の特殊階数は $\binom{N+2}{3}-1$ であり、
この大きさの単元 minor がある。行消去後の残る成分は、
完成上では $\widehat R W_h$ という剰余非零の解があるため全て零となる。
忠実平坦性により $R$ 上でも零なので、
$\mathcal L_R=\ker D_{Q_R}\subset\operatorname{Sym}^3\Lambda$ は直線部分束である。
その完成は $\widehat R W_h$ であり、判別式は特殊ファイバーで非零である。
\end{proof}

この系の格子は選んだ形式比較に依存し得る。
自然な次数一の束への正準的な延長、全 O-同値に沿う正準的な trait 模型、
$h=0$ を含む大域モジュライ射、比較 kernel の延長・cocycle 降下は別問題である。
また形式線型座標変更を許した射影族の代数化であり、
固定座標での $W_h$ の各係数の代数性を主張しない。


\subsection{固定生成元の形式 tensor の very general な非代数性}

\begin{proposition}[多項式係数形式級数の very general な超越性]
\label{prop:nfm-very-general-transcendence}
$A=\mathbb A^s_{\mathbb C}$ を係数空間とし、
\[
 f(\mathbf c,h)=\sum_{m\ge0}f_m(\mathbf c)h^m
       \in\mathbb C[\mathbf c][[h]]
\]
とする。ある $\mathbf c_0\in A(\mathbb C)$ について
$f(\mathbf c_0,h)$ が $\mathbb C(h)$ 上超越的なら、次が成立する。
\begin{enumerate}
\item $f$ は $\mathbb C(\mathbf c,h)$ 上超越的である。
\item $f(\mathbf c,h)$ が $\mathbb C(h)$ 上代数的となる点全体は、
可算個の真の Zariski 閉集合の和である。
従って任意の非空 Zariski 開 $U\subset A$ の very general な点で超越的である。
\item 上記の超越性は、それぞれの有理関数体の有限代数拡大によっては解消されない。
\end{enumerate}
\end{proposition}

\begin{proof}
$d\ge0,e\ge1$ を固定し、$M=(d+1)(e+1)$ 個の形式級数
\[
 h^if(\mathbf c,h)^j\qquad(0\le i\le d,\ 0\le j\le e)
\]
の係数を列とする、無限行・$M$ 列の行列 $B_{d,e}$ を取る。
各成分は $\mathbb C[\mathbf c]$ に属する。
全ての $M$ 次小行列式の零点集合を $Z_{d,e}$ とすると、これは Zariski 閉である。
Noether 性から、その定義には有限個の多項式で十分である。

$\mathbf c_0$ ではこれらの列は $\mathbb C$ 上一次独立である。
有限個の列が独立なら、有限個の行を選んですでに階数 $M$ となるので、
$\mathbf c_0$ で非零の $M$ 次小行列式が存在する。
従って $Z_{d,e}$ は真の閉集合であり、同じ小行列式は一般点でも非零である。
一般点での列の独立性を全ての $d,e$ に適用すれば第一の主張を得る。

一方、$f(\mathbf c,h)$ の $\mathbb C(h)$ 上の代数性は、
分母を払うことにより、ある $d,e$ について上記の列が一次従属となることと同値である。
従って代数的となる点全体は
$\bigcup_{d\ge0,e\ge1}Z_{d,e}$ に等しい。
既に各 $Z_{d,e}$ が真と分かっているので、第二の主張が従う。
最後の主張は代数拡大の推移性による。
\end{proof}

\begin{lemma}[logarithmic な境界値]
\label{lem:nfm-logarithmic-specialization-transcendental}
\[
 \ell(h)=\log\frac{2+h}{2-h}=2\operatorname{arctanh}(h/2)
       \in\mathbb C[[h]]
\]
は $\mathbb C(h)$ 上超越的である。
\end{lemma}

\begin{proof}
代数的なら、有限関数体拡大 $\mathbb C(h,\ell)$ の滑らかな射影曲線上で
\[
 d\ell=\frac{4}{4-h^2}\,dh
\]
が成立する。$h=2$ 上の点の分岐指数を $r\ge1$ とすると、
右辺の留数は $-r$ である。
しかし有理関数の完全微分の留数は全て零なので矛盾する。
\end{proof}

\begin{remark}[固定された規格化への適用範囲]
固定生成元と規格化を選んだ形式 Poisson tensor の一つの座標係数が
$\mathbb C[\mathbf c][[h]]$ に属し、ある境界点で上の $\ell(h)$ に等しいなら、
命題~\ref{prop:nfm-very-general-transcendence} をその係数へ適用できる。
境界点自身の三次が滑らかである必要はない。
ただし、係数の多項式性を境界まで保証すること、および境界値の等式を
全ての次数で証明することが適用の前提であり、有限 jet の一致では不十分である。

結論は一般点および very general な点での超越性であり、
全ての点で超越性を持つ単一の Zariski 開の存在を主張しない。
また、同じ係数空間内の非空開に関する主張であり、
全ての intrinsic な既約成分を被覆する主張ではない。
有限代数的な底変更や代数的に可逆な座標変更は、
固定されたテンソルの超越性を解消しない。
任意の形式座標変更、形式的な尺度変更、gauge 変更、
あるいは三次超曲面の射影同型類の代数化には、この障害をそのまま適用できない。
\end{remark}


\begin{theorem}[一般次元の正規化 tensor の超越性]
\label{thm:nfm-normalized-tensor-transcendence}
$N=n+2\ge5$、$V=\langle x,y,z,u_1,\ldots,u_{n-1}\rangle$、
$P_0=\partial_x\wedge\partial_y\wedge\partial_z$ とし、
$\mathcal A=\Sym^3V$ を全三次係数の affine 空間とする。
$F\in\mathcal A$ に対して
\[
 [z_i,z_j]=h\operatorname{sym}\bigl(\Pi_F(dz_i,dz_j)\bigr),
 \qquad \Pi_F=\iota_{dF}P_0
\]
の固定 harmonic-angle・生成元固定の逆 tensor を $\pi_h(F)$ と書く。
滑らかな全三次、滑らかな active 制限、および上の模型の存在条件を満たす
任意の非空係数開を $\mathcal A^{\mathrm{good}}$ とする。
このとき、$\pi_h(F)$ の係数の全てが $\mathbb C(h)$ 上代数的となる
$F\in\mathcal A^{\mathrm{good}}(\mathbb C)$ は、
可算個の真の閉集合の和に含まれる。
従って very general な $F$ の tensor は $\mathbb C(h)$ 上代数的でない。
普遍係数の関数体上でも非代数的であり、有限代数拡大によって解消しない。
\end{theorem}

\begin{proof}
式\eqref{eq:nfm-normalized-tensor-recursion}を全 affine 係数空間上で使う。
この再帰の各段階は、既に得た係数の和・積・座標微分と、
固定した複素グラフ重みだけからなる。
$h$ の各次数でグラフは有限個であり、$q_{ij}=\Pi_F(dz_i,dz_j)$ は
$F$ の係数に線型なので、帰納的に
\[
 \pi_m(F)\in\mathbb C[\mathcal A]\otimes
      \bigl(\Sym^2V\otimes\textstyle\bigwedge^2V^*\bigr)
\]
である。係数空間の判別式などで割る操作はない。
従って、境界で PBW 性や MC 恒等式を主張しなくても、
再帰で得る係数列自体は全 $\mathcal A$ 上に定義される。
模型が存在する良い開では、補題\ref{lem:nfm-normalized-star-model}の一意性により
これは実際の逆 tensor である。

$\pi_h(F)(dx,dy)$ における単項式 $xy$ の係数を $f(F,h)$ と置く。
これは $\mathbb C[\mathcal A][[h]]$ の元である。
$F_0=xyz$ では、補題\ref{lem:nfm-log-canonical-inverse}により
\[
 f(F_0,h)=\log\frac{2+h}{2-h}.
\]
この恒等式は Lindberg--Pym\cite[Section~1.6; Theorem~6.9, Remark~6.10]{NFMLindbergPym}
の原生成元での交換比から得る全次数の恒等式であり、有限 jet の一致ではない。
前補題と命題\ref{prop:nfm-very-general-transcendence}を適用すれば結論を得る。
$F_0$ は良い開に属さないが、各有界次数の代数関係条件が
全係数空間の真の閉集合であることを検出するためだけに使っている。
その閉集合の任意の非空開への制限も真なので、滑らかな一般部分への主張が従う。
\end{proof}

\begin{remark}[代数的な Poisson 射との両立]
この障害は指定した $h$ と次数一生成元を保つ実際の $\pi_h$ の係数に対するものである。
代数的係数の可逆線型変更の後に全 tensor 係数が代数的なら、
逆変換によって元の係数も代数的となるため、そのような変更でも障害は消えない。
一方、形式行列による座標変更、尺度を忘れた Poisson 直線、
あるいはその underlying cubic の射影同型類には、同じ結論は及ばない。
本節の代数的な残余射は後者のデータを復元する。
また、この定理は $W_h$ の特定の係数の超越性を主張せず、
実際の tensor の尺度を選ぶ代数的な局所切断一般を排除するものでもない。
その切断を元の固定生成元、または代数的に可逆な変更で得た frame で
代数的に表示でき、そこで $\pi_h$ と一致する場合は排除される。
形式的な Nat の同定を経た一致だけは、この定理からは排除しない。
\end{remark}

% !TEX root = ../nc-cubic-n-fold.tex
\section{点曲線による残余 Poisson 構造の規格化}
\label{sec:nfm-residual-point-scale}

残余方程式の直線と中心作用から得た射に、量子側の楕円点曲線が持つ微分を組み合わせる。
全 $n\ge3$ で active 曲線の同型類を比較し、符号を除いた標準尺度と、
標準二重被覆上の尺度付き残余射を構成する。
比較先の三次は前節の $X_{\mathrm{res}}$ のままである。
本節の残余圏への適用は、前節で明記した固定 MF・window・生成性・飽和性の入力の下にある。

\subsection{残余三次の active 切断と楕円点曲線}

以下の主張は任意の $n\geq3$ に一様であり、
定理\ref{thm:nfm-projective-residual-morphism}および
定理\ref{thm:nfm-split-corrected-residual}の固定入力を用いる。
中心変数を零にする操作について、新たな MF 圏の基底変更定理は仮定しない。
必要な比較を正規化 tensor の普遍再帰の段階で行う。

\begin{lemma}[active 制限と Hessian 三次の全次数 Casimir 性]
\label{lem:nfm-active-restriction-all-orders}
$R$ を可換複素代数、$U=\langle u_1,\ldots,u_r\rangle$ とし、
$F\in R[x,y,z,U]_3$、$g=F|_{U=0}$ と置く。
$P_0=\partial_x\wedge\partial_y\wedge\partial_z$ に関する関係式
$[a,b]=h\operatorname{sym}(\iota_{dF}P_0)(a,b)$ から、
式\eqref{eq:nfm-normalized-tensor-recursion}で得る形式 tensor を $\pi_h(F)$ とする。
三変数だけで同じ再帰を行ったものを $\pi_h^{\mathrm{act}}(g)$ と書くと
\[
 \pi_h(F)|_{U=0}=\pi_h^{\mathrm{act}}(g).
\]
さらに普遍的な形式多項式恒等式
\begin{equation}\label{eq:nfm-active-hessian-casimir}
 \iota_{dJ_g}\pi_h^{\mathrm{act}}(g)=0,
 \qquad J_g=g+\frac{h^2}{24}\det\operatorname{Hess}(g)
\end{equation}
が成立する。従って、これらの恒等式は任意の係数変更と両立する。
\end{lemma}

\begin{proof}
補題\ref{lem:nfm-formal-linear-center}の帰納法により、
全ての tensor 頂点から出る辺は active 方向を微分する。
実際、元の二次関係では $q_{u_i a}=0$ であり、低次の tensor が active なら、
$u_i$ を入力とする正次数グラフは全て零になる。
再帰式から次の係数についても $\pi_m(du_i,-)=0$ となるので、
この消滅と以下の制限との可換性を任意係数環上で同時に帰納できる。
固定次数の再帰は加減乗算、active 微分、および Weyl 対称化の有限和であり、
これらはいずれも $U=0$ への制限と可換する。
$h$ の次数に関する帰納法で第一の等式を得る。
この議論では純中心変数の個数に依存する項は生じない。

第二の等式も各 $h$ 係数について $g$ の係数の多項式恒等式である。
まず $g_\lambda=S/3+\lambda xyz$、$S=x^3+y^3+z^3$ を考える。
定理\ref{thm:nfm-split-corrected-residual}の証明では、
生成元固定模型の active potential が
$a(s)S/3+b(s)xyz$ であり、同じ active 座標において
\[
 \frac{b(s)}{a(s)}
 =\frac{\lambda+s^2(4+\lambda^3)/12}
        {1-s^2\lambda^2/4}
\]
であることを示した。
従ってこの potential は
$g_\lambda+s^2\det\operatorname{Hess}(g_\lambda)/24$ の形式単元倍である。
ここで使うのは単なる $j$ の一致ではない。
同証明の共通の非零な二次先頭係数が、$j$ の分岐点を含めて
Hesse 係数の分枝を一意に定めることまで用いている。
Casimir 条件から\eqref{eq:nfm-active-hessian-casimir}を得る。

$g=p g_\lambda$ については $s=ph$ とすれば
$J_g=p\bigl(g_\lambda+s^2\det\operatorname{Hess}(g_\lambda)/24\bigr)$ である。
生成元固定再帰の線型共変性により、この恒等式は
$\operatorname{SL}_3$ による座標変更後にも成立する。
定理\ref{thm:nfm-correction-degree-nine}の証明で用いた支配射
\[
 \operatorname{SL}_3\times\mathbb A^2\longrightarrow
 \operatorname{Sym}^3\mathbf C^3{}^*,\qquad
 (A,p,q)\longmapsto A\cdot(pS/3+qxyz)
\]
から、各 $h$ 係数の多項式は三次係数空間の稠密開で零となる。
従って普遍係数環上で零である。
この多項式恒等式を任意の $R$ へ引き戻せばよい。
三変数再帰への制限と一般次元の MF/window 比較の適用は別の操作であり、
$n\geq3$ の圏比較を $n=1$ に代入してはいない。
\end{proof}

\begin{proposition}[形式残余三次の active 部分の厳密な決定]
\label{prop:nfm-residual-active-point-cubic}
定理\ref{thm:nfm-residual-formal-poisson}の模型で
$g=F|_{U=0}$ は滑らかな平面三次を定めるとする。
純中心三次 $G$ を中心元に加えてよい。
このときある形式単元 $b_h\in1+hR[[h]]$ によって
\begin{equation}\label{eq:nfm-residual-active-point-cubic}
 W_h|_{U=0}=b_h\left(g+\frac{h^2}{24}\det\operatorname{Hess}(g)\right)
 =b_hJ_g
\end{equation}
となる。従って、生成元固定形式座標における active 平面切断は、
量子側の楕円点曲線と全次数で同じ平面三次である。
\end{proposition}

\begin{proof}
$J_g$ の active 勾配は $g$ の正則勾配の形式変形なので、
補題\ref{lem:nfm-relative-casimir-tensor}を三変数で適用できる。
前補題の Casimir 恒等式から
\[
 \pi_h^{\mathrm{act}}(g)=c_h\,\iota_{dJ_g}P_0,
 \qquad c_h=h+O(h^2)
\]
を得る。$c_h$ は active 座標に依存しない。
一方、全変数の形式模型では
$\pi_h(F)=a_h\iota_{dW_h}P_0$、$a_h=h+O(h^2)$ であり、
$a_h$ は中心変数を含む全座標に依存しない。
$U=0$ に制限して $h$ を除くと
\[
 d(W_h|_{U=0})
 =\frac{c_h/h}{a_h/h}\,dJ_g.
\]
両辺は斉次三次式の微分であるから Euler 恒等式によって
\eqref{eq:nfm-residual-active-point-cubic}を得る。
純中心三次は制限で消え、形式 tensor 自体も選んだ中心三次に依存しない。
点曲線の方程式は命題\ref{prop:nfm-correction-point-cubic}である。
なお $h=0$ では、ここで延長するのは一般ファイバーの楕円成分の滑らかな閉包である。
可換多項式環の点スキーム全体を三次曲線と同定する主張ではない。
\end{proof}

\begin{theorem}[点曲線の $j$ を保つ残余射]
\label{thm:nfm-residual-point-j}
$\mathfrak M_n^{\mathrm{res}}$ を適切な非空稠密開に縮め、
量子側の楕円点曲線 $E_{\mathrm{pt}}$ と
\[
 C_{\mathrm{res}}=X_{\mathrm{res}}\cap\mathbf P(K_{\mathrm{res}})
\]
がともに滑らかであるとする。
全 O-同値に自然な正則関数として
\begin{equation}\label{eq:nfm-residual-point-j}
 j(E_{\mathrm{pt}})=j(C_{\mathrm{res}}),\qquad
 j_{\mathrm{pt}}=j_{\mathrm{act}}\circ\Psi_{\mathrm{res}}
\end{equation}
が成立する。
特に、$\Psi_{\mathrm{res}}$ の各幾何学的ファイバー上で
楕円点曲線の幾何学的同型類は一定である。
\end{theorem}

\begin{proof}
相対点関手の降下により $E_{\mathrm{pt}}$ は元の stack 上の曲線族である。
$C_{\mathrm{res}}$ も方程式直線と support の降下から定まる。
両者の周囲平面には Brauer--Severi 平面を許す。
genus-one 曲線の $j$ は Jacobian の $j$ として定義でき、原点を要しない。

補題\ref{lem:nfm-dominant-central-family}の支配的な中心 $h$ 族に引き戻す。
系\ref{cor:nfm-projective-formal-compatibility}の相対比較により、
$C_{\mathrm{res}}$ は形式的な次数一同定の下で
$V(W_h|_{U=0})$ に移る。
前命題によってこれは $V(J_g)$ であり、点曲線に一致する。
従って両方の $j$ はこの支配族の $k((h))$ 上で等しい。
支配性による関数体の単射から元の stack の一般点で等しく、
整な開上の正則関数なので全域で等しい。
全 O-降下はそれぞれの曲線族の降下から従う。
代数閉体上で滑らかな genus-one 曲線の同型類は $j$ で決まるため、
ファイバーに関する結論を得る。
\end{proof}

\begin{proposition}[平面同型の十八重被覆と Jacobian 同型の二重被覆]
\label{prop:nfm-point-residual-isom-cover}
前定理の開でさらに $j\notin\{0,1728\}$ とする。
$J_E=\operatorname{Pic}^0(E_{\mathrm{pt}})$、
$J_C=\operatorname{Pic}^0(C_{\mathrm{res}})$ と置くと
\[
 \mathfrak D_J=\underline{\operatorname{Isom}}_{\mathrm{gp}}(J_E,J_C)
       \longrightarrow\mathfrak M_n^{\mathrm{res}}
\]
は標準的な有限 étale 階数二の $\mu_2$-torsor である。
また、周囲平面も含めた同型の層
\[
 \mathfrak I=\underline{\operatorname{Isom}}_{\mathrm{plane}}
  \bigl((\mathbf P_{\mathrm{pt}},E_{\mathrm{pt}}),
        (\mathbf P_{\mathrm{res}},C_{\mathrm{res}})\bigr)
\]
は有限 étale 階数十八である。
曲線同型が誘導する共変な Jacobian 同型による射
$\mathfrak I\to\mathfrak D_J$ は $J_E[3]$-torsor であり、
従って
\[
 \mathfrak I/J_E[3]\simeq\mathfrak D_J.
\]
ここで ``plane'' は次数三偏極の\emph{類}を保つ同型を意味し、
$\mathcal O(1)$ の同型やその scalar を追加データとして選んではいない。
これらは元の stack に対して表現可能な有限 étale 射である。
\end{proposition}

\begin{proof}
全て étale 局所的に検証できる。
Jacobian を短い Weierstrass 式
$y^2=x^3+Ax+B$ と $y^2=x^3+A'x+B'$ で表す。
$j\ne0,1728$ なので $A,B,A',B'$ は単元である。
$j$ の等式は
$(A'/A)^3=(B'/B)^2$ を与える。
\[
 u^2=\frac{B'/B}{A'/A}
\]
と置く有限 étale 階数二の被覆上で
$u^4=A'/A$、$u^6=B'/B$ となり、
$(x,y)\mapsto(u^2x,u^3y)$ は group 同型である。
この二つが全ての group 同型であり、その差は反転である。
従って $\mathfrak D_J$ は表示に依存せず貼り合わさる $\mu_2$-torsor となる。

次に $\mathfrak D_J$ 上で Jacobian 同型 $\theta$ を固定する。
étale 局所的に両 genus-one 曲線に原点を選べば、
$\theta$ を誘導する曲線同型 $f_0:E_{\mathrm{pt}}\to C_{\mathrm{res}}$ がある。
さらに周囲の Brauer--Severi 平面を split し、次数三線束を取る。
$\theta$ を誘導する別の曲線同型は $f_0\circ t_q$ と一意に書ける。
次数三線束 $L$ に対する
$q\mapsto t_q^*L\otimes L^{-1}$ は符号を除いて $[3]$ である。
従って二つの偏極類を一致させる条件は $[3]q=a$ という形であり、
その解は有限 étale 階数九の $J_E[3]$-torsor をなす。
偏極類が一致すれば完全線型系によって周囲平面の同型へ一意に延長する。
基底から引いた線束による差は射影束を変えず、余分な $\mathbf G_m$ は生じない。
これらの同定は étale 重なり上で自然なので、
Brauer--Severi 平面と元の moduli stack に降下する。
有限 étale 性、階数、および表示した商の記述が従う。
\end{proof}

\subsection{尺度線束と active 三次の Hodge 線束}

以下では $X_{\mathrm{res}}$ の active 平面切断が滑らかな開を取る。
これは定理\ref{thm:nfm-projective-residual-morphism}の構成で既に取った非空開に含まれる。
表示 gerbe 上で
\[
 C_{\mathrm{res}}
   =X_{\mathrm{res}}\cap\mathbf P(K_{\mathrm{res}})
\]
と置く。方程式直線 $\mathcal L$ の
$\operatorname{Sym}^3K_{\mathrm{res}}^*$ への制限は直線部分束である。
射影 ambient と方程式直線の降下から、$C_{\mathrm{res}}$ は元の stack 上の
滑らかな genus-one 曲線族へ降下する。
ここで原点、flex、相対次数三線束の大域的な選択は仮定しない。

\begin{proposition}[尺度と Hodge 線束の標準的な双対性]
\label{prop:nfm-scale-hodge-line}
$f:C_{\mathrm{res}}\to\mathfrak M_n^{\mathrm{res}}$ を上の曲線族とし、
$\lambda_{\mathrm{res}}=f_*\omega_{C_{\mathrm{res}}/\mathfrak M_n^{\mathrm{res}}}$ と置く。
任意基底変更および全 O-同値に自然な同型
\[
 \lambda_{\mathrm{res}}\simeq
 (\mathcal L\otimes\det K_{\mathrm{res}})^*
 =\mathscr T^*
\]
がある。従って、残余 Poisson 直線の尺度を選ぶことは、active 三次の
Hodge 線束の双対を自明化することに等しい。
\end{proposition}

\begin{proof}
fpqc 局所的に表示を選び、$K=K_{\mathrm{res}}$ と略記する。
$\mathbf P(K)$ は $K$ の直線を parametrise する規約なので、
Euler 完全列から
\[
 \omega_{\mathbf P(K)/S}
   \simeq\mathcal O_{\mathbf P(K)}(-3)\otimes(\det K)^*
\]
である。曲線は $\mathcal O(3)\otimes\mathcal L^*$ の正則切断で定義され、
相対随伴により
\[
 \omega_{C/S}\simeq f^*(\mathcal L\otimes\det K)^*.
\]
滑らかな平面三次は幾何学的に連結な genus-one 曲線であり、
$f_*\mathcal O_C=\mathcal O_S$ は任意底変更と両立する。
射影公式によって所要の同型を得る。

この計算は方程式の生成元と $K$ の基底を使わない。
また、表示の変更 $\mathcal L\mapsto\mathcal L\otimes M^3$、
$K\mapsto K\otimes M^{-1}$ に対して右辺の重みは零である。
従って同型は表示 gerbe から元の stack へ降下する。
基底変更との両立は相対随伴、$f_*\mathcal O_C=\mathcal O_S$ とその底変更から従う。
\end{proof}

局所的には、$G\in\mathcal L$ と $P\in\det K$ の非零生成元を選び、
$P$ に双対な体積形式を $\nu_P$、Euler ベクトル場を $E$ とすれば、
対応する Hodge 基底は
\[
 \operatorname{Res}_{C}
       \frac{\iota_E\nu_P}{G|_K}
\]
である。$G\mapsto aG$、$P\mapsto bP$ の下でこの基底は $(ab)^{-1}$ 倍となる。
Poisson tensor $\iota_{dG}P$ が $ab$ 倍となることと双対である。

\begin{theorem}[尺度を正規化する一様な有限 étale 被覆]
\label{thm:nfm-modular-scale-covers}
上の開には標準的な可逆切断
\[
 \Delta\in\Gamma(\mathfrak M_n^{\mathrm{res}},\mathscr T^{-12})
\]
があり、$\mathscr T^{12}$ を自明化する。
したがって
\[
 \mathfrak M_n^{(12)}
   =\{t\in\mathscr T^\times:\Delta(t^{\otimes12})=1\}
   \longrightarrow\mathfrak M_n^{\mathrm{res}}
\]
は有限 étale な $\mu_{12}$-torsor であり、その上では標準的な尺度付きの lift
\[
 \widetilde\Psi_{\Delta}:
 \mathfrak M_n^{(12)}\longrightarrow\mathfrak P_n^{\mathrm{dec}}
\]
を持つ。

さらに非空稠密開 $c_4c_6\ne0$、すなわち $j\notin\{0,1728\}$ では、
\[
 b=c_6/c_4\in\Gamma(\mathfrak M_n^{\mathrm{res}},\mathscr T^{-2})^\times
\]
を使い、同じ結論を有限 étale 階数二の標準的な被覆
\[
 \mathfrak M_n^{(2)}
   =\{t\in\mathscr T^\times:b(t^{\otimes2})=1\}
\]
上で得る。全ての構成は全 O-同値と任意基底変更に両立する。
この被覆の二つの局所切断は互いに反対符号の tensor を与える。
\end{theorem}

\begin{proof}
genus-one 族の Hodge 線束上の標準 modular forms を使う。
原点を持たない場合にも、fpqc 局所的に曲線の切断を取り、これを原点に選べる。
二つの原点の変更は平行移動であり、平行移動は正則一形式に恒等に作用する。
従って以下の切断は原点に依存せず降下する。

標数零で局所的な短い Weierstrass 式
$y^2=x^3+Ax+B$ と $\eta=dx/(2y)$ を取ると
\[
 c_4=(-48A)\eta^4,\qquad
 c_6=(-864B)\eta^6,\qquad
 \Delta=-16(4A^3+27B^2)\eta^{12}
\]
はそれぞれ $\lambda^4,\lambda^6,\lambda^{12}$ の切断である。
座標変更 $x'=u^2x$、$y'=u^3y$ の下で
$A'=u^4A$、$B'=u^6B$、$\eta'=u^{-1}\eta$ となるので、
これらは独立な局所表示から同じ切断を与える。
また $c_4^3-c_6^2=1728\Delta$、$j=c_4^3/\Delta$ であり、
滑らかな locus では $\Delta$ が至る所非零である。
判別式による $\lambda^{12}$ の標準的な自明化は既知であり、
\cite[Section~1, equations~(1-1)--(1-5), pp.~87--88]{NFMFultonOlsson}
にこの変換則が明記されている。

前命題で $\lambda=\mathscr T^*$ と同定する。
$\mathscr T$ の局所基底 $e$ に関して $\Delta=d e^{-12}$ と書けば、
被覆は $dz^{12}=1$ で与えられる。
$d,z,12$ は可逆なので有限 étale 階数十二であり、$t=ze$ が標準切断となる。
$c_4c_6\ne0$ の開では同様に $b=d_2e^{-2}$ と書き、
$d_2z^2=1$ が有限 étale 階数二の被覆を与える。
尺度 torsor の記述により、それぞれの $t$ が所要の lift を与える。

最後に、この一般開が非空であることを確認する。
固定した active 平面への全三次式の制限は全射であり、
任意の滑らかな ternary cubic $g$ は
$g(x,y,z)+\sum u_i^3$ という滑らかな全三次へ延長できる。
従って射影 Poisson 標的上の active 曲線の $j$ は一定でない。
$\Psi_{\mathrm{res}}$ の支配性から $c_4c_6\ne0$ の逆像も非空稠密開である。
\end{proof}

\begin{remark}[正規化と標準尺度の区別]
この定理は、元の stack 上で $\mathscr T$ の標準切断を構成したという意味ではない。
被覆上の切断は modular な正規化を満たすが、固定 formality が指定する
$\pi_h$ の実際の尺度との一致も含まない。
また被覆の次数二や十二は、この尺度正規化のための次数であり、
残余射の一般次数や比較核を定義する被覆の次数ではない。

曲線の同型類だけから微分の非零基底を functorial に選ぶ方法はない。
実際、楕円曲線の反転 $[-1]$ は Hodge 線に $-1$ で作用する。
三次偏極を加えても、幾何学的には flex を原点に取れば、この反転は偏極を保つ。
しかし全三次や非可換圏の追加データはこの反転を自己同型として持つとは限らない。
従ってこの観察を、元の $\mathfrak M_n^{\mathrm{res}}$ 上での
標準尺度の不存在へ読み替えてはならない。
\end{remark}

\subsection{点曲線の尺度から定まる符号付き残余射}

良い一ブロック O-marked stack の非空開をさらに縮め、
Hessian 補正 $\Phi$ と $\Psi_{\mathrm{res}}$ がともに定義され、
楕円点曲線 $E_{\mathrm{pt}}$ と残余 active 曲線 $C_{\mathrm{res}}$ が滑らかで、
両者の $j$ 不変量が $0,1728$ を避けるようにする。
さらに点 shift が固定点を持たない相対点関手の良い開に含める。
この開を $\mathfrak U$ と書く。
ここでの $\Phi$ は定理\ref{thm:nfm-algebraic-hessian-correction}の
尺度を含む代数的な射であり、形式復元と同一視してはいない。

\Needspace{6\baselineskip}
\begin{lemma}[楕円点曲線の標準微分]
\label{lem:nfm-point-hodge-scale}
$E_{\mathrm{pt}}$ の Hodge 線束を $\lambda_{\mathrm{pt}}$ と書くと、
全 O-同値と任意基底変更に両立する標準切断
\[
 \tau_{\mathrm{pt}}\in
 \Gamma(\mathfrak U,\lambda_{\mathrm{pt}}^*)^\times
\]
がある。局所的な Hessian 補正の表示 $\Phi(H,P)=(H^\sharp,P)$ の下では
\[
 \tau_{\mathrm{pt}}=(H^\sharp|_K)\otimes P
\]
である。従って $\lambda_{\mathrm{pt}}$ は標準的な非零一形式
$\omega_{\mathrm{pt}}=\tau_{\mathrm{pt}}^{-1}$ を持つ。
\end{lemma}

\begin{proof}
命題\ref{prop:nfm-correction-point-cubic}により、
$H^\sharp|_K$ は $E_{\mathrm{pt}}$ の方程式そのものである。
命題\ref{prop:nfm-scale-hodge-line}の相対随伴計算をこの active 曲線へ適用する。
$H^\sharp\mapsto aH^\sharp$ と $P\mapsto a^{-1}P$ の下で表示した tensor は不変であり、
座標変更と表示 gerbe の重みも相殺する。
従って $\Phi$ の全 O-同値への降下から所要の切断の降下を得る。
局所随伴および $\Phi$ の構成は任意基底変更と両立する。
\end{proof}

残余側の尺度線束を $\mathscr T=\lambda_{\mathrm{res}}^*$ と書く。
前節の modular forms を用い、
\begin{equation}\label{eq:nfm-point-residual-squared-scale}
 \begin{split}
 b_{\mathrm{res}}&=c_{6,\mathrm{res}}/c_{4,\mathrm{res}}
       \in\Gamma(\mathfrak U,\mathscr T^{-2})^\times,\\
 b_{\mathrm{pt}}&=
 (c_{6,\mathrm{pt}}/c_{4,\mathrm{pt}})
          (\tau_{\mathrm{pt}}^{\otimes2})
       \in\Gamma(\mathfrak U,\mathcal O)^\times,\\
 \eta_{\mathrm{res}}&=b_{\mathrm{pt}}b_{\mathrm{res}}^{-1}
       \in\Gamma(\mathfrak U,\mathscr T^{\otimes2})^\times
 \end{split}
\end{equation}
と定める。$(c_6/c_4)$ の同じ規約を両側で用いるため、
この比は両側の modular forms に共通の数値規約にも依存しない。

\begin{theorem}[符号を除いた標準尺度と支配的な残余射]
\label{thm:nfm-canonical-signed-residual-map}
定理\ref{thm:nfm-projective-residual-morphism}の固定入力の下で、
任意の $n\ge3$ について、式\eqref{eq:nfm-point-residual-squared-scale}は
標準的な代数的射
\[
 \Psi_{\mathrm{res}}^{\pm}:\mathfrak U\longrightarrow
       [\mathfrak P_n^{\mathrm{dec}}/\mu_2]
\]
を与える。ここで $-1\in\mu_2$ は Poisson tensor の符号を変える。
この射は全 O-同値に降下し、任意基底変更と両立し、
尺度を忘れると $\Psi_{\mathrm{res}}|_{\mathfrak U}$ となる。
さらに支配的かつ一般有限であり、source と target を非空開へ制限すれば
有限 \'etale 全射にできる。

標準的な有限 \'etale 階数二の被覆
\begin{equation}\label{eq:nfm-point-normalized-scale-cover}
 \mathfrak D=
 \{t\in\mathscr T^\times:t^{\otimes2}=\eta_{\mathrm{res}}\}
 \longrightarrow\mathfrak U
\end{equation}
上では、tautological な切断 $t$ により尺度を含む標準射
\[
 \widetilde\Psi_{\mathrm{res}}:\mathfrak D
       \longrightarrow\mathfrak P_n^{\mathrm{dec}}
\]
を得る。これは $\Psi_{\mathrm{res}}^{\pm}$ の
$\mathfrak P_n^{\mathrm{dec}}\to[\mathfrak P_n^{\mathrm{dec}}/\mu_2]$
に沿う基底変更である。
従ってこれも支配的・一般有限であり、適切な非空開上で有限 \'etale 全射となる。
\end{theorem}

\begin{proof}
射影標的上の普遍尺度線束を $\mathscr T_0$ と書く。
符号商は $(\mathscr T_0^{\otimes2})^\times$ であるから、
切断 $\eta_{\mathrm{res}}$ が所要の射を定める。
全 O-同値と基底変更への両立は、Hodge 線束、modular forms、
$\tau_{\mathrm{pt}}$ および $\Psi_{\mathrm{res}}$ の自然性から従う。
平方根の被覆は局所的に $t^2=a$、$a$ は単元で与えられるので、
標数零では有限 \'etale 階数二である。

支配性を示す。命題\ref{prop:nfm-projective-residual-dominance}と同じ
支配的な中心 $h$ 族を取り、$F+G$ と support を一般に動かす。
まず射影標的を自明な慣性を持つ稠密 scheme 開 $B_0$ へ制限する。
そのような開は、一般の滑らかな三次の射影自己同型が自明であること
\cite[Theorem~1.5 and the preceding paragraph]{NFMPoonenAutomorphisms}から存在する。
既存の射影残余射の支配性により source の逆像も非空であり、
一般の特殊対 $(F+G,K)$ は $B_0$ に属する。
$g=F|_K$ とする。補題\ref{lem:point-relative-casimirs}により、
局所表示の primitive 三次の active 部分を $c(h)g$ と書ける。
ここで $c(h)\in k((h))^\times$ であり、方程式の生成元を替えれば $c(h)=1$ としてよい。
元の二次関係が $h\iota_{dg}P_0$ なので、
対応する determinant 生成元は $(h/c(h))P_0$ となる。
命題\ref{prop:nfm-correction-point-cubic}の Hessian 公式を使うと、
補正後の active 方程式は $c(h)J_g$ であり、二つの因子 $c(h)$ は相殺する。
この議論は支配族の中心持上げ $\Omega_h+G$ の選択によらず、
点曲線の方程式と尺度は厳密に
\[
 J_g=g+\frac{h^2}{24}\det\operatorname{Hess}(g),\qquad
 \tau_{\mathrm{pt}}=h\,(J_g\otimes P_0)
\]
と書ける。従って $b_{\mathrm{pt}}=h^2v_h$ であり、$v_h$ は
$h=0$ でも非零な形式単元である。
一方、相対 MF flow による形式同定の下で残余値は
$(V(W_h),K)$ となり、$W_h\equiv F+G\pmod h$ である。
Hodge 線束の随伴同定をこの形式族に適用すると、
$b_{\mathrm{res}}$ も特殊点で非零な基底を与える。
そのため $\eta_{\mathrm{res}}$ は $\mathscr T_0^{\otimes2}$ の全空間へ延長し、
そこで $h^2$ と形式単元の積となる。
特殊点は、射影標的の一般点における零切断である。

$B_0$ 上の $\mathscr T_0^{\otimes2}$ の全空間における
$\Psi_{\mathrm{res}}^{\pm}$ の像の閉包を $Z$ とする。
source は既約なので $Z$ は既約である。
上の形式族から $Z$ は零切断の一般点を含み、従って零切断全体を含む。
また像は零切断の外にあるので、$Z$ は零切断と等しくない。
零切断は整な全空間の既約 Cartier 因子なので、$Z$ は全空間に等しい。
これで支配性を得る。
ここでは次数一空間の形式同定を代数的と仮定していない。
像の閉包と零切断は射影 frame 変更に不変な intrinsic な対象であり、
形式同定はその閉集合条件と特殊化を変えない。
実際、$W_h$ の有限個の三次係数と上の尺度の平方は $k[[h]]$ に属し、
全空間への実際の $\operatorname{Spec}k[[h]]$ の射を与える。
その一般点だけを intrinsic な像と形式 frame で同定すれば十分である。

両 moduli の次元は
$\binom{n+2}{3}+3n-2$ なので、支配性から一般有限性を得る。
一般の滑らかな三次 $n$-fold の射影自己同型が自明である非空開が存在する
\cite[Theorem~1.5 and the preceding paragraph]{NFMPoonenAutomorphisms}。
source は $\mathfrak P_n^{\mathrm{good}}$ と同型なので、source と target の
自明な慣性を持つ非空開へ同時に制限できる。
さらに非空開を取れば schemes の間の一般有限射として扱える。
射を準有限かつ分離的とし、Zariski の主定理による有限化の中で
source の scheme 理論的閉包を取る。
境界は一般点の上では空であり、その有限な閉像を除けば元の射が有限になる。
標数零なので分岐 locus を除けば有限 \'etale 全射となる。
この結論を基底変更すれば、$\widetilde\Psi_{\mathrm{res}}$ にも同じ結論を得る。
\end{proof}

この支配性の証明は、点曲線と残余 active 曲線の $j$ の一致を必要としない。
一致を用いると、上の二重被覆には次のより直接的な意味がある。

\begin{proposition}[二重被覆の Jacobian による解釈]
\label{prop:nfm-point-normalized-jacobian-cover}
定理\ref{thm:nfm-residual-point-j}の $j$ の一致の下で、
式\eqref{eq:nfm-point-normalized-scale-cover}の被覆は、
命題\ref{prop:nfm-point-residual-isom-cover}の $\mathfrak D_J|_{\mathfrak U}$ と標準的に同型である。
すなわち
\[
 \mathfrak D\simeq
 \operatorname{Isom}_{\mathrm{gp}}
 \bigl(\operatorname{Jac}(E_{\mathrm{pt}}),
       \operatorname{Jac}(C_{\mathrm{res}})\bigr)
\]
がある。右辺は、両楕円曲線の原点を保つ群スキーム同型の sheaf である。
tautological な同型 $\theta$ の下で、残余の尺度は
\[
 t=\tau_{\mathrm{pt}}\circ\theta^*,\qquad
 \omega_{\mathrm{res}}=(\theta^*)^{-1}\omega_{\mathrm{pt}}
\]
で与えられる。従って二つの局所同型 $\theta$ と $[-1]\circ\theta$ は
互いに反対符号の残余 Poisson tensor を与える。
元の開上でも、尺度の平方に関する恒等式
\[
 c_{4,\mathrm{res}}(\eta_{\mathrm{res}}^{\otimes2})
   =c_{4,\mathrm{pt}}(\tau_{\mathrm{pt}}^{\otimes4}),\qquad
 c_{6,\mathrm{res}}(\eta_{\mathrm{res}}^{\otimes3})
   =c_{6,\mathrm{pt}}(\tau_{\mathrm{pt}}^{\otimes6})
\]
が成り立つ。すなわち、点曲線の微分から得る両 Weierstrass 不変量を保つ。
\end{proposition}

\begin{proof}
genus-one 曲線とその Jacobian の Hodge 線束は標準的に同型である。
実際、fpqc 局所的に曲線の切断を取り、原点として Jacobian と同定すればよい。
原点の変更は平行移動であり、不変微分への作用は恒等なので、同定は降下する。

$j\notin\{0,1728\}$ の楕円曲線の原点を保つ自己同型群は
$\mu_2$ であり、同じ $j$ を持つ楕円族は \'etale 局所的に同型である。
従って表示した Isom sheaf は有限 \'etale な $\mu_2$-torsor となる。
これは短い Weierstrass 式から直接確認でき、moduli の記述としては
\cite[proof of Corollary~2.4, p.~91]{NFMFultonOlsson}
の $\mathcal M_{1,1}\to\mathbb A^1_j$ がこの開上で
$\mu_2$-gerbe となる事実である。

この torsor 上で $t=\tau_{\mathrm{pt}}\circ\theta^*$ と定める。
modular forms の同型への自然性により
\[
 (c_{6,\mathrm{res}}/c_{4,\mathrm{res}})(t^{\otimes2})
 =b_{\mathrm{pt}}
\]
となり、従って $t^{\otimes2}=\eta_{\mathrm{res}}$ である。
$[-1]$ は不変微分に $-1$ で作用するので、これは二つの
$\mu_2$-torsor の同変射であり、よって同型である。
最後の二つの恒等式も modular forms の同型への自然性から従い、
符号に不変なので元の開へ降下する。
この証明は対象曲線の原点や相対次数三線束の選択を使わない。
\end{proof}

\begin{remark}[他の尺度正規化と比較の範囲]
modular な二重被覆
$\{t:b_{\mathrm{res}}(t^2)=1\}$ と、この点曲線により正規化された被覆
$\mathfrak D=\{t:b_{\mathrm{res}}(t^2)=b_{\mathrm{pt}}\}$ は区別する。
両者の差は $b_{\mathrm{pt}}$ の平方根の Kummer torsor である。
前者は射影標的だけから得られる有限個の尺度を選ぶのに対し、後者は
量子側の点曲線が持つ尺度を用い、相対次元一の方向を支配的に写す。

Jacobian の同型は、genus-one torsor である曲線自体の同型や、
その平面偏極の同型を意味しない。
また本構成は、固定 formality による $\pi_h$ との全次数の尺度の一致、
元の $\mathfrak U$ 上の符号選択、一般次数の計算を主張しない。
比較核の有限 \'etale 被覆を $\mathfrak D$ に引き戻せば
同じ $X_{\mathrm{res}}$ の Kuznetsov 成分との比較を保つが、
$\mathfrak D$ 自身への比較核の降下が新たに得られたわけではない。
\end{remark}

\begin{corollary}[点 shift の二重被覆上での移送]
\label{cor:nfm-point-shift-residual-transfer}
上の固定点なしの開で、点曲線上の shift $\sigma_{\mathrm{pt}}$ は、
一意な切断 $p_{\mathrm{pt}}\in J_E$ による平行移動である。
二重被覆 $\mathfrak D$ 上で
\[
 p_{\mathrm{res}}=\theta(p_{\mathrm{pt}}),\qquad
 \sigma_{\mathrm{res}}=t_{p_{\mathrm{res}}}
\]
と置くと、十八重被覆上の全ての tautological な平面同型 $f$ は
$f\sigma_{\mathrm{pt}}f^{-1}=\sigma_{\mathrm{res}}$ を満たす。
deck 変換は $\sigma_{\mathrm{res}}$ をその逆へ移す。
\end{corollary}

\begin{proof}
良い点関手開では $\sigma_{\mathrm{pt}}$ は固定点を持たない。
幾何学的ファイバー上で、genus-one 曲線の自己同型を原点の選択により
$t_p\circ\alpha$ と書く。$\alpha\ne1$ なら $1-\alpha$ は非零 isogeny であり、
固定点方程式は解を持つ。従って $\alpha=1$ である。
相対 Jacobian の group 自己同型の層はこの開で étale なので、
この恒等性は族上でも成立する。
曲線はその Jacobian の torsor であり、Jacobian に恒等を誘導する自己同型は
一意な切断による平行移動である。
残る主張は $f$ の誘導する Jacobian 同型が $\theta$ であることと、
平行移動同士が可換することから従う。
\end{proof}


\begin{proposition}[点曲線で規格化した形式 tensor]
\label{prop:nfm-point-scale-formal-comparison}
生成元固定模型で、$W_h|_{U=0}=b_hJ_g$ と書く。
点曲線と残余 active 曲線の恒等同型を与える形式分枝上では
\begin{equation}\label{eq:nfm-point-scale-formal-leading}
 \tau_{\mathrm{res}}^{\mathrm{pt}}
    =\frac{h}{b_h}\,W_h\otimes P_0,
 \qquad
 \Pi_{\mathrm{res}}^{\mathrm{pt}}
    =\frac{h}{c_h}\,\pi_h(F),
 \qquad
 \pi_h^{\mathrm{act}}(g)=c_h\iota_{dJ_g}P_0.
\end{equation}
特に $\tau_{\mathrm{res}}^{\mathrm{pt}}=h(F+G)\otimes P_0+O(h^2)$ である。
\end{proposition}

\begin{proof}
点曲線の尺度は $hJ_g\otimes P_0$ であり、
相対随伴で $W_h|_K=b_hJ_g$ に対する微分を同じ微分へ合わせると、
残余三次全体の尺度は $(h/b_h)W_h\otimes P_0$ となる。
命題\ref{prop:nfm-residual-active-point-cubic}の証明で
$a_hb_h=c_h$ なので、$\pi_h=a_h\iota_{dW_h}P_0$ と比較すれば第二式を得る。
最初の非零項は $b_0=1$、$W_0=F+G$ から従う。
\end{proof}

この規格化は固定 harmonic tensor の尺度と一般には異なる。
例えば命題\ref{prop:nfm-active-fifth-covariant}の純 active 五次公式、
命題\ref{prop:nfm-normalized-parity}の偶奇性と $J_g$ の Casimir 性から
\[
 c_h=h+\frac{S(g)}{80}h^5+O(h^7)
\]
であり、$S(g)\ne0$ の一般係数で $h/c_h\ne1$ となる。
ここで $S$ は $S(xyz)=1$ で規格化した Aronhold 四次不変式である。
この比を代数的な単元とも仮定しない。
従って、点曲線による代数的射は
定理\ref{thm:nfm-normalized-tensor-transcendence}の固定 tensor の超越性と両立する。
移送した $\sigma_{\mathrm{res}}$ は今回曲線上に定めた自己同型であり、
Kuznetsov rotation の制限を同定したものではない。
定常点成分の同定、比較核の一意性や大域降下も含まない。
十八、二、および補正射 $\Phi$ の九はそれぞれ異なる被覆・射の次数であり、
$\Psi_{\mathrm{res}}^\pm$ の一般次数は未決定である。

\begin{remark}[元の開上の非標準的な full lift]
\label{rem:nfm-existential-full-residual-map}
符号を標準的に選ぶ問題とは別に、元の O-marked stack のさらに小さい非空開上で、
尺度を含む支配的・一般有限な lift を\emph{選ぶ}ことはできる。
source と射影標的を自明な慣性を持つ整な scheme 開へ制限し、関数体を $L/K$ とする。
命題\ref{prop:nfm-projective-residual-dominance}より
$\operatorname{trdeg}_K L=1$ である。
標的の尺度線束を一般点で自明化し、$K$ 上超越的な $f\in L$ を選べば、
$K(z)\hookrightarrow L$、$z\mapsto f$ は有限体拡大となる。
この有理的な尺度切断を非空開へ広げればよい。
前定理と同じ有限化により、さらに有限 étale 全射にもできる。
underlying cubic は同じ $X_{\mathrm{res}}$ なので被覆上の Ku 比較も保つ。
ただし $f$ を $f^m$ に替えると一般次数は $m$ 倍となり、
この存在論的構成は標準次数、点曲線の微分との整合、
固定 formality の尺度との一致を与えない。
\end{remark}

% !TEX root = ../nc-cubic-n-fold.tex
\section{点の平行移動と比較核の Picard 類}
\label{sec:nfm-kernel-point-class}

最高次線束型の比較核について、スカラー障害は次数零 Picard 類
$m_P$ の表現障害である。
本節では楕円曲線の rotation の共役式を求め、
可換特殊ファイバーを含む非可換 active 商の相対 THCR/window1 比較を構成する。
これにより、固定入力と良い一ブロックの仮定の下で、
点の平行移動と $m_P$ を結ぶ追加型の非空性とスカラー障害の具体式を証明する。
代数的な active 同値核の assignment や、full 核の同型類の選択を
無条件に構成するものではない。

\subsection{平行移動・線束と rotation の共役}

楕円曲線の点 $b$ を次数零線束 $P_b=\mathcal O_E(b-0)$ と同一視する。
以下の式はまず代数閉体上で証明し、必要な等式を忠実平坦降下して用いる。
右加群の twisted homogeneous coordinate ring の規約は
\cite[\S3, equation~(3.1)]{NFMStaffordVdB}に合わせる。
すなわち $\sigma$ が点加群の shift を表すとき、
coherent sheaf 上の次数 shift は
\begin{equation}
\label{eq:nfm-right-thcr-shift}
 s=\sigma_*\circ(-\otimes L)
\end{equation}
である。特に $s(k(x))=k(\sigma(x))$ となる。
$\sigma^*$ と取り違えると、以下の平行移動の符号が逆になる。

\begin{lemma}[楕円曲線の rotation の共役式]
\label{lem:nfm-elliptic-rotation-conjugacy}
$E$ を標数零の代数閉体上の楕円曲線、$L$ を次数三の線束とする。
\[
 U_{a,b}=t_a^*\circ(-\otimes P_b),\qquad
 \rho=T_{\mathcal O_E}\circ(-\otimes L)
\]
と置く。関手の同型類について
\begin{equation}
\label{eq:nfm-elliptic-rotation-conjugacy}
 \rho U_{a,b}\rho^{-1}\simeq U_{-2a-b,\,3a+b}
\end{equation}
が成り立つ。
$\sigma=t_p$ とし、右加群の規約に従って
$\rho_p=T_{\mathcal O_E}\circ\sigma_*\circ(-\otimes L)$ と置くと、
\begin{equation}
\label{eq:nfm-elliptic-point-intertwiner}
 U_{a,b}\rho_pU_{a,b}^{-1}\simeq\rho
 \quad\Longleftrightarrow\quad
 3a=0,\qquad b=p.
\end{equation}
特に $\rho$ と可換する $U_{a,b}$ は $U_{a,0}$、$a\in E[3]$ に限られ、
これらは次数零 Picard 類を変えない。
\end{lemma}

\begin{proof}
平行移動は次数零 Picard 類に自明に作用するので、
$U_{a,b}$ の同型類の積は $(a,b)$ の加法である。
$t_a^*L\otimes L^{-1}\simeq P_{-3a}$ より
\[
 \operatorname{Ad}_{-\otimes L}(a,b)=(a,b+3a).
\]
また
\[
 \operatorname{Ad}_{T_{\mathcal O_E}}(a,b)=(a-b,b)
\]
である。後式の graph と線束を確認する。
twist の Fourier--Mukai 核は
$\operatorname{Cone}(\mathcal O_{E\times E}\to\mathcal O_\Delta)
=\mathcal O_{E\times E}(-\Delta)[1]$ であり、
$T_{\mathcal O_E}^{-1}k(x)=\mathcal O_E(x)$ となる。
従って共役 $T_{\mathcal O_E}U_{a,b}T_{\mathcal O_E}^{-1}$ は
$k(x)$ を $k(x-a+b)$ へ送り、$\mathcal O_E$ を $P_b$ へ送る。
点層を shift なしの点層へ送る同値の核は、その台である graph 上の線束である。
この場合の graph は $t_{a-b}$ の pullback、線束は $P_b$ なので、所要の式を得る。
これは上記の twist 核を使った畳込み計算とも一致する。

二つの共役作用の積は
\[
 A=\begin{pmatrix}-2&-1\\3&1\end{pmatrix}
\]
であり、\eqref{eq:nfm-elliptic-rotation-conjugacy}を与える。
平行移動は $\mathcal O_E$ を保つので、その pullback は
$T_{\mathcal O_E}$ と可換する。
従って $\rho_p=U_{-p,0}\rho$ であり、intertwining の条件は
\[
 (I-A)\binom ab=\binom p0,
 \qquad\text{すなわち}\qquad
 3a+b=p,\quad -3a=0
\]
に等しい。これで全ての主張が従う。
\end{proof}

この補題は、全自己同値群の分類ではない。
平行移動と次数零線束の tensor からなる部分について、
rotation を保ったまま Picard 類を調整できる範囲を決定している。
特に、この部分の自己同値だけを用いて任意の $m_P$ を零にすることはできない。

\begin{lemma}[右楕円座標環の window1 と内部 shift]
\label{lem:nfm-twisted-elliptic-window}
$A=B(E,L,\sigma)$ を右 twisted homogeneous coordinate ring とする。
$\operatorname{qgr}A\simeq\operatorname{coh}E$ の下で
$\pi A\leftrightarrow\mathcal O_E$、$(1)\leftrightarrow s$ は
\eqref{eq:nfm-right-thcr-shift}に従うものとする。
$A$ に Gorenstein parameter 零の Orlov 同値が適用でき、
window1 を $\Phi_1=q_{\mathrm{sg}}R\Gamma_{\ge1}$ と取る。
このとき
\begin{equation}
\label{eq:nfm-twisted-active-window}
 \kappa_A\longleftrightarrow\mathcal O_E[1],\qquad
 \Phi_1^{-1}(1)\Phi_1=T_{\mathcal O_E}\circ s.
\end{equation}
\end{lemma}

\begin{proof}
この window の存在については
\cite[Theorem~2.5(iii)]{NFMOrlovSingularities}を用いる。
$m>0$ に対して $s^m\mathcal O_E$ は次数 $3m$ の線束なので、
その高次 cohomology は消える。
従って $R\Gamma_{\ge1}\mathcal O_E=A_{\ge1}$ であり、
$0\to A_{\ge1}\to A\to k\to0$ を特異圏へ写すと第一式を得る。
任意の $F$ に対し
\[
 (R\Gamma_{\ge1}F)(1)=R\Gamma_{\ge0}(sF)
\]
であり、次数零の切断は三角形
\[
 R\Gamma_{\ge1}(sF)\longrightarrow R\Gamma_{\ge0}(sF)
 \longrightarrow R\operatorname{Hom}(\mathcal O_E,sF)\otimes k
 \longrightarrow
\]
を与える。
特異圏へ写し、第一式を用いて三角形を回転すると、評価写像の三角形
\[
 R\operatorname{Hom}(\mathcal O_E,sF)\otimes\mathcal O_E
 \longrightarrow sF\longrightarrow\Phi_1^{-1}(1)\Phi_1(F)
 \longrightarrow
\]
を得る。この cone は $T_{\mathcal O_E}(sF)$ である。
従って剰余体と内部 shift の双方について同じ window の規約を得た。
この切断の議論は $A$ の可換性を用いない。
\end{proof}

\subsection{可換特殊ファイバーを含む相対楕円座標環と window1}

以下では $k$ を標数零の体、$R=k[[h]]$、$K=k((h))$ とする。
$g\in k[x,y,z]_3$ は滑らかな平面三次を定め、
\begin{equation}\label{eq:nfm-active-dvr-algebra}
 D_R=R\langle x,y,z\rangle/
 \bigl([y,z]-h\operatorname{sym}g_x,
       [z,x]-h\operatorname{sym}g_y,
       [x,y]-h\operatorname{sym}g_z\bigr)
\end{equation}
の各次数が有限自由で、剰余が $k[x,y,z]$ であると仮定する。
ここで $\operatorname{sym}$ は二次単項式の対称化である。
$\omega\in Z(D_R)_3$ は $\omega\bmod h=g$ を満たすものとし、
$A_R=D_R/(\omega)$ と置く。
この次数別自由性と指定中心元は、支配的な一ブロック族の active 制限で用いる PBW 入力である。

\begin{proposition}[指定中心商の相対 THCR 表示]
\label{prop:nfm-relative-active-thcr}
上の仮定の下で
\[
 J_g=g+\frac{h^2}{24}\det\operatorname{Hess}(g),\qquad
 E_R=V(J_g)\subset\mathbf P^2_R,\qquad L=\mathcal O_{E_R}(1)
\]
と置く。$E_R$ は滑らかな固有 genus-one 曲線族である。
右点加群の関係は、閉ファイバーで恒等となる自己同型
$\sigma:E_R\xrightarrow{\sim}E_R$ の graph を含み、
その一般ファイバーには固定点がない。
さらに
\begin{equation}\label{eq:nfm-relative-active-thcr}
 A_R\xrightarrow{\sim}
 B(E_R,L,\sigma)
 =\bigoplus_{m\ge0}H^0(E_R,L_m),\qquad
 L_m=\bigotimes_{j=0}^{m-1}\sigma^{j*}L
\end{equation}
が次数一の座標切断から定まり、
\[
 \operatorname{qgr}A_R\simeq\operatorname{coh}E_R,
 \qquad (1)\longleftrightarrow s=\sigma_*\circ(-\otimes L)
\]
となる。
代数同型の各次数は $R$-有限自由加群の同型であり、
剰余体・分数体への底変更後にも同じ点対応と THCR 表示を与える。
\end{proposition}

\begin{proof}
$u=(x,y,z)^{\mathsf t}$、$H(u)=\operatorname{Hess}(g)(u)$ とし、
$[u]_\times v=u\times v$ と置く。
右点関係の行列は
\[
 M_h(u)=[u]_\times-\frac h2H(u)
\]
である。対称行列 $H$ と三次の Euler 恒等式から
\begin{equation}\label{eq:nfm-active-point-matrix-determinant}
 \det M_h(u)=-\frac h2u^{\mathsf t}H(u)u
               -\frac{h^3}{8}\det H(u)=-3hJ_g(u).
\end{equation}
$J_g\bmod h=g$ なので、各次数の乗法 $J_g$ は剰余で単射である。
単元小行列式によってその余核は $R$-自由となり、$E_R$ は $R$-flat である。
滑らかでない locus の $\operatorname{Spec}R$ への像は proper 性により閉であり、
閉点を含まないので空である。これで滑らかな genus-one 族を得る。

$E_R$ 上で $\det M_h=0$ であり、閉ファイバーの $M_0(u)=[u]_\times$ は
各点で階数二である。同じ proper 性により階数一以下の locus は空であり、
$\ker M_h$ は直線部分束となる。その直線が定める射を
$\sigma:E_R\to\mathbf P^2_R$ と書く。
三次の偏微分の対称性より
\[
 M_{-h}(v)u=-M_h(u)v.
\]
従って $v=\sigma(u)$ に対し $hJ_g(v)=0$ であり、
$E_R$ の $R$-flatness から $J_g(v)=0$ である。
同じ構成を $-h$ で行うと逆射を得るため、$\sigma$ は $E_R$ の自己同型である。
$\ker M_0(u)=ku$ なので閉ファイバーで恒等となる。
一般ファイバーで $\sigma(u)=u$ なら
$M_h(u)u=-h\nabla g(u)=0$ となり、$V(g)$ の滑らかさに反する。

忠実平坦な有限非分岐拡大後、閉ファイバーの閉点を持ち上げて $E_R$ に原点を選べる。
原点を保つ楕円曲線自己同型の非分岐性から、
$\sigma\bmod h=\mathrm{id}$ は $\sigma=t_p$、$p\bmod h=0$ を意味する。
$E_R[3]$ は有限 \'etale なので、$3p=0$ なら $p=0$ となる。
一般ファイバーに固定点がないことから $p\ne0$ であり、
従って $\sigma_K^3\ne\mathrm{id}$ である。この不等式は原点の選択に依存しない。

$m>0$ について $\deg L_m=3m$ なので、$H^1(E_R,L_m)=0$、
$H^0(E_R,L_m)$ は階数 $3m$ の自由 $R$-加群で、形成は底変更に両立する。
積を $a\cdot b=a\otimes\sigma^{m*}b$（$\deg a=m$）と定める。
次数一の座標切断は上の graph 上の三つの関係を満たすので、
$D_R\to B(E_R,L,\sigma)$ を与える。

指定した $\omega$ が零に写ることを確認する。
一般ファイバーの関数体上で $L$ を有理的に自明化すると、
THCR は $K(E_K)[z;\sigma^*]$ に埋め込まれる。
$\omega$ の像が非零な $fz^3$ なら、次数一の切断 $az,bz$ との可換性から
\[
 f\sigma^{3*}(a)=a\sigma^*(f),\qquad
 f\sigma^{3*}(b)=b\sigma^*(f),\qquad
 \sigma^{3*}(a/b)=a/b
\]
を得る。$L$ は平面埋込みを与えるので、このような比は $K(E_K)$ を生成する。
従って $\sigma_K^3=\mathrm{id}$ となって矛盾する。
THCR の次数三は $R$-torsionfree なので、$\omega$ の像は $R$ 上でも零である。
この議論では中心三次元の一般公式やその scalar 規格化を用いない。

$\omega\bmod h=g$ の乗法は単射である。
再び各次数の単元小行列式により、$\omega$ は正則で、
$(A_R)_m$ は $m=0$ で階数一、$m>0$ で階数 $3m$ の自由 $R$-加群となる。
得られた $A_R\to B(E_R,L,\sigma)$ の剰余は
\[
 k[x,y,z]/(g)\xrightarrow{\sim}
 \bigoplus_{m\ge0}H^0(E_0,\mathcal O_{E_0}(m))
\]
という通常の平面三次の座標環の同型である。
次数ごとの Nakayama と等しい自由階数から
式\eqref{eq:nfm-relative-active-thcr}を得る。

相対的な ample pair の条件も確認する。
任意の coherent 層 $\mathcal F$ に対し $\mathcal G_m=\mathcal F\otimes L_m$ と置く。
曲線族の cohomological dimension は一なので、
\[
 0\to\mathcal G_m[h]\to\mathcal G_m\xrightarrow h h\mathcal G_m\to0,
 \qquad
 0\to h\mathcal G_m\to\mathcal G_m\to\mathcal G_m/h\mathcal G_m\to0
\]
から
\[
 H^1(E_R,\mathcal G_m)/hH^1(E_R,\mathcal G_m)
 \simeq H^1(E_0,(\mathcal F/h\mathcal F)\otimes\mathcal O_{E_0}(m))
\]
となる。右辺は $m\gg0$ で零であり、左辺の $H^1$ は有限 $R$-加群なので零となる。
同じ完全列から $H^0(\mathcal G_m)\to H^0(\mathcal G_m/h\mathcal G_m)$ は全射となる。
閉ファイバーの大きな twist の生成切断を持ち上げ、
評価写像の余核の proper support を用いれば $\mathcal G_m$ 自身も大域生成される。
$s^{-m}\mathcal O=L_m^{-1}$ なので、これらは $(\mathcal O,s)$ の ample-pair 条件である。
$\sigma^{-m*}L_m$ も剰余が $\mathcal O_{E_0}(m)$ であり、同じ議論は反対側にも適用できる。
一般 noetherian 底上の ample-pair 定理
\cite[Proposition~7.2 and Lemma~7.6]{NFMKeelerAmple}により qgr 同値を得る。
右加群での shift は\cite[Section~3, equation~(3.1)]{NFMStaffordVdB}の規約で
$s=\sigma_*\circ(-\otimes L)$ である。
\end{proof}

\begin{lemma}[三変数の相対 Gorenstein 条件]
\label{lem:nfm-active-dvr-gorenstein}
前命題の $D_R,A_R$ は左右 noetherian であり、
\[
 \operatorname{gldim}D_R\le4,\qquad
 R\underline{\operatorname{Hom}}_{D_R}(R,D_R)\simeq R(3)[-3],
 \qquad
 R\underline{\operatorname{Hom}}_{A_R}(R,A_R)\simeq R[-2]
\]
が左右で成立する。$A_R$ の左右 injective dimension は高々三であり、
Burke--Stevenson の意味で parameter 零の Gorenstein 環で、条件 $\chi$ を満たす。
これらの有限相対分解は剰余体と分数体への底変更に両立する。
\end{lemma}

\begin{proof}
入力B\cite{NFNCCubicB}の \texttt{local\_model.tex}、
補題 \texttt{lp:lem:noetherian}, \texttt{lp:lem:resolution} と
命題 \texttt{lp:prop:koszul}, \texttt{lp:prop:gorenstein} の次数別の証明を用いる。
三変数への適用は同入力 \texttt{explicit\_example.tex} の
補題 \texttt{gre:lem:elliptic-mf} にも明記されている。
必要な点を本設定で確認する。

$h$-随伴次数環は中央多項式拡大 $k[x,y,z,t]$ の商なので noetherian である。
斉次イデアルの記号の有限生成元を持ち上げ、
各内部次数の有限 $R$-加群に Krull 交叉定理を用いると $D_R$ の noetherian 性を得る。
任意のイデアルは内部次数の最高次項による有限降下で扱える。
多項式環の augmentation と対角の Koszul 分解を、各次数の分裂全射を用いて持ち上げる。
各段階の核は次数ごとに自由であり、最後の核は Nakayama によって零となる。
得られる二つの長さ三の分解の項は、それぞれ
$D_R(-j)^{\binom3j}$ と $D_R^e(-j)^{\binom3j}$ である。
対角分解と $\operatorname{gldim}R=1$ から $\operatorname{gldim}D_R\le4$ を得る。

augmentation 分解の双対を各内部次数で調べる。
剰余の Koszul 双対複体は内部次数 $-3$、cohomological degree 三以外で完全であり、
有界有限自由複体の単元成分の消去により、この完全性は $R$ 上へ持ち上がる。
内部次数 $-3$ の最後の項には $R$ だけが残るので、最初の導来 Hom 公式を得る。
中心正則列
$0\to D_R(-3)\xrightarrow\omega D_R\to A_R\to0$ による change of rings は
\[
 R\underline{\operatorname{Hom}}_{D_R}(M,D_R)
 \simeq R\underline{\operatorname{Hom}}_{A_R}(M,A_R)(3)[-1]
\]
を与える。$M=R$ から第二式、一般の $M$ から injective dimension の評価を得る。
全て左右で同じであり、$A_R$ は noetherian な $D_R$ の商である。
$R$-flatness と\cite[Definition~5.1 and Lemma~6.2]{NFMBurkeStevenson}から条件 $\chi$ が従う。
分解の各内部次数は $R$ 上で分裂しているので、述べた底変更にも両立する。
\end{proof}

\begin{proposition}[source の相対 window1 と二つの体への底変更]
\label{prop:nfm-relative-active-window-basechange}
命題\ref{prop:nfm-relative-active-thcr}の仮定の下で、
window1 と余核は $R$-線型 dg Morita 同値
\begin{equation}\label{eq:nfm-relative-active-window-proved}
 \operatorname{MF}^{\mathbf Z}(D_R,\omega)
 \simeq D^b_{\mathrm{sg,dg}}(\operatorname{gr}A_R)^{\mathrm{idem}}
 \simeq\operatorname{Perf}(E_R)
\end{equation}
を与える。これらは $R\to k$ および $R\to K$ の導来底変更に両立し、
augmentation stabilization と内部 shift はそれぞれ
\[
 \kappa_{A_R}\longleftrightarrow\mathcal O_{E_R}[1],\qquad
 (1)\longleftrightarrow
 T_{\mathcal O_{E_R}}\circ\sigma_*\circ(-\otimes L)
\]
となる。指定周期は選んだ $\omega$ の余核二拡張から移し、任意の scalar 倍で取り替えない。
\end{proposition}

\begin{proof}
前補題は\cite[Theorem~6.4(3)]{NFMBurkeStevenson}の全ての環論的仮定を満たす。
従って
\[
 \Phi_{1,R}=q_{\mathrm{sg}}R\Gamma_{\ge1}:
 D^b(\operatorname{coh}E_R)\longrightarrow D^b_{\mathrm{sg}}(\operatorname{gr}A_R)
\]
は同値である。
$R\Gamma_{\ge1}$ が有界有限生成の対象に値を持つことは、
同論文 Proposition~4.18 とその直前の有限生成性の議論による。
有限 cohomological dimension と $\chi$ は同論文 Lemmas~6.1--6.2 による。
derived section、内部次数切断、dg 商を合成した標準 enhancement の quasi-functor を用い、
全ての shift に対する完全忠実性から dg Morita 同値を得る。
$E_R$ は正則な二次元 scheme なので $D^b(\operatorname{coh}E_R)=\operatorname{Perf}(E_R)$ である。
余核による最初の dg 比較は入力Bの
\texttt{gre:lem:elliptic-mf} と \texttt{tw:lem:stabilization} の
exact 圏・有限 totalization による構成を使う。
その仮定である ambient の有限大域次元、中心元の正則性、商の有限 injective dimension は全て前補題で確認済みである。

底変更を関手そのものについて確認する。
$F\in\operatorname{Perf}(E_R)$ に対し、derived section の内部次数 $m$ 成分は
\[
 (R\Gamma_*F)_m=R\Gamma(E_R,s^mF)
   \simeq R\Gamma(E_R,F\otimes L_m)\quad(m\ge0)
\]
である。後者は $R$-perfect であり、その自然な比較写像は任意の底変更で擬同型となる
\cite[Lemma~36.30.4, Tag~0B91]{NFMStacksPerfectBasechange}。
この比較は切断の積に自然なので、$\ell=k,K$ について
\begin{equation}\label{eq:nfm-active-window-section-basechange}
 (R\Gamma_{\ge1}F)\otimes_R^{\mathbf L}\ell
 \xrightarrow{\sim}
 R\Gamma_{\ge1}(F\otimes_R^{\mathbf L}\ell)
 \quad\text{in }D(\operatorname{Gr}A_\ell)
\end{equation}
を得る。これは内部次数ごとの擬同型を組み合わせたものであり、
無限個の内部次数を有限個に切り捨てる議論ではない。
$F$ の固定された有限振幅と曲線族の cohomological dimension 一により、
cohomological な振幅は内部次数に一様である。
有限生成性は先に確認した Proposition~4.18 による。

$-\otimes_R^{\mathbf L}\ell$ は Tor 振幅高々一であり、
有界有限生成 graded 複体を同種の複体へ、perfect $A_R$-複体を perfect $A_\ell$-複体へ送る。
従って特異 dg 商への実際の底変更関手があり、
式\eqref{eq:nfm-active-window-section-basechange}は $\Phi_1$ の底変更図式を与える。
有限自由 MF の行列は、閉ファイバーで積が非零多項式 $g$ の乗法となるため単射である。
各内部次数で単元小行列式を使えば、その余核は $R$-自由である。
よって余核、周期分解、有限 totalization も $k,K$ への底変更と可換する。

これで二つの体上の比較との可換図式を、係数変更関手から直接構成した。
曲線側では $\operatorname{Perf}(E_R)\otimes_R^{\mathbf L}\ell
\simeq\operatorname{Perf}(E_\ell)$ が成り立つ。
実際、ample な $L$ の全ての整数冪は $R$ 上に定義され、
それらとその引戻しはそれぞれの perfect 圏を厚生成する。
これらの間の導来 Hom の底変更は、双対との tensor 積へ
Tag~0B91 を適用すれば分かる。
従って自然な係数変更は生成対象間で完全忠実で、ファイバー側も生成する。
各辺の同値と上の図式から、MF 側の自然な係数変更も Morita 同値を与える。
任意の特異圏の dg 商が任意の非平坦底変更に可換する、という一般則は用いていない。

最後に $R\Gamma_{\ge1}\mathcal O_{E_R}=(A_R)_{\ge1}$ と
$0\to(A_R)_{\ge1}\to A_R\to R\to0$ から augmentation の式を得る。
補題\ref{lem:nfm-twisted-elliptic-window}の切断三角形を $R$ 上で用いれば、
内部 shift は相対評価写像の cone
$T_{\mathcal O_{E_R}}\circ s$ へ移る。
指定周期の比較は、入力B \texttt{tw:eq:period-extension} の余核二拡張をそのまま使う。
\end{proof}

\subsection{恒等から変形した相対比較}

\begin{lemma}[形式的な楕円比較が定める Picard 類]
\label{lem:nfm-formal-elliptic-point-line}
$R$ を剰余標数零の完備 DVR とし、$E/R$ を滑らかな固有 genus-one 曲線、
$L$ を相対次数三の実際の線束とする。
$\sigma=t_p$ は閉ファイバーで恒等となる平行移動とする。
同じ $E/R$ 上に perfect な相対 Fourier--Mukai 核で与えられる同値 $F$ があり、
その核の閉ファイバーは $\mathcal O_\Delta$ とする。
さらに
\begin{equation}
\label{eq:nfm-formal-elliptic-intertwining}
 F\bigl(T_{\mathcal O_E}\circ\sigma_*\circ(-\otimes L)\bigr)F^{-1}
 \simeq T_{\mathcal O_E}\circ(-\otimes L)
\end{equation}
を仮定する。
このとき $F(\mathcal O_E)$ は次数零線束 $M$ であり、
\begin{equation}
\label{eq:nfm-formal-point-line-class}
 [M]=p\quad\text{in }\operatorname{Pic}^0_{E/R}(R)
\end{equation}
となる。
\end{lemma}

\begin{proof}
閉ファイバーの閉点を取り、有限非分岐拡大後にその点を $R$ 上の切断へ持ち上げる。
従って忠実平坦降下のため、原点を持つ楕円曲線の場合を証明すればよい。
核を $Q$ と書く。
第一射影 $E\times_R E\to E$ に関して $Q$ は relative perfect であり、
閉ファイバーの各点上の導来制限は長さ一の点層である。
relative Tor amplitude が $[0,0]$ となる locus は開で、閉ファイバー全体を含む。
その補集合の $\operatorname{Spec}R$ への像は proper 性により閉であり、
閉点を含まないので空である。
従って $Q$ は第一射影について平坦な層で、各ファイバーの長さは一である。
同じことが第二射影についても成り立つ。

長さ一の層の族の台は section であり、その上の層は線束である。
両射影についてこの記述を用いると、$Q$ は相対自己同型 $f$ の graph 上の線束となる。
閉ファイバーで $f$ は恒等、線束は自明であり、cohomological shift は零である。
原点を保つ楕円曲線自己同型の非分岐性から、
閉ファイバーで恒等となる $f$ は平行移動である。
従って $F$ の同型類は $U_{a,b}$ と書け、$a,b$ は閉点で零となる。
線束の次数は局所定数なので零である。

一般ファイバーの代数閉包上で
補題\ref{lem:nfm-elliptic-rotation-conjugacy}を適用すると、
\eqref{eq:nfm-formal-elliptic-intertwining}は $3a=0$、$b=p$ を与える。
これらは相対 Jacobian の切断の等式なので $R$ 上へ降りる。
$E[3]$ は有限 \'etale であり、閉点で零となる $a$ は零である。
従って $F(\mathcal O_E)$ の Picard 類は $p$ となる。
この等式は原点の選択に依存しないため、最初の非分岐拡大からも降下する。
\end{proof}

この補題は、perfect な複体が閉ファイバーで線束となることだけから
線束の Picard 類まで一定であると主張するものではない。
その類を決定する追加条件は rotation の intertwining であり、
一般には $M$ は自明線束ではなく $p$ を表す。

\subsection{点の平行移動を指定した full 比較核の型}

定理\ref{thm:nfm-kernel-active-line-bound}の split 表示、
$\kappa_{\mathrm{nc}}$、$q_Y$、window1 を固定する。
その底を $S$ と書く。
系\ref{cor:nfm-semiclassical-full-residual}の半古典規格化により、
命題\ref{prop:nfm-point-normalized-jacobian-cover}の Jacobian 同型の切断
$\theta_{\mathrm{sc}}:J_E\xrightarrow{\sim}J_C$ が選ばれる。
命題\ref{prop:nfm-point-residual-isom-cover}の平面同型被覆をこの切断で引き戻し、
\[
 \mathfrak I_{\mathrm{sc}}
   =\mathfrak I\times_{\mathfrak D_J,\theta_{\mathrm{sc}}}S
   \longrightarrow S
\]
と置く。これは階数九の有限 \'etale な $J_E[3]$-torsor である。
この被覆上では
\[
 \phi:E_{\mathrm{pt}}\xrightarrow{\sim}C_{\mathrm{res}},\qquad
 \operatorname{Jac}(\phi)=\theta_{\mathrm{sc}},\qquad
 p_{\mathrm{res}}=\theta_{\mathrm{sc}}(p_{\mathrm{pt}})
\]
となる。以下では $\theta=\theta_{\mathrm{sc}}$ と固定する。
$p_{\mathrm{pt}}$ は系\ref{cor:nfm-point-shift-residual-transfer}の切断である。
最高次線束型の核 $P$ に、さらに
\begin{equation}
\label{eq:nfm-kernel-point-line-type}
 m_P=p_{\mathrm{res}}
\end{equation}
を課す。
この条件は $\phi$ の Jacobian 上の作用だけに依存し、
$p_{\mathrm{res}}$ 自体は $S$ 上で定義される。
従って平面同型の三捩れの選択を忘れても条件は定義できるが、
そのことから元の底上の full 核の同型類が選べるとは結論しない。

支配的な完備族に、前小節の相対構成を適用する。
その係数体を $k$、$R=k[[h]]$ とする。
中心一次変数を $U=0$ とした二次関係は、$g=F|_{U=0}$ だけに依存し、
式\eqref{eq:nfm-active-dvr-algebra}に一致する。
中心変数を含まない PBW 単項式が商の基底となるので、
$D_R=B_h/(U)$ は各次数有限自由で剰余は $k[x,y,z]$ である。
選んだ中心三次の像 $\omega=\Omega_h\bmod(U)$ は中心であり、$\omega\bmod h=g$ である。
従って相対点曲線 $E_R=V(J_g)$ と $L=\mathcal O_{E_R}(1)$ に対し、
命題\ref{prop:nfm-relative-active-thcr}と
命題\ref{prop:nfm-relative-active-window-basechange}が適用でき、
\begin{equation}
\label{eq:nfm-relative-elliptic-window-comparison}
 \operatorname{MF}^{\mathbf Z}(B_h/(U),\Omega_h\bmod U)
 \simeq\operatorname{Perf}(E_R)
\end{equation}
を閉ファイバー・分数体への導来底変更と両立して得る。
この比較は augmentation stabilization を $\mathcal O_{E_R}[1]$ へ送り、
内部 shift を $T_{\mathcal O_{E_R}}\circ\sigma_*\circ(-\otimes L)$ へ送る。
閉ファイバーでは可換側と同じ window1 を用いる。
ここで相対 THCR/window1 は追加仮定ではなく、
PBW 自由性・中心元・滑らかな $g$ から証明した帰結である。

\begin{proposition}[点の平行移動を指定した核の型の非空性]
\label{prop:nfm-kernel-point-line-nonempty}
固定入力Bと良い一ブロックの仮定を置く。
このとき、非空開の有限 \'etale 全射上に、平面同型 $\phi$ と
rotation・指定周期を保つ full 比較核 $P$ を取り、
\[
 \operatorname{Jac}(\phi)=\theta_{\mathrm{sc}}
\]
および
最高次線束条件\eqref{eq:nfm-kernel-active-line-type}と
\eqref{eq:nfm-kernel-point-line-type}を同時に満たすようにできる。
従来の三つの Euler 条件\eqref{eq:nfm-kernel-numerical-type}も同時に課せる。
これは少なくとも一つの支配的な被覆上での非空性であり、
全ての平面同型の選択や全ての最高次線束型の核に対する主張ではない。
\end{proposition}

\begin{proof}
命題\ref{prop:nfm-completed-central-restriction}の相対 CF 比較を
active 商へ制限する。
標的の方程式は $b_hJ_g$ であり、underlying curve は $E_R=V(J_g)$ である。
potential の単元倍を扱うときは、既存の比較と同じ指定周期を保つ。
構成\eqref{eq:nfm-relative-elliptic-window-comparison}と可換側の window1 により、
この比較は同じ $E_R$ 上の $R$-線型同値 $F$ となる。
滑らかな固有曲線の perfect 圏の間のこの比較は perfect な相対核で実現される。
選んだ CF 比較は閉ファイバーで恒等であり、二つの window もそこで一致するので、
その核の閉ファイバーは $\mathcal O_\Delta$ である。
内部 shift の適合性と
補題\ref{lem:nfm-twisted-elliptic-window}から
\eqref{eq:nfm-formal-elliptic-intertwining}が成立する。
従って補題\ref{lem:nfm-formal-elliptic-point-line}により
\[
 F(\mathcal O_{E_R})=M_R,\qquad [M_R]=p_{\mathrm{pt}}.
\]

ここから分数体 $\widehat K=k((h))$ へ移る。
命題\ref{prop:nfm-mf-linear-restriction}の graded purity と
全 Taylor 成分における中心係数商との可換性から
\[
 q_Y\Phi_P(\kappa_{\mathrm{nc}})
 \simeq
 \bigoplus_{j=0}^{n-1}
 \rho^{-j}(M_{\widehat K})[j+1]^{\oplus\binom{n-1}{j}}
\]
を得る。
補題\ref{lem:nfm-active-line-rotation}により、その最高次 cohomology は
$M_{\widehat K}$ である。
形式模型と代数的残余三次との一次空間・方程式の同定を選ぶと、
中心作用の保存により active 平面も同定され、平面同型 $\phi$ を得る。
曲線上への pullback と pushforward は Picard 類と平行移動の対応を保つので、
代数的残余曲線の frame では
$m_P=\operatorname{Jac}(\phi)(p_{\mathrm{pt}})$ となる。

この $\phi$ は $\mathfrak I_{\mathrm{sc}}$ に属することも確認する。
active frame では $\phi$ は二つの $J_g$ の恒等同型である。
命題\ref{prop:nfm-point-scale-formal-comparison}の正確な尺度公式により、
その Jacobian 同型から定まる尺度は
\[
 \tau_{\mathrm{pt}}\circ\operatorname{Jac}(\phi)^*
       =\frac{h}{b_h}\,W_h\otimes P_0,
       \qquad W_h|_{U=0}=b_hJ_g
\]
である。支配的な中心族の記法では、その先頭項は
$h(F+G)\otimes P_0$ となる。
命題\ref{prop:nfm-point-normalized-jacobian-cover}は
Jacobian 同型と尺度被覆の切断を一対一に対応させる。
系\ref{cor:nfm-semiclassical-full-residual}の二つの切断のうち、
この先頭項を持つものは $\theta_{\mathrm{sc}}$ に対応するものだけである。
従って形式体上で
$\operatorname{Jac}(\phi)=\theta_{\mathrm{sc}}$ であり、
$\phi$ は $\mathfrak I_{\mathrm{sc}}$ の点を与える。

あとは命題\ref{prop:nfm-kernel-active-line-nonempty}の有限表示の議論を用いる。
元の関数体を $K$ とし、full 核・逆核・単位・余単位・rotation・周期と
平面同型 $\phi\in\mathfrak I_{\mathrm{sc}}$ を
$K$ 上有限生成な整域 $D$ へ降ろす。
ここでは単に $\mathfrak I$ への点としてではなく、
有限表示の $\mathfrak I_{\mathrm{sc}}$ への点として降ろすので、
Jacobian 同型と半古典符号の一致も保存される。
必要な開への縮小により、最高次 cohomology の可逆層性・次数零・底変更適合性を得る。
相対 Picard 類 $m_P$ と $p_{\mathrm{res}}$ は相対 Jacobian の切断である。
形式体への忠実平坦拡大後に両者が等しいので、
$\operatorname{Frac}(D)$ 上で既に等しい。
相対 Jacobian の分離性により、この等式は両切断が定義された非空開上に成立する。
従ってその開の閉点へ特殊化すれば、有限分離拡大 $K'/K$ 上で
所要の追加型を満たす full 核と $\phi$ を得る。
元の表示 chart の非空開へ広げ、正規化と分岐 locus の除去を行えば、
有限 \'etale 被覆上の主張となる。
三つの perfect mixed Hom を同じ有限データに含めれば、
従来の Euler 条件も保存される。

この特殊化で保持したのは full 核とその最高次線束の Picard 等式である。
形式的な active 比較核そのものや、その coherent な中心線型性を
代数的に降下させたという主張は用いていない。
\end{proof}

\begin{corollary}[選択済みの核の障害を測る点の平行移動]
\label{cor:nfm-kernel-point-picard-obstruction}
split 表示と上記の $\theta=\theta_{\mathrm{sc}}$ を固定し、
\eqref{eq:nfm-kernel-point-line-type}を満たす full 核の同型類が
既に整合的に選ばれているとする。
$C=C_{\mathrm{res}}$ とし、その実現 gerbe の類を $\alpha$、
次数零 Picard 類の表現障害を
$\delta_C(m)=m^*\beta_{\mathrm{Pic},C}$ と書くと、
\begin{equation}
\label{eq:nfm-kernel-point-brauer-class}
 \alpha=\delta_C(p_{\mathrm{res}}).
\end{equation}
従って $p_{\mathrm{res}}$ が実際の次数零線束で表されれば、比較核は降下する。
また定理\ref{thm:nfm-kernel-tripling-cover}の三分点被覆を
$p_{\mathrm{res}}$ に適用すれば、その階数九の有限 \'etale 被覆上で
最高次線束を規格化した full 比較核を得る。
$p_{\mathrm{res}}\in3\operatorname{Jac}(C/S)(S)$ なら
系\ref{cor:nfm-kernel-picard-mod-three}により $\alpha=0$ である。
\end{corollary}

\begin{proof}
系\ref{cor:nfm-kernel-picard-obstruction}により
$\alpha=\delta_C(m_P)$ であるから、
\eqref{eq:nfm-kernel-point-line-type}を代入すればよい。
三分点被覆と規格化については引用した定理と系を用いる。
\end{proof}

\begin{remark}[適用範囲と残る問題]
最高次線束型の全ての核について $m_P=p_{\mathrm{res}}$ を示したわけではない。
また、追加型の非空性と、元の底でその full 核の同型類を整合的に選ぶことは別である。
本節では source の相対 THCR/window1 と必要な二つの体への底変更を構成したが、
表示変更に対する型の内在性、full 核の同型類の整合、
および代数的 active 核の構成はなお別の問題である。
一方、$\theta$ の符号は半古典規格化により既に選ばれており、
本節の非空な型はその $\theta_{\mathrm{sc}}$ に一致する。
既に選ばれた追加型に対する\eqref{eq:nfm-kernel-point-brauer-class}は
Picard 障害を具体的な点へ評価した式であり、代数的 active 核の存在を前提としない。

固定点のない $\sigma$ については $p_{\mathrm{pt}}\ne0$ である。
従ってこの追加型の最高次線束は幾何的に非自明であり、
$m_P=0$ という十分条件だけでその核を降下させる方針は適切でない。
実際の線束による $p_{\mathrm{res}}$ の表現、または三分可能性を調べる必要がある。
ここで点の平行移動を full Ku の rotation 自体と同一視しているのではなく、
構成した active window における rotation 共役を利用している。
\end{remark}

% !TEX root = ../nc-cubic-n-fold.tex
\section{残余比較のための Poisson 補正}
\label{sec:nfm-poisson-correction}

比較障害の後に残る問題は、残余圏が復元する三次と同じ三次を選ぶ
Poisson 対を、一般次元の族として構成することである。
まず全 O-同値に降下する代数的補正を与え、次に固定 formality による
全次数の形式構成と比較する。最後に分離型の族では、固定入力の下で
代数的な補正後三次への全次数の残余比較を示す。
三つの構成の適用範囲を区別し、一般混合係数での全次数の一致を仮定しない。

\subsection{Hessian と trace による代数的な補正}
\label{sec:nfm-algebraic-correction}

有限 jet の障害は、周囲 Poisson tensor を変えることで取り除ける。
ここではその二次 jet を持つ、座標に依存しない代数的な処方を構成する。
この処方の全次数が残余復元と一致するとは主張しない。

$R$ を可換 $\mathbf C$-代数、$V$ を階数 $N=n+2$ の有限局所自由加群、
$W=V^*$ とする。分解可能 Poisson 対を局所的に $(H,P)$ と表し、
$K=\operatorname{supp}(P)\subset W$、$U=\operatorname{ann}(K)\subset V$ と置く。
$H$ の定める三次超曲面と、その $\mathbb P(K)$ への制限がともに滑らかとする。
$S=\operatorname{Sym}U$ および
\[
 T_{H,K}=\operatorname{Sym}V/(\partial_kH:k\in K),\qquad
 \tau_{H,K}=\tfrac18\operatorname{Tr}_{T_{H,K}/S}
\]
と置く。ここでも trace には商への写像を合成する。
active な勾配の最高次部分が三変数の正則列をなすので、
$T_{H,K}$ は $S$ 上階数八の有限局所自由代数である。
この性質と trace は任意の基底変更と両立する。

$K$ 方向の二階微分は
$D_K^2H\in\operatorname{Sym}^2K^*\otimes V$ であり、その行列式は
$(\det K^*)^{\otimes2}\otimes\operatorname{Sym}^3V$ に属する。
従って次の三次式は $K$ の座標と体積形式の選択に依存しない。
\begin{equation}\label{eq:nfm-intrinsic-hessian-correction}
 \mathcal D_P(H)=\langle\det(D_K^2H),P\otimes P\rangle,
 \qquad
 \mathcal C_P(H)=\frac1{24}\mathcal D_P(H)
       +\frac1{12}\tau_{H,K}(\mathcal D_P(H)).
\end{equation}
例えば $P=t\,\partial_x\wedge\partial_y\wedge\partial_z$ なら
$\mathcal D_P(H)=t^2\det\operatorname{Hess}_{x,y,z}(H)$ である。
trace 項は $\operatorname{Sym}^3U$ に属する。

\begin{theorem}[良い Poisson 対の代数的 Hessian 補正]
\label{thm:nfm-algebraic-hessian-correction}
active な平面三次も滑らかな $\mathfrak P_n^{\mathrm{dec}}$ の開を
$\mathfrak P_n^{\mathrm{ell}}$ と書く。
次の処方は、出力の三次と active な平面三次がともに滑らかな非空開上で、
代数的な射
\[
 \Phi:\mathfrak P_n^{\mathrm{ell}}\dashrightarrow
             \mathfrak P_n^{\mathrm{ell}},\qquad
 (H,P)\longmapsto(H^\sharp,P),\quad H^\sharp=H+\mathcal C_P(H)
\]
を定める。これは任意の基底変更、方程式の生成元変更、
全ての射影座標同型と両立する。
補正後の強い lift は厳密に
\begin{equation}\label{eq:nfm-corrected-algebraic-tensor}
 \Pi^\sharp=\iota_{dH^\sharp}P
 =\Pi+\frac1{24}\iota_{d\mathcal D_P(H)}P
\end{equation}
であり、同じ $K$ と $U$ を持つ Poisson tensor である。

この有理射は支配的で、非空開への制限は \emph{\'etale} である。
さらに定理\ref{thm:nfm-correction-degree-nine}で、始域と終域を適切に縮めると
有限 \emph{\'etale} で階数 $9$ となることを示す。
従って定理\ref{thm:poisson-good-open-isomorphism}を介し、
$\mathfrak M_n^{\mathrm{sm}}$ の非空開から Poisson 対のモジュライへの射を与える。
必要なら定義域をさらに縮め、出力も $\mathfrak P_n^{\mathrm{good}}$ に入るようにできる。
ただしこのことは、元の非可換圏とその出力の Kuznetsov 成分の
全次数の同値を意味しない。
\end{theorem}

\begin{proof}
三次方程式を $H\mapsto cH$ とすると、強い lift の表示は
$P\mapsto c^{-1}P$ となる。
Hessian 行列式は $c^3$ 倍、$P^{\otimes2}$ は $c^{-2}$ 倍になるから
\[
 \mathcal D_{c^{-1}P}(cH)=c\mathcal D_P(H).
\]
active な勾配イデアルは変わらず、trace も変わらない。
従って $\mathcal C_{c^{-1}P}(cH)=c\mathcal C_P(H)$ であり、
出力 $H^\sharp\otimes P$ は方程式の生成元に依存しない。
Hessian の縮約と有限局所自由代数の trace は線型座標変換に共変である。
特に active 変数に中心一次式を加える frame 変更も含み、
$V\twoheadrightarrow K^*$ の分裂を選ぶ必要はない。
次数スカラーの重みも相殺するので、射影群の商 stack へ降下する。
有限局所自由性と trace の基底変更により、非平坦・非被約の場合も含めて
任意の基底変更と両立する。

局所的に $K$ の座標を選ぶと、$\Pi^\sharp$ は三変数の Jacobian 括弧で、
残りの変数は中心である。従って Jacobi 恒等式は多項式恒等式として成立し、
$H^\sharp$ はその Casimir である。
trace 項の微分は $K$ 方向で零なので、式
\eqref{eq:nfm-corrected-algebraic-tensor}を得る。
出力が滑らかなら、系\ref{cor:poisson-support-recovery}により一次中心は
ちょうど $U$ である。

定義域の非空性と支配性を一様に確認する。
滑らかな三次と active な平面三次の開
\[
 B\subset\mathbb P(\operatorname{Sym}^3V)\times\operatorname{Gr}(3,W)
\]
上で、$\mathcal L_H\otimes\det\mathcal K$ の全空間 $\overline Z$ を取り、
零切断も許す。ただし零切断では $K$ を忘れない。
式\eqref{eq:nfm-intrinsic-hessian-correction}は $P=0$ でも定義でき、
補正は線束方向の二次以上である。
したがって $\Phi$ は零切断を含む開で定義され、零切断を各点ごとに固定し、
そこでの微分は恒等である。
$\overline Z$ は滑らかなので、$\Phi$ は零切断の近傍で \emph{\'etale} となる。
この近傍の零切断を除いた部分は非空で、像は開かつ稠密である。
よって元の非零部分上の有理射は支配的であり、非空開で \emph{\'etale} である。
同じ群による同変商に移すと、この性質は stack 上にも降りる。

$\mathfrak P_n^{\mathrm{good}}$ は既約な
$\mathfrak P_n^{\mathrm{ell}}$ の非空開である。
支配性により、その逆像と定義域中の
$\mathfrak P_n^{\mathrm{good}}$ の交わりも非空開となる。
そこに既存の O-marked moduli との同型を適用すればよい。
\end{proof}

\begin{corollary}[flag を残した零切断の形式近傍]
\label{cor:nfm-correction-formal-automorphism}
上の $\overline Z$ の零切断に沿う形式完備化上で、$\Phi$ は
恒等を剰余に持つ形式自己同型である。
同じことは射影群の商にも成立する。
\end{corollary}
\begin{proof}
零切断を固定する \emph{\'etale} 射に対し、零切断上の恒等写像の
持上げは各冪零厚みで一意である。
これらを順次持ち上げると、任意の有限厚み上で相互に逆な写像を得る。
その逆極限が形式逆写像を与える。一意性により群作用とも両立する。
これは明示的な補正 $\Phi$ 自身の形式可逆性であり、
formality の中心写像との全次数の一致ではない。
\end{proof}

\begin{proposition}[楕円点曲線との厳密な整合]
\label{prop:nfm-correction-point-cubic}
良い一ブロック型の ambient 二次代数 $B_\Pi$ を考える。
$g=H|_{\mathbb P(K)}$ とすると、補正後の制限は
\[
 H^\sharp|_{\mathbb P(K)}
 =g+\frac1{24}\mathcal D_P(g)
\]
であり、これは $B_\Pi$ の active な楕円点曲線の方程式である。
従って補正射は、回復した support 平面の中で、
量子側の楕円点曲線と同じ平面三次を選ぶ。
この主張は点曲線の shift を保持した対応全体や、
定常点成分の同定を含まない。
\end{proposition}

三生成元の二次 Calabi--Yau 代数の点曲線を Hessian によって記述することは、
Mori--Smith\cite[Theorems~1.2, 7.2]{MoriSmith}の既知の結果である。
以下の証明は本稿の係数規約を確認するもので、その点曲線公式自体の新規性を主張しない。
本節ではこれを混合係数の trace と組み合わせ、
全次元の Poisson 対および O-marked moduli 上の代数的補正に用いる。

\begin{proof}
trace 項は純中心三次なので、$\mathbb P(K)$ への制限で零になる。
局所座標で $P=t\,\partial_x\wedge\partial_y\wedge\partial_z$ と書き、
$p=(x,y,z)^{\mathsf t}$ とする。二次関係の双線型化は
\[
 \left(C(p)-\frac t2\operatorname{Hess}(g)(p)\right)q=0,
 \qquad C(p)q=p\times q
\]
である。対称行列 $M$ に対する三次恒等式
\[
 \det(C(p)+sM)=s\,p^{\mathsf t}Mp+s^3\det M
\]
と $p^{\mathsf t}\operatorname{Hess}(g)(p)p=6g(p)$ より、
その行列式は
\[
 -3t\left(g+\frac{t^2}{24}\det\operatorname{Hess}(g)\right)
\]
である。非零の $P$ を持つ相対チャートでは $t$ は単元に取れる。
行列式恒等式は、$\operatorname{adj}C(p)=pp^{\mathsf t}$ と
転置による $s$ の偶数次係数の消滅からも直接分かる。
この三次が滑らかな開では、係数行列の階数は曲線上の全ての幾何点で二である。
階数が一以下なら全ての余因子が消え、行列式の全偏微分も消えるので、
滑らかさに反する。さらに $Mq=0$ に $\operatorname{adj}(M)$ を掛けると
$\det(M)q=0$ となり、射影座標の局所的な可逆成分から incidence は
行列式曲線をスキームとして経由する。
可逆な二次小行列を使えば、任意基底上で核は直線部分束となる。
従って核は任意基底変更と両立する直線部分束であり、
双線型関係の incidence はこの曲線からの射のグラフとなる。
左右の変数を交換すると $t$ の符号が変わるだけなので、
右射影の三次行列式も同じ曲線を与え、その核も直線部分束である。
従って両射影は同型で、incidence は曲線の自己同型のグラフである。
これにより三次が楕円点曲線であることも任意基底上で確認できる。
$t=0$ で得られる三次はこの族の極限であり、
可換多項式環の点スキームそのものと同一視しない。
\end{proof}

\begin{corollary}[残余復元の二次 jet を実現する代数的処方]
\label{cor:nfm-corrected-residual-jet}
命題\ref{prop:nfm-relative-residual-second-jet}の規約と入力を使う。
Hessian 正規化中心元の trace--Hessian 出力に $\Phi$ を適用すると、
\[
 H^\sharp
 =F+\frac{h^2}{24}\mathscr H_F+h^4\tau_F(\Psi_\lambda(F)),
 \qquad H^\sharp\equiv W^{\mathrm{res}}\pmod{h^4}.
\]
その強い lift は
\[
 \Pi^\sharp=h\Pi_F+\frac{h^3}{24}\Pi_{\mathscr H_F}
\]
である。従って前節の二次 Jacobi 障害を、この代数的な補正により除去できる。
\end{corollary}
\begin{proof}
trace--Hessian 出力と $F$ の差は純中心三次なので、
active な勾配、Hessian、および $\tau_F$ は変わらない。
その強い lift の表示は $P=hP_0$ である。
従って $\mathcal D_{hP_0}(H^{\mathrm{TH}})=h^2\mathscr H_F$ となり、
二つの trace 項が相殺して表示した式を得る。
tensor の式は純中心項の縮約が零であることから従う。
残余二次 jet と以下の系\ref{cor:nfm-mixed-even-jet}の偶奇性により、
合同の剰余は $h^4$ まで改善する。
この計算は $h\ne0$ の対に対する代数的処方を当該 trait に沿って延長したもので、
可換特殊点全体の O-marked moduli に非零 Poisson tensor を付与するものではない。
\end{proof}

\begin{remark}[残る比較問題]
$\Phi$ は特定の formality の全次数の逆写像として定義したものではない。
ここで示したのは、任意基底の代数的構成、全 O-同値への降下、
楕円点曲線との一致、および指定中心元に対する残余三次の二次 jet との一致である。
実際の formal Poisson tensor は scalar formal unit による規格化も持ち得る。
従ってその tensor との比較には、potential と support の比較に加えて尺度を追跡する。
全次数での $\operatorname{Ku}(V(H^\sharp))$ との同値、
またはこの処方の一意性は、二次 jet の一致からは従わない。
\end{remark}

\subsection{一般混合係数における補正射の次数}
\label{sec:nfm-correction-degree}

Hessian 補正の零切断に沿う形式逆写像は、大域の有理逆写像ではない。
この相違は分離型の族に限らず、任意次元の一般混合係数で定量化できる。
以下の次数は、同じ射影座標群による商 stack の間の表現可能射の次数である。
粗モジュライの関数体の次数を、慣性群の比較なしに言い換えてはいない。
Hesse pencil と同変性による還元は、射影的な純 Hessian 写像の次数三を扱う
Ciliberto--Ottaviani\cite[Section~5, Theorem~5.7]{NFMCilibertoOttaviani}にも現れる。
本稿で数えるのは $f\mapsto f+\det\operatorname{Hess}(f)/24$ とその混合拡張であり、
純 Hessian 写像の二回合成の次数九とも異なる。

\begin{theorem}[補正射の一般次数]
\label{thm:nfm-correction-degree-nine}
任意の $n\geq3$ に対し、代数的 Hessian--trace 補正
\[
 \Phi:\mathfrak P_n^{\mathrm{ell}}\dashrightarrow
       \mathfrak P_n^{\mathrm{ell}}
\]
の一般次数は $9$ である。
より正確には、始域と終域に非空開部分 stack を取ると、$\Phi$ は
表現可能な有限 \emph{\'etale} 射で、階数は $9$ となる。
両方の開を $\mathfrak P_n^{\mathrm{good}}$ に含めてよい。
従って既存の $\mathfrak M_n^{\mathrm{sm}}\simeq
\mathfrak P_n^{\mathrm{good}}$ を介しても、同じ次数の射を得る。
特に、この補正射自身は有理同型ではない。
\end{theorem}

\begin{proof}
まず support $K$ を固定し、その体積三重ベクトル $P_0$ を選ぶ。
局所表示 $P=sP_0$ に対し、方程式の生成元変更によって
$(H,sP_0)$ を $(F,P_0)$、$F=sH$ と同一視する。
この表示における補正は
\begin{equation}\label{eq:nfm-normalized-correction-map}
 F\longmapsto F+\frac1{24}\Delta_F
                  +\frac1{12}\tau_F(\Delta_F),
 \qquad \Delta_F=\det\operatorname{Hess}_K(F).
\end{equation}
trace は $F$ の純中心部分に依存しない。

\smallskip\noindent
\textbf{active な三次の次数。}
三次元の体積座標 $(x,y,z)$ における多項式写像
\[
 \phi_3:\operatorname{Sym}^3\mathbf C^3{}^*
   \longrightarrow\operatorname{Sym}^3\mathbf C^3{}^*,
 \qquad f\longmapsto f+\tfrac1{24}\det\operatorname{Hess}(f)
\]
を考える。$g_{p,q}=p(x^3+y^3+z^3)/3+qxyz$ とおくと
\[
 \phi_3(g_{p,q})=g_{\psi_1(p,q),\psi_2(p,q)},\qquad
 \psi(p,q)=\left(p(1-q^2/4),\ q+p^3/3+q^3/12\right).
\]
この $\psi$ は、射影平面の射
\begin{equation}\label{eq:nfm-projective-correction-chart}
 [p:q:t]\longmapsto
 [3p(4t^2-q^2):12qt^2+4p^3+q^3:12t^3]
\end{equation}
へ拡張する。$t\ne0$ では最後の成分が零でなく、$t=0$ では
$pq^2=4p^3+q^3=0$ の非零解はないので基点はない。
この射は $\mathcal O(1)$ を $\mathcal O(3)$ に引き戻す。
従って有限かつ次数 $3^2=9$ である。
無限遠直線の逆像は同じ無限遠直線なので、$\psi$ 自身も有限であり、
滑らかな二次元の始域と終域の間の有限射として平坦、階数 $9$ となる。

ここで
\[
 \vartheta:\operatorname{SL}_3\times\mathbb A^2
    \longrightarrow\operatorname{Sym}^3\mathbf C^3{}^*,
 \qquad (A,p,q)\longmapsto A\cdot g_{p,q}
\]
は支配的かつ一般に有限である。
実際、$(1,1,0)$ における微分は、$\mathfrak{sl}_3$ の非対角成分から
六つの $x_i^2x_j$、対角成分から二つの $x_i^3$ の差、$p,q$ から
$x^3+y^3+z^3$ と $xyz$ を与える。これらは十次元の三次式空間の基底である。
Hessian は $\operatorname{SL}_3$ 同変なので
\[
 \vartheta\circ(1\times\psi)=\phi_3\circ\vartheta.
\]
一般次数の乗法性により
\[
 (\deg\vartheta)\,9=(\deg\phi_3)(\deg\vartheta),
 \qquad \deg\phi_3=9.
\]
この論証は Hesse frame の有限同定の個数を仮定していない。

\smallskip\noindent
\textbf{混合係数の復元。}
補空間を一時的に選び、active 次数で
\[
 F=f_3+f_2+f_1+f_0,\qquad
 f_i\in\operatorname{Sym}^iK^*\otimes
                    \operatorname{Sym}^{3-i}U
\]
と書く。$\operatorname{Hess}_K(F)$ は $f_3+f_2$ だけに依存する。
従って $\Delta_F=\delta_3+\delta_2+\delta_1+\delta_0$ と分けると
\[
 \delta_3=\det\operatorname{Hess}(f_3),\qquad
 \delta_2=\operatorname{tr}\bigl(
   \operatorname{adj}(\operatorname{Hess}(f_3))
                       \operatorname{Hess}_K(f_2)\bigr),
\]
また $\delta_1,\delta_0$ は $f_3,f_2$ だけから決まる。
式\eqref{eq:nfm-normalized-correction-map}は次の三角形の形になる。
\begin{align*}
 f'_3&=\phi_3(f_3),\\
 f'_2&=(L_{f_3}\otimes1_U)f_2,\\
 f'_1&=f_1+\delta_1/24,\\
 f'_0&=f_0+\delta_0/24+\tau_F(\Delta_F)/12,
\end{align*}
ここで
\[
 L_f(b)=b+\tfrac1{24}\operatorname{tr}\bigl(
   \operatorname{adj}(\operatorname{Hess}(f))\operatorname{Hess}(b)\bigr)
 \quad (b\in\operatorname{Sym}^2K^*).
\]
$L_f$ は一般の $f$ で可逆である。
例えば $f=g_{p,q}$ の場合、順序 $(x^2,yz)$、$(y^2,zx)$、$(z^2,xy)$
では同じ二次正方行列の三つの直和になり、各行列は
\[
 \begin{pmatrix}
  1-q^2/12&-pq/6\\ p^2/3&1+q^2/12
 \end{pmatrix}.
 \qquad
 \det L_{g_{p,q}}
   =\left(1+\frac{8p^3q-q^4}{144}\right)^3.
\]
特に $g_{1,0}$ で行列式は $1$ である。
一般の $f'_3$ には九つの $f_3$ があり、それらすべてで $L_{f_3}$ を
可逆にするよう終域の開を縮めてよい。
各 $f_3$ を選ぶと、上式から $f_2$、次に $f_1$、最後に $f_0$ が一意に決まる。
最後の trace は $f_0$ に依存しないため、そこで新たな方程式は発生しない。
よって任意の $\dim U=n-1$ に対して一般次数は同じく $9$ である。

\smallskip\noindent
\textbf{開の選択と商。}
ここまでの計算は有限型既約な係数空間の一般点の計算である。
一般に有限な写像では、始域の真の閉集合は終域を支配できない。
従って full cubic の特異点条件や good 開の補集合を除いても、
終域を縮めれば九つの一般分枝はすべて残る。
active 三次の滑らかさにより、trace を定める次数八の代数は有限局所自由で、
trace は正則である。全三次の滑らかさ、補正後の滑らかさ、
および既存の good 条件もそれぞれ非空開である。
一般に有限な写像では、取り除いた真の閉集合の像の閉包は終域を覆わない。
したがって有限性の開を取り、分岐像およびこれらの除外集合の像を取り除くと、
有限 \emph{\'etale} 階数 $9$ の射が得られる。
階数は一般の混合係数から求めており、特殊な分離型部分空間の次数を
そのまま全空間へ移したのではない。

表示 $(H,P)$ の方程式スカラーを商にした係数空間は、
$\operatorname{Gr}(3,W)$ 上の
$\operatorname{Sym}^3V\otimes\det\mathcal K$ の非零部分の開と見られる。
実際、その点は $H\otimes P$ である。
$\operatorname{PGL}(W)$ は support $K$ に推移的に作用し、
$K$ を固定して $P_0$ を選ぶと上の $F\otimes P_0$ 表示を得る。
$K$ の安定化群の作用は、体積の変化を $F$ の倍率へ吸収した作用であり、
始域と終域で同じである。従って support を固定した計算の次数 $9$ は、
support と体積 frame を動かした係数空間でも変わらない。
構成は同じ射影座標群 $G$ に関して同変であり、
有限 \emph{\'etale} 性を持つ終域の開を $G$ 不変に選べる。
一般に $G$ 同変な $X\to Y$ に対して
\[
 [X/G]\mathbin{\times}_{[Y/G]}Y\simeq X
\]
である。従って $[X/G]\to[Y/G]$ は表現可能であり、
有限 \emph{\'etale} 性と階数 $9$ はこの忠実平坦被覆から降下する。
以上により主張を得る。
\end{proof}

\begin{remark}[Hesse pencil と楕円曲線の乗法写像]
\label{rem:nfm-degree-not-multiplication}
active な入力と出力は同じ Hesse pencil に属するが、一般には同じ楕円曲線ではない。
例えば $q=0$、$p\ne0$ の入力曲線は Fermat 曲線である一方、出力の
Hesse 係数は $q'/p'=p^2/3$ であり、その $j$ 不変量は非定数である。
従って上の次数 $9$ を、入力の楕円曲線上の $[3]$ と同一視できない。
Hesse pencil とその $j$ 関数については
\cite[Section~2, Lemma~2.1, equations~(4)--(5)]{NFMArtebaniDolgachev}を参照。
Sklyanin 中心三次を別の superpotential と見たときに現れる
$(E,p)\mapsto(E,[-2]p)$\cite[Section~2]{NFMDLL}とも異なる。
後者は同じ $E$ を保つ。
今回の次数定理は、形式逆量子化の代数化や残余圏の全次数比較を主張していない。
\end{remark}

\subsection{残余復元と両立する形式 Poisson 対}
\label{sec:nfm-corrected-formal-poisson}

前節の障害は、量子代数の二次関係から直読する tensor と、
formality の逆対応が与える tensor を区別することで解消される。
本節では harmonic-angle formality を一つ固定し、形式パラメータ $h$ と
次数一生成元を保つ模型を用いる。
次数付き二次代数と二次 Poisson tensor の形式対応そのものは
Kontsevich\cite[Section~2.3, Proposition~6]{NFMKontsevichAlg}の既知結果である。
ここで明らかにするのは、一ブロックの場合の中心一次空間の保存、
残余三次による全次数の tensor の決定、および前節の障害を打ち消す補正である。
構成は形式的なものであり、非零の複素数 $h$ への収束・代数化とは区別する。

\subsubsection{次数一を固定する模型と有限係数の再帰}

$R$ を複素代数、$V_R$ を階数 $N$ の有限自由加群とし、
$S_R=\operatorname{Sym}_R V_R$ と置く。
次数一生成元を $z_1,\ldots,z_N$ とする。
次数別に $h$ 進完備な二次代数
\[
 B_h=T_R(V_R)[[h]]\Big/
 \left(z_i z_j-z_j z_i-h\operatorname{sym}(q_{ij})\right),
 \qquad q_{ij}\in\operatorname{Sym}^2_R V_R[[h]],\quad q_{ji}=-q_{ij},
\]
を考える。完備化は各内部次数で行い、最後にその直和を取る。
PBW 条件とは、この各内部次数が $S_R[[h]]$ と同じ有限局所自由加群になることをいう。
以下で存在を用いるのは、固定入力論文Bの
補題\texttt{cft:lem:model}が適用できる完備模型、
または後述の被約有限型の普遍係数模型からの基底変更である。

\begin{lemma}[生成元を固定する正規化の一意性]
\label{lem:nfm-normalized-star-model}
固定した affine-equivariant な Kontsevich 星積\cite[Section~6.3]{NFMKontsevich}について、
$h$ を固定し、全ての次数一生成元を固定する厳密同型
\[
 \phi:B_h\xrightarrow{\sim}(S_R[[h]],*_{\pi_h})
\]
が存在するなら、二次係数の formal Poisson tensor $\pi_h$ と $\phi$ は一意である。
もとの formality 同型が次数を保ち剰余で恒等なら、
形式的な線型座標変換によりこの正規化を行える。
従って正規化された対応は、係数に依存する可逆な線型座標変換にも共変である。
\end{lemma}

\begin{proof}
もとの同型の次数一部分は、剰余で恒等な可逆行列である。
その逆による線型座標変換を標的へ合成する。
グラフ写像の線型共変性により、標的は変換した tensor の同じ Kontsevich 星積であり、
全生成元が固定される。

二つの正規化された tensor が初めて $h^m$ で異なり、
$\pi'_h-\pi_h=h^mD\pmod{h^{m+1}}$ とする。
二頂点以上のグラフにこの差を入れると、他の tensor 入力が $h$ の倍数なので
$h^{m+1}$ の倍数になる。従って星積の差は
\[
 f*_{\pi'_h}g-f*_{\pi_h}g
 =\frac{h^m}{2}D(df,dg)\pmod{h^{m+1}}.
\]
同じ二次関係を両標的で評価すると、交換子の差は
$h^mD(dz_i,dz_j)$ である。
関係式の右辺には既に $h$ が掛かっているため、その差は $h^{m+1}$ に入る。
よって全ての $i,j$ について $D(dz_i,dz_j)=0$ であり、$D=0$ となる。
これは任意の最初の相違を排除する。
$B_h$ は次数一で生成されるので、tensor が等しければ $\phi$ も等しい。
座標変換した模型も同じ生成元固定条件を満たすため、最後の共変性は一意性から従う。
\end{proof}

この議論は構成的である。$\pi_{<m}=\sum_{j<m}h^j\pi_j$ が得られたとき、
\begin{equation}
\label{eq:nfm-normalized-tensor-recursion}
 \pi_m(dz_i,dz_j)=
 -[h^m]\left(
 z_i*_{\pi_{<m}}z_j-z_j*_{\pi_{<m}}z_i
 -h\operatorname{sym}_{*_{\pi_{<m}}}(q_{ij})\right)
\end{equation}
で次の係数が定まる。右辺の星積は、この段階ではグラフによる双微分式として評価する。
最終的な Maurer--Cartan 恒等式は、formality による存在と一意性から確認する。
各固定次数で使うグラフは有限個であり、右辺は二次多項式である。
従って $\pi_m$ は、$q_{ij}$ の有限個の係数と固定したグラフ重みの有限式である。

\begin{proposition}[普遍係数上での正則性と基底変更]
\label{prop:nfm-normalized-model-basechange}
$T$ を被約な有限型複素スキームとし、その上に上記の次数別 PBW 二次代数族と
形式中心三次 $\Omega_h$ があるとする。
各複素閉点で固定入力の星積模型が存在すると仮定する。
このとき正規化された $\pi_h,\phi$ および
\[
 W_h=C_{\pi_h}^{-1}\phi(\Omega_h)
\]
は、各 $h$ 係数が $T$ 上正則な形式族として一意に構成される。
その形成は $T$ からの任意の基底変更と両立する。
基底変更先は非被約でも非平坦でもよい。
\end{proposition}

\begin{proof}
局所枠で\eqref{eq:nfm-normalized-tensor-recursion}を使う。
各係数は有限式なので正則であり、前補題の共変性により局所枠の間で貼り合わさる。
各閉点で再帰は formality の正規化された解と一致する。
従って $[\pi_h,\pi_h]=0$ の各係数は、被約有限型複素スキームの全閉点で零となり、
$T$ 上で零である。これで得られた tensor の星積は結合的である。
再帰により二次関係が成立するので、生成元を固定する代数準同型 $\phi$ が存在する。
各内部次数でこれは剰余が恒等な有限自由加群間の行列であり、
$h$ 進幾何級数によって逆を持つ。

$C_{\pi_h}$ は次数三の有限局所自由加群上で $1+O(h)$ である。
従ってその逆も係数wise正則であり、$W_h$ が正則に得られる。
中心性と formality の接写像の chain 恒等式から
$[\pi_h,W_h]=0$ となる。これは固定入力の補題\texttt{cft:lem:center}と同じ
$h$ 進の HKR 単射性の議論である。
以上は全て有限係数の恒等式と $h$ 進の行列逆であるため、
任意の基底変更後も成立する。
ここで非被約基底について用いたのは、既に構成した恒等式の基底変更であり、
その閉点だけを調べて恒等式を証明したのではない。
\end{proof}

\subsubsection{中心一次空間と一ブロック support の保存}

$U_R=\langle u_1,\ldots,u_r\rangle\subset V_R$ を中心一次空間とし、
$V_R=\langle x,y,z\rangle\oplus U_R$、$P=\partial_x\wedge\partial_y\wedge\partial_z$ とする。
この分解は証明用の局所座標であり、結論はその選択に依存しない。

\begin{lemma}[線型中心の厳密な保存]
\label{lem:nfm-formal-linear-center}
上の生成元固定模型で $u\in U_R$ が $B_h$ の中心に属するなら、
$\iota_{du}\pi_h=0$ である。
さらに任意の $f\in\operatorname{Sym}U_R[[h]]$ について
\[
 C_{\pi_h}(f)=f,
 \qquad f*_{\pi_h}a=fa=a*_{\pi_h}f
\]
が成立する。従って補正後の tensor も、元の中心変数方向を微分しない。
\end{lemma}

\begin{proof}
$h$ の次数について帰納法を使う。
$h^m$ 未満の全 tensor 係数が $U_R$ を消すと仮定する。
これらの係数のみを使う正の次数のグラフで、境界入力に線型関数 $u$ を置く。
$u$ へ辺が入れば、その辺の添字は $U_R$ 方向であり、
出発頂点の tensor のその成分は零である。
辺が一本も入らない場合のグラフの総寄与は、単位の正規化によって零である。
従って $u* a-a*u$ の $h^m$ 係数に寄与するのは
HKR 項 $\pi_m(du,da)$ だけである。
$\phi$ は $u$ を固定し、$u$ は中心元なのでこの交換子は零である。
$a=z_i$ として全成分の消滅を得る。帰納法により $\pi_h$ 全体が相対的となる。

これ以降、$f(U)$ への全 active 微分は零である。
星積の正次数グラフも同じ理由で零となる。
中心写像のグラフで関数頂点 $f(U)$ に辺が入れば零であり、
一本も入らなければその関数頂点は孤立する。
孤立頂点を忘れると積分形式は必要な次元より低い配置空間からの引戻しとなるため、
正次数の重みは零である。これは $C_{\pi_h}(1)=1$ のグラフによる証明でもある。
よって $C_{\pi_h}(f)=f$ を得る。
\end{proof}

\begin{lemma}[三次 Casimir による相対二次 tensor の決定]
\label{lem:nfm-relative-casimir-tensor}
$W_h\in R[[h]][x,y,z,u_1,\ldots,u_r]_3$ とし、
その active 偏微分 $W_x,W_y,W_z$ の Koszul 複体が普遍的に完全であるとする。
二次係数の相対 bivector $Q$ が $\iota_{dW_h}Q=0$ を満たすなら、
一意な $a\in R[[h]]$ によって
\[
 Q=a\,\iota_{dW_h}P
\]
と書ける。$a$ は座標変数に依存しない。
\end{lemma}

\begin{proof}
$Q=Q^{yz}\partial_y\wedge\partial_z+
Q^{zx}\partial_z\wedge\partial_x+Q^{xy}\partial_x\wedge\partial_y$ とし、
$\alpha=Q^{yz}dx+Q^{zx}dy+Q^{xy}dz$ と置く。
Casimir 条件は $d_KW_h\wedge\alpha=0$ である。
三つの偏微分による wedge Koszul 複体の一次完全性から
$\alpha=a\,d_KW_h$ を得る。
$\alpha$ と $d_KW_h$ の係数次数はいずれも二なので、
次数零成分を取れば $a$ は $R[[h]]$ に属する。
次数零の Koszul 単射性により一意である。
全ての座標を内部次数一と数えているため、$a$ に中心変数の依存も残らない。
\end{proof}

active 制限 $g=F|_{U=0}$ が滑らかな平面三次を定める場合、
必要な Koszul 完全性は、active 勾配による八次有限局所自由代数とその次数別分裂完全性から従う。
$h$ 進変形については各固定内部次数で持上げ、完備性により極限を取る。
この議論は、その普遍係数模型からの任意の基底変更と両立する。

\begin{theorem}[残余三次に適合する形式 Poisson 対]
\label{thm:nfm-residual-formal-poisson}
良い一ブロックの PBW 中心三次族に対し、上の固定 formality による
生成元固定模型を取る。量子関係の先頭を
$[a,b]=h\operatorname{sym}\Pi_F(a,b)$、$\Pi_F=\iota_{dF}P$ とし、
$\Omega_h\bmod h=F$ とする。
$V(F)$ と active な平面三次がともに基底上滑らかなら、一意な形式スカラー $a_h=h+O(h^2)$ が存在して
\begin{equation}
\label{eq:nfm-residual-formal-factor}
 \pi_h=a_h\,\iota_{dW_h}P,
 \qquad W_h=C_{\pi_h}^{-1}\phi(\Omega_h)
\end{equation}
となる。従って $\pi_h/h$ の線は $W_h$ と元の rank 三 support によって決まり、
$h$ を逆にした後の $(V(W_h),\pi_h|_{V(W_h)})$ は非零の一ブロック Poisson 対である。

固定入力の完備 MF 移送と相対 window 比較が適用できる場合には、同じ $W_h$ に対し
\[
 \mathcal K_{B_h/(\Omega_h)}\simeq\operatorname{Ku}(V(W_h))
\]
を得る。従ってこの形式 Poisson 対では、残余比較の対象となる三次超曲面が一致する。
この圏同値は入力の結論を適用したものであり、本節で完備 MF 定理の全証明を
独立に再検証したという意味ではない。
\end{theorem}

\begin{proof}
中心逆像の構成から $[\pi_h,W_h]=0$ であり、前補題から $\pi_h$ は相対 bivector である。
Koszul 補題を適用して\eqref{eq:nfm-residual-formal-factor}を得る。
$\pi_h\equiv h\Pi_F\pmod{h^2}$ と $W_h\bmod h=F$ を比較すると
$a_h\equiv h\pmod{h^2}$ である。
従って $a_h/h$ は単元であり、$h$ の逆元を許すと tensor は非零となる。
Jacobi 恒等式は formality 模型の性質から従うが、
相対三変数の Jacobian tensor の表示からも直接従う。
$W_h$ は滑らかな $F$ の形式変形なので、対応する各 Artin 切断とその形式的極限は滑らかな射影三次の変形である。
最後は $W_h$ の定義が固定入力の完備 MF 移送で用いる中心逆像そのものであることによる。
\end{proof}

$\pi_h/h$ の次数一 Poisson 中心は厳密に $U_R[[h]]$ である。
実際、余分な active 線型中心があれば、その定数 active 勾配が
$d_KW_h$ と wedge 零となる。
同じ Koszul 完全性は、それが係数次数 $-2$ の倍数であることを要求するため零である。
従って元の中心一次空間と rank 三 support は、形式変形の途中でも保存される。

\begin{corollary}[形式中心と Casimir 空間の線型同定]
\label{cor:nfm-formal-center-casimir}
上の仮定の下で、補正された $\pi_h$ は $B_h$ だけから定まり、
選んだ中心三次には依存しない。中心逆像は線型同型
\[
 C_{\pi_h}^{-1}\phi:
 Z(B_h)_3\xrightarrow{\sim}
 \operatorname{Cas}(\pi_h)_3
 =R[[h]]W_h\oplus\operatorname{Sym}^3 U_R[[h]]
\]
を与え、純中心三次の部分加群上で恒等である。
特に $\Omega_h$ を $b\Omega_h+G(U)$ に替えると、
$W_h$ は $bW_h+G(U)$ に変わり、$\pi_h$ は変わらない。
ここでは $h$ 進形式中心を扱っており、$h=0$ での通常の中心や、
有限 Artin 切断で新たに現れる $h$-torsion の中心元と同一視しない。
\end{corollary}

\begin{proof}
$\pi_h,\phi$ の構成は二次関係だけを使う。
formality の接写像の chain 恒等式と、その vector 場上での最低次の単射性により、
$C_{\pi_h}$ は Casimir と星積中心の間の双射である。
次数を保つので次数三に制限でき、純中心三次上の恒等性は
補題\ref{lem:nfm-formal-linear-center}による。
$h$ は $R[[h]]$ 上非零因子、$a_h/h$ は単元なので、
三次 $H$ が Casimir である条件は
$d_KH\wedge d_KW_h=0$ である。
Koszul 完全性と次数により $d_KH=c\,d_KW_h$、$c\in R[[h]]$ を得る。
標数零であるため $H-cW_h$ は active 変数に依存せず、純中心三次になる。
和が直和であることも active 勾配の Koszul 単射性から従う。
\end{proof}

\subsubsection{三次 tensor のスカラー消滅}

命題\ref{prop:nfm-relative-residual-second-jet}と同じ中心元の正規化を使い、
$F$ の係数を $h$ に関して定数とする。
$\mathscr H_F=\det\operatorname{Hess}_{x,y,z}(F)$ と書く。
二次の星積の対称性と既知の残余二次 jet を
\eqref{eq:nfm-residual-formal-factor}へ代入すると、
$a_h=h(1+c_2h^2+O(h^3))$ かつ
\begin{equation}
\label{eq:nfm-corrected-poisson-third}
 \pi_h=h\Pi_F+h^3\left(\frac1{24}\Pi_{\mathscr H_F}
                  +c_2\Pi_F\right)+O(h^4)
\end{equation}
となる。Casimir 条件に加えて普遍係数の多項式性と不変式を使うことで、
この $c_2$ を決定できる。

\begin{lemma}[三変数三次形式の二次不変式]
\label{lem:nfm-no-quadratic-cubic-invariant}
三変数三次形式の係数に関する斉次二次の
$\operatorname{SL}_3$ 不変多項式は零である。
\end{lemma}

\begin{proof}
$g=\sum_{i+j+k=3}a_{ijk}x^iy^jz^k$ と書く。
行列式一の対角行列に対する重みを比較すると、不変な斉次二次式は
\[
 Q=Aa_{210}a_{012}+Ba_{201}a_{021}
   +Ca_{120}a_{102}+Da_{111}^2
\]
という形に限られる。実際、二つの係数の添字の和は
$(2,2,2)$ でなければならない。
せん断 $x\mapsto x+ty$ と $x\mapsto x+tz$ に対する
$t=0$ での微分を、それぞれ $L_y,L_z$ とする。
\[
\begin{aligned}
 L_yQ={}&3Aa_{300}a_{012}
       +(A+2C)a_{210}a_{102}+(B+4D)a_{201}a_{111},\\
 L_zQ={}&3Ba_{300}a_{021}
       +(B+2C)a_{201}a_{120}+(A+4D)a_{210}a_{111}.
\end{aligned}
\]
これらはともに零である。まず $A=B=0$ を得て、続いて
$C=D=0$ を得る。
\end{proof}

\begin{proposition}[三次の補正スカラーの消滅]
\label{prop:nfm-third-tensor-scalar-zero}
命題\ref{prop:nfm-relative-residual-second-jet}の正規化と
定理\ref{thm:nfm-residual-formal-poisson}の仮定を用いる。
$F$ の係数は $h$ に依存せず、二次関係を厳密に
$[z_i,z_j]=h\operatorname{sym}(\Pi_F(z_i,z_j))$ とする。
harmonic-angle formality による生成元固定模型では
\begin{equation}
\label{eq:nfm-third-tensor-no-scalar}
 \pi_h=h\Pi_F+\frac{h^3}{24}\Pi_{\mathscr H_F}+O(h^4),
 \qquad a_h=h+O(h^4).
\end{equation}
特に\eqref{eq:nfm-corrected-poisson-third}の $c_2$ は零である。
この恒等式は普遍係数模型からの任意の基底変更と両立する。
\end{proposition}

\begin{proof}
まず\eqref{eq:nfm-normalized-tensor-recursion}を普遍係数で評価する。
関係の係数が $h$ に依存しない場合、$\pi_m$ はその係数に関する
斉次次数 $m$ の多項式である。
これは帰納法による。星積の $h^m$ 係数の各グラフでは、
tensor 入力の $h$ 次数の和が $m$ である。
右辺 $h\operatorname{sym}_*(\Pi_F(z_i,z_j))$ の寄与では、
tensor 入力の次数の和が $m-1$ であり、さらに関係の係数が一つ入る。
従ってどちらの寄与も係数次数 $m$ を持つ。
固定次数でグラフは有限個であり、再帰に係数による除算はない。
またこの多項式構成は線型座標変更に共変である。

中心変数を含まない普遍三次形式
\[
 g=\sum_{|\alpha|=3}a_\alpha x^\alpha,
 \qquad A=\C[a_\alpha:|\alpha|=3],
 \qquad P=\partial_x\wedge\partial_y\wedge\partial_z
\]
を考え、
\[
 D=\pi_3(g)-\frac1{24}\Pi_{\det\operatorname{Hess}(g)}
\]
と置く。$D$ の係数は $A$ の斉次三次式である。
滑らかな三次形式の稠密開では、既存の Casimir の分解と二次 jet により
$D=c\Pi_g$ である。この主張は、必要なら
$g+\sum_j u_j^3$ という滑らかな分離型模型を用いて適用できる。

$v$ を $\Pi_g$ の二次単項式係数を並べたベクトル、$d$ を
$D$ の同じ係数ベクトルとする。
$d_i v_j-d_j v_i$ は稠密開で零なので $A$ 内で恒等的に零である。
従って $\operatorname{Frac}(A)$ 上で $d=cv$ となる。
ここで $c=p/q$ を互いに素な分子・分母で書くと、$cv_i\in A$ より
$q$ は全ての $v_i$ を割る。
一方、$v$ の非零成分は普遍係数 $a_\alpha$ の非零整数倍であり、
全ての $a_\alpha$ が現れる。従ってそれらの最大公約元は $1$ である。
$A$ は一意分解整域なので $q$ は単元、すなわち $c\in A$ を得る。
この分母の除去は、判別式の補集合上の正則性だけからは従わない点に注意する。

斉次性から $c$ は係数次数二である。
$P$ は $\operatorname{SL}_3$ に不変であり、再帰と Hessian は共変なので、
$d=cv$ におけるスカラーの一意性から $c$ は
$\operatorname{SL}_3$ 不変である。
補題\ref{lem:nfm-no-quadratic-cubic-invariant}により $c=0$ となる。

次に中心変数と任意の混合係数を許す。
補題\ref{lem:nfm-formal-linear-center}により全ての tensor 入力は
active 方向だけを持つ。従ってグラフと再帰の全微分は active 方向であり、
$U=0$ への制限と可換である。
$g=F|_{U=0}$ とすれば
\[
 \pi_3(F)|_{U=0}=\pi_3(g),\qquad
 \mathscr H_F|_{U=0}=\det\operatorname{Hess}(g).
\]
既存の分解
$\pi_3(F)-\Pi_{\mathscr H_F}/24=c_2(F)\Pi_F$ のスカラーは
中心変数を含む全座標に依存しない。
制限を取ると $c_2(F)\Pi_g=0$ となり、active 勾配の
次数零 Koszul 単射性から $c_2(F)=0$ を得る。
以上の普遍恒等式を基底変更すれば、非被約基底についても同じ結論を得る。
最後に $W_h=F+h^2\mathscr H_F/24+O(h^3)$ と
$\pi_h=a_h\iota_{dW_h}P$ を比較すれば $a_h=h+O(h^4)$ となる。
\end{proof}

\begin{remark}[規格化と適用範囲]
この証明が用いる中心元の情報は二次 jet だけである。
高次の中心元を Hesse 座標で記述した際の選択は、
$D=\pi_3-\Pi_{\mathscr H}/24$ の普遍多項式性と共変性に影響しない。
一方、一般の formality への変更、$h$ の再パラメータ化、
または $h$ に依存する二次関係を追加した場合には、
同じ係数式をそのまま主張しない。

偶奇性を用いた剰余の改善は、以下の系\ref{cor:nfm-mixed-even-jet}で証明する。
\end{remark}

\subsubsection{反対積による全次数の偶奇性}

ここでは harmonic-angle の Kontsevich formality を固定する。
一般の angle form や、任意に gauge 変換した formality に同じ偶奇性を
無条件に主張するものではない。
\begin{lemma}[harmonic-angle グラフの反転]
\label{lem:nfm-harmonic-parity}
同じ bivector $\pi$ を $k$ 個入れた星積の項を $B_k(\pi;f,g)$ と書く。
このとき
\[
 B_k(\pi;g,f)=(-1)^kB_k(\pi;f,g),\qquad
 f*_{-\pi}g=g*_{\pi}f.
\]
さらに関数上の接 formality 写像について
\[
 C_{-\pi}(f)=C_{\pi}(f)
\]
が成立する。最後の等式は各奇数 $k$ に対して
$\mathcal U_{k+1}(f,\pi,\ldots,\pi)=0$ であることから従う。
\end{lemma}

\begin{proof}
Kontsevich の配置空間と向きの規約を用いる
\cite[Sections~2.2 and~6.2]{NFMKontsevich}。
上半平面の反射 $r(z)=-\overline z$ は harmonic-angle 一形式を
$r^*d\varphi=-d\varphi$ に変える。
星積では境界の二点を $0,1$ に固定した表示を使い、
$z\mapsto1-\overline z$ により両境界点のラベルを交換する。
内部点は $k$ 個なので配置空間の向きは $(-1)^k$ 倍される。
辺は $2k$ 本であり角形式の積には符号が生じない。
従って両入力を交換したグラフの重みは $(-1)^k$ 倍になり、第一式を得る。
多重線型性から第二式も従う。

中心写像の $k$ 次項では、内部点は bivector の $k$ 頂点と関数の一頂点の
合計 $k+1$ 個、境界点は零個、辺は $2k$ 本である。
配置空間 $C_{k+1,0}$ の次元は $2k$ である。
商を取る前の配置空間の向きは反射により $(-1)^{k+1}$ 倍される一方、
実平行移動と正の拡大の二次元群の向きは $-1$ 倍される。
従って商の向きは $(-1)^k$ 倍される。
ここでは境界ラベルの交換はなく、同じグラフに戻る。
角形式は $2k$ 個なのでその積は不変であり、
重み $w_\Gamma$ は $w_\Gamma=(-1)^k w_\Gamma$ を満たす。
標数零では奇数 $k$ の重みは零である。
これは二次 tensor の線型性・二次性を仮定しないグラフの対称性であり、
偶数次数の中心グラフ全てが消えるという主張ではない。
\end{proof}

\begin{proposition}[生成元固定模型の偶奇性]
\label{prop:nfm-normalized-parity}
補題\ref{lem:nfm-normalized-star-model}の存在仮定の下で、
二次関係の $q_{ij}$ は $h$ に依存しないとする。
すると正規化した tensor は
\[
 \pi_{-h}=-\pi_h
\]
を満たす。Weyl 対称化を $\sigma_h$ とすれば、線型写像
$T_h=\phi_h\sigma_h$ も $T_{-h}=T_h$ を満たす。
特に $\Omega_h=\sigma_h(F_h)$、$F_{-h}=F_h$ という規格化の形式中心元に対し、
\[
 W_h=C_{\pi_h}^{-1}\phi_h(\Omega_h)
 \quad\text{は}\quad W_{-h}=W_h
\]
を満たす。
\end{proposition}

\begin{proof}
生成元を固定して語順を逆にする写像は
$B_{-h}\simeq B_h^{\mathrm{op}}$ を与える。
前補題により $*_{-\pi_h}$ は $*_{\pi_h}$ の反対積である。
従って $B_{-h}$ の生成元固定模型を二通りに書き、
補題\ref{lem:nfm-normalized-star-model}の一意性を適用すれば
$\pi_{-h}=-\pi_h$ となる。

Weyl 対称化された語の総和は語順を逆にしても不変である。
このため同じ反対積の比較で $T_{-h}=T_h$ を得る。
また $C_{\pi_{-h}}=C_{-\pi_h}=C_{\pi_h}$ なので、
$F_h$ が偶数級数なら $C_{\pi_h}^{-1}T_h(F_h)$ も偶数級数になる。
全ての等式は各係数の有限グラフ恒等式であり、
命題\ref{prop:nfm-normalized-model-basechange}で得た普遍族からの
任意の基底変更後も成立する。
\end{proof}

\begin{corollary}[全混合係数における jet の改善]
\label{cor:nfm-mixed-even-jet}
命題\ref{prop:nfm-relative-residual-second-jet}の全混合 Hesse 族で、
同じ規格化中心元 $\widetilde\Omega_{F,h}$ を用いる。
このとき
\[
 W_h=F+\frac{h^2}{24}\mathscr H_F+h^4W_4(F)+O(h^6),
 \qquad
 W_h\equiv H^\sharp\pmod{h^4}.
\]
また形式スカラー $a_h$ も奇数級数である。
命題\ref{prop:nfm-third-tensor-scalar-zero}の $c_2=0$ と合わせると
\[
 \pi_h=h\Pi_F+\frac{h^3}{24}\Pi_{\mathscr H_F}+O(h^5)
       \equiv\Pi^\sharp\pmod{h^5}.
\]
これは全 $n\ge3$、全混合係数についての合同である。
次に決定すべき potential の係数は $W_4$ であり、
$h^3$ の混合補正は存在しない。
\end{corollary}

\begin{proof}
指定中心元は
$\sigma_h(F-h^2\mathscr H_F/12+h^4\Psi_\lambda(F))$
と書かれ、括弧内は $h$ の偶数多項式である。
前命題と既知の二次 jet を合わせる。
代数的補正は厳密に
$H^\sharp=F+h^2\mathscr H_F/24+h^4\tau_F(\Psi_\lambda(F))$
であるため、第一の合同が従う。
$\pi_h=a_h\Pi_{W_h}$ に $h\mapsto-h$ を施すと
$(a_{-h}+a_h)\Pi_{W_h}=0$ となる。
相対勾配 Koszul 複体の完全性からスカラーの作用は単射であり、
$a_{-h}=-a_h$ を得る。
最後は $c_2=0$ と tensor の奇数性による。
ここでは $W_4=\tau_F(\Psi_\lambda(F))$ とも、
有限式 $H^\sharp$ が全次数の中心逆像であるとも結論していない。
\end{proof}

\begin{remark}[中心次数の filtration]
\label{rem:nfm-central-filtration}
相対グラフの全ての微分方向は $x,y,z$ であり、中心変数を微分しない。
従って $I=(U)$ による filtration を保ち、$U=0$ への制限は
生成元固定再帰、Weyl 対称化の像、中心写像およびその逆と可換である。
さらに各グラフの中心次数は、各頂点に選んだ係数の中心次数の和である。
従って tensor の係数を $I^d/I^{d+1}$ で求めるために、
元の関係の中心次数 $d$ を超える係数は不要である。
これは $\pi_h$ の境界部分を純三変数の計算から求める根拠となるが、
その情報のみから正の中心次数を持つ全混合係数は決まらない。
中心三次自身の選択に正の中心次数の係数が含まれる場合には、
もちろんその指定係数も入力に含める。
\end{remark}

\subsubsection{五次の尺度補正と代数的処方の適用範囲}

\begin{lemma}[log-canonical 模型の厳密な逆 tensor]
\label{lem:nfm-log-canonical-inverse}
$g=xyz$ とし、二次関係を
$[x,y]=h(xy+yx)/2$ およびその巡回置換とする。
harmonic-angle formality の生成元固定模型では
\[
 \pi_h(g)=\log\left(\frac{1+h/2}{1-h/2}\right)\Pi_g
       =\left(h+\frac{h^3}{12}+\frac{h^5}{80}+O(h^7)\right)\Pi_g.
\]
\end{lemma}

\begin{proof}
Kontsevich の元の星積による
$\sigma=\sum_{i<j}\lambda_{ij}x_ix_j\partial_i\wedge\partial_j$
の量子化では、原座標に関する交換比が
$q_{ij}(t\sigma)=\exp(t\lambda_{ij})$ となる
\cite[Section~1.6; Theorem~6.9 and Remark~6.10, pp.~59--60]{NFMLindbergPym}。
ここで交換比は、トーラス同変性による
$x_i\star x_j=w_{ij}x_ix_j$ の係数の比 $w_{ij}/w_{ji}$ である。
従って抽象的な量子トーラスの同型による生成元変更は不要である。

$q_h=(1+h/2)/(1-h/2)$、$\ell(h)=\log q_h$ と置き、
この恒等式を $t\mapsto\ell(h)$ で形式的に基底変更する。
$\ell(h)\Pi_{xyz}$ の星積は、原生成元で
$x\star y=q_h y\star x$ とその巡回置換を満たす。
これは指定した $h$ を保った二次関係と同値である。
単元係数の $q$-交換関係には順序単項式の PBW 基底があり、
星積側も同じ単項式の単元倍を基底に持つ。
従って生成元を固定する準同型は、各内部次数で剰余が恒等な
$\C[[h]]$-線型同型である。
補題\ref{lem:nfm-normalized-star-model}の一意性が所望の式を与える。
ここで $\ell(h)$ は標的 tensor の係数であり、
比較の際に指定されたパラメータ $h$ 自体を取り替えたのではない。
\end{proof}

\begin{proposition}[純 active 族の五次 tensor]
\label{prop:nfm-active-fifth-covariant}
普遍三変数三次形式 $g$ に対する生成元固定 tensor を
$\pi_h(g)=\sum_{m\ge1}h^m\pi_m(g)$ とする。
$x_1=x,x_2=y,x_3=z$、$\varepsilon^{123}=1$ と規約を取り、
$t_{ijk}=\partial_i\partial_j\partial_k g/6$ とし、
\[
 S(g)=54\,t_{abc}t_{def}t_{ghi}t_{jkl}
             \varepsilon^{adg}\varepsilon^{bej}
             \varepsilon^{chk}\varepsilon^{fil}
\]
と定める。添字は $1,2,3$ を走り、反復添字について和を取る。
この $S$ は古典的な四次 Aronhold 不変式であり、
四つの volume の縮約による表示も既知である
\cite[Section~3, after Remark~3.3]{NFMOttavianiAronhold}。
ここでは $S(xyz)=1$ と規格化する。
harmonic-angle formality の生成元固定模型では
\[
 \pi_5(g)=\frac1{80}S(g)\Pi_g
\]
となる。
\end{proposition}

\begin{proof}
$\pi_m(g)$ と volume によって双対な一次形式を $\alpha_m(g)$ とする。
滑らかな $g$ に純中心 Fermat 三次を加えた分離型を用いると、
Casimir の分解と $U=0$ への制限により
$\sum_m h^m\alpha_m=a_h\,dW_h$ であり、$a_h$ は座標に依存しない。
従って $d\alpha_m=0$ である。
これは普遍係数に関する多項式恒等式なので、全ての $g$ で成立する。
$E=x\partial_x+y\partial_y+z\partial_z$ に対し
\[
 G_m(g)=\frac13\iota_E\alpha_m(g)
\]
と置くと、Cartan の公式と内部次数により $dG_m=\alpha_m$ となる。
従って $G_m$ は係数次数 $m$、変数次数三の
$\operatorname{SL}_3$ 共変式である。

この型の共変式の次元は
$\operatorname{Sym}^m(\operatorname{Sym}^3\C^3)$ における
最高ウェイト $(m+2,m-1,m-1)$ の重複度に等しい。
重みの個数を
\[
 p_m(a,b,c)=[t^mX^aY^bZ^c]
     \prod_{i+j+k=3}(1-tX^iY^jZ^k)^{-1}
\]
と置く。$\delta=(2,1,0)$ とすれば、Weyl の交代和より
最高ウェイト $\lambda$ の重複度は
\[
 \sum_{\sigma\in S_3}\operatorname{sgn}(\sigma)
             p_m(\lambda+\delta-\sigma\delta)
\]
である。$m=5$、$\lambda=(7,4,4)$ では、この六項は
\[
 41-35-30+23+21-19=1
\]
となる。
一方、上に定義した $S$ は明示的な volume 縮約により不変であり、
$S(xyz)=1$ なので非零である。
従って $g\mapsto S(g)g$ は当該一次元空間を張り、
$G_5(g)=\kappa S(g)g$ を得る。
補題\ref{lem:nfm-log-canonical-inverse}と $S(xyz)=1$ により
$\kappa=1/80$ と決まる。
\end{proof}

\begin{corollary}[混合係数の制限と代数的 tensor 補正の限界]
\label{cor:nfm-fifth-restriction-obstruction}
任意の混合三次 $F$ に対し $g=F|_{U=0}$ と置くと、
\[
 \pi_5(F)|_{U=0}=\frac1{80}S(g)\Pi_g
\]
となる。従って先の Hessian 規格化の族に対する代数的補正
\[
 \pi_h^\sharp=h\Pi_F+\frac{h^3}{24}\Pi_{\mathscr H_F}
\]
は、$S(g)\ne0$ なる稠密開では、固定した $h$ と生成元の下で
実際の形式 tensor $\pi_h$ と $h^5$ の係数から異なる。
この開は全ての $n\ge3$ で非空である。
\end{corollary}

\begin{proof}
線型中心の保存により全グラフの微分は active 方向であるため、
$U=0$ への制限と再帰は可換である。
従って純 active の公式を制限に適用できる。
代数的 tensor 補正の $h^5$ 係数は零であるが、
上の制限は $S(g)\ne0$ で零でない。
五次より低い次数の一致は系\ref{cor:nfm-mixed-even-jet}による。
実際、$g=(x^3+y^3+z^3)/3+xyz$ は滑らかで $S(g)=-7$ であり、
$F=g+\sum_j u_j^3$ は全次元で滑らかな分離型を与える。
普遍係数空間の既約性から非消滅条件は稠密開を定める。
\end{proof}

\begin{remark}[規格化の検算と混合係数]
上の四つの $\varepsilon$ の縮約は $6^4$ 個の有理数項で厳密に展開できる。
得られる $S$ は通常の単項式係数に関する整数係数二十五項であり、
\[
 S\left(\frac{x^3+y^3+z^3}{3}+\lambda xyz\right)
       =\lambda(\lambda^3-8)
\]
を満たす。
共変式の同じ重複度を $m=1,\ldots,8$ で計算すると
$1,0,1,0,1,0,2,0$ となる。
これらは上の有限生成関数による整数計数であり、
Kontsevich グラフ重みや残余同値を計算したものではない。

この結果は純 active 係数についての還元である。
中心変数を含む一般の混合三次について、
$\pi_5(F)=S(F|_{U=0})\Pi_F/80$ と結論することはできない。
active 方向の局所 jet による別の共変式が混合項に寄与する可能性は残る。
また系\ref{cor:nfm-fifth-restriction-obstruction}は固定された tensor の不一致であり、
tensor のスカラーを忘れた比較、残余三次の四次 jet、
または Kuznetsov 成分の非同値までは結論しない。
\end{remark}


\subsubsection{五次の混合項と残余四次項}

次に、五次 tensor の全混合係数を二つの普遍定数へ還元する。
まず共変式の空間を決定し、続いて既知の graph weight を使って二定数を評価する。
active 微分について
\[
 T_{ijk}=F_{ijk},\qquad M_{ij}=F_{ij},\qquad q_i=F_i
\]
と置く。$T$ は座標 $x,y,z,U$ に依存せず、$F|_{U=0}=g$ だけから決まる。
全添字は $1,2,3$ を走り、$\varepsilon^{123}=1$ とする。
次の完全縮約を定義する。
\begin{align*}
 S(T)&=\frac1{24}
 T_{adg}T_{bej}T_{chk}T_{fil}
 \varepsilon^{abc}\varepsilon^{def}\varepsilon^{ghi}\varepsilon^{jkl},\\
 I(F)&=\frac1{24}
 T_{adg}T_{bej}M_{cf}M_{hk}M_{il}
 \varepsilon^{abc}\varepsilon^{def}\varepsilon^{ghi}\varepsilon^{jkl},\\
 J(F)&=\frac1{24}
 T_{adg}T_{bej}T_{chk}M_{fi}q_l
 \varepsilon^{abc}\varepsilon^{def}\varepsilon^{ghi}\varepsilon^{jkl}.
\end{align*}
$S$ は $S(xyz)=1$ の規格化の四次 Aronhold 不変式であり、
$S(g_\lambda)=\lambda(\lambda^3-8)$ である。
$I(F),J(F)$ は全変数に関する三次式である。

\begin{proposition}[五次の混合項の二定数への還元]
\label{prop:nfm-mixed-fifth-reduction}
固定 harmonic-angle formality に依存する二つの普遍定数
$\beta,\gamma\in\mathbf C$ があり、全 $n\ge3$ の全混合係数について
\[
 \pi_5(F)=\Pi_{Q_5(F)},\qquad
 Q_5(F)=\frac1{80}S(g)F
       +\beta\bigl(I(F)-S(g)F\bigr)
       +\gamma\bigl(J(F)-S(g)F\bigr).
\]
ここで $\pi_h=h\Pi_F+h^3\Pi_{\mathscr H_F}/24+h^5\pi_5+O(h^7)$ である。
係数 $1/80$ は別途得た純三変数の五次 tensor の計算を用いる。
この入力を使わない場合も、右辺は三つの普遍定数
$\alpha S(g)F+\beta I(F)+\gamma J(F)$ の勾配へ還元できる。

\end{proposition}

\begin{proof}
生成元固定再帰の $\pi_5$ は、関係 $\Pi_F$ の係数に関する
斉次次数五の微分多項式であり、active 方向だけを微分する。
再帰の Weyl 対称化も局所 jet で書けることを確認しておく。
任意の二次式 $r$ に対して
\[
 \operatorname{sym}_*(r)=r+\frac12\sum_{a,b}
      (\partial_a\partial_b r)(z_a*z_b-z_az_b)
\]
である。$r=\Pi_F(z_i,z_j)$ なら、その active 二階微分は $T$ の成分である。
片方の添字が中心なら括弧内は零となる。
active 添字の場合、正次数の単位正規化グラフは各線型入力をちょうど一回微分するので、
括弧内も下位 tensor の jet だけで書け、裸の座標は残らない。
従って帰納的に再帰の全項が以下の局所 jet 型に入る。
各係数をある点で評価すると、使える jet は $q,M,T$ だけである。
$F$ 自身は二次関係には現れない。
$q,M,T$ の全座標次数はそれぞれ $2,1,0$ であり、
$\pi_5$ の係数は全座標に関する次数二の多項式である。
従って現れ得る型は $qT^4$ と $M^2T^3$ の二つに限られる。
固定体積形式により $\pi_5$ を一形式 $\alpha_5$ と同一視する。
$\operatorname{SL}_3$ 共変性と determinant weight を合わせると、
これら二型の共変式空間の次元はそれぞれ $1,2$ である。

この小さい表現論の計算は、次の有限係数計算で直接確認できる。
$H_{d,k}(X,Y,Z)$ を
\[
 H_{d,k}=[t^k]\prod_{a+b+c=d}(1-tX^aY^bZ^c)^{-1}
\]
と置く。$\lambda$ の重複度は
$\rho=(2,1,0)$ として次の式で与えられる。
\[
 \sum_{w\in\mathfrak S_3}\operatorname{sgn}(w)
 [X^{\lambda_1+\rho_1-(w\rho)_1}
  Y^{\lambda_2+\rho_2-(w\rho)_2}
  Z^{\lambda_3+\rho_3-(w\rho)_3}]H.
\]
二型について $\lambda=(5,4,4)$ を使うと
\[
\begin{array}{c|c|c}
H&\text{六つの符号付き係数}&\text{和}\\\hline
H_{1,1}H_{3,4}&67,-58,-55,41,42,-36&1\\
H_{2,2}H_{3,3}&158,-136,-126,94,92,-80&2
\end{array}
\]
となる。

三つの一形式 $S(g)d_KF,d_KI(F),d_KJ(F)$ はこの三次元空間に属し、
一次独立である。具体的には $T$ を $g_1$ の三階微分とし、
ある点で
\[
 q=(1,0,0),\qquad
 M=\begin{pmatrix}1&1&0\\1&0&0\\0&0&0\end{pmatrix}
\]
とすると、三つの一形式を列に持つ行列は
\[
 \begin{pmatrix}
 -7&-2/3&-10/3\\
 0&-2/3&-1\\
 0&5/3&0
 \end{pmatrix},\qquad \det=-35/3\ne0
\]
となる。ここで $q,M,T$ は独立な Taylor 係数として指定できる。
実際、一つの中心変数 $u_1$ を用いて
$F=T_{ijk}x_ix_jx_k/6+u_1M_{ij}x_ix_j/2+u_1^2q_ix_i+cu_1^3$
とし、$(x,u_1)=(0,1)$ で評価すればよい。
したがってこの独立性は全ての $n\ge3$ に有効である。
したがって三つの勾配は共変式空間全体を張る。
このため $\alpha_5=d_K(\alpha SF+\beta I+\gamma J)$ と書ける。
全 graph の有効添字は三つの active 方向に限られるため、
三定数は中心変数の個数に依存しない。

純三変数の場合、$I(g),J(g)$ はいずれも
$S(g)g$ に等しい。これも小さい共変式計算で確認できる。
$H_{3,5}$ に $\lambda=(7,4,4)$ を適用した六項は
$41,-35,-30,23,21,-19$ で、和は一である。
従って次数五・三次出力の共変式は $S(g)g$ の一次元空間である。
$g=g_1$、$(x,y,z)=(1,0,0)$ で
$S=-7$、$g=1/3$、$I=J=-7/3$ を計算すれば比例定数は一となる。
純 active 計算 $\pi_5(g)=S(g)\Pi_g/80$ から
$\alpha+\beta+\gamma=1/80$ が従い、表示された形を得る。
以上の係数計算と縮約の評価は、保存した検証コードでも有理数として確認できる。
局所 jet の多項式恒等式なので、普遍模型からの任意の基底変更とも両立する。
\end{proof}

\begin{remark}[分離型で消える二方向]
\label{rem:nfm-mixed-fifth-potential}
$I-SF$ と $J-SF$ は $U=0$ に制限すると零となる。
より一般に $F=g(x+a(U))+G(U)$ という中心平行移動された分離型では、
両者は純中心三次 $-S(g)G(U)$ に等しく、その active 勾配は零である。
従ってこの分離型の計算だけでは二つの普遍定数を区別できない。
以下では二つの混合 Taylor jet で定数を決定し、
さらに中心逆像の局所性から四次 potential とその幾何的比較類を求める。
\end{remark}

% !TEX root = ../nc-cubic.tex

\begin{proposition}[混合五次 tensor の普遍定数]
\label{prop:nfm-mixed-fifth-explicit}
命題\ref{prop:nfm-mixed-fifth-reduction}の harmonic-angle 規格化において
\[
 \beta=\frac1{60},\qquad \gamma=\frac1{240}.
\]
従って全ての $n\ge3$ の混合三次 $F$ に対し
\[
 \boxed{\quad\pi_5(F)=\frac1{240}\Pi_{\,4I(F)+J(F)-2S(g)F}\quad}
\]
となる。この式は普遍係数の多項式恒等式であり、既存の普遍 PBW 模型からの
任意の基底変更と両立する。
\end{proposition}

\begin{proof}
先の二定数への還元により、二つの独立な局所 jet で正規化再帰を評価すれば足りる。
一つの中心変数 $u$ を用いた斉次三次
\[
 F_A=xyz+\frac12uz^2,\qquad F_B=xyz+u^2z
\]
を採用する。$u=1$ とし、$p=(x,y,z)=(0,0,1)$ で評価する。
いずれも $g=xyz$、$S(g)=1$ であり、完全縮約の定義から
\[
\begin{array}{c|cc}
 &I(F)-S(g)F&J(F)-S(g)F\\\hline
 F_A&-2uz^2/3&-uz^2/6\\
 F_B&-u^2z&-2u^2z/3
\end{array}
\]
を得る。従って求める二つの成分は
\[
 \pi_5(F_A)^{xy}(p)=\frac1{80}-\frac{4\beta+\gamma}{3},\qquad
 \pi_5(F_B)^{xy}(p)=\frac1{80}-\beta-\frac23\gamma.
\]
ここでは特異な二つの三次を独立 jet の評価に用いている。
一般対象の主張を特異な例から直接推定するのではなく、
既に証明した普遍多項式の二定数表示に代入している。

$P=\Pi_F$、$R=\Pi_{\det\operatorname{Hess}_K(F)}/24$ とする。
$B_k(P^{\otimes k};f,g)$ を、先頭項が $fg+\frac12 P(df,dg)$ となる
Kontsevich 星積の $k$ 頂点部分とし、$B_k[P^{k-1},R]$ は
各頂点への $R$ の挿入の和を表す。
$q^{ij}=P(dx_i,dx_j)$ とすると、生成元固定再帰の五次部分は
\begin{align*}
 \pi_5^{ij}
 ={}&-2\left(B_5[P^5;x_i,x_j]+B_3[P^2,R;x_i,x_j]\right)\\
 &+\frac12\sum_{a,b=1}^3 \partial_a\partial_b q^{ij}
 \left(B_4[P^4;x_a,x_b]+B_2[P,R;x_a,x_b]\right).
\end{align*}
偶奇性により奇数頂点部分は交換子に二倍で現れ、
右辺の最後の項は二次式の Weyl 対称化から生じる。

この有限式の評価には、Banks--Panzer--Pym\cite[Corollary~1.3; Section~6]{NFMBanksPanzerPym}の harmonic weights を
Buring\cite[Sections~3.4--3.6]{NFMBuringGraphs}がまとめた \texttt{star5.txt} と、同じ著者資料の
\texttt{kontsevich\_2\_k.txt},
\texttt{weights\_kontsevich\_2\_k.txt} ($1\le k\le5$) を用いる。
規約は Buring の学位論文 Section~3.4.1--3.4.2, pp.~74--76 にある。
前者の $k$ 次係数は $fg+hP(df,dg)$ の規約であり、
本稿の規約では $2^k$ で割る。
頂点置換と outgoing edge の交換をまとめた multiplicity は
$m(\Gamma)=2^k k!/|\operatorname{Aut}\Gamma|$ なので、
元の weight $w(\Gamma)$ に対する本稿の係数は
$w(\Gamma)/|\operatorname{Aut}\Gamma|$ となる。
この二種類のデータの一致も、有理数演算と有限置換の列挙で確認できる。

線型入力では二つの sink の入次数がともに一でなければ項が消え、
二次 tensor では aerial vertex の入次数が三以上なら項が消える。
これらの条件で残る非零項数は $k=1,2,3,4,5$ の順に
$1,1,1,19,51$ である。
上の式の三項の評価は次の表になる。
最後の列 $W$ は $\frac12\sum_{a,b}\partial_{ab}q^{xy}(B_4+B_2)$ を表す。
\[
\begin{array}{c|ccc|c}
 F&B_5^{xy}&B_3[P^2,R]^{xy}&W&-2(B_5+B_3)+W\\\hline
 F_A&1/480&0&-1/144&-1/90\\
 F_B&1/720&0&-1/240&-1/144
\end{array}
\]
全ての項は保存した73項の有限 graph certificate と\texttt{verify\_mixed\_fifth.py}で有理数として評価される。
この計算は既知の graph weight を入力とし、その積分自体を再計算してはいない。
元の graph--weight 対との符号・自己同型群を含む照合、
公刊の二変数 log-canonical 星積の五次までの係数との照合も行う。
従って
\[
 4\beta+\gamma=\frac{17}{240},\qquad
 \beta+\frac23\gamma=\frac7{360},
\]
であり、これを解くと $\beta=1/60$、$\gamma=1/240$ を得る。
\end{proof}


\subsubsection{中心逆像の四次項と一つの幾何的障害}

以下では指定した Hessian 正規化中心元を
\[
 \widetilde\Omega_h=\operatorname{sym}
 \left(F-\frac{h^2}{12}\mathscr H_F+h^4\Psi_\lambda(F)\right)
\]
とする。$S,I,J$ の規約は命題\ref{prop:nfm-mixed-fifth-reduction}と同じである。
この中心元の規格化は四次のスカラーを決めるために必要であり、
任意の単元倍に対して同じスカラー公式を主張するものではない。

\begin{lemma}[四次中心逆像の局所共変式]
\label{lem:nfm-fourth-center-locality}
固定した harmonic-angle 模型に依存する二定数 $d,e\in\C$ があって、
全混合係数に対して
\[
 W_4(F)=\Psi_\lambda(F)+d I(F)+e J(F)
\]
となる。二定数は $n\ge3$ に依存しない。
\end{lemma}

\begin{proof}
$R_{\pi_h}=\phi\circ\operatorname{sym}$ と置く。
これは可換多項式を同じ生成元の星積で完全対称化する写像である。
Kontsevich グラフと完全対称化は active な affine 座標変更に共変であり、
$R_{\pi_h}(1)=C_{\pi_h}(1)=1$ である。
したがって正次数の係数をある点で評価すると、
入力関数を少なくとも一回微分した局所 jet の多項式で表せる。
具体的には、その点を原点へ移して入力多項式を Taylor 展開すればよい。
定数項への作用は単位の保存で零となり、
残る有限個の単項式の星積は有限個のグラフで計算される。
グラフの微分添字は全て active 方向であるので、
純中心関数への正次数の作用も零である。

$h^4$ の係数には、$R$ と $C^{-1}$ の四次項のほか、
$\pi_3=\Pi_{\mathscr H_F}/24$ を二次項へ入れた項と、
$-h^2\mathscr H_F/12$ への二次作用が含まれる。
いずれも $F$ の係数次数は五、全座標次数は三である。
中心元の $h^4\Psi_\lambda(F)$ 自身への正次数の作用は
$h^4$ の係数には寄与しない。
従って $W_4-\Psi_\lambda(F)$ は、
$q_i=F_i$, $M_{ij}=F_{ij}$, $T_{ijk}=F_{ijk}$
だけの局所多項式であり、微分されない $F$ は現れない。
次数から可能な型は $qMT^3$ と $M^3T^2$ に限られる。

この二型の $\operatorname{SL}_3$ スカラー共変式の空間はいずれも一次元である。
命題\ref{prop:nfm-mixed-fifth-reduction}の Weyl 係数公式を
最高重み $(4,4,4)$ に適用すると、六項はそれぞれ
\[
 \begin{aligned}
 H_{1,1}H_{2,1}H_{3,3}&:
       &&156,-136,-136,100,109,-92,\\
 H_{2,3}H_{3,2}&:
       &&122,-105,-105,78,84,-73
 \end{aligned}
\]
で、和はともに一である。
非零の $J,I$ がそれぞれを張るので表示を得る。
有効な微分添字は三つの active 方向に限られ、二定数は $n$ に依存しない。
\end{proof}

\begin{lemma}[Hessian 写像と Aronhold 不変式による表示]
\label{lem:nfm-classical-mixed-covariants}
$\mathscr H(F)=\det\operatorname{Hess}_K(F)$ とし、
$T=D_K^3F$, $M=D_K^2F$, $q=d_KF$ と置く。
$S(T)$ は上の規格化の Aronhold 四次式である。
このとき
\[
 \begin{aligned}
 I(F)
 &=\frac1{12}\operatorname{tr}
       \bigl(\operatorname{adj}(M)\operatorname{Hess}_K(\mathscr H(F))\bigr)
   =\frac1{12}D\mathscr H_F[\mathscr H(F)],\\
 J(F)&=\frac14(dS)_T[M\odot q],\qquad
 (M\odot q)_{ijk}
       =\frac{M_{ij}q_k+M_{ik}q_j+M_{jk}q_i}{3}.
 \end{aligned}
\]
第二式の微分は対称三階 tensor $T$ の空間上で取り、
$M\odot q$ の対称化には表示した平均を用いる。
従って $I,J$ は Hessian 写像の微分と Aronhold 四次式の微分によって
計算でき、独立した新しい不変式を導入する必要はない。
\end{lemma}

\begin{proof}
第一式の trace は $M^3T^2$ 型、第二式の右辺は $qMT^3$ 型の
スカラー共変式である。
それぞれの共変式空間は前の計数により一次元なので、
$I,J$ との比例定数を一つの非零評価で決められる。
$F=xyz$ と $(x,y,z)=(1,1,1)$ を取る。
このとき $\mathscr H(F)=2F$ であり、
\[
 \operatorname{tr}\bigl(\operatorname{adj}(M)
           \operatorname{Hess}(\mathscr H(F))\bigr)
       =2\operatorname{tr}(\operatorname{adj}(M)M)=12.
\]
一方 $I=1$ である。
また $T_{123}=1$ と添字対称性だけが非零で、
$q=(1,1,1)$、$M_{ii}=0$, $M_{ij}=1$ ($i\ne j$) となる。
$S(T)=1$、$(M\odot q)_{123}=1$ であり、
四次式を微分すると $(dS)_T[M\odot q]=4$、$J=1$ となる。
これで両比例定数が定まる。
Hessian 写像の方向微分の表示は、行列式の微分公式による。
\end{proof}

体積を局所的に固定しない表示では、同じ公式で
$\mathscr H$ を $\mathscr H_P$、$S$ を $S_P$ に置き換え、
方向微分では $P$ を固定する。
これは各式の四つの $P$ の重みと一致する。

\begin{lemma}[Hesse 中心元の四次恒等式]
\label{lem:nfm-fourth-center-symbol}
\[
 F=g_\lambda+Axy+Byz+Czx+Dx+Ey+Hz+G
\]
とし、$A,B,C$ は中心一次式、$D,E,H$ は中心二次式、$G$ は中心三次式とする。
$S=\lambda(\lambda^3-8)$ と書く。
指定中心元の四次 Weyl symbol は恒等的に
\[
 \Psi_\lambda(F)=-\frac{S}{80}F-\frac1{40}I(F)
                         +\frac1{60}J(F)+\mathcal C_4(U),
\]
ただし
\[
 \mathcal C_4(U)=
  \frac{\lambda}{60}(A^3+B^3+C^3)
  +\frac{\lambda^2}{40}ABC
  +\frac{2-\lambda^3}{180}(AH+BD+CE)
  +\frac{S}{80}G
\]
である。
\end{lemma}

\begin{proof}
固定入力論文A\cite{NCCubicA}の
系\texttt{cor:pg-hesse-hessian-form}、式\texttt{eq:pg-hesse-hessian-closed}
にある有限中心公式を $h^4$ について展開する。
$Q=Axy+Byz+Czx$, $Z_2=Az^2+Bx^2+Cy^2$,
$L=Dx+Ey+Hz$, $R_2=B^2x+C^2y+A^2z$,
$R_{11}=ACx+ABy+BCz$ と置けば、その係数は
\[
 \Psi_\lambda(F)=
 -\frac S{48}g_\lambda-\frac{\lambda^4}{48}Q
 +\frac{2-\lambda^3}{36}Z_2-\frac S{144}L
 -\frac{\lambda^2}{24}R_2-\frac{\lambda^3+4}{72}R_{11}.
\]
$T_{111}=T_{222}=T_{333}=2$, $T_{123}=\lambda$
および添字対称性を、$I,J$ の完全縮約へ代入して整理すると上の恒等式を得る。
これは全係数を独立記号とした有限多項式恒等式であり、
保存した検証器は左辺と右辺の全係数差が零であることを確認する。
ここで検算したのは固定入力の中心公式の Weyl symbol であり、
非可換中心性そのものをこの計算で証明したという意味ではない。
\end{proof}



\begin{theorem}[混合残余 potential の普遍四次係数]
\label{thm:nfm-mixed-fourth-potential}
五次 tensor の係数 $\beta=1/60$, $\gamma=1/240$ を用いると、
指定中心元に対して全 $n\ge3$ で
\[
 \boxed{\quad W_4=\Psi_\lambda(F)+\frac1{24}I(F)-\frac1{80}J(F)\quad}
\]
であり、同値に
\[
 W_4=-\frac S{80}F+\frac1{60}I+\frac1{240}J+\mathcal C_4(U).
\]
さらに Casimir 分解のスカラーは
\[
 a_h=h\left(1+\frac{S}{240}h^4\right)+O(h^7)
\]
である。従って先の表示における純中心項 $G_4$ は
$\mathcal C_4(U)$ に等しい。
\end{theorem}

\begin{proof}
局所共変式の補題によって $W_4=\Psi+dI+eJ$ と書く。
まず $U=0$ へ制限する。
$I(g)=J(g)=Sg$, $\Psi(g)=-Sg/48$ および
$\pi_5(g)=S\Pi_g/80$ を用いると
\[
 a_4=S(1/30-d-e)
\]
となる。$a_4$ は座標変数に依存しないので、この式は混合係数にもそのまま有効である。
前補題の恒等式を微分し、
$\pi_5=a_4\Pi_F+\Pi_{W_4}$ と比較する。
$S\,d_KF$, $d_KI$, $d_KJ$ は既に確認した独立な三つの共変式なので
\[
 d-\frac1{40}=\beta=\frac1{60},\qquad
 e+\frac1{60}=\gamma=\frac1{240}
\]
を得る。従って $d=1/24$, $e=-1/80$, $a_4=S/240$ である。
最後の純中心項の同定も、前補題の恒等式から従う。
\end{proof}

\begin{corollary}[四次の幾何的比較類]
\label{cor:nfm-fourth-jacobi-class}
$\tau_F$ を従来の八次 active Jacobi 代数の正規化 trace とする。
全変数の Jacobi 環で
\[
 \boxed{\quad
 [W_4-\tau_F(\Psi_\lambda(F))]
 =\frac1{60}[I(F)]+\frac1{40}[\tau_F(I(F))]
             +\frac S{80}[\tau_F(F)] .\quad}
\]
特に四次の幾何的比較には $\gamma$ と $a_4$ を別途評価する必要はなく、
純中心項の別のグラフ積分も不要である。
\end{corollary}

\begin{proof}
$J$ の縮約は $q_l=F_l$ を一つ含むので、
$J=\sum_l B_l F_l$ と書ける。各 $B_l$ は全座標の一次式である。
従って $[J]=0$ かつ $\tau_F(J)=0$。
また Euler の公式から全 Jacobi 環では $[F]=0$ である。
中心 symbol の恒等式に $\tau_F$ を適用して
\[
 \mathcal C_4-\tau_F(\Psi)
       =\frac S{80}\tau_F(F)+\frac1{40}\tau_F(I)
\]
を得る。これを四次 potential の公式へ代入すればよい。
\end{proof}

\begin{proposition}[Fermat active 切断上の混合公式と一般次元の非消滅]
\label{prop:nfm-fourth-uniform-witness}
$\lambda=0$ の上記全混合族では
\[
 [W_4-\tau_F(\Psi_0(F))]
 =\frac19[ACx+ABy+BCz+AH+BD+CE].
\]
従って各 $n\ge3$ の滑らかな全混合係数空間で、
この四次比較類が非零となる Zariski 開は非空である。
\end{proposition}

\begin{proof}
$R_{11}=ACx+ABy+BCz$, $K=AH+BD+CE$ と書く。
縮約式を $\lambda=0$ で展開すると
\[
 I=-\frac83(Az^2+Bx^2+Cy^2)+\frac43R_{11},\qquad
 J=-\frac23(AF_z+BF_x+CF_y).
\]
$F_x=x^2+Ay+Cz+D$ 等を使えば
\[
 I+\frac83(AF_z+BF_x+CF_y)
        =\frac{20}{3}R_{11}+\frac83K.
\]
八つの平方自由単項式を基底とする active Jacobi 代数では
$\tau_F(x)=\tau_F(y)=\tau_F(z)=0$ なので $\tau_F(I)=8K/3$ である。
前系を適用すると所望の $[R_{11}+K]/9$ を得る。

非消滅には、全次元に共通して
\[
 F_*=\frac13\left(x^3+y^3+z^3+u^3+v^3+\sum_j w_j^3\right)
                    +u xy+v yz
\]
を使える。$n=3$ では $w_j$ はなく、それ以外では $n-3$ 個を加える。
この三次は滑らかである。実際、全偏微分が零なら $w_j=0$。
$y=0$ なら残りも全て零なので、非零解があれば $y=1$ と規格化できる。
そのとき $x^2+u=u^2+x=0$ から $xu\in\{0,1\}$、
$z^2+v=v^2+z=0$ から $zv\in\{0,1\}$ となるが、
$F_y=1+xu+zv$ は零にならない。

三次多項式の空間上の線型汎関数 $\ell$ を
\[
 \ell(vx^2)=\ell(uz^2)=1,\qquad \ell(uvy)=-1
\]
とし、その他の三次単項式では零と定める。
これは全ての $z_i\partial_{z_j}F_*$ と全ての純中心三次を消す。
非自明な確認は
$vF_x=vx^2+uvy$ と $uF_z=uz^2+uvy$ の二つだけである。
従って全 Jacobi 環上の汎関数へ降り、
\[
 \ell\left(W_4-\tau_F(\Psi_0(F_*))\right)
      =\ell(uvy)/9=-1/9\ne0
\]
を得る。実際の比較類は滑らかな係数開上の Jacobi 束の正則切断なので、
非消滅は Zariski 開である。
ここでは純中心三次の像でさらに商を取った束の一定階数性を仮定していない。
\end{proof}

\begin{corollary}[Hessian 補正後にも残る標識付き比較の障害]
\label{cor:nfm-corrected-fourth-obstruction}
各 $n\ge3$ について、前命題の四次類が非零となる開の複素点で、
指定した中心量子化・PBW・formality の入力を満たすものを取る。
残余三次 $V(W_h)$ と有限補正三次 $V(H^\sharp)$ の間には、
同じ $h$ と特殊ファイバーの同定を保つ
$\C[h]/(h^5)$ 上の同型は存在しない。
固定入力の完備 MF 移送と相対 window 比較を適用する場合には、
量子残余圏と $\operatorname{Ku}(V(H^\sharp))$ の間にも、
同じ標識を保つ $\C[[h]]$-線型 dg Morita 同値は存在しない。
後者は $n=4$ を含む全 $n\ge3$ で成立する入力付きの主張である。
\end{corollary}

\begin{proof}
二つの三次は $h^4$ を法として同じである。
小拡大 $\C[h]/(h^5)\to\C[h]/(h^4)$ 上の持上げの差は
$h^4[W_4-\tau_F(\Psi)]$ である。
滑らかな三次の $H^0(T)=0$ により、特殊ファイバー上で恒等な同型は、
逐次の小拡大を通じて $h^4$ を法として恒等となる。
従ってそのような同型が存在すれば四次の差が Jacobi 環で零となり、矛盾する。

圏についても、命題\ref{prop:nf-hh}の $HH^1=0$ と
実際の幾何的 Kodaira--Spencer 写像
$(J_F)_3\hookrightarrow HH^2(\operatorname{Ku}(V(F)))$ の単射性を使う。
共通の $h^4$ 切断の持上げの差は、非零の四次類の像である。
特殊ファイバー上恒等な同値の曖昧さは $HH^1=0$ により逐次消えるため、
この差を同値で消せない。最後に、剰余で恒等かつ Artin 切断と両立する
固定入力の MF・window 同値を用いる。
四次元の幾何的変形の全射性は必要ない。
\end{proof}

この結論は、一般ファイバー上で標識を忘れた三次超曲面または
Kuznetsov 成分の非同値を意味しない。
また、Hessian 補正 $\Phi$ の全 O-降下と次数九の定理はそのまま有効であり、
ここで否定したのは $\Phi$ の三次と残余三次の全次数での標識付き一致である。

\subsubsection{四次まで適合する代数的 Poisson 射}

四次の差も座標不変に打ち消せる。
$P\in\det K$ を使い、$S_P(H),I_P(H),J_P(H)$ を
先の完全縮約の四つの $\varepsilon$ を $P$ に置き換えて定義する。
局所体積枠で $P=sP_0$ なら、それぞれ
$s^4S(H|_K),s^4I(H),s^4J(H)$ である。
$\tau=\tau_{H,K}$ と書き、
\begin{align}
 \mathcal E_P(H)
  &=\frac{I_P(H)}{60}+\frac{J_P(H)}{240}
     -\frac{S_P(H)}{80}\bigl(H-\tau(H)\bigr)
     +\frac{\tau(I_P(H))}{40},\label{eq:nfm-fourth-invariant-correction}\\
 H^{[4]}&=H^\sharp+\mathcal E_P(H),&
 P^{[4]}&=\left(1+\frac{S_P(H)}{240}\right)P
 \label{eq:nfm-fourth-corrected-pair}
\end{align}
と置く。

\begin{theorem}[全 O-同値へ降下する四次補正]
\label{thm:nfm-fourth-algebraic-correction}
$H,H|_K,H^{[4]},H^{[4]}|_K$ が滑らかで
$1+S_P(H)/240$ が可逆な非空開上で、
\[
 \Phi^{[4]}:(H,P)\longmapsto(H^{[4]},P^{[4]})
\]
は分解可能な三次 Poisson 対の代数的射を定める。
この射は任意の基底変更・方程式の生成元の変更・全射影同型と両立し、
既存の良い O-marked moduli の対応する開へ降下する。

指定中心元の trace--Hessian 像
$H=H^{\mathrm{TH}}_{F,h}$ と $P=hP_0$ に沿って、
固定 harmonic-angle 模型では全混合係数・全 $n\ge3$ に対して
\[
 H^{[4]}\equiv W_h\pmod{h^6},\qquad
 \iota_{dH^{[4]}}P^{[4]}\equiv\pi_h\pmod{h^7}
\]
が成立する。この有限式の射と、その全次数での残余比較は別の主張であり、
後者の収束・代数化・同値核の降下はここでは結論しない。
\end{theorem}

\begin{proof}
四つの縮約は $K$ の枠によらず、$\tau$ は任意底変更と両立する。
$(H,P)$ を同じ Poisson 対を表す $(cH,c^{-1}P)$ に替えると、
$S_P(H)$ は不変、$I_P(H),J_P(H)$ は $c$ 倍になる。
従って $\mathcal E_P(H)$ と $H^{[4]}$ は $c$ 倍、
$P^{[4]}$ は $c^{-1}$ 倍となる。全射影同型への共変性も同じ縮約から従う。
また、$H$ に純中心三次 $G(U)$ を加えても、
$S,I,J$ と active 臨界代数は変わらず、
$H-\tau(H)$ も変わらない。従って $\mathcal E_P$ はこの変更に不変である。

$P^{[4]}$ は $P$ の座標に依存しない単元倍なので、
出力は同じ rank 三 support を持つ Jacobian Poisson tensor を定める。
滑らかさと単元条件は開条件であり、$P=sP_0$ の $s=0$ では
$H^{[4]}=H$、単元は一となる。滑らかな三次と active 制限を出発点に取れば、
この形式 trait の一般点で必要な開条件を満たす。
既存の good 一ブロック開との交わりも非空である。
表示上の射を全同値へ降下する議論は
定理\ref{thm:nfm-algebraic-hessian-correction}と同じである。

jet の比較では $H^{\mathrm{TH}}_{F,h}-F$ が純中心三次であることを使う。
そのため
\[
 \mathcal E_{hP_0}(H^{\mathrm{TH}}_{F,h})
 =h^4\left(\frac I{60}+\frac J{240}
          -\frac S{80}(F-\tau_F(F))+\frac{\tau_F(I)}{40}\right)
 =h^4\bigl(W_4-\tau_F(\Psi)\bigr).
\]
$H^\sharp=F+h^2\mathscr H_F/24+h^4\tau_F(\Psi)$ と偶奇性から
最初の合同を得る。
さらに $P^{[4]}=h(1+Sh^4/240)P_0$ と
定理\ref{thm:nfm-mixed-fourth-potential}の $a_h$ の式を
$\pi_h=a_h\iota_{dW_h}P_0$ へ代入すれば二つ目の合同を得る。
\end{proof}

\begin{remark}[有限補正の二つの役割]
従来の $\Phi$ は量子側の楕円点曲線との厳密な一致と一般次数九を持つ。
新しい $\Phi^{[4]}$ についてここで証明したのは、全混合係数でのより高い残余 jet の一致である。
実際、純 active 三次では $\mathcal E_{hP_0}(g)=h^4S(g)g/120$ となるため、
新しい有限式に従来の点曲線との厳密な一致をそのまま移してはいない。
新しい射の一般次数も上の議論では決定していない。
\end{remark}


有限式
\[
 W_h^{[2]}=F+\frac{h^2}{24}\mathscr H_F,\qquad
 \pi_h^{[2]}=h\iota_{dW_h^{[2]}}P
\]
は全ての $h$ に対して Jacobi 恒等式と Casimir 条件を満たす。
しかし、固定 formality の実際の tensor との差は一般に五次から現れる。
四次 potential と新たな有限補正は上で決定したが、
三次の単元倍を許した全次数比較はなお別の課題である。

補正の座標共変性は内在的に書ける。
$K$ を rank 三の active 方向、$P\in\det K$ として
\[
 \mathscr H_{P}(F)
 =\bigl(\det\operatorname{Hess}_{K}F\bigr)(P\otimes P)
\]
と定める。局所的に $P=e_1\wedge e_2\wedge e_3$ と書けば
$\mathscr H_P(F)=\det(e_i e_jF)$ である。
従ってこれは補空間やその基底の選択に依存せず、
全線型座標変換に共変である。また
\[
 \mathscr H_{bP}(aF)=a^3b^2\mathscr H_P(F).
\]
特に $(F,P)\mapsto(aF,a^{-1}P)$ という方程式直線の局所自明化変更では、
$F+h^2\mathscr H_P(F)/24$ は $a$ 倍され、対応する tensor は変わらない。
従って有限 jet の補正は、方程式の代表や volume の選択から独立に定義される。

\begin{remark}[形式的な降下と残る代数化]
固定 formality と $h$ を保つ次数付き中心三次対の線型同型は、
正規化の一意性により形式 Poisson 対の同型を誘導する。
中心元を単元 $b$ 倍すると $W_h$ は $b$ 倍され、$\pi_h$ は変わらない。
従って構成は中心直線についても自然であり、
形式的な表示 groupoid 上では座標変更と cocycle に両立する。
この主張には formal Poisson tensor を単に候補として選ぶ曖昧さはない。

一方、ここで固定した formality 自体を変えた場合の比較、
形式級数の収束・代数化、$h$ の標識を忘れる操作、
および全 intrinsic O-marked 圏族への降下は別問題である。
Kontsevich\cite[Section~2.3, Conjecture~1 and Questions~1]{NFMKontsevichAlg}の
形式対応の収束・周期に関する問題とも区別する。
本節から有限 $h$ の一般点における代数的モジュライ射を宣言しない。
既存の全 O-同値比較を用いる場合も、形式変形とその標識を保つ同値が
表示 groupoid の同値へ移るという適用条件を確認してから降下する。
\end{remark}

\subsection{点曲線が決める split 族の厳密補正}

ここでは量子化の尺度 $s$ と Hesse 係数 $\lambda$ を独立に扱う。
点曲線の Hessian 表示は Mori--Smith\cite[Theorems~1.2, 7.2]{MoriSmith}
に既にある結果であり、新規性を主張する対象ではない。
本小節の用途は、その Hessian 表示を残余復元と結び付け、
補正後の Poisson 対の尺度と定義体を明記することである。

$\lambda\in\C$ を $1+\lambda^3\ne0$ なる定数とし、
$R=\C[s]_{(s)}$、$\widehat R=\C[[s]]$、$K=\C(s)$ と置く。
$S=x^3+y^3+z^3$、$P=xyz$、$g_\lambda=S/3+\lambda P$ と書く。
$D_R$ は
\[
 \begin{aligned}
 [y,z]&=s\left(x^2+\frac\lambda2(yz+zy)\right),\\
 [z,x]&=s\left(y^2+\frac\lambda2(zx+xz)\right),\\
 [x,y]&=s\left(z^2+\frac\lambda2(xy+yx)\right).
 \end{aligned}
\]
で定まる代数とする。$r=n-1\ge2$ とし、中心変数を加えて
$B_R=D_R[u_1,\ldots,u_r]$ と置く。
$G\in\C[u_1,\ldots,u_r]_3$ の偏微分の共通零点は原点だけと仮定する。
特に $G=\sum u_i^3/3$ を全ての $n$ に一様に使える。
\[
 \begin{aligned}
 a_{\lambda,s}&=1+\frac{s^2\lambda^2}{2}
                       +\frac{s^4\lambda}{6}-\frac{s^4\lambda^4}{48},\\
 b_{\lambda,s}&=\lambda-\frac{s^2\lambda^3}{6}-\frac{2s^2}{3}
                       +\frac{s^4\lambda^2}{6}-\frac{s^4\lambda^5}{48},\\
 \Omega_D&=a_{\lambda,s}S/3+b_{\lambda,s}\operatorname{sym}(P),
 \qquad w=\Omega_D+G.
 \end{aligned}
\]
これらは既存の Hesse 中心公式の二パラメータ表示である。
$r$ 個の中心変数を付加しても、各次数の有限自由性と中心性は保存される。

\begin{proposition}[点曲線の厳密な Hessian 補正]
\label{prop:nfm-hesse-point-correction}
$\mathscr H_\lambda=\det\operatorname{Hess}_{x,y,z}(g_\lambda)$ と置き、
\[
 \begin{aligned}
 A_{\lambda,s}&=1-\frac{s^2\lambda^2}{4},\\
 B_{\lambda,s}&=\lambda+\frac{s^2(4+\lambda^3)}{12},\\
 J_{\lambda,s}&=g_\lambda+\frac{s^2}{24}\mathscr H_\lambda
       =A_{\lambda,s}S/3+B_{\lambda,s}P,\\
 \nu(\lambda,s)&=\frac{B_{\lambda,s}}{A_{\lambda,s}}
       =\frac{\lambda+s^2(4+\lambda^3)/12}{1-s^2\lambda^2/4}
 \end{aligned}
\]
とする。一般点の active な点曲線は $V(J_{\lambda,s})$ である。
この曲線の $R$ 上の滑らかな閉包は $V(S+3\nu(\lambda,s)P)$ であり、
$s=0$ のファイバーは $V(g_\lambda)$ である。
ただし $s=0$ の可換代数自体の点スキームは $\mathbb P^2$ であり、
この閉包と混同しない。

Trace--Hessian 出力は $H^{\mathrm{TH}}=g_\lambda+G$ であり、
補正後の対
\[
 H^{\mathrm{corr}}=J_{\lambda,s}+G,
 \qquad
 \Pi^{\mathrm{corr}}
 =s\,\iota_{dH^{\mathrm{corr}}}
       (\partial_x\wedge\partial_y\wedge\partial_z)
\]
は $K$ 上の滑らかな三次超曲面と非零 Poisson 構造を定める。
その active な平面切断は、もとの代数の点曲線と埋込みも含めて一致する。
\end{proposition}

\begin{proof}
Sklyanin 係数は
$a=1-s\lambda/2$、$b=-1-s\lambda/2$、$c=-s$ である。
De Laet--Le Bruyn\cite[Section~2, preprint p.~3]{NFMDLL}の中心元
（本稿の基礎節の式\texttt{gre:eq:known-center}も参照）を
$C=xyz+yzx+zxy$、$D=yxz+xzy+zyx$ により
\[
 a(b^3-c^3)C+b(c^3-a^3)D+c(a^3-b^3)S
\]
と書き、関係式 $C-D=sS+(s\lambda/2)(C+D)$ を代入すると
$-9s\Omega_D$ を得る。
一般点は滑らかな楕円型の Sklyanin 代数であり、
剰余は三変数多項式環である。
両者の Hilbert 関数が一致するので、DVR 上の構造定理により全次数が自由である。
一般点の中心性は次数四の torsionfree 性により $R$ 上へ戻る。
この次数別自由性と多項式環への特殊化により、固定入力の
$\texttt{gcd:prop:automatic-as}$ の仮定も満たされる。
従って一般点の $B$ は両側 noetherian な Koszul AS 正則代数であり、
大域次元と Gorenstein parameter は $N=n+2$ である。
また $w\bmod s=g_\lambda+G\ne0$ なので $w$ は正則で、中心商も次数別に自由である。

点曲線の行列式は
\[
 -\frac{a^3+b^3+c^3}{3abc}
 =\frac{\lambda+s^2(4+\lambda^3)/12}{1-s^2\lambda^2/4}
\]
を与える。直接計算でも
\[
 \mathscr H_\lambda=(8+2\lambda^3)P-2\lambda^2S
\]
であり、表示した $J_{\lambda,s}$ を得る。
$A_{\lambda,s}$ と $1+\nu(\lambda,s)^3$ は $R$ の単元なので曲線は滑らかである。
点自己同型を与える切断は $(a:b:c)$ であり、$s=0$ では $(1:-1:0)$ に特殊化する。

active な Jacobi 環で $x^3=y^3=z^3=-\lambda P$ であるから
\[
 b_{\lambda,s}-\lambda a_{\lambda,s}
 =-\frac23s^2(1+\lambda^3),\qquad
 \nu_{s\Pi_{g_\lambda}}(\overline g_\lambda)
 =8s^2(1+\lambda^3)[P].
\]
正次数の active 元の乗法 trace は零であり、純中心成分は $G$ となる。
従って $-12$ 規格化の Trace--Hessian 出力は $g_\lambda+G$ である。
$H^{\mathrm{corr}}\bmod s=g_\lambda+G$ は滑らかなので、
$R$ 上の族と一般ファイバーは滑らかである。
Jacobian 形の三変数 Poisson tensor は全ての $u_i$ を Casimir とし、
$H^{\mathrm{corr}}$ も Casimir であるから、表示した対を得る。
\end{proof}

\begin{theorem}[split 族での残余圏と補正後三次の比較]
\label{thm:nfm-split-corrected-residual}
上の族について、固定入力論文Bの完備 MF 移送、相対 window、
固定対象の局所化および有限降下、coherent 残余圏の飽和性の入力を用いる。
特に最後の入力は $\texttt{crs:prop:residual-smooth}$ と
$\texttt{crs:thm:coherent}$ である。
このとき、完備残余復元の Hesse 係数は全次数で $\nu(\lambda,s)$ に等しい。
さらに、ある有限代数拡大 $K'/K$ 上で
\[
 (\mathcal K_{B/(w)},\rho)
 \simeq
 (\operatorname{Ku}(V(H^{\mathrm{corr}})),\rho_{\mathrm{corr}})
\]
という dg Morita 同値が存在する。
同値は $w$ と表示した $H^{\mathrm{corr}}$ に由来する指定三次周期を保存する。
これは任意の $n\ge3$ と任意の定数 $\lambda\in\C$、$1+\lambda^3\ne0$ に成立する。
\end{theorem}

\begin{proof}
固定入力の一パラメータ式へ単に記号を代入するのではなく、
その構成を本族について確認する。
まず三変数の $D_{\widehat R}$ と $\Omega_D$ に完備移送を適用する。
$g_\lambda$ は滑らかで、上で次数別自由性と中心性を確認したので、その仮定を満たす。
巡回置換と $(x,y,z)\mapsto(x,\omega y,\omega^2z)$ が生成する有限 Heisenberg 群は、
全ての $s$ に対して関係式と $\Omega_D$ を保存する。
ここでも harmonic-angle formality と、次数一生成元を固定する模型を使う。
補題\ref{lem:nfm-normalized-star-model}の一意性と線型共変性により、
この模型は自動的に Heisenberg 群と同変である。
従って固定入力の $\texttt{gre:prop:hesse-normalform}$ と同じ不変三次空間に入り、
active potential は
\[
 W_{\mathrm{act}}=a(s)S/3+b(s)P,
 \qquad a(0)=1,\quad b(0)=\lambda
\]
となる。中心多項式変数への自明な延長により、全 potential は
$W=W_{\mathrm{act}}+G$ である。
active 座標だけを $a(s)^{-1/3}$ 倍して得る Hesse 係数を
$\mu(s)=b(s)/a(s)$ と書く。
普遍二次 jet 計算により
\[
 a(s)=1-\frac{s^2\lambda^2}{4}+O(s^3),\qquad
 b(s)=\lambda+\frac{s^2(4+\lambda^3)}{12}+O(s^3),
\]
従って
\[
 \mu(s)=\lambda+\frac{1+\lambda^3}{3}s^2+O(s^3).
\]

$D_R/(\Omega_D)$ の相対 THCR 表示も同じ証明で得られる。
すなわち、滑らかな点曲線と平行移動切断を $R$ 上で取り、
その $\mathcal O(1)$ の twisted tensor 冪の切断環への写像を次数一から構成する。
各正次数の切断加群は有限自由であり、一般点で二次関係と $\Omega_D$ が消えるので、
torsionfree 性によりこの写像は $R$ 上で $D_R/(\Omega_D)$ を経由する。
また $\Omega_D\bmod s=g_\lambda$ は非零因子であるから、
次数別 Nakayama と分裂単射によって商の各次数も有限自由となる。点曲線と平行移動切断の剰余は
$(V(g_\lambda),(1:-1:0))$ であり、剰余の商写像は通常の斉次座標環の同型である。
各次数で同じ階数の有限自由加群の間の写像なので Nakayama により同型となる。
相対 ampleness と三変数の指数零の MF 比較は固定入力の
$\texttt{gre:prop:elliptic-quotient}$、$\texttt{gre:lem:elliptic-mf}$ の証明通りである。
ここでは Fermat という条件を使わず、$V(g_\lambda)$ の滑らかさだけを使う。
従って $\C((s))$ 上で $E_\mu$ と $E_\nu$ の導来同値を得る。
両者は $(1:-1:0)$ を原点に持つ楕円曲線なので、
Antieau--Krashen--Ward\cite[Theorem~2.2]{NFMAKW}により同型となり、
$j(\mu)=j(\nu)$ である。

この等式から分枝を一意に決める。
$j(z)$ は $z=\lambda$ で正則な非定数有理関数なので、ある整数 $e\ge1$ と
形式的可逆変数 $\psi(z)=z+O(z^2)$、非零定数 $c$ により
\[
 j(\lambda+z)-j(\lambda)=c\,\psi(z)^e
\]
と書ける。$\mu-\lambda$ と $\nu-\lambda$ はともに
$(1+\lambda^3)s^2/3+O(s^3)$ で始まり、その係数は非零である。
$j$ の等式から二つの $\psi$ の値の比は $e$ 乗根となり、
その剰余は一だから比は一である。
$\psi$ の形式逆により $\mu=\nu$ を得る。
この証明は $j'(\lambda)=0$ の場合も含む。

従って全 potential は $a(s)g_\nu+G$ である。
active 座標だけを $(A_{\lambda,s}/a(s))^{1/3}$ 倍すれば、
厳密に $H^{\mathrm{corr}}=A_{\lambda,s}g_\nu+G$ へ移る。
この根は剰余が一であり、底の $s$、中心変数、周期変数 $t$ を固定する。
従って内部 shift と指定三次周期を保つ完備 MF 比較が得られる。

$H^{\mathrm{corr}}$ 自体は $K$ 上の多項式である。
固定入力の一般次元の固定対象への制限、局所化、生成性および有限降下を
この既に指定した代数的標的へ適用すると、ある有限 $K'/K$ 上で
まず線束生成残余部分について同値を得る。
ここで用いる局所化と有限降下は $\texttt{rog:lem:localization}$ と
$\texttt{rog:lem:descent}$ である。
全 coherent 残余圏への移行には、さらに上記の残余 smoothness と
coherent 飽和性の入力を用いる。これにより表示した $\mathcal K$ 全体の同値となる。
完備化上の形式座標変換そのものが $K'$ 上に代数化されることは使わない。
また、特殊点を含む $R$ 上の代数的相対圏同値はここでは結論しない。
\end{proof}

\begin{remark}[尺度と既存の一パラメータ例]
$A=A_{\lambda,s}$ と書く。
$A$ の三乗根を加え、active 座標に $x,y,z\mapsto A^{-1/3}(x,y,z)$ を施すと、
補正後の対は
\[
 \left(g_\nu+G,\ sA\,\Pi_{g_\nu}\right)
\]
へ移る。三次式を規格化するだけで tensor を $s\Pi_{g_\nu}$ と書くのは正しくない。
三重ベクトルは逆 determinant によって $A$ 倍されるためである。
補正前の表示 $(H^{\mathrm{corr}},\Pi^{\mathrm{corr}})$ は既に $K$ 上にあり、
三乗根は規格化した表示との同型のためにだけ必要となる。

前節の $\lambda=s=h$ の族では
$A=1-h^4/4$、$\nu=\lambda_{\mathrm{res}}(h)$ となる。
従って補正後の代数的三次は、前節で独立に決定した残余三次と
有限三乗根拡大後に厳密な座標変換で一致する。
この場合の残余比較は、前節の固定入力の exact formula からも直接従う。

上の $\Pi^{\mathrm{corr}}$ は、点曲線と Hessian 共変式によって選んだ
代数的 Poisson tensor である。
固定した Kontsevich 比較の formal Poisson tensor の全次数と同一だとは証明していない。
後者の尺度や gauge の全高次比較、一般混合係数での残余比較、
全 O-同値に対する残余同値の降下は別の課題として残る。
\end{remark}


\begin{thebibliography}{99}

\bibitem{NCCubicA}
\emph{滑らかな三次超曲面のPoisson構造：正規形・退化軌跡・代数的量子化},
\texttt{nc-cubic} 論文A、未公表研究稿 A-v61、2026年11月14日。
\href{https://github.com/uedakazushi/nc-cubic/blob/eabe601c1341777742964584dde5314c0bcf64a9/papers/paper_a.tex}{固定 commit の主原稿}。
支持節は \texttt{papers/support/papers/sections/poisson\_geometry.tex}
の命題 \texttt{prop:pg-represent}, \texttt{prop:pg-jacobi-extension}、
定理 \texttt{thm:pg-normal}、系 \texttt{cor:pg-generic}、および
\texttt{general\_dimension\_universal\_quantization.tex}
の定理 \texttt{thm:pg-all-dimensions}、命題 \texttt{prop:pg-ad-relative-chart}。

\bibitem{NFNCCubicB}
\emph{中心三次非可換変形の一般点における可換復元},
\texttt{nc-cubic} 論文B、未公表研究稿 B-v59、2026年11月14日。
\href{https://github.com/uedakazushi/nc-cubic/blob/eabe601c1341777742964584dde5314c0bcf64a9/papers/paper_b.tex}{固定 commit の主原稿}。
\texttt{papers/generic\_recovery/} の
\texttt{main\_result.tex}, \texttt{generic\_descent.tex},
\texttt{coherent\_saturation.tex}, \texttt{field\_model.tex} を参照。

\bibitem{Threefold}
\emph{非可換三次三様体のAbel圏モジュライ：二次包絡とKuznetsov残余射},
未公表研究稿 v22、2026年11月20日。
\href{https://github.com/uedakazushi/nc-cubic-threefold-moduli/blob/6c35d2c84671c2538e61836a6bf3a511392eacce/papers/nc-cubic-threefold-moduli.tex}{固定 commit の主原稿}。
点スキームは \texttt{papers/sections/point\_scheme\_poisson.tex}、
全 Hom の比較は \texttt{jacobian\_abel\_comparison.tex} を参照。

\bibitem{NFMEG}
P.~Etingof and V.~Ginzburg,
\emph{Noncommutative del Pezzo surfaces and Calabi--Yau algebras},
J. Eur. Math. Soc. \textbf{12} (2010), 1371--1416.
\url{https://doi.org/10.4171/JEMS/235}.
\url{https://ems.press/content/serial-article-files/31746},
Theorem~3.3.2, Section~9.

\bibitem{MoriSmith}
I.~Mori and S.~P.~Smith,
\emph{The classification of 3-Calabi--Yau algebras with 3 generators and 3 quadratic relations},
Math. Z. \textbf{287} (2017), 215--241.
\url{https://doi.org/10.1007/s00209-016-1824-5}.
\url{https://arxiv.org/abs/1502.07403v2}, Theorem~4.2.

\bibitem{LBSV}
L.~Le Bruyn, S.~P.~Smith and M.~Van den Bergh,
\emph{Central extensions of three dimensional Artin--Schelter regular algebras},
Math. Z. \textbf{222} (1996), 171--212.
\url{https://doi.org/10.1007/PL00004532}.
\url{https://sites.math.washington.edu/~smith/Research/Central-extns.pdf},
Theorem~4.2.2, Proposition~4.2.5, Section~4.3.

\bibitem{Mayanskiy}
E.~Mayanskiy,
\emph{Poisson cohomology of two Fano threefolds},
J. Algebra \textbf{424} (2015), 21--45.
\url{https://doi.org/10.1016/j.jalgebra.2014.08.049}.
\url{https://arxiv.org/abs/1303.5916v1}, Theorem~1, Section~2.2.
Theorem~1 は Loray--Pereira--Touzet の分類に帰属されている。
スキーム構造には Section~2.2 の Schouten 計算を用いる。

\bibitem{BCMStandardMonomial}
N.~C.~Bonala, O.~Clarke and F.~Mohammadi,
\emph{Standard monomial theory and toric degenerations of Richardson varieties in the Grassmannian},
\url{https://arxiv.org/abs/2103.16208v1}, 2021,
Section~2.4, Theorem~2.3.

\bibitem{NFMStacksProjective}
The Stacks Project Authors,
\emph{The Stacks Project}, Section~27.21, Projective bundles, Tag~01OA.
\url{https://stacks.math.columbia.edu/tag/01OA}.

\bibitem{NFMStacksAffineDescent}
The Stacks Project Authors,
\emph{The Stacks Project}, Lemma~35.37.1, Tag~0245.
\url{https://stacks.math.columbia.edu/tag/0245}.

\bibitem{NFLZSerre}
X.~Lin and S.~Zhang,
\emph{Serre algebra, matrix factorization and categorical Torelli theorem for hypersurfaces},
Math. Ann. \textbf{391} (2025), 163--177.
\url{https://doi.org/10.1007/s00208-024-02915-8}.
定理番号は \url{https://arxiv.org/abs/2310.09927v2}
の Theorem~3.5, Proposition~3.7, Theorems~4.1 and~5.1 による。

\bibitem{NFBFKKernels}
M.~Ballard, D.~Favero and L.~Katzarkov,
\emph{A category of kernels for equivariant factorizations and its implications for Hodge theory},
Publ. Math. Inst. Hautes \'Etudes Sci. \textbf{120} (2014), 1--111.
\url{https://doi.org/10.1007/s10240-013-0059-9}.
\url{https://numdam.org/item/10.1007/s10240-013-0059-9.pdf},
Theorem~5.39 の一般式とその証明、pp.~82--84。

\bibitem{NFPirozhkov}
D.~Pirozhkov,
\emph{A categorical Torelli theorem for hypersurfaces},
Bull. London Math. Soc. \textbf{56} (2024), no.~10, 3075--3089.
\url{https://doi.org/10.1112/blms.13117}.
本文の Lemmas~5.1 and~5.2、Theorem~5.3 は
\url{https://arxiv.org/abs/2208.13604v1} による。

\bibitem{NFHuybrechtsRennemo}
D.~Huybrechts and J.~V.~Rennemo,
\emph{Hochschild cohomology versus the Jacobian ring and the Torelli theorem for cubic fourfolds},
Algebraic Geometry \textbf{6} (2019), no.~1, 76--99.
\url{https://doi.org/10.14231/AG-2019-005}.
\url{https://content.algebraicgeometry.nl/2019-1/2019-1-005.pdf},
Proposition~3.5 and Remark~3.8.

\bibitem{NFDAGX}
J.~Lurie,
\emph{Derived Algebraic Geometry X: Formal Moduli Problems}, November 5, 2011.
\url{https://www.math.ias.edu/~lurie/papers/DAG-X.pdf},
Section~2, Theorems~5.1.9 and~5.3.16, Proposition~5.3.21.

\bibitem{NFBKP}
A.~Blanc, L.~Katzarkov and P.~Pandit,
\emph{Generators in formal deformations of categories},
Compos. Math. \textbf{154} (2018), 2055--2089.
\url{https://doi.org/10.1112/S0010437X18007303}.
\url{https://arxiv.org/abs/1705.00655v1}, Sections~3.3 and~4.3.

\bibitem{NFKuznetsovFCY}
A.~Kuznetsov,
\emph{Calabi--Yau and fractional Calabi--Yau categories},
J. reine angew. Math. \textbf{753} (2019), 239--267.
\url{https://doi.org/10.1515/crelle-2017-0004}.
\url{https://arxiv.org/abs/1509.07657v2},
Section~3.3, Corollaries~3.18 and~4.1, Lemma~3.19.

\bibitem{NFBFKOrlov}
M.~Ballard, D.~Favero and L.~Katzarkov,
\emph{Orlov spectra: bounds and gaps},
Invent. Math. \textbf{189} (2012), 359--430.
\url{https://doi.org/10.1007/s00222-011-0367-y}.
定理番号は \url{https://arxiv.org/abs/1012.0864v3} の Theorem~5.15 による。

\bibitem{NFKuznetsovCubicFourfolds}
A.~Kuznetsov,
\emph{Derived categories of cubic fourfolds},
in \emph{Cohomological and Geometric Approaches to Rationality Problems},
Progress in Mathematics \textbf{282}, Birkh\"auser, 2010, 219--243.
\url{https://arxiv.org/abs/0808.3351}, Theorem~3.1.
著者稿 \url{https://cims.nyu.edu/~tschinke/books/rationality/kuznetsov.pdf}.

\bibitem{NFMGabrielQuotient}
P.~Gabriel, \emph{Des cat\'egories ab\'eliennes},
Bull. Soc. Math. France \textbf{90} (1962), 323--448,
Chapitre~III, \S1, p.~365.
\url{https://doi.org/10.24033/bsmf.1583}.
\url{https://www.numdam.org/article/BSMF_1962__90__323_0.pdf}.

\bibitem{Fraenkel}
A.~McMillan Fraenkel,
\emph{Extensions of Poisson structures on singular hypersurfaces},
\url{https://arxiv.org/abs/1310.6083v2}, 2013,
Proposition~4.1, Theorem~4.3.

\bibitem{Carles}
R.~Carles,
\emph{D\'eformations dans les sch\'emas d\'efinis par les identit\'es de Jacobi},
Ann. Math. Blaise Pascal \textbf{3} (1996), no.~2, 33--62.
\url{https://doi.org/10.5802/ambp.67}.
\url{https://www.numdam.org/item/AMBP_1996__3_2_33_0.pdf},
Section~8, Lemma~8.35.

\bibitem{CarlesMarquez}
R.~Carles and M.~C.~M\'arquez,
\emph{Different methods for the study of obstructions in the schemes of Jacobi},
Ann. Inst. Fourier \textbf{61} (2011), no.~2, 453--490.
\url{https://doi.org/10.5802/aif.2620}.
\url{https://www.numdam.org/item/AIF_2011__61_2_453_0.pdf},
Section~3, Theorem~3.2, Propositions~3.9 and~3.10.

\end{thebibliography}

\end{document}
